\documentclass[11pt]{article}
\usepackage[a4paper,left=30mm,right=34mm,top=22mm,bottom=22mm]{geometry}
\usepackage{amssymb}
\usepackage{fontspec,needspace,xcolor}
\usepackage[unicode,hidelinks]{hyperref}
\usepackage{polyglossia}
\setmainlanguage{english}
\setotherlanguages{latin,german}
\setmainfont{Linux Libertine O}[Ligatures=TeX]
\newfontfamily\SansFont{Linux Biolinum O}
\newcommand{\SansBold}[1]{{\SansFont\bfseries #1}}
\newcommand{\Sperr}[1]{{\addfontfeature{LetterSpace=15}#1}}
% a bar over letters or figures ({ov:...} in the page files): the thousands of a Roman numeral, the barred zero
\newcommand{\Ovl}[1]{\ensuremath{\overline{\mbox{#1}}}}
% a value with a small label printed above it ({st:LABEL|VALUE} in the page files)
\newcommand{\Stk}[2]{\begin{tabular}[b]{@{}c@{}}{\tiny #1}\\[-3pt]#2\end{tabular}}
% a fraction printed built up ({fr:NUM|DEN} in the page files)
\newcommand{\Frac}[2]{\ensuremath{\displaystyle\frac{\mbox{#1}}{\mbox{#2}}}}
% a case fraction, small figures over a bar within the line ({cf:NUM|DEN} in the page files)
\newcommand{\CaseFrac}[2]{\ensuremath{\textstyle\frac{\mbox{\scriptsize #1}}{\mbox{\scriptsize #2}}}}
% a square root, the radical sign with its bar over the radicand ({rt:...} in the page files)
\newcommand{\Root}[1]{\ensuremath{\sqrt{\mbox{#1}}}}
% two words in small type one over the other within the line ({sk:TOP|BOTTOM} in the page files)
\newcommand{\SmallStack}[2]{\begin{tabular}[c]{@{}c@{}}\footnotesize #1\\[-4pt]\footnotesize #2\end{tabular}}
\hypersetup{pdftitle={Die astronomischen Tafeln des al-Khwarizmi, ed. Suter 1914: the editor's matter},pdfauthor={ed. H. Suter (Copenhagen 1914)}}
\setlength{\parindent}{0pt}\setlength{\parskip}{0pt}
\tracinglostchars=2
\newcommand{\FrameHead}[4]{\par\bigskip\Needspace{8\baselineskip}\penalty-300\pdfbookmark[1]{JP2 #1 (#2)}{fr-#1}\hypertarget{AK07-J#3}{}%
 {\large\bfseries #4}\enspace{\footnotesize\color{gray}JP2 #1}\par\smallskip}
\newcommand{\LnStrut}{\vrule height 0.95em depth 0.35em width 0pt}
% a note in the right or the left margin: the folios of the manuscripts
\newcommand{\Marg}[1]{\rlap{\hspace{\linewidth}\hspace{3mm}\footnotesize\color{gray}#1}}
\newcommand{\MargL}[1]{\llap{\footnotesize\color{gray}#1\hspace{12mm}}}
% one printed line, with its number in the left margin
\newcommand{\EdL}[4]{\par\noindent\hypertarget{#2}{}\llap{\footnotesize\color{gray}#1\hspace{4mm}}#4\parbox[t]{\linewidth}{\raggedright\LnStrut #3\par}\par}
\newcommand{\EdLp}[4]{\EdL{#1}{#2}{\hspace*{6mm}#3}{#4}}
\newcommand{\EdLr}[4]{\par\noindent\hypertarget{#2}{}\llap{\footnotesize\color{gray}#1\hspace{4mm}}#4\parbox[t]{\linewidth}{\raggedleft\LnStrut #3\par}\par}
\newcommand{\EdLc}[4]{\par\noindent\hypertarget{#2}{}\llap{\footnotesize\color{gray}#1\hspace{4mm}}#4\parbox[t]{\linewidth}{\centering\LnStrut #3\par}\par}
\newcommand{\EdLh}[4]{\par\noindent\hypertarget{#2}{}\llap{\footnotesize\color{gray}#1\hspace{4mm}}#4\parbox[t]{\linewidth}{\centering\large\LnStrut #3\par}\par\smallskip}
% a line of the contents: the title with its leader, the pages, and (in the chapters of the commentary) the tables
\newcommand{\EdLt}[6]{\par\noindent\hypertarget{#2}{}\llap{\footnotesize\color{gray}#1\hspace{4mm}}#6\parbox[t]{\dimexpr\linewidth-30mm}{\LnStrut #3\dotfill\par}\makebox[12mm][r]{#4}\makebox[18mm][c]{#5}\par}
\newcommand{\EdLtp}[5]{\par\noindent\hypertarget{#2}{}\llap{\footnotesize\color{gray}#1\hspace{4mm}}#5\parbox[t]{\dimexpr\linewidth-30mm}{\LnStrut #3\dotfill\par}\makebox[30mm][r]{#4}\par}
\newcommand{\ColHead}[1]{\par\smallskip\noindent{\footnotesize\color{gray}#1}\par}
% a printed line of the footnotes
\newcommand{\EdN}[3]{\par\noindent\hypertarget{#2}{}\llap{\footnotesize\color{gray}#1\hspace{4mm}}\parbox[t]{\linewidth}{\raggedright\footnotesize\LnStrut #3\par}\par}
\newcommand{\NoteRule}{\par\smallskip\noindent\rule{30mm}{0.4pt}\par}
\newcommand{\Withheld}[1]{\par\smallskip{\small\color{gray}[#1]}\par}
\begin{document}

\thispagestyle{empty}
\begin{center}\LARGE Die astronomischen Tafeln des Muḥammed ibn Mūsā al-Khwārizmī\par\vspace{3mm}\large in der Bearbeitung des Maslama ibn Aḥmed al-Madjrīṭī und der lateinischen Uebersetzung des Athelhard von Bath\par\vspace{4mm}\normalsize edited by H. Suter on the preliminary work of A. Bjørnbo and R. Besthorn\\(Copenhagen: A. F. Høst \& Søn, 1914)\par\vspace{3mm}the Wellcome Library copy, Internet Archive b24872489, JP2 images 0000–0295\par\vspace{4mm}\large the editor’s German matter and the pages without text, line by line\end{center}\vspace{6mm}
\noindent The Arabic \textit{Zīj al-Sindhind} of Muḥammad ibn Mūsā al-Khwārizmī is lost. Suter’s book edits a derivative line of it: the recension of Maslama ibn Aḥmad al-Majrīṭī, in the Latin translation ascribed to Adelard of Bath (twelfth century). It is not al-Khwārizmī’s text. The book is transcribed in two files: the Latin text and tables of the recension, with the editor’s apparatus, and Suter’s German matter, his foreword, introduction, commentary, corrections, glossary and index, with the pages that carry no text.
\par\smallskip\noindent Every printed line of the book is one line here, in the order of the images of the Wellcome copy (Internet Archive b24872489), with its number in the margin and its anchor AK07-J\#\#\#\#-L\#\# (JP2 image, line); the lines of the footnotes follow below a short rule with the anchors N\#\#. The text is transcribed letter for letter as printed, with its punctuation, italics and spaced type, and the words divided at the ends of the lines. The folios of the manuscripts that the edition prints in the margin are set in the margin. Bookplates, stamps, shelfmarks and notes written in this copy are recorded in the page records, not reproduced. Where the print is in error, the transcription keeps it; the critical notes at the end, which are not part of the transcription, propose the correction.
\par\smallskip\noindent The Latin text and Suter’s edition are in the public domain, and this transcription is dedicated to the public domain (CC0 1.0). It is AI-integrated work (Claude); no person has proofread it. This file holds 4916 printed lines and 387 lines of the footnotes from 142 images.\par\clearpage
\FrameHead{0000}{matter of the digitisation}{0000}{matter of the digitisation}
\Withheld{matter of the digitisation: The colour card of the digitisation: a chart of 24 colour patches («24ColorCard») in a holder labelled «APRIL2013», a white patch, a card marked «LEFT», and an angled ruler in inches and centimetres, on a black ground. Not part of the book; the navigation PDF of the Internet Archive omits this image.}
\FrameHead{0001}{binding}{0001}{binding}
\Withheld{the binding: The front cover of the binding, of dark, textured cloth, without lettering on the side; photographed in low light.}
\FrameHead{0002}{library matter}{0002}{library matter}
\Withheld{library matter only: The bookplate «Ex Libris R. M. Burrell», with the drawing of a turbaned figure seated cross-legged and reading a book; below it, in pencil, the shelfmark «(2) AJB.CA.AA2»; at the foot, a barcode label «22501959681».; Nothing of the edition is printed on this page.}
\FrameHead{0003}{matter of the digitisation}{0003}{matter of the digitisation}
\Withheld{matter of the digitisation: A price in pencil, «£170», in the lower middle of the leaf.; A flyleaf without printed matter of the edition, over which the digitisation has set its notice: the drawing of a temple front, «Digitized by the Internet Archive in 2016» and «https://archive.org/details/b24872489». Not part of the book.}
\FrameHead{0004}{blank}{0004}{blank}
\Withheld{blank}
\FrameHead{0005}{p. [I]}{0005}{p. [I]}
\EdLc{1}{AK07-J0005-L01}{DIE ASTRONOMISCHEN TAFELN}{}
\EdLc{2}{AK07-J0005-L02}{\textsc{des}}{}
\EdLc{3}{AK07-J0005-L03}{MUḤAMMED IBN MŪSĀ AL-KHWĀRIZMĪ}{}
\EdLc{4}{AK07-J0005-L04}{\textsc{in der bearbeitung des}}{}
\EdLc{5}{AK07-J0005-L05}{MASLAMA IBN AḤMED AL-MADJRĪṬĪ}{}
\EdLc{6}{AK07-J0005-L06}{\textsc{und der latein. uebersetzung des}}{}
\EdLc{7}{AK07-J0005-L07}{ATHELHARD \textsc{von} BATH}{}
\EdLc{8}{AK07-J0005-L08}{\textsc{auf grund der vorarbeiten von}}{}
\EdLc{9}{AK07-J0005-L09}{A. BJØRNBO† \textsc{und} R. BESTHORN}{}
\EdLc{10}{AK07-J0005-L10}{\textsc{in kopenhagen}}{}
\EdLc{11}{AK07-J0005-L11}{\textsc{herausgegeben und kommentiert von}}{}
\EdLc{12}{AK07-J0005-L12}{H. SUTER}{}
\EdLc{13}{AK07-J0005-L13}{\textsc{in zürich}}{}
\EdLc{14}{AK07-J0005-L14}{D. K\textsc{gl}. D\textsc{anske} V\textsc{idensk}. S\textsc{elsk}. S\textsc{krifter}. 7. R\textsc{ække}. \textsc{historisk og filosofisk} A\textsc{fd}. III. 1}{}
\EdLc{15}{AK07-J0005-L15}{KØBENHAVN}{}
\EdLc{16}{AK07-J0005-L16}{HOVEDKOMMISSIONÆR: ANDR. FRED. HØST \& SØN, KGL. HOF-BOGHANDEL}{}
\EdLc{17}{AK07-J0005-L17}{BIANCO LUNOS BOGTRYKKERI}{}
\EdLc{18}{AK07-J0005-L18}{1914}{}
\FrameHead{0006}{p. [II]}{0006}{p. [II]}
\Withheld{library matter only: In pencil, the shelfmark «(2) AJB.CA.AA2», at mid-height on the left; at the lower left, the round stamp of the Wellcome Institute library, «WELLCOME» above, «LIBRARY» across the middle, «INSTITUTE» below.; Nothing is printed on the back of the title page; large stains of damp run down from the top, as on the title page (JP2 0005).}
\FrameHead{0007}{p. [III]}{0007}{p. [III]}
\EdLh{1}{AK07-J0007-L01}{INHALT}{}
\EdLr{2}{AK07-J0007-L02}{Pag.}{}
\EdLtp{3}{AK07-J0007-L03}{Vorwort}{V}{}
\EdLtp{4}{AK07-J0007-L04}{Einleitung}{VII—XXV}{}
\EdLtp{5}{AK07-J0007-L05}{Text der lateinischen Übersetzung}{1—31}{}
\EdLtp{6}{AK07-J0007-L06}{Kommentar}{32—106}{}
\EdLtp{7}{AK07-J0007-L07}{Tafeln}{109—231}{}
\EdLtp{8}{AK07-J0007-L08}{Nachträge und Berichtigungen}{233—237}{}
\EdLtp{9}{AK07-J0007-L09}{Glossar}{241—248}{}
\EdLtp{10}{AK07-J0007-L10}{Namen- und Sachregister}{249—255}{}
\EdLh{11}{AK07-J0007-L11}{Kapitel des Kommentars}{}
\EdLc{12}{AK07-J0007-L12}{(denen des lateinischen Textes entsprechend)}{}
\EdLr{13}{AK07-J0007-L13}{Pag. Tab.}{}
\EdLt{14}{AK07-J0007-L14}{Einleitung des Textes}{32}{}{}
\EdLt{15}{AK07-J0007-L15}{1. Beschreibung des arabischen Jahres}{34}{}{}
\EdLt{16}{AK07-J0007-L16}{2. Beschreibung des römischen Jahres}{34}{1}{}
\EdLt{17}{AK07-J0007-L17}{3. Über die Überschriften der Tafel-Kolonnen}{36}{}{}
\EdLt{18}{AK07-J0007-L18}{4. Über den Beginn der arabischen Monate}{36}{2—2\textsuperscript{b}}{}
\EdLt{19}{AK07-J0007-L19}{5. Über die Auffindung des römischen Jahres}{38}{3—3\textsuperscript{a}}{}
\EdLt{20}{AK07-J0007-L20}{6. Über die Einteilung der Sphäre}{41}{}{}
\EdLt{21}{AK07-J0007-L21}{7. Über die mittlere Bewegung der Planeten}{41}{4—20}{}
\EdLt{22}{AK07-J0007-L22}{8. Über die Auffindung des (wahren) Ortes der Sonne}{43}{21—26}{}
\EdLt{23}{AK07-J0007-L23}{9. Über die Auffindung des (wahren) Ortes des Mondes}{45}{21—26}{}
\EdLt{24}{AK07-J0007-L24}{10. Bestimmung des (wahren) Ortes der drei obern Planeten}{46}{27—44}{}
\EdLt{25}{AK07-J0007-L25}{11. Bestimmung des (wahren) Ortes der zwei untern Planeten}{54}{45—56}{}
\EdLt{26}{AK07-J0007-L26}{12. Bestimmung der Länge der Mondknoten}{55}{19—20}{}
\EdLt{27}{AK07-J0007-L27}{13. Über die Stationen und das Vor- und Rückwärtsschreiten der Planeten}{55}{27—56}{}
\EdLt{28}{AK07-J0007-L28}{14. Über die Zeit der Rechtläufigkeit und Rückläufigkeit}{56}{27—56}{}
\EdLt{29}{AK07-J0007-L29}{15. Über die Deklination der Sonne}{58}{21—26}{}
\EdLt{30}{AK07-J0007-L30}{16. Über die Breite des Mondes}{58}{19—26}{}
\EdLt{31}{AK07-J0007-L31}{17. Über die Breite der Planeten}{60}{27—56}{}
\EdLt{32}{AK07-J0007-L32}{18. Über Apogeum und Perigeum der Planeten}{64}{}{}
\EdLt{33}{AK07-J0007-L33}{19. Über die Knotenpunkte der Planeten}{64}{}{}
\EdLt{34}{AK07-J0007-L34}{20. u. 21. Über die tägliche Bewegung des Mondes und der Sonne}{65}{}{}
\EdLt{35}{AK07-J0007-L35}{22. Über das Erscheinen des Mondes am Abend des 29. Tages (des Monats)}{66}{57\textsuperscript{a}}{}
\FrameHead{0008}{p. [IV]}{0008}{p. [IV]}
\EdLr{1}{AK07-J0008-L01}{Pag, Tab.}{}
\EdLt{2}{AK07-J0008-L02}{23. Über die Auffindung des Sinus aus dem Bogen und umgekehrt}{68}{58—58\textsuperscript{a}}{}
\EdLt{3}{AK07-J0008-L03}{24. Über die Auffindung der Breite eines beliebigen Ortes}{71}{}{}
\EdLt{4}{AK07-J0008-L04}{25. Die Aufgänge der Tierkreiszeichen in der geraden Sphäre}{72}{59—59\textsuperscript{b}}{}
\EdL{5}{AK07-J0008-L05}{26. u. 26\textsuperscript{a}. Mit wie viel Graden (des Aequators) an einem beliebigen Orte jedes Tierkreis-}{}
\EdLt{6}{AK07-J0008-L06}{zeichen aufsteigt}{73}{59—59\textsuperscript{b}}{}
\EdLt{7}{AK07-J0008-L07}{27. Bestimmung der Grösse der Zeitstunde eines beliebigen Tages für einen beliebigen Ort}{76}{}{}
\EdLt{8}{AK07-J0008-L08}{28.—28\textsuperscript{b}. Wie man aus der Höhe der Sonne die beiden Schatten bestimmt u. umgekehrt}{76}{60}{}
\EdLt{9}{AK07-J0008-L09}{29. Über elbuht und seine Auffindung}{78}{61—66}{}
\EdLt{10}{AK07-J0008-L10}{30. u. 30\textsuperscript{a}. Über die scheinbaren Sonnen- und Monddurchmesser}{78}{61—66}{}
\EdLt{11}{AK07-J0008-L11}{Die Tafeln der Tagesgleichung (Zeitgleichung)}{80}{67—68}{}
\EdLt{12}{AK07-J0008-L12}{31. Über Konjunktion und Opposition von Sonne und Mond}{81}{69—72}{}
\EdLt{13}{AK07-J0008-L13}{33. Wie die Verfinsterung des Mondes bestimmt wird}{86}{73—76}{}
\EdLt{14}{AK07-J0008-L14}{34. Über die Parallaxen des Mondes in Länge und Breite}{91}{77—77\textsuperscript{a}}{}
\EdLt{15}{AK07-J0008-L15}{35. Über die Verfinsterung der Sonne}{94}{78}{}
\EdLt{16}{AK07-J0008-L16}{32. u. 36. Über die Verteilung der zwölf astrologischen Häuser}{96}{79—90}{}
\EdLt{17}{AK07-J0008-L17}{37. Über die Örter der Strahlenwerfung}{98}{91—114}{}
\EdLt{18}{AK07-J0008-L18}{Über die Umwandlung der Geburtsjahre}{102}{115}{}
\EdLt{19}{AK07-J0008-L19}{Über die Domicilia, Exaltationes, Decani, etc. der Planeten}{103}{116}{}
\FrameHead{0009}{p. [V]}{0009}{p. [V]}
\EdLh{1}{AK07-J0009-L01}{VORWORT}{}
\EdLp{2}{AK07-J0009-L02}{Die Herausgabe dieser astronomischen Tafeln wurde von dem leider zu früh}{}
\EdL{3}{AK07-J0009-L03}{verstorbenen, hervorragenden dänischen Gelehrten A. \Sperr{Björnbo} vor etwa fünf Jahren}{}
\EdL{4}{AK07-J0009-L04}{begonnen. Nach seinem im Oktober 1911 erfolgten Tode wurde mir sein Manuskript}{}
\EdL{5}{AK07-J0009-L05}{von der hohen kgl. dänischen Gesellschaft der Wissenschaften übermittelt, nachdem}{}
\EdL{6}{AK07-J0009-L06}{ich mich bereit erklärt hatte, die von A. \Sperr{Björnbo} begonnene Arbeit zu Ende zu}{}
\EdL{7}{AK07-J0009-L07}{führen. Dieses Manuskript enthielt den Text dieses Werkes mit den Lesarten, wie}{}
\EdL{8}{AK07-J0009-L08}{sie sich aus der Kollation von drei Handschriften (\textit{Oxford}, \textit{Chartres} und \textit{Paris}) ergeben}{}
\EdL{9}{AK07-J0009-L09}{hatten, ebenso die Tafeln mit den Lesarten dieser Handschriften, sowie die Übertragung}{}
\EdL{10}{AK07-J0009-L10}{der Mehrzahl der darin auftretenden arabischen Ausdrücke ins Lateinische durch den}{}
\EdL{11}{AK07-J0009-L11}{verdienten dänischen Arabisten, Herrn Dr. R. \Sperr{Besthorn} in Kopenhagen. A. \Sperr{Björnbo}}{}
\EdL{12}{AK07-J0009-L12}{hatte noch kurz vor seinem Tode vernommen, dass noch eine vierte Handschrift}{}
\EdL{13}{AK07-J0009-L13}{dieser \Sperr{Athelhard}’schen Übersetzung in der Nationalbibliothek zu \textit{Madrid} vorhanden}{}
\EdL{14}{AK07-J0009-L14}{sei\,; mir blieb also \Sperr{erstens} die Aufgabe, auch diese Handschrift noch zur Verglei-}{}
\EdL{15}{AK07-J0009-L15}{chung herbeizuziehen\,; \Sperr{zweitens} musste ich auch das Ms. von \textit{Chartres} noch einmal}{}
\EdL{16}{AK07-J0009-L16}{zur Revision kommen lassen, indem A. \Sperr{Björnbo} auf verschiedenen Blättern seines}{}
\EdL{17}{AK07-J0009-L17}{Manuskriptes auf eine nochmalige Vergleichung mit dieser Handschrift hingewiesen}{}
\EdL{18}{AK07-J0009-L18}{hatte\,; \Sperr{drittens} war es notwendig, zu dem kurzen, oft schwer verständlichen Text}{}
\EdL{19}{AK07-J0009-L19}{einen Kommentar auszuarbeiten, und der Arbeit eine Einleitung, ein Glossar, sowie}{}
\EdL{20}{AK07-J0009-L20}{ein Namen- u. Sachregister beizufügen. Ende Juli 1913 reichte ich der kgl. dänischen}{}
\EdL{21}{AK07-J0009-L21}{Gesellschaft der Wissenschaften meine Arbeit zur Prüfung ein, im December desselben}{}
\EdL{22}{AK07-J0009-L22}{Jahres beschloss sie, dieselbe in ihren \textit{Skrifter} zu veröffentlichen, wofür ich ihr hier}{}
\EdL{23}{AK07-J0009-L23}{meinen verbindlichen Dank ausspreche.}{}
\FrameHead{0010}{p. VI}{0010}{p. VI}
\EdLp{1}{AK07-J0010-L01}{Zu Dank verpflichtet bin ich ebenfalls den Herren Professoren Dr. D. \Sperr{Andersen}}{}
\EdL{2}{AK07-J0010-L02}{in Kopenhagen, und Dr. G. \Sperr{Billeter} in Zürich für die gütige Hilfe, die sie mir bei}{}
\EdL{3}{AK07-J0010-L03}{der Korrektur geleistet haben. Leider trifft mein Dank an einen dritten Gelehrten}{}
\EdL{4}{AK07-J0010-L04}{diesen nicht mehr unter den Lebenden\,: Herr C. F. \Sperr{Pechüle}, Observator an der}{}
\EdL{5}{AK07-J0010-L05}{kgl. Sternwarte in Kopenhagen, der im Mai dieses Jahres gestorben ist, hat mich auf}{}
\EdL{6}{AK07-J0010-L06}{einige Punkte im Kommentar aufmerksam gemacht, die eine etwas andere Auffassung}{}
\EdL{7}{AK07-J0010-L07}{und Darstellung verlangten\,; ich habe seine Bemerkungen im Nachtrag 1) in Berück-}{}
\EdL{8}{AK07-J0010-L08}{sichtigung gezogen.}{}
\EdLp{9}{AK07-J0010-L09}{Der Unterzeichnete hegt die Hoffnung, dass die Veröffentlichung dieses Werkes}{}
\EdL{10}{AK07-J0010-L10}{dazu beitragen möge, die Kenntnis der ältern arabischen Astronomie zu fördern, und}{}
\EdL{11}{AK07-J0010-L11}{besonders das Verständnis der astronomischen Tafeln des abendländischen Mittel-}{}
\EdL{12}{AK07-J0010-L12}{alters, die sehr wahrscheinlich zum grossen Teile diese Athelhard’sche Übersetzung}{}
\EdL{13}{AK07-J0010-L13}{zur Grundlage haben, zu erleichtern.}{}
\EdLp{14}{AK07-J0010-L14}{Küsnacht, bei Zürich, im Juli 1914.}{}
\EdLr{15}{AK07-J0010-L15}{\textit{H. Suter.}}{}
\FrameHead{0011}{p. [VII]}{0011}{p. [VII]}
\EdLh{1}{AK07-J0011-L01}{EINLEITUNG}{}
\EdLp{2}{AK07-J0011-L02}{Schon im Jahre 1846 hat der hervorragende französische Gelehrte M. \Sperr{Chasles}}{}
\EdL{3}{AK07-J0011-L03}{in einem Artikel der \textit{Comptes rendus de l’Académie des Sciences}\textsuperscript{1}) auf die Wichtigkeit}{}
\EdL{4}{AK07-J0011-L04}{der \Sperr{Athelhard}’schen Übersetzung der Tafeln \Sperr{Al-Khwārizmīs} für die Geschichte}{}
\EdL{5}{AK07-J0011-L05}{der Astronomie aufmerksam gemacht. Er hat aber das Werk nicht eingehend stu-}{}
\EdL{6}{AK07-J0011-L06}{diert, da er keine Herausgabe desselben beabsichtigt hatte, deshalb halten auch die}{}
\EdL{7}{AK07-J0011-L07}{meisten Folgerungen, die er aus seinem Studium gezogen hat, heute vor der Kritik}{}
\EdL{8}{AK07-J0011-L08}{nicht mehr stand. \Sperr{Chasles} wusste nämlich damals noch nicht, dass verschiedene}{}
\EdL{9}{AK07-J0011-L09}{arabische Autoren dem spanisch-arabischen Mathematiker und Astronomen \Sperr{Maslama}}{}
\EdL{10}{AK07-J0011-L10}{\Sperr{al-Madjrīṭī}\textsuperscript{2}) eine Bearbeitung, oder besser Umarbeitung der Tafeln \Sperr{Al-Khwāriz-}}{}
\EdL{11}{AK07-J0011-L11}{\Sperr{mīs} zuschreiben. Diese Umarbeitung bestand hauptsächlich, wie die arabischen}{}
\EdL{12}{AK07-J0011-L12}{Biographen \Sperr{Ibn al Ḳifṭī} und \Sperr{Ibn Abī Uṣaibiʿa} nach dem \Sperr{Ḳāḍī Ibn Ṣāʿid}\textsuperscript{3})}{}
\EdL{13}{AK07-J0011-L13}{berichten, in der Umwandlung der nach der \textit{Aera} \Sperr{Yezdegirds} verfassten Tafeln in}{}
\EdL{14}{AK07-J0011-L14}{solche, die auf die \Sperr{muhammedanische} \textit{Aera} sich stützten, dann auch, wie einige}{}
\EdL{15}{AK07-J0011-L15}{Stellen des Textes und einige Tafeln beweisen, in teilweiser Ersetzung des Meridians}{}
\EdL{16}{AK07-J0011-L16}{von \textit{Arin}, den \Sperr{Al-Khwārizmī} zu Grunde gelegt hatte, durch denjenigen von \textit{Cordova},}{}
\EdL{17}{AK07-J0011-L17}{dem Wohnorte \Sperr{Al-Madjrīṭīs}.}{}
\EdLp{18}{AK07-J0011-L18}{Aus eben diesen Textstellen und Tafeln (vergl. auch den Kommentar) folgt nun}{}
\EdL{19}{AK07-J0011-L19}{mit beinahe absoluter Sicherheit, dass wir in dieser Übersetzung \Sperr{Athelhards}\textsuperscript{4}) nicht}{}
\EdL{20}{AK07-J0011-L20}{diejenige des ursprünglichen astronomischen Werkes \Sperr{Al-Khwārizmīs}, sondern}{}
\EdL{21}{AK07-J0011-L21}{der Umarbeitung desselben durch \Sperr{Maslama al-Madjrīṭī} vor uns haben. Auf}{}
\EdL{22}{AK07-J0011-L22}{diese Tatsache hat, so viel ich weiss, zuerst C. H. \Sperr{Haskins} hingedeutet in seiner}{}
\EdL{23}{AK07-J0011-L23}{Abhandlung über \textit{Adelard of Bath}\textsuperscript{5})\,; ob auch A. \Sperr{Björnbo} dieselbe erkannt habe,}{}
\EdL{24}{AK07-J0011-L24}{ergibt sich aus seinen hinterlassenen Papieren nicht, ist aber nicht wahrscheinlich,}{}
\EdL{25}{AK07-J0011-L25}{sonst hätte er wohl in seiner Abhandlung \textit{Al-Chwarizmi’s trigonometriske Tavler}\textsuperscript{6})}{}
\EdL{26}{AK07-J0011-L26}{darauf aufmerksam gemacht. Hieraus folgt natürlich, dass wir in Bezug auf Schlüsse,}{}
\NoteRule
\EdN{N1}{AK07-J0011-N01}{\textsuperscript{1}) T. XXIII, p. 848 ff.}
\EdN{N2}{AK07-J0011-N02}{\textsuperscript{2}) Vergl. H. \Sperr{Suter}, \textit{Die Mathematiker u. Astronomen der Araber}, etc. in Abhandlungen zur}
\EdN{N3}{AK07-J0011-N03}{Geschichte der mathem. Wissensch., X (1900), p. 76—77.}
\EdN{N4}{AK07-J0011-N04}{\textsuperscript{3}) Ibid. p. 106—107.}
\EdN{N5}{AK07-J0011-N05}{\textsuperscript{4}) Auch geschrieben \Sperr{Adhelhard}, \Sperr{Adelhard}, \Sperr{Adelard}, etc.}
\EdN{N6}{AK07-J0011-N06}{\textsuperscript{5}) In \textit{The English Historical Review}, July 1911, p. 493.}
\EdN{N7}{AK07-J0011-N07}{\textsuperscript{6}) In der \textit{Festskrift til H. G. Zeuthen}, København 1909.}
\FrameHead{0012}{p. VIII}{0012}{p. VIII}
\EdL{1}{AK07-J0012-L01}{die man über die Abhängigkeit der ältern arabischen Astronomie von der indischen}{}
\EdL{2}{AK07-J0012-L02}{aus diesem Werke gezogen hat und etwa noch ziehen möchte, sehr vorsichtig sein}{}
\EdL{3}{AK07-J0012-L03}{müssen, denn es ist sehr schwierig, ja in den meisten Fällen unmöglich zu ent-}{}
\EdL{4}{AK07-J0012-L04}{scheiden, was in dem vorliegenden Werke noch von \Sperr{Al-Khwārizmī} herrührt, und}{}
\EdL{5}{AK07-J0012-L05}{was \Sperr{Maslama} zuzuschreiben ist. Das Werk verliert aber dadurch keineswegs an}{}
\EdL{6}{AK07-J0012-L06}{Interesse, denn erstens wird der Hauptteil des ursprünglichen Werkes in der Bear-}{}
\EdL{7}{AK07-J0012-L07}{beitung wohl noch vorhanden sein, und zweitens werden allfällige Einflüsse, die}{}
\EdL{8}{AK07-J0012-L08}{dieses Werk auf die abendländische Astronomie des Mittelalters ausgeübt hat,}{}
\EdL{9}{AK07-J0012-L09}{jedenfalls dieser Bearbeitung des spanisch-arabischen Astronomen zu verdanken sein,}{}
\EdL{10}{AK07-J0012-L10}{denn das ursprüngliche Werk ist sehr wahrscheinlich dem christlichen Abendlande}{}
\EdL{11}{AK07-J0012-L11}{nie bekannt geworden.}{}
\EdLp{12}{AK07-J0012-L12}{Bevor wir auf die Besprechung der Manuskripte eintreten, die dieser Ausgabe}{}
\EdL{13}{AK07-J0012-L13}{zu Grunde liegen, wird es wohl am Platze sein, einige kurze biographische Notizen}{}
\EdL{14}{AK07-J0012-L14}{über die arabischen Gelehrten beizufügen, die hier in erster Linie in Frage kommen\,;}{}
\EdL{15}{AK07-J0012-L15}{für weiteres Material verweise ich auf meine oben (p. VII) zitierte Arbeit.}{}
\EdLp{16}{AK07-J0012-L16}{1. \Sperr{Muḥammed b. Mūsā al-Khwārizmī} (auch \Sperr{al-Khuwārizmī} ge-}{}
\EdL{17}{AK07-J0012-L17}{schrieben) \Sperr{Abū ʿAbdallāh}\textsuperscript{1}), gehörte zu den Astronomen, die in \Sperr{Al-Ma’mūns}}{}
\EdL{18}{AK07-J0012-L18}{(reg. als Khalife von 813—833) Diensten standen\,; er war nach Abstammung ein}{}
\EdL{19}{AK07-J0012-L19}{Perser, sein Geburts- und Todesjahr sind nicht bekannt. Er verfasste im Auftrage}{}
\EdL{20}{AK07-J0012-L20}{\Sperr{Al-Ma’mūns} einen Auszug aus einem indischen \textit{Siddhānta} (= astronomisches Werk,}{}
\EdL{21}{AK07-J0012-L21}{System), den \Sperr{Ibrāhīm} (nach Andern \Sperr{Muh. b. Ibrāhīm}) \Sperr{al-Fazārī} aus dem}{}
\EdL{22}{AK07-J0012-L22}{Indischen ins Arabische übersetzt hatte, und der bei den Arabern „der \textit{grosse Sind-}}{}
\EdL{23}{AK07-J0012-L23}{\textit{hind}“ genannt wurde, der Auszug hiess dann im Gegensatz dazu „der \textit{kleine Sindhind}“.}{}
\EdL{24}{AK07-J0012-L24}{Ebenso verfasste er astronomische Tafeln (\textit{zīdj}), die sich auf diesen indischen \textit{Sidd-}}{}
\EdL{25}{AK07-J0012-L25}{\textit{hānta} stützten, und die in zwei Ausgaben erschienen sein sollen.\textsuperscript{2}) Er bestimmte}{}
\EdL{26}{AK07-J0012-L26}{darin die mittlern Örter der Planeten nach der persischen \textit{Aera} \Sperr{Yezdegirds} und}{}
\EdL{27}{AK07-J0012-L27}{nach dem Meridian von \textit{Arin} (vergl. Kommentar). Diese Tafeln sind also nicht mehr}{}
\EdL{28}{AK07-J0012-L28}{vorhanden, dagegen ihre Umarbeitung durch \Sperr{Maslama al-Madjrīṭī} in der latei-}{}
\EdL{29}{AK07-J0012-L29}{nischen Übersetzung des \Sperr{Athelhard von Bath}. — Ferner verfasste er ein Buch}{}
\EdL{30}{AK07-J0012-L30}{über die \textit{indische Rechnungsweise}, arabisch ebenfalls nicht mehr vorhanden, in}{}
\EdL{31}{AK07-J0012-L31}{lateinischer Übersetzung (Übersetzer nicht festgestellt) herausgegeben von B. \Sperr{Bon-}}{}
\EdL{32}{AK07-J0012-L32}{\Sperr{compagni}\textsuperscript{3})\,; ein solches über die \textit{Algebra} (\textit{fī’l-djebr wa’l-muḳābala}), ins Lateinische}{}
\EdL{33}{AK07-J0012-L33}{übersetzt von \Sperr{Gerhard von Cremona} und \Sperr{Robert von Chester} (\textit{Cestrensis}),}{}
\EdL{34}{AK07-J0012-L34}{arabisch mit englischer Übersetzung herausgegeben von F. \Sperr{Rosen}\textsuperscript{4})\,; eine \textit{Geographie}}{}
\NoteRule
\EdN{N1}{AK07-J0012-N01}{\textsuperscript{1}) Als sein Beiname (\textit{Kunya}) wird auch \textit{Abū Djaʿfar} genannt, doch halte ich \textit{Abū ʿAbdallāh}}
\EdN{N2}{AK07-J0012-N02}{für den richtigen, der erstere ist wahrscheinlich aus einer Verwechslung unseres Astronomen mit dem}
\EdN{N3}{AK07-J0012-N03}{etwas später lebenden \Sperr{Muh. b. Mūsā b. Shākir} (gest. 873), der diesen Beinamen hatte, ent-}
\EdN{N4}{AK07-J0012-N04}{standen.}
\EdN{N5}{AK07-J0012-N05}{\textsuperscript{2}) Die Quellen stimmen hier nicht ganz überein, nach andern wäre es auch möglich, dass diese}
\EdN{N6}{AK07-J0012-N06}{Tafeln selbst der Auszug aus dem grossen \textit{Sindhind} gewesen wären, diese Frage ist nicht mehr zu}
\EdN{N7}{AK07-J0012-N07}{entscheiden.}
\EdN{N8}{AK07-J0012-N08}{\textsuperscript{3}) \textit{Trattati d’aritmetica}, I, Roma, 1857.}
\EdN{N9}{AK07-J0012-N09}{\textsuperscript{4}) \textit{Algebra of Mohammed ben Musa, edited and translated} by \Sperr{Frederic Rosen}, London, 1831.}
\FrameHead{0013}{p. IX}{0013}{p. IX}
\EdL{1}{AK07-J0013-L01}{(\textit{ṣūrat al-arḍ}), die eine ziemlich selbständige Bearbeitung der Geographie des \Sperr{Ptole-}}{}
\EdL{2}{AK07-J0013-L02}{\Sperr{mäus} darstellt, in Strassburg noch arabisch vorhanden\textsuperscript{1})\,; eine \textit{Chronik} (\textit{kitāb al-}}{}
\EdL{3}{AK07-J0013-L03}{\textit{ta’rīkh}), und noch drei Schriften über Sonnenuhren und Astrolabien, die alle nicht}{}
\EdL{4}{AK07-J0013-L04}{mehr vorhanden sind. Die „indische Rechenkunst“ und die „Algebra“ sind die}{}
\EdL{5}{AK07-J0013-L05}{Werke, denen man hauptsächlich die erste Bekanntschaft mit diesen Disziplinen im}{}
\EdL{6}{AK07-J0013-L06}{christlichen Abendlande zu verdanken hat\,; vielleicht, wir können dies nicht mit}{}
\EdL{7}{AK07-J0013-L07}{Sicherheit behaupten, bildeten auch seine astronomischen Tafeln in der Bearbeitung}{}
\EdL{8}{AK07-J0013-L08}{\Sperr{Maslamas} die Grundlage für die späteren astronomischen Arbeiten im west-}{}
\EdL{9}{AK07-J0013-L09}{arabischen und christlichen Mittelalter.}{}
\EdLp{10}{AK07-J0013-L10}{2. \Sperr{Maslama b. Aḥmed al-Madjrīṭī, Abū’l-Ḳāsim}, aus Cordova, lebte}{}
\EdL{11}{AK07-J0013-L11}{zur Zeit \Sperr{Ḥakems} II. (961—976) und \Sperr{Hishāms} II. (976—1013) und starb 1007/08.}{}
\EdL{12}{AK07-J0013-L12}{Er war zu seiner Zeit der bedeutendste Mathematiker und Astronom des arabischen}{}
\EdL{13}{AK07-J0013-L13}{Spaniens, seine Hauptwerke sind\,: Über das \textit{kaufmännische} oder \textit{Geschäfts-Rechnen}}{}
\EdL{14}{AK07-J0013-L14}{(\textit{al-muʿāmalāt}), weder arabisch noch in Übersetzung vorhanden\,; über \textit{Konstruktion}}{}
\EdL{15}{AK07-J0013-L15}{\textit{und Gebrauch des Astrolabiums}, arabisch noch vorhanden im Escurial\,; Übersetzung}{}
\EdL{16}{AK07-J0013-L16}{und Kommentar des \textit{Planisphärium} des \Sperr{Ptolemäus}, ins Lateinische übersetzt von}{}
\EdL{17}{AK07-J0013-L17}{\Sperr{Hermannus Dalmata}, gedruckt zu Basel 1536 und Venedig 1558\,; ein \textit{Auszug}}{}
\EdL{18}{AK07-J0013-L18}{aus den \textit{Tafeln} \Sperr{Al-Battānīs}, nicht mehr vorhanden\,; die \textit{Umarbeitung der Tafeln}}{}
\EdL{19}{AK07-J0013-L19}{\Sperr{Al-Khwārizmīs}, nur noch in der lateinischen Übersetzung des \Sperr{Athelhard von}}{}
\EdL{20}{AK07-J0013-L20}{\Sperr{Bath} vorhanden, die wir hier der Öffentlichkeit übergeben.}{}
\EdLp{21}{AK07-J0013-L21}{Wir geben hier auch noch die Stelle in wörtlicher Übersetzung wieder, die der}{}
\EdL{22}{AK07-J0013-L22}{spanisch-arabische Historiker \Sperr{Ibn Ṣāʿid} (gest. 1070 in Toledo)\textsuperscript{2}) über die Bearbeitung}{}
\EdL{23}{AK07-J0013-L23}{der Tafeln \Sperr{Al-Khwārizmīs} durch \Sperr{Maslama} uns überliefert hat\textsuperscript{3})\,: „Er beschäftigte}{}
\EdL{24}{AK07-J0013-L24}{sich mit der Bearbeitung der Tafeln \Sperr{Al-Khwārizmīs}\,; er wandelte die persische}{}
\EdL{25}{AK07-J0013-L25}{Zeitrechnung derselben in die arabische um, und bestimmte die mittlern Örter der}{}
\EdL{26}{AK07-J0013-L26}{Planeten für den Beginn der \textit{Hidjra}\,; er fügte zu dem Werke noch weitere schöne}{}
\EdL{27}{AK07-J0013-L27}{Tafeln hinzu, in denen er jedoch ihm (d. h. \Sperr{Khwārizmī}) genau folgte und nicht}{}
\EdL{28}{AK07-J0013-L28}{auf die Fehler in seinem Werke aufmerksam machte\,; ich habe auf diesen Umstand}{}
\EdL{29}{AK07-J0013-L29}{hingewiesen in meinem Buche „über die Verbesserung der Bewegung der Planeten}{}
\EdL{30}{AK07-J0013-L30}{und die Belehrung über die Irrtümer der Astronomen“.“}{}
\EdLp{31}{AK07-J0013-L31}{Diese Stelle beweist uns also, dass nicht alle Tafeln in dieser Bearbeitung von}{}
\EdL{32}{AK07-J0013-L32}{\Sperr{Al-Khwārizmī} herstammen\,; welches aber diejenigen sind, die \Sperr{Maslama} hinzu-}{}
\EdL{33}{AK07-J0013-L33}{gefügt hat, ist schwer zu entscheiden, wir haben auf diesen Punkt schon oben auf-}{}
\EdL{34}{AK07-J0013-L34}{merksam gemacht und kommen auf denselben noch im Kommentar zurück. — Mit}{}
\EdL{35}{AK07-J0013-L35}{der Bemerkung, \Sperr{Maslama} habe die Fehler \Sperr{Al-Khwārizmīs} nicht bemerkt, oder}{}
\EdL{36}{AK07-J0013-L36}{nicht auf dieselben aufmerksam gemacht, meint er wohl in erster Linie die Nicht-}{}
\EdL{37}{AK07-J0013-L37}{beachtung der Praezession der Aequinoctialpunkte.}{}
\NoteRule
\EdN{N1}{AK07-J0013-N01}{\textsuperscript{1}) Vergl. C. A. \Sperr{Nallino}\,: \textit{Al-Ḫuwārizmī e il suo rifacimento della geografia di Tolomeo}, in den}
\EdN{N2}{AK07-J0013-N02}{Memorie della Reale Accademia dei Lincei (Classe di scienze morali, storiche e filologiche), Serie 5,}
\EdN{N3}{AK07-J0013-N03}{Vol. II., P. I., Roma, 1894.}
\EdN{N4}{AK07-J0013-N04}{\textsuperscript{2}) Vergl. meine oben (p. VII) zitierte Arbeit, p. 106.}
\EdN{N5}{AK07-J0013-N05}{\textsuperscript{3}) Aus der Ausgabe seines Werkes \textit{Ṭabaḳāt al-umam} (= Klassen der Völker) durch L. \Sperr{Cheikho},}
\EdN{N6}{AK07-J0013-N06}{S. J., Beirut 1912, p. 69.}
\FrameHead{0014}{p. X}{0014}{p. X}
\EdLp{1}{AK07-J0014-L01}{Leider ist von den arabischen Texten \Sperr{Al-Khwārizmīs} sowohl als \Sperr{Maslamas},}{}
\EdL{2}{AK07-J0014-L02}{soweit unsre heutige Kenntnis der arabischen Handschriften geht, nichts mehr vor-}{}
\EdL{3}{AK07-J0014-L03}{handen\,; wir geben freilich die Hoffnung noch nicht auf, dass von der Bearbeitung}{}
\EdL{4}{AK07-J0014-L04}{\Sperr{Maslamas} noch einmal in spanischen Bibliotheken etwas gefunden werden könnte,}{}
\EdL{5}{AK07-J0014-L05}{ist doch auch erst vor drei Jahren zum ersten mal durch C. H. \Sperr{Haskins} auf das}{}
\EdL{6}{AK07-J0014-L06}{Vorhandensein einer Handschrift der \Sperr{Athelhard}’schen Übersetzung in Madrid}{}
\EdL{7}{AK07-J0014-L07}{aufmerksam gemacht worden.\textsuperscript{1})}{}
\EdLp{8}{AK07-J0014-L08}{Was die von A. \Sperr{Björnbo} und mir benutzten Handschriften anbetrifft, so ist}{}
\EdL{9}{AK07-J0014-L09}{eine genaue Beschreibung derselben mir nicht möglich, da mir nur diejenige von}{}
\EdL{10}{AK07-J0014-L10}{\textit{Chartres}, die ich noch einmal kommen lassen musste, vorgelegen hat, von denjenigen}{}
\EdL{11}{AK07-J0014-L11}{von \textit{Oxford} und \textit{Madrid} sind mir nur Photographien der Athelhard’schen Über-}{}
\EdL{12}{AK07-J0014-L12}{setzung zu Gesichte gekommen, die sehr unvollständige Handschrift der \textit{Mazarin}’schen}{}
\EdL{13}{AK07-J0014-L13}{Bibliothek in \textit{Paris} habe ich nie gesehen. Ich muss mich daher darauf beschränken,}{}
\EdL{14}{AK07-J0014-L14}{nach den mir vorliegenden Photographien und nach den Katalogen der \textit{Bibliothèque}}{}
\EdL{15}{AK07-J0014-L15}{\textit{publique} von \textit{Chartres} und der \textit{Bibliothèque Mazarine}\textsuperscript{2}) in \textit{Paris} einige wenige Notizen}{}
\EdL{16}{AK07-J0014-L16}{über diese Handschriften zu geben.}{}
\EdLp{17}{AK07-J0014-L17}{Die Manuskripte von \textit{Chartres}, Bibliothèque publique, (C) und \textit{Oxford}, Bodleian}{}
\EdL{18}{AK07-J0014-L18}{Library, (O) sind beide Pergament-Handschriften aus dem 12. Jahrh., nach den An-}{}
\EdL{19}{AK07-J0014-L19}{gaben der Verfasser der Kataloge\,; das erstere ist bezeichnet mit Nr. 214 (173)\textsuperscript{3}),}{}
\EdL{20}{AK07-J0014-L20}{das zweite mit Cod. Auct. F. I. 9. Sie sind im allgemeinen gut geschrieben und leicht}{}
\EdL{21}{AK07-J0014-L21}{zu lesen, C etwas besser als O. Der beschriebene Blattraum misst bei C 26/18 cm,}{}
\EdL{22}{AK07-J0014-L22}{bei O (nach den Photographien) im Durchschnitt 15/10 cm. Beide Handschriften}{}
\EdL{23}{AK07-J0014-L23}{sind ziemlich vollständig, doch enthält C mehrere grössere Stellen, die in O fehlen,}{}
\EdL{24}{AK07-J0014-L24}{und auch umgekehrt (s. Text). Unsere Abhandlung bildet in C den Schluss von}{}
\EdL{25}{AK07-J0014-L25}{fol. 41\textsuperscript{r} bis 102\textsuperscript{r}, in O nimmt sie fol. 99\textsuperscript{v} bis 159\textsuperscript{v} ein, fol. 159 ist aber aus Versehen}{}
\EdL{26}{AK07-J0014-L26}{ans Ende verschoben worden, es sollte zwischen 99 und 100 stehen, so dass die}{}
\EdL{27}{AK07-J0014-L27}{Arbeit eigentlich mit fol. 158 abschliesst.}{}
\EdLp{28}{AK07-J0014-L28}{C und O enthalten noch einige andere Abhandlungen, von denen ich nur die}{}
\EdL{29}{AK07-J0014-L29}{wichtigsten nennen will\,: In C befinden sich noch die erste Hälfte der Abhandlung}{}
\EdL{30}{AK07-J0014-L30}{über das Astrolabium von \Sperr{Hermannus Contractus}, und Fragmente der Geometrie}{}
\EdL{31}{AK07-J0014-L31}{\Sperr{Gerberts} und des \textit{Liber Embadorum} von \Sperr{Savasorda} (nach \Sperr{Chasles})\,; in O kom-}{}
\EdL{32}{AK07-J0014-L32}{men noch vor\,: der \textit{Computus} \Sperr{Gerlandi} und der \textit{Abacus} \Sperr{Bernelini} (\textit{auctore juniore}}{}
\EdL{33}{AK07-J0014-L33}{\textit{Bernhelmo\,!}).}{}
\EdLp{34}{AK07-J0014-L34}{Das Manuskript der Bibliothèque Mazarine in \textit{Paris} (M)\textsuperscript{4}), bezeichnet mit}{}
\EdL{35}{AK07-J0014-L35}{Nr. 3642 (1258)\textsuperscript{5}) ist ganz unvollständig\,; statt 37 Kapitel im Texte enthält es deren}{}
\EdL{36}{AK07-J0014-L36}{nur die ersten 21, und diese bilden die kleinere Hälfte des Textes, und statt c. 120}{}
\NoteRule
\EdN{N1}{AK07-J0014-N01}{\textsuperscript{1}) \textit{English Historical Review}, July 1911, p. 493—494.}
\EdN{N2}{AK07-J0014-N02}{\textsuperscript{2}) Beide im Jahre 1890 zu Paris erschienen.}
\EdN{N3}{AK07-J0014-N03}{\textsuperscript{3}) \textit{Catalogue général des mss. des biblioth. publ. de France}, Départements, Tome XI, p. 109—110.}
\EdN{N4}{AK07-J0014-N04}{\textsuperscript{4}) \Sperr{Björnbo} hatte für diese Handschrift den Buchstaben M gewählt\,; ich wollte in dieser}
\EdN{N5}{AK07-J0014-N05}{Hinsicht keine Änderungen mehr vornehmen, und habe daher für das Ms. der \textit{Biblioteca Nacional}}
\EdN{N6}{AK07-J0014-N06}{zu Madrid den Buchstaben N gewählt.}
\EdN{N7}{AK07-J0014-N07}{\textsuperscript{5}) \textit{Catalogue général des mss. des biblioth. publ. de France}, Départements, Tom. III, p. 151– 152.}
\FrameHead{0015}{p. XI}{0015}{p. XI}
\EdL{1}{AK07-J0015-L01}{Tafeln hat es deren nur 9. Es ist ebenfalls eine Pergament-Handschrift aus dem}{}
\EdL{2}{AK07-J0015-L02}{13. Jahrh., der beschriebene Blattraum misst durchschnittlich 30/23 cm. Es haben}{}
\EdL{3}{AK07-J0015-L03}{verschiedene Hände daran gearbeitet, Björnbo hat dies durch die Bezeichnungen}{}
\EdL{4}{AK07-J0015-L04}{M\textsubscript{1}, M\textsubscript{2}, etc. ausgedrückt.\textsuperscript{1}) Es ist sehr wahrscheinlich, dass M eine unvollständige}{}
\EdL{5}{AK07-J0015-L05}{und sehr flüchtige Abschrift von C sei, da auch fehlerhafte Lesarten meistens über-}{}
\EdL{6}{AK07-J0015-L06}{einstimmen. Unsere Abhandlung nimmt in diesem Ms. fol. 82\textsuperscript{r}—89\textsuperscript{v} ein. Neben}{}
\EdL{7}{AK07-J0015-L07}{ihr befinden sich darin unter anderem noch\,: \Sperr{Thebit}, \textit{de iis quae indigent expositione}}{}
\EdL{8}{AK07-J0015-L08}{\textit{antequam legatur Almagesti}, unvollständig\,; \Sperr{Jordanus}, \textit{de pondere}\,; der Traktat, der}{}
\EdL{9}{AK07-J0015-L09}{mit \textit{communis et consuetus} . . . beginnt\textsuperscript{2})\,; \textit{liber alchorismi arithmeticae qui editus est a}}{}
\EdL{10}{AK07-J0015-L10}{\textit{magistro Johanne}.}{}
\EdLp{11}{AK07-J0015-L11}{Mit der Handschrift der Biblioteca Nacional in \textit{Madrid} (N), bezeichnet mit}{}
\EdL{12}{AK07-J0015-L12}{Nr. 10016, die A. \Sperr{Björnbo} leider nicht mehr benutzen konnte, haben wir uns etwas}{}
\EdL{13}{AK07-J0015-L13}{ausführlicher zu beschäftigen, da sie von besonderer historischer Wichtigkeit ist.}{}
\EdL{14}{AK07-J0015-L14}{Sie ist eine Pergament-Handschrift aus dem 13. Jahrh. nach den Angaben des Kata-}{}
\EdL{15}{AK07-J0015-L15}{loges, die mir durch die Bemühungen des Herrn \Sperr{Miguel Asín}, Prof. des Arabischen}{}
\EdL{16}{AK07-J0015-L16}{an der Universität Madrid, dem ich hier öffentlich für seine liebenswürdige Hilfe,}{}
\EdL{17}{AK07-J0015-L17}{besonders auch zur Erlangung von Photographien der Handschrift, meinen verbind-}{}
\EdL{18}{AK07-J0015-L18}{lichsten Dank ausspreche, zur Kenntnis gekommen sind. Der beschriebene Blattraum}{}
\EdL{19}{AK07-J0015-L19}{misst nach den Photographien im Text 16/10 cm, in den Tafeln 17/12 cm. Die Hand-}{}
\EdL{20}{AK07-J0015-L20}{schrift ist sehr schön und gleichmässig geschrieben, und enthält wenig Abkürzungen,}{}
\EdL{21}{AK07-J0015-L21}{die Initialen sind in schwarz, rot und grün ausgeführt, auch die Zahlen der ver-}{}
\EdL{22}{AK07-J0015-L22}{schiedenen Kolonnen der Tafeln und die Titel derselben wechseln oft (nicht durch-}{}
\EdL{23}{AK07-J0015-L23}{weg) in schwarzer, roter und grüner Tinte mit einander ab\textsuperscript{3}). Die Handschrift}{}
\EdL{24}{AK07-J0015-L24}{befand sich ursprünglich in der Kathedrale von \textit{Toledo}. Herr Prof. \Sperr{Asín} bemerkte}{}
\EdL{25}{AK07-J0015-L25}{mir in einem Briefe, dass es wahrscheinlich sei, dass dieselbe aus Italien nach Spa-}{}
\EdL{26}{AK07-J0015-L26}{nien gekommen sei, weil sie einen Teil der Sammlung des Cardinals \Sperr{Celada} gebildet}{}
\EdL{27}{AK07-J0015-L27}{habe, die in Italien zusammengestellt worden sei. Wir glauben doch eher, dass sie}{}
\EdL{28}{AK07-J0015-L28}{erst in Spanien in diese Sammlung geraten sei, die wohl in Italien begonnen, aber}{}
\EdL{29}{AK07-J0015-L29}{vielleicht in Spanien fortgesetzt worden sein mag.}{}
\EdLp{30}{AK07-J0015-L30}{Auch diese schöne Handschrift ist leider unvollständig, von den 37 Kapiteln}{}
\EdL{31}{AK07-J0015-L31}{des Textes enthält sie nur 1—5, 22 u. 23, 31—37, also noch weniger als M, dafür}{}
\EdL{32}{AK07-J0015-L32}{aber von den c. 125 Tafeln deren 110 (es fehlen nur die Nummern\,: 4—8, 21, 57\textsuperscript{b}, 59,}{}
\EdL{33}{AK07-J0015-L33}{59\textsuperscript{a}, 59\textsuperscript{b}, 60, 61, 67, 68, 115, 116).}{}
\EdLp{34}{AK07-J0015-L34}{Im ersten Satze (Titel) des Textes im Ms. N führt sich \Sperr{Rodbertus Cestrensis},}{}
\EdL{35}{AK07-J0015-L35}{d. i. \Sperr{Robert von Chester}, mit den Worten \textit{et per rodbertum cestrensem ordine digestus},}{}
\NoteRule
\EdN{N1}{AK07-J0015-N01}{\textsuperscript{1}) So scheint z. B. die erste Tafel der Gleichung Saturns (Tafel 27) in drei Abschriften von}
\EdN{N2}{AK07-J0015-N02}{vərschiedener Hand (M\textsubscript{1}, M\textsubscript{2}, M\textsubscript{3}) auf fol. 86 u. 87 vorhanden zu sein, die zweite Tafel der Gleichung}
\EdN{N3}{AK07-J0015-N03}{dieses Planeten (Tafel 28) in zwei Abschriften (M\textsubscript{1} u. M\textsubscript{2}) auf denselben Blättern\,; die folgenden}
\EdN{N4}{AK07-J0015-N04}{Tafeln fehlen alle.}
\EdN{N5}{AK07-J0015-N05}{\textsuperscript{2}) Vergl. G. \Sperr{Eneström}, in der \textit{Biblioth. mathemat.}, 3. Folge, Bd. 8, p. 135 ff.}
\EdN{N6}{AK07-J0015-N06}{\textsuperscript{3}) Dies wurde für die Photographie der Tafeln zum Verhängnis, die mit grüner und roter Tinte}
\EdN{N7}{AK07-J0015-N07}{geschriebenen Titel und Zahlen wurden im Bilde so schwach, dass die beste Loupe sogar an einigen}
\EdN{N8}{AK07-J0015-N08}{Stellen versagt hat.}
\FrameHead{0016}{p. XII}{0016}{p. XII}
\EdL{1}{AK07-J0016-L01}{die in den übrigen drei Handschriften fehlen, als „Ordner“ der Übersetzung des \Sperr{Athel-}}{}
\EdL{2}{AK07-J0016-L02}{\Sperr{hard von Bath} ein. Diese ordnende Hand spürt man nun allerdings in verschiedener}{}
\EdL{3}{AK07-J0016-L03}{Hinsicht im Texte, wie wir unten noch sehen werden, nicht aber in der Reihenfolge}{}
\EdL{4}{AK07-J0016-L04}{der Tafeln, da zeigt sich im Ms. N eine arge Unordnung, der wir in C und O nicht}{}
\EdL{5}{AK07-J0016-L05}{begegnen, aber jedenfalls dürfen wir diese nicht \Sperr{Robert von Chester} in die Schuhe}{}
\EdL{6}{AK07-J0016-L06}{schieben, sondern spätern Abschreibern oder Bindern der Handschrift.}{}
\EdLp{7}{AK07-J0016-L07}{L. C. \Sperr{Karpinski}\textsuperscript{1}) betrachtet \Sperr{Robert von Chester} identisch mit \Sperr{Robertus}}{}
\EdL{8}{AK07-J0016-L08}{\Sperr{Ketenensis}, oder \Sperr{Retinensis} (\Sperr{Robert von Reading}), der den Korān zuerst}{}
\EdL{9}{AK07-J0016-L09}{ins Lateinische übersetzt (1143) und sich dann, wie er selbst bemerkt\textsuperscript{2}), dem Studium}{}
\EdL{10}{AK07-J0016-L10}{und der Übersetzung astronomisch-astrologischer Werke der Araber zugewandt hat.}{}
\EdL{11}{AK07-J0016-L11}{Bekanntlich existiert von ihm auch eine Übersetzung der \textit{Algebra} \Sperr{Al-Khwārizmīs}\textsuperscript{3}).}{}
\EdL{12}{AK07-J0016-L12}{— Es ist schade, dass das Resultat dieser ordnenden Hand \Sperr{Robert von Chesters}}{}
\EdL{13}{AK07-J0016-L13}{uns leider wieder ungeordnet und unvollständig zugekommen ist.}{}
\EdLp{14}{AK07-J0016-L14}{Unsere Abhandlung nimmt in der Handschrift N fol. 8\textsuperscript{r}—68\textsuperscript{v} ein. Ausserdem}{}
\EdL{15}{AK07-J0016-L15}{finden sich in diesem Codex noch eine Reihe von meist anonymen Texten und}{}
\EdL{16}{AK07-J0016-L16}{Tafeln mit astronomischem und astrologischem Inhalt, von denen uns nur die un-}{}
\EdL{17}{AK07-J0016-L17}{mittelbar nach unserer Abhandlung folgenden (fol. 69—84, d. h. bis zum Schluss)}{}
\EdL{18}{AK07-J0016-L18}{interessieren\,; es sind Tafeln über verschiedene astronomische Gebiete, besonders}{}
\EdL{19}{AK07-J0016-L19}{über Finsternisse, verfasst von dem Engländer \Sperr{Roger von Hereford} (c. 1150—1220)\textsuperscript{4})}{}
\EdL{20}{AK07-J0016-L20}{für die Breite und Länge dieses Ortes, die zu 52° bezw. 24° angegeben werden. So}{}
\EdL{21}{AK07-J0016-L21}{heisst es z. B. fol. 71\,: \textit{Sunt autem hae tabulae compositae super Herefordiam a magistro}}{}
\EdL{22}{AK07-J0016-L22}{\Sperr{Rogerio Herefordensi}, \textit{sicut et praemissa tabula de eclipsi solis et regulae prae-}}{}
\EdL{23}{AK07-J0016-L23}{\textit{scriptae. Est autem longitudo Herefordiae ab occidente arcus XXIV graduum, latitudo}}{}
\EdL{24}{AK07-J0016-L24}{\textit{vero LII}\textsuperscript{5}). — Fol. 73 folgt\,: \textit{Tabula medii motus solis super mediam noctem Herefordiae}}{}
\EdL{25}{AK07-J0016-L25}{\textit{secundum annos domini}. — Diese Tafeln können als Ergänzungen der Übersetzung}{}
\EdL{26}{AK07-J0016-L26}{\Sperr{Athelhards} aufgefasst werden, denn gerade diese Partie des Werkes über die}{}
\EdL{27}{AK07-J0016-L27}{Finsternisse lässt zu wünschen übrig, indem die Tafeln der \textit{Ascensiones obliquae} für}{}
\EdL{28}{AK07-J0016-L28}{\textit{Cordova} fehlen\,; dieses hat nun wohl \Sperr{Roger von Hereford} nachgeholt, aber nicht}{}
\EdL{29}{AK07-J0016-L29}{für Cordova, sondern für seinen Heimatsort.}{}
\EdLp{30}{AK07-J0016-L30}{In welcher Beziehung stehen nun diese Handschriften zu einander\,? Wie ich oben}{}
\EdL{31}{AK07-J0016-L31}{schon erwähnt habe, ist M höchst wahrscheinlich eine flüchtige Abschrift von C. Was}{}
\EdL{32}{AK07-J0016-L32}{N anbetrifft, so sind wir der Ansicht, \Sperr{Robert von Chester} habe sowohl C als O,}{}
\EdL{33}{AK07-J0016-L33}{bezw. ihre Archetypen, vor sich gehabt\,; im Texte folgt er fast wörtlich C, nur wenige}{}
\EdL{34}{AK07-J0016-L34}{Abweichungen kommen vor\,; in den Zahlen der Tafeln folgt er bald der Lesart von C,}{}
\EdL{35}{AK07-J0016-L35}{bald der von O, doch ist das erste häufiger. Die verbessernde Hand \Sperr{Roberts} zeigt}{}
\EdL{36}{AK07-J0016-L36}{sich vor allem darin, dass er die vielen arabischen Wörter, die C und O unübersetzt}{}
\EdL{37}{AK07-J0016-L37}{oder neben der Übersetzung beibehalten haben, mit ganz wenigen Ausnahmen durch}{}
\NoteRule
\EdN{N1}{AK07-J0016-N01}{\textsuperscript{1}) \textit{Biblioth. mathem.} 3. Folge, 11. Bd. (1911) p. 127.}
\EdN{N2}{AK07-J0016-N02}{\textsuperscript{2}) In der Einleitung zu seiner Korān-Übersetzung.}
\EdN{N3}{AK07-J0016-N03}{\textsuperscript{3}) \textit{Biblioth. mathem.} ibid. p. 125—131.}
\EdN{N4}{AK07-J0016-N04}{\textsuperscript{4}) Vergl. meine Abhandlung\,: \textit{Die Mathematik auf den Universitäten des Mittelalters}, p. 21 (Wis-}
\EdN{N5}{AK07-J0016-N05}{senschaftl. Beilage zum Programm der Kantonsschule in Zürich, 1887).}
\EdN{N6}{AK07-J0016-N06}{\textsuperscript{5}) An einer andern Stelle ist die Breite zu 51° 30′ angegeben.}
\FrameHead{0017}{p. XIII}{0017}{p. XIII}
\EdL{1}{AK07-J0017-L01}{lateinische ersetzt hat\,; so hat er für \textit{dakaicae} stets \textit{minuta}, für \textit{eltheniae secunda}, für}{}
\EdL{2}{AK07-J0017-L02}{\textit{elwacat medius cursus}, für \textit{tadil coaequatio} (nicht \textit{aequatio}), für \textit{elgeib sinus}, für \textit{elfelek}}{}
\EdL{3}{AK07-J0017-L03}{\textit{elmustakim circulus aequatorialis} oder \textit{aequinoctialis}, etc.}{}
\EdLp{4}{AK07-J0017-L04}{Am schwierigsten ist das Verhältnis von C zu O und die Beziehung beider}{}
\EdL{5}{AK07-J0017-L05}{zum Autographen \Sperr{Athelhards} zu erklären. \Sperr{Björnbo}, der hierüber keine schrift-}{}
\EdL{6}{AK07-J0017-L06}{lichen Angaben hinterlassen hat, war nach einer brieflichen Mitteilung von Herrn}{}
\EdL{7}{AK07-J0017-L07}{Prof. \Sperr{Heiberg} in Kopenhagen an mich der Ansicht, „dass nur Abschriften zweier}{}
\EdL{8}{AK07-J0017-L08}{nicht vollendeten Arbeitsexemplare des Übersetzers vorlägen, so dass die Restitution}{}
\EdL{9}{AK07-J0017-L09}{eines Urtextes unmöglich wäre\,; vielmehr müssten in vielen Fällen beide Texte}{}
\EdL{10}{AK07-J0017-L10}{neben einander gedruckt werden.“ So stark weichen nun doch die beiden Texte}{}
\EdL{11}{AK07-J0017-L11}{C und O nicht von einander ab, dass dies notwendig wäre, \Sperr{Björnbo} hat selbst nur}{}
\EdL{12}{AK07-J0017-L12}{in einem einzigen Kapitel (Nr. 7) dies getan, und wir sind ihm hierin gefolgt. Darin}{}
\EdL{13}{AK07-J0017-L13}{hat er recht, dass eine Restitution des Urtextes unmöglich ist, und dass sehr wahr-}{}
\EdL{14}{AK07-J0017-L14}{scheinlich Text und Tafeln auch in C und O unvollständig sind, wenn auch in}{}
\EdL{15}{AK07-J0017-L15}{geringem Masse. Die Abweichungen sind aber immerhin derart, dass sie nicht vom}{}
\EdL{16}{AK07-J0017-L16}{blossen Abschreiben der einen Handschrift von der andern herrühren können\,; darum}{}
\EdL{17}{AK07-J0017-L17}{hat sich uns die Vermutung aufgedrängt, der Schreiber von O oder von C, welcher}{}
\EdL{18}{AK07-J0017-L18}{von beiden können wir nicht entscheiden, habe neben der Übersetzung \Sperr{Athelhards}}{}
\EdL{19}{AK07-J0017-L19}{noch eine andere vor sich gehabt, und habe nun bald aus der einen, bald aus der}{}
\EdL{20}{AK07-J0017-L20}{andern den Text genommen, der ihm als der bessere erschien. In der Tat ist es}{}
\EdL{21}{AK07-J0017-L21}{sehr wahrscheinlich, dass auch \Sperr{Hermannus Dalmata}, Zeitgenosse und Freund des}{}
\EdL{22}{AK07-J0017-L22}{\Sperr{Robert v. Chester}, die (von \Sperr{Maslama} umgeänderten) Tafeln \Sperr{Al-Khwārizmīs}}{}
\EdL{23}{AK07-J0017-L23}{übersetzt habe. An einer Stelle der Übersetzung des von \Sperr{Maslama} kommentierten}{}
\EdL{24}{AK07-J0017-L24}{\textit{Planisphärium} des \Sperr{Ptolemäus}, die lange Zeit fälschlich dem \Sperr{Rudolf von Brügge}}{}
\EdL{25}{AK07-J0017-L25}{zugeschrieben wurde, sagt \Sperr{Hermannus}\,: \textit{Quem locum a Ptolemaeo minus diligenter}}{}
\EdL{26}{AK07-J0017-L26}{\textit{perspectum cum} \Sperr{Albatene} \textit{miratur et} \Sperr{Alchoarismus}, \textit{quorum hunc quidem opera}}{}
\EdL{27}{AK07-J0017-L27}{\textit{nostra Latium habet, illius vero commodissima translatio} \Sperr{Roberti} \textit{mei industria, Latinae}}{}
\EdL{28}{AK07-J0017-L28}{\textit{orationis thesaurum accumulat}\textsuperscript{1}). Sollte damit vielleicht die verschiedene Wiedergabe}{}
\EdL{29}{AK07-J0017-L29}{des arabischen Buchstabens \textit{schīn} in den beiden Mss. C und O durch \textit{sh} (englisch),}{}
\EdL{30}{AK07-J0017-L30}{bezw. \textit{sc} (dalmat.-italien.\,?) in Zusammenhang stehen\,? Diese Frage weiter zu ver-}{}
\EdL{31}{AK07-J0017-L31}{folgen und wenn möglich zu entscheiden, muss ich den Herren Philologen überlassen.}{}
\EdL{32}{AK07-J0017-L32}{Es tritt dann auch die Frage noch hinzu, ob \Sperr{Hermannus Dalmata} auch die}{}
\EdL{33}{AK07-J0017-L33}{Übersetzung \Sperr{Athelhards} zur Vergleichung herbeigezogen habe\,; als Freund von}{}
\EdL{34}{AK07-J0017-L34}{\Sperr{Robert v. Chester}, dem „Ordner“ derselben, wird er wohl mit dieser Übersetzung}{}
\EdL{35}{AK07-J0017-L35}{bekannt gewesen sein.}{}
\EdLp{36}{AK07-J0017-L36}{Da ich die Photographien des Ms. N aus Madrid erst erhielt (hauptsächlich}{}
\EdL{37}{AK07-J0017-L37}{durch meine Schuld), als der Text und die Tafeln schon fertig für den Druck}{}
\NoteRule
\EdN{N1}{AK07-J0017-N01}{\textsuperscript{1}) Vergl. \Sperr{Wüstenfeld}, \textit{Die Übersetzungen arab. Werke in das Lateinische seit dem 11. Jahrh.},}
\EdN{N2}{AK07-J0017-N02}{p. 52 (in Abhandlungen der k. Gesellsch. d. Wissensch. zu Göttingen, histor.-philolog. Kl. XXII, 2,}
\EdN{N3}{AK07-J0017-N03}{1877)\,; \Sperr{Steinschneider}, \textit{Die europäischen Übersetzungen aus dem Arabischen}, A, p. 34 (in Sitzungs-}
\EdN{N4}{AK07-J0017-N04}{berichte der k. Akademie d. Wissensch. in Wien, philos. histor. Kl. CXLIX, 1904)\,; \textit{Al-Battānī}, edid.}
\EdN{N5}{AK07-J0017-N05}{\Sperr{Nallino}, I, Praefatio, p. XLIX.}
\FrameHead{0018}{p. XIV}{0018}{p. XIV}
\EdL{1}{AK07-J0018-L01}{geschrieben waren, habe ich, um nicht zu viel Korrekturen in den Fussnoten}{}
\EdL{2}{AK07-J0018-L02}{anbringen zu müssen, folgenden Weg eingeschlagen\,: In den Fällen, wo N mit C}{}
\EdL{3}{AK07-J0018-L03}{übereinstimmt, und diese bilden die Mehrzahl, habe ich das Ms. N meistens gar}{}
\EdL{4}{AK07-J0018-L04}{nicht erwähnt\,; wo also N in den Fussnoten fehlt, ist seine Lesart gleich derjenigen}{}
\EdL{5}{AK07-J0018-L05}{von C. — Da N ziemlich viel Abweichungen von C und O in den Kolonnen-Titeln}{}
\EdL{6}{AK07-J0018-L06}{aufweist, so habe ich solche für die beiden Hauptgruppen von Tafeln im folgenden}{}
\EdL{7}{AK07-J0018-L07}{zusammengestellt, um sie nicht immer wieder bei jeder Tafel in den Noten angeben}{}
\EdL{8}{AK07-J0018-L08}{zu müssen\,:}{}
\EdLp{9}{AK07-J0018-L09}{In Tafel 9 bis 20, die die mittlere Bewegung der Planeten und des Drachen-}{}
\EdL{10}{AK07-J0018-L10}{kopfes der Mondbahn enthalten, hat N}{}
\EdL{11}{AK07-J0018-L11}{für \textit{Medialitas} stets \textit{Medius cursus}}{}
\EdL{12}{AK07-J0018-L12}{„ \textit{plani} „ \textit{residui}}{}
\EdL{13}{AK07-J0018-L13}{„ \textit{Heza} „ \textit{Argumentum}}{}
\EdL{14}{AK07-J0018-L14}{„ \textit{partes horarum} bisweilen \textit{minuta horarum}}{}
\EdLp{15}{AK07-J0018-L15}{Die Monatsnamen sind in N fast durchweg (Abweichungen in Klammer) so}{}
\EdL{16}{AK07-J0018-L16}{geschrieben\,:}{}
\ColHead{left column}
\EdL{17}{AK07-J0018-L17}{\textit{Almuharram}}{}
\EdL{18}{AK07-J0018-L18}{\textit{Szafar} (\textit{Szaffar, Szaphar})}{}
\EdL{19}{AK07-J0018-L19}{\textit{Rabbe}\textsuperscript{1}) (\textit{Rabe})}{}
\EdL{20}{AK07-J0018-L20}{\textit{Rabbe}\textsuperscript{1}) (\textit{Rabe})}{}
\EdL{21}{AK07-J0018-L21}{\textit{Jumedi}\textsuperscript{1})}{}
\EdL{22}{AK07-J0018-L22}{\textit{Jumedi}\textsuperscript{1})}{}
\ColHead{right column}
\EdL{23}{AK07-J0018-L23}{\textit{Rageb}}{}
\EdL{24}{AK07-J0018-L24}{\textit{Saaben} (\textit{Szaaben})}{}
\EdL{25}{AK07-J0018-L25}{\textit{Ramadhan} (\textit{Ramadham})}{}
\EdL{26}{AK07-J0018-L26}{\textit{Sewil}}{}
\EdL{27}{AK07-J0018-L27}{\textit{Dulkidda} (\textit{Dulkadda})}{}
\EdL{28}{AK07-J0018-L28}{\textit{Dulheigga} (\textit{Dulheiga, Dulhegga})}{}
\ColHead{full width}
\EdLp{29}{AK07-J0018-L29}{In Tafel 22 bis 56 (21 fehlt in N), die die Gleichungen der Sonne, des Mondes}{}
\EdL{30}{AK07-J0018-L30}{und der fünf Planeten enthalten, hat N}{}
\EdL{31}{AK07-J0018-L31}{für \textit{Examinatio} stets \textit{Coaequatio}}{}
\EdL{32}{AK07-J0018-L32}{„ \textit{Catar} (\textit{Semita, Pagina, Via}) „ \textit{Linea}}{}
\EdL{33}{AK07-J0018-L33}{„ \textit{Computatio minuenda} meistens \textit{Computatio diminuta}}{}
\EdL{34}{AK07-J0018-L34}{„ \textit{Computatio augenda} „ \textit{Computatio aucta}}{}
\EdL{35}{AK07-J0018-L35}{„ \textit{Sublimatio examinata} stets \textit{Sublimatio definita}}{}
\EdL{36}{AK07-J0018-L36}{„ \textit{Heza} „ \textit{Argumentum}}{}
\EdLp{37}{AK07-J0018-L37}{Das Ms. N enthält nun noch einige grössere Zuzätze und Randglossen, die ich}{}
\EdL{38}{AK07-J0018-L38}{nicht gut in die Fussnoten aufnehmen konnte, da sie fast alle eines Kommentars}{}
\EdL{39}{AK07-J0018-L39}{bedürfen\,; ich gebe sie hier in der Einleitung in der Reihenfolge, wie sie sich im}{}
\EdL{40}{AK07-J0018-L40}{Texte und in den Tafeln vorfinden\,:}{}
\EdLp{41}{AK07-J0018-L41}{1) Auf fol. 8\textsuperscript{r} kommen nach den Worten \textit{omnia quae hic} (in unserm Text p. 2),}{}
\EdL{42}{AK07-J0018-L42}{die den Schluss der ersten Spalte bilden sollten, noch folgende vier Zeilen in gleicher}{}
\EdL{43}{AK07-J0018-L43}{Schrift, die offenbar nicht zu unserm Text, sondern wahrscheinlich zu dem Text}{}
\EdL{44}{AK07-J0018-L44}{von 7\textsuperscript{v} gehören\,: \textit{Ab aequinoctio usque ad solstitium augmentatio sive diminutio dierum}}{}
\EdL{45}{AK07-J0018-L45}{\textit{et noctium duorum minutorum et XL secundorum computatur}.}{}
\NoteRule
\EdN{N1}{AK07-J0018-N01}{\textsuperscript{1}) Ohne die unterscheidenden Nummern I und II.}
\FrameHead{0019}{p. XV}{0019}{p. XV}
\EdLp{1}{AK07-J0019-L01}{2) Am Schlusse der zweiten Spalte desselben fol. (in unserm Text p. 2) steht}{}
\EdL{2}{AK07-J0019-L02}{nach den Worten \textit{quandoque in hieme interdum in aestate accidat} folgende Tabelle}{}
\EdL{3}{AK07-J0019-L03}{der arabischen Monate ohne Angabe der Zahl der Tage\,:}{}
\EdL{4}{AK07-J0019-L04}{\parbox[t]{\dimexpr0.4900\linewidth-0.5em\relax}{\raggedright\small\mbox{}\textit{Almuharram}}\hspace{0.5em}\parbox[t]{\dimexpr0.4900\linewidth-0.5em\relax}{\raggedright\small\mbox{}. . . . . . . . .}\hspace{0.5em}}{}
\EdL{5}{AK07-J0019-L05}{\parbox[t]{\dimexpr0.4900\linewidth-0.5em\relax}{\raggedright\small\mbox{}\textit{Szaphar}}\hspace{0.5em}\parbox[t]{\dimexpr0.4900\linewidth-0.5em\relax}{\raggedright\small\mbox{}. . . . . . . . .}\hspace{0.5em}}{}
\EdL{6}{AK07-J0019-L06}{\parbox[t]{\dimexpr0.4900\linewidth-0.5em\relax}{\raggedright\small\mbox{}\textit{Rabbe}}\hspace{0.5em}\parbox[t]{\dimexpr0.4900\linewidth-0.5em\relax}{\raggedright\small\mbox{}. . . . . . . . .}\hspace{0.5em}}{}
\EdL{7}{AK07-J0019-L07}{\parbox[t]{\dimexpr0.4900\linewidth-0.5em\relax}{\raggedright\small\mbox{}\textit{Rabbe}}\hspace{0.5em}\parbox[t]{\dimexpr0.4900\linewidth-0.5em\relax}{\raggedright\small\mbox{}. . . . . . . . .}\hspace{0.5em}}{}
\EdL{8}{AK07-J0019-L08}{\parbox[t]{\dimexpr0.4900\linewidth-0.5em\relax}{\raggedright\small\mbox{}\textit{Jumedi}}\hspace{0.5em}\parbox[t]{\dimexpr0.4900\linewidth-0.5em\relax}{\raggedright\small\mbox{}. . . . . . . . .}\hspace{0.5em}}{}
\EdL{9}{AK07-J0019-L09}{\parbox[t]{\dimexpr0.4900\linewidth-0.5em\relax}{\raggedright\small\mbox{}\textit{Jumedi}}\hspace{0.5em}\parbox[t]{\dimexpr0.4900\linewidth-0.5em\relax}{\raggedright\small\mbox{}. . . . . . . . .}\hspace{0.5em}}{}
\EdL{10}{AK07-J0019-L10}{\parbox[t]{\dimexpr0.4900\linewidth-0.5em\relax}{\raggedright\small\mbox{}\textit{Rageb}}\hspace{0.5em}\parbox[t]{\dimexpr0.4900\linewidth-0.5em\relax}{\raggedright\small\mbox{}. . . . . . . . .}\hspace{0.5em}}{}
\EdL{11}{AK07-J0019-L11}{\parbox[t]{\dimexpr0.4900\linewidth-0.5em\relax}{\raggedright\small\mbox{}\textit{Saaben}}\hspace{0.5em}\parbox[t]{\dimexpr0.4900\linewidth-0.5em\relax}{\raggedright\small\mbox{}. . . . . . . . .}\hspace{0.5em}}{}
\EdL{12}{AK07-J0019-L12}{\parbox[t]{\dimexpr0.4900\linewidth-0.5em\relax}{\raggedright\small\mbox{}\textit{Ramadhan}}\hspace{0.5em}\parbox[t]{\dimexpr0.4900\linewidth-0.5em\relax}{\raggedright\small\mbox{}. . . . . . . . .}\hspace{0.5em}}{}
\EdL{13}{AK07-J0019-L13}{\parbox[t]{\dimexpr0.4900\linewidth-0.5em\relax}{\raggedright\small\mbox{}\textit{Sewil}}\hspace{0.5em}\parbox[t]{\dimexpr0.4900\linewidth-0.5em\relax}{\raggedright\small\mbox{}. . . . . . . . .}\hspace{0.5em}}{}
\EdL{14}{AK07-J0019-L14}{\parbox[t]{\dimexpr0.4900\linewidth-0.5em\relax}{\raggedright\small\mbox{}\textit{Dulkidda}}\hspace{0.5em}\parbox[t]{\dimexpr0.4900\linewidth-0.5em\relax}{\raggedright\small\mbox{}. . . . . . . . .}\hspace{0.5em}}{}
\EdL{15}{AK07-J0019-L15}{\parbox[t]{\dimexpr0.4900\linewidth-0.5em\relax}{\raggedright\small\mbox{}\textit{Dulheiga}}\hspace{0.5em}\parbox[t]{\dimexpr0.4900\linewidth-0.5em\relax}{\raggedright\small\mbox{}. . . . . . . . .}\hspace{0.5em}}{}
\EdLp{16}{AK07-J0019-L16}{3) Die Tafel 3, die in den Mss. C und O je eine Seite in Anspruch nimmt, ist in}{}
\EdL{17}{AK07-J0019-L17}{N auf zwei Seiten (9\textsuperscript{r} und 9\textsuperscript{v}) verteilt\,; es kommen nämlich hier einige Kolonnen}{}
\EdL{18}{AK07-J0019-L18}{hinzu, die in C und O fehlen, die aber in den \textit{Tabulae spuriae}, die \Sperr{Nallino} der Aus-}{}
\EdL{19}{AK07-J0019-L19}{gabe \Sperr{Al-Battānīs} hinzugefügt hat (II, p. 300—307) vorkommen. Es lassen sich also}{}
\EdL{20}{AK07-J0019-L20}{alle die unechten Tafeln \Sperr{Al-Battānīs} auf die \Sperr{Maslama}’sche Bearbeitung der Tafeln}{}
\EdL{21}{AK07-J0019-L21}{\Sperr{Al-Khwārizmīs} zurückführen, und damit hat sich die Vermutung, die \Sperr{Nallino}}{}
\EdL{22}{AK07-J0019-L22}{über diese Tafeln ausgesprochen hat (II, p. 308—309\,: \textit{Haud improbabile est tabulas}}{}
\EdL{23}{AK07-J0019-L23}{\textit{p. 302—307 ex opere Maslamah quoque, ut de p. 300—301 affirmatur, exscriptas esse, etc.})}{}
\EdL{24}{AK07-J0019-L24}{bewahrheitet. Die in Ms. N hinzukommenden Kolonnen sind folgende\,:}{}
\EdLp{25}{AK07-J0019-L25}{A) Auf fol. 9\textsubscript{r} kommt nach der Kolonne \textit{Anni Alexandri collecti}, die in unserer}{}
\EdL{26}{AK07-J0019-L26}{Tafel 3 die 2. Kol. ist, eine Kolonne mit der Überschrift \textit{Anni iazdeiardi} (\textit{Yezdegirdi})}{}
\EdL{27}{AK07-J0019-L27}{\textit{post aeram arabum}\,; diese stimmt mit der Kolonne 4 bei \Sperr{Al-Battānī} (II, 304), über-}{}
\EdL{28}{AK07-J0019-L28}{schrieben\,: \textit{Anni Persarum sine intercalatione, etc.}, und zwar bis auf die Radix, die}{}
\EdL{29}{AK07-J0019-L29}{dort zu 9\textsuperscript{a}11\textsuperscript{m}9\textsuperscript{d} angegeben wird (s. auch Tafel 1, wo sich auch die unrichtige Radix}{}
\EdL{30}{AK07-J0019-L30}{9\textsuperscript{a}11\textsuperscript{m}4\textsuperscript{d} findet).}{}
\EdLp{31}{AK07-J0019-L31}{B) Auf fol. 9\textsuperscript{v} befindet sich die der vorigen entsprechende Kolonne für die \textit{anni}}{}
\EdL{32}{AK07-J0019-L32}{\textit{plani} oder \textit{residui}, sie steht nach der Kolonne \textit{Anni Alexandri plani}, die in unserer}{}
\EdL{33}{AK07-J0019-L33}{Tafel 3 die 4. Kol. ist\,; sie stimmt mit der Kolonne 4 bei \Sperr{Al-Battānī} (II, 305),}{}
\EdL{34}{AK07-J0019-L34}{überschrieben\,: \textit{Anni singuli Persarum sine intercalatione}, und zwar bis auf die Tage}{}
\EdL{35}{AK07-J0019-L35}{für \textit{annus} 27, wo unsere Tafel 19, \Sperr{Al-Battānī} aber 18 hat.}{}
\EdLp{36}{AK07-J0019-L36}{C) Auf fol. 9\textsuperscript{r} befindet sich unmittelbar nach der Kolonne A für die \textit{Anni iazdeiardi}}{}
\EdL{37}{AK07-J0019-L37}{eine solche für die \textit{Anni Diocletiani completi}, sie stimmt im ganzen mit der Kol. 3}{}
\EdL{38}{AK07-J0019-L38}{bei \Sperr{Al-Battānī} (II, 304), überschrieben \textit{Anni Coptorum intercalationem habentes etc.}}{}
\EdL{39}{AK07-J0019-L39}{Bei verschiedenen arabischen Astronomen wurde die \textit{Aera Diocletiani} (Beginn am}{}
\EdL{40}{AK07-J0019-L40}{29. Aug. 284 n. Chr.) auch \textit{Aera Coptorum} genannt, mit welchem Namen sonst ver-}{}
\FrameHead{0020}{p. XVI}{0020}{p. XVI}
\EdL{1}{AK07-J0020-L01}{schiedene andere Aeren, so die \textit{Aera Nabonassaris}, diejenige \textit{Philippi Aridaei}, und von}{}
\EdL{2}{AK07-J0020-L02}{\Sperr{Al-Battānī} auch die \textit{Aera Augusti} oder \textit{Alexandrina}, über deren Beginn keine volle}{}
\EdL{3}{AK07-J0020-L03}{Übereinstimmung herrscht, bezeichnet wurden\textsuperscript{1}). Die Kolonne in N zeigt verschiedene}{}
\EdL{4}{AK07-J0020-L04}{Abweichungen von der bei \Sperr{Al-Battānī}\,: Unter a) sollte es heissen 483 statt 488,}{}
\EdL{5}{AK07-J0020-L05}{unter b) 977 statt 997\,; dann sollten natürlich die vier letzten Zeilen so lauten\,:}{}
\EdL{6}{AK07-J0020-L06}{\parbox[t]{\dimexpr0.2450\linewidth-0.5em\relax}{\raggedright\small\mbox{}920.}\hspace{0.5em}\parbox[t]{\dimexpr0.2450\linewidth-0.5em\relax}{\raggedright\small\mbox{}0.}\hspace{0.5em}\parbox[t]{\dimexpr0.2450\linewidth-0.5em\relax}{\raggedright\small\mbox{}25.}\hspace{0.5em}\parbox[t]{\dimexpr0.2450\linewidth-0.5em\relax}{\raggedright\small\mbox{}\textsuperscript{1}/\textsubscript{4}}\hspace{0.5em}}{}
\EdL{7}{AK07-J0020-L07}{\parbox[t]{\dimexpr0.2450\linewidth-0.5em\relax}{\raggedright\small\mbox{}949.}\hspace{0.5em}\parbox[t]{\dimexpr0.2450\linewidth-0.5em\relax}{\raggedright\small\mbox{}1.}\hspace{0.5em}\parbox[t]{\dimexpr0.2450\linewidth-0.5em\relax}{\raggedright\small\mbox{}4.}\hspace{0.5em}\parbox[t]{\dimexpr0.2450\linewidth-0.5em\relax}{\raggedright\small\mbox{}0}\hspace{0.5em}}{}
\EdL{8}{AK07-J0020-L08}{\parbox[t]{\dimexpr0.2450\linewidth-0.5em\relax}{\raggedright\small\mbox{}978.}\hspace{0.5em}\parbox[t]{\dimexpr0.2450\linewidth-0.5em\relax}{\raggedright\small\mbox{}2.}\hspace{0.5em}\parbox[t]{\dimexpr0.2450\linewidth-0.5em\relax}{\raggedright\small\mbox{}13.}\hspace{0.5em}\parbox[t]{\dimexpr0.2450\linewidth-0.5em\relax}{\raggedright\small\mbox{}\textsuperscript{3}/\textsubscript{4}}\hspace{0.5em}}{}
\EdL{9}{AK07-J0020-L09}{\parbox[t]{\dimexpr0.2450\linewidth-0.5em\relax}{\raggedright\small\mbox{}1007.}\hspace{0.5em}\parbox[t]{\dimexpr0.2450\linewidth-0.5em\relax}{\raggedright\small\mbox{}4.}\hspace{0.5em}\parbox[t]{\dimexpr0.2450\linewidth-0.5em\relax}{\raggedright\small\mbox{}21.}\hspace{0.5em}\parbox[t]{\dimexpr0.2450\linewidth-0.5em\relax}{\raggedright\small\mbox{}\textsuperscript{1}/\textsubscript{2}}\hspace{0.5em}}{}
\EdLp{10}{AK07-J0020-L10}{Auch \Sperr{Al-Battānī} hat 919. 12 statt 920. 0, und die letzte Zeile fehlt bei ihm.}{}
\EdL{11}{AK07-J0020-L11}{Die Zahlen der Kol. \textit{Dies} stimmen mit denen bei \Sperr{Al-Battānī}, wenn wir die 25 der}{}
\EdL{12}{AK07-J0020-L12}{zweiten Zeile von oben weglassen, und die übrigen Zahlen der Kolonne alle um eine}{}
\EdL{13}{AK07-J0020-L13}{Zeile nach oben rücken, bis auf die beiden untersten Zahlen (13 u. 21).}{}
\EdL{14}{AK07-J0020-L14}{\parbox[t]{\dimexpr0.3015\linewidth-0.5em\relax}{\raggedright\small\mbox{}A.}\hspace{0.5em}\parbox[t]{\dimexpr0.3015\linewidth-0.5em\relax}{\raggedright\small\mbox{}B.}\hspace{0.5em}\parbox[t]{\dimexpr0.3769\linewidth-0.5em\relax}{\raggedright\small\mbox{}C.}\hspace{0.5em}}{}
\EdL{15}{AK07-J0020-L15}{\parbox[t]{\dimexpr0.0754\linewidth-0.5em\relax}{\raggedright\small\mbox{}Anni collecti}\hspace{0.5em}\parbox[t]{\dimexpr0.2262\linewidth-0.5em\relax}{\raggedright\small\mbox{}Anni iazdeiardi post aeram arabum}\hspace{0.5em}\parbox[t]{\dimexpr0.0754\linewidth-0.5em\relax}{\raggedright\small\mbox{}Anni residui}\hspace{0.5em}\parbox[t]{\dimexpr0.2262\linewidth-0.5em\relax}{\raggedright\small\mbox{}Anni iezdeiardi}\hspace{0.5em}\parbox[t]{\dimexpr0.0754\linewidth-0.5em\relax}{\raggedright\small\mbox{}Anni collecti}\hspace{0.5em}\parbox[t]{\dimexpr0.3015\linewidth-0.5em\relax}{\raggedright\small\mbox{}Anni diocletiani completi}\hspace{0.5em}}{}
\EdL{16}{AK07-J0020-L16}{\parbox[t]{\dimexpr0.0754\linewidth-0.5em\relax}{\raggedright\small\mbox{}}\hspace{0.5em}\parbox[t]{\dimexpr0.0754\linewidth-0.5em\relax}{\raggedright\small\mbox{}Anni}\hspace{0.5em}\parbox[t]{\dimexpr0.0754\linewidth-0.5em\relax}{\raggedright\small\mbox{}Menses}\hspace{0.5em}\parbox[t]{\dimexpr0.1508\linewidth-0.5em\relax}{\raggedright\small\mbox{}Dies}\hspace{0.5em}\parbox[t]{\dimexpr0.0754\linewidth-0.5em\relax}{\raggedright\small\mbox{}Anni}\hspace{0.5em}\parbox[t]{\dimexpr0.0754\linewidth-0.5em\relax}{\raggedright\small\mbox{}Menses}\hspace{0.5em}\parbox[t]{\dimexpr0.1508\linewidth-0.5em\relax}{\raggedright\small\mbox{}Dies}\hspace{0.5em}\parbox[t]{\dimexpr0.0754\linewidth-0.5em\relax}{\raggedright\small\mbox{}Anni}\hspace{0.5em}\parbox[t]{\dimexpr0.0754\linewidth-0.5em\relax}{\raggedright\small\mbox{}Menses}\hspace{0.5em}\parbox[t]{\dimexpr0.0754\linewidth-0.5em\relax}{\raggedright\small\mbox{}Dies}\hspace{0.5em}\parbox[t]{\dimexpr0.0754\linewidth-0.5em\relax}{\raggedright\small\mbox{}Partes}\hspace{0.5em}}{}
\EdL{17}{AK07-J0020-L17}{\parbox[t]{\dimexpr0.0754\linewidth-0.5em\relax}{\raggedright\small\mbox{}Radix}\hspace{0.5em}\parbox[t]{\dimexpr0.0754\linewidth-0.5em\relax}{\raggedright\small\mbox{}9}\hspace{0.5em}\parbox[t]{\dimexpr0.0754\linewidth-0.5em\relax}{\raggedright\small\mbox{}11}\hspace{0.5em}\parbox[t]{\dimexpr0.0754\linewidth-0.5em\relax}{\raggedright\small\mbox{}4}\hspace{0.5em}\parbox[t]{\dimexpr0.0754\linewidth-0.5em\relax}{\raggedright\small\mbox{}1}\hspace{0.5em}\parbox[t]{\dimexpr0.0754\linewidth-0.5em\relax}{\raggedright\small\mbox{}0}\hspace{0.5em}\parbox[t]{\dimexpr0.0754\linewidth-0.5em\relax}{\raggedright\small\mbox{}11}\hspace{0.5em}\parbox[t]{\dimexpr0.0754\linewidth-0.5em\relax}{\raggedright\small\mbox{}24}\hspace{0.5em}\parbox[t]{\dimexpr0.0754\linewidth-0.5em\relax}{\raggedright\small\mbox{}Radix}\hspace{0.5em}\parbox[t]{\dimexpr0.0754\linewidth-0.5em\relax}{\raggedright\small\mbox{}337}\hspace{0.5em}\parbox[t]{\dimexpr0.0754\linewidth-0.5em\relax}{\raggedright\small\mbox{}10}\hspace{0.5em}\parbox[t]{\dimexpr0.0754\linewidth-0.5em\relax}{\raggedright\small\mbox{}19}\hspace{0.5em}\parbox[t]{\dimexpr0.0754\linewidth-0.5em\relax}{\raggedright\small\mbox{}\textsuperscript{3}/\textsubscript{4}}\hspace{0.5em}}{}
\EdL{18}{AK07-J0020-L18}{\parbox[t]{\dimexpr0.0754\linewidth-0.5em\relax}{\raggedright\small\mbox{}30}\hspace{0.5em}\parbox[t]{\dimexpr0.0754\linewidth-0.5em\relax}{\raggedright\small\mbox{}19}\hspace{0.5em}\parbox[t]{\dimexpr0.0754\linewidth-0.5em\relax}{\raggedright\small\mbox{}2}\hspace{0.5em}\parbox[t]{\dimexpr0.0754\linewidth-0.5em\relax}{\raggedright\small\mbox{}12}\hspace{0.5em}\parbox[t]{\dimexpr0.0754\linewidth-0.5em\relax}{\raggedright\small\mbox{}2}\hspace{0.5em}\parbox[t]{\dimexpr0.0754\linewidth-0.5em\relax}{\raggedright\small\mbox{}1}\hspace{0.5em}\parbox[t]{\dimexpr0.0754\linewidth-0.5em\relax}{\raggedright\small\mbox{}11}\hspace{0.5em}\parbox[t]{\dimexpr0.0754\linewidth-0.5em\relax}{\raggedright\small\mbox{}14}\hspace{0.5em}\parbox[t]{\dimexpr0.0754\linewidth-0.5em\relax}{\raggedright\small\mbox{}30}\hspace{0.5em}\parbox[t]{\dimexpr0.0754\linewidth-0.5em\relax}{\raggedright\small\mbox{}366}\hspace{0.5em}\parbox[t]{\dimexpr0.0754\linewidth-0.5em\relax}{\raggedright\small\mbox{}11}\hspace{0.5em}\parbox[t]{\dimexpr0.0754\linewidth-0.5em\relax}{\raggedright\small\mbox{}25}\hspace{0.5em}\parbox[t]{\dimexpr0.0754\linewidth-0.5em\relax}{\raggedright\small\mbox{}\textsuperscript{1}/\textsubscript{2}}\hspace{0.5em}}{}
\EdL{19}{AK07-J0020-L19}{\parbox[t]{\dimexpr0.0754\linewidth-0.5em\relax}{\raggedright\small\mbox{}60}\hspace{0.5em}\parbox[t]{\dimexpr0.0754\linewidth-0.5em\relax}{\raggedright\small\mbox{}48}\hspace{0.5em}\parbox[t]{\dimexpr0.0754\linewidth-0.5em\relax}{\raggedright\small\mbox{}3}\hspace{0.5em}\parbox[t]{\dimexpr0.0754\linewidth-0.5em\relax}{\raggedright\small\mbox{}28}\hspace{0.5em}\parbox[t]{\dimexpr0.0754\linewidth-0.5em\relax}{\raggedright\small\mbox{}3}\hspace{0.5em}\parbox[t]{\dimexpr0.0754\linewidth-0.5em\relax}{\raggedright\small\mbox{}2}\hspace{0.5em}\parbox[t]{\dimexpr0.0754\linewidth-0.5em\relax}{\raggedright\small\mbox{}11}\hspace{0.5em}\parbox[t]{\dimexpr0.0754\linewidth-0.5em\relax}{\raggedright\small\mbox{}3}\hspace{0.5em}\parbox[t]{\dimexpr0.0754\linewidth-0.5em\relax}{\raggedright\small\mbox{}60}\hspace{0.5em}\parbox[t]{\dimexpr0.0754\linewidth-0.5em\relax}{\raggedright\small\mbox{}396}\hspace{0.5em}\parbox[t]{\dimexpr0.0754\linewidth-0.5em\relax}{\raggedright\small\mbox{}1}\hspace{0.5em}\parbox[t]{\dimexpr0.0754\linewidth-0.5em\relax}{\raggedright\small\mbox{}28}\hspace{0.5em}\parbox[t]{\dimexpr0.0754\linewidth-0.5em\relax}{\raggedright\small\mbox{}0}\hspace{0.5em}}{}
\EdL{20}{AK07-J0020-L20}{\parbox[t]{\dimexpr0.0754\linewidth-0.5em\relax}{\raggedright\small\mbox{}90}\hspace{0.5em}\parbox[t]{\dimexpr0.0754\linewidth-0.5em\relax}{\raggedright\small\mbox{}77}\hspace{0.5em}\parbox[t]{\dimexpr0.0754\linewidth-0.5em\relax}{\raggedright\small\mbox{}5}\hspace{0.5em}\parbox[t]{\dimexpr0.0754\linewidth-0.5em\relax}{\raggedright\small\mbox{}14}\hspace{0.5em}\parbox[t]{\dimexpr0.0754\linewidth-0.5em\relax}{\raggedright\small\mbox{}4}\hspace{0.5em}\parbox[t]{\dimexpr0.0754\linewidth-0.5em\relax}{\raggedright\small\mbox{}3}\hspace{0.5em}\parbox[t]{\dimexpr0.0754\linewidth-0.5em\relax}{\raggedright\small\mbox{}10}\hspace{0.5em}\parbox[t]{\dimexpr0.0754\linewidth-0.5em\relax}{\raggedright\small\mbox{}22}\hspace{0.5em}\parbox[t]{\dimexpr0.0754\linewidth-0.5em\relax}{\raggedright\small\mbox{}90}\hspace{0.5em}\parbox[t]{\dimexpr0.0754\linewidth-0.5em\relax}{\raggedright\small\mbox{}425}\hspace{0.5em}\parbox[t]{\dimexpr0.0754\linewidth-0.5em\relax}{\raggedright\small\mbox{}2}\hspace{0.5em}\parbox[t]{\dimexpr0.0754\linewidth-0.5em\relax}{\raggedright\small\mbox{}2}\hspace{0.5em}\parbox[t]{\dimexpr0.0754\linewidth-0.5em\relax}{\raggedright\small\mbox{}\textsuperscript{3}/\textsubscript{4}}\hspace{0.5em}}{}
\EdL{21}{AK07-J0020-L21}{\parbox[t]{\dimexpr0.0754\linewidth-0.5em\relax}{\raggedright\small\mbox{}120}\hspace{0.5em}\parbox[t]{\dimexpr0.0754\linewidth-0.5em\relax}{\raggedright\small\mbox{}106}\hspace{0.5em}\parbox[t]{\dimexpr0.0754\linewidth-0.5em\relax}{\raggedright\small\mbox{}7}\hspace{0.5em}\parbox[t]{\dimexpr0.0754\linewidth-0.5em\relax}{\raggedright\small\mbox{}0}\hspace{0.5em}\parbox[t]{\dimexpr0.0754\linewidth-0.5em\relax}{\raggedright\small\mbox{}5}\hspace{0.5em}\parbox[t]{\dimexpr0.0754\linewidth-0.5em\relax}{\raggedright\small\mbox{}4}\hspace{0.5em}\parbox[t]{\dimexpr0.0754\linewidth-0.5em\relax}{\raggedright\small\mbox{}10}\hspace{0.5em}\parbox[t]{\dimexpr0.0754\linewidth-0.5em\relax}{\raggedright\small\mbox{}12}\hspace{0.5em}\parbox[t]{\dimexpr0.0754\linewidth-0.5em\relax}{\raggedright\small\mbox{}120}\hspace{0.5em}\parbox[t]{\dimexpr0.0754\linewidth-0.5em\relax}{\raggedright\small\mbox{}454}\hspace{0.5em}\parbox[t]{\dimexpr0.0754\linewidth-0.5em\relax}{\raggedright\small\mbox{}3}\hspace{0.5em}\parbox[t]{\dimexpr0.0754\linewidth-0.5em\relax}{\raggedright\small\mbox{}10}\hspace{0.5em}\parbox[t]{\dimexpr0.0754\linewidth-0.5em\relax}{\raggedright\small\mbox{}\textsuperscript{1}/\textsubscript{2}}\hspace{0.5em}}{}
\EdL{22}{AK07-J0020-L22}{\parbox[t]{\dimexpr0.0754\linewidth-0.5em\relax}{\raggedright\small\mbox{}150}\hspace{0.5em}\parbox[t]{\dimexpr0.0754\linewidth-0.5em\relax}{\raggedright\small\mbox{}135}\hspace{0.5em}\parbox[t]{\dimexpr0.0754\linewidth-0.5em\relax}{\raggedright\small\mbox{}8}\hspace{0.5em}\parbox[t]{\dimexpr0.0754\linewidth-0.5em\relax}{\raggedright\small\mbox{}16}\hspace{0.5em}\parbox[t]{\dimexpr0.0754\linewidth-0.5em\relax}{\raggedright\small\mbox{}6}\hspace{0.5em}\parbox[t]{\dimexpr0.0754\linewidth-0.5em\relax}{\raggedright\small\mbox{}5}\hspace{0.5em}\parbox[t]{\dimexpr0.0754\linewidth-0.5em\relax}{\raggedright\small\mbox{}10}\hspace{0.5em}\parbox[t]{\dimexpr0.0754\linewidth-0.5em\relax}{\raggedright\small\mbox{}1}\hspace{0.5em}\parbox[t]{\dimexpr0.0754\linewidth-0.5em\relax}{\raggedright\small\mbox{}150}\hspace{0.5em}\parbox[t]{\dimexpr0.0754\linewidth-0.5em\relax}{\raggedright\small\mbox{}488\textsuperscript{a})}\hspace{0.5em}\parbox[t]{\dimexpr0.0754\linewidth-0.5em\relax}{\raggedright\small\mbox{}4}\hspace{0.5em}\parbox[t]{\dimexpr0.0754\linewidth-0.5em\relax}{\raggedright\small\mbox{}19}\hspace{0.5em}\parbox[t]{\dimexpr0.0754\linewidth-0.5em\relax}{\raggedright\small\mbox{}\textsuperscript{1}/\textsubscript{4}}\hspace{0.5em}}{}
\EdL{23}{AK07-J0020-L23}{\parbox[t]{\dimexpr0.0754\linewidth-0.5em\relax}{\raggedright\small\mbox{}180}\hspace{0.5em}\parbox[t]{\dimexpr0.0754\linewidth-0.5em\relax}{\raggedright\small\mbox{}164}\hspace{0.5em}\parbox[t]{\dimexpr0.0754\linewidth-0.5em\relax}{\raggedright\small\mbox{}10}\hspace{0.5em}\parbox[t]{\dimexpr0.0754\linewidth-0.5em\relax}{\raggedright\small\mbox{}2}\hspace{0.5em}\parbox[t]{\dimexpr0.0754\linewidth-0.5em\relax}{\raggedright\small\mbox{}7}\hspace{0.5em}\parbox[t]{\dimexpr0.0754\linewidth-0.5em\relax}{\raggedright\small\mbox{}6}\hspace{0.5em}\parbox[t]{\dimexpr0.0754\linewidth-0.5em\relax}{\raggedright\small\mbox{}9}\hspace{0.5em}\parbox[t]{\dimexpr0.0754\linewidth-0.5em\relax}{\raggedright\small\mbox{}21}\hspace{0.5em}\parbox[t]{\dimexpr0.0754\linewidth-0.5em\relax}{\raggedright\small\mbox{}180}\hspace{0.5em}\parbox[t]{\dimexpr0.0754\linewidth-0.5em\relax}{\raggedright\small\mbox{}512}\hspace{0.5em}\parbox[t]{\dimexpr0.0754\linewidth-0.5em\relax}{\raggedright\small\mbox{}6}\hspace{0.5em}\parbox[t]{\dimexpr0.0754\linewidth-0.5em\relax}{\raggedright\small\mbox{}28}\hspace{0.5em}\parbox[t]{\dimexpr0.0754\linewidth-0.5em\relax}{\raggedright\small\mbox{}0}\hspace{0.5em}}{}
\EdL{24}{AK07-J0020-L24}{\parbox[t]{\dimexpr0.0754\linewidth-0.5em\relax}{\raggedright\small\mbox{}210}\hspace{0.5em}\parbox[t]{\dimexpr0.0754\linewidth-0.5em\relax}{\raggedright\small\mbox{}193}\hspace{0.5em}\parbox[t]{\dimexpr0.0754\linewidth-0.5em\relax}{\raggedright\small\mbox{}11}\hspace{0.5em}\parbox[t]{\dimexpr0.0754\linewidth-0.5em\relax}{\raggedright\small\mbox{}18}\hspace{0.5em}\parbox[t]{\dimexpr0.0754\linewidth-0.5em\relax}{\raggedright\small\mbox{}8}\hspace{0.5em}\parbox[t]{\dimexpr0.0754\linewidth-0.5em\relax}{\raggedright\small\mbox{}7}\hspace{0.5em}\parbox[t]{\dimexpr0.0754\linewidth-0.5em\relax}{\raggedright\small\mbox{}9}\hspace{0.5em}\parbox[t]{\dimexpr0.0754\linewidth-0.5em\relax}{\raggedright\small\mbox{}10}\hspace{0.5em}\parbox[t]{\dimexpr0.0754\linewidth-0.5em\relax}{\raggedright\small\mbox{}210}\hspace{0.5em}\parbox[t]{\dimexpr0.0754\linewidth-0.5em\relax}{\raggedright\small\mbox{}541}\hspace{0.5em}\parbox[t]{\dimexpr0.0754\linewidth-0.5em\relax}{\raggedright\small\mbox{}7}\hspace{0.5em}\parbox[t]{\dimexpr0.0754\linewidth-0.5em\relax}{\raggedright\small\mbox{}7}\hspace{0.5em}\parbox[t]{\dimexpr0.0754\linewidth-0.5em\relax}{\raggedright\small\mbox{}\textsuperscript{3}/\textsubscript{4}}\hspace{0.5em}}{}
\EdL{25}{AK07-J0020-L25}{\parbox[t]{\dimexpr0.0754\linewidth-0.5em\relax}{\raggedright\small\mbox{}240}\hspace{0.5em}\parbox[t]{\dimexpr0.0754\linewidth-0.5em\relax}{\raggedright\small\mbox{}223}\hspace{0.5em}\parbox[t]{\dimexpr0.0754\linewidth-0.5em\relax}{\raggedright\small\mbox{}0}\hspace{0.5em}\parbox[t]{\dimexpr0.0754\linewidth-0.5em\relax}{\raggedright\small\mbox{}29}\hspace{0.5em}\parbox[t]{\dimexpr0.0754\linewidth-0.5em\relax}{\raggedright\small\mbox{}9}\hspace{0.5em}\parbox[t]{\dimexpr0.0754\linewidth-0.5em\relax}{\raggedright\small\mbox{}8}\hspace{0.5em}\parbox[t]{\dimexpr0.0754\linewidth-0.5em\relax}{\raggedright\small\mbox{}8}\hspace{0.5em}\parbox[t]{\dimexpr0.0754\linewidth-0.5em\relax}{\raggedright\small\mbox{}29}\hspace{0.5em}\parbox[t]{\dimexpr0.0754\linewidth-0.5em\relax}{\raggedright\small\mbox{}240}\hspace{0.5em}\parbox[t]{\dimexpr0.0754\linewidth-0.5em\relax}{\raggedright\small\mbox{}570}\hspace{0.5em}\parbox[t]{\dimexpr0.0754\linewidth-0.5em\relax}{\raggedright\small\mbox{}8}\hspace{0.5em}\parbox[t]{\dimexpr0.0754\linewidth-0.5em\relax}{\raggedright\small\mbox{}15}\hspace{0.5em}\parbox[t]{\dimexpr0.0754\linewidth-0.5em\relax}{\raggedright\small\mbox{}\textsuperscript{1}/\textsubscript{2}}\hspace{0.5em}}{}
\EdL{26}{AK07-J0020-L26}{\parbox[t]{\dimexpr0.0754\linewidth-0.5em\relax}{\raggedright\small\mbox{}270}\hspace{0.5em}\parbox[t]{\dimexpr0.0754\linewidth-0.5em\relax}{\raggedright\small\mbox{}252}\hspace{0.5em}\parbox[t]{\dimexpr0.0754\linewidth-0.5em\relax}{\raggedright\small\mbox{}2}\hspace{0.5em}\parbox[t]{\dimexpr0.0754\linewidth-0.5em\relax}{\raggedright\small\mbox{}15}\hspace{0.5em}\parbox[t]{\dimexpr0.0754\linewidth-0.5em\relax}{\raggedright\small\mbox{}10}\hspace{0.5em}\parbox[t]{\dimexpr0.0754\linewidth-0.5em\relax}{\raggedright\small\mbox{}9}\hspace{0.5em}\parbox[t]{\dimexpr0.0754\linewidth-0.5em\relax}{\raggedright\small\mbox{}8}\hspace{0.5em}\parbox[t]{\dimexpr0.0754\linewidth-0.5em\relax}{\raggedright\small\mbox{}19}\hspace{0.5em}\parbox[t]{\dimexpr0.0754\linewidth-0.5em\relax}{\raggedright\small\mbox{}270}\hspace{0.5em}\parbox[t]{\dimexpr0.0754\linewidth-0.5em\relax}{\raggedright\small\mbox{}599}\hspace{0.5em}\parbox[t]{\dimexpr0.0754\linewidth-0.5em\relax}{\raggedright\small\mbox{}10}\hspace{0.5em}\parbox[t]{\dimexpr0.0754\linewidth-0.5em\relax}{\raggedright\small\mbox{}24}\hspace{0.5em}\parbox[t]{\dimexpr0.0754\linewidth-0.5em\relax}{\raggedright\small\mbox{}\textsuperscript{1}/\textsubscript{4}}\hspace{0.5em}}{}
\EdL{27}{AK07-J0020-L27}{\parbox[t]{\dimexpr0.0754\linewidth-0.5em\relax}{\raggedright\small\mbox{}300}\hspace{0.5em}\parbox[t]{\dimexpr0.0754\linewidth-0.5em\relax}{\raggedright\small\mbox{}281}\hspace{0.5em}\parbox[t]{\dimexpr0.0754\linewidth-0.5em\relax}{\raggedright\small\mbox{}4}\hspace{0.5em}\parbox[t]{\dimexpr0.0754\linewidth-0.5em\relax}{\raggedright\small\mbox{}1}\hspace{0.5em}\parbox[t]{\dimexpr0.0754\linewidth-0.5em\relax}{\raggedright\small\mbox{}11}\hspace{0.5em}\parbox[t]{\dimexpr0.0754\linewidth-0.5em\relax}{\raggedright\small\mbox{}10}\hspace{0.5em}\parbox[t]{\dimexpr0.0754\linewidth-0.5em\relax}{\raggedright\small\mbox{}8}\hspace{0.5em}\parbox[t]{\dimexpr0.0754\linewidth-0.5em\relax}{\raggedright\small\mbox{}8}\hspace{0.5em}\parbox[t]{\dimexpr0.0754\linewidth-0.5em\relax}{\raggedright\small\mbox{}300}\hspace{0.5em}\parbox[t]{\dimexpr0.0754\linewidth-0.5em\relax}{\raggedright\small\mbox{}628}\hspace{0.5em}\parbox[t]{\dimexpr0.0754\linewidth-0.5em\relax}{\raggedright\small\mbox{}11}\hspace{0.5em}\parbox[t]{\dimexpr0.0754\linewidth-0.5em\relax}{\raggedright\small\mbox{}3}\hspace{0.5em}\parbox[t]{\dimexpr0.0754\linewidth-0.5em\relax}{\raggedright\small\mbox{}0}\hspace{0.5em}}{}
\EdL{28}{AK07-J0020-L28}{\parbox[t]{\dimexpr0.0754\linewidth-0.5em\relax}{\raggedright\small\mbox{}330}\hspace{0.5em}\parbox[t]{\dimexpr0.0754\linewidth-0.5em\relax}{\raggedright\small\mbox{}310}\hspace{0.5em}\parbox[t]{\dimexpr0.0754\linewidth-0.5em\relax}{\raggedright\small\mbox{}5}\hspace{0.5em}\parbox[t]{\dimexpr0.0754\linewidth-0.5em\relax}{\raggedright\small\mbox{}17}\hspace{0.5em}\parbox[t]{\dimexpr0.0754\linewidth-0.5em\relax}{\raggedright\small\mbox{}12}\hspace{0.5em}\parbox[t]{\dimexpr0.0754\linewidth-0.5em\relax}{\raggedright\small\mbox{}11}\hspace{0.5em}\parbox[t]{\dimexpr0.0754\linewidth-0.5em\relax}{\raggedright\small\mbox{}7}\hspace{0.5em}\parbox[t]{\dimexpr0.0754\linewidth-0.5em\relax}{\raggedright\small\mbox{}27}\hspace{0.5em}\parbox[t]{\dimexpr0.0754\linewidth-0.5em\relax}{\raggedright\small\mbox{}330}\hspace{0.5em}\parbox[t]{\dimexpr0.0754\linewidth-0.5em\relax}{\raggedright\small\mbox{}658}\hspace{0.5em}\parbox[t]{\dimexpr0.0754\linewidth-0.5em\relax}{\raggedright\small\mbox{}0}\hspace{0.5em}\parbox[t]{\dimexpr0.0754\linewidth-0.5em\relax}{\raggedright\small\mbox{}12}\hspace{0.5em}\parbox[t]{\dimexpr0.0754\linewidth-0.5em\relax}{\raggedright\small\mbox{}\textsuperscript{1}/\textsubscript{2}}\hspace{0.5em}}{}
\EdL{29}{AK07-J0020-L29}{\parbox[t]{\dimexpr0.0754\linewidth-0.5em\relax}{\raggedright\small\mbox{}360}\hspace{0.5em}\parbox[t]{\dimexpr0.0754\linewidth-0.5em\relax}{\raggedright\small\mbox{}339}\hspace{0.5em}\parbox[t]{\dimexpr0.0754\linewidth-0.5em\relax}{\raggedright\small\mbox{}7}\hspace{0.5em}\parbox[t]{\dimexpr0.0754\linewidth-0.5em\relax}{\raggedright\small\mbox{}3}\hspace{0.5em}\parbox[t]{\dimexpr0.0754\linewidth-0.5em\relax}{\raggedright\small\mbox{}13}\hspace{0.5em}\parbox[t]{\dimexpr0.0754\linewidth-0.5em\relax}{\raggedright\small\mbox{}12}\hspace{0.5em}\parbox[t]{\dimexpr0.0754\linewidth-0.5em\relax}{\raggedright\small\mbox{}7}\hspace{0.5em}\parbox[t]{\dimexpr0.0754\linewidth-0.5em\relax}{\raggedright\small\mbox{}17}\hspace{0.5em}\parbox[t]{\dimexpr0.0754\linewidth-0.5em\relax}{\raggedright\small\mbox{}360}\hspace{0.5em}\parbox[t]{\dimexpr0.0754\linewidth-0.5em\relax}{\raggedright\small\mbox{}687}\hspace{0.5em}\parbox[t]{\dimexpr0.0754\linewidth-0.5em\relax}{\raggedright\small\mbox{}1}\hspace{0.5em}\parbox[t]{\dimexpr0.0754\linewidth-0.5em\relax}{\raggedright\small\mbox{}15}\hspace{0.5em}\parbox[t]{\dimexpr0.0754\linewidth-0.5em\relax}{\raggedright\small\mbox{}\textsuperscript{1}/\textsubscript{4}}\hspace{0.5em}}{}
\EdL{30}{AK07-J0020-L30}{\parbox[t]{\dimexpr0.0754\linewidth-0.5em\relax}{\raggedright\small\mbox{}390}\hspace{0.5em}\parbox[t]{\dimexpr0.0754\linewidth-0.5em\relax}{\raggedright\small\mbox{}368}\hspace{0.5em}\parbox[t]{\dimexpr0.0754\linewidth-0.5em\relax}{\raggedright\small\mbox{}8}\hspace{0.5em}\parbox[t]{\dimexpr0.0754\linewidth-0.5em\relax}{\raggedright\small\mbox{}19}\hspace{0.5em}\parbox[t]{\dimexpr0.0754\linewidth-0.5em\relax}{\raggedright\small\mbox{}14}\hspace{0.5em}\parbox[t]{\dimexpr0.0754\linewidth-0.5em\relax}{\raggedright\small\mbox{}13}\hspace{0.5em}\parbox[t]{\dimexpr0.0754\linewidth-0.5em\relax}{\raggedright\small\mbox{}7}\hspace{0.5em}\parbox[t]{\dimexpr0.0754\linewidth-0.5em\relax}{\raggedright\small\mbox{}6}\hspace{0.5em}\parbox[t]{\dimexpr0.0754\linewidth-0.5em\relax}{\raggedright\small\mbox{}390}\hspace{0.5em}\parbox[t]{\dimexpr0.0754\linewidth-0.5em\relax}{\raggedright\small\mbox{}716}\hspace{0.5em}\parbox[t]{\dimexpr0.0754\linewidth-0.5em\relax}{\raggedright\small\mbox{}3}\hspace{0.5em}\parbox[t]{\dimexpr0.0754\linewidth-0.5em\relax}{\raggedright\small\mbox{}24}\hspace{0.5em}\parbox[t]{\dimexpr0.0754\linewidth-0.5em\relax}{\raggedright\small\mbox{}0}\hspace{0.5em}}{}
\EdL{31}{AK07-J0020-L31}{\parbox[t]{\dimexpr0.0754\linewidth-0.5em\relax}{\raggedright\small\mbox{}420}\hspace{0.5em}\parbox[t]{\dimexpr0.0754\linewidth-0.5em\relax}{\raggedright\small\mbox{}397}\hspace{0.5em}\parbox[t]{\dimexpr0.0754\linewidth-0.5em\relax}{\raggedright\small\mbox{}10}\hspace{0.5em}\parbox[t]{\dimexpr0.0754\linewidth-0.5em\relax}{\raggedright\small\mbox{}5}\hspace{0.5em}\parbox[t]{\dimexpr0.0754\linewidth-0.5em\relax}{\raggedright\small\mbox{}15}\hspace{0.5em}\parbox[t]{\dimexpr0.0754\linewidth-0.5em\relax}{\raggedright\small\mbox{}14}\hspace{0.5em}\parbox[t]{\dimexpr0.0754\linewidth-0.5em\relax}{\raggedright\small\mbox{}6}\hspace{0.5em}\parbox[t]{\dimexpr0.0754\linewidth-0.5em\relax}{\raggedright\small\mbox{}25}\hspace{0.5em}\parbox[t]{\dimexpr0.0754\linewidth-0.5em\relax}{\raggedright\small\mbox{}420}\hspace{0.5em}\parbox[t]{\dimexpr0.0754\linewidth-0.5em\relax}{\raggedright\small\mbox{}745}\hspace{0.5em}\parbox[t]{\dimexpr0.0754\linewidth-0.5em\relax}{\raggedright\small\mbox{}4}\hspace{0.5em}\parbox[t]{\dimexpr0.0754\linewidth-0.5em\relax}{\raggedright\small\mbox{}3}\hspace{0.5em}\parbox[t]{\dimexpr0.0754\linewidth-0.5em\relax}{\raggedright\small\mbox{}\textsuperscript{3}/\textsubscript{4}}\hspace{0.5em}}{}
\EdL{32}{AK07-J0020-L32}{\parbox[t]{\dimexpr0.0754\linewidth-0.5em\relax}{\raggedright\small\mbox{}450}\hspace{0.5em}\parbox[t]{\dimexpr0.0754\linewidth-0.5em\relax}{\raggedright\small\mbox{}426}\hspace{0.5em}\parbox[t]{\dimexpr0.0754\linewidth-0.5em\relax}{\raggedright\small\mbox{}11}\hspace{0.5em}\parbox[t]{\dimexpr0.0754\linewidth-0.5em\relax}{\raggedright\small\mbox{}21}\hspace{0.5em}\parbox[t]{\dimexpr0.0754\linewidth-0.5em\relax}{\raggedright\small\mbox{}16}\hspace{0.5em}\parbox[t]{\dimexpr0.0754\linewidth-0.5em\relax}{\raggedright\small\mbox{}15}\hspace{0.5em}\parbox[t]{\dimexpr0.0754\linewidth-0.5em\relax}{\raggedright\small\mbox{}6}\hspace{0.5em}\parbox[t]{\dimexpr0.0754\linewidth-0.5em\relax}{\raggedright\small\mbox{}15}\hspace{0.5em}\parbox[t]{\dimexpr0.0754\linewidth-0.5em\relax}{\raggedright\small\mbox{}450}\hspace{0.5em}\parbox[t]{\dimexpr0.0754\linewidth-0.5em\relax}{\raggedright\small\mbox{}774}\hspace{0.5em}\parbox[t]{\dimexpr0.0754\linewidth-0.5em\relax}{\raggedright\small\mbox{}5}\hspace{0.5em}\parbox[t]{\dimexpr0.0754\linewidth-0.5em\relax}{\raggedright\small\mbox{}11}\hspace{0.5em}\parbox[t]{\dimexpr0.0754\linewidth-0.5em\relax}{\raggedright\small\mbox{}\textsuperscript{1}/\textsubscript{2}}\hspace{0.5em}}{}
\EdL{33}{AK07-J0020-L33}{\parbox[t]{\dimexpr0.0754\linewidth-0.5em\relax}{\raggedright\small\mbox{}480}\hspace{0.5em}\parbox[t]{\dimexpr0.0754\linewidth-0.5em\relax}{\raggedright\small\mbox{}456}\hspace{0.5em}\parbox[t]{\dimexpr0.0754\linewidth-0.5em\relax}{\raggedright\small\mbox{}1}\hspace{0.5em}\parbox[t]{\dimexpr0.0754\linewidth-0.5em\relax}{\raggedright\small\mbox{}2}\hspace{0.5em}\parbox[t]{\dimexpr0.0754\linewidth-0.5em\relax}{\raggedright\small\mbox{}17}\hspace{0.5em}\parbox[t]{\dimexpr0.0754\linewidth-0.5em\relax}{\raggedright\small\mbox{}16}\hspace{0.5em}\parbox[t]{\dimexpr0.0754\linewidth-0.5em\relax}{\raggedright\small\mbox{}6}\hspace{0.5em}\parbox[t]{\dimexpr0.0754\linewidth-0.5em\relax}{\raggedright\small\mbox{}4}\hspace{0.5em}\parbox[t]{\dimexpr0.0754\linewidth-0.5em\relax}{\raggedright\small\mbox{}480}\hspace{0.5em}\parbox[t]{\dimexpr0.0754\linewidth-0.5em\relax}{\raggedright\small\mbox{}803}\hspace{0.5em}\parbox[t]{\dimexpr0.0754\linewidth-0.5em\relax}{\raggedright\small\mbox{}6}\hspace{0.5em}\parbox[t]{\dimexpr0.0754\linewidth-0.5em\relax}{\raggedright\small\mbox{}20}\hspace{0.5em}\parbox[t]{\dimexpr0.0754\linewidth-0.5em\relax}{\raggedright\small\mbox{}\textsuperscript{1}/\textsubscript{4}}\hspace{0.5em}}{}
\EdL{34}{AK07-J0020-L34}{\parbox[t]{\dimexpr0.0754\linewidth-0.5em\relax}{\raggedright\small\mbox{}510}\hspace{0.5em}\parbox[t]{\dimexpr0.0754\linewidth-0.5em\relax}{\raggedright\small\mbox{}485}\hspace{0.5em}\parbox[t]{\dimexpr0.0754\linewidth-0.5em\relax}{\raggedright\small\mbox{}2}\hspace{0.5em}\parbox[t]{\dimexpr0.0754\linewidth-0.5em\relax}{\raggedright\small\mbox{}18}\hspace{0.5em}\parbox[t]{\dimexpr0.0754\linewidth-0.5em\relax}{\raggedright\small\mbox{}18}\hspace{0.5em}\parbox[t]{\dimexpr0.0754\linewidth-0.5em\relax}{\raggedright\small\mbox{}17}\hspace{0.5em}\parbox[t]{\dimexpr0.0754\linewidth-0.5em\relax}{\raggedright\small\mbox{}5}\hspace{0.5em}\parbox[t]{\dimexpr0.0754\linewidth-0.5em\relax}{\raggedright\small\mbox{}24}\hspace{0.5em}\parbox[t]{\dimexpr0.0754\linewidth-0.5em\relax}{\raggedright\small\mbox{}510}\hspace{0.5em}\parbox[t]{\dimexpr0.0754\linewidth-0.5em\relax}{\raggedright\small\mbox{}832}\hspace{0.5em}\parbox[t]{\dimexpr0.0754\linewidth-0.5em\relax}{\raggedright\small\mbox{}8}\hspace{0.5em}\parbox[t]{\dimexpr0.0754\linewidth-0.5em\relax}{\raggedright\small\mbox{}28}\hspace{0.5em}\parbox[t]{\dimexpr0.0754\linewidth-0.5em\relax}{\raggedright\small\mbox{}0}\hspace{0.5em}}{}
\EdL{35}{AK07-J0020-L35}{\parbox[t]{\dimexpr0.0754\linewidth-0.5em\relax}{\raggedright\small\mbox{}540}\hspace{0.5em}\parbox[t]{\dimexpr0.0754\linewidth-0.5em\relax}{\raggedright\small\mbox{}514}\hspace{0.5em}\parbox[t]{\dimexpr0.0754\linewidth-0.5em\relax}{\raggedright\small\mbox{}4}\hspace{0.5em}\parbox[t]{\dimexpr0.0754\linewidth-0.5em\relax}{\raggedright\small\mbox{}4}\hspace{0.5em}\parbox[t]{\dimexpr0.0754\linewidth-0.5em\relax}{\raggedright\small\mbox{}19}\hspace{0.5em}\parbox[t]{\dimexpr0.0754\linewidth-0.5em\relax}{\raggedright\small\mbox{}18}\hspace{0.5em}\parbox[t]{\dimexpr0.0754\linewidth-0.5em\relax}{\raggedright\small\mbox{}5}\hspace{0.5em}\parbox[t]{\dimexpr0.0754\linewidth-0.5em\relax}{\raggedright\small\mbox{}13}\hspace{0.5em}\parbox[t]{\dimexpr0.0754\linewidth-0.5em\relax}{\raggedright\small\mbox{}540}\hspace{0.5em}\parbox[t]{\dimexpr0.0754\linewidth-0.5em\relax}{\raggedright\small\mbox{}861}\hspace{0.5em}\parbox[t]{\dimexpr0.0754\linewidth-0.5em\relax}{\raggedright\small\mbox{}9}\hspace{0.5em}\parbox[t]{\dimexpr0.0754\linewidth-0.5em\relax}{\raggedright\small\mbox{}8}\hspace{0.5em}\parbox[t]{\dimexpr0.0754\linewidth-0.5em\relax}{\raggedright\small\mbox{}\textsuperscript{3}/\textsubscript{4}}\hspace{0.5em}}{}
\EdL{36}{AK07-J0020-L36}{\parbox[t]{\dimexpr0.0754\linewidth-0.5em\relax}{\raggedright\small\mbox{}570}\hspace{0.5em}\parbox[t]{\dimexpr0.0754\linewidth-0.5em\relax}{\raggedright\small\mbox{}543}\hspace{0.5em}\parbox[t]{\dimexpr0.0754\linewidth-0.5em\relax}{\raggedright\small\mbox{}5}\hspace{0.5em}\parbox[t]{\dimexpr0.0754\linewidth-0.5em\relax}{\raggedright\small\mbox{}20}\hspace{0.5em}\parbox[t]{\dimexpr0.0754\linewidth-0.5em\relax}{\raggedright\small\mbox{}20}\hspace{0.5em}\parbox[t]{\dimexpr0.0754\linewidth-0.5em\relax}{\raggedright\small\mbox{}19}\hspace{0.5em}\parbox[t]{\dimexpr0.0754\linewidth-0.5em\relax}{\raggedright\small\mbox{}5}\hspace{0.5em}\parbox[t]{\dimexpr0.0754\linewidth-0.5em\relax}{\raggedright\small\mbox{}2}\hspace{0.5em}\parbox[t]{\dimexpr0.0754\linewidth-0.5em\relax}{\raggedright\small\mbox{}570}\hspace{0.5em}\parbox[t]{\dimexpr0.0754\linewidth-0.5em\relax}{\raggedright\small\mbox{}890}\hspace{0.5em}\parbox[t]{\dimexpr0.0754\linewidth-0.5em\relax}{\raggedright\small\mbox{}10}\hspace{0.5em}\parbox[t]{\dimexpr0.0754\linewidth-0.5em\relax}{\raggedright\small\mbox{}16}\hspace{0.5em}\parbox[t]{\dimexpr0.0754\linewidth-0.5em\relax}{\raggedright\small\mbox{}\textsuperscript{1}/\textsubscript{2}}\hspace{0.5em}}{}
\EdL{37}{AK07-J0020-L37}{\parbox[t]{\dimexpr0.0754\linewidth-0.5em\relax}{\raggedright\small\mbox{}600}\hspace{0.5em}\parbox[t]{\dimexpr0.0754\linewidth-0.5em\relax}{\raggedright\small\mbox{}572}\hspace{0.5em}\parbox[t]{\dimexpr0.0754\linewidth-0.5em\relax}{\raggedright\small\mbox{}7}\hspace{0.5em}\parbox[t]{\dimexpr0.0754\linewidth-0.5em\relax}{\raggedright\small\mbox{}6}\hspace{0.5em}\parbox[t]{\dimexpr0.0754\linewidth-0.5em\relax}{\raggedright\small\mbox{}21}\hspace{0.5em}\parbox[t]{\dimexpr0.0754\linewidth-0.5em\relax}{\raggedright\small\mbox{}20}\hspace{0.5em}\parbox[t]{\dimexpr0.0754\linewidth-0.5em\relax}{\raggedright\small\mbox{}4}\hspace{0.5em}\parbox[t]{\dimexpr0.0754\linewidth-0.5em\relax}{\raggedright\small\mbox{}22}\hspace{0.5em}\parbox[t]{\dimexpr0.0754\linewidth-0.5em\relax}{\raggedright\small\mbox{}600}\hspace{0.5em}\parbox[t]{\dimexpr0.0754\linewidth-0.5em\relax}{\raggedright\small\mbox{}919}\hspace{0.5em}\parbox[t]{\dimexpr0.0754\linewidth-0.5em\relax}{\raggedright\small\mbox{}12}\hspace{0.5em}\parbox[t]{\dimexpr0.0754\linewidth-0.5em\relax}{\raggedright\small\mbox{}25}\hspace{0.5em}\parbox[t]{\dimexpr0.0754\linewidth-0.5em\relax}{\raggedright\small\mbox{}\textsuperscript{1}/\textsubscript{4}}\hspace{0.5em}}{}
\EdL{38}{AK07-J0020-L38}{\parbox[t]{\dimexpr0.0754\linewidth-0.5em\relax}{\raggedright\small\mbox{}630}\hspace{0.5em}\parbox[t]{\dimexpr0.0754\linewidth-0.5em\relax}{\raggedright\small\mbox{}601}\hspace{0.5em}\parbox[t]{\dimexpr0.0754\linewidth-0.5em\relax}{\raggedright\small\mbox{}8}\hspace{0.5em}\parbox[t]{\dimexpr0.0754\linewidth-0.5em\relax}{\raggedright\small\mbox{}22}\hspace{0.5em}\parbox[t]{\dimexpr0.0754\linewidth-0.5em\relax}{\raggedright\small\mbox{}22}\hspace{0.5em}\parbox[t]{\dimexpr0.0754\linewidth-0.5em\relax}{\raggedright\small\mbox{}21}\hspace{0.5em}\parbox[t]{\dimexpr0.0754\linewidth-0.5em\relax}{\raggedright\small\mbox{}4}\hspace{0.5em}\parbox[t]{\dimexpr0.0754\linewidth-0.5em\relax}{\raggedright\small\mbox{}11}\hspace{0.5em}\parbox[t]{\dimexpr0.0754\linewidth-0.5em\relax}{\raggedright\small\mbox{}630}\hspace{0.5em}\parbox[t]{\dimexpr0.0754\linewidth-0.5em\relax}{\raggedright\small\mbox{}948}\hspace{0.5em}\parbox[t]{\dimexpr0.0754\linewidth-0.5em\relax}{\raggedright\small\mbox{}13}\hspace{0.5em}\parbox[t]{\dimexpr0.0754\linewidth-0.5em\relax}{\raggedright\small\mbox{}4}\hspace{0.5em}\parbox[t]{\dimexpr0.0754\linewidth-0.5em\relax}{\raggedright\small\mbox{}0}\hspace{0.5em}}{}
\EdL{39}{AK07-J0020-L39}{\parbox[t]{\dimexpr0.0754\linewidth-0.5em\relax}{\raggedright\small\mbox{}660}\hspace{0.5em}\parbox[t]{\dimexpr0.0754\linewidth-0.5em\relax}{\raggedright\small\mbox{}630}\hspace{0.5em}\parbox[t]{\dimexpr0.0754\linewidth-0.5em\relax}{\raggedright\small\mbox{}10}\hspace{0.5em}\parbox[t]{\dimexpr0.0754\linewidth-0.5em\relax}{\raggedright\small\mbox{}8}\hspace{0.5em}\parbox[t]{\dimexpr0.0754\linewidth-0.5em\relax}{\raggedright\small\mbox{}23}\hspace{0.5em}\parbox[t]{\dimexpr0.0754\linewidth-0.5em\relax}{\raggedright\small\mbox{}22}\hspace{0.5em}\parbox[t]{\dimexpr0.0754\linewidth-0.5em\relax}{\raggedright\small\mbox{}4}\hspace{0.5em}\parbox[t]{\dimexpr0.0754\linewidth-0.5em\relax}{\raggedright\small\mbox{}0}\hspace{0.5em}\parbox[t]{\dimexpr0.0754\linewidth-0.5em\relax}{\raggedright\small\mbox{}660}\hspace{0.5em}\parbox[t]{\dimexpr0.0754\linewidth-0.5em\relax}{\raggedright\small\mbox{}997\textsuperscript{b})}\hspace{0.5em}\parbox[t]{\dimexpr0.0754\linewidth-0.5em\relax}{\raggedright\small\mbox{}14}\hspace{0.5em}\parbox[t]{\dimexpr0.0754\linewidth-0.5em\relax}{\raggedright\small\mbox{}13}\hspace{0.5em}\parbox[t]{\dimexpr0.0754\linewidth-0.5em\relax}{\raggedright\small\mbox{}\textsuperscript{3}/\textsubscript{4}}\hspace{0.5em}}{}
\EdL{40}{AK07-J0020-L40}{\parbox[t]{\dimexpr0.0754\linewidth-0.5em\relax}{\raggedright\small\mbox{}690}\hspace{0.5em}\parbox[t]{\dimexpr0.0754\linewidth-0.5em\relax}{\raggedright\small\mbox{}659}\hspace{0.5em}\parbox[t]{\dimexpr0.0754\linewidth-0.5em\relax}{\raggedright\small\mbox{}11}\hspace{0.5em}\parbox[t]{\dimexpr0.0754\linewidth-0.5em\relax}{\raggedright\small\mbox{}24}\hspace{0.5em}\parbox[t]{\dimexpr0.0754\linewidth-0.5em\relax}{\raggedright\small\mbox{}24}\hspace{0.5em}\parbox[t]{\dimexpr0.0754\linewidth-0.5em\relax}{\raggedright\small\mbox{}23}\hspace{0.5em}\parbox[t]{\dimexpr0.0754\linewidth-0.5em\relax}{\raggedright\small\mbox{}3}\hspace{0.5em}\parbox[t]{\dimexpr0.0754\linewidth-0.5em\relax}{\raggedright\small\mbox{}20}\hspace{0.5em}\parbox[t]{\dimexpr0.0754\linewidth-0.5em\relax}{\raggedright\small\mbox{}690}\hspace{0.5em}\parbox[t]{\dimexpr0.0754\linewidth-0.5em\relax}{\raggedright\small\mbox{}1006}\hspace{0.5em}\parbox[t]{\dimexpr0.0754\linewidth-0.5em\relax}{\raggedright\small\mbox{}16}\hspace{0.5em}\parbox[t]{\dimexpr0.0754\linewidth-0.5em\relax}{\raggedright\small\mbox{}21}\hspace{0.5em}\parbox[t]{\dimexpr0.0754\linewidth-0.5em\relax}{\raggedright\small\mbox{}\textsuperscript{1}/\textsubscript{2}}\hspace{0.5em}}{}
\EdL{41}{AK07-J0020-L41}{\parbox[t]{\dimexpr0.3015\linewidth-0.5em\relax}{\raggedright\small\mbox{}}\hspace{0.5em}\parbox[t]{\dimexpr0.0754\linewidth-0.5em\relax}{\raggedright\small\mbox{}25}\hspace{0.5em}\parbox[t]{\dimexpr0.0754\linewidth-0.5em\relax}{\raggedright\small\mbox{}24}\hspace{0.5em}\parbox[t]{\dimexpr0.0754\linewidth-0.5em\relax}{\raggedright\small\mbox{}3}\hspace{0.5em}\parbox[t]{\dimexpr0.4523\linewidth-0.5em\relax}{\raggedright\small\mbox{}9}\hspace{0.5em}}{}
\EdL{42}{AK07-J0020-L42}{\parbox[t]{\dimexpr0.3015\linewidth-0.5em\relax}{\raggedright\small\mbox{}}\hspace{0.5em}\parbox[t]{\dimexpr0.0754\linewidth-0.5em\relax}{\raggedright\small\mbox{}26}\hspace{0.5em}\parbox[t]{\dimexpr0.0754\linewidth-0.5em\relax}{\raggedright\small\mbox{}25}\hspace{0.5em}\parbox[t]{\dimexpr0.0754\linewidth-0.5em\relax}{\raggedright\small\mbox{}2}\hspace{0.5em}\parbox[t]{\dimexpr0.4523\linewidth-0.5em\relax}{\raggedright\small\mbox{}29}\hspace{0.5em}}{}
\EdL{43}{AK07-J0020-L43}{\parbox[t]{\dimexpr0.3015\linewidth-0.5em\relax}{\raggedright\small\mbox{}}\hspace{0.5em}\parbox[t]{\dimexpr0.0754\linewidth-0.5em\relax}{\raggedright\small\mbox{}27}\hspace{0.5em}\parbox[t]{\dimexpr0.0754\linewidth-0.5em\relax}{\raggedright\small\mbox{}26}\hspace{0.5em}\parbox[t]{\dimexpr0.0754\linewidth-0.5em\relax}{\raggedright\small\mbox{}2}\hspace{0.5em}\parbox[t]{\dimexpr0.4523\linewidth-0.5em\relax}{\raggedright\small\mbox{}19}\hspace{0.5em}}{}
\EdL{44}{AK07-J0020-L44}{\parbox[t]{\dimexpr0.3015\linewidth-0.5em\relax}{\raggedright\small\mbox{}}\hspace{0.5em}\parbox[t]{\dimexpr0.0754\linewidth-0.5em\relax}{\raggedright\small\mbox{}28}\hspace{0.5em}\parbox[t]{\dimexpr0.0754\linewidth-0.5em\relax}{\raggedright\small\mbox{}27}\hspace{0.5em}\parbox[t]{\dimexpr0.0754\linewidth-0.5em\relax}{\raggedright\small\mbox{}2}\hspace{0.5em}\parbox[t]{\dimexpr0.4523\linewidth-0.5em\relax}{\raggedright\small\mbox{}7}\hspace{0.5em}}{}
\EdL{45}{AK07-J0020-L45}{\parbox[t]{\dimexpr0.3015\linewidth-0.5em\relax}{\raggedright\small\mbox{}}\hspace{0.5em}\parbox[t]{\dimexpr0.0754\linewidth-0.5em\relax}{\raggedright\small\mbox{}29}\hspace{0.5em}\parbox[t]{\dimexpr0.0754\linewidth-0.5em\relax}{\raggedright\small\mbox{}28}\hspace{0.5em}\parbox[t]{\dimexpr0.0754\linewidth-0.5em\relax}{\raggedright\small\mbox{}1}\hspace{0.5em}\parbox[t]{\dimexpr0.4523\linewidth-0.5em\relax}{\raggedright\small\mbox{}27}\hspace{0.5em}}{}
\EdL{46}{AK07-J0020-L46}{\parbox[t]{\dimexpr0.3015\linewidth-0.5em\relax}{\raggedright\small\mbox{}}\hspace{0.5em}\parbox[t]{\dimexpr0.0754\linewidth-0.5em\relax}{\raggedright\small\mbox{}30}\hspace{0.5em}\parbox[t]{\dimexpr0.0754\linewidth-0.5em\relax}{\raggedright\small\mbox{}29}\hspace{0.5em}\parbox[t]{\dimexpr0.0754\linewidth-0.5em\relax}{\raggedright\small\mbox{}1}\hspace{0.5em}\parbox[t]{\dimexpr0.4523\linewidth-0.5em\relax}{\raggedright\small\mbox{}16}\hspace{0.5em}}{}
\NoteRule
\EdN{N1}{AK07-J0020-N01}{\textsuperscript{1}) Vergl. \textit{Al-Battānī}, ed. \Sperr{Nallino}, I, 244—245.}
\FrameHead{0021}{p. XVII}{0021}{p. XVII}
\EdLp{1}{AK07-J0021-L01}{4) Am Fusse von fol. 9\textsuperscript{r} findet sich eine kleine Tafel, die die Verwandlung der}{}
\EdL{2}{AK07-J0021-L02}{\textit{arabischen} Jahre in die \textit{christlichen} ermöglicht, aber nur für die \textit{anni collecti} von 30}{}
\EdL{3}{AK07-J0021-L03}{bis 660\,; sie ist also unvollständig, vielleicht ist sie ein Überrest der im Kommentar}{}
\EdL{4}{AK07-J0021-L04}{p. 36 erwähnten, in den übrigen Handschriften fehlenden Tafel, sie ist allerdings nur}{}
\EdL{5}{AK07-J0021-L05}{ein Scholium einer spätern Hand, ich gebe sie hier aber ebenfalls\,:}{}
\EdL{6}{AK07-J0021-L06}{\parbox[t]{\dimexpr0.1633\linewidth-0.5em\relax}{\raggedright\small\mbox{}Notae}\hspace{0.5em}\parbox[t]{\dimexpr0.1633\linewidth-0.5em\relax}{\raggedright\small\mbox{}Anni arab.}\hspace{0.5em}\parbox[t]{\dimexpr0.1633\linewidth-0.5em\relax}{\raggedright\small\mbox{}Anni dom. collecti}\hspace{0.5em}\parbox[t]{\dimexpr0.1633\linewidth-0.5em\relax}{\raggedright\small\mbox{}Menses}\hspace{0.5em}\parbox[t]{\dimexpr0.1633\linewidth-0.5em\relax}{\raggedright\small\mbox{}Dies}\hspace{0.5em}\parbox[t]{\dimexpr0.1633\linewidth-0.5em\relax}{\raggedright\small\mbox{}Fract.}\hspace{0.5em}}{}
\EdL{7}{AK07-J0021-L07}{\parbox[t]{\dimexpr0.1633\linewidth-0.5em\relax}{\raggedright\small\mbox{}5}\hspace{0.5em}\parbox[t]{\dimexpr0.1633\linewidth-0.5em\relax}{\raggedright\small\mbox{}30}\hspace{0.5em}\parbox[t]{\dimexpr0.1633\linewidth-0.5em\relax}{\raggedright\small\mbox{}650}\hspace{0.5em}\parbox[t]{\dimexpr0.1633\linewidth-0.5em\relax}{\raggedright\small\mbox{}7}\hspace{0.5em}\parbox[t]{\dimexpr0.1633\linewidth-0.5em\relax}{\raggedright\small\mbox{}23}\hspace{0.5em}\parbox[t]{\dimexpr0.1633\linewidth-0.5em\relax}{\raggedright\small\mbox{}3}\hspace{0.5em}}{}
\EdL{8}{AK07-J0021-L08}{\parbox[t]{\dimexpr0.1633\linewidth-0.5em\relax}{\raggedright\small\mbox{}4}\hspace{0.5em}\parbox[t]{\dimexpr0.1633\linewidth-0.5em\relax}{\raggedright\small\mbox{}60}\hspace{0.5em}\parbox[t]{\dimexpr0.1633\linewidth-0.5em\relax}{\raggedright\small\mbox{}679}\hspace{0.5em}\parbox[t]{\dimexpr0.1633\linewidth-0.5em\relax}{\raggedright\small\mbox{}9}\hspace{0.5em}\parbox[t]{\dimexpr0.1633\linewidth-0.5em\relax}{\raggedright\small\mbox{}2}\hspace{0.5em}\parbox[t]{\dimexpr0.1633\linewidth-0.5em\relax}{\raggedright\small\mbox{}2}\hspace{0.5em}}{}
\EdL{9}{AK07-J0021-L09}{\parbox[t]{\dimexpr0.1633\linewidth-0.5em\relax}{\raggedright\small\mbox{}1}\hspace{0.5em}\parbox[t]{\dimexpr0.1633\linewidth-0.5em\relax}{\raggedright\small\mbox{}90}\hspace{0.5em}\parbox[t]{\dimexpr0.1633\linewidth-0.5em\relax}{\raggedright\small\mbox{}708}\hspace{0.5em}\parbox[t]{\dimexpr0.1633\linewidth-0.5em\relax}{\raggedright\small\mbox{}10}\hspace{0.5em}\parbox[t]{\dimexpr0.1633\linewidth-0.5em\relax}{\raggedright\small\mbox{}11}\hspace{0.5em}\parbox[t]{\dimexpr0.1633\linewidth-0.5em\relax}{\raggedright\small\mbox{}1}\hspace{0.5em}}{}
\EdL{10}{AK07-J0021-L10}{\parbox[t]{\dimexpr0.1633\linewidth-0.5em\relax}{\raggedright\small\mbox{}5}\hspace{0.5em}\parbox[t]{\dimexpr0.1633\linewidth-0.5em\relax}{\raggedright\small\mbox{}120}\hspace{0.5em}\parbox[t]{\dimexpr0.1633\linewidth-0.5em\relax}{\raggedright\small\mbox{}737}\hspace{0.5em}\parbox[t]{\dimexpr0.1633\linewidth-0.5em\relax}{\raggedright\small\mbox{}11}\hspace{0.5em}\parbox[t]{\dimexpr0.1633\linewidth-0.5em\relax}{\raggedright\small\mbox{}20}\hspace{0.5em}\parbox[t]{\dimexpr0.1633\linewidth-0.5em\relax}{\raggedright\small\mbox{}0}\hspace{0.5em}}{}
\EdL{11}{AK07-J0021-L11}{\parbox[t]{\dimexpr0.1633\linewidth-0.5em\relax}{\raggedright\small\mbox{}4}\hspace{0.5em}\parbox[t]{\dimexpr0.1633\linewidth-0.5em\relax}{\raggedright\small\mbox{}150}\hspace{0.5em}\parbox[t]{\dimexpr0.1633\linewidth-0.5em\relax}{\raggedright\small\mbox{}767}\hspace{0.5em}\parbox[t]{\dimexpr0.1633\linewidth-0.5em\relax}{\raggedright\small\mbox{}0}\hspace{0.5em}\parbox[t]{\dimexpr0.1633\linewidth-0.5em\relax}{\raggedright\small\mbox{}28}\hspace{0.5em}\parbox[t]{\dimexpr0.1633\linewidth-0.5em\relax}{\raggedright\small\mbox{}2}\hspace{0.5em}}{}
\EdL{12}{AK07-J0021-L12}{\parbox[t]{\dimexpr0.1633\linewidth-0.5em\relax}{\raggedright\small\mbox{}2}\hspace{0.5em}\parbox[t]{\dimexpr0.1633\linewidth-0.5em\relax}{\raggedright\small\mbox{}180}\hspace{0.5em}\parbox[t]{\dimexpr0.1633\linewidth-0.5em\relax}{\raggedright\small\mbox{}796}\hspace{0.5em}\parbox[t]{\dimexpr0.1633\linewidth-0.5em\relax}{\raggedright\small\mbox{}2}\hspace{0.5em}\parbox[t]{\dimexpr0.1633\linewidth-0.5em\relax}{\raggedright\small\mbox{}2}\hspace{0.5em}\parbox[t]{\dimexpr0.1633\linewidth-0.5em\relax}{\raggedright\small\mbox{}1}\hspace{0.5em}}{}
\EdL{13}{AK07-J0021-L13}{\parbox[t]{\dimexpr0.1633\linewidth-0.5em\relax}{\raggedright\small\mbox{}7}\hspace{0.5em}\parbox[t]{\dimexpr0.1633\linewidth-0.5em\relax}{\raggedright\small\mbox{}210}\hspace{0.5em}\parbox[t]{\dimexpr0.1633\linewidth-0.5em\relax}{\raggedright\small\mbox{}825}\hspace{0.5em}\parbox[t]{\dimexpr0.1633\linewidth-0.5em\relax}{\raggedright\small\mbox{}3}\hspace{0.5em}\parbox[t]{\dimexpr0.1633\linewidth-0.5em\relax}{\raggedright\small\mbox{}11}\hspace{0.5em}\parbox[t]{\dimexpr0.1633\linewidth-0.5em\relax}{\raggedright\small\mbox{}0}\hspace{0.5em}}{}
\EdL{14}{AK07-J0021-L14}{\parbox[t]{\dimexpr0.1633\linewidth-0.5em\relax}{\raggedright\small\mbox{}5}\hspace{0.5em}\parbox[t]{\dimexpr0.1633\linewidth-0.5em\relax}{\raggedright\small\mbox{}240}\hspace{0.5em}\parbox[t]{\dimexpr0.1633\linewidth-0.5em\relax}{\raggedright\small\mbox{}824\textsuperscript{a})}\hspace{0.5em}\parbox[t]{\dimexpr0.1633\linewidth-0.5em\relax}{\raggedright\small\mbox{}4}\hspace{0.5em}\parbox[t]{\dimexpr0.1633\linewidth-0.5em\relax}{\raggedright\small\mbox{}29}\hspace{0.5em}\parbox[t]{\dimexpr0.1633\linewidth-0.5em\relax}{\raggedright\small\mbox{}3}\hspace{0.5em}}{}
\EdL{15}{AK07-J0021-L15}{\parbox[t]{\dimexpr0.1633\linewidth-0.5em\relax}{\raggedright\small\mbox{}4}\hspace{0.5em}\parbox[t]{\dimexpr0.1633\linewidth-0.5em\relax}{\raggedright\small\mbox{}270}\hspace{0.5em}\parbox[t]{\dimexpr0.1633\linewidth-0.5em\relax}{\raggedright\small\mbox{}883}\hspace{0.5em}\parbox[t]{\dimexpr0.1633\linewidth-0.5em\relax}{\raggedright\small\mbox{}5}\hspace{0.5em}\parbox[t]{\dimexpr0.1633\linewidth-0.5em\relax}{\raggedright\small\mbox{}28}\hspace{0.5em}\parbox[t]{\dimexpr0.1633\linewidth-0.5em\relax}{\raggedright\small\mbox{}2}\hspace{0.5em}}{}
\EdL{16}{AK07-J0021-L16}{\parbox[t]{\dimexpr0.1633\linewidth-0.5em\relax}{\raggedright\small\mbox{}1}\hspace{0.5em}\parbox[t]{\dimexpr0.1633\linewidth-0.5em\relax}{\raggedright\small\mbox{}300}\hspace{0.5em}\parbox[t]{\dimexpr0.1633\linewidth-0.5em\relax}{\raggedright\small\mbox{}912}\hspace{0.5em}\parbox[t]{\dimexpr0.1633\linewidth-0.5em\relax}{\raggedright\small\mbox{}7}\hspace{0.5em}\parbox[t]{\dimexpr0.1633\linewidth-0.5em\relax}{\raggedright\small\mbox{}7}\hspace{0.5em}\parbox[t]{\dimexpr0.1633\linewidth-0.5em\relax}{\raggedright\small\mbox{}1}\hspace{0.5em}}{}
\EdL{17}{AK07-J0021-L17}{\parbox[t]{\dimexpr0.1633\linewidth-0.5em\relax}{\raggedright\small\mbox{}5}\hspace{0.5em}\parbox[t]{\dimexpr0.1633\linewidth-0.5em\relax}{\raggedright\small\mbox{}330}\hspace{0.5em}\parbox[t]{\dimexpr0.1633\linewidth-0.5em\relax}{\raggedright\small\mbox{}941}\hspace{0.5em}\parbox[t]{\dimexpr0.1633\linewidth-0.5em\relax}{\raggedright\small\mbox{}8}\hspace{0.5em}\parbox[t]{\dimexpr0.1633\linewidth-0.5em\relax}{\raggedright\small\mbox{}16}\hspace{0.5em}\parbox[t]{\dimexpr0.1633\linewidth-0.5em\relax}{\raggedright\small\mbox{}0}\hspace{0.5em}}{}
\EdL{18}{AK07-J0021-L18}{\parbox[t]{\dimexpr0.1633\linewidth-0.5em\relax}{\raggedright\small\mbox{}4}\hspace{0.5em}\parbox[t]{\dimexpr0.1633\linewidth-0.5em\relax}{\raggedright\small\mbox{}360}\hspace{0.5em}\parbox[t]{\dimexpr0.1633\linewidth-0.5em\relax}{\raggedright\small\mbox{}970}\hspace{0.5em}\parbox[t]{\dimexpr0.1633\linewidth-0.5em\relax}{\raggedright\small\mbox{}9}\hspace{0.5em}\parbox[t]{\dimexpr0.1633\linewidth-0.5em\relax}{\raggedright\small\mbox{}24}\hspace{0.5em}\parbox[t]{\dimexpr0.1633\linewidth-0.5em\relax}{\raggedright\small\mbox{}3}\hspace{0.5em}}{}
\EdL{19}{AK07-J0021-L19}{\parbox[t]{\dimexpr0.1633\linewidth-0.5em\relax}{\raggedright\small\mbox{}2}\hspace{0.5em}\parbox[t]{\dimexpr0.1633\linewidth-0.5em\relax}{\raggedright\small\mbox{}390}\hspace{0.5em}\parbox[t]{\dimexpr0.1633\linewidth-0.5em\relax}{\raggedright\small\mbox{}999}\hspace{0.5em}\parbox[t]{\dimexpr0.1633\linewidth-0.5em\relax}{\raggedright\small\mbox{}11}\hspace{0.5em}\parbox[t]{\dimexpr0.1633\linewidth-0.5em\relax}{\raggedright\small\mbox{}3}\hspace{0.5em}\parbox[t]{\dimexpr0.1633\linewidth-0.5em\relax}{\raggedright\small\mbox{}2}\hspace{0.5em}}{}
\EdL{20}{AK07-J0021-L20}{\parbox[t]{\dimexpr0.1633\linewidth-0.5em\relax}{\raggedright\small\mbox{}7}\hspace{0.5em}\parbox[t]{\dimexpr0.1633\linewidth-0.5em\relax}{\raggedright\small\mbox{}420}\hspace{0.5em}\parbox[t]{\dimexpr0.1633\linewidth-0.5em\relax}{\raggedright\small\mbox{}1029}\hspace{0.5em}\parbox[t]{\dimexpr0.1633\linewidth-0.5em\relax}{\raggedright\small\mbox{}0}\hspace{0.5em}\parbox[t]{\dimexpr0.1633\linewidth-0.5em\relax}{\raggedright\small\mbox{}7}\hspace{0.5em}\parbox[t]{\dimexpr0.1633\linewidth-0.5em\relax}{\raggedright\small\mbox{}0}\hspace{0.5em}}{}
\EdL{21}{AK07-J0021-L21}{\parbox[t]{\dimexpr0.1633\linewidth-0.5em\relax}{\raggedright\small\mbox{}5}\hspace{0.5em}\parbox[t]{\dimexpr0.1633\linewidth-0.5em\relax}{\raggedright\small\mbox{}450}\hspace{0.5em}\parbox[t]{\dimexpr0.1633\linewidth-0.5em\relax}{\raggedright\small\mbox{}1058}\hspace{0.5em}\parbox[t]{\dimexpr0.1633\linewidth-0.5em\relax}{\raggedright\small\mbox{}1}\hspace{0.5em}\parbox[t]{\dimexpr0.1633\linewidth-0.5em\relax}{\raggedright\small\mbox{}15}\hspace{0.5em}\parbox[t]{\dimexpr0.1633\linewidth-0.5em\relax}{\raggedright\small\mbox{}3}\hspace{0.5em}}{}
\EdL{22}{AK07-J0021-L22}{\parbox[t]{\dimexpr0.1633\linewidth-0.5em\relax}{\raggedright\small\mbox{}3}\hspace{0.5em}\parbox[t]{\dimexpr0.1633\linewidth-0.5em\relax}{\raggedright\small\mbox{}480}\hspace{0.5em}\parbox[t]{\dimexpr0.1633\linewidth-0.5em\relax}{\raggedright\small\mbox{}1087}\hspace{0.5em}\parbox[t]{\dimexpr0.1633\linewidth-0.5em\relax}{\raggedright\small\mbox{}2}\hspace{0.5em}\parbox[t]{\dimexpr0.1633\linewidth-0.5em\relax}{\raggedright\small\mbox{}24}\hspace{0.5em}\parbox[t]{\dimexpr0.1633\linewidth-0.5em\relax}{\raggedright\small\mbox{}2}\hspace{0.5em}}{}
\EdL{23}{AK07-J0021-L23}{\parbox[t]{\dimexpr0.1633\linewidth-0.5em\relax}{\raggedright\small\mbox{}1}\hspace{0.5em}\parbox[t]{\dimexpr0.1633\linewidth-0.5em\relax}{\raggedright\small\mbox{}510}\hspace{0.5em}\parbox[t]{\dimexpr0.1633\linewidth-0.5em\relax}{\raggedright\small\mbox{}1116}\hspace{0.5em}\parbox[t]{\dimexpr0.1633\linewidth-0.5em\relax}{\raggedright\small\mbox{}4}\hspace{0.5em}\parbox[t]{\dimexpr0.1633\linewidth-0.5em\relax}{\raggedright\small\mbox{}3}\hspace{0.5em}\parbox[t]{\dimexpr0.1633\linewidth-0.5em\relax}{\raggedright\small\mbox{}1}\hspace{0.5em}}{}
\EdL{24}{AK07-J0021-L24}{\parbox[t]{\dimexpr0.1633\linewidth-0.5em\relax}{\raggedright\small\mbox{}5}\hspace{0.5em}\parbox[t]{\dimexpr0.1633\linewidth-0.5em\relax}{\raggedright\small\mbox{}540}\hspace{0.5em}\parbox[t]{\dimexpr0.1633\linewidth-0.5em\relax}{\raggedright\small\mbox{}1145}\hspace{0.5em}\parbox[t]{\dimexpr0.1633\linewidth-0.5em\relax}{\raggedright\small\mbox{}5}\hspace{0.5em}\parbox[t]{\dimexpr0.1633\linewidth-0.5em\relax}{\raggedright\small\mbox{}12}\hspace{0.5em}\parbox[t]{\dimexpr0.1633\linewidth-0.5em\relax}{\raggedright\small\mbox{}0}\hspace{0.5em}}{}
\EdL{25}{AK07-J0021-L25}{\parbox[t]{\dimexpr0.1633\linewidth-0.5em\relax}{\raggedright\small\mbox{}4}\hspace{0.5em}\parbox[t]{\dimexpr0.1633\linewidth-0.5em\relax}{\raggedright\small\mbox{}570}\hspace{0.5em}\parbox[t]{\dimexpr0.1633\linewidth-0.5em\relax}{\raggedright\small\mbox{}1174}\hspace{0.5em}\parbox[t]{\dimexpr0.1633\linewidth-0.5em\relax}{\raggedright\small\mbox{}6}\hspace{0.5em}\parbox[t]{\dimexpr0.1633\linewidth-0.5em\relax}{\raggedright\small\mbox{}20}\hspace{0.5em}\parbox[t]{\dimexpr0.1633\linewidth-0.5em\relax}{\raggedright\small\mbox{}3}\hspace{0.5em}}{}
\EdL{26}{AK07-J0021-L26}{\parbox[t]{\dimexpr0.1633\linewidth-0.5em\relax}{\raggedright\small\mbox{}2}\hspace{0.5em}\parbox[t]{\dimexpr0.1633\linewidth-0.5em\relax}{\raggedright\small\mbox{}600}\hspace{0.5em}\parbox[t]{\dimexpr0.1633\linewidth-0.5em\relax}{\raggedright\small\mbox{}1203}\hspace{0.5em}\parbox[t]{\dimexpr0.1633\linewidth-0.5em\relax}{\raggedright\small\mbox{}7}\hspace{0.5em}\parbox[t]{\dimexpr0.1633\linewidth-0.5em\relax}{\raggedright\small\mbox{}29}\hspace{0.5em}\parbox[t]{\dimexpr0.1633\linewidth-0.5em\relax}{\raggedright\small\mbox{}2?}\hspace{0.5em}}{}
\EdL{27}{AK07-J0021-L27}{\parbox[t]{\dimexpr0.1633\linewidth-0.5em\relax}{\raggedright\small\mbox{}7}\hspace{0.5em}\parbox[t]{\dimexpr0.1633\linewidth-0.5em\relax}{\raggedright\small\mbox{}630}\hspace{0.5em}\parbox[t]{\dimexpr0.1633\linewidth-0.5em\relax}{\raggedright\small\mbox{}1232}\hspace{0.5em}\parbox[t]{\dimexpr0.1633\linewidth-0.5em\relax}{\raggedright\small\mbox{}9}\hspace{0.5em}\parbox[t]{\dimexpr0.1633\linewidth-0.5em\relax}{\raggedright\small\mbox{}7}\hspace{0.5em}\parbox[t]{\dimexpr0.1633\linewidth-0.5em\relax}{\raggedright\small\mbox{}1}\hspace{0.5em}}{}
\EdL{28}{AK07-J0021-L28}{\parbox[t]{\dimexpr0.1633\linewidth-0.5em\relax}{\raggedright\small\mbox{}5}\hspace{0.5em}\parbox[t]{\dimexpr0.1633\linewidth-0.5em\relax}{\raggedright\small\mbox{}660}\hspace{0.5em}\parbox[t]{\dimexpr0.1633\linewidth-0.5em\relax}{\raggedright\small\mbox{}1261}\hspace{0.5em}\parbox[t]{\dimexpr0.1633\linewidth-0.5em\relax}{\raggedright\small\mbox{}10}\hspace{0.5em}\parbox[t]{\dimexpr0.1633\linewidth-0.5em\relax}{\raggedright\small\mbox{}17}\hspace{0.5em}\parbox[t]{\dimexpr0.1633\linewidth-0.5em\relax}{\raggedright\small\mbox{}0}\hspace{0.5em}}{}
\EdL{29}{AK07-J0021-L29}{\parbox[t]{\dimexpr0.1633\linewidth-0.5em\relax}{\raggedright\small\mbox{}}\hspace{0.5em}\parbox[t]{\dimexpr0.1633\linewidth-0.5em\relax}{\raggedright\small\mbox{}690}\hspace{0.5em}\parbox[t]{\dimexpr0.1633\linewidth-0.5em\relax}{\raggedright\small\mbox{}1290}\hspace{0.5em}\parbox[t]{\dimexpr0.1633\linewidth-0.5em\relax}{\raggedright\small\mbox{}11}\hspace{0.5em}\parbox[t]{\dimexpr0.1633\linewidth-0.5em\relax}{\raggedright\small\mbox{}25}\hspace{0.5em}\parbox[t]{\dimexpr0.1633\linewidth-0.5em\relax}{\raggedright\small\mbox{}3}\hspace{0.5em}}{}
\EdL{30}{AK07-J0021-L30}{\parbox[t]{\dimexpr0.1633\linewidth-0.5em\relax}{\raggedright\small\mbox{}}\hspace{0.5em}\parbox[t]{\dimexpr0.1633\linewidth-0.5em\relax}{\raggedright\small\mbox{}720}\hspace{0.5em}\parbox[t]{\dimexpr0.1633\linewidth-0.5em\relax}{\raggedright\small\mbox{}1320}\hspace{0.5em}\parbox[t]{\dimexpr0.1633\linewidth-0.5em\relax}{\raggedright\small\mbox{}1}\hspace{0.5em}\parbox[t]{\dimexpr0.1633\linewidth-0.5em\relax}{\raggedright\small\mbox{}4}\hspace{0.5em}\parbox[t]{\dimexpr0.1633\linewidth-0.5em\relax}{\raggedright\small\mbox{}2}\hspace{0.5em}}{}
\EdL{31}{AK07-J0021-L31}{\parbox[t]{\dimexpr0.1633\linewidth-0.5em\relax}{\raggedright\small\mbox{}}\hspace{0.5em}\parbox[t]{\dimexpr0.1633\linewidth-0.5em\relax}{\raggedright\small\mbox{}750}\hspace{0.5em}\parbox[t]{\dimexpr0.1633\linewidth-0.5em\relax}{\raggedright\small\mbox{}1349}\hspace{0.5em}\parbox[t]{\dimexpr0.1633\linewidth-0.5em\relax}{\raggedright\small\mbox{}2}\hspace{0.5em}\parbox[t]{\dimexpr0.1633\linewidth-0.5em\relax}{\raggedright\small\mbox{}13}\hspace{0.5em}\parbox[t]{\dimexpr0.1633\linewidth-0.5em\relax}{\raggedright\small\mbox{}1}\hspace{0.5em}}{}
\EdLp{32}{AK07-J0021-L32}{Zu dieser Tafel haben wir zu bemerken\,: Unter a) sollte es heissen 854\,; die}{}
\EdL{33}{AK07-J0021-L33}{Zahlen der letzten Kolonne bedeuten „Viertel“\,; die Zahlen der Kolonne \textit{Dies} sind}{}
\EdL{34}{AK07-J0021-L34}{öfters fehlerhaft, man kann sie nach den heutigen Tafeln leicht verbessern\,; auch die}{}
\EdL{35}{AK07-J0021-L35}{\textit{Fractiones} stimmen nicht durchweg. Die letzten drei Zeilen für die Jahre 690, 720 u.}{}
\EdL{36}{AK07-J0021-L36}{750 sind von einer neuern Hand hinzugefügt, und zwar in arabischen Ziffern, die mir}{}
\EdL{37}{AK07-J0021-L37}{diejenigen des 14. bis 15. Jahrhunderts zu sein scheinen\textsuperscript{1}).}{}
\EdLp{38}{AK07-J0021-L38}{5) Am Rande von fol. 9\textsuperscript{r} steht folgende Note von 2. Hand\,: \textit{Era (Aera) quoque} (?)}{}
\EdL{39}{AK07-J0021-L39}{\textit{cesaris augusti tempore posita est\,: dicta autem era ex quo omnis orbis es (aes) reddere}}{}
\EdL{40}{AK07-J0021-L40}{\textit{professus est reipublicae. Era secundum ysidorum adcescit (sic!) — Bissextus autem}}{}
\NoteRule
\EdN{N1}{AK07-J0021-N01}{\textsuperscript{1}) Die Ziffern der 2., 4. u. 6. Kolonne waren für mich wegen zu schwacher photographischer}
\EdN{N2}{AK07-J0021-N02}{Abbildung unlesbar, ich verdanke eine Abschrift derselben aus dem Original Herrn \Sperr{Gonzalez},}
\EdN{N3}{AK07-J0021-N03}{Archivar in Madrid, einem Schüler von Herrn Prof. \Sperr{Asín}, aufs beste.}
\FrameHead{0022}{p. XVIII}{0022}{p. XVIII}
\EdL{1}{AK07-J0022-L01}{\textit{accescitur}\textsuperscript{1}) \textit{a sexto. Nonas marcii usque ad diem pridie kalendarum ianuarii apponitur}}{}
\EdL{2}{AK07-J0022-L02}{\textit{in cursu lunae} (?)}{}
\EdLp{3}{AK07-J0022-L03}{Die Worte \textit{dicta autem era ex quo . . . reipublicae} sind in der Tat ein Zitat aus}{}
\EdL{4}{AK07-J0022-L04}{\Sperr{Isidor v. Sevilla}, das auch von L. \Sperr{Ideler} in seinen Versuchen, das Wort \textit{aera}}{}
\EdL{5}{AK07-J0022-L05}{zu erklären, angeführt wird\textsuperscript{2}). Die Erklärung von \textit{bissextus} ist kaum zu verstehen,}{}
\EdL{6}{AK07-J0022-L06}{und wohl auch unrichtig aufgefasst vom Schreiber. Bekanntlich kommt es davon}{}
\EdL{7}{AK07-J0022-L07}{her, dass man in Schaltjahren \textit{zwei mal} gesagt hat\,: \textit{Sextus ante Kalendas Martii}, also}{}
\EdL{8}{AK07-J0022-L08}{den 24. Februar \textit{zwei mal} gezählt hat. Der Schreiber nimmt statt der \textit{Kalendae Martii}}{}
\EdL{9}{AK07-J0022-L09}{die \textit{Nonae Martii}, und meint dann wohl mit dem folgenden, dass der Schalttag an}{}
\EdL{10}{AK07-J0022-L10}{verschiedenen Stellen eingeschaltet worden sei, sogar auch am Ende des Monats}{}
\EdL{11}{AK07-J0022-L11}{Dezember, was, wie wir unten sehen werden, in der sog. spanischen Aera der Fall}{}
\EdL{12}{AK07-J0022-L12}{war. — Was mit den Worten \textit{in cursu lunae} gesagt werden will, verstehen wir nicht.}{}
\EdLp{13}{AK07-J0022-L13}{6) Auf fol. 11\textsuperscript{v} ist zu den Worten \textit{ideo quotiens duodecim menses in annum redu-}}{}
\EdL{14}{AK07-J0022-L14}{\textit{cimus, totiens ab ipsa reductionis summa} (in unserm Text p. 6) folgende Randnote}{}
\EdL{15}{AK07-J0022-L15}{von 2. Hand beigefügt\,: \textit{id est a mensibus et diebus residuis ex quibus fit annus qui}}{}
\EdL{16}{AK07-J0022-L16}{\textit{aliis annis adiicitur\,; non enim complebitur annus ex XII mensibus XXX dierum nisi}}{}
\EdL{17}{AK07-J0022-L17}{\textit{addatur V dies et quarta\,; menses enim tantum XXX dierum hic dantur} (?), \textit{unde erit eis}}{}
\EdL{18}{AK07-J0022-L18}{\textit{quandoque addendum quandoque subtrahendum ut vere romani menses habeantur.}}{}
\EdLp{19}{AK07-J0022-L19}{7) Am untern Rand von fol. 18\textsuperscript{v} (Tafel 15) findet sich folgende Note\,: \textit{Adiectiones}}{}
\EdL{20}{AK07-J0022-L20}{\textit{Veneris et Mercurii sunt adiectio solis\,; hae autem adiectiones omnes iuxta civitatem Lon-}}{}
\EdL{21}{AK07-J0022-L21}{\textit{doniae in hoc libro computantur et mediis cursibus planetarum adiciuntur} (sic). — Hiermit}{}
\EdL{22}{AK07-J0022-L22}{will der Verfasser dieser Note, und dies ist ohne Zweifel \Sperr{Robert von Chester}\textsuperscript{3}),}{}
\EdL{23}{AK07-J0022-L23}{sagen, dass die mittlern Längen von Venus und Merkur die gleichen seien wie die}{}
\EdL{24}{AK07-J0022-L24}{der Sonne, und dass er für diese und die drei obern Planeten die \textit{adiectiones} für}{}
\EdL{25}{AK07-J0022-L25}{den Meridian von London berechnet habe. Unter diesen \textit{adiectiones} ist folgendes zu}{}
\EdL{26}{AK07-J0022-L26}{verstehen\,: Die nach den Tafeln 9—20 für Sonne, Mond und die Planeten berechneten}{}
\EdL{27}{AK07-J0022-L27}{mittlern Längen gelten für den Meridian von \textit{Arin} (\textit{Udjain}), für London muss der}{}
\EdL{28}{AK07-J0022-L28}{mittlern Länge jedes dieser Himmelskörper eine bestimmte Grösse beigefügt werden,}{}
\EdL{29}{AK07-J0022-L29}{und zwar die mittlere Bewegung (\textit{motus medius}) desselben für c. 5\textsuperscript{h} (denn Udjain ist}{}
\EdL{30}{AK07-J0022-L30}{c. 75° östlich von London), diese zu addierende Grösse nennt \Sperr{Robert von Chester}}{}
\EdL{31}{AK07-J0022-L31}{\textit{adiectio}. — Er hat nun diese \textit{adiectiones} an verschiedenen Stellen angegeben (jeweilen}{}
\EdL{32}{AK07-J0022-L32}{am untern Rand des Blattes) und zwar\,: für Saturn auf fol. 24\textsuperscript{v} (Tafel 9) zu 18″,}{}
\EdL{33}{AK07-J0022-L33}{für Jupiter auf fol. 28\textsuperscript{v} (Tafel 11) zu 45″, für Mars auf fol. 14\textsuperscript{v} (Tafel 13) zu 4′ 49″,}{}
\EdL{34}{AK07-J0022-L34}{für den Drachenkopf auf fol. 23\textsuperscript{v} (Tafel 19) zu 30″\,; die Tafeln für die mittlere}{}
\EdL{35}{AK07-J0022-L35}{Bewegung der Sonne und des Mondes fehlen in N. Diese Zahlen stimmen aber}{}
\EdL{36}{AK07-J0022-L36}{nicht mit unsern Tafeln für einen Zeitunterschied von 5\textsuperscript{h}, \Sperr{Robert von Chester}}{}
\EdL{37}{AK07-J0022-L37}{scheint diesen nicht ganz 4\textsuperscript{h} angenommen zu haben, die angenähert richtigen Zahlen}{}
\NoteRule
\EdN{N1}{AK07-J0022-N01}{\textsuperscript{1}) Wohl für \textit{accersitur} = \textit{arcessitur} (= herholen, hernehmen).}
\EdN{N2}{AK07-J0022-N02}{\textsuperscript{2}) \textit{Handbuch der mathem. u. technischen Chronologie}, Berlin 1825—26, II, 426—430.}
\EdN{N3}{AK07-J0022-N03}{\textsuperscript{3}) Denn dieser hat nach dem Zeugnis des \Sperr{Hermannus Dalmata} das Werk \Sperr{Al-Battānīs}}
\EdN{N4}{AK07-J0022-N04}{ins Lateinische übersetzt, und soll auch die Tafeln desselben für den Meridian von London um-}
\EdN{N5}{AK07-J0022-N05}{gerechnet haben. Vergl. oben p. XIII, und \Sperr{Al-Battānī}, ed. \Sperr{Nallino} I, \textit{Praefatio}, XXXIV–XXXV,}
\EdN{N6}{AK07-J0022-N06}{wo aber noch \Sperr{Rudolphus Brugensis} statt \Sperr{Hermannus Dalmata} steht.}
\FrameHead{0023}{p. XIX}{0023}{p. XIX}
\EdL{1}{AK07-J0023-L01}{sind nach den Tafeln 10, 12, 14 u. 20\,: Saturn 25″, Jupiter 1′ 2″, Mars 6′ 33″, Drachen-}{}
\EdL{2}{AK07-J0023-L02}{kopf 40″ (für die Sonne 12′ 19″, für den Mond 2° 44′ 42″).}{}
\EdLp{3}{AK07-J0023-L03}{8) Auf fol. 23\textsuperscript{v} (Tafel 19) fügt N in der linken Ecke unten folgende in den}{}
\EdL{4}{AK07-J0023-L04}{Sekunden und teilweise auch in den Minuten von den unmittelbar oberhalb stehenden}{}
\EdL{5}{AK07-J0023-L05}{Daten der \textit{Medialitas draconis} abweichende Zahlen hinzu\,:}{}
\EdL{6}{AK07-J0023-L06}{\parbox[t]{\dimexpr0.1960\linewidth-0.5em\relax}{\raggedright\small\mbox{}collecti}\hspace{0.5em}\parbox[t]{\dimexpr0.7840\linewidth-0.5em\relax}{\raggedright\small\mbox{}medio die Cordubae}\hspace{0.5em}}{}
\EdL{7}{AK07-J0023-L07}{\parbox[t]{\dimexpr0.1960\linewidth-0.5em\relax}{\raggedright\small\mbox{}360}\hspace{0.5em}\parbox[t]{\dimexpr0.1960\linewidth-0.5em\relax}{\raggedright\small\mbox{}5}\hspace{0.5em}\parbox[t]{\dimexpr0.1960\linewidth-0.5em\relax}{\raggedright\small\mbox{}5}\hspace{0.5em}\parbox[t]{\dimexpr0.1960\linewidth-0.5em\relax}{\raggedright\small\mbox{}20}\hspace{0.5em}\parbox[t]{\dimexpr0.1960\linewidth-0.5em\relax}{\raggedright\small\mbox{}42}\hspace{0.5em}}{}
\EdL{8}{AK07-J0023-L08}{\parbox[t]{\dimexpr0.1960\linewidth-0.5em\relax}{\raggedright\small\mbox{}390}\hspace{0.5em}\parbox[t]{\dimexpr0.1960\linewidth-0.5em\relax}{\raggedright\small\mbox{}11}\hspace{0.5em}\parbox[t]{\dimexpr0.1960\linewidth-0.5em\relax}{\raggedright\small\mbox{}28}\hspace{0.5em}\parbox[t]{\dimexpr0.1960\linewidth-0.5em\relax}{\raggedright\small\mbox{}48}\hspace{0.5em}\parbox[t]{\dimexpr0.1960\linewidth-0.5em\relax}{\raggedright\small\mbox{}24}\hspace{0.5em}}{}
\EdL{9}{AK07-J0023-L09}{\parbox[t]{\dimexpr0.1960\linewidth-0.5em\relax}{\raggedright\small\mbox{}420}\hspace{0.5em}\parbox[t]{\dimexpr0.1960\linewidth-0.5em\relax}{\raggedright\small\mbox{}6}\hspace{0.5em}\parbox[t]{\dimexpr0.1960\linewidth-0.5em\relax}{\raggedright\small\mbox{}22}\hspace{0.5em}\parbox[t]{\dimexpr0.1960\linewidth-0.5em\relax}{\raggedright\small\mbox{}16}\hspace{0.5em}\parbox[t]{\dimexpr0.1960\linewidth-0.5em\relax}{\raggedright\small\mbox{}50\,?}\hspace{0.5em}}{}
\EdL{10}{AK07-J0023-L10}{\parbox[t]{\dimexpr0.1960\linewidth-0.5em\relax}{\raggedright\small\mbox{}450}\hspace{0.5em}\parbox[t]{\dimexpr0.1960\linewidth-0.5em\relax}{\raggedright\small\mbox{}1}\hspace{0.5em}\parbox[t]{\dimexpr0.1960\linewidth-0.5em\relax}{\raggedright\small\mbox{}15}\hspace{0.5em}\parbox[t]{\dimexpr0.1960\linewidth-0.5em\relax}{\raggedright\small\mbox{}43}\hspace{0.5em}\parbox[t]{\dimexpr0.1960\linewidth-0.5em\relax}{\raggedright\small\mbox{}49}\hspace{0.5em}}{}
\EdL{11}{AK07-J0023-L11}{\parbox[t]{\dimexpr0.1960\linewidth-0.5em\relax}{\raggedright\small\mbox{}480}\hspace{0.5em}\parbox[t]{\dimexpr0.1960\linewidth-0.5em\relax}{\raggedright\small\mbox{}8}\hspace{0.5em}\parbox[t]{\dimexpr0.1960\linewidth-0.5em\relax}{\raggedright\small\mbox{}9}\hspace{0.5em}\parbox[t]{\dimexpr0.1960\linewidth-0.5em\relax}{\raggedright\small\mbox{}11}\hspace{0.5em}\parbox[t]{\dimexpr0.1960\linewidth-0.5em\relax}{\raggedright\small\mbox{}32}\hspace{0.5em}}{}
\EdL{12}{AK07-J0023-L12}{\parbox[t]{\dimexpr0.1960\linewidth-0.5em\relax}{\raggedright\small\mbox{}510}\hspace{0.5em}\parbox[t]{\dimexpr0.1960\linewidth-0.5em\relax}{\raggedright\small\mbox{}3}\hspace{0.5em}\parbox[t]{\dimexpr0.1960\linewidth-0.5em\relax}{\raggedright\small\mbox{}2}\hspace{0.5em}\parbox[t]{\dimexpr0.1960\linewidth-0.5em\relax}{\raggedright\small\mbox{}39}\hspace{0.5em}\parbox[t]{\dimexpr0.1960\linewidth-0.5em\relax}{\raggedright\small\mbox{}14}\hspace{0.5em}}{}
\EdL{13}{AK07-J0023-L13}{\parbox[t]{\dimexpr0.1960\linewidth-0.5em\relax}{\raggedright\small\mbox{}540}\hspace{0.5em}\parbox[t]{\dimexpr0.1960\linewidth-0.5em\relax}{\raggedright\small\mbox{}9}\hspace{0.5em}\parbox[t]{\dimexpr0.1960\linewidth-0.5em\relax}{\raggedright\small\mbox{}26}\hspace{0.5em}\parbox[t]{\dimexpr0.1960\linewidth-0.5em\relax}{\raggedright\small\mbox{}6}\hspace{0.5em}\parbox[t]{\dimexpr0.1960\linewidth-0.5em\relax}{\raggedright\small\mbox{}57}\hspace{0.5em}}{}
\EdL{14}{AK07-J0023-L14}{\parbox[t]{\dimexpr0.1960\linewidth-0.5em\relax}{\raggedright\small\mbox{}570}\hspace{0.5em}\parbox[t]{\dimexpr0.1960\linewidth-0.5em\relax}{\raggedright\small\mbox{}4}\hspace{0.5em}\parbox[t]{\dimexpr0.1960\linewidth-0.5em\relax}{\raggedright\small\mbox{}19}\hspace{0.5em}\parbox[t]{\dimexpr0.1960\linewidth-0.5em\relax}{\raggedright\small\mbox{}34}\hspace{0.5em}\parbox[t]{\dimexpr0.1960\linewidth-0.5em\relax}{\raggedright\small\mbox{}39}\hspace{0.5em}}{}
\EdLp{15}{AK07-J0023-L15}{9) Auf fol. 24\textsuperscript{r} (Tafel 20) hat N in der letzten Kolonne (die Tertien lässt es}{}
\EdL{16}{AK07-J0023-L16}{weg) folgende Zahlen für die Sekunden\,: 0, 0, 0, 0, 1, 1, 1, 1, 2, 2, 2, 3, 3, 3, 4, 4, 4, 4, 5,}{}
\EdL{17}{AK07-J0023-L17}{5, 5, 6, 6, 6, 6, 7, 7, 7, 8, 8. Man sieht, dass verschiedene dieser Zahlen, auch wenn}{}
\EdL{18}{AK07-J0023-L18}{die abgeworfenen Tertien berücksichtigt werden, nicht mit denen der Mss. C u. O}{}
\EdL{19}{AK07-J0023-L19}{übereinstimmen.}{}
\EdLp{20}{AK07-J0023-L20}{10) Auf fol. 42\textsuperscript{r} ist nach den Worten\,: \textit{Sed et aequalibus idem e contrario addendo}}{}
\EdL{21}{AK07-J0023-L21}{\textit{inaequales fient} (in unserm Text p. 25) von 2. Hand hinzugefügt\,: \textit{Verus dies naturalis}}{}
\EdL{22}{AK07-J0023-L22}{\textit{medius est XXIIII horarum, id est revolutio CCCLX graduum et LIX min. et V sec. Unde}}{}
\EdL{23}{AK07-J0023-L23}{(od. \textit{una}) \textit{hora est XV grad. et II. min. et XXVIII sec. Sed quoniam gradus aequinoctialis}}{}
\EdL{24}{AK07-J0023-L24}{\textit{lineae, secundum quam haec determinata} (?) \textit{sunt, et zodiaci in quo sol movetur, non}}{}
\EdL{25}{AK07-J0023-L25}{\textit{similiter transeunt meridianam lineam, sed et sol aliquando citius aliquando tardius}}{}
\EdL{26}{AK07-J0023-L26}{\textit{contra celum movetur, necesse est ergo naturali diei medio aliquando addi, aliquando}}{}
\EdL{27}{AK07-J0023-L27}{\textit{subtrahi. Attamen haec tabula sic facta est, quod semper subtrahitur a die media, sed}}{}
\EdL{28}{AK07-J0023-L28}{\textit{aliquando sequenti die maius aliquando minus quam priori. Sed quando maius tunc}}{}
\EdL{29}{AK07-J0023-L29}{\textit{ille dies aequatus a meridiana linea in meridianam erit minor die media, si autem minus}}{}
\EdL{30}{AK07-J0023-L30}{\textit{tunc necesse est esse maiorem\,; maior enim subtractio minuit de medio, minor vero sub-}}{}
\EdL{31}{AK07-J0023-L31}{\textit{tractio de medio et aequali residuum relinquit maius quam faceret subtractio maior.}}{}
\EdL{32}{AK07-J0023-L32}{\textit{Exempli gratia\,: si hodie a revolutione sperae}\textsuperscript{1}) \textit{et LIX min. et VIII sec. subtraherentur}}{}
\EdL{33}{AK07-J0023-L33}{\textit{II min. et VIII sec., et in sequenti die III min. et VIII sec., esset tunc dies aequatus}}{}
\EdL{34}{AK07-J0023-L34}{\textit{revolutio sperae et LVIII min., quod est minus die medio\,; item si priori die III min. et}}{}
\EdL{35}{AK07-J0023-L35}{\textit{VIII sec. subtraherentur et sequenti die tantum II min. et VIII sec., erit dies aequatus}}{}
\EdL{36}{AK07-J0023-L36}{\textit{revolutio sperae et LX min., quod est maius die medio.}}{}
\EdLp{37}{AK07-J0023-L37}{11) In den Tafeln der Konjunktion und Opposition von Sonne und Mond für}{}
\EdL{38}{AK07-J0023-L38}{die \textit{anni collecti} (fol. 42\textsuperscript{v} u. 43\textsuperscript{r}, Tafeln 69 u. 70) führt N noch die Daten für die Jahre}{}
\NoteRule
\EdN{N1}{AK07-J0023-N01}{\textsuperscript{1}) \textit{revolutio sperae} = 360°.}
\FrameHead{0024}{p. XX}{0024}{p. XX}
\EdL{1}{AK07-J0024-L01}{541 und 571 an, die in C und O fehlen\,; die erstern scheinen von gleicher Hand zu}{}
\EdL{2}{AK07-J0024-L02}{sein, wie die vorhergehenden, die letztern von einer andern Hand. Da die Tafel}{}
\EdL{3}{AK07-J0024-L03}{stets von jedermann fortgesetzt werden kann, so verzichten wir auf die Wiedergabe}{}
\EdL{4}{AK07-J0024-L04}{dieser Zeilen.}{}
\EdLp{5}{AK07-J0024-L05}{12) Auf fol. 47\textsuperscript{v} (in unserm Text p. 29) ist nach den Worten \textit{erit et ipsa meri-}}{}
\EdL{6}{AK07-J0024-L06}{\textit{dionalis} von 2. Hand (?) hinzugefügt\,: \textit{Haec circa meridiem, alibi autem\,: Inventa vera}}{}
\EdL{7}{AK07-J0024-L07}{\textit{lunae latitudine nadair}\textsuperscript{1}) \textit{gradus coniunctionis ponatur in astrolabio super horam con-}}{}
\EdL{8}{AK07-J0024-L08}{\textit{iunctiones et videatur quot gradibus distet gradus coniunctionis a centro minimi circuli}}{}
\EdL{9}{AK07-J0024-L09}{\textit{almucantarat, et si fuit vera lunae latitudo septemtrionalis, auferatur ab illis gradibus,}}{}
\EdL{10}{AK07-J0024-L10}{\textit{si meridiana, addatur\,; et per id fiat introitus in tabulas latitudinis per aspectum.}}{}
\EdLp{11}{AK07-J0024-L11}{13) Auf fol. 49\textsuperscript{r} (in unserm Text p. 29) ist nach den Worten\,: \textit{horas coniunctionis}}{}
\EdL{12}{AK07-J0024-L12}{\textit{quaesitae ponet} von 2. Hand hinzugefügt\,: \textit{Introitus harum tabularum hic est\,: primo}}{}
\EdL{13}{AK07-J0024-L13}{\textit{videatur coniunctionis distantia solis et lunae a capite vel cauda et si fuerit luna in}}{}
\EdL{14}{AK07-J0024-L14}{\textit{longiori longitudine sui brevis circuli}\textsuperscript{2}) \textit{intretur in tabulas priores, aliter in posteriores}}{}
\EdL{15}{AK07-J0024-L15}{\textit{et sumatur de digitis quod opponitur distantiae praedictae. Exinde sumatur diversitas}}{}
\EdL{16}{AK07-J0024-L16}{\textit{aspectus lunae et videatur quot digitis ipsa aequipolleat. Et si fuerit luna tunc septem-}}{}
\EdL{17}{AK07-J0024-L17}{\textit{trionalis a sole addantur illi digiti super inventos digitos in tabula, si meridiana, demantur}}{}
\EdL{18}{AK07-J0024-L18}{\textit{si poterit, aliter enim non erit eclipsis\,; et sic inventum erit quantitas eclipsis. Et tunc}}{}
\EdL{19}{AK07-J0024-L19}{\textit{cum illo invento in tabula sumantur minuta casus in eius directo et erit mora usque ad}}{}
\EdL{20}{AK07-J0024-L20}{\textit{mediam eclipsim. In diversitate aspectus quot fuerint digiti sic inveniatur\,: Per tabulas}}{}
\EdL{21}{AK07-J0024-L21}{\textit{elbut inveniatur quantitas solaris corporis et dividatur in XII aequas partes quia semper}}{}
\EdL{22}{AK07-J0024-L22}{\textit{duodecima diametri solis digitus dicitur. Videatur quotiens eius duodecima in diversitate}}{}
\EdL{23}{AK07-J0024-L23}{\textit{aspectus contineatur, et tot fient ex ea digiti. Si ergo inventum fuerit plus XII digitis}}{}
\EdL{24}{AK07-J0024-L24}{\textit{et XLVIII min. in longitudine propiori vel plus X et XLVIII min. in longiori, et luna}}{}
\EdL{25}{AK07-J0024-L25}{\textit{fuerit septemtrionalis a sole, auferantur praedicti digiti de invento et residuum clarum}}{}
\EdL{26}{AK07-J0024-L26}{\textit{in sole relinquitur. Si vero luna fuerit meridiana et inventum secundum se obscurabit}}{}
\EdL{27}{AK07-J0024-L27}{\textit{de sole, haec secundum eos quibus sol semper meridianus est\,; quia cum fuerit septem-}}{}
\EdL{28}{AK07-J0024-L28}{\textit{trionalis contra operari oportebit.}}{}
\EdLp{29}{AK07-J0024-L29}{Hierauf folgt von 3. Hand\,: \textit{Sed et in maiori elongatione ante caudam vel post}}{}
\EdL{30}{AK07-J0024-L30}{\textit{caput, quia usque ad XXIII gradus fit solis eclipsis. Et tunc sciatur vera lunae latitudo}}{}
\EdL{31}{AK07-J0024-L31}{\textit{et diversitas aspectus latitudinis eius\,; et auferatur minus de maiori, et per residuum erit}}{}
\EdL{32}{AK07-J0024-L32}{\textit{quantitas scita eclipsis, cuius digiti per solis et residui quantitatem, ut praedictum est,}}{}
\EdL{33}{AK07-J0024-L33}{\textit{investigentur, et per eos mora in tabula. Iungantur autem medietates}\textsuperscript{3}) \textit{diametrorum}}{}
\EdL{34}{AK07-J0024-L34}{\textit{solis et lunae, et si residuum praedictum fuerit eis aequale vel maius, non erit eclipsis,}}{}
\EdL{35}{AK07-J0024-L35}{\textit{si nichilum fuerit residuum, totalis, si vero residuum minus fuerit, erit eclipsis partialis.}}{}
\EdLp{36}{AK07-J0024-L36}{Wir haben im Kommentar zu Kap. 35 bemerkt, dass dieses sehr kurz und}{}
\EdL{37}{AK07-J0024-L37}{jedenfalls unvollständig sei. Dies scheint auch der Verfasser des ersten der obigen}{}
\EdL{38}{AK07-J0024-L38}{Zusätze gefühlt zu haben, und glaubte daher, eine bessere, verständlichere Darstellung}{}
\EdL{39}{AK07-J0024-L39}{der Bestimmung einer Sonnenfinsternis geben zu müssen. Allein seine Regel ist,}{}
\NoteRule
\EdN{N1}{AK07-J0024-N01}{\textsuperscript{1}) Wahrscheinlich \textit{nadir}.}
\EdN{N2}{AK07-J0024-N02}{\textsuperscript{2}) \textit{brevis circuli} = \textit{epicycli}.}
\EdN{N3}{AK07-J0024-N03}{\textsuperscript{3}) Im Text \textit{mediates}.}
\FrameHead{0025}{p. XXI}{0025}{p. XXI}
\EdL{1}{AK07-J0025-L01}{wenn man sie mit derjenigen \Sperr{Al-Battānīs} vergleicht, wohl auch nicht richtig, zum}{}
\EdL{2}{AK07-J0025-L02}{mindesten ebenfalls unvollständig, was eine spätere Hand veranlasst hat, eine zweite}{}
\EdL{3}{AK07-J0025-L03}{Berichtigung hinzuzufügen.}{}
\EdLp{4}{AK07-J0025-L04}{Im Text und in den Tafeln kommen einige ganz verdorbene \textit{arabische} Wörter und}{}
\EdL{5}{AK07-J0025-L05}{Sätze vor, die ich an dieser Stelle zu erklären versuchen möchte, da ich dieselben}{}
\EdL{6}{AK07-J0025-L06}{im Glossar nicht gut unter einem bestimmen Stichwort unterbringen konnte. Wie}{}
\EdL{7}{AK07-J0025-L07}{sie in das Werk gekommen sind, ist mit Ausnahme der persischen Namen für die}{}
\EdL{8}{AK07-J0025-L08}{Planeten schwer zu entscheiden, mir scheinen sie blosse Spielereien des Übersetzers,}{}
\EdL{9}{AK07-J0025-L09}{oder eines Abschreibers zu sein\,; es ist kaum anzunehmen, dass sie von \Sperr{Maslama}}{}
\EdL{10}{AK07-J0025-L10}{\Sperr{al-Madjrīṭī}, oder gar von \Sperr{Al-Khwārizmī} herstammen. Für die freundliche}{}
\EdL{11}{AK07-J0025-L11}{Hilfe, die mir bei den schwierigen Stellen die Herren Proff. \Sperr{Nallino} in Palermo und}{}
\EdL{12}{AK07-J0025-L12}{\Sperr{Seybold} in Tübingen geleistet haben, spreche ich ihnen auch an dieser Stelle}{}
\EdL{13}{AK07-J0025-L13}{meinen verbindlichen Dank aus.}{}
\EdLp{14}{AK07-J0025-L14}{Im Ms. O steht auf fol. 103\textsuperscript{v} (Tafel 7) folgende Note am obern Rand\,: \textit{kulelach.}}{}
\EdL{15}{AK07-J0025-L15}{\textit{vveli. kemithl. elhileli. fikemelen. wanuczeni.} Dies soll wohl heissen\,: \textit{kull al-aḥwāli}}{}
\EdL{16}{AK07-J0025-L16}{\textit{kamithl al-hilāli fī kamālin wanuḳṣānin,} und bedeutet frei übersetzt\,: Alle Verände-}{}
\EdL{17}{AK07-J0025-L17}{rungen (des Lebens, der Natur) verhalten sich wie der Mond, (sie äussern sich) in}{}
\EdL{18}{AK07-J0025-L18}{einer Zunahme (bis zur Vollständigkeit), und in einer (darauf folgenden) Abnahme.}{}
\EdLp{19}{AK07-J0025-L19}{In demselben Ms. steht auf fol. 104\textsuperscript{v} (Tafel 9) über dem Kolonnentitel \textit{Heza Saturni}}{}
\EdL{20}{AK07-J0025-L20}{\textit{supradictum} folgende Note\,: \textit{Nota quod hoc argumentum felek eltedewir servit. bifelek}}{}
\EdL{21}{AK07-J0025-L21}{\textit{eltedewir lieztekim asmir.} — \textit{Felek eltedewir} = \textit{falak al-tadwīr} ist der \textit{Epizykel} der}{}
\EdL{22}{AK07-J0025-L22}{Planeten\,; \textit{lieztekim asmir} könnte, wie \Sperr{Nallino} vermutet, heissen\,: \textit{liyastakīm as-sair},}{}
\EdL{23}{AK07-J0025-L23}{oder \textit{al-masīr}, d. h. dass er den Lauf (des Planeten) bestimme, ordne, reguliere\,; in}{}
\EdL{24}{AK07-J0025-L24}{der Tat ist der Epizykel zu diesem Zwecke eingeführt worden.}{}
\EdLp{25}{AK07-J0025-L25}{In dem gleichen Ms. befindet sich auf fol. 127\textsuperscript{v} (Tafel 55) am obern Rand}{}
\EdL{26}{AK07-J0025-L26}{folgende Note\,: \textit{Cullel. kavvekib. kumna. fiscerafehu. ille Otari. ahine unxi ademu.}}{}
\EdL{27}{AK07-J0025-L27}{\Sperr{Nallino} vermutete, das Ganze möchte ein Vers zur Ehre oder Lobpreisung irgend}{}
\EdL{28}{AK07-J0025-L28}{einer hohen Persönlichkeit (eines Fürsten) sein, und sollte lauten\,: \textit{Kullu’l-kawākibi}}{}
\EdL{29}{AK07-J0025-L29}{\textit{ḳumna fī sharafihi illā ʿOṭārida,} d. h. alle Planeten erhoben sich zu seiner Ehre}{}
\EdL{30}{AK07-J0025-L30}{ausser Merkur\,; die drei Wörter aber \textit{ahine unxi ademu} blieben ihm unklar. — Bei}{}
\EdL{31}{AK07-J0025-L31}{einem Besuche im Herbst 1913 bei Herrn Prof. \Sperr{Seybold} in Tübingen legte ich ihm}{}
\EdL{32}{AK07-J0025-L32}{diese Stelle ebenfalls vor, sie erinnerte ihn sofort an einige Ḳorānstellen\textsuperscript{1}), wo von}{}
\EdL{33}{AK07-J0025-L33}{den Engeln die Rede ist, die alle ausser \textit{Iblīs} (dem Satan) vor \textit{Adam} niederfielen,}{}
\EdL{34}{AK07-J0025-L34}{und er stellte daher die wohl mit Recht begründete Vermutung auf, das \textit{a} von \textit{ahine}}{}
\EdL{35}{AK07-J0025-L35}{sei wohl verdorben aus \textit{d} und gehöre ans Ende von \textit{Otari}, das \textit{x} in \textit{unxi} sei nach}{}
\EdL{36}{AK07-J0025-L36}{spanischer Transskriptionsart als \textit{sh} zu lesen, so dass also die drei letzten Wörter}{}
\EdL{37}{AK07-J0025-L37}{zu lesen wären\,: \textit{ḥīna unshi’a ādamu}, also der ganze Satz in halb klassischer, halb}{}
\EdL{38}{AK07-J0025-L38}{vulgärer Aussprache\,: \textit{kull al-kawākib ḳumna fī sharafihi illā ʿOṭārid ḥīna unshi’a}}{}
\EdL{39}{AK07-J0025-L39}{\textit{ādamu}, d. h. alle Planeten erhoben sich zu seiner (Adams) Ehre ausser Merkur, zur}{}
\EdL{40}{AK07-J0025-L40}{Zeit da Adam erschaffen wurde. — Dass hier \textit{Merkur} die Stelle des \textit{Satans} (\textit{Iblīs})}{}
\EdL{41}{AK07-J0025-L41}{vertritt, hat seinen Grund wohl darin, dass im Mittelalter der griechiche \textit{Hermes} mit}{}
\NoteRule
\EdN{N1}{AK07-J0025-N01}{\textsuperscript{1}) Vergl. die Suren II, 32\,; VII, 11\,; XX, 116.}
\FrameHead{0026}{p. XXII}{0026}{p. XXII}
\EdL{1}{AK07-J0026-L01}{dem \textit{Hermes Trismegistos} oder \textit{Thoth} der Aegypter identifiziert worden ist, und dieser}{}
\EdL{2}{AK07-J0026-L02}{war der Schutzgott der Alchemisten, Astrologen und Magier. In den Mirakel- und}{}
\EdL{3}{AK07-J0026-L03}{Zauberbüchern des christlichen Mittelalters wird als einer der sieben Fürsten der}{}
\EdL{4}{AK07-J0026-L04}{Unterwelt \textit{Ophiel} genannt\,; dieses Wort ist sehr wahrscheinlich zusammengesetzt aus}{}
\EdL{5}{AK07-J0026-L05}{dem griech. \textit{ophis} = Schlange, und dem hebr. \textit{el} = Gott, also Schlangengott =}{}
\EdL{6}{AK07-J0026-L06}{\textit{Hermes}\,; verbunden mit dem \textit{megistos} des Wortes \textit{Trismegistos} wurde \textit{Ophiel} zu \textit{Megi-}}{}
\EdL{7}{AK07-J0026-L07}{\textit{stophiel,} und hieraus entstand das verderbte oder absichtlich entstellte \textit{Mephistopheles.}}{}
\EdL{8}{AK07-J0026-L08}{Es ist dies wohl die beste Etymologie für dieses Wort\textsuperscript{1}). .}{}
\EdLp{9}{AK07-J0026-L09}{Wieder im Ms. O stehen am obern Rand von fol. 132\textsuperscript{v} (in unserm Text p. 20)}{}
\EdL{10}{AK07-J0026-L10}{die vier Wörter\,: \textit{eltabeb, helhacib, elhaazem, el-munegim}\,; das erste, zweite und vierte}{}
\EdL{11}{AK07-J0026-L11}{sind wohl zu lesen\,: \textit{al-ṭabīb, al-ḥāsib, al-munaddjim}, und bedeuten dann\,: der Arzt,}{}
\EdL{12}{AK07-J0026-L12}{der Rechner, der Astronom\,; das dritte weiss ich nicht zu deuten. Was der Schreiber}{}
\EdL{13}{AK07-J0026-L13}{mit diesen Wörtern sagen will, ist mir unklar, vielleicht meint er, dass das folgende}{}
\EdL{14}{AK07-J0026-L14}{Kapitel\,: \textit{Quot gradibus in qualibet regione quodlibet oriatur signum} für diese Berufs-}{}
\EdL{15}{AK07-J0026-L15}{arten von Wichtigkeit sei.}{}
\EdLp{16}{AK07-J0026-L16}{Im Ms. C stehen auf fol. 83\textsuperscript{r} (p. 27 unseres Textes) nach den Worten \textit{et accipies}}{}
\EdL{17}{AK07-J0026-L17}{\textit{quod contra invenies} die folgenden\,: \textit{wakhod mebibelche}\,; \Sperr{Nallino} glaubt und wohl}{}
\EdL{18}{AK07-J0026-L18}{mit Recht, das letzte Wort müsse gelesen werden \textit{mekibelche} und dann hiesse der}{}
\EdL{19}{AK07-J0026-L19}{Satz\,: \textit{wakhudh mā ḳibālahā}\textsuperscript{2}), d. h. und nimm was ihm (d. h. dem Argument) gegen-}{}
\EdL{20}{AK07-J0026-L20}{über steht\,; also einfach der arabische Text zu der vorhergehenden Übersetzung}{}
\EdL{21}{AK07-J0026-L21}{\textit{et accipies quod contra invenies.}}{}
\EdLp{22}{AK07-J0026-L22}{In den Titeln der Tafeln 27, 33, 39 u. 45 steht in beiden Mss. C u. O jeweilen}{}
\EdL{23}{AK07-J0026-L23}{nach dem \textit{arabischen} Namen des Planeten noch der \textit{persische}, nur auf Tafel 51 fehlt}{}
\EdL{24}{AK07-J0026-L24}{der persische für Merkur. Diejenigen für Saturn (\textit{Keiwen} = \textit{Kaiwān}), Jupiter (\textit{Hormi}}{}
\EdL{25}{AK07-J0026-L25}{= \textit{Hormuz}) und Mars (\textit{Baheram} = \textit{Bahrām}) habe ich wohl erkannt, die für Venus\textsuperscript{3})}{}
\EdL{26}{AK07-J0026-L26}{aber sind so verdorben, dass mir die Deutung nicht gelingen wollte. Auf fol. 122\textsuperscript{v}}{}
\EdL{27}{AK07-J0026-L27}{(Tafel 45) des Ms. O stehen nach dem arabischen Namen für Venus (\textit{elzohara}) folgende}{}
\EdL{28}{AK07-J0026-L28}{Worte\,: \textit{beied hakht eneid,} auf fol. 65\textsuperscript{r} des Ms. C steht nach demselben arabischen}{}
\EdL{29}{AK07-J0026-L29}{Worte\,: \textit{bei ed halt eneid.} Ich vermutete wohl im letzten Worte \textit{eneid} den mir bekannten}{}
\EdL{30}{AK07-J0026-L30}{persischen Namen für Venus \textit{nāhīd}, mich störte aber das \textit{e} vor \textit{neid}, ich wusste}{}
\EdL{31}{AK07-J0026-L31}{nämlich nicht, dass der persische Name auch \textit{anāhīd} geschrieben wird\,; gänzlich}{}
\EdL{32}{AK07-J0026-L32}{unverständlich aber war mir \textit{beied hakht} oder \textit{bei ed halt.} Da führte mich auch}{}
\EdL{33}{AK07-J0026-L33}{wieder das Wissen \Sperr{Nallinos} in diesen Dingen auf den richtigen Weg\,: \textit{beied hakht}}{}
\EdL{34}{AK07-J0026-L34}{\textit{eneid} ist verdorben aus \textit{bēdukht anāhīd}\,; dies sind Namen zweier alter iranischer}{}
\EdL{35}{AK07-J0026-L35}{weiblicher Gottheiten, die später dann auf den Planeten Venus übertragen worden}{}
\EdL{36}{AK07-J0026-L36}{sind. — Es ist sehr wahrscheinlich, dass diese persischen Namen für die Planeten}{}
\EdL{37}{AK07-J0026-L37}{von \Sperr{Al-Khwārizmī} herstammen, denn \Sperr{Maslama} wird kaum persisch verstanden}{}
\EdL{38}{AK07-J0026-L38}{haben.}{}
\NoteRule
\EdN{N1}{AK07-J0026-N01}{\textsuperscript{1}) Vergl. J. \Sperr{Goebel} in \textit{Beilage zur Allgem. Zeitung}, 1905, Nr. 195.}
\EdN{N2}{AK07-J0026-N02}{\textsuperscript{2}) Für den Leser, dem das Arabische unbekannt ist, bemerken wir, dass die Westaraber das}
\EdN{N3}{AK07-J0026-N03}{\textit{ā} wie \textit{ē} ausgesprochen haben, und jetzt in Marokko noch so sprechen (\textit{Fēs} für \textit{Fās}).}
\EdN{N4}{AK07-J0026-N04}{\textsuperscript{3}) Für diesen Planeten hatten die Perser zwei verschiedene Namen.}
\FrameHead{0027}{p. XXIII}{0027}{p. XXIII}
\EdLp{1}{AK07-J0027-L01}{Über die Grundsätze, die mich bei der Wiedergabe des Textes und der Tafeln}{}
\EdL{2}{AK07-J0027-L02}{geleitet haben, muss ich folgendes bemerken\,:}{}
\EdLp{3}{AK07-J0027-L03}{Im Texte sind die Wörter wiedergegeben, wie sie in den Handschriften stehen,}{}
\EdL{4}{AK07-J0027-L04}{mit den Ausnahmen, dass erstens für die Genitiv- und Dativendung der ersten Deklin.}{}
\EdL{5}{AK07-J0027-L05}{Sing. das heutige \textit{ae} statt des mittelalterlichen \textit{e} geschrieben worden ist, zweitens}{}
\EdL{6}{AK07-J0027-L06}{für \textit{u}, wenn es vor einem Vokal steht, also \textit{v} gelesen wird, stets dieser Buchstabe}{}
\EdL{7}{AK07-J0027-L07}{gesetzt worden ist\,; drittens ist in beiden Mss. für \textit{t} vor \textit{i} und folgendem Vokal oft}{}
\EdL{8}{AK07-J0027-L08}{(nicht immer) \textit{c} gesetzt, wir haben immer \textit{t} geschrieben.}{}
\EdLp{9}{AK07-J0027-L09}{Der Text und die Tafeln enthalten nur römische Ziffern, mit Ausnahme einiger}{}
\EdL{10}{AK07-J0027-L10}{Zusätze aus späterer Zeit in O und N. Im Text habe ich die römischen Ziffern}{}
\EdL{11}{AK07-J0027-L11}{beibehalten mit Ausnahme des Zusatzes von späterer Hand zu Kap. 21, der sehr viele}{}
\EdL{12}{AK07-J0027-L12}{und grössere Zahlen enthält und nur in C vorkommt. In den Tafeln habe ich, um}{}
\EdL{13}{AK07-J0027-L13}{den Druck und das Lesen zu erleichtern, alle Ziffern durch arabische ersetzt.}{}
\EdLp{14}{AK07-J0027-L14}{Für die Null kommen im Text sowohl als in den Tafeln drei Zeichen vor\,:}{}
\EdL{15}{AK07-J0027-L15}{C und O haben bald τ (= \textit{teca})\textsuperscript{1}), bald θ (= \textit{theta}\,?), N hat stets \Ovl{0}\textsuperscript{2}). Interessant ist,}{}
\EdL{16}{AK07-J0027-L16}{dass in N an verschiedenen Stellen, von denen ich in den Fussnoten einige erwähnt}{}
\EdL{17}{AK07-J0027-L17}{habe, für 38 (einige male auch für 28) geschrieben steht XXX0 (bezw. XX0)\,; um die}{}
\EdL{18}{AK07-J0027-L18}{Zahl in römischen Ziffern nicht zu gross werden zu lassen, wurde also VIII durch}{}
\EdL{19}{AK07-J0027-L19}{0 ersetzt, es ist dies jedenfalls der Anfangsbuchstabe von \textit{octo}.}{}
\EdLp{20}{AK07-J0027-L20}{Was die Tafeln anbetrifft, so weisen alle Handschriften mit Ausnahme von N}{}
\EdL{21}{AK07-J0027-L21}{in den Titeln und Kolonnenüberschriften eine unglaubliche Willkür und Inkonsequenz}{}
\EdL{22}{AK07-J0027-L22}{auf\,; bald steht sowohl der arabische als auch der lateinische Titel, bald nur der}{}
\EdL{23}{AK07-J0027-L23}{arabische, bald nur der lateinische, bald gar kein Titel. Auch die Kolonnenüber-}{}
\EdL{24}{AK07-J0027-L24}{schriften lauten oft verschieden, obgleich sie dasselbe bezeichnen sollen, so steht z. B}{}
\EdL{25}{AK07-J0027-L25}{im gleichen Ms. in der einen Tafel \textit{menses arabici}, in der andern \textit{menses arabum}, in}{}
\EdL{26}{AK07-J0027-L26}{der einen \textit{in annis collectis}, in der andern bloss \textit{in collectis}, in der einen \textit{via}, in der}{}
\EdL{27}{AK07-J0027-L27}{andern \textit{semita}, in einer dritten \textit{catar} (= \textit{linea}), in der einen \textit{recedentis}, in der andern}{}
\EdL{28}{AK07-J0027-L28}{\textit{descendentis}, u. s. w.\,; oft sind auch Kolonnentitel ganz weggelassen, besonders in C.}{}
\EdL{29}{AK07-J0027-L29}{Das grösste Durcheinander herrscht in den Bezeichnungen der kleinern Winkelteile\,:}{}
\EdL{30}{AK07-J0027-L30}{die eine Tafel hat \textit{Min. Sec.}, die andere \textit{Dak.}\textsuperscript{3}) \textit{Then.}\textsuperscript{4}), eine dritte \textit{Dak. Sec.}, eine}{}
\EdL{31}{AK07-J0027-L31}{vierte (besonders in C, aber nicht durchweg) \textit{P.}\textsuperscript{5}) \textit{M.}\textsuperscript{6}), und zwar diese Namen bald}{}
\EdL{32}{AK07-J0027-L32}{so abgekürzt, bald vollständig ausgeschrieben. Nur in der Bezeichnung der Tier-}{}
\EdL{33}{AK07-J0027-L33}{kreiszeichen (\textit{Signa}) und Grade (\textit{Gradus}) herrscht Übereinstimmung, nur sind sie}{}
\NoteRule
\EdN{N1}{AK07-J0027-N01}{\textsuperscript{1}) Vergl. M. \Sperr{Curtze}, \textit{Petri Philomeni de Dacia in Algorismum vulgarem Johannis de Sacro-}}
\EdN{N2}{AK07-J0027-N02}{\textit{bosco Commentarius}, Hauniae 1897, p. 2 u. 26.}
\EdN{N3}{AK07-J0027-N03}{\textsuperscript{2}) τ und \Ovl{0} befinden sich auch in dem \textit{Prologus N. Ocreati in Helceph ad Adelardum Batensem}}
\EdN{N4}{AK07-J0027-N04}{\textit{magistrum suum} (herausgeg. von C. \Sperr{Henry} in den Abhandlungen z. Geschichte der mathem. Wissen-}
\EdN{N5}{AK07-J0027-N05}{schaften III, p. 131 ff.). Ebenso steht im Ms. des Escurial, das \Sperr{Nallino} zu seiner Ausgabe \Sperr{Al-}}
\EdN{N6}{AK07-J0027-N06}{\Sperr{Battānīs} benutzt hat, in den Tafeln, wenn z. B. Minuten fehlen, dafür das Zeichen \Ovl{0}.}
\EdN{N7}{AK07-J0027-N07}{\textsuperscript{3}) \textbf{\textit{Dak.}} = \textbf{\textit{Dakaicae}} (v. Gloss. arab.).}
\EdN{N8}{AK07-J0027-N08}{\textsuperscript{4}) \textbf{\textit{Then.}} = \textbf{\textit{Theniae}} (v. Gloss. arab.).}
\EdN{N9}{AK07-J0027-N09}{\textsuperscript{5}) \textbf{\textit{P.}} = \textbf{\textit{Puncta}} für \textbf{\textit{Min.}}}
\EdN{N10}{AK07-J0027-N10}{\textsuperscript{6}) \textbf{\textit{M.}} = \textbf{\textit{Minuta}} für \textit{Sec.}}
\FrameHead{0028}{p. XXIV}{0028}{p. XXIV}
\EdL{1}{AK07-J0028-L01}{bisweilen mehr, bisweilen weniger abgekürzt. In diesem Punkte habe ich die Sache}{}
\EdL{2}{AK07-J0028-L02}{möglichst vereinfacht und stets Sig. Gr. Min. Sec. geschrieben.}{}
\EdLp{3}{AK07-J0028-L03}{Die Festsetzung des Textes und der kritische Apparat sind im grossen und}{}
\EdL{4}{AK07-J0028-L04}{ganzen das Werk \Sperr{Björnbos}. Dieser Gelehrte war, wie ich aus einem Briefe des}{}
\EdL{5}{AK07-J0028-L05}{Hrn. Prof. \Sperr{Heiberg} an mich erfahren habe, der Ansicht, dass der Text von C der}{}
\EdL{6}{AK07-J0028-L06}{beste sei, und darum hat er auch diesen im allgemeinen zu Grunde gelegt. Eine}{}
\EdL{7}{AK07-J0028-L07}{nochmalige Vergleichung von C u. O\textsuperscript{1}) hat mir aber gezeigt, dass an verschiedenen}{}
\EdL{8}{AK07-J0028-L08}{Stellen die Lesart von O die treffendere sei (einige male auch die von N), ich habe}{}
\EdL{9}{AK07-J0028-L09}{mir daher erlaubt, in solchen Fällen die bessere Lesart in den Text aufzunehmen.}{}
\EdL{10}{AK07-J0028-L10}{Ist in den Fussnoten die Lesart eines Ms. nicht angegeben, so stimmt sie mit der-}{}
\EdL{11}{AK07-J0028-L11}{jenigen des Textes überein.}{}
\EdLp{12}{AK07-J0028-L12}{Der Kommentar, der auf den Text folgt, ist ganz meine Arbeit. Ich fragte}{}
\EdL{13}{AK07-J0028-L13}{mich zuerst, ob es wohl nötig wäre, dem Text eine deutsche Übersetzung folgen zu}{}
\EdL{14}{AK07-J0028-L14}{lassen, da der Gegenstand oft sehr kurz und unklar dargelegt ist\,; allein ich sagte}{}
\EdL{15}{AK07-J0028-L15}{mir, dass damit dem Leser doch nur halb gedient wäre, denn wer mit der alten}{}
\EdL{16}{AK07-J0028-L16}{und mittelalterlichen Astronomie nicht bekannt ist, könnte mit einer blossen Über-}{}
\EdL{17}{AK07-J0028-L17}{setzung nicht viel anfangen, es müsste auch einer solchen ein Kommentar beigegeben}{}
\EdL{18}{AK07-J0028-L18}{werden\,; darum habe ich auf die Übersetzung verzichtet, und mich mit einem}{}
\EdL{19}{AK07-J0028-L19}{Kommentar begnügt, den ich dadurch nutzbarer zu machen suchte, dass ich zu den}{}
\EdL{20}{AK07-J0028-L20}{hauptsächlichsten astronomischen Problemen, die in diesem Werke behandelt werden,}{}
\EdL{21}{AK07-J0028-L21}{Beispiele gegeben habe, die nach den Regeln des Textes gelöst sind. — Im Kommentar}{}
\EdL{22}{AK07-J0028-L22}{habe ich die Titel der einzelnen Kapitel des Textes ziemlich wörtlich übersetzt\,; die}{}
\EdL{23}{AK07-J0028-L23}{Kapitel sind behufs bequemer Verweisung nummeriert, was in den Handschriften}{}
\EdL{24}{AK07-J0028-L24}{nicht der Fall ist.}{}
\EdLp{25}{AK07-J0028-L25}{Ich muss gestehen, dass dieser Kommentar eine schwierige und mühevolle}{}
\EdL{26}{AK07-J0028-L26}{Arbeit für mich war, und wer die einzelnen Kapitel des Textes durchliest, wird dies}{}
\EdL{27}{AK07-J0028-L27}{wohl auch erklärlich finden. Da hat mir nun die ausgezeichnete Ausgabe \Sperr{Al-Bat-}}{}
\EdL{28}{AK07-J0028-L28}{\Sperr{tānīs} und der Kommentar dazu von C. A. \Sperr{Nallino} und G. \Sperr{Schiaparelli}\textsuperscript{2}) die}{}
\EdL{29}{AK07-J0028-L29}{wertvollsten Dienste geleistet, ohne diese Ausgabe wäre mir die Kommentierung des}{}
\EdL{30}{AK07-J0028-L30}{vorliegenden Werkes kaum möglich gewesen\,; es sind aber trotz dieser Hilfe noch}{}
\EdL{31}{AK07-J0028-L31}{verschiedene Punkte unklar geblieben, zu deren Aufhellung vielleicht, wie ich hoffe,}{}
\EdL{32}{AK07-J0028-L32}{in diesen Fragen bewanderte Gelehrte ihre Unterstützung leihen werden.}{}
\EdLp{33}{AK07-J0028-L33}{Ich habe mir erlaubt, im Kommentar für die drei am meisten auftretenden}{}
\EdL{34}{AK07-J0028-L34}{Eigennamen Abkürzungen anzuwenden, da dadurch auch der Druck wesentlich}{}
\EdL{35}{AK07-J0028-L35}{erleichtert wird\,; so schreibe ich für \Sperr{Al-Khwārizmī} stets \Sperr{Kh.}, für \Sperr{Al-Battānī}}{}
\EdL{36}{AK07-J0028-L36}{\Sperr{Batt.}, für \Sperr{Ptolemäus} \Sperr{Ptol.}}{}
\NoteRule
\EdN{N1}{AK07-J0028-N01}{\textsuperscript{1}) \Sperr{Björnbo} hat dies auch noch beabsichtigt gehabt, was verschiedene Notizen in seinem Ms.}
\EdN{N2}{AK07-J0028-N02}{beweisen\,; leider konnte er aber eine zweite Lesung nicht mehr ausführen, ich war daher öfters}
\EdN{N3}{AK07-J0028-N03}{genötigt, an seiner Wiedergabe der Lesarten Änderungen vorzunehmen.}
\EdN{N4}{AK07-J0028-N04}{\textsuperscript{2}) Ich zitiere dieses Werk, das von 1899—1907 in drei Teilen unter dem Titel\,: \textit{Al-Battānī}}
\EdN{N5}{AK07-J0028-N05}{\textit{sive Albatenii Opus astronomicum, edit. a} C. A. \Sperr{Nallino}, in den \textit{Pubblicazioni del R. Osservatorio}}
\EdN{N6}{AK07-J0028-N06}{\textit{di Brera in Milano}, N. XL, erschienen ist, künftig stets mit \Sperr{Batt.}, ed. \Sperr{Nallino}, oder auch}
\EdN{N7}{AK07-J0028-N07}{nur \Sperr{Batt.}}
\FrameHead{0029}{p. XXV}{0029}{p. XXV}
\EdLp{1}{AK07-J0029-L01}{Das Werk schliessen zwei Indices ab, ein \textit{arabisches Glossar} und ein allgemeines}{}
\EdL{2}{AK07-J0029-L02}{\textit{Namen- und Sachregister}. Der Grundstock des erstern ist von Dr. R. O. \Sperr{Besthorn}}{}
\EdL{3}{AK07-J0029-L03}{in Kopenhagen verfasst, dem mit J. L. \Sperr{Heiberg} zusammen wir die Herausgabe der}{}
\EdL{4}{AK07-J0029-L04}{ältesten arabischen Euklid-Übersetzung verdanken. Verschiedenes ist nachher von}{}
\EdL{5}{AK07-J0029-L05}{mir ergänzt worden, so besonders die richtigen Formen der ägyptischen und per-}{}
\EdL{6}{AK07-J0029-L06}{sischen Monatsnamen. Die arabischen Namen gebe ich nur in Transskription, zuerst}{}
\EdL{7}{AK07-J0029-L07}{wie sie im Texte stehen, und dann in der richtigen Form\,; hierbei halte ich mich}{}
\EdL{8}{AK07-J0029-L08}{an die Regeln der \textit{Enzyklopädie des Islām}, die gegenwärtig unter der Leitung von}{}
\EdL{9}{AK07-J0029-L09}{Prof. \Sperr{Houtsma} in Utrecht u. A. erscheint\,; für den Leser will ich hier nur anführen,}{}
\EdL{10}{AK07-J0029-L10}{dass \textit{dj} wie \textit{dsch}, \textit{z} wie weiches \textit{s}, \textit{kh} wie alemannisches \textit{ch}, \textit{y} wie \textit{j}, \textit{dh} u. \textit{th} wie}{}
\EdL{11}{AK07-J0029-L11}{weiches u. hartes englisches \textit{th} auszusprechen sind\,; die übrigen Zeichen bedürfen}{}
\EdL{12}{AK07-J0029-L12}{keiner Erklärung. Die Bedeutung der arabischen Wörter habe ich mit Zustimmung}{}
\EdL{13}{AK07-J0029-L13}{des Hrn. Dr. \Sperr{Besthorn}, der sie ursprünglich in lateinischer Sprache beigefügt hat,}{}
\EdL{14}{AK07-J0029-L14}{in deutscher Sprache wiedergegeben. — Was die Benützung des Glossars anbetrifft,}{}
\EdL{15}{AK07-J0029-L15}{so mache ich darauf aufmerksam, dass bei allen Wörtern, die mit dem Artikel \textit{al}}{}
\EdL{16}{AK07-J0029-L16}{oder \textit{el} beginnen, dieser wegzulassen ist, man suche also \textit{elmuht} unter \textit{muht}, \textit{eltheniae}}{}
\EdL{17}{AK07-J0029-L17}{unter \textit{theniae}, etc. — Bei sehr häufig wiederkehrenden Wörtern, wie \textit{dakaicae}}{}
\EdL{18}{AK07-J0029-L18}{(= minutae), \textit{theniae} (= secundae) und einigen andern, lasse ich die Verweisung}{}
\EdL{19}{AK07-J0029-L19}{auf die Seiten weg, mit Ausnahme derjenigen, wo ihre Bedeutung erklärt wird.}{}
\EdLp{20}{AK07-J0029-L20}{Für beide Indices gilt folgendes\,: Zahlen ohne vorangehenden Buchstaben be-}{}
\EdL{21}{AK07-J0029-L21}{zeichnen die Seitenzahl der Einleitung, des Textes und des Kommentars, solche mit}{}
\EdL{22}{AK07-J0029-L22}{vorangehendem T. die Nummern der Tafeln, N in Klammer bedeutet „Note“.}{}
\FrameHead{0030}{p. [XXVI]}{0030}{p. [XXVI]}
\Withheld{blank}
\FrameHead{0031}{p. [XXVII]}{0031}{p. [XXVII]}
\EdLc{1}{AK07-J0031-L01}{TEXT}{}
\FrameHead{0032}{p. [XXVIII]}{0032}{p. [XXVIII]}
\Withheld{blank}
\FrameHead{0064}{p. [32]}{0064}{p. [32]}
\EdLh{1}{AK07-J0064-L01}{KOMMENTAR}{}
\EdLc{2}{AK07-J0064-L02}{\textbf{Einleitung des Textes.}}{}
\EdLc{3}{AK07-J0064-L03}{(Text p. 1—2.)}{}
\EdLp{4}{AK07-J0064-L04}{Das arabische Wort \textit{ezich} = \textit{elzīdj} (gesprochen \textit{ez-zīdj}) ist, wie noch verschiedene}{}
\EdL{5}{AK07-J0064-L05}{andere arabische technische Ausdrücke, nicht übersetzt\,; es ist ursprünglich ein per-}{}
\EdL{6}{AK07-J0064-L06}{sisches Wort (\textit{zīg}) und bedeutet die Grundfäden (Zettel) des Zeuges, in die der Weber}{}
\EdL{7}{AK07-J0064-L07}{oder Sticker den Einschlag oder die Stickerei hineinwirkt, dann diesem entsprechend}{}
\EdL{8}{AK07-J0064-L08}{„astronomisches Tafelwerk“. Solche gab es in der arabischen astronomischen Litte-}{}
\EdL{9}{AK07-J0064-L09}{ratur eine grosse Zahl\,: Die ältesten waren die nach den indischen \textit{Siddhāntas} (s. Ein-}{}
\EdL{10}{AK07-J0064-L10}{leitung) verfassten Tafeln von \Sperr{Ibrāhīm al-Fazārī} und \Sperr{Yaʿḳūb b. Ṭāriḳ}\,; dann}{}
\EdL{11}{AK07-J0064-L11}{kommen die beiden Ausgaben des \textit{Sindhind} (arab. Schreibweise des indischen \textit{Siddhānta})}{}
\EdL{12}{AK07-J0064-L12}{des \Sperr{Muh. b. Mūsā al-Khwārizmī}, die wahrscheinlich in den Jahren 800—830}{}
\EdL{13}{AK07-J0064-L13}{n. Chr. verfasst worden sind\,; hierauf die „erprobten Tafeln“ der Astronomen \Sperr{Al-}}{}
\EdL{14}{AK07-J0064-L14}{\Sperr{Ma’mūns} (\Sperr{Yaḥyā b. Abī Manṣūr}, \Sperr{Sened b. ʿAlī}, \Sperr{ʿAbbās b. Saʿīd al-Djauharī}}{}
\EdL{15}{AK07-J0064-L15}{u. a.)\textsuperscript{1}). — Es ist möglich, dass die spätere Ausgabe des \textit{Sindhind} von \Sperr{Kh.} schon}{}
\EdL{16}{AK07-J0064-L16}{etwas beeinflusst gewesen sein mag durch die Beobachtungen der Astronomen}{}
\EdL{17}{AK07-J0064-L17}{\Sperr{Al-Ma’mūns}, und sich deshalb etwas weiter von den indischen Vorlagen entfernt}{}
\EdL{18}{AK07-J0064-L18}{hat als die erste Ausgabe, und dass wahrscheinlich, wie wir aus dem Inhalt des}{}
\EdL{19}{AK07-J0064-L19}{vorliegenden Werkes schliessen möchten, \Sperr{Maslama al-Madjrīṭī} (s. Einleitung) diese}{}
\EdL{20}{AK07-J0064-L20}{zweite Ausgabe einer Umarbeitung unterzogen hat. Doch lässt sich selbstverständlich}{}
\EdL{21}{AK07-J0064-L21}{hierüber nichts sicheres sagen, wir stellen nur Vermutungen auf, die uns als berechtigt}{}
\EdL{22}{AK07-J0064-L22}{erscheinen. — Welchen der verschiedenen indischen \textit{Siddhāntas} \Sperr{Kh.} seinen Tafeln zu}{}
\EdL{23}{AK07-J0064-L23}{Grunde gelegt habe, ist ebenfalls nicht sicher zu entscheiden\,; bis jetzt hat man es}{}
\EdL{24}{AK07-J0064-L24}{als ziemlich wahrscheinlich betrachtet, dass es der \textit{Brahma-Siddhānta} (verfasst von}{}
\EdL{25}{AK07-J0064-L25}{\Sperr{Brahmagupta} in der ersten Hälfte des 7. Jahrh. n. Chr.) gewesen sei\,; wir werden}{}
\EdL{26}{AK07-J0064-L26}{aber weiter unten sehen, dass auch andere \textit{Siddhāntas} in Frage kommen können,}{}
\EdL{27}{AK07-J0064-L27}{besonders der \textit{Sūrya-Siddhānta}. Dies würde dann darauf hindeuten, dass vielleicht}{}
\EdL{28}{AK07-J0064-L28}{\Sperr{Kh.} auch die persischen „Tafeln des Shāh“ (\textit{zīdj al-shāh}, oder \textit{zīdj al-shahriyār})\textsuperscript{2}) zu}{}
\EdL{29}{AK07-J0064-L29}{Rate gezogen hat, die im 2. Jahrh. der \textit{Hidjra} ins Arabische übersetzt worden sind,}{}
\NoteRule
\EdN{N1}{AK07-J0064-N01}{\textsuperscript{1}) Vergl. für diese Astronomen\,: H. \Sperr{Suter}, \textit{Die Mathematiker und Astronomen der Araber}, etc.,}
\EdN{N2}{AK07-J0064-N02}{p. 3—14, und unsere Einleitung zu dieser Ausgabe.}
\EdN{N3}{AK07-J0064-N03}{\textsuperscript{2}) Vergl. \Sperr{Batt.}, ed. \Sperr{Nallino}, I, 218.}
\FrameHead{0065}{p. 33}{0065}{p. 33}
\EdL{1}{AK07-J0065-L01}{und sich wahrscheinlich auf den \textit{Sūrya-Siddhānta} stützten. Besonders der Wert des}{}
\EdL{2}{AK07-J0065-L02}{Apogeums der Sonne und die Art der Berechnung der Gleichungen der Planeten}{}
\EdL{3}{AK07-J0065-L03}{könnten aus diesen persischen Tafeln genommen sein, sagt doch \Sperr{Ibn al-Ḳifṭī}\textsuperscript{1}) in}{}
\EdL{4}{AK07-J0065-L04}{seiner Biographie \Sperr{Al-Fazārīs} nach \Sperr{Ibn al-Adamī}\textsuperscript{2}) folgendes\,: „Er (\Sperr{Kh.}) stützte}{}
\EdL{5}{AK07-J0065-L05}{sich in seinen Tafeln auf die mittlern Örter des \textit{Sindhind}, wich aber von ihm ab}{}
\EdL{6}{AK07-J0065-L06}{in den Gleichungen (der Planeten) und in der Schiefe (der Ekliptik)\,; er setzte seine}{}
\EdL{7}{AK07-J0065-L07}{Gleichungen fest nach der Methode der \Sperr{Perser}, und die Deklination der Sonne}{}
\EdL{8}{AK07-J0065-L08}{nach der Methode des \Sperr{Ptolemäus}.“ Da die letztere Behauptung, wie wir unten}{}
\EdL{9}{AK07-J0065-L09}{sehen werden, richtig ist, so mag wohl auch die erste wahr sein, nachweisen können}{}
\EdL{10}{AK07-J0065-L10}{wir dies nicht, da diese persischen Tafeln nicht mehr vorhanden sind.}{}
\EdLp{11}{AK07-J0065-L11}{Nach der Einleitung des Textes beziehen sich die Daten der Tafeln für die}{}
\EdL{12}{AK07-J0065-L12}{mittlern und wahren Örter von Sonne, Mond und Planeten (oder wie die Araber}{}
\EdL{13}{AK07-J0065-L13}{sagen, der sieben Planeten) auf den Meridian von \textit{Arin} (falsche Schreibweise für}{}
\EdL{14}{AK07-J0065-L14}{\textit{Azīn} oder \textit{Uzain}, und dieses wieder die arabische Form für das indische \textit{Udjain}, in}{}
\EdL{15}{AK07-J0065-L15}{Sanskrit \textit{Ujjayini}), im nördlichen Vorder-Indien gelegen, (g. B. 23° 11′, g. L. 75° 52′ Gr.)}{}
\EdL{16}{AK07-J0065-L16}{Was die Zeit anbetrifft, worauf sich die \textit{Radix} der verschiedenen Tafeln jeweilen}{}
\EdL{17}{AK07-J0065-L17}{bezieht, so ist als Anfangszeitpunkt angegeben der Mittag des 1. \textit{Muḥarram} des Jahres 1}{}
\EdL{18}{AK07-J0065-L18}{der Hidjra (d. i. der 15. Juli 622 n. Chr.)\,; z. B. in Tafel 4 steht als Radix für die}{}
\EdL{19}{AK07-J0065-L19}{Sonne 3\textsuperscript{s} 23° 25′ 48″ (= 113° 25′ 48″), d. h. dies war die mittlere Länge der Sonne am}{}
\EdL{20}{AK07-J0065-L20}{Mittag (astron. Beginn des Tages bei den meisten arab. Astronomen) des 1. \textit{Muḥarram}}{}
\EdL{21}{AK07-J0065-L21}{des Jahres 1 der \textit{Hidjra} für den Meridian von \textit{Arin}. \Sperr{Maslama al-Madjrīṭī} hat also,}{}
\EdL{22}{AK07-J0065-L22}{wie wir schon in der Einleitung bemerkt haben, die \textit{Aera} \Sperr{Yezdegirds} (Beginn am}{}
\EdL{23}{AK07-J0065-L23}{16. Juni 632 n. Chr.), die \Sperr{Kh.} zu Grunde gelegt hatte, durch die muhammedanische}{}
\EdL{24}{AK07-J0065-L24}{ersetzt, den Meridian von \textit{Arin} aber hat er beibehalten\,; dies gilt aber nur für die}{}
\EdL{25}{AK07-J0065-L25}{Tafeln der mittlern und wahren Bewegung der Sonne, des Mondes und der fünf}{}
\EdL{26}{AK07-J0065-L26}{Planeten\,; wir werden später sehen, dass \Sperr{Maslama} eine Reihe von Tafeln (besonders}{}
\EdL{27}{AK07-J0065-L27}{die der Konjunktionen, Oppositionen und Finsternisse) auf den Meridian von \textit{Cordova}}{}
\EdL{28}{AK07-J0065-L28}{bezogen hat. Es ist also die Möglichkeit nicht ausgeschlossen, dass diese Tafeln}{}
\EdL{29}{AK07-J0065-L29}{ursprünglich nicht im Werke \Sperr{Kh.}’s enthalten waren\,; hierauf hindeutende Angaben}{}
\EdL{30}{AK07-J0065-L30}{hat ja, wie wir in der Einleitung gesehen haben, der spanische Gelehrte \Sperr{Ibn Ṣāʿid}}{}
\EdL{31}{AK07-J0065-L31}{gemacht, und \Sperr{Ibn al-Ḳifṭī}\textsuperscript{3}) hat sie wiederholt\,; doch wäre in diesem Falle das}{}
\EdL{32}{AK07-J0065-L32}{Werk \Sperr{Kh.}’s auch gar zu unbedeutend gewesen, denn gerade diese Tafeln gehören}{}
\EdL{33}{AK07-J0065-L33}{zu den wichtigsten eines astronomischen Werkes.}{}
\EdLp{34}{AK07-J0065-L34}{In dieser Einleitung tritt auch der einzige Eigenname auf, den ausser denen}{}
\EdL{35}{AK07-J0065-L35}{des Verfassers und des Übersetzers diese Tafeln enthalten, der Name „\Sperr{Ptolomaeus}“,}{}
\EdL{36}{AK07-J0065-L36}{Man sieht also, was sich auch aus dem Studium der Tafeln ergibt, dass \Sperr{Kh.} schon,}{}
\EdL{37}{AK07-J0065-L37}{wie dann die spätern Astronomen, insbesondere \Sperr{Batt.}, sich auf den \textit{Almagest} des}{}
\EdL{38}{AK07-J0065-L38}{\Sperr{Ptolemäus} stützte, und also auch durch diesen wesentlich beeinflusst worden ist,}{}
\EdL{39}{AK07-J0065-L39}{was sich im Laufe unserer Betrachtung an verschiedenen Stellen zeigen wird.}{}
\NoteRule
\EdN{N1}{AK07-J0065-N01}{\textsuperscript{1}) Vergl. die Ausgabe von \Sperr{Lippert}, Leipzig, 1903, p. 270—271.}
\EdN{N2}{AK07-J0065-N02}{\textsuperscript{2}) Vergl. H. \Sperr{Suter}, \textit{Die Mathematiker und Astronomen der Araber}, etc. p. 44.}
\EdN{N3}{AK07-J0065-N03}{\textsuperscript{3}) Edid. \Sperr{Lippert}, p. 326.}
\FrameHead{0066}{p. 34}{0066}{p. 34}
\EdLc{1}{AK07-J0066-L01}{\textbf{Kap. 1. Beschreibung des arabischen Jahres.}}{}
\EdLc{2}{AK07-J0066-L02}{(Text p. 2.)}{}
\EdLp{3}{AK07-J0066-L03}{Das arabische Mondjahr hat ungefähr 354\textsuperscript{11}/\textsubscript{30} Tage\,; es ist in 12 Monate mit}{}
\EdL{4}{AK07-J0066-L04}{abwechselnd 30 und 29 Tagen geteilt\,; um den Bruchteil eines Tages auszugleichen,}{}
\EdL{5}{AK07-J0066-L05}{werden bisweilen dem letzten Monat, der sonst 29 Tage zählt, 30 Tage gegeben\,; in}{}
\EdL{6}{AK07-J0066-L06}{welchen Jahren dies geschieht, sagt \Sperr{Kh.} nicht, es sind dies das 2., 5., 7., 10., 13., 16.,}{}
\EdL{7}{AK07-J0066-L07}{18., 21., 24., 26. u. 29. in einem Zyklus von 30 Jahren. Ein solches Jahr mit 355}{}
\EdL{8}{AK07-J0066-L08}{Tagen heisst \textit{sana kabīsa} = Schaltjahr. Die arabischen Monate heissen (\Sperr{Kh.} nennt}{}
\EdL{9}{AK07-J0066-L09}{im Text bloss die zwei ersten und den letzten)\,:}{}
\EdL{10}{AK07-J0066-L10}{\parbox[t]{\dimexpr0.2450\linewidth-0.5em\relax}{\raggedright\small\mbox{}1) \textit{Muḥarram}}\hspace{0.5em}\parbox[t]{\dimexpr0.2450\linewidth-0.5em\relax}{\raggedright\small\mbox{}= 30 Tage}\hspace{0.5em}\parbox[t]{\dimexpr0.2450\linewidth-0.5em\relax}{\raggedright\small\mbox{}7) \textit{Radjab}}\hspace{0.5em}\parbox[t]{\dimexpr0.2450\linewidth-0.5em\relax}{\raggedright\small\mbox{}= 30 Tage}\hspace{0.5em}}{}
\EdL{11}{AK07-J0066-L11}{\parbox[t]{\dimexpr0.2450\linewidth-0.5em\relax}{\raggedright\small\mbox{}2) \textit{Ṣafar}}\hspace{0.5em}\parbox[t]{\dimexpr0.2450\linewidth-0.5em\relax}{\raggedright\small\mbox{}= 29 „}\hspace{0.5em}\parbox[t]{\dimexpr0.2450\linewidth-0.5em\relax}{\raggedright\small\mbox{}8) \textit{Shaʿbān}}\hspace{0.5em}\parbox[t]{\dimexpr0.2450\linewidth-0.5em\relax}{\raggedright\small\mbox{}= 29 „}\hspace{0.5em}}{}
\EdL{12}{AK07-J0066-L12}{\parbox[t]{\dimexpr0.2450\linewidth-0.5em\relax}{\raggedright\small\mbox{}3) \textit{Rabīʿ I}}\hspace{0.5em}\parbox[t]{\dimexpr0.2450\linewidth-0.5em\relax}{\raggedright\small\mbox{}= 30 „}\hspace{0.5em}\parbox[t]{\dimexpr0.2450\linewidth-0.5em\relax}{\raggedright\small\mbox{}9) \textit{Ramaḍān}}\hspace{0.5em}\parbox[t]{\dimexpr0.2450\linewidth-0.5em\relax}{\raggedright\small\mbox{}= 30 „}\hspace{0.5em}}{}
\EdL{13}{AK07-J0066-L13}{\parbox[t]{\dimexpr0.2450\linewidth-0.5em\relax}{\raggedright\small\mbox{}4) \textit{Rabīʿ II}}\hspace{0.5em}\parbox[t]{\dimexpr0.2450\linewidth-0.5em\relax}{\raggedright\small\mbox{}= 29 „}\hspace{0.5em}\parbox[t]{\dimexpr0.2450\linewidth-0.5em\relax}{\raggedright\small\mbox{}10) \textit{Shawwāl}}\hspace{0.5em}\parbox[t]{\dimexpr0.2450\linewidth-0.5em\relax}{\raggedright\small\mbox{}= 29 „}\hspace{0.5em}}{}
\EdL{14}{AK07-J0066-L14}{\parbox[t]{\dimexpr0.2450\linewidth-0.5em\relax}{\raggedright\small\mbox{}5) \textit{Djumādā I}}\hspace{0.5em}\parbox[t]{\dimexpr0.2450\linewidth-0.5em\relax}{\raggedright\small\mbox{}= 30 „}\hspace{0.5em}\parbox[t]{\dimexpr0.2450\linewidth-0.5em\relax}{\raggedright\small\mbox{}11) \textit{Dhū’l-Ḳaʿda}}\hspace{0.5em}\parbox[t]{\dimexpr0.2450\linewidth-0.5em\relax}{\raggedright\small\mbox{}= 30 „}\hspace{0.5em}}{}
\EdL{15}{AK07-J0066-L15}{\parbox[t]{\dimexpr0.2450\linewidth-0.5em\relax}{\raggedright\small\mbox{}6) \textit{Djumādā II}}\hspace{0.5em}\parbox[t]{\dimexpr0.2450\linewidth-0.5em\relax}{\raggedright\small\mbox{}= 29 „}\hspace{0.5em}\parbox[t]{\dimexpr0.2450\linewidth-0.5em\relax}{\raggedright\small\mbox{}12) \textit{Dhū’l-Ḥidjdja}}\hspace{0.5em}\parbox[t]{\dimexpr0.2450\linewidth-0.5em\relax}{\raggedright\small\mbox{}= 29 „ (in Schalt-}\hspace{0.5em}}{}
\EdLr{16}{AK07-J0066-L16}{jahren 30)}{}
\EdLp{17}{AK07-J0066-L17}{Diese Namen sind im Text und in den Tafeln oft sehr entstellt\,; zu bemerken}{}
\EdL{18}{AK07-J0066-L18}{ist besonders, dass \textit{dj} meistens durch \textit{i} wiedergegeben wird, so wird der letzte Monat}{}
\EdL{19}{AK07-J0066-L19}{fast durchweg \textit{Dulheia} geschrieben\,; \textit{sh} wird oft, besonders in Ms. O durch \textit{sc} aus-}{}
\EdL{20}{AK07-J0066-L20}{gedrückt (s. auch Einleitung).}{}
\EdLp{21}{AK07-J0066-L21}{Am Schlusse des Kapitels bemerkt \Sperr{Kh.}, dass aus der Länge des arabischen}{}
\EdL{22}{AK07-J0066-L22}{Mondjahres im Verhältnis zum Sonnenjahr der Umstand hervorgehe, dass der Anfang}{}
\EdL{23}{AK07-J0066-L23}{des arabischen Jahres in verschiedene Jahreszeiten, bald in den Winter, bald in den}{}
\EdL{24}{AK07-J0066-L24}{Sommer falle.}{}
\EdLc{25}{AK07-J0066-L25}{\textbf{Kap. 2. Beschreibung des römischen Jahres.}}{}
\EdLc{26}{AK07-J0066-L26}{(Text p. 2—4, Tafel 1.)}{}
\EdLp{27}{AK07-J0066-L27}{Mit diesem Namen bezeichnen die arabischen Astronomen gewöhnlich die}{}
\EdL{28}{AK07-J0066-L28}{Jahre der Aera \textit{Dhū’l-Ḳarnains} (d. h. \Sperr{Alexanders des Grossen}), d. i. der sog.}{}
\EdL{29}{AK07-J0066-L29}{\textit{Seleucidischen} Aera, die am 1. Oktober 312 v. Chr. ihren Anfang nahm, und deren}{}
\EdL{30}{AK07-J0066-L30}{Jahre zu 365\textsuperscript{1}/\textsubscript{4} Tagen wie in der \textit{julianischen} Zeitrechnung angenommen wurden.}{}
\EdL{31}{AK07-J0066-L31}{Interessant ist die Bemerkung unseres Textes, dass der Schalttag jeweilen (also alle}{}
\EdL{32}{AK07-J0066-L32}{4 Jahre) am Ende des Dezembers angefügt werde. \Sperr{Nallino}\textsuperscript{1}) bezeichnet dies als}{}
\EdL{33}{AK07-J0066-L33}{eine irrtümliche Ansicht\,; er kannte aber diese Stelle unserer Tafeln noch nicht, nach}{}
\EdL{34}{AK07-J0066-L34}{dieser scheint es ziemlich sicher zu sein, dass wenigstens bei den christlichen Spaniern}{}
\EdL{35}{AK07-J0066-L35}{des 10. Jahrh. (zur Zeit \Sperr{Maslamas}) diese \textit{Intercalation} des Schalttages wirklich im}{}
\EdL{36}{AK07-J0066-L36}{Dezember stattgefunden hat, statt im Februar. Es ergibt sich dies auch in unserer}{}
\EdL{37}{AK07-J0066-L37}{Tafel 3\textsuperscript{a} aus den \textit{Notulae} oder \textit{Signa}, die vom Dezember auf den Januar in den}{}
\EdL{38}{AK07-J0066-L38}{Schaltjahren um 4 zunehmen, in den übrigen Jahren aber um 3. Die Spanier kannten}{}
\EdL{39}{AK07-J0066-L39}{ja auch eine besondere Zeitrechnung (s. unten), die 38 Jahre vor der unsrigen begann,}{}
\EdL{40}{AK07-J0066-L40}{so konnten sie wohl auch eine eigene Intercalation haben, die nur bei ihnen Giltig-}{}
\NoteRule
\EdN{N1}{AK07-J0066-N01}{\textsuperscript{1}) \Sperr{Batt.}, ed. \Sperr{Nallino}, II, 314—316.}
\FrameHead{0067}{p. 35}{0067}{p. 35}
\EdL{1}{AK07-J0067-L01}{keit hatte, im Osten aber unbekannt war\,; es sagen ja auch die Verfasser der \textit{Alphon-}}{}
\EdL{2}{AK07-J0067-L02}{\textit{sinischen} Tafeln (nach \Sperr{Nallino})\,: \textit{Dissentiunt de mense cui dies bissextus addere debeant\,;}}{}
\EdL{3}{AK07-J0067-L03}{\textit{alii enim mensi Decembri eum addunt, alii Februario.} Dasselbe sagen auch spätere}{}
\EdL{4}{AK07-J0067-L04}{westarabische Autoren.}{}
\EdLp{5}{AK07-J0067-L05}{Im letzten Alinea dieses Kapitels ist auf eine Tafel verwiesen, die die Zeiten}{}
\EdL{6}{AK07-J0067-L06}{angibt, die seit dem Eintreten wichtiger Ereignisse, wie z. B. der Sündflut, oder der}{}
\EdL{7}{AK07-J0067-L07}{Thronbesteigung (bezw. des Todes) hervorragender Herrscher bis zu einem gewissen}{}
\EdL{8}{AK07-J0067-L08}{Zeitpunkt verflossen sind\,; diese Tafel steht in den Manuskripten nicht an erster Stelle,}{}
\EdL{9}{AK07-J0067-L09}{weil sie aber im Texte zuerst erwähnt wird, haben wir sie auch als Tafel 1 gesetzt.}{}
\EdL{10}{AK07-J0067-L10}{Diese Tafel wird auch in dem in eckige Klammern eingeschlossenen Zusatz (p. 3—4)}{}
\EdL{11}{AK07-J0067-L11}{erwähnt, der nur in den Mss. CMN vorkommt. Das dieser Zusatz im arabischen}{}
\EdL{12}{AK07-J0067-L12}{Ms. \Sperr{Kh.}’s nicht gestanden, auch von \Sperr{Maslama} nicht hinzugefügt worden sein kann,}{}
\EdL{13}{AK07-J0067-L13}{beweisen gewisse Sätze desselben, wie z. B.\,: \textit{nos vero latinorum studemus utilitati, ideo}}{}
\EdL{14}{AK07-J0067-L14}{\textit{et in primis tabulam suppositam praenotavimus, per quam ab eo quidem anno quo hic}}{}
\EdL{15}{AK07-J0067-L15}{\textit{liber in nostrum sermonem translatus est} . . . . Dieser Zusatz ist also als Kommentar}{}
\EdL{16}{AK07-J0067-L16}{oder Scholium des Übersetzers oder eines Abschreibers zu betrachten, und weist uns}{}
\EdL{17}{AK07-J0067-L17}{vielleicht auch darauf hin, dass die Mss. CMN von einer andern Übersetzung her-}{}
\EdL{18}{AK07-J0067-L18}{stammen könnten als O (vergl. meine Einleitung).}{}
\EdLp{19}{AK07-J0067-L19}{In diesem Zusatz ist die Zeit, die von der Sündflut bis auf \textit{Dhū’l-Ḳarnain},}{}
\EdL{20}{AK07-J0067-L20}{d. h. bis zum Beginn der \textit{Seleucidischen Aera} verflossen war, zu IIDCCXCIII Jahren,}{}
\EdL{21}{AK07-J0067-L21}{II Monaten, V Tagen angegeben\,; das II vor D sollte geschrieben sein \Ovl{II} und bedeutet}{}
\EdL{22}{AK07-J0067-L22}{dann MM, die Zeit beträgt also 2793 J. 2 M. 5 T. — Die Tafel 1 ist sehr ungeordnet}{}
\EdL{23}{AK07-J0067-L23}{in Bezug auf die verschiedenen Aeren, ich gebe im folgenden die wichtigsten in}{}
\EdL{24}{AK07-J0067-L24}{direkter Reihenfolge an\,:}{}
\EdL{25}{AK07-J0067-L25}{\parbox[t]{\dimexpr0.4200\linewidth-0.5em\relax}{\raggedright\small\mbox{}Von \Sperr{Adam} bis zum \textit{Diluvium}}\hspace{0.5em}\parbox[t]{\dimexpr0.5600\linewidth-0.5em\relax}{\raggedright\small\mbox{}2280 Jahre}\hspace{0.5em}}{}
\EdL{26}{AK07-J0067-L26}{\parbox[t]{\dimexpr0.4200\linewidth-0.5em\relax}{\raggedright\small\mbox{}Vom \textit{Diluvium} bis auf \Sperr{Nabukodonossor}\textsuperscript{1})}\hspace{0.5em}\parbox[t]{\dimexpr0.1400\linewidth-0.5em\relax}{\raggedright\small\mbox{}2356 „}\hspace{0.5em}\parbox[t]{\dimexpr0.1400\linewidth-0.5em\relax}{\raggedright\small\mbox{}7 Mon.}\hspace{0.5em}\parbox[t]{\dimexpr0.2800\linewidth-0.5em\relax}{\raggedright\small\mbox{}23 Tg.}\hspace{0.5em}}{}
\EdL{27}{AK07-J0067-L27}{\parbox[t]{\dimexpr0.4200\linewidth-0.5em\relax}{\raggedright\small\mbox{}Von \Sperr{Nabukodonossor} bis auf \Sperr{Alexander}}\hspace{0.5em}\parbox[t]{\dimexpr0.1400\linewidth-0.5em\relax}{\raggedright\small\mbox{}436 „}\hspace{0.5em}\parbox[t]{\dimexpr0.1400\linewidth-0.5em\relax}{\raggedright\small\mbox{}6 „}\hspace{0.5em}\parbox[t]{\dimexpr0.1400\linewidth-0.5em\relax}{\raggedright\small\mbox{}12\textsuperscript{1}/\textsubscript{2} „}\hspace{0.5em}\parbox[t]{\dimexpr0.1400\linewidth-0.5em\relax}{\raggedright\small\mbox{}(?)}\hspace{0.5em}}{}
\EdL{28}{AK07-J0067-L28}{\parbox[t]{\dimexpr0.4200\linewidth-0.5em\relax}{\raggedright\small\mbox{}Von \Sperr{Alexander} bis auf \Sperr{Christus}}\hspace{0.5em}\parbox[t]{\dimexpr0.1400\linewidth-0.5em\relax}{\raggedright\small\mbox{}311 „}\hspace{0.5em}\parbox[t]{\dimexpr0.4200\linewidth-0.5em\relax}{\raggedright\small\mbox{}3 „}\hspace{0.5em}}{}
\EdL{29}{AK07-J0067-L29}{\parbox[t]{\dimexpr0.4200\linewidth-0.5em\relax}{\raggedright\small\mbox{}Von \Sperr{Christus} bis auf die \textit{Hidjra}}\hspace{0.5em}\parbox[t]{\dimexpr0.1400\linewidth-0.5em\relax}{\raggedright\small\mbox{}621 „}\hspace{0.5em}\parbox[t]{\dimexpr0.1400\linewidth-0.5em\relax}{\raggedright\small\mbox{}6 „}\hspace{0.5em}\parbox[t]{\dimexpr0.2800\linewidth-0.5em\relax}{\raggedright\small\mbox{}15 „}\hspace{0.5em}}{}
\EdL{30}{AK07-J0067-L30}{\parbox[t]{\dimexpr0.4200\linewidth-0.5em\relax}{\raggedright\small\mbox{}Von der \textit{Hidjra} bis auf \Sperr{Yezdegird}}\hspace{0.5em}\parbox[t]{\dimexpr0.1400\linewidth-0.5em\relax}{\raggedright\small\mbox{}9 „}\hspace{0.5em}\parbox[t]{\dimexpr0.1400\linewidth-0.5em\relax}{\raggedright\small\mbox{}11 „}\hspace{0.5em}\parbox[t]{\dimexpr0.2800\linewidth-0.5em\relax}{\raggedright\small\mbox{}4 „}\hspace{0.5em}}{}
\EdLp{31}{AK07-J0067-L31}{In dieser Tafel 1 ist auch zweimal eine \textit{Aera azofar} genannt\,; es ist dies eine}{}
\EdL{32}{AK07-J0067-L32}{in Spanien vor und nach der Zeit \textit{Maslamas} unter der christlichen Bevölkerung}{}
\EdL{33}{AK07-J0067-L33}{gebräuchliche Aera, die genau 38 Jahre vor unserer christlichen Zeitrechnung ihren}{}
\EdL{34}{AK07-J0067-L34}{Anfang hatte, wie es auch unsere Tafel 1 in der fünft- und viertletzten Zeile richtig}{}
\EdL{35}{AK07-J0067-L35}{angibt\,; nach der sechst- und siebtletzten Zeile aber wäre der Unterschied beider}{}
\EdL{36}{AK07-J0067-L36}{Aeren 38 Jahre und 2 Tage, hier wird wohl ein Fehler vorliegen. — Diese spanische}{}
\EdL{37}{AK07-J0067-L37}{Aera hiess bei den Arabern \textit{ta’rīkh al-ṣufr} (gesprochen \textit{aṣ-ṣufr})\,; ob dies mit \textit{ṣifr}}{}
\NoteRule
\EdN{N1}{AK07-J0067-N01}{\textsuperscript{1}) Dies ist die von \Sperr{Ptol.} gebrauchte \textit{Aera Nabonassars} (\Sperr{Batt.} u. a. nennen ihn \textit{Bukhtanassar}}
\EdN{N2}{AK07-J0067-N02}{\textit{primus}), am 26. Februar 747 v. Chr. beginnend\,; dies stimmt allerdings nicht ganz mit den obigen}
\EdN{N3}{AK07-J0067-N03}{Angaben.}
\FrameHead{0068}{p. 36}{0068}{p. 36}
\EdL{1}{AK07-J0068-L01}{(= leer, Null) in Verbindung stehe, wie \Sperr{Ideler}\textsuperscript{1}) vermutet, ist sehr unwahr-}{}
\EdL{2}{AK07-J0068-L02}{scheinlich.\textsuperscript{2})}{}
\EdLp{3}{AK07-J0068-L03}{Am Schlusse dieses Zusatzes wird eine zweite Tafel erwähnt, mit welcher die}{}
\EdL{4}{AK07-J0068-L04}{Jahre der christlichen Zeitrechnung (es steht aber neben \textit{anni domini} auch \textit{anni romani}}{}
\EdL{5}{AK07-J0068-L05}{im Text) in arabische verwandelt werden können\,; die Tafel fehlt in allen Mss. (s.}{}
\EdL{6}{AK07-J0068-L06}{Einleitung), es befindet sich bloss die Tafel 3 in CO und N, die betitelt ist\,: \textit{Pagina}}{}
\EdL{7}{AK07-J0068-L07}{\textit{reducendi annos Alexandri ad elhigeram}, die aber in erster Linie dazu dient, arabische}{}
\EdL{8}{AK07-J0068-L08}{Jahre in solche der Seleucidischen Aera zu verwandeln, allerdings kann sie für das}{}
\EdL{9}{AK07-J0068-L09}{umgekehrte auch dienen, wir kommen auf sie unten zurück.}{}
\EdLc{10}{AK07-J0068-L10}{\textbf{Kap. 3. Über die Überschriften der Tafel-Kolonnen.}}{}
\EdLc{11}{AK07-J0068-L11}{(Text p. 4—5.)}{}
\EdLp{12}{AK07-J0068-L12}{So muss man den Titel \textit{de numerorum praescriptionibus} übersetzen. Der Ver-}{}
\EdL{13}{AK07-J0068-L13}{fasser erklärt hier, was man unter \textit{anni collecti}, \textit{anni plani}, \textit{gradus}, \textit{dakaicae}, etc. zu}{}
\EdL{14}{AK07-J0068-L14}{verstehen habe, und wie man beim Gebrauche der Tafeln (er führt als Beispiel die-}{}
\EdL{15}{AK07-J0068-L15}{jenigen für die Bestimmung der mittlern Örter der Planeten an) zu verfahren habe.}{}
\EdL{16}{AK07-J0068-L16}{Ich habe zu diesem Kapitel folgende Bemerkungen hinzuzufügen\,: In \textit{anni plani} ist}{}
\EdL{17}{AK07-J0068-L17}{\textit{planum} die Übersetzung des arabischen \textit{mabsūṭ} = ausgebreitet\,; statt \textit{planum} findet}{}
\EdL{18}{AK07-J0068-L18}{man daher auch (in C) \textit{expansum}, \Sperr{wir} würden sagen \textit{anni singuli}, N hat \textit{anni residui}.}{}
\EdL{19}{AK07-J0068-L19}{— Die erste Kolonne der Tafeln\textsuperscript{3}) für die Bestimmung der mittlern Örter der Planeten}{}
\EdL{20}{AK07-J0068-L20}{enthält gewöhnlich den Weg des Planeten von der \textit{Radix} an in je 30 Jahren\,; dann}{}
\EdL{21}{AK07-J0068-L21}{folgt die Kolonne für die einzelnen Jahre von 1 bis 30, dann die Kolonne für die}{}
\EdL{22}{AK07-J0068-L22}{Monate u. s. f. Für das Wort „Minute“ wird eigentümlicherweise öfters in C und O}{}
\EdL{23}{AK07-J0068-L23}{das arabische \textit{daḳīḳa}, und zwar in der Pluralform \textit{daḳā’iḳ} gebraucht, diese wird}{}
\EdL{24}{AK07-J0068-L24}{lateinisch durch \textit{dakaica} wiedergegeben, und dieses dann wie ein echt lateinisches}{}
\EdL{25}{AK07-J0068-L25}{Wort dekliniert (z. B. \textit{dakaicarum}). Dagegen wird, besonders in Ms. C, für „Sekunde“}{}
\EdL{26}{AK07-J0068-L26}{die arabische Singularform \textit{thenia} = \textit{thāniya} gesetzt, und von dieser der Plural \textit{theniae}}{}
\EdL{27}{AK07-J0068-L27}{gebildet, statt die arabische Pluralform \textit{thawānī} zu setzen. Im übrigen hat das Ms. C}{}
\EdL{28}{AK07-J0068-L28}{meistens \textit{puncta} statt \textit{minutae}\textsuperscript{4}), und \textit{minuta} statt \textit{secundae}\textsuperscript{4}). — Auf die Willkür und}{}
\EdL{29}{AK07-J0068-L29}{Inkonsequenz, die in diesen Kolonnen-Überschriften herrscht, habe ich schon in der}{}
\EdL{30}{AK07-J0068-L30}{Einleitung aufmerksam gemacht.}{}
\EdLc{31}{AK07-J0068-L31}{\textbf{Kap. 4. Über den Beginn der arabischen Monate.}}{}
\EdLc{32}{AK07-J0068-L32}{(Text p. 5, Tafeln 2—2\textsuperscript{b}.)}{}
\EdLp{33}{AK07-J0068-L33}{Dieses Kapitel, zu dem Tafel 2 gehört, die in allen vier Mss., in O aber auf}{}
\EdL{34}{AK07-J0068-L34}{der letzten Seite (159\textsuperscript{r}) sich befindet, lehrt den Wochentag (\textit{feria}) bestimmen, mit}{}
\EdL{35}{AK07-J0068-L35}{dem jeder beliebige arabische Monat in jedem beliebigen Jahre beginnt. Ich gebe,}{}
\NoteRule
\EdN{N1}{AK07-J0068-N01}{\textsuperscript{1}) \textit{Handbuch der math. u. techn. Chronologie}, Berlin 1826, II, 423—425.}
\EdN{N2}{AK07-J0068-N02}{\textsuperscript{2}) Vergl. auch \Sperr{Dozy}, \textit{Supplément aux dictionnaires arabes}, s. v. \textit{ṣafar} u. \textit{aṣfar}.}
\EdN{N3}{AK07-J0068-N03}{\textsuperscript{3}) Nach der Kolonne der Jahre.}
\EdN{N4}{AK07-J0068-N04}{\textsuperscript{4}) Ich schreibe \textit{minutae} und \textit{secundae} statt \textit{minuta} und \textit{secunda}, indem ich diese Wörter als}
\EdN{N5}{AK07-J0068-N05}{Adjectiva zu \textit{partes} auffasse, das richtige wäre eigentlich (\textit{partes}) \textit{primae}, (\textit{partes}) \textit{secundae}.}
\FrameHead{0069}{p. 37}{0069}{p. 37}
\EdL{1}{AK07-J0069-L01}{um die Regeln des Textes, die oft etwas schwer verständlich sind, klar zu machen,}{}
\EdL{2}{AK07-J0069-L02}{folgende zwei Beispiele\,:}{}
\EdLp{3}{AK07-J0069-L03}{1. An welchem Wochentage begann der \textit{Ramaḍān} des Jahres 529 d. H.\,?}{}
\EdL{4}{AK07-J0069-L04}{\parbox[t]{\dimexpr0.6533\linewidth-0.5em\relax}{\raggedright\small\mbox{}für 510 Jahre ist die Nota (d. h. Wochentag-Nummer)\textsuperscript{1})}\hspace{0.5em}\parbox[t]{\dimexpr0.3267\linewidth-0.5em\relax}{\raggedright\small\mbox{}= 1}\hspace{0.5em}}{}
\EdL{5}{AK07-J0069-L05}{\parbox[t]{\dimexpr0.6533\linewidth-0.5em\relax}{\raggedright\small\mbox{}„ 19 „ „ „ „}\hspace{0.5em}\parbox[t]{\dimexpr0.3267\linewidth-0.5em\relax}{\raggedright\small\mbox{}= 6}\hspace{0.5em}}{}
\EdL{6}{AK07-J0069-L06}{\parbox[t]{\dimexpr0.6533\linewidth-0.5em\relax}{\raggedright\small\mbox{}„ den Monat Ramaḍān ist die Nota}\hspace{0.5em}\parbox[t]{\dimexpr0.3267\linewidth-0.5em\relax}{\raggedright\small\mbox{}= 6}\hspace{0.5em}}{}
\EdL{7}{AK07-J0069-L07}{\parbox[t]{\dimexpr0.6533\linewidth-0.5em\relax}{\raggedright\small\mbox{}}\hspace{0.5em}\parbox[t]{\dimexpr0.3267\linewidth-0.5em\relax}{\raggedright\small\mbox{}Summa = 13}\hspace{0.5em}}{}
\EdLp{8}{AK07-J0069-L08}{Von dieser Summe wird 7 subtrahiert (bezw. sie wird durch 7 geteilt), der Rest}{}
\EdL{9}{AK07-J0069-L09}{ist 6, d. h. der 1. \textit{Ramaḍān} war ein \textit{Freitag}.}{}
\EdLp{10}{AK07-J0069-L10}{2. Dieses Kapitel enthält am Schlusse die Angabe, dass der 26. Januar des J. 1126}{}
\EdL{11}{AK07-J0069-L11}{(julian. Zeitrechnung) auf den 1. \textit{Muḥarram} des J. 520 d. H. fiel, und ein \textit{Dienstag} (\textit{feria}}{}
\EdL{12}{AK07-J0069-L12}{\textit{tertia}) war\,; wir wollen auch die Richtigkeit dieser Angabe aus Tafel 2 nachweisen\,:}{}
\EdL{13}{AK07-J0069-L13}{\parbox[t]{\dimexpr0.6533\linewidth-0.5em\relax}{\raggedright\small\mbox{}für 510 Jahre ist die Nota}\hspace{0.5em}\parbox[t]{\dimexpr0.3267\linewidth-0.5em\relax}{\raggedright\small\mbox{}= 1}\hspace{0.5em}}{}
\EdL{14}{AK07-J0069-L14}{\parbox[t]{\dimexpr0.6533\linewidth-0.5em\relax}{\raggedright\small\mbox{}„ 10 „ „ „ „}\hspace{0.5em}\parbox[t]{\dimexpr0.3267\linewidth-0.5em\relax}{\raggedright\small\mbox{}= 1}\hspace{0.5em}}{}
\EdL{15}{AK07-J0069-L15}{\parbox[t]{\dimexpr0.6533\linewidth-0.5em\relax}{\raggedright\small\mbox{}„ den Monat Muḥarram ist die Nota}\hspace{0.5em}\parbox[t]{\dimexpr0.3267\linewidth-0.5em\relax}{\raggedright\small\mbox{}= 1}\hspace{0.5em}}{}
\EdL{16}{AK07-J0069-L16}{\parbox[t]{\dimexpr0.6533\linewidth-0.5em\relax}{\raggedright\small\mbox{}}\hspace{0.5em}\parbox[t]{\dimexpr0.3267\linewidth-0.5em\relax}{\raggedright\small\mbox{}Summa = 3}\hspace{0.5em}}{}
\EdL{17}{AK07-J0069-L17}{d. h. der 1. \textit{Muḥarram} d. J. 520 war ein \textit{Dienstag}. — Diese Angabe findet sich bloss}{}
\EdL{18}{AK07-J0069-L18}{in Ms. O und will wahrscheinlich sagen, dass die Übersetzung dieses Werkes ins}{}
\EdL{19}{AK07-J0069-L19}{Lateinische an diesem Tage entweder begonnen, oder vollendet worden sei, demnach}{}
\EdL{20}{AK07-J0069-L20}{läge der Schluss nahe, dass das Ms. O, wenn nicht den Autographen \Sperr{Athelhards},}{}
\EdL{21}{AK07-J0069-L21}{so doch eine bald nachher gemachte Abschrift desselben darstellen könnte.}{}
\EdLp{22}{AK07-J0069-L22}{In Tafel 2 steht die Bezeichnung \textit{keb.} (= \textit{kebice} = Schaltjahr) mehrmals an}{}
\EdL{23}{AK07-J0069-L23}{unrichtiger Stelle, auch fehlt ein Schaltjahr, indem die Zahl derselben innerhalb 30}{}
\EdL{24}{AK07-J0069-L24}{Jahren 11 beträgt.\textsuperscript{2})}{}
\EdLp{25}{AK07-J0069-L25}{Ich schliesse hier noch die Erwähnung zweier Tafeln an, die im Texte gar nicht}{}
\EdL{26}{AK07-J0069-L26}{genannt sind, und die ich mit 2\textsuperscript{a} und 2\textsuperscript{b} bezeichnet habe. Die erste Tafel lehrt,}{}
\EdL{27}{AK07-J0069-L27}{mit welchem Wochentage die Monate der \textit{Aera} \Sperr{Yezdegirds} beginnen\,; man wird}{}
\EdL{28}{AK07-J0069-L28}{leicht den Wochentag bestimmen können, wenn wir bemerken, dass diese Tafel in}{}
\EdL{29}{AK07-J0069-L29}{der ersten Kolonne der Abteilung, die mit \textit{Signa vel Notae} überschrieben ist, jeweilen}{}
\EdL{30}{AK07-J0069-L30}{den Wochentag des Beginns der Monate des 1. Jahres der Aera angibt\,; so begann}{}
\EdL{31}{AK07-J0069-L31}{der 1. \textit{Affrordinme} des Jahres 1 mit einem Dienstag (3), der 1. \textit{Azdihestme} mit einem}{}
\EdL{32}{AK07-J0069-L32}{Donnerstag (5), u. s. w. Für jedes folgende Jahr verschiebt sich der Wochentag um}{}
\EdL{33}{AK07-J0069-L33}{eins, da das persische Jahr 365 Tage zählte. Da jeder persische Monat 30 Tage hatte,}{}
\EdL{34}{AK07-J0069-L34}{ausser dem \textit{Abenme}, in welchem die 5 Tage über 360 hinaus eingeschaltet wurden,}{}
\EdL{35}{AK07-J0069-L35}{wodurch er auf 35 Tage, also 5 volle Wochen kam, hat der folgende Monat \textit{Adarme}}{}
\EdL{36}{AK07-J0069-L36}{dieselben \textit{Notae} wie \textit{Abenme}.\textsuperscript{3}) — Die Tafel 2\textsuperscript{b} enthält die \textit{ägyptischen} (\textit{koptischen})}{}
\NoteRule
\EdN{N1}{AK07-J0069-N01}{\textsuperscript{1}) 1 = Sonntag, 2 = Montag, u. s. f. . . . 7 = Samstag.}
\EdN{N2}{AK07-J0069-N02}{\textsuperscript{2}) Siehe oben p. 34.}
\EdN{N3}{AK07-J0069-N03}{\textsuperscript{3}) Eine Anleitung, wie man für jedes beliebige Jahr dieser Aera den Wochentag des Anfangs}
\EdN{N4}{AK07-J0069-N04}{jedes Monats bestimme, giebt \Sperr{Al-Bīrūnī} in seiner \textit{Chronologie} (engl. Übersetzung von \Sperr{Sachau},}
\EdN{N5}{AK07-J0069-N05}{p. 184, arab. Text p. 203).}
\FrameHead{0070}{p. 38}{0070}{p. 38}
\EdL{1}{AK07-J0070-L01}{Monate, die auch \Sperr{Ptol.} für seine astronomischen Angaben gebraucht hat, mit ihren}{}
\EdL{2}{AK07-J0070-L02}{\textit{Signa} oder \textit{Notae}. — Die persischen sowohl als koptischen Namen sind in den beiden}{}
\EdL{3}{AK07-J0070-L03}{Mss. M und N, in denen allein die beiden Tafeln 2\textsuperscript{a} und 2\textsuperscript{b} stehen, ziemlich entstellt,}{}
\EdL{4}{AK07-J0070-L04}{wir haben die richtigen nach \Sperr{Nallino}\textsuperscript{1}) und \Sperr{Al-Bīrūnīs} \textit{Chronologie}\textsuperscript{2}) in das}{}
\EdL{5}{AK07-J0070-L05}{Glossar aufgenommen.}{}
\EdLc{6}{AK07-J0070-L06}{\textbf{Kap. 5. Über die Auffindung des römischen Jahres.}}{}
\EdLc{7}{AK07-J0070-L07}{(Text p. 6—7, Tafeln 3 u. 3\textsuperscript{a}.)}{}
\EdLp{8}{AK07-J0070-L08}{Der Titel heisst eigentlich\,: \textit{De anni bissextilis inventione}. Es handelt sich aber}{}
\EdL{9}{AK07-J0070-L09}{in der Tat um die Verwandlung der Jahre der \textit{Hidjra} in Jahre der \textit{Seleucidischen Aera},}{}
\EdL{10}{AK07-J0070-L10}{oder um die Auffindung des \textit{römischen} Datums, das einem gegebenen \textit{arabischen} Datum}{}
\EdL{11}{AK07-J0070-L11}{entspricht. Die Jahre der römischen Aera wurden nämlich auch \textit{anni bissextiles} (Schalt-}{}
\EdL{12}{AK07-J0070-L12}{jahre) genannt, was etwas eigentümlich erscheint, da ja die Araber auch Schaltjahre}{}
\EdL{13}{AK07-J0070-L13}{hatten\,; es steht dieses \textit{anni bissextiles} aber im Gegensatz zu den \textit{anni vagi} der Aegypter,}{}
\EdL{14}{AK07-J0070-L14}{die keine \textit{Intercalation} kannten. Die zugehörige Tafel 3 ist betitelt\,: \textit{Pagina reducendi}}{}
\EdL{15}{AK07-J0070-L15}{\textit{annos Alexandri ad elhigeram}, was eigentlich unrichtig ist. Die Tafel ist etwas anders}{}
\EdL{16}{AK07-J0070-L16}{eingerichtet als die entsprechenden Tafeln bei \Sperr{Batt.}\textsuperscript{3}), die für den Anfang jedes Jahres}{}
\EdL{17}{AK07-J0070-L17}{der \textit{Hidjra} von 1 bis 600 das zugehörige Datum der \textit{römischen Aera} geben\,; die Tafel}{}
\EdL{18}{AK07-J0070-L18}{\Sperr{Kh.}’s nimmt viel weniger Raum ein, erfordert aber noch Nebenrechnungen. Wir}{}
\EdL{19}{AK07-J0070-L19}{geben drei Beispiele für diese Reduktion\,:}{}
\EdLp{20}{AK07-J0070-L20}{1. Welchem Tage der \textit{römischen Aera} entspricht der \textit{1. Muḥarram} des Jahres}{}
\EdL{21}{AK07-J0070-L21}{520 d. H.\,?}{}
\EdL{22}{AK07-J0070-L22}{\parbox[t]{\dimexpr0.3267\linewidth-0.5em\relax}{\raggedright\small\mbox{}Am Ende des 510. Jahres d. H. waren verflossen}\hspace{0.5em}\parbox[t]{\dimexpr0.1633\linewidth-0.5em\relax}{\raggedright\small\mbox{}1427 J.}\hspace{0.5em}\parbox[t]{\dimexpr0.1633\linewidth-0.5em\relax}{\raggedright\small\mbox{}7 M.}\hspace{0.5em}\parbox[t]{\dimexpr0.1633\linewidth-0.5em\relax}{\raggedright\small\mbox{}5\textsuperscript{1}/\textsubscript{4} T.}\hspace{0.5em}\parbox[t]{\dimexpr0.1633\linewidth-0.5em\relax}{\raggedright\small\mbox{}der röm. Aera}\hspace{0.5em}}{}
\EdL{23}{AK07-J0070-L23}{\parbox[t]{\dimexpr0.3267\linewidth-0.5em\relax}{\raggedright\small\mbox{}9 muhammed. Jahre sind}\hspace{0.5em}\parbox[t]{\dimexpr0.1633\linewidth-0.5em\relax}{\raggedright\small\mbox{}= 8 „}\hspace{0.5em}\parbox[t]{\dimexpr0.1633\linewidth-0.5em\relax}{\raggedright\small\mbox{}8 „}\hspace{0.5em}\parbox[t]{\dimexpr0.1633\linewidth-0.5em\relax}{\raggedright\small\mbox{}27 „}\hspace{0.5em}\parbox[t]{\dimexpr0.1633\linewidth-0.5em\relax}{\raggedright\small\mbox{}„ „ „}\hspace{0.5em}}{}
\EdL{24}{AK07-J0070-L24}{\parbox[t]{\dimexpr0.1633\linewidth-0.5em\relax}{\raggedright\small\mbox{}}\hspace{0.5em}\parbox[t]{\dimexpr0.1633\linewidth-0.5em\relax}{\raggedright\small\mbox{}Summa\,:}\hspace{0.5em}\parbox[t]{\dimexpr0.1633\linewidth-0.5em\relax}{\raggedright\small\mbox{}1436 „}\hspace{0.5em}\parbox[t]{\dimexpr0.1633\linewidth-0.5em\relax}{\raggedright\small\mbox{}3 „}\hspace{0.5em}\parbox[t]{\dimexpr0.3267\linewidth-0.5em\relax}{\raggedright\small\mbox{}32\textsuperscript{1}/\textsubscript{4} „}\hspace{0.5em}}{}
\EdLp{25}{AK07-J0070-L25}{Da eine einmalige Verwandlung von 12 Monaten zu einem Jahre stattgefunden}{}
\EdL{26}{AK07-J0070-L26}{hat, so hat man von dieser Summe 5\textsuperscript{1}/\textsubscript{4} Tage zu subtrahieren (nimmt man nämlich}{}
\EdL{27}{AK07-J0070-L27}{für 12 Monate ein Jahr, so hat man 5\textsuperscript{1}/\textsubscript{4} Tage zu viel genommen)\textsuperscript{4}), dies gibt 1436 J.}{}
\EdL{28}{AK07-J0070-L28}{3 M. 27 T.\,; da unter den 3 Monaten vom Oktober an zwei Monate mit 31 Tagen sich}{}
\EdL{29}{AK07-J0070-L29}{befinden, so muss man von der Zahl der Tage noch 2 subtrahieren, dies gibt 25 Tage,}{}
\EdL{30}{AK07-J0070-L30}{d. h. seit dem Beginn der römischen Aera (1. Okt. 312 v. Chr.) sind 1436 J. 3 M. 25 T.}{}
\EdL{31}{AK07-J0070-L31}{verflossen, d. h. der \textit{1. Muḥarram} 520 d. H. entspricht dem \textit{26. Kānūn II.} des Jahres}{}
\EdL{32}{AK07-J0070-L32}{1437, d. i. dem \textit{26. Januar} des J. 1126 n. Chr.}{}
\EdLp{33}{AK07-J0070-L33}{2. Welchem Tage der \textit{römischen Aera} entspricht der \textit{5. Radjab} des Jahres}{}
\EdL{34}{AK07-J0070-L34}{413 d. H.\,?}{}
\NoteRule
\EdN{N1}{AK07-J0070-N01}{\textsuperscript{1}) \Sperr{Batt.}, ed. \Sperr{Nallino}, I, 67 u. 243.}
\EdN{N2}{AK07-J0070-N02}{\textsuperscript{2}) Englische Übersetzung v. \Sperr{Sachau}, p. 82, u. 199—217, arab. Text p. 215—233.}
\EdN{N3}{AK07-J0070-N03}{\textsuperscript{3}) Ed. \Sperr{Nallino}, II, 9—18.}
\EdN{N4}{AK07-J0070-N04}{\textsuperscript{4}) Da die Monate nur zu 30 Tagen gerechnet sind.}
\FrameHead{0071}{p. 39}{0071}{p. 39}
\EdL{1}{AK07-J0071-L01}{\parbox[t]{\dimexpr0.3267\linewidth-0.5em\relax}{\raggedright\small\mbox{}Am Ende des 390. Jahres d. H. waren verflossen}\hspace{0.5em}\parbox[t]{\dimexpr0.1633\linewidth-0.5em\relax}{\raggedright\small\mbox{}1311 J.}\hspace{0.5em}\parbox[t]{\dimexpr0.1633\linewidth-0.5em\relax}{\raggedright\small\mbox{}2 M.}\hspace{0.5em}\parbox[t]{\dimexpr0.1633\linewidth-0.5em\relax}{\raggedright\small\mbox{}\textsuperscript{1}/\textsubscript{4} T.}\hspace{0.5em}\parbox[t]{\dimexpr0.1633\linewidth-0.5em\relax}{\raggedright\small\mbox{}der röm. Aera}\hspace{0.5em}}{}
\EdL{2}{AK07-J0071-L02}{\parbox[t]{\dimexpr0.3267\linewidth-0.5em\relax}{\raggedright\small\mbox{}22 muhammed. Jahre sind}\hspace{0.5em}\parbox[t]{\dimexpr0.1633\linewidth-0.5em\relax}{\raggedright\small\mbox{}= 21 „}\hspace{0.5em}\parbox[t]{\dimexpr0.1633\linewidth-0.5em\relax}{\raggedright\small\mbox{}4 „}\hspace{0.5em}\parbox[t]{\dimexpr0.1633\linewidth-0.5em\relax}{\raggedright\small\mbox{}5\textsuperscript{3}/\textsubscript{4} „}\hspace{0.5em}\parbox[t]{\dimexpr0.1633\linewidth-0.5em\relax}{\raggedright\small\mbox{}„ „ „}\hspace{0.5em}}{}
\EdL{3}{AK07-J0071-L03}{\parbox[t]{\dimexpr0.3267\linewidth-0.5em\relax}{\raggedright\small\mbox{}die 6 ersten muhammed. Monate sind}\hspace{0.5em}\parbox[t]{\dimexpr0.1633\linewidth-0.5em\relax}{\raggedright\small\mbox{}=}\hspace{0.5em}\parbox[t]{\dimexpr0.1633\linewidth-0.5em\relax}{\raggedright\small\mbox{}5 „}\hspace{0.5em}\parbox[t]{\dimexpr0.1633\linewidth-0.5em\relax}{\raggedright\small\mbox{}27 „}\hspace{0.5em}\parbox[t]{\dimexpr0.1633\linewidth-0.5em\relax}{\raggedright\small\mbox{}„ „ „}\hspace{0.5em}}{}
\EdL{4}{AK07-J0071-L04}{\parbox[t]{\dimexpr0.6533\linewidth-0.5em\relax}{\raggedright\small\mbox{}dazu im Radjab noch}\hspace{0.5em}\parbox[t]{\dimexpr0.3267\linewidth-0.5em\relax}{\raggedright\small\mbox{}4 „}\hspace{0.5em}}{}
\EdL{5}{AK07-J0071-L05}{\parbox[t]{\dimexpr0.1633\linewidth-0.5em\relax}{\raggedright\small\mbox{}}\hspace{0.5em}\parbox[t]{\dimexpr0.1633\linewidth-0.5em\relax}{\raggedright\small\mbox{}Summa\,:}\hspace{0.5em}\parbox[t]{\dimexpr0.1633\linewidth-0.5em\relax}{\raggedright\small\mbox{}1333 „}\hspace{0.5em}\parbox[t]{\dimexpr0.1633\linewidth-0.5em\relax}{\raggedright\small\mbox{}0 „}\hspace{0.5em}\parbox[t]{\dimexpr0.3267\linewidth-0.5em\relax}{\raggedright\small\mbox{}7 „}\hspace{0.5em}}{}
\EdLp{6}{AK07-J0071-L06}{Da wieder eine einmalige Verwandlung von 12 Monaten in ein Jahr stattgefunden}{}
\EdL{7}{AK07-J0071-L07}{hat, muss man von dieser Summe 5\textsuperscript{1}/\textsubscript{4} Tage abziehen, dies gibt 1333 J. 0 M. 2 T.,}{}
\EdL{8}{AK07-J0071-L08}{d. h. seit dem Beginn der römischen Aera sind bis zum 5. Radjab 413 d. H. 1333 J.}{}
\EdL{9}{AK07-J0071-L09}{0 M. 2 T. verflossen, oder der 5. Radjab 413 entspricht dem 3. \textit{Tishrīn I.} des Jahres}{}
\EdL{10}{AK07-J0071-L10}{1334, d. i. dem \textit{3. Oktober} 1022 n. Chr.\textsuperscript{1})}{}
\EdLp{11}{AK07-J0071-L11}{3. Welchem Tage der \textit{römischen Aera} entspricht der \textit{3. Shawwāl} des Jahres}{}
\EdL{12}{AK07-J0071-L12}{399 d. H.\,?}{}
\EdL{13}{AK07-J0071-L13}{\parbox[t]{\dimexpr0.3267\linewidth-0.5em\relax}{\raggedright\small\mbox{}Am Ende des 390. Jahres d. H. waren verflossen}\hspace{0.5em}\parbox[t]{\dimexpr0.1633\linewidth-0.5em\relax}{\raggedright\small\mbox{}1311 J.}\hspace{0.5em}\parbox[t]{\dimexpr0.1633\linewidth-0.5em\relax}{\raggedright\small\mbox{}2 M.}\hspace{0.5em}\parbox[t]{\dimexpr0.1633\linewidth-0.5em\relax}{\raggedright\small\mbox{}\textsuperscript{1}/\textsubscript{4} T.}\hspace{0.5em}\parbox[t]{\dimexpr0.1633\linewidth-0.5em\relax}{\raggedright\small\mbox{}der röm. Aera}\hspace{0.5em}}{}
\EdL{14}{AK07-J0071-L14}{\parbox[t]{\dimexpr0.3267\linewidth-0.5em\relax}{\raggedright\small\mbox{}8 muhammed. Jahre sind}\hspace{0.5em}\parbox[t]{\dimexpr0.1633\linewidth-0.5em\relax}{\raggedright\small\mbox{}= 7 „}\hspace{0.5em}\parbox[t]{\dimexpr0.1633\linewidth-0.5em\relax}{\raggedright\small\mbox{}9 „}\hspace{0.5em}\parbox[t]{\dimexpr0.1633\linewidth-0.5em\relax}{\raggedright\small\mbox{}8\textsuperscript{1}/\textsubscript{4} „}\hspace{0.5em}\parbox[t]{\dimexpr0.1633\linewidth-0.5em\relax}{\raggedright\small\mbox{}„ „ „}\hspace{0.5em}}{}
\EdL{15}{AK07-J0071-L15}{\parbox[t]{\dimexpr0.3267\linewidth-0.5em\relax}{\raggedright\small\mbox{}die 9 ersten muhammed. Monate sind}\hspace{0.5em}\parbox[t]{\dimexpr0.1633\linewidth-0.5em\relax}{\raggedright\small\mbox{}=}\hspace{0.5em}\parbox[t]{\dimexpr0.1633\linewidth-0.5em\relax}{\raggedright\small\mbox{}8 „}\hspace{0.5em}\parbox[t]{\dimexpr0.1633\linewidth-0.5em\relax}{\raggedright\small\mbox{}26 „}\hspace{0.5em}\parbox[t]{\dimexpr0.1633\linewidth-0.5em\relax}{\raggedright\small\mbox{}„ „ „}\hspace{0.5em}}{}
\EdL{16}{AK07-J0071-L16}{\parbox[t]{\dimexpr0.6533\linewidth-0.5em\relax}{\raggedright\small\mbox{}dazu im Shawwāl noch}\hspace{0.5em}\parbox[t]{\dimexpr0.3267\linewidth-0.5em\relax}{\raggedright\small\mbox{}2 „}\hspace{0.5em}}{}
\EdL{17}{AK07-J0071-L17}{\parbox[t]{\dimexpr0.1633\linewidth-0.5em\relax}{\raggedright\small\mbox{}}\hspace{0.5em}\parbox[t]{\dimexpr0.1633\linewidth-0.5em\relax}{\raggedright\small\mbox{}Summa\,:}\hspace{0.5em}\parbox[t]{\dimexpr0.1633\linewidth-0.5em\relax}{\raggedright\small\mbox{}1319 „}\hspace{0.5em}\parbox[t]{\dimexpr0.1633\linewidth-0.5em\relax}{\raggedright\small\mbox{}7 „}\hspace{0.5em}\parbox[t]{\dimexpr0.3267\linewidth-0.5em\relax}{\raggedright\small\mbox{}36\textsuperscript{1}/\textsubscript{2} „}\hspace{0.5em}}{}
\EdLp{18}{AK07-J0071-L18}{Da wieder eine einmalige Verwandlung von 12 Monaten in ein Jahr statt-}{}
\EdL{19}{AK07-J0071-L19}{gefunden hat, muss man von dieser Summe 5\textsuperscript{1}/\textsubscript{4} Tage abziehen, dies gibt 1319 J.}{}
\EdL{20}{AK07-J0071-L20}{7 M. 31\textsuperscript{1}/\textsubscript{4} T.\,; da unter den 7 Monaten vom Oktober an vier Monate mit 31 Tagen}{}
\EdL{21}{AK07-J0071-L21}{sich befinden, so muss man von der Zahl der Tage noch 4 subtrahieren, dies gibt}{}
\EdL{22}{AK07-J0071-L22}{1319 J. 7 M. 27\textsuperscript{1}/\textsubscript{4} T.\,; da ferner unter diesen 7 Monaten der Februar sich befindet,}{}
\EdL{23}{AK07-J0071-L23}{muss man zu der Zahl der Tage noch 2 addieren, dies gibt 1319 J. 7 M. 29\textsuperscript{1}/\textsubscript{4} T.,}{}
\EdL{24}{AK07-J0071-L24}{d. h. seit dem Beginn der römischen Aera sind bis zum 3. Shawwāl 399 d. H. 1319 J.}{}
\EdL{25}{AK07-J0071-L25}{7 M. 29 T. verflossen, oder der \textit{3. Shawwāl} 399 entspricht dem \textit{30. Hiyār} (richtige}{}
\EdL{26}{AK07-J0071-L26}{Form \textit{Ayyār}) 1320, d. i. dem \textit{30. Mai} 1009.}{}
\EdLp{27}{AK07-J0071-L27}{Nach den heutigen Vergleichungstabellen erhält man für das erste Beispiel den}{}
\EdL{28}{AK07-J0071-L28}{27. Januar, für das zweite den 4. Oktober, und für das dritte den 31. Mai\,; \Sperr{Kh.} oder}{}
\EdL{29}{AK07-J0071-L29}{\Sperr{Maslama} hat also wie \Sperr{Batt.} den Beginn der muhammedanischen Aera auf den}{}
\EdL{30}{AK07-J0071-L30}{15. Juli 622 statt auf den 16. angesetzt.}{}
\EdLp{31}{AK07-J0071-L31}{Die Tafel 3\textsuperscript{a} dient dazu zu bestimmen, ob das gefundene römische Jahr ein}{}
\EdL{32}{AK07-J0071-L32}{Schaltjahr sei oder nicht, und mit welchem Wochentage es und seine Monate be-}{}
\EdL{33}{AK07-J0071-L33}{ginnen. Die mit den Monatsnamen überschriebenen Kolonnen enthalten die \textit{Signa}}{}
\EdL{34}{AK07-J0071-L34}{(\textit{Notae}). Es ist zu bemerken, dass man, wenn man den Beginn eines Monats eines}{}
\EdL{35}{AK07-J0071-L35}{bestimmten Jahres in der \textit{julianischen} Zeitrechnung haben will, zu dieser Jahreszahl}{}
\EdL{36}{AK07-J0071-L36}{noch 311 (312) zu addieren hat.\textsuperscript{2}) Ferner ist zu berücksichtigen, dass im julianischen}{}
\NoteRule
\EdN{N1}{AK07-J0071-N01}{\textsuperscript{1}) Fällt das Datum in die Monate Oktober bis Dezember, so muss man von der römischen}
\EdN{N2}{AK07-J0071-N02}{Jahreszahl 312 abziehen, um die christliche (julianische) Jahreszahl zu erhalten, sonst 311, wofür}
\EdN{N3}{AK07-J0071-N03}{der Leser die Erklärung leicht finden wird.}
\EdN{N4}{AK07-J0071-N04}{\textsuperscript{2}) Vergl. die vorhergehende Note.}
\FrameHead{0072}{p. 40}{0072}{p. 40}
\EdL{1}{AK07-J0072-L01}{Kalender nach Verfluss von 28 Jahren die Tage eines Monats wieder auf dieselben}{}
\EdL{2}{AK07-J0072-L02}{Wochentage fallen wie vorher, dass also in Tafel 3\textsuperscript{a} unter \textit{Anni romanorum plani}}{}
\EdL{3}{AK07-J0072-L03}{nur die Reste verzeichnet sind, die sich bei der Division der Jahreszahl durch 28}{}
\EdL{4}{AK07-J0072-L04}{ergeben, statt des Restes 0 steht die Zahl 28 am Ende der Kolonne. Um diese Divi-}{}
\EdL{5}{AK07-J0072-L05}{sion jeweilen zu ersparen, befindet sich in Tafel 3 eine Kolonne (die letzte) mit den}{}
\EdL{6}{AK07-J0072-L06}{Vielfachen von 28, dieselben Vielfachen in bedeutend grösserer Zahl findet man auch}{}
\EdL{7}{AK07-J0072-L07}{auf Tafel 1. Ich gebe drei Beispiele\,:}{}
\EdLp{8}{AK07-J0072-L08}{1. Welcher Wochentag war der 5. März (syr. \textit{ādhār}, in unsern Tafeln geschrieben}{}
\EdL{9}{AK07-J0072-L09}{adar) des Jahres 1363 der \textit{römischen} Aera\,?}{}
\EdLp{10}{AK07-J0072-L10}{In Tafel 1 findet man die nächste durch 28 teilbare Zahl unter 1363 gleich}{}
\EdL{11}{AK07-J0072-L11}{1344, dieses von 1363 abgezogen, bleibt 19, in Tafel 3\textsuperscript{a} findet man zu 19 in der}{}
\EdL{12}{AK07-J0072-L12}{Kolonne „\textit{Mart.}“ überschrieben die \textit{Nota} 1, d. h. der 1. März fiel auf einen Sonntag,}{}
\EdL{13}{AK07-J0072-L13}{der 5. also auf einen \textit{Donnerstag}. Dieses Jahr der römischen Aera war ein Schaltjahr,}{}
\EdL{14}{AK07-J0072-L14}{denn in der Kolonne \textit{Anni bissextiles} überschrieben steht \textit{bissext.}, in der Tat, denn}{}
\EdL{15}{AK07-J0072-L15}{es entspricht dem Jahre 1052 der julianischen Zeitrechnung.}{}
\EdLp{16}{AK07-J0072-L16}{2. Welcher Wochentag war der 26. Januar 1126 der \textit{julianischen} Zeitrechnung\,?}{}
\EdL{17}{AK07-J0072-L17}{Wir addiren zu 1126 die Zahl 311, dies giebt 1437, in Tafel 3 findet man die nächste}{}
\EdL{18}{AK07-J0072-L18}{durch 28 teilbare Zahl gleich 1428, diese von 1437 abgezogen, bleibt 9, in Tafel 3\textsuperscript{a}}{}
\EdL{19}{AK07-J0072-L19}{findet man zu 9 in der Kolonne Januar überschrieben die \textit{Nota} 6, d. h. der 1. Januar}{}
\EdL{20}{AK07-J0072-L20}{fiel auf einen Freitag, also auch der 8., 15., 22., der 26. also auf einen \textit{Dienstag}.}{}
\EdL{21}{AK07-J0072-L21}{Dieses Resultat entspricht der Angabe, die sich p. 5 des Textes befindet (siehe auch}{}
\EdL{22}{AK07-J0072-L22}{oben p. 37).}{}
\EdLp{23}{AK07-J0072-L23}{3. Welcher Wochentag war der 10. Oktober des Jahres 1912 der \textit{julianischen}}{}
\EdL{24}{AK07-J0072-L24}{Zeitrechnung\,?}{}
\EdLp{25}{AK07-J0072-L25}{Wir addieren zu 1912 die Zahl 312 (siehe p. 39), dies gibt 2224, die nächste}{}
\EdL{26}{AK07-J0072-L26}{durch 28 teilbare Zahl ist 2212 (diese Zahl steht nicht mehr in unsern Tafeln), Rest}{}
\EdL{27}{AK07-J0072-L27}{12, dazu gehört in Tafel 3\textsuperscript{a} unter Oktober die Note 2, also fiel der 1. Oktober auf}{}
\EdL{28}{AK07-J0072-L28}{einen Montag, ebenso der 8., der 10. also auf einen \textit{Mittwoch}.}{}
\EdLp{29}{AK07-J0072-L29}{Für die richtigen Lesarten der syrischen Monatsnamen, die in den Überschriften}{}
\EdL{30}{AK07-J0072-L30}{der einzelnen Kolonnen an oberster Stelle stehen, verweise ich auf das Glossar.}{}
\EdLp{31}{AK07-J0072-L31}{Diese sechs ersten Tafeln (1, 2, 2\textsuperscript{a}, 2\textsuperscript{b}, 3, 3\textsuperscript{a}) entsprechen, einzelne ganz, andere}{}
\EdL{32}{AK07-J0072-L32}{nur teilweise den unechten, d. h. nicht von \Sperr{Batt.} stammenden Tafeln, die \Sperr{Nallino}}{}
\EdL{33}{AK07-J0072-L33}{am Ende seiner Ausgabe (II, p. 300—305) angefügt hat, und die jedenfalls, was bei}{}
\EdL{34}{AK07-J0072-L34}{einigen direkt angegeben ist, von \Sperr{Maslama al-Madjrīṭī} herstammen. Dies bildete}{}
\EdL{35}{AK07-J0072-L35}{natürlich einen wesentlichen Grund für unsere Vermutung, dass wir in diesem Werke}{}
\EdL{36}{AK07-J0072-L36}{nicht die ursprünglichen Tafeln \Sperr{Kh.}’s, sondern die Bearbeitung \Sperr{Maslamas} vor uns}{}
\EdL{37}{AK07-J0072-L37}{haben. Die Entsprechung dieser Tafeln ist folgende\,:}{}
\EdL{38}{AK07-J0072-L38}{\parbox[t]{\dimexpr0.3267\linewidth-0.5em\relax}{\raggedright\small\mbox{}\Sperr{Batt.} p. 301 (1. Hälfte)}\hspace{0.5em}\parbox[t]{\dimexpr0.3267\linewidth-0.5em\relax}{\raggedright\small\mbox{}=}\hspace{0.5em}\parbox[t]{\dimexpr0.3267\linewidth-0.5em\relax}{\raggedright\small\mbox{}\Sperr{Kh.} Tafel 2}\hspace{0.5em}}{}
\EdL{39}{AK07-J0072-L39}{\parbox[t]{\dimexpr0.3267\linewidth-0.5em\relax}{\raggedright\small\mbox{}„ p. 301 (2. Hälfte)}\hspace{0.5em}\parbox[t]{\dimexpr0.3267\linewidth-0.5em\relax}{\raggedright\small\mbox{}=}\hspace{0.5em}\parbox[t]{\dimexpr0.3267\linewidth-0.5em\relax}{\raggedright\small\mbox{}„ „ 2\textsuperscript{a}}\hspace{0.5em}}{}
\EdL{40}{AK07-J0072-L40}{\parbox[t]{\dimexpr0.3267\linewidth-0.5em\relax}{\raggedright\small\mbox{}„ p. 302}\hspace{0.5em}\parbox[t]{\dimexpr0.3267\linewidth-0.5em\relax}{\raggedright\small\mbox{}=}\hspace{0.5em}\parbox[t]{\dimexpr0.3267\linewidth-0.5em\relax}{\raggedright\small\mbox{}„ „ 3\textsuperscript{a} und linke Hälfte von 1}\hspace{0.5em}}{}
\FrameHead{0073}{p. 41}{0073}{p. 41}
\EdL{1}{AK07-J0073-L01}{\parbox[t]{\dimexpr0.3267\linewidth-0.5em\relax}{\raggedright\small\mbox{}\Sperr{Batt.} p. 303}\hspace{0.5em}\parbox[t]{\dimexpr0.3267\linewidth-0.5em\relax}{\raggedright\small\mbox{}=}\hspace{0.5em}\parbox[t]{\dimexpr0.3267\linewidth-0.5em\relax}{\raggedright\small\mbox{}\Sperr{Kh.} Tafel 2\textsuperscript{b}}\hspace{0.5em}}{}
\EdL{2}{AK07-J0073-L02}{\parbox[t]{\dimexpr0.3267\linewidth-0.5em\relax}{\raggedright\small\mbox{}„ p. 304—305}\hspace{0.5em}\parbox[t]{\dimexpr0.3267\linewidth-0.5em\relax}{\raggedright\small\mbox{}=}\hspace{0.5em}\parbox[t]{\dimexpr0.3267\linewidth-0.5em\relax}{\raggedright\small\mbox{}„ „ 3 (Hier sind aber nur die \textit{anni Romanorum}}\hspace{0.5em}}{}
\EdL{3}{AK07-J0073-L03}{\parbox[t]{\dimexpr0.6533\linewidth-0.5em\relax}{\raggedright\small\mbox{}}\hspace{0.5em}\parbox[t]{\dimexpr0.3267\linewidth-0.5em\relax}{\raggedright\small\mbox{}aufgenommen, die \textit{anni Coptorum} und}\hspace{0.5em}}{}
\EdL{4}{AK07-J0073-L04}{\parbox[t]{\dimexpr0.6533\linewidth-0.5em\relax}{\raggedright\small\mbox{}}\hspace{0.5em}\parbox[t]{\dimexpr0.3267\linewidth-0.5em\relax}{\raggedright\small\mbox{}\textit{Persarum} nicht (s. Einleitung)).}\hspace{0.5em}}{}
\EdLc{5}{AK07-J0073-L05}{\textbf{Kap. 6. Über die Einteilung der Sphäre.}}{}
\EdLc{6}{AK07-J0073-L06}{(Text p. 7.)}{}
\EdLp{7}{AK07-J0073-L07}{Dieses kurze Kapitel bedarf keines Kommentars.}{}
\EdLc{8}{AK07-J0073-L08}{\textbf{Kap. 7. Über die mittlere Bewegung der Planeten.}}{}
\EdLc{9}{AK07-J0073-L09}{(Text p. 7—9, Tafeln 4—20.)}{}
\EdLp{10}{AK07-J0073-L10}{Dieses Kapitel mit seinen Tafeln nimmt in jedem astronomischen Werk der}{}
\EdL{11}{AK07-J0073-L11}{Griechen und Araber eine Hauptstelle ein. Mit ihrer Hilfe berechnet man die mittlern}{}
\EdL{12}{AK07-J0073-L12}{Örter der Sonne, des Mondes und der fünf Planeten, und zwar zuerst mit Hilfe der}{}
\EdL{13}{AK07-J0073-L13}{Kolonne \textit{anni collecti} überschrieben die Örter von je 30 zu 30 Jahren, dann für die}{}
\EdL{14}{AK07-J0073-L14}{\textit{anni plani} oder \textit{singuli}, hierauf für die Monate, Tage und Stunden, ja sogar für je}{}
\EdL{15}{AK07-J0073-L15}{2 zu 2 Minuten\,; man kann also für jeden dieser sieben Himmelskörper seine mittlere}{}
\EdL{16}{AK07-J0073-L16}{Länge, d. h. seinen Abstand vom Frühlingspunkt für jeden beliebigen Zeitpunkt bis}{}
\EdL{17}{AK07-J0073-L17}{auf zwei Minuten hinunter bestimmen. Da für die Bestimmung der Breite des Mondes}{}
\EdL{18}{AK07-J0073-L18}{und auch für die Astrologie die Stellung der Knoten der Mondbahn von Bedeutung}{}
\EdL{19}{AK07-J0073-L19}{war, so wird neben Sonne, Mond und den fünf Planeten auch die mittlere Bewegung}{}
\EdL{20}{AK07-J0073-L20}{des nördlichen oder aufsteigenden Mondknotens (arab. \textit{al-ʿaḳd al-shamālī}) angegeben.}{}
\EdL{21}{AK07-J0073-L21}{In unsern Tafeln wird er mit dem persisch-arabischen Wort \textit{gavzehar} (= \textit{djauzahar}),}{}
\EdL{22}{AK07-J0073-L22}{lateinisch durch \textit{draco} (besser wäre \textit{caput draconis} = Drachenkopf) bezeichnet.}{}
\EdLp{23}{AK07-J0073-L23}{Bei \Sperr{Batt.} heisst die mittlere Beweguug \textit{al-masīr al-wasaṭ} (= der mittlere Lauf),}{}
\EdL{24}{AK07-J0073-L24}{oder auch nur \textit{al-wasaṭ} (das Mittel), oder sogar nur \textit{al-masīr} (Lauf, Weg, Reise)\,; in}{}
\EdL{25}{AK07-J0073-L25}{unsern Tafeln steht teils das arabische \textit{wasaṭ}, geschrieben \textit{wazat} und \textit{wacat}, teils}{}
\EdL{26}{AK07-J0073-L26}{\textit{medialitas}, oder \textit{mediocritas}, teils \textit{medius cursus}.}{}
\EdLp{27}{AK07-J0073-L27}{Der mittlere Ort ist jeweilen angegeben in Tierkreiszeichen (\textit{Signa}), Graden,}{}
\EdL{28}{AK07-J0073-L28}{Minuten und Sekunden\,; \Sperr{Batt.} hat die Signa in Grade verwandelt, und vernachlässigt}{}
\EdL{29}{AK07-J0073-L29}{die Sekunden bei den fünf Planeten, während er sie bei der Sonne und beim Monde}{}
\EdL{30}{AK07-J0073-L30}{ebenfalls angibt.}{}
\EdLp{31}{AK07-J0073-L31}{Die \textit{Radix}, die im Anfang der \textit{anni collecti} verzeichnet ist, gibt, wie oben schon}{}
\EdL{32}{AK07-J0073-L32}{bemerkt wurde, die Länge der Planeten am Mittag des 1. \textit{Muḥarram} des Jahres 1 d. H.}{}
\EdL{33}{AK07-J0073-L33}{für den Meridian von \textit{Udjain} an\,; aber wenn man auch den Längenunterschied}{}
\EdL{34}{AK07-J0073-L34}{zwischen \textit{Udjain} (c. 75° 52′ Gr.) und \textit{Raḳḳa} (c. 39° 4′ Gr.) in Rechnung bringt, so stimmt}{}
\EdL{35}{AK07-J0073-L35}{die Radix doch nicht mit derjenigen \Sperr{Batt.}’s. Unsere Tafeln haben z. B. für die Sonne}{}
\EdL{36}{AK07-J0073-L36}{die Radix 113° 25′ 48″, \Sperr{Batt.} dagegen 113° 58′ 4″\,; im Anfang des Jahres 511 d. H.}{}
\EdL{37}{AK07-J0073-L37}{macht der Unterschied c. 9° 15′ aus\,; dies sind jedenfalls die 9°, von denen \Sperr{Ibn ʿEzrā}}{}
\EdL{38}{AK07-J0073-L38}{in der Vorrede zur Übersetzung der \textit{Gründe der Tafeln des} \Sperr{Khwārizmī} spricht\textsuperscript{1})\,;}{}
\NoteRule
\EdN{N1}{AK07-J0073-N01}{\textsuperscript{1}) Vergl. \Sperr{Steinschneider}\,: \textit{Zur Geschichte der Übersetzungen aus dem Indischen ins Arabische},}
\EdN{N2}{AK07-J0073-N02}{in Zeitschr. d. deutschen morgenländ. Gesellsch. 24. Bd., p. 355.}
\FrameHead{0074}{p. 42}{0074}{p. 42}
\EdL{1}{AK07-J0074-L01}{auch bei den Längen der Planeten finden wir diesen Unterschied von 8\textsuperscript{1}/\textsubscript{2} bis 9\textsuperscript{1}/\textsubscript{2}°}{}
\EdL{2}{AK07-J0074-L02}{wieder, mit Ausnahme des Mars, wo der Unterschied 11\textsuperscript{1}/\textsubscript{2}° ist\,; wahrscheinlich liegen}{}
\EdL{3}{AK07-J0074-L03}{hier Fehler der Abschreiber vor. Diese Unterschiede rühren von der Nichtbeachtung}{}
\EdL{4}{AK07-J0074-L04}{der \textit{Praezession} des Frühlingspunktes her\,; in der Tat legte \Sperr{Kh.} für die Bestimmung}{}
\EdL{5}{AK07-J0074-L05}{des mittlern Laufes der Sonne das \textit{Sternjahr} und nicht das \textit{tropische} Jahr, wie \Sperr{Batt.}}{}
\EdL{6}{AK07-J0074-L06}{zu Grunde\,; wenn wir nämlich nach Tafel 4 und 5 die Zeit bestimmen, die die Sonne}{}
\EdL{7}{AK07-J0074-L07}{braucht, um einen ganzen Umkreis (= 360°) zu durchlaufen, so finden wir\,:}{}
\EdL{8}{AK07-J0074-L08}{\parbox[t]{\dimexpr0.2450\linewidth-0.5em\relax}{\raggedright\small\mbox{}In 12 · 30 =}\hspace{0.5em}\parbox[t]{\dimexpr0.2450\linewidth-0.5em\relax}{\raggedright\small\mbox{}360 Tagen}\hspace{0.5em}\parbox[t]{\dimexpr0.2450\linewidth-0.5em\relax}{\raggedright\small\mbox{}= 12 · (29° 34′ 5″) =}\hspace{0.5em}\parbox[t]{\dimexpr0.2450\linewidth-0.5em\relax}{\raggedright\small\mbox{}354° 49′}\hspace{0.5em}}{}
\EdL{9}{AK07-J0074-L09}{\parbox[t]{\dimexpr0.2450\linewidth-0.5em\relax}{\raggedright\small\mbox{}„}\hspace{0.5em}\parbox[t]{\dimexpr0.2450\linewidth-0.5em\relax}{\raggedright\small\mbox{}5 „}\hspace{0.5em}\parbox[t]{\dimexpr0.2450\linewidth-0.5em\relax}{\raggedright\small\mbox{}=}\hspace{0.5em}\parbox[t]{\dimexpr0.2450\linewidth-0.5em\relax}{\raggedright\small\mbox{}4° 55′ 41″}\hspace{0.5em}}{}
\EdL{10}{AK07-J0074-L10}{\parbox[t]{\dimexpr0.2450\linewidth-0.5em\relax}{\raggedright\small\mbox{}„}\hspace{0.5em}\parbox[t]{\dimexpr0.2450\linewidth-0.5em\relax}{\raggedright\small\mbox{}6 Stdn.}\hspace{0.5em}\parbox[t]{\dimexpr0.2450\linewidth-0.5em\relax}{\raggedright\small\mbox{}=}\hspace{0.5em}\parbox[t]{\dimexpr0.2450\linewidth-0.5em\relax}{\raggedright\small\mbox{}14′ 47″}\hspace{0.5em}}{}
\EdL{11}{AK07-J0074-L11}{\parbox[t]{\dimexpr0.2450\linewidth-0.5em\relax}{\raggedright\small\mbox{}„}\hspace{0.5em}\parbox[t]{\dimexpr0.2450\linewidth-0.5em\relax}{\raggedright\small\mbox{}13 Min.}\hspace{0.5em}\parbox[t]{\dimexpr0.2450\linewidth-0.5em\relax}{\raggedright\small\mbox{}=}\hspace{0.5em}\parbox[t]{\dimexpr0.2450\linewidth-0.5em\relax}{\raggedright\small\mbox{}32″}\hspace{0.5em}}{}
\EdL{12}{AK07-J0074-L12}{\parbox[t]{\dimexpr0.2450\linewidth-0.5em\relax}{\raggedright\small\mbox{}also in}\hspace{0.5em}\parbox[t]{\dimexpr0.2450\linewidth-0.5em\relax}{\raggedright\small\mbox{}365\textsuperscript{d} 6\textsuperscript{h} 13\textsuperscript{m}}\hspace{0.5em}\parbox[t]{\dimexpr0.2450\linewidth-0.5em\relax}{\raggedright\small\mbox{}=}\hspace{0.5em}\parbox[t]{\dimexpr0.2450\linewidth-0.5em\relax}{\raggedright\small\mbox{}360° 0′ 0″}\hspace{0.5em}}{}
\EdL{13}{AK07-J0074-L13}{\Sperr{Kh.} nahm also die Länge des Sternjahres zu 365\textsuperscript{d} 6\textsuperscript{h} 13\textsuperscript{m} an\,; wir finden im \textit{Sūrya-}}{}
\EdL{14}{AK07-J0074-L14}{\textit{Siddhānta}\textsuperscript{1}) für die Länge des Sternjahres 365\textsuperscript{d} 6\textsuperscript{h} 12,6\textsuperscript{m}, was ungefähr mit unserer}{}
\EdL{15}{AK07-J0074-L15}{Berechnung stimmt\textsuperscript{2}). — \Sperr{Kh.} erwähnt das tropische Jahr gar nicht in seinem Werke,}{}
\EdL{16}{AK07-J0074-L16}{und es scheint in der Tat, dass auch die Inder bis und mit \Sperr{Brahmagupta} es nicht}{}
\EdL{17}{AK07-J0074-L17}{gekannt haben. Das, was im \textit{Sūrya-Siddhānta}\textsuperscript{3}) über die Praezession gesagt ist, soll}{}
\EdL{18}{AK07-J0074-L18}{nach \Sperr{Burgess} und \Sperr{Thibaut}\textsuperscript{4}) eine spätere Einschiebung sein\,; dagegen lehrte aller-}{}
\EdL{19}{AK07-J0074-L19}{dings Āryabhaṭṭa am Ende des 5. Jahrh. n. Chr. die \textit{Trepidation} der Fixsterne, bezw.}{}
\EdL{20}{AK07-J0074-L20}{des Frühlingspunktes.}{}
\EdLp{21}{AK07-J0074-L21}{Als Beispiele für den Gebrauch der Tafeln geben wir die Berechnung des}{}
\EdL{22}{AK07-J0074-L22}{mittlern Ortes der Sonne und des Mondes für Mittag den \textit{15. Ṣafar} des J. 540 d. H.}{}
\EdL{23}{AK07-J0074-L23}{\parbox[t]{\dimexpr0.2450\linewidth-0.5em\relax}{\raggedright\small\mbox{}\textit{1. Sonne}\,:}\hspace{0.5em}\parbox[t]{\dimexpr0.2450\linewidth-0.5em\relax}{\raggedright\small\mbox{}Länge am Ende des 510. Jahres}\hspace{0.5em}\parbox[t]{\dimexpr0.2450\linewidth-0.5em\relax}{\raggedright\small\mbox{}=}\hspace{0.5em}\parbox[t]{\dimexpr0.2450\linewidth-0.5em\relax}{\raggedright\small\mbox{}38° 35′ 45″}\hspace{0.5em}}{}
\EdL{24}{AK07-J0074-L24}{\parbox[t]{\dimexpr0.2450\linewidth-0.5em\relax}{\raggedright\small\mbox{}}\hspace{0.5em}\parbox[t]{\dimexpr0.2450\linewidth-0.5em\relax}{\raggedright\small\mbox{}Zunahme in 29 Jahren}\hspace{0.5em}\parbox[t]{\dimexpr0.2450\linewidth-0.5em\relax}{\raggedright\small\mbox{}=}\hspace{0.5em}\parbox[t]{\dimexpr0.2450\linewidth-0.5em\relax}{\raggedright\small\mbox{}49° 2′ 50″}\hspace{0.5em}}{}
\EdL{25}{AK07-J0074-L25}{\parbox[t]{\dimexpr0.2450\linewidth-0.5em\relax}{\raggedright\small\mbox{}}\hspace{0.5em}\parbox[t]{\dimexpr0.2450\linewidth-0.5em\relax}{\raggedright\small\mbox{}„ „ 1 Monat}\hspace{0.5em}\parbox[t]{\dimexpr0.2450\linewidth-0.5em\relax}{\raggedright\small\mbox{}=}\hspace{0.5em}\parbox[t]{\dimexpr0.2450\linewidth-0.5em\relax}{\raggedright\small\mbox{}29° 34′ 5″}\hspace{0.5em}}{}
\EdL{26}{AK07-J0074-L26}{\parbox[t]{\dimexpr0.2450\linewidth-0.5em\relax}{\raggedright\small\mbox{}}\hspace{0.5em}\parbox[t]{\dimexpr0.2450\linewidth-0.5em\relax}{\raggedright\small\mbox{}„ „ 14 Tagen}\hspace{0.5em}\parbox[t]{\dimexpr0.2450\linewidth-0.5em\relax}{\raggedright\small\mbox{}=}\hspace{0.5em}\parbox[t]{\dimexpr0.2450\linewidth-0.5em\relax}{\raggedright\small\mbox{}13° 47′ 54″}\hspace{0.5em}}{}
\EdL{27}{AK07-J0074-L27}{\parbox[t]{\dimexpr0.2450\linewidth-0.5em\relax}{\raggedright\small\mbox{}}\hspace{0.5em}\parbox[t]{\dimexpr0.2450\linewidth-0.5em\relax}{\raggedright\small\mbox{}Summa}\hspace{0.5em}\parbox[t]{\dimexpr0.2450\linewidth-0.5em\relax}{\raggedright\small\mbox{}=}\hspace{0.5em}\parbox[t]{\dimexpr0.2450\linewidth-0.5em\relax}{\raggedright\small\mbox{}131° 0′ 34″}\hspace{0.5em}}{}
\EdL{28}{AK07-J0074-L28}{Also war die mittlere Länge der Sonne zur genannten Zeit = 131° 0′ 34″.}{}
\EdL{29}{AK07-J0074-L29}{\parbox[t]{\dimexpr0.2450\linewidth-0.5em\relax}{\raggedright\small\mbox{}\textit{2. Mond}\,:}\hspace{0.5em}\parbox[t]{\dimexpr0.2450\linewidth-0.5em\relax}{\raggedright\small\mbox{}Länge am Ende des 510. Jahres}\hspace{0.5em}\parbox[t]{\dimexpr0.2450\linewidth-0.5em\relax}{\raggedright\small\mbox{}=}\hspace{0.5em}\parbox[t]{\dimexpr0.2450\linewidth-0.5em\relax}{\raggedright\small\mbox{}40° 56′ 38″}\hspace{0.5em}}{}
\EdL{30}{AK07-J0074-L30}{\parbox[t]{\dimexpr0.2450\linewidth-0.5em\relax}{\raggedright\small\mbox{}}\hspace{0.5em}\parbox[t]{\dimexpr0.2450\linewidth-0.5em\relax}{\raggedright\small\mbox{}Zunahme in 29 Jahren}\hspace{0.5em}\parbox[t]{\dimexpr0.2450\linewidth-0.5em\relax}{\raggedright\small\mbox{}=}\hspace{0.5em}\parbox[t]{\dimexpr0.2450\linewidth-0.5em\relax}{\raggedright\small\mbox{}53° 24′ 16″}\hspace{0.5em}}{}
\EdL{31}{AK07-J0074-L31}{\parbox[t]{\dimexpr0.2450\linewidth-0.5em\relax}{\raggedright\small\mbox{}}\hspace{0.5em}\parbox[t]{\dimexpr0.2450\linewidth-0.5em\relax}{\raggedright\small\mbox{}„ „ 1 Monat}\hspace{0.5em}\parbox[t]{\dimexpr0.2450\linewidth-0.5em\relax}{\raggedright\small\mbox{}=}\hspace{0.5em}\parbox[t]{\dimexpr0.2450\linewidth-0.5em\relax}{\raggedright\small\mbox{}35° 17′ 26″}\hspace{0.5em}}{}
\EdL{32}{AK07-J0074-L32}{\parbox[t]{\dimexpr0.2450\linewidth-0.5em\relax}{\raggedright\small\mbox{}}\hspace{0.5em}\parbox[t]{\dimexpr0.2450\linewidth-0.5em\relax}{\raggedright\small\mbox{}„ „ 14 Tagen}\hspace{0.5em}\parbox[t]{\dimexpr0.2450\linewidth-0.5em\relax}{\raggedright\small\mbox{}=}\hspace{0.5em}\parbox[t]{\dimexpr0.2450\linewidth-0.5em\relax}{\raggedright\small\mbox{}184° 28′ 8″}\hspace{0.5em}}{}
\EdL{33}{AK07-J0074-L33}{\parbox[t]{\dimexpr0.2450\linewidth-0.5em\relax}{\raggedright\small\mbox{}}\hspace{0.5em}\parbox[t]{\dimexpr0.2450\linewidth-0.5em\relax}{\raggedright\small\mbox{}Summa\,:}\hspace{0.5em}\parbox[t]{\dimexpr0.2450\linewidth-0.5em\relax}{\raggedright\small\mbox{}=}\hspace{0.5em}\parbox[t]{\dimexpr0.2450\linewidth-0.5em\relax}{\raggedright\small\mbox{}314° 6′ 28″}\hspace{0.5em}}{}
\EdL{34}{AK07-J0074-L34}{Also war die mittlere Länge des Mondes zur genannten Zeit = 314° 6′ 28″.}{}
\NoteRule
\EdN{N1}{AK07-J0074-N01}{\textsuperscript{1}) E. \Sperr{Burgess}, \textit{Translation of the S. S.}, im Journal of the American Oriental Society, Vol. VI,}
\EdN{N2}{AK07-J0074-N02}{1860, p. 161.}
\EdN{N3}{AK07-J0074-N03}{\textsuperscript{2}) S. unten Kap. 21 u. 22.}
\EdN{N4}{AK07-J0074-N04}{\textsuperscript{3}) Transl. \Sperr{Burgess}, l. c. p. 244—249.}
\EdN{N5}{AK07-J0074-N05}{\textsuperscript{4}) \textit{Astronomie, Astrologie u. Mathematik}, im Grundriss der Indo-Arischen Philologie u. Alter-}
\EdN{N6}{AK07-J0074-N06}{tumskunde, 3. Bd. 9. Heft, p. 45–46.}
\FrameHead{0075}{p. 43}{0075}{p. 43}
\EdLp{1}{AK07-J0075-L01}{Es wird im Text am Schlusse der Regel bemerkt, wie man zu verfahren habe,}{}
\EdL{2}{AK07-J0075-L02}{wenn der Ort, für den man die mittlere Länge eines Planeten bestimmen will, östlich}{}
\EdL{3}{AK07-J0075-L03}{oder westlich von \textit{Arin} liege.}{}
\EdLp{4}{AK07-J0075-L04}{Am Schlusse des Kapitels findet sich ein Zusatz, der in C und O nur am Rande,}{}
\EdL{5}{AK07-J0075-L05}{dagegen in M im Texte steht, und der sich über das Verhältnis von Graden und}{}
\EdL{6}{AK07-J0075-L06}{Meilen ausspricht\,; es wird angegeben, dass auf der Erde auf einen Grad 66\textsuperscript{2}/\textsubscript{3} Meilen}{}
\EdL{7}{AK07-J0075-L07}{kommen, dass also der Umfang der Erde 24000 Meilen sei\,; anfänglich ist bemerkt,}{}
\EdL{8}{AK07-J0075-L08}{dass nach den Chaldäern 4000 Kameelschritte eine Meile ausmachen.}{}
\EdLp{9}{AK07-J0075-L09}{Diese Meile, von denen 66\textsuperscript{2}/\textsubscript{3} auf einen Grad gehen sollen, ist die von den}{}
\EdL{10}{AK07-J0075-L10}{arabischen Geographen und Astronomen in ihren Angaben gewöhnlich gebrauchte}{}
\EdL{11}{AK07-J0075-L11}{Meile zu 4000 Ellen (\textit{dhirāʿ}), die Elle zu 24 Fingern (\textit{aṣbaʿ} oder \textit{iṣbaʿ})\,; nach \Sperr{Nallino}\textsuperscript{1})}{}
\EdL{12}{AK07-J0075-L12}{war eine solche Meile c. 1973 m. Nach der Gradmessung unter \Sperr{al-Ma’mūn} kam}{}
\EdL{13}{AK07-J0075-L13}{dann neben diesem Werte auch der bessere von 56\textsuperscript{2}/\textsubscript{3} Meilen = 1° in Gebrauch.}{}
\EdLp{14}{AK07-J0075-L14}{M. \Sperr{Chasles}\textsuperscript{2}) schliesst aus obiger Stelle, dass schon die \textit{Chaldäer} (Babylonier)}{}
\EdL{15}{AK07-J0075-L15}{den nahezu richtigen Umfang der Erde gekannt hätten. Wir wollen dies nicht}{}
\EdL{16}{AK07-J0075-L16}{absolut bestreiten, aber erstens ist die Angabe, dass es für eine Meile von c. 2000}{}
\EdL{17}{AK07-J0075-L17}{Meter 4000 Kameelschritte brauche, schon etwas verdächtig, und zweitens sind hier}{}
\EdL{18}{AK07-J0075-L18}{unter \textit{Caldei} höchst wahrscheinlich die \textit{Araber} verstanden, denn wir halten diese}{}
\EdL{19}{AK07-J0075-L19}{Stelle für eine Einschiebung des Übersetzers, oder von Abschreibern\,; p. 11 des}{}
\EdL{20}{AK07-J0075-L20}{Textes kommt nämlich in einer Stelle, die unzweifelhaft eine Einschiebung ist, noch-}{}
\EdL{21}{AK07-J0075-L21}{mals das Wort \textit{Caldei} vor, und dort bezeichnet es sicher die Araber\,; das Ms. O hat}{}
\EdL{22}{AK07-J0075-L22}{auch in der Tat statt „a Caldeis“ „arabice“.}{}
\EdLc{23}{AK07-J0075-L23}{\textbf{Kap. 8. Über die Auffindung des (wahren) Ortes der Sonne.}}{}
\EdLc{24}{AK07-J0075-L24}{(Text p. 9—10, Tafel 21—26.)}{}
\EdLp{25}{AK07-J0075-L25}{Um den wahren Ort der Sonne aus dem mittlern zu bestimmen, muss man}{}
\EdL{26}{AK07-J0075-L26}{nach der Vorschrift dieses Kapitels von der mittlern Länge der Sonne die Länge des}{}
\EdL{27}{AK07-J0075-L27}{\textit{Apogeums} derselben, die in diesen Tafeln zu 77° 55′ angenommen wird\textsuperscript{3}), subtrahieren,}{}
\EdL{28}{AK07-J0075-L28}{dies gibt die sogenannte \textit{Anomalie} der Sonne\,; zu dieser als Argument sucht man in}{}
\EdL{29}{AK07-J0075-L29}{den Tafeln 21—26 die sogenannte \textit{Gleichung} der Sonne\,; diese Gleichung muss man}{}
\EdL{30}{AK07-J0075-L30}{von der mittlern Länge der Sonne subtrahieren, bezw. zu ihr addieren, je nachdem}{}
\EdL{31}{AK07-J0075-L31}{die Anomalie kleiner oder grösser als 180° ist, das Resultat gibt die \textit{wahre Länge}}{}
\EdL{32}{AK07-J0075-L32}{der Sonne.}{}
\EdLp{33}{AK07-J0075-L33}{\textit{Beispiel}\,: Im vorigen Kapitel fanden wir die mittlere Länge der Sonne für den}{}
\EdL{34}{AK07-J0075-L34}{Mittag des \textit{15. Ṣafar} des Jahres 540 d. H. = 131° 0′ 34″. Subtrahieren wir davon das}{}
\NoteRule
\EdN{N1}{AK07-J0075-N01}{\textsuperscript{1}) \textit{Il valore metrico del grado di meridiano secondo i geografi arabi}, im Cosmos di Guido Cora,}
\EdN{N2}{AK07-J0075-N02}{Vol. XI, 1892—93, fascic. I—IV, u. \textit{Estratto}, Torino, 1893.}
\EdN{N3}{AK07-J0075-N03}{\textsuperscript{2}) \textit{Recherches sur l’astronomie ancienne}, in Comptes-rendus de l’Académie des Sciences, t. XXIII,}
\EdN{N4}{AK07-J0075-N04}{Paris, 1846, p. 851.}
\EdN{N5}{AK07-J0075-N05}{\textsuperscript{3}) Diese Zahl kann, wenn die Zahlen durch arab. Buchstaben ausgedrückt sind, leicht mit}
\EdN{N6}{AK07-J0075-N06}{77° 15′ verwechselt werden, und dann würde der Wert mit dem der spätern Redaktion des \textit{Sūrya-}}
\EdN{N7}{AK07-J0075-N07}{\textit{Siddhānta} und des arab. \textit{Sindhind} (nach dem Zeugnis \Sperr{Ibn Hibintās}) stimmen (vergl. \Sperr{Batt.}, ed.}
\EdN{N8}{AK07-J0075-N08}{\Sperr{Nallino}, I, 218).}
\FrameHead{0076}{p. 44}{0076}{p. 44}
\EdL{1}{AK07-J0076-L01}{\textit{Apogeum} = 77° 55′, so erhalten wir die \textit{Anomalie} = 53° 5′ 34″, mit diesem Argument}{}
\EdL{2}{AK07-J0076-L02}{tritt man in die Tafel 22 ein und findet hier für die \textit{Gleichung} der Sonne mit Hilfe}{}
\EdL{3}{AK07-J0076-L03}{der Interpolation 1° 46′, diese Gleichung von der mittlern Länge subtrahiert, gibt}{}
\EdL{4}{AK07-J0076-L04}{als \textit{wahren Ort} der Sonne 129° 14′ 34″, d. h. die Sonne befand sich in diesem Zeit-}{}
\EdL{5}{AK07-J0076-L05}{punkt in 9° 14′ 34″ des Löwen.}{}
\EdLp{6}{AK07-J0076-L06}{Die Erklärung, wie man bei allfälligen Minuten der Anomalie zu verfahren habe,}{}
\EdL{7}{AK07-J0076-L07}{ist in unserm Texte (p. 9—10\,: \textit{Quod si ab argumento partes aliquae restiterint} . . .) ver-}{}
\EdL{8}{AK07-J0076-L08}{dorben\,; es ist natürlich dieselbe Berechnung von Proportionalteilen auszuführen,}{}
\EdL{9}{AK07-J0076-L09}{wie bei unsern logarith.-trigonometrischen Tafeln.\textsuperscript{1})}{}
\EdLp{10}{AK07-J0076-L10}{Wir geben im folgenden die geometrische Erklärung zu dem obigen Verfahren}{}
\EdL{11}{AK07-J0076-L11}{der Bestimmung der wahren Länge der Sonne nach \Sperr{Ptol.-Batt.}, und zwar nach}{}
\EdL{12}{AK07-J0076-L12}{der exzentrischen Hypothese, d. h. die Sonne bewege sich in einem zur Ekliptik}{}
\EdL{13}{AK07-J0076-L13}{exzentrischen Kreise, nicht auf einem Epizykel. \Sperr{Ptol.} hat übrigens schon gezeigt,}{}
\EdL{14}{AK07-J0076-L14}{dass man genau zu demselben Resultate komme mit Hilfe eines Epizykels.}{}
\EdLp{15}{AK07-J0076-L15}{Es sei \textit{E} der Mittelpunkt der Ekliptik, also der Erdmittelpunkt, \textit{M} der Mittel-}{}
\EdL{16}{AK07-J0076-L16}{punkt des exzentrischen Kreises, in dem sich die Sonne gleichförmig bewegen soll.}{}
\EdL{17}{AK07-J0076-L17}{\textit{A} ist also das Apogeum, \textit{P} das Perigeum der Sonne\,;}{}
\EdL{18}{AK07-J0076-L18}{diese letztere befinde sich im Punkte \textit{S}, \textit{F} sei der}{}
\EdL{19}{AK07-J0076-L19}{Frühlingspunkt. Nun ist \ensuremath{\sphericalangle} \textit{FMS} die mittlere Länge}{}
\EdL{20}{AK07-J0076-L20}{der Sonne, hievon zieht man die Länge des Apo-}{}
\EdL{21}{AK07-J0076-L21}{geums, d. h. den \ensuremath{\sphericalangle} \textit{FMA} ab, und hat dann den}{}
\EdL{22}{AK07-J0076-L22}{\ensuremath{\sphericalangle} \textit{AMS}, die sogenannte \textit{mittlere Anomalie}, in unsern}{}
\EdL{23}{AK07-J0076-L23}{Tafeln gewöhnlich \textit{Argument} genannt. Wir wollen}{}
\EdL{24}{AK07-J0076-L24}{aber den \ensuremath{\sphericalangle} \textit{AES} haben, unter dem von der Erde (\textit{E})}{}
\EdL{25}{AK07-J0076-L25}{aus gesehen, der Bogen \textit{AS} erscheint, den die Sonne}{}
\EdL{26}{AK07-J0076-L26}{vom Apogeum an durchlaufen hat\,; diesen Winkel erhalten wir, wenn wir von der}{}
\EdL{27}{AK07-J0076-L27}{mittlern Anomalie den \ensuremath{\sphericalangle} \textit{α}, die sogenannte \textit{Gleichung} der Sonne, subtrahieren. Der}{}
\EdL{28}{AK07-J0076-L28}{\ensuremath{\sphericalangle} \textit{AES}, oder auch \textit{AMG}, wenn \textit{MG} ∥ \textit{ES}, heisst nun die \textit{wahre Anomalie}, addiert}{}
\EdL{29}{AK07-J0076-L29}{man dazu wieder \ensuremath{\sphericalangle} \textit{AMF}, d. h. die Länge des Apogeums, so hat man den \ensuremath{\sphericalangle} \textit{FMG}}{}
\EdL{30}{AK07-J0076-L30}{oder Bogen \textit{FG}, d. h. die \textit{wahre Länge} der Sonne. — In unserm Text wird die}{}
\EdL{31}{AK07-J0076-L31}{Gleichung \textit{α} direkt von der mittlern Länge der Sonne abgezogen, was natürlich auf}{}
\EdL{32}{AK07-J0076-L32}{dasselbe hinauskommt.}{}
\EdLp{33}{AK07-J0076-L33}{Wie man für jeden Bogen \textit{AS}, oder für jede mittlere Anomalie, die Gleichung}{}
\EdL{34}{AK07-J0076-L34}{berechnet, will ich hier nicht ausführlich auseinandersetzen, ich bemerke nur, dass}{}
\EdL{35}{AK07-J0076-L35}{\Sperr{Ptol.}\textsuperscript{2}) und \Sperr{Batt.}\textsuperscript{3}) sie mittels des pythagor. Lehrsatzes und trigonometrischer Be-}{}
\EdL{36}{AK07-J0076-L36}{ziehungen aus den Dreiecken \textit{HME} und \textit{HES} berechnen, \Sperr{al-Bīrūnī}\textsuperscript{4}) aber, indem er}{}
\EdL{37}{AK07-J0076-L37}{mit Hilfe des Cosinus-Satzes aus dem Dreieck \textit{EMS} zuerst \textit{ES}, und dann aus dem}{}
\EdL{38}{AK07-J0076-L38}{rechtwinkligen Dreieck \textit{HES} den \ensuremath{\sphericalangle} \textit{α} berechnet.}{}
\NoteRule
\EdN{N1}{AK07-J0076-N01}{\textsuperscript{1}) Vergl. auch \Sperr{Batt.}, ed. \Sperr{Nallino}, I, 72, Z. 9—13.}
\EdN{N2}{AK07-J0076-N02}{\textsuperscript{2}) \textit{Almagest} III, 5 (ed. \Sperr{Heiberg}, I, 240—255, deutsche Übers. von K. \Sperr{Manitius}, I, 173—180).}
\EdN{N3}{AK07-J0076-N03}{\textsuperscript{3}) Ed. \Sperr{Nallino}, I, 48, 219—220.}
\EdN{N4}{AK07-J0076-N04}{\textsuperscript{4}) Vergl. \textit{Biblioth. mathem.} \textbf{11}\textsubscript{3}, 1910, p. 45 u. 71.}
\FrameHead{0077}{p. 45}{0077}{p. 45}
\EdLp{1}{AK07-J0077-L01}{Die Werte der \textit{Gleichung} der Sonne stimmen übrigens mit denjenigen \Sperr{Batt.}’s}{}
\EdL{2}{AK07-J0077-L02}{(ed. \Sperr{Nallino}, II, 78—83), die nach der Formel\,: tg aeq. = \Frac{\textit{e} sin \textit{m}}{1 + \textit{e} cos \textit{m}} berechnet sind,}{}
\EdL{3}{AK07-J0077-L03}{nicht überein, schon für 1° der mittlern Anomalie ist der Unterschied 16″. Dies kann}{}
\EdL{4}{AK07-J0077-L04}{nur von einem andern Werte der Exzentrizität \textit{e} der Sonnenbahn herrühren\,; \Sperr{Batt.}}{}
\EdL{5}{AK07-J0077-L05}{nimmt diese (also in voriger Figur \textit{ME}) zu 2\textsuperscript{p} 4\textsuperscript{3}/\textsubscript{4}′ (wenn \textit{MA} = 60\textsuperscript{p} gesetzt wird) an\,;}{}
\EdL{6}{AK07-J0077-L06}{bei \Sperr{Kh.} findet man keine Angabe für die Exzentrizität der Sonne\,; aus dem Maximum}{}
\EdL{7}{AK07-J0077-L07}{der Gleichung (2° 14′) bei 90° finden wir \textit{e} = 2\textsuperscript{p} 20′, bei andern Proben, die wir}{}
\EdL{8}{AK07-J0077-L08}{gemacht haben, fanden wir sie = 2\textsuperscript{p} 22′\,; der \textit{Sūrya-Siddhānta} hat nach \Sperr{Burgess}\textsuperscript{1})}{}
\EdL{9}{AK07-J0077-L09}{eine maximale Gleichung von 2° 10′ 31″, dies entspricht einer Exzentrizität \textit{e} = 2\textsuperscript{p} 17′\,;}{}
\EdL{10}{AK07-J0077-L10}{\Sperr{Ptol.} nimmt eine Exzentrizität von 2\textsuperscript{p} 29′ 30″ an, rundet sie aber bisweilen zu 2\textsuperscript{p} 30′ auf.}{}
\EdL{11}{AK07-J0077-L11}{Wir sehen also, dass \Sperr{Kh.} weder dem \textit{Sūrya-Siddhānta} noch \Sperr{Ptol.} folgt, vielleicht}{}
\EdL{12}{AK07-J0077-L12}{entnahm er seinen Wert den Werken \Sperr{Brahmaguptas}, oder den persischen Tafeln}{}
\EdL{13}{AK07-J0077-L13}{(s. oben p. 32).}{}
\EdLc{14}{AK07-J0077-L14}{\textbf{Kap. 9. Über die Auffindung des (wahren) Ortes des Mondes.}}{}
\EdLc{15}{AK07-J0077-L15}{(Text p. 10, Tafeln 21—26.)}{}
\EdLp{16}{AK07-J0077-L16}{In den Tafeln 6—8 findet man neben der mittlern Länge des Mondes auch}{}
\EdL{17}{AK07-J0077-L17}{seine Anomalie auf analoge Weise\,; die hiezu dienende Kolonne ist in unsern Tafeln}{}
\EdL{18}{AK07-J0077-L18}{überschrieben\,: \textit{Argumentum lunae}. Die Bestimmung \Sperr{Kh.}’s ist gegenüber derjenigen}{}
\EdL{19}{AK07-J0077-L19}{\Sperr{Batt.}’s sehr einfach, der erstere entnimmt die Gleichung des Mondes einer einzigen}{}
\EdL{20}{AK07-J0077-L20}{Kolonne, während der letztere deren vier braucht, indem er die Ungleichheiten der}{}
\EdL{21}{AK07-J0077-L21}{Mondbewegung, die sogenannte \textit{Evectio} und \textit{Prosneusis} berücksichtigt.\textsuperscript{2}) Wir wollen}{}
\EdL{22}{AK07-J0077-L22}{nach unsern Regeln und Tafeln den wahren Ort des Mondes für den Mittag des}{}
\EdL{23}{AK07-J0077-L23}{15. Ṣafar des Jahres 540 d. H. bestimmen.}{}
\EdL{24}{AK07-J0077-L24}{\parbox[t]{\dimexpr0.2450\linewidth-0.5em\relax}{\raggedright\small\mbox{}\textit{Beispiel}\,:}\hspace{0.5em}\parbox[t]{\dimexpr0.2450\linewidth-0.5em\relax}{\raggedright\small\mbox{}Die \textit{Anomalie} beträgt am Ende des 510. Jahres}\hspace{0.5em}\parbox[t]{\dimexpr0.2450\linewidth-0.5em\relax}{\raggedright\small\mbox{}=}\hspace{0.5em}\parbox[t]{\dimexpr0.2450\linewidth-0.5em\relax}{\raggedright\small\mbox{}61° 52′ 40″}\hspace{0.5em}}{}
\EdL{25}{AK07-J0077-L25}{\parbox[t]{\dimexpr0.2450\linewidth-0.5em\relax}{\raggedright\small\mbox{}}\hspace{0.5em}\parbox[t]{\dimexpr0.2450\linewidth-0.5em\relax}{\raggedright\small\mbox{}Zunahme für 29 Jahre}\hspace{0.5em}\parbox[t]{\dimexpr0.2450\linewidth-0.5em\relax}{\raggedright\small\mbox{}=}\hspace{0.5em}\parbox[t]{\dimexpr0.2450\linewidth-0.5em\relax}{\raggedright\small\mbox{}348° 56′ 57″}\hspace{0.5em}}{}
\EdL{26}{AK07-J0077-L26}{\parbox[t]{\dimexpr0.2450\linewidth-0.5em\relax}{\raggedright\small\mbox{}}\hspace{0.5em}\parbox[t]{\dimexpr0.2450\linewidth-0.5em\relax}{\raggedright\small\mbox{}„ „ den 1. Monat Muḥarram}\hspace{0.5em}\parbox[t]{\dimexpr0.2450\linewidth-0.5em\relax}{\raggedright\small\mbox{}=}\hspace{0.5em}\parbox[t]{\dimexpr0.2450\linewidth-0.5em\relax}{\raggedright\small\mbox{}31° 56′ 59″}\hspace{0.5em}}{}
\EdL{27}{AK07-J0077-L27}{\parbox[t]{\dimexpr0.2450\linewidth-0.5em\relax}{\raggedright\small\mbox{}}\hspace{0.5em}\parbox[t]{\dimexpr0.2450\linewidth-0.5em\relax}{\raggedright\small\mbox{}„ „ 14 Tage}\hspace{0.5em}\parbox[t]{\dimexpr0.2450\linewidth-0.5em\relax}{\raggedright\small\mbox{}=}\hspace{0.5em}\parbox[t]{\dimexpr0.2450\linewidth-0.5em\relax}{\raggedright\small\mbox{}182° 54′ 36″}\hspace{0.5em}}{}
\EdL{28}{AK07-J0077-L28}{\parbox[t]{\dimexpr0.2450\linewidth-0.5em\relax}{\raggedright\small\mbox{}}\hspace{0.5em}\parbox[t]{\dimexpr0.4900\linewidth-0.5em\relax}{\raggedright\small\mbox{}Summa\,:}\hspace{0.5em}\parbox[t]{\dimexpr0.2450\linewidth-0.5em\relax}{\raggedright\small\mbox{}625° 41′ 12″}\hspace{0.5em}}{}
\EdL{29}{AK07-J0077-L29}{\parbox[t]{\dimexpr0.7350\linewidth-0.5em\relax}{\raggedright\small\mbox{}}\hspace{0.5em}\parbox[t]{\dimexpr0.2450\linewidth-0.5em\relax}{\raggedright\small\mbox{}— 360°}\hspace{0.5em}}{}
\EdL{30}{AK07-J0077-L30}{\parbox[t]{\dimexpr0.2450\linewidth-0.5em\relax}{\raggedright\small\mbox{}}\hspace{0.5em}\parbox[t]{\dimexpr0.2450\linewidth-0.5em\relax}{\raggedright\small\mbox{}\textit{Anomalie}}\hspace{0.5em}\parbox[t]{\dimexpr0.2450\linewidth-0.5em\relax}{\raggedright\small\mbox{}=}\hspace{0.5em}\parbox[t]{\dimexpr0.2450\linewidth-0.5em\relax}{\raggedright\small\mbox{}265° 41′ 12″}\hspace{0.5em}}{}
\EdLp{31}{AK07-J0077-L31}{Dafür findet man in der Tafel 24 die \textit{Gleichung} 4° 55′ 7″\textsuperscript{3}), diese wird zu der}{}
\EdL{32}{AK07-J0077-L32}{mittlern Länge addiert, die wir oben = 314° 6′ 28″ gefunden haben, weil die Anomalie}{}
\EdL{33}{AK07-J0077-L33}{grösser als 180° ist, dies gibt für den \textit{wahren Ort} des Mondes 319° 1′ 35″, d. h. der}{}
\EdL{34}{AK07-J0077-L34}{Mond befand sich zu der gegebenen Zeit in 19° 1′ 35″ des Wassermanns.}{}
\EdLp{35}{AK07-J0077-L35}{Nach den Tafeln \Sperr{Batt.}’s findet man 330° 41′, dieselben geben nämlich für die}{}
\NoteRule
\EdN{N1}{AK07-J0077-N01}{\textsuperscript{1}) l. c. p. 220.}
\EdN{N2}{AK07-J0077-N02}{\textsuperscript{2}) Man vergleiche für die Bestimmung \Sperr{Batt.}’s die Regeln des Textes bei \Sperr{Nallino} I, 50—54}
\EdN{N3}{AK07-J0077-N03}{u. 75, und die Erklärungen dazu in II, 223—226.}
\EdN{N4}{AK07-J0077-N04}{\textsuperscript{3}) Für 265° Grad findet man 4° 54′ 52″ und für die 41′ ergeben sich noch 15″.}
\FrameHead{0078}{p. 46}{0078}{p. 46}
\EdL{1}{AK07-J0078-L01}{mittlere Länge am Ende des 510. Jahres 50° 44′ 20″ an, statt wie unsere Tafeln 40° 56′ 38″,}{}
\EdL{2}{AK07-J0078-L02}{und die Gleichung wird 6° 19′ 57″ gefunden statt 4° 55′ 7″.}{}
\EdLp{3}{AK07-J0078-L03}{Wie die Mondtafeln \Sperr{Kh.}’s berechnet sind, ist mir nicht klar geworden\,; die}{}
\EdL{4}{AK07-J0078-L04}{grösste Gleichung beträgt bei 90° mittlerer Anomalie 4° 56′, dies entspricht einem}{}
\EdL{5}{AK07-J0078-L05}{Epizykelradius von c. 5\textsuperscript{p} 10′, anstatt 5\textsuperscript{p} 15′, wie \Sperr{Batt.} hat\,; berechnet man aber die}{}
\EdL{6}{AK07-J0078-L06}{Gleichung für andere Anomalien (\textit{m}) nach der Formel}{}
\EdLc{7}{AK07-J0078-L07}{tg aeq. anom. = \Frac{\textit{e} sin \textit{m}}{1 + \textit{e} cos \textit{m}},}{}
\EdL{8}{AK07-J0078-L08}{worin \textit{e} also = \Frac{5\textsuperscript{p} 10′}{60\textsuperscript{p}} = 0,08611 zu setzen ist, so findet man keine Übereinstimmung,}{}
\EdL{9}{AK07-J0078-L09}{es müssen mithin eine der oben genannten, oder beide Ungleichheiten doch wohl}{}
\EdL{10}{AK07-J0078-L10}{berücksichtigt worden sein, in welcher Weise aber konnten wir nicht herausfinden.}{}
\EdLp{11}{AK07-J0078-L11}{In den Tafel- und Kolonnenüberschriften finden wir für das arabische \textit{Taʿdīl}}{}
\EdL{12}{AK07-J0078-L12}{= \textit{Gleichung} bald das arabische Wort in der Form \textit{Tadil}, bald das lateinische \textit{Exa-}}{}
\EdL{13}{AK07-J0078-L13}{\textit{minatio}, \textit{Aequatio} kommt nicht vor, statt dieses Wortes hat N \textit{Coaequatio}. Das Wort}{}
\EdL{14}{AK07-J0078-L14}{\textit{elheza}, auch \textit{elheca} geschrieben, ist das arabische \textit{al-khāṣṣa} (auch \textit{al-ḥāṣṣa} geschrieben)}{}
\EdL{15}{AK07-J0078-L15}{= die (dem Planeten) eigentümliche (Bewegung), richtig wäre\,: \textit{al-ḥaraka al-khāṣṣa}}{}
\EdL{16}{AK07-J0078-L16}{= die (ihm) eigentümliche Bewegung, d. i. die sogenannte \textit{Anomalie}, d. h. der Bogen,}{}
\EdL{17}{AK07-J0078-L17}{den der Mond oder der Planet auf dem Epizykel vom Apogeum desselben aus bis}{}
\EdL{18}{AK07-J0078-L18}{zu einem bestimmten Zeitpunkt durchlaufen hat. Dieses Wort wird in unserm}{}
\EdL{19}{AK07-J0078-L19}{Werke, besonders in Ms. C, mit \textit{Argumentum} übersetzt\,; es wird verwechselt mit \textit{al-ḥiṣṣa}}{}
\EdL{20}{AK07-J0078-L20}{= Anteil, Loos, d. h. das Argument, mit dem man in eine Tafel eintritt.}{}
\EdLc{21}{AK07-J0078-L21}{\textbf{Kap. 10. Bestimmung des (wahren) Ortes der drei obern Planeten.}}{}
\EdLc{22}{AK07-J0078-L22}{(Text p. 10—11, Tafeln 27—44.)}{}
\EdLp{23}{AK07-J0078-L23}{Betrachten wir zuerst die \textit{Apogeen} der Planeten und der Sonne, so finden wir}{}
\EdL{24}{AK07-J0078-L24}{für dieselben im Text p. 14—15 und in den Tafeln 27—44 folgende Werte\,:}{}
\EdL{25}{AK07-J0078-L25}{\parbox[t]{\dimexpr0.1633\linewidth-0.5em\relax}{\raggedright\small\mbox{}Saturn}\hspace{0.5em}\parbox[t]{\dimexpr0.1633\linewidth-0.5em\relax}{\raggedright\small\mbox{}Jupiter}\hspace{0.5em}\parbox[t]{\dimexpr0.1633\linewidth-0.5em\relax}{\raggedright\small\mbox{}Mars}\hspace{0.5em}\parbox[t]{\dimexpr0.1633\linewidth-0.5em\relax}{\raggedright\small\mbox{}Venus}\hspace{0.5em}\parbox[t]{\dimexpr0.1633\linewidth-0.5em\relax}{\raggedright\small\mbox{}Merkur}\hspace{0.5em}\parbox[t]{\dimexpr0.1633\linewidth-0.5em\relax}{\raggedright\small\mbox{}Sonne}\hspace{0.5em}}{}
\EdL{26}{AK07-J0078-L26}{\parbox[t]{\dimexpr0.1633\linewidth-0.5em\relax}{\raggedright\small\mbox{}244° 55′}\hspace{0.5em}\parbox[t]{\dimexpr0.1633\linewidth-0.5em\relax}{\raggedright\small\mbox{}172° 32′}\hspace{0.5em}\parbox[t]{\dimexpr0.1633\linewidth-0.5em\relax}{\raggedright\small\mbox{}128° 24′}\hspace{0.5em}\parbox[t]{\dimexpr0.1633\linewidth-0.5em\relax}{\raggedright\small\mbox{}81° 15′}\hspace{0.5em}\parbox[t]{\dimexpr0.1633\linewidth-0.5em\relax}{\raggedright\small\mbox{}224° 54′}\hspace{0.5em}\parbox[t]{\dimexpr0.1633\linewidth-0.5em\relax}{\raggedright\small\mbox{}77° 55′.}\hspace{0.5em}}{}
\EdLp{27}{AK07-J0078-L27}{Diese Zahlen stimmen, auch wenn die \textit{Praezession} berücksichtigt wird, weder}{}
\EdL{28}{AK07-J0078-L28}{mit denen des \Sperr{Ptol.} noch mit denen des \Sperr{Batt.}, auch nicht mit den Werten im}{}
\EdL{29}{AK07-J0078-L29}{\textit{Sūrya-Siddhānta} und in der \textit{Panchasiddhāntikā} überein.\textsuperscript{1}) Besonders weichen die}{}
\EdL{30}{AK07-J0078-L30}{Werte für Jupiter, Merkur und Venus in den genannten Werken wesentlich von}{}
\EdL{31}{AK07-J0078-L31}{einander ab\,; der Wert für das Apogeum der Sonne stimmt, wenn wir für 55′ 15′}{}
\EdL{32}{AK07-J0078-L32}{lesen, wie wir oben (p. 43) schon bemerkt haben, mit dem Werte, den der arabische}{}
\EdL{33}{AK07-J0078-L33}{\textit{Sindhind} nach dem Zeugnis \Sperr{Ibn Hibintās} gegeben hat\,; ebenso der Wert für}{}
\EdL{34}{AK07-J0078-L34}{Jupiter mit demjenigen der „erprobten Tafeln“ nach \Sperr{Ibn Yūnus}\textsuperscript{2}). Am stärksten}{}
\EdL{35}{AK07-J0078-L35}{weicht der Wert für Merkur von denen der übrigen Tafeln ab, in denen er nicht}{}
\EdL{36}{AK07-J0078-L36}{über 202° hinaus geht, doch ist davon der \textit{Sūrya-Siddhānta} auszunehmen, in welchem}{}
\NoteRule
\EdN{N1}{AK07-J0078-N01}{\textsuperscript{1}) Vergl. \Sperr{Burgess}, l. c. p. 216 u. 427\,; ebenso \textit{the Panchasiddhāntikā of Varāha Mihira}, ed.}
\EdN{N2}{AK07-J0078-N02}{by G. \Sperr{Thibaut}, etc., Benares 1889, p. 93 der engl. Übersetzung.}
\EdN{N3}{AK07-J0078-N03}{\textsuperscript{2}) \Sperr{Batt.}, ed. \Sperr{Nallino}, I, 218 u. 239—242.}
\FrameHead{0079}{p. 47}{0079}{p. 47}
\EdL{1}{AK07-J0079-L01}{das Apogeum des Merkur zu 220° 28′ angegeben wird, immerhin noch c. 4\textsuperscript{1}/\textsubscript{2}° geringer}{}
\EdL{2}{AK07-J0079-L02}{als bei \Sperr{Kh.} Die Übereinstimmung des Jupiter-Apogeums bei \Sperr{Kh.} und den „erprobten}{}
\EdL{3}{AK07-J0079-L03}{Tafeln“ bestärkt uns in der oben (p. 32) schon ausgesprochenen Ansicht, dass gewisse}{}
\EdL{4}{AK07-J0079-L04}{Angaben der Tafeln \Sperr{Kh.}’s den „erprobten Tafeln“ der Astronomen \Sperr{Al-Ma’mūns}}{}
\EdL{5}{AK07-J0079-L05}{entnommen sein könnten. Allerdings weicht \Sperr{Kh.} wieder (wenn man sich wenigstens}{}
\EdL{6}{AK07-J0079-L06}{auf die Berichte des \Sperr{Ibn Yūnus} über die „erprobten Tafeln“ verlassen kann) darin}{}
\EdL{7}{AK07-J0079-L07}{von diesen Tafeln ab, dass er die Länge des Apogeums der Venus nicht gleich der}{}
\EdL{8}{AK07-J0079-L08}{des Apogeums der Sonne annimmt.}{}
\EdLp{9}{AK07-J0079-L09}{Ich gebe im folgenden zum bessern Verständnis dieser Fragen eine kurze}{}
\EdL{10}{AK07-J0079-L10}{Darstellung der Bewegung eines obern Planeten nach \Sperr{Ptol.} und \Sperr{Batt.}, ich schliesse}{}
\EdL{11}{AK07-J0079-L11}{mich möglichst nahe an \Sperr{Nallino} an.\textsuperscript{1})}{}
\EdLp{12}{AK07-J0079-L12}{Es wird angenommen, der Planet bewege sich auf einem \textit{Epizykel}, dessen}{}
\EdL{13}{AK07-J0079-L13}{Centrum\textsuperscript{2}) \textit{H} sich auf dem in Bezug auf die Erde \textit{E} exzentrischen Kreise mit dem}{}
\EdL{14}{AK07-J0079-L14}{Mittelpunkt \textit{M}, dem sogenannten \textit{Deferens}, mit}{}
\EdL{15}{AK07-J0079-L15}{ungleicher Winkelgeschwindigkeit bewege, und}{}
\EdL{16}{AK07-J0079-L16}{zwar sowohl in Bezug auf den Punkt \textit{M} als}{}
\EdL{17}{AK07-J0079-L17}{auf den Punkt \textit{E}. Das Centrum \textit{H} bewege sich}{}
\EdL{18}{AK07-J0079-L18}{aber gleichförmig auf einem dritten fingierten}{}
\EdL{19}{AK07-J0079-L19}{Kreis (ich lasse diesen in der Zeichnung weg)}{}
\EdL{20}{AK07-J0079-L20}{mit dem Centrum \textit{E′} (\textit{E′M} = \textit{ME}), dem so-}{}
\EdL{21}{AK07-J0079-L21}{genannten \textit{Aequans}, oder besser, die Bewegung}{}
\EdL{22}{AK07-J0079-L22}{des Punktes \textit{H} erscheine von diesem Punkt \textit{E′}}{}
\EdL{23}{AK07-J0079-L23}{aus gleichförmig. Sowohl der Planet auf dem}{}
\EdL{24}{AK07-J0079-L24}{Epizykel, als auch das Centrum des Epizykels}{}
\EdL{25}{AK07-J0079-L25}{auf dem Deferens bewegen sich von Westen}{}
\EdL{26}{AK07-J0079-L26}{nach Osten. Die obern Planeten vollenden}{}
\EdL{27}{AK07-J0079-L27}{ihren Umkreis im \textit{Epizykel} in einem siderischen}{}
\EdL{28}{AK07-J0079-L28}{Jahr, die Centren ihrer Epizykeln ihren Um-}{}
\EdL{29}{AK07-J0079-L29}{kreis im \textit{Deferens} in verschiedenen Zeiten. \textit{EM}}{}
\EdL{30}{AK07-J0079-L30}{ist die Exzentrizität der Planetenbahn, \textit{EE′} die doppelte. Das \textit{wahre Apogeum} des}{}
\EdL{31}{AK07-J0079-L31}{Epizykels ist in jeder Lage des Punktes \textit{H} der Punkt, in welchem die Verlängerung}{}
\EdL{32}{AK07-J0079-L32}{von \textit{EH} die Ekliptik trifft (in unserer Fig. also \textit{R}), das \textit{mittlere Apogeum} der Punkt,}{}
\EdL{33}{AK07-J0079-L33}{in welchem die Verlängerung von \textit{E′H} die Ekliptik trifft (in unserer Fig. also \textit{Q}).}{}
\EdL{34}{AK07-J0079-L34}{Beide Punkte fallen zusammen, wenn der Mittelpunkt des Epizykels im Apogeum \textit{G}}{}
\EdL{35}{AK07-J0079-L35}{oder im Perigeum \textit{K} des Deferens sich befindet. \textit{AFB} sei die Ekliptik, das Apogeum \textit{G}}{}
\EdL{36}{AK07-J0079-L36}{habe in der Ekliptik die Länge \textit{FA}, das Centrum \textit{H} des Epizykels hat nach seiner}{}
\EdL{37}{AK07-J0079-L37}{mittlern Bewegung um das Centrum \textit{E′} des Aequans die Länge \textit{FQ}, nach seiner}{}
\EdL{38}{AK07-J0079-L38}{wahren Bewegung aber aus dem Centrum der Erde \textit{E} gesehen die Länge \textit{FR}\,; die}{}
\EdL{39}{AK07-J0079-L39}{Differenz zwischen beiden, der mittlern und der wahren Länge, d. h. der Bogen \textit{RQ},}{}
\NoteRule
\EdN{N1}{AK07-J0079-N01}{\textsuperscript{1}) Ibid. II, 238—241.}
\EdN{N2}{AK07-J0079-N02}{\textsuperscript{2}) Da in der Folge oft deutsche mit lateinischen Bezeichnungen abwechseln, schreibe ich der}
\EdN{N3}{AK07-J0079-N03}{Einheitlichkeit wegen stets „Centrum“ statt „Zentrum“.}
\FrameHead{0080}{p. 48}{0080}{p. 48}
\EdL{1}{AK07-J0080-L01}{wird von \Sperr{Batt.} die \textit{Gleichung der Anomalie und des Centrums} oder auch nur \textit{Gleichung}}{}
\EdL{2}{AK07-J0080-L02}{\textit{des Centrums} genannt, sie ist auch die Differenz zwischen dem \textit{mittlern} und \textit{wahren}}{}
\EdL{3}{AK07-J0080-L03}{\textit{Apogeum} des Epizykels. — Diese Gleichung wird durch den \ensuremath{\sphericalangle} \textit{E′HE} gemessen\,; zieht}{}
\EdL{4}{AK07-J0080-L04}{man \textit{MS} und \textit{EN} senkrecht auf die Verlängerung von \textit{QHE′}, so hat man aus dem}{}
\EdL{5}{AK07-J0080-L05}{rechtwinkligen Dreieck \textit{ENH}\,:}{}
\EdLc{6}{AK07-J0080-L06}{sin \textit{E′HE} = \Frac{\textit{EN}}{\textit{EH}}, oder tg \textit{E′HE} = \Frac{\textit{EN}}{\textit{HN}}.}{}
\EdLp{7}{AK07-J0080-L07}{Es sei nun \textit{m} = Bogen \textit{AQ}, d. h. die Differenz zwischen der mittlern Länge des}{}
\EdL{8}{AK07-J0080-L08}{Planeten und derjenigen des Apogeums des exzentrischen Kreises, welcher Bogen}{}
\EdL{9}{AK07-J0080-L09}{durch den \ensuremath{\sphericalangle} \textit{AE′Q} gemessen wird, und von den Arabern das \textit{Centrum} (\textit{al-markaz})}{}
\EdL{10}{AK07-J0080-L10}{genannt wurde, weil er aus dem Centrum des \textit{Aequans} gemessen gedacht wird\,; ferner}{}
\EdL{11}{AK07-J0080-L11}{sei \textit{e} die Exzentrizität \textit{EM} = \textit{E′M}, \textit{ρ} = \textit{HM} = Radius des Deferens, so findet man}{}
\EdL{12}{AK07-J0080-L12}{leicht\,:}{}
\EdLc{13}{AK07-J0080-L13}{tg aeq. centri = tg \textit{E′HE} = \Frac{2\textit{e} sin \textit{m}}{\Root{\textit{ρ}\textsuperscript{2} − (\Frac{\textit{e} sin \textit{m}}{\textit{R}})\textsuperscript{2}} + \Frac{\textit{e} cos \textit{m}}{\textit{R}}}\textsuperscript{1}),}{}
\EdL{14}{AK07-J0080-L14}{oder \textit{R} = 1, und \Frac{\textit{e}}{\textit{ρ}} = \textit{e}\textsubscript{1} gesetzt\,:}{}
\EdLc{15}{AK07-J0080-L15}{tg \textit{E′HE} = \Frac{2\textit{e}\textsubscript{1} sin \textit{m}}{\Root{1 − (\textit{e}\textsubscript{1} sin \textit{m})\textsuperscript{2}} + \textit{e}\textsubscript{1} cos \textit{m}}\textsuperscript{2}).}{}
\EdLp{16}{AK07-J0080-L16}{Die \textit{Gleichung der Anomalie und des Centrums} mit Hilfe dieser Formeln bestimmt}{}
\EdL{17}{AK07-J0080-L17}{ergibt die Werte, die bei \Sperr{Batt.} in den Tafeln der Gleichungen der Planeten (ed.}{}
\EdL{18}{AK07-J0080-L18}{\Sperr{Nallino} II, 108—137) in Kolonne III stehen. Es muss aber an der Länge der Pla-}{}
\EdL{19}{AK07-J0080-L19}{neten noch eine andere Korrektur vorgenommen werden, die abhängig ist von der}{}
\EdL{20}{AK07-J0080-L20}{Stellung des Planeten im Epizykel. Der Planet sei in \textit{P} (s. folg. Fig.), es wird der}{}
\EdL{21}{AK07-J0080-L21}{Bogen \textit{RT}, oder der \ensuremath{\sphericalangle} \textit{RET} gesucht\,; zieht man \textit{PS} senkrecht auf \textit{ER}, so ist\,:}{}
\EdLc{22}{AK07-J0080-L22}{sin \textit{RET} = \Frac{\textit{PS}}{\textit{EP}} = \Frac{\textit{PS}}{Geozentr. Dist. d. Plan.},}{}
\EdL{23}{AK07-J0080-L23}{aber \textit{PS} = \textit{PH} · sin \textit{RHP} = \textit{r} sin \textit{ψ}, wo \textit{r} der Radius des Epizykels, in Dezimalteilen}{}
\EdL{24}{AK07-J0080-L24}{des Radius des Deferens oder des Aequans (\textit{MH} = 1) ausgedrückt ist, und \textit{ψ} den Winkel}{}
\EdL{25}{AK07-J0080-L25}{(Bogen) bedeutet, den der Planet vom Apogeum des Epizykels an bis zum Ort \textit{P}}{}
\EdL{26}{AK07-J0080-L26}{beschrieben hat, d. h. die \textit{Anomalie}. Die Distanz \textit{EP} des Planeten von der Erde ist}{}
\EdL{27}{AK07-J0080-L27}{aber = \Root{(\textit{EH} + \textit{HS})\textsuperscript{2} + \textit{PS}\textsuperscript{2}} = \Root{(Dist. centri epic. + \textit{r} cos \textit{ψ})\textsuperscript{2} + (\textit{r} sin \textit{ψ})\textsuperscript{2}}, also}{}
\EdLc{28}{AK07-J0080-L28}{sin \textit{RET} = \Frac{\textit{r} sin \textit{ψ}}{\Root{(Dist. centri epic. + \textit{r} cos \textit{ψ})\textsuperscript{2} + (\textit{r} sin \textit{ψ})\textsuperscript{2}}}.}{}
\EdL{29}{AK07-J0080-L29}{Dieser Winkel wird die \textit{Gleichung der Anomalie} genannt, er ist der Unterschied zwischen}{}
\EdL{30}{AK07-J0080-L30}{der wahren Länge des Centrums \textit{H} des Epizykels und der wahren Länge des Planeten}{}
\EdL{31}{AK07-J0080-L31}{in der Ekliptik. Durch die Tangens ausgedrückt, findet man einfacher:}{}
\NoteRule
\EdN{N1}{AK07-J0080-N01}{\textsuperscript{1}) \textit{R} bedeutet hier 60\textsuperscript{p}, da die trigonom. Funktionen der Araber sich immer auf einen Kreis}
\EdN{N2}{AK07-J0080-N02}{mit dem Radius 60 beziehen\,; in der zweiten Formel ist der Radius des Kreises = 1 statt = 60}
\EdN{N3}{AK07-J0080-N03}{angenommen.}
\EdN{N4}{AK07-J0080-N04}{\textsuperscript{2}) Wollte man den \ensuremath{\sphericalangle} \textit{E′HE} mit Hilfe des Sinus bestimmen, was sehr wahrscheinlich \Sperr{Batt.}}
\EdN{N5}{AK07-J0080-N05}{getan hat, so würde die Formel wesentlich komplizierter (vergl. \Sperr{Batt.} ed. \Sperr{Nallino}, II, 240).}
\FrameHead{0081}{p. 49}{0081}{p. 49}
\EdLc{1}{AK07-J0081-L01}{tg \textit{RET} = \Frac{\textit{PS}}{\textit{ES}} = \Frac{\textit{PS}}{\textit{EH} + \textit{HS}} = \Frac{\textit{r} sin \textit{ψ}}{Dist. centri epic. + \textit{r} cos \textit{ψ}}.}{}
\EdLp{2}{AK07-J0081-L02}{Diese \textit{Gleichung der Anomalie} ändert sich, wie man sieht, mit der variabeln}{}
\EdL{3}{AK07-J0081-L03}{Distanz des Centrums des Epizykels von der Erde\,; \Sperr{Batt.} hat in Kol. VI die Gleichung}{}
\EdL{4}{AK07-J0081-L04}{der Anomalie für eine mittlere Distanz}{}
\EdL{5}{AK07-J0081-L05}{gegeben, Kol. V u. VII dienen dazu, die Kor-}{}
\EdL{6}{AK07-J0081-L06}{rekturen für die Abweichung der übrigen}{}
\EdL{7}{AK07-J0081-L07}{Distanzen von der mittlern zu bestimmen\,;}{}
\EdL{8}{AK07-J0081-L08}{auf Einzelheiten können wir hier nicht}{}
\EdL{9}{AK07-J0081-L09}{eintreten (s. das Beispiel unten).}{}
\EdLp{10}{AK07-J0081-L10}{Was nun die Tafeln \Sperr{Kh.}’s anbetrifft,}{}
\EdL{11}{AK07-J0081-L11}{so lässt er in erster Linie die Kol. V u. VII}{}
\EdL{12}{AK07-J0081-L12}{bei \Sperr{Batt.}, durch welche die von der geo-}{}
\EdL{13}{AK07-J0081-L13}{zentrischen Distanz des Epizykelcentrums}{}
\EdL{14}{AK07-J0081-L14}{abhängige Gleichung der Anomalie korri-}{}
\EdL{15}{AK07-J0081-L15}{giert wird, weg. Dagegen hat \Sperr{Kh.} eine}{}
\EdL{16}{AK07-J0081-L16}{Kolonne, überschrieben mit \textit{Sublimatio}}{}
\EdL{17}{AK07-J0081-L17}{\textit{Saturni examinata}, oder besser \textit{Apogeum}}{}
\EdL{18}{AK07-J0081-L18}{\textit{Saturni aequatum}, deren Werte zwischen}{}
\EdL{19}{AK07-J0081-L19}{8\textsuperscript{s} 4° 55′ (= 244° 55′ = wirkliches Apogeum}{}
\EdL{20}{AK07-J0081-L20}{des Saturns) und 8\textsuperscript{s} 2° 3′ (= 242° 3′) variieren.}{}
\EdL{21}{AK07-J0081-L21}{Dies ist nun wohl der Punkt, auf den sich die Angaben bei einigen arabischen}{}
\EdL{22}{AK07-J0081-L22}{Autoren über die Tafeln \Sperr{Kh.}’s beziehen. Ich habe oben (p. 33) schon angeführt, dass}{}
\EdL{23}{AK07-J0081-L23}{\Sperr{Ibn al-Adamī} berichtet hat, dass \Sperr{Kh.} seine Gleichungen (der Planeten) nach der}{}
\EdL{24}{AK07-J0081-L24}{Methode der Perser (also wahrscheinlich nach den Tafeln des \textit{Shāh}) berechnet habe\,;}{}
\EdL{25}{AK07-J0081-L25}{ich muss nun hier noch hinzufügen, dass \Sperr{Abū Naṣr Manṣūr b. ʿAlī b. ʿIrāḳ}}{}
\EdL{26}{AK07-J0081-L26}{(c. 950—1000)\textsuperscript{1}) eine Abhandlung an \Sperr{Al-Bīrūnī} gerichtet hat „über die Gründe der}{}
\EdL{27}{AK07-J0081-L27}{\textit{Halbierung} deı \textit{Gleichung} bei den Astronomen die dem Sindhind folgten“.\textsuperscript{2}) Wie}{}
\EdL{28}{AK07-J0081-L28}{verhält es sich nun mit dieser \textit{Halbierung} der \textit{Gleichung}\,? Schon M. \Sperr{Chasles} hat}{}
\EdL{29}{AK07-J0081-L29}{in der oben (p. I der Einleitung) zitierten Abhandlung diese Frage berührt, aber}{}
\EdL{30}{AK07-J0081-L30}{seine Darstellung ist in gewissen Punkten nicht ganz klar und auch nicht ganz}{}
\EdL{31}{AK07-J0081-L31}{richtig. Er sagt p. 849\,: \textit{Former un apogée fictif, c’est comme si l’on donnait une équa-}}{}
\EdL{32}{AK07-J0081-L32}{\textit{tion fictive à l’apogée réel. Et en effet, j’ai trouvé depuis cette équation ou la manière}}{}
\EdL{33}{AK07-J0081-L33}{\textit{de la former dans d’autres tables que j’ai reconnues être imitées du Sindhind et dans}}{}
\EdL{34}{AK07-J0081-L34}{\textit{lesquelles ces traces manifestes d’une astronomie étrangère avaient échappé à l’attention}}{}
\EdL{35}{AK07-J0081-L35}{\textit{des astronomes historiens. L’équation en question est précisement la moitié de l’équation}}{}
\EdL{36}{AK07-J0081-L36}{\textit{du centre, c’est ainsi que l’auteur la désigne. C’est la prosthaphérèse de longitude de}}{}
\EdL{37}{AK07-J0081-L37}{\textit{Ptolemée.} — Es ist nun wohl richtig, dass \Sperr{Kh.} auch das \textit{Apogeum} korrigiert, und nicht}{}
\EdL{38}{AK07-J0081-L38}{nur die \textit{Anomalie} und das \textit{Centrum} des Planeten\,; es ist auch möglich, dass diese}{}
\EdL{39}{AK07-J0081-L39}{Korrektur des Apogeums ein Ersatz für den von den Indern, bezw. Persern fallen-}{}
\NoteRule
\EdN{N1}{AK07-J0081-N01}{\textsuperscript{1}) Vergl. \Sperr{Suter}, in \textit{Abhandlgn. zur Gesch. der mathem. Wissensch.} X, 81.}
\EdN{N2}{AK07-J0081-N02}{\textsuperscript{2}) Vergl. \Sperr{Albērūnī}, \textit{Chronologie orientalischer Völker}, ed. \Sperr{Sachau}, Leipzig, 1878, p. XXXXVII.}
\FrameHead{0082}{p. 50}{0082}{p. 50}
\EdL{1}{AK07-J0082-L01}{gelassenen \textit{Aequans} des \Sperr{Ptol.} sein soll\,; \Sperr{Chasles} irrt sich aber, oder verschreibt}{}
\EdL{2}{AK07-J0082-L02}{sich, wenn er sagt\,: \textit{L’équation en question est précisément la moitié de l’équation du}}{}
\EdL{3}{AK07-J0082-L03}{\textit{centre}\,; wir haben gefunden, dass die in Frage kommende Gleichung (des Apogeums)}{}
\EdL{4}{AK07-J0082-L04}{bei \Sperr{Kh.} die \textit{Hälfte der Gleichung der Anomalie} ist. Die Sache verhält sich so\,:}{}
\EdLp{5}{AK07-J0082-L05}{Die Gleichung der Anomalie (in unsern Tafeln \textit{Examinatio argumenti}) wird}{}
\EdL{6}{AK07-J0082-L06}{halbiert\,; die eine Hälfte wird als Gleichung (Korrektur) des \textit{Apogeums} von diesem}{}
\EdL{7}{AK07-J0082-L07}{abgezogen, und ergibt so das korrigierte Apogeum, das in der zweiten Hauptkolonne}{}
\EdL{8}{AK07-J0082-L08}{(vergl. z. B. Tafel 27 für Saturn), überschrieben \textit{Sublimatio examinata}, verzeichnet ist\,; die}{}
\EdL{9}{AK07-J0082-L09}{andere Hälfte wird wahrscheinlich zur Korrektur der \textit{Gleichung} des \textit{Centrums} verwendet,}{}
\EdL{10}{AK07-J0082-L10}{in welcher Weise aber, haben wir nicht herausgefunden. Wir haben verschiedene}{}
\EdL{11}{AK07-J0082-L11}{Proben gemacht, und dabei diese Hälfte das eine mal gleich, das andere mal viel}{}
\EdL{12}{AK07-J0082-L12}{grösser, ein drittes mal bedeutend kleiner wie die Gleichung des Centrums gefunden.}{}
\EdL{13}{AK07-J0082-L13}{Ich gebe hier zwei von diesen Proben, indem ich die Zahlen bei \Sperr{Kh.} mit den ent-}{}
\EdL{14}{AK07-J0082-L14}{sprechenden bei \Sperr{Batt.} vergleiche, die von denen des \Sperr{Ptol.} wenig abweichen.}{}
\EdLp{15}{AK07-J0082-L15}{Wir nehmen zur Vergleichung\,: 1) die Daten beider Tafeln, die zum Argument}{}
\EdL{16}{AK07-J0082-L16}{70° des Saturns (bei \Sperr{Kh.} Tafel 29, bei \Sperr{Batt.} II, 110) gehören\,:}{}
\EdL{17}{AK07-J0082-L17}{\parbox[t]{\dimexpr0.4900\linewidth-0.5em\relax}{\raggedright\small\mbox{}\Sperr{Batt.}}\hspace{0.5em}\parbox[t]{\dimexpr0.2450\linewidth-0.5em\relax}{\raggedright\small\mbox{}Aequat. (anomal. et) centri}\hspace{0.5em}\parbox[t]{\dimexpr0.2450\linewidth-0.5em\relax}{\raggedright\small\mbox{}Aequat. anomaliae}\hspace{0.5em}}{}
\EdL{18}{AK07-J0082-L18}{\parbox[t]{\dimexpr0.4900\linewidth-0.5em\relax}{\raggedright\small\mbox{}}\hspace{0.5em}\parbox[t]{\dimexpr0.2450\linewidth-0.5em\relax}{\raggedright\small\mbox{}6° 2′}\hspace{0.5em}\parbox[t]{\dimexpr0.2450\linewidth-0.5em\relax}{\raggedright\small\mbox{}5° 37′}\hspace{0.5em}}{}
\EdL{19}{AK07-J0082-L19}{\parbox[t]{\dimexpr0.2450\linewidth-0.5em\relax}{\raggedright\small\mbox{}\Sperr{Kh.}}\hspace{0.5em}\parbox[t]{\dimexpr0.2450\linewidth-0.5em\relax}{\raggedright\small\mbox{}Sublim. examin.}\hspace{0.5em}\parbox[t]{\dimexpr0.2450\linewidth-0.5em\relax}{\raggedright\small\mbox{}Examin. centri}\hspace{0.5em}\parbox[t]{\dimexpr0.2450\linewidth-0.5em\relax}{\raggedright\small\mbox{}Examin. argumenti (anomaliae)}\hspace{0.5em}}{}
\EdL{20}{AK07-J0082-L20}{\parbox[t]{\dimexpr0.2450\linewidth-0.5em\relax}{\raggedright\small\mbox{}}\hspace{0.5em}\parbox[t]{\dimexpr0.2450\linewidth-0.5em\relax}{\raggedright\small\mbox{}242° 20′}\hspace{0.5em}\parbox[t]{\dimexpr0.2450\linewidth-0.5em\relax}{\raggedright\small\mbox{}8° 1′}\hspace{0.5em}\parbox[t]{\dimexpr0.2450\linewidth-0.5em\relax}{\raggedright\small\mbox{}5° 10′}\hspace{0.5em}}{}
\EdL{21}{AK07-J0082-L21}{Halbieren wir die \textit{Gleichung der Anomalie} bei \Sperr{Kh.}, also 5° 10′, so erhalten wir 2° 35′,}{}
\EdL{22}{AK07-J0082-L22}{diese von dem wirklichen Apogeum des Saturns, das p. 14 des Textes zu 244° 55′}{}
\EdL{23}{AK07-J0082-L23}{angegeben ist, subtrahiert, gibt die \textit{Sublimatio examinata} = 242° 20′\,; diese Hälfte zur}{}
\EdL{24}{AK07-J0082-L24}{Gleichung des Centrums bei \Sperr{Batt.} = 6° 2′ addiert, gibt 8° 37′, was von der Gleichung}{}
\EdL{25}{AK07-J0082-L25}{des Centrums bei \Sperr{Kh.}, nämlich 8° 1′, um 36′ abweicht.}{}
\EdLp{26}{AK07-J0082-L26}{2) die Daten beider Tafeln, die zum Argument 70° des Jupiters (bei \Sperr{Kh.} Tafel 35,}{}
\EdL{27}{AK07-J0082-L27}{bei \Sperr{Batt.} II, 116) gehören\,:}{}
\EdL{28}{AK07-J0082-L28}{\parbox[t]{\dimexpr0.4900\linewidth-0.5em\relax}{\raggedright\small\mbox{}\Sperr{Batt.}}\hspace{0.5em}\parbox[t]{\dimexpr0.2450\linewidth-0.5em\relax}{\raggedright\small\mbox{}Aequat. (anomal. et) centri}\hspace{0.5em}\parbox[t]{\dimexpr0.2450\linewidth-0.5em\relax}{\raggedright\small\mbox{}Aequat. anomaliae}\hspace{0.5em}}{}
\EdL{29}{AK07-J0082-L29}{\parbox[t]{\dimexpr0.4900\linewidth-0.5em\relax}{\raggedright\small\mbox{}}\hspace{0.5em}\parbox[t]{\dimexpr0.2450\linewidth-0.5em\relax}{\raggedright\small\mbox{}4° 52′}\hspace{0.5em}\parbox[t]{\dimexpr0.2450\linewidth-0.5em\relax}{\raggedright\small\mbox{}9° 36′}\hspace{0.5em}}{}
\EdL{30}{AK07-J0082-L30}{\parbox[t]{\dimexpr0.2450\linewidth-0.5em\relax}{\raggedright\small\mbox{}\Sperr{Kh.}}\hspace{0.5em}\parbox[t]{\dimexpr0.2450\linewidth-0.5em\relax}{\raggedright\small\mbox{}Sublim. exam.}\hspace{0.5em}\parbox[t]{\dimexpr0.2450\linewidth-0.5em\relax}{\raggedright\small\mbox{}Examin. centri}\hspace{0.5em}\parbox[t]{\dimexpr0.2450\linewidth-0.5em\relax}{\raggedright\small\mbox{}Examin. argumenti (anomaliae)}\hspace{0.5em}}{}
\EdL{31}{AK07-J0082-L31}{\parbox[t]{\dimexpr0.2450\linewidth-0.5em\relax}{\raggedright\small\mbox{}}\hspace{0.5em}\parbox[t]{\dimexpr0.2450\linewidth-0.5em\relax}{\raggedright\small\mbox{}167° 50′}\hspace{0.5em}\parbox[t]{\dimexpr0.2450\linewidth-0.5em\relax}{\raggedright\small\mbox{}4° 44′}\hspace{0.5em}\parbox[t]{\dimexpr0.2450\linewidth-0.5em\relax}{\raggedright\small\mbox{}9° 25′}\hspace{0.5em}}{}
\EdL{32}{AK07-J0082-L32}{Halbiert man die \textit{Gleichung der Anomalie} bei \Sperr{Kh.}, also 9° 25′, so erhält man 4° 42′,}{}
\EdL{33}{AK07-J0082-L33}{diese von dem wirklichen Apogeum des Jupiters, das p. 14 des Textes zu 172° 32′}{}
\EdL{34}{AK07-J0082-L34}{angegeben ist, subtrahiert, gibt die \textit{Sublimatio examinata} =167° 50′\,; diese Hälfte ist}{}
\EdL{35}{AK07-J0082-L35}{hier nur um 2′ kleiner als die Gleichung des Centrums bei \Sperr{Kh.} (vergl. Nachträge).}{}
\EdLp{36}{AK07-J0082-L36}{Die Gründe dieses Verfahrens sind uns, wie schon gesagt, nicht klar geworden,}{}
\EdL{37}{AK07-J0082-L37}{wir haben nur soviel herausgefunden, dass man ungefähr auf denselben wahren Ort}{}
\EdL{38}{AK07-J0082-L38}{des Planeten kommt, ob man nach der Methode \Sperr{Batt.}’s, oder nach der \Sperr{Kh.}’s vorgehe,}{}
\EdL{39}{AK07-J0082-L39}{vorausgesetzt, dass man bei letzterem die \textit{Praezession} berücksichtige, was er eben}{}
\EdL{40}{AK07-J0082-L40}{nicht getan hat. Sehr bedauert aber haben wir bei unsrer Kommentierungs-Arbeit,}{}
\FrameHead{0083}{p. 51}{0083}{p. 51}
\EdL{1}{AK07-J0083-L01}{und besonders bei diesem Kapitel über die wahren Örter der Planeten, dass die}{}
\EdL{2}{AK07-J0083-L02}{„Gründe der Tafeln des \Sperr{Kh.}“ von \Sperr{Al-Bīrūnī}\textsuperscript{1}), in der hebräischen Übersetzung des}{}
\EdL{3}{AK07-J0083-L03}{\Sperr{Ibn ʿEzrā} noch existierend, noch nicht zur Veröffentlichung gelangt sind.}{}
\EdLp{4}{AK07-J0083-L04}{Sehen wir noch zu, was der \textit{Sūrya-Siddhānta} über die Bestimmung des wahren}{}
\EdL{5}{AK07-J0083-L05}{Ortes der Planeten enthält. Seine Regeln\textsuperscript{2}) sind wesentlich komplizierter als diejenigen}{}
\EdL{6}{AK07-J0083-L06}{\Sperr{Kh.}’s, woraus wir wohl den Schluss ziehen dürfen, dass schon die persischen Tafeln des}{}
\EdL{7}{AK07-J0083-L07}{\textit{Shāh} eine Vereinfachung der indischen Methode vorgenommen hatten. Ich führe hier}{}
\EdL{8}{AK07-J0083-L08}{nur die Verse 43 u. 44 (des II. Buches) in der engl. Übersetzung von E. \Sperr{Burgess}\textsuperscript{3}) an\,:}{}
\EdLp{9}{AK07-J0083-L09}{43. \textit{The process of correction for the apsis is the only one required for the sun and}}{}
\EdL{10}{AK07-J0083-L10}{\textit{moon\,; for Mars and the other planets are prescribed that for the conjunction, that for}}{}
\EdL{11}{AK07-J0083-L11}{\textit{the apsis, again that for the apsis and that for the conjunction — four in succession.}}{}
\EdLp{12}{AK07-J0083-L12}{44. \textit{To the mean place of the planet apply half the equation of the conjunction,}}{}
\EdL{13}{AK07-J0083-L13}{\textit{likewise half the equation of the apsis\,; to the mean place of the planet apply the whole}}{}
\EdL{14}{AK07-J0083-L14}{\textit{equation of the apsis and also that of the conjunction.}}{}
\EdLp{15}{AK07-J0083-L15}{Mit diesen Regeln würde man ohne Kommentar wohl nicht viel anfangen können,}{}
\EdL{16}{AK07-J0083-L16}{wir sehen nur soviel daraus, dass es sich dabei auch um \textit{Halbierung} der \textit{Gleichung}}{}
\EdL{17}{AK07-J0083-L17}{handelt. \Sperr{Burgess} kommentiert diese Verse in folgender Weise\,:}{}
\EdLp{18}{AK07-J0083-L18}{\textit{The phraseology of the text is not entirely explicit, and would bear, if taken alone,}}{}
\EdL{19}{AK07-J0083-L19}{\textit{a different interpretation from that which the commentary puts upon it, and which the}}{}
\EdL{20}{AK07-J0083-L20}{\textit{rules to be given later show to be its true meaning\,; this is as follows\,: first calculate}}{}
\EdL{21}{AK07-J0083-L21}{\textit{from the mean place of the planet the equation of the conjunction, and apply the half}}{}
\EdL{22}{AK07-J0083-L22}{\textit{of it to the mean place\,; from the position thus obtained calculate the equation of the}}{}
\EdL{23}{AK07-J0083-L23}{\textit{apsis, and apply half of it to the longitude as already once equated\,; from this result}}{}
\EdL{24}{AK07-J0083-L24}{\textit{find once more the equation of the apsis and apply it to the original mean place of the}}{}
\EdL{25}{AK07-J0083-L25}{\textit{planet\,; and finally calculate from and apply to this last place the whole equation of}}{}
\EdL{26}{AK07-J0083-L26}{\textit{the conjunction.}}{}
\EdLp{27}{AK07-J0083-L27}{\Sperr{Burgess} gibt dann Beispiele für die Berechnung der wahren Örter der fünf}{}
\EdL{28}{AK07-J0083-L28}{Planeten nach diesen Regeln. Wir haben uns lange mit der Vergleichung dieser}{}
\EdL{29}{AK07-J0083-L29}{Berechnungen mit den nach den Regeln \Sperr{Kh.}’s ausgeführten beschäftigt, aber den}{}
\EdL{30}{AK07-J0083-L30}{wahren Zusammenhang nicht herausfinden können. Im \textit{Sūrya-Siddhānta} wird zwei}{}
\EdL{31}{AK07-J0083-L31}{mal eine Gleichung halbiert\,: zuerst die \textit{Gleichung der Konjunktion} (bei uns \textit{Gleichung}}{}
\EdL{32}{AK07-J0083-L32}{\textit{des Centrums}), und die Hälfte dieser Gleichung der mittlern Länge des Planeten hin-}{}
\EdL{33}{AK07-J0083-L33}{zugefügt\,; hierauf wird die \textit{Gleichung der Apsis} (des Epizykels, bei uns \textit{Gleichung der}}{}
\EdL{34}{AK07-J0083-L34}{\textit{Anomalie}) halbiert, und die Hälfte wieder der neuen, korrigierten Länge des Planeten}{}
\EdL{35}{AK07-J0083-L35}{hinzugefügt\,; dann wird von dieser zum zweiten mal korrigierten Länge wieder die}{}
\EdL{36}{AK07-J0083-L36}{\textit{Gleichung der Anomalie} bestimmt und diese zur ursprünglichen mittlern Länge}{}
\EdL{37}{AK07-J0083-L37}{addiert, endlich wird von diesem Resultate die \textit{Gleichung des Centrums} bestimmt und}{}
\EdL{38}{AK07-J0083-L38}{dieselbe zu dem vorhin erhaltenen Resultate addiert, so hat man die \textit{wahre Länge}}{}
\EdL{39}{AK07-J0083-L39}{des Planeten. \Sperr{Burgess} hat nach diesen Regeln den wahren Ort des Saturns für}{}
\NoteRule
\EdN{N1}{AK07-J0083-N01}{\textsuperscript{1}) Oder \Sperr{Aḥmed b. al-Muthannā}, vergl. \textit{Biblioth. mathem.} 3. Folge, XII, p. 282.}
\EdN{N2}{AK07-J0083-N02}{\textsuperscript{2}) \textit{Transl.} \Sperr{Burgess}, p. 211—216.}
\EdN{N3}{AK07-J0083-N03}{\textsuperscript{3}) Ibid. p. 214—215.}
\FrameHead{0084}{p. 52}{0084}{p. 52}
\EdL{1}{AK07-J0084-L01}{den 1. Jan. 1860 für den Meridian von Washington bestimmt und denselben um c. 11°}{}
\EdL{2}{AK07-J0084-L02}{grösser als die mittlere Länge gefunden, wir haben ihn nach den Regeln \Sperr{Kh.}’s und}{}
\EdL{3}{AK07-J0084-L03}{mit denselben Angaben, wie sie \Sperr{Burgess} machte, nur um c. 3° grösser gefunden\,;}{}
\EdL{4}{AK07-J0084-L04}{dies ist ein so starker Unterschied, dass man fast annehmen muss, \Sperr{Burgess} habe}{}
\EdL{5}{AK07-J0084-L05}{die Regeln des \textit{Sūrya-Siddhānta} nicht richtig aufgefasst oder angewandt. In jedem}{}
\EdL{6}{AK07-J0084-L06}{Falle aber müssen, wie schon angedeutet worden ist, die Regeln der persischen Tafeln}{}
\EdL{7}{AK07-J0084-L07}{des \textit{Shāh}, wenn ihnen \Sperr{Kh.} wirklich gefolgt ist, ziemlich von denjenigen des \textit{Sūrya-}}{}
\EdL{8}{AK07-J0084-L08}{\textit{Siddhānta} abgewichen sein\,; von den \textit{zwei} Halbierungen von Gleichungen, die}{}
\EdL{9}{AK07-J0084-L09}{\Sperr{Burgess} in den indischen Tafeln gefunden hat, ist bei \Sperr{Kh.} unserer Ansicht nach}{}
\EdL{10}{AK07-J0084-L10}{nur \textit{eine} geblieben.}{}
\EdLp{11}{AK07-J0084-L11}{Im folgenden geben wir als Beispiel für die Bestimmung des \textit{wahren Ortes} eines}{}
\EdL{12}{AK07-J0084-L12}{obern Planeten diejenige für \textit{Saturn} und zwar sowohl nach den Tafeln \Sperr{Batt.}’s als}{}
\EdL{13}{AK07-J0084-L13}{auch nach denen \Sperr{Kh.}’s.}{}
\EdLp{14}{AK07-J0084-L14}{\textit{Berechnung des wahren Ortes des Saturns} für den Anfang (Mittag) des 15. Ṣafar}{}
\EdL{15}{AK07-J0084-L15}{des Jahres 540 d. H.}{}
\EdLc{16}{AK07-J0084-L16}{1. Nach \Sperr{Al-Battānī}.}{}
\EdL{17}{AK07-J0084-L17}{\parbox[t]{\dimexpr0.2450\linewidth-0.5em\relax}{\raggedright\small\mbox{}Mittlere Länge der Sonne}\hspace{0.5em}\parbox[t]{\dimexpr0.2450\linewidth-0.5em\relax}{\raggedright\small\mbox{}(nach den Tafeln II, 19—23)}\hspace{0.5em}\parbox[t]{\dimexpr0.2450\linewidth-0.5em\relax}{\raggedright\small\mbox{}=}\hspace{0.5em}\parbox[t]{\dimexpr0.2450\linewidth-0.5em\relax}{\raggedright\small\mbox{}140° 46′}\hspace{0.5em}}{}
\EdL{18}{AK07-J0084-L18}{\parbox[t]{\dimexpr0.2450\linewidth-0.5em\relax}{\raggedright\small\mbox{}„ „ des Saturns}\hspace{0.5em}\parbox[t]{\dimexpr0.2450\linewidth-0.5em\relax}{\raggedright\small\mbox{}( „ „ „ II, 24—28)}\hspace{0.5em}\parbox[t]{\dimexpr0.2450\linewidth-0.5em\relax}{\raggedright\small\mbox{}=}\hspace{0.5em}\parbox[t]{\dimexpr0.2450\linewidth-0.5em\relax}{\raggedright\small\mbox{}36° 13′ (396° 13′)}\hspace{0.5em}}{}
\EdL{19}{AK07-J0084-L19}{\parbox[t]{\dimexpr0.2450\linewidth-0.5em\relax}{\raggedright\small\mbox{}}\hspace{0.5em}\parbox[t]{\dimexpr0.4900\linewidth-0.5em\relax}{\raggedright\small\mbox{}Anomalie des Saturns\,:}\hspace{0.5em}\parbox[t]{\dimexpr0.2450\linewidth-0.5em\relax}{\raggedright\small\mbox{}104° 33′}\hspace{0.5em}}{}
\EdL{20}{AK07-J0084-L20}{\parbox[t]{\dimexpr0.3267\linewidth-0.5em\relax}{\raggedright\small\mbox{}Mittlere Länge des Saturns}\hspace{0.5em}\parbox[t]{\dimexpr0.3267\linewidth-0.5em\relax}{\raggedright\small\mbox{}=}\hspace{0.5em}\parbox[t]{\dimexpr0.3267\linewidth-0.5em\relax}{\raggedright\small\mbox{}396° 13′}\hspace{0.5em}}{}
\EdL{21}{AK07-J0084-L21}{\parbox[t]{\dimexpr0.3267\linewidth-0.5em\relax}{\raggedright\small\mbox{}Apogeum (im J. 540)\textsuperscript{1})}\hspace{0.5em}\parbox[t]{\dimexpr0.3267\linewidth-0.5em\relax}{\raggedright\small\mbox{}=}\hspace{0.5em}\parbox[t]{\dimexpr0.3267\linewidth-0.5em\relax}{\raggedright\small\mbox{}248° 30′}\hspace{0.5em}}{}
\EdL{22}{AK07-J0084-L22}{\parbox[t]{\dimexpr0.3267\linewidth-0.5em\relax}{\raggedright\small\mbox{}Centrum}\hspace{0.5em}\parbox[t]{\dimexpr0.3267\linewidth-0.5em\relax}{\raggedright\small\mbox{}=}\hspace{0.5em}\parbox[t]{\dimexpr0.3267\linewidth-0.5em\relax}{\raggedright\small\mbox{}147° 43′}\hspace{0.5em}}{}
\EdL{23}{AK07-J0084-L23}{\parbox[t]{\dimexpr0.3267\linewidth-0.5em\relax}{\raggedright\small\mbox{}Gleichung des Centrums}\hspace{0.5em}\parbox[t]{\dimexpr0.3267\linewidth-0.5em\relax}{\raggedright\small\mbox{}=}\hspace{0.5em}\parbox[t]{\dimexpr0.3267\linewidth-0.5em\relax}{\raggedright\small\mbox{}3° 37′ 40″\textsuperscript{2}) (nach Tafel II, 112)}\hspace{0.5em}}{}
\EdL{24}{AK07-J0084-L24}{\parbox[t]{\dimexpr0.3267\linewidth-0.5em\relax}{\raggedright\small\mbox{}Centrum aequatum}\hspace{0.5em}\parbox[t]{\dimexpr0.3267\linewidth-0.5em\relax}{\raggedright\small\mbox{}=}\hspace{0.5em}\parbox[t]{\dimexpr0.3267\linewidth-0.5em\relax}{\raggedright\small\mbox{}144° 5′ 20″}\hspace{0.5em}}{}
\EdL{25}{AK07-J0084-L25}{\parbox[t]{\dimexpr0.3267\linewidth-0.5em\relax}{\raggedright\small\mbox{}Anomalie}\hspace{0.5em}\parbox[t]{\dimexpr0.3267\linewidth-0.5em\relax}{\raggedright\small\mbox{}=}\hspace{0.5em}\parbox[t]{\dimexpr0.3267\linewidth-0.5em\relax}{\raggedright\small\mbox{}104° 33′}\hspace{0.5em}}{}
\EdL{26}{AK07-J0084-L26}{\parbox[t]{\dimexpr0.3267\linewidth-0.5em\relax}{\raggedright\small\mbox{}Gleichung des Centrums}\hspace{0.5em}\parbox[t]{\dimexpr0.3267\linewidth-0.5em\relax}{\raggedright\small\mbox{}=}\hspace{0.5em}\parbox[t]{\dimexpr0.3267\linewidth-0.5em\relax}{\raggedright\small\mbox{}3° 37′ 40″\textsuperscript{2})}\hspace{0.5em}}{}
\EdL{27}{AK07-J0084-L27}{\parbox[t]{\dimexpr0.3267\linewidth-0.5em\relax}{\raggedright\small\mbox{}Anomalia aequata}\hspace{0.5em}\parbox[t]{\dimexpr0.3267\linewidth-0.5em\relax}{\raggedright\small\mbox{}=}\hspace{0.5em}\parbox[t]{\dimexpr0.3267\linewidth-0.5em\relax}{\raggedright\small\mbox{}108° 10′ 40″}\hspace{0.5em}}{}
\EdLp{28}{AK07-J0084-L28}{Der Anomal. aeq. von 108° 10′ 40″ entsprechen (in Tafel II, 111) in den Kolonnen}{}
\EdL{29}{AK07-J0084-L29}{VI und VII die Zahlen 6° 6′ und 0° 25′, in der Kolonne IV die Zahl 20′, es muss nun}{}
\EdL{30}{AK07-J0084-L30}{0° 25′ im Verhältnis 20′ ∶ 60′ verkleinert werden, dies gibt 8′ 20″\textsuperscript{3})\,; dies addiert man}{}
\EdL{31}{AK07-J0084-L31}{zu der Zahl der VI. Kolonne, also zu 6° 6′, dies gibt 6° 14′ 20″, dies zum Centrum}{}
\EdL{32}{AK07-J0084-L32}{aequat. addiert\textsuperscript{4}), gibt\,:}{}
\NoteRule
\EdN{N1}{AK07-J0084-N01}{\textsuperscript{1}) Hier ist die Praezession von 266 bis 540 d. H. in Berücksichtigung gezogen, das Apogeum}
\EdN{N2}{AK07-J0084-N02}{ist also um (\Frac{274}{68})° = 4° 2′ grösser als es \Sperr{Batt.} angibt (244° 28′).}
\EdN{N3}{AK07-J0084-N03}{\textsuperscript{2}) Ist das Centrum grösser als 180°, so muss die Gleichung des Centrums zum letztern addiert,}
\EdN{N4}{AK07-J0084-N04}{dagegen von der Anomalie subtrahiert werden.}
\EdN{N5}{AK07-J0084-N05}{\textsuperscript{3}) Sind die Differenzen, die dem Centrum aequat. (hier 144° 5′) in der IV. Kolonne entsprechen,}
\EdN{N6}{AK07-J0084-N06}{abnehmend, so muss man die Zahlen nehmen, die der Anomal. aequata in den Kol. VI u. V ent-}
\EdN{N7}{AK07-J0084-N07}{sprechen.}
\EdN{N8}{AK07-J0084-N08}{\textsuperscript{4}) Zu subtrahieren, wenn die Anomal. aeq. grösser als 180° ist.}
\FrameHead{0085}{p. 53}{0085}{p. 53}
\EdL{1}{AK07-J0085-L01}{\parbox[t]{\dimexpr0.1633\linewidth-0.5em\relax}{\raggedright\small\mbox{}}\hspace{0.5em}\parbox[t]{\dimexpr0.3267\linewidth-0.5em\relax}{\raggedright\small\mbox{}Centrum aequat.}\hspace{0.5em}\parbox[t]{\dimexpr0.1633\linewidth-0.5em\relax}{\raggedright\small\mbox{}=}\hspace{0.5em}\parbox[t]{\dimexpr0.3267\linewidth-0.5em\relax}{\raggedright\small\mbox{}144° 5′ 20″}\hspace{0.5em}}{}
\EdL{2}{AK07-J0085-L02}{\parbox[t]{\dimexpr0.1633\linewidth-0.5em\relax}{\raggedright\small\mbox{}}\hspace{0.5em}\parbox[t]{\dimexpr0.3267\linewidth-0.5em\relax}{\raggedright\small\mbox{}Aequat. anomal.corr.}\hspace{0.5em}\parbox[t]{\dimexpr0.1633\linewidth-0.5em\relax}{\raggedright\small\mbox{}=}\hspace{0.5em}\parbox[t]{\dimexpr0.3267\linewidth-0.5em\relax}{\raggedright\small\mbox{}6° 14′ 20″}\hspace{0.5em}}{}
\EdL{3}{AK07-J0085-L03}{\parbox[t]{\dimexpr0.6533\linewidth-0.5em\relax}{\raggedright\small\mbox{}}\hspace{0.5em}\parbox[t]{\dimexpr0.3267\linewidth-0.5em\relax}{\raggedright\small\mbox{}150° 19′ 40″}\hspace{0.5em}}{}
\EdL{4}{AK07-J0085-L04}{\parbox[t]{\dimexpr0.3267\linewidth-0.5em\relax}{\raggedright\small\mbox{}}\hspace{0.5em}\parbox[t]{\dimexpr0.1633\linewidth-0.5em\relax}{\raggedright\small\mbox{}Apogeum}\hspace{0.5em}\parbox[t]{\dimexpr0.1633\linewidth-0.5em\relax}{\raggedright\small\mbox{}=}\hspace{0.5em}\parbox[t]{\dimexpr0.3267\linewidth-0.5em\relax}{\raggedright\small\mbox{}248° 30′}\hspace{0.5em}}{}
\EdL{5}{AK07-J0085-L05}{\parbox[t]{\dimexpr0.3267\linewidth-0.5em\relax}{\raggedright\small\mbox{}}\hspace{0.5em}\parbox[t]{\dimexpr0.1633\linewidth-0.5em\relax}{\raggedright\small\mbox{}addiert\,:}\hspace{0.5em}\parbox[t]{\dimexpr0.1633\linewidth-0.5em\relax}{\raggedright\small\mbox{}=}\hspace{0.5em}\parbox[t]{\dimexpr0.3267\linewidth-0.5em\relax}{\raggedright\small\mbox{}398° 50′}\hspace{0.5em}}{}
\EdL{6}{AK07-J0085-L06}{\parbox[t]{\dimexpr0.6533\linewidth-0.5em\relax}{\raggedright\small\mbox{}}\hspace{0.5em}\parbox[t]{\dimexpr0.3267\linewidth-0.5em\relax}{\raggedright\small\mbox{}—360°}\hspace{0.5em}}{}
\EdL{7}{AK07-J0085-L07}{\parbox[t]{\dimexpr0.4900\linewidth-0.5em\relax}{\raggedright\small\mbox{}}\hspace{0.5em}\parbox[t]{\dimexpr0.1633\linewidth-0.5em\relax}{\raggedright\small\mbox{}=}\hspace{0.5em}\parbox[t]{\dimexpr0.3267\linewidth-0.5em\relax}{\raggedright\small\mbox{}38° 50′ = \textit{wahre Länge}}\hspace{0.5em}}{}
\EdL{8}{AK07-J0085-L08}{d. h. Saturn befand sich zu der gegebenen Zeit in 8° 50′ des Stiers.}{}
\EdLc{9}{AK07-J0085-L09}{2. Nach \Sperr{Al-Khwārizmī}.}{}
\EdL{10}{AK07-J0085-L10}{\parbox[t]{\dimexpr0.4900\linewidth-0.5em\relax}{\raggedright\small\mbox{}Mittlere Länge der Sonne (nach den Tafeln 4 u. 5, s. oben p. 42)}\hspace{0.5em}\parbox[t]{\dimexpr0.1633\linewidth-0.5em\relax}{\raggedright\small\mbox{}=}\hspace{0.5em}\parbox[t]{\dimexpr0.3267\linewidth-0.5em\relax}{\raggedright\small\mbox{}131° 1′\textsuperscript{1})}\hspace{0.5em}}{}
\EdL{11}{AK07-J0085-L11}{\parbox[t]{\dimexpr0.4900\linewidth-0.5em\relax}{\raggedright\small\mbox{}„ „ des Saturns (nach den Tafeln 9 u. 10)}\hspace{0.5em}\parbox[t]{\dimexpr0.1633\linewidth-0.5em\relax}{\raggedright\small\mbox{}=}\hspace{0.5em}\parbox[t]{\dimexpr0.3267\linewidth-0.5em\relax}{\raggedright\small\mbox{}26° 33′ (386° 33′)}\hspace{0.5em}}{}
\EdL{12}{AK07-J0085-L12}{\parbox[t]{\dimexpr0.1633\linewidth-0.5em\relax}{\raggedright\small\mbox{}}\hspace{0.5em}\parbox[t]{\dimexpr0.3267\linewidth-0.5em\relax}{\raggedright\small\mbox{}Anomalie (Argumentum)}\hspace{0.5em}\parbox[t]{\dimexpr0.1633\linewidth-0.5em\relax}{\raggedright\small\mbox{}=}\hspace{0.5em}\parbox[t]{\dimexpr0.3267\linewidth-0.5em\relax}{\raggedright\small\mbox{}104° 28′}\hspace{0.5em}}{}
\EdL{13}{AK07-J0085-L13}{\parbox[t]{\dimexpr0.1633\linewidth-0.5em\relax}{\raggedright\small\mbox{}}\hspace{0.5em}\parbox[t]{\dimexpr0.3267\linewidth-0.5em\relax}{\raggedright\small\mbox{}Mittlere Länge des Saturns}\hspace{0.5em}\parbox[t]{\dimexpr0.1633\linewidth-0.5em\relax}{\raggedright\small\mbox{}=}\hspace{0.5em}\parbox[t]{\dimexpr0.3267\linewidth-0.5em\relax}{\raggedright\small\mbox{}386° 33′}\hspace{0.5em}}{}
\EdL{14}{AK07-J0085-L14}{\parbox[t]{\dimexpr0.1633\linewidth-0.5em\relax}{\raggedright\small\mbox{}}\hspace{0.5em}\parbox[t]{\dimexpr0.3267\linewidth-0.5em\relax}{\raggedright\small\mbox{}Sublimatio (Apogeum) definita (für 104° 28′)}\hspace{0.5em}\parbox[t]{\dimexpr0.1633\linewidth-0.5em\relax}{\raggedright\small\mbox{}=}\hspace{0.5em}\parbox[t]{\dimexpr0.3267\linewidth-0.5em\relax}{\raggedright\small\mbox{}242° 5′\textsuperscript{2})}\hspace{0.5em}}{}
\EdL{15}{AK07-J0085-L15}{\parbox[t]{\dimexpr0.3267\linewidth-0.5em\relax}{\raggedright\small\mbox{}}\hspace{0.5em}\parbox[t]{\dimexpr0.1633\linewidth-0.5em\relax}{\raggedright\small\mbox{}Centrum}\hspace{0.5em}\parbox[t]{\dimexpr0.1633\linewidth-0.5em\relax}{\raggedright\small\mbox{}=}\hspace{0.5em}\parbox[t]{\dimexpr0.3267\linewidth-0.5em\relax}{\raggedright\small\mbox{}144° 28′}\hspace{0.5em}}{}
\EdL{16}{AK07-J0085-L16}{\parbox[t]{\dimexpr0.1633\linewidth-0.5em\relax}{\raggedright\small\mbox{}}\hspace{0.5em}\parbox[t]{\dimexpr0.3267\linewidth-0.5em\relax}{\raggedright\small\mbox{}Hierzu gehört (nach Tafel 31) die Examinatio centri}\hspace{0.5em}\parbox[t]{\dimexpr0.1633\linewidth-0.5em\relax}{\raggedright\small\mbox{}=}\hspace{0.5em}\parbox[t]{\dimexpr0.3267\linewidth-0.5em\relax}{\raggedright\small\mbox{}4° 56′ 30″\textsuperscript{3})}\hspace{0.5em}}{}
\EdL{17}{AK07-J0085-L17}{\parbox[t]{\dimexpr0.3267\linewidth-0.5em\relax}{\raggedright\small\mbox{}}\hspace{0.5em}\parbox[t]{\dimexpr0.1633\linewidth-0.5em\relax}{\raggedright\small\mbox{}subtrahiert\,:}\hspace{0.5em}\parbox[t]{\dimexpr0.1633\linewidth-0.5em\relax}{\raggedright\small\mbox{}=}\hspace{0.5em}\parbox[t]{\dimexpr0.3267\linewidth-0.5em\relax}{\raggedright\small\mbox{}139° 31′ 30″}\hspace{0.5em}}{}
\EdLp{18}{AK07-J0085-L18}{Dies ist das Centrum aequatum oder examinatum.}{}
\EdL{19}{AK07-J0085-L19}{\parbox[t]{\dimexpr0.1633\linewidth-0.5em\relax}{\raggedright\small\mbox{}}\hspace{0.5em}\parbox[t]{\dimexpr0.3267\linewidth-0.5em\relax}{\raggedright\small\mbox{}Anomalie}\hspace{0.5em}\parbox[t]{\dimexpr0.1633\linewidth-0.5em\relax}{\raggedright\small\mbox{}=}\hspace{0.5em}\parbox[t]{\dimexpr0.3267\linewidth-0.5em\relax}{\raggedright\small\mbox{}104° 28′}\hspace{0.5em}}{}
\EdL{20}{AK07-J0085-L20}{\parbox[t]{\dimexpr0.1633\linewidth-0.5em\relax}{\raggedright\small\mbox{}}\hspace{0.5em}\parbox[t]{\dimexpr0.3267\linewidth-0.5em\relax}{\raggedright\small\mbox{}Examinatio centri}\hspace{0.5em}\parbox[t]{\dimexpr0.1633\linewidth-0.5em\relax}{\raggedright\small\mbox{}=}\hspace{0.5em}\parbox[t]{\dimexpr0.3267\linewidth-0.5em\relax}{\raggedright\small\mbox{}4° 56′ 30″\textsuperscript{3})}\hspace{0.5em}}{}
\EdL{21}{AK07-J0085-L21}{\parbox[t]{\dimexpr0.3267\linewidth-0.5em\relax}{\raggedright\small\mbox{}}\hspace{0.5em}\parbox[t]{\dimexpr0.3267\linewidth-0.5em\relax}{\raggedright\small\mbox{}addiert\,:}\hspace{0.5em}\parbox[t]{\dimexpr0.3267\linewidth-0.5em\relax}{\raggedright\small\mbox{}109° 24′ 30″}\hspace{0.5em}}{}
\EdLp{22}{AK07-J0085-L22}{Dies ist das Argumentum (Anomalia) examinatum. Zu diesem wird in Tafel 30}{}
\EdL{23}{AK07-J0085-L23}{die Examinatio argumenti (\textit{elheza}), d. i. die Gleichung der Anomalie gesucht, sie ist}{}
\EdL{24}{AK07-J0085-L24}{= 5° 35′ 35″, diese ist zum Centrum examinatum = 139° 31′ 30″ zu addieren\textsuperscript{4}), dies}{}
\EdL{25}{AK07-J0085-L25}{gibt 145° 7′\,: diese Grösse nennt \Sperr{Kh.} das Centrum ultimum\,; addiert man dazu wieder}{}
\EdL{26}{AK07-J0085-L26}{die Sublimatio examinata = 242° 5′, so hat man die \textit{wahre Länge} des Saturns für die}{}
\EdL{27}{AK07-J0085-L27}{gegebene Zeit, also = 387° 12′ = 27° 12′, d. h. der Saturn befand sich in 27° 12′ des}{}
\EdL{28}{AK07-J0085-L28}{Widders.}{}
\NoteRule
\EdN{N1}{AK07-J0085-N01}{\textsuperscript{1}) Die Sekunden lassen wir weg, wenn sie unter 30 sind, und runden sie zu 1′ auf, wenn sie}
\EdN{N2}{AK07-J0085-N02}{über 30 sind.}
\EdN{N3}{AK07-J0085-N03}{\textsuperscript{2}) Nach den Regeln des Textes muss man, wenn die Anomalie grösser als 180° ist, zu der}
\EdN{N4}{AK07-J0085-N04}{Sublim. examin. noch die Examin. argumenti (\textit{elheza}) addieren, und die Summe von der mittlern}
\EdN{N5}{AK07-J0085-N05}{Länge des Planeten subtrahieren.}
\EdN{N6}{AK07-J0085-N06}{\textsuperscript{3}) Ist das Centrum grösser als 180°, so muss die Examin. centri zum Centrum addiert, dagegen}
\EdN{N7}{AK07-J0085-N07}{von der Anomalie subtrahiert werden.}
\EdN{N8}{AK07-J0085-N08}{\textsuperscript{4}) Wenn das Argumentum (Anomalia) examinatum kleiner als 180° ist, dagegen vom Centrum}
\EdN{N9}{AK07-J0085-N09}{examinatum zu subtrahieren, wenn das Argumentum examinatum grösser als 180° ist.}
\FrameHead{0086}{p. 54}{0086}{p. 54}
\EdLp{1}{AK07-J0086-L01}{Der Fehler von c. 11° gegenüber der Bestimmung des wahren Ortes nach \Sperr{Batt.}}{}
\EdL{2}{AK07-J0086-L02}{rührt in erster Linie von der kleinern mittlern Länge des Saturns bei \Sperr{Kh.} her, die}{}
\EdL{3}{AK07-J0086-L03}{in Folge der Nichtbeachtung der Praezession bei \Sperr{Kh.} um c. 10° kleiner ist als bei}{}
\EdL{4}{AK07-J0086-L04}{\Sperr{Batt.}, in zweiter Linie aber auch von abweichenden Werten in den Gleichungen}{}
\EdL{5}{AK07-J0086-L05}{des Centrums und der Anomalie. Berücksichtigt man auch bei \Sperr{Kh.} die Praezession,}{}
\EdL{6}{AK07-J0086-L06}{verfährt aber im übrigen nach seiner Methode, so erhält man für den wahren Ort}{}
\EdL{7}{AK07-J0086-L07}{des Saturns c. 43′ weniger als nach der Methode \Sperr{Batt.}’s.}{}
\EdLc{8}{AK07-J0086-L08}{\textbf{Kap. 11. Bestimmung des (wahren) Ortes der zwei untern Planeten.}}{}
\EdLc{9}{AK07-J0086-L09}{(Text p. 12, Tafeln 45—56.)}{}
\EdLp{10}{AK07-J0086-L10}{\Sperr{Kh.} bemerkt, dass der wahre Ort von Venus und Merkur auf die gleiche Weise}{}
\EdL{11}{AK07-J0086-L11}{bestimmt werde, wie bei den obern Planeten, dass die mittlere Länge für diese}{}
\EdL{12}{AK07-J0086-L12}{beiden Planeten aber dieselbe sei wie die der Sonne, und dass in den Tafeln für}{}
\EdL{13}{AK07-J0086-L13}{die mittlere Bahn dieser Planeten (Tafeln 15 u. 16) statt der mittlern Länge die}{}
\EdL{14}{AK07-J0086-L14}{Anomalie angegeben sei\textsuperscript{1}). In ähnlicher Weise drückt sich auch \Sperr{Batt.}\textsuperscript{2}) aus.}{}
\EdLp{15}{AK07-J0086-L15}{Wir geben auch hier ein Beispiel für den Planeten Venus und verweisen für die}{}
\EdL{16}{AK07-J0086-L16}{etwas von der der obern Planeten abweichende Theorie und Berechnung der Glei-}{}
\EdL{17}{AK07-J0086-L17}{chungen (besonders für Merkur) auf \Sperr{Batt.}\textsuperscript{3})}{}
\EdLp{18}{AK07-J0086-L18}{\textit{Bestimmung des wahren Ortes der Venus} für Mittag den 15. Ṣafar des Jahres}{}
\EdL{19}{AK07-J0086-L19}{540 d. H.}{}
\EdL{20}{AK07-J0086-L20}{\parbox[t]{\dimexpr0.8167\linewidth-0.5em\relax}{\raggedright\small\mbox{}Anomalie der Venus am Ende des 510. Jahres (nach Tafel 15 u. 16)}\hspace{0.5em}\parbox[t]{\dimexpr0.1633\linewidth-0.5em\relax}{\raggedright\small\mbox{}232° 45′}\hspace{0.5em}}{}
\EdL{21}{AK07-J0086-L21}{\parbox[t]{\dimexpr0.4900\linewidth-0.5em\relax}{\raggedright\small\mbox{}}\hspace{0.5em}\parbox[t]{\dimexpr0.3267\linewidth-0.5em\relax}{\raggedright\small\mbox{}Zunahme in 29 Jahren}\hspace{0.5em}\parbox[t]{\dimexpr0.1633\linewidth-0.5em\relax}{\raggedright\small\mbox{}216° 15′}\hspace{0.5em}}{}
\EdL{22}{AK07-J0086-L22}{\parbox[t]{\dimexpr0.4900\linewidth-0.5em\relax}{\raggedright\small\mbox{}}\hspace{0.5em}\parbox[t]{\dimexpr0.3267\linewidth-0.5em\relax}{\raggedright\small\mbox{}„ „ 1 Monat}\hspace{0.5em}\parbox[t]{\dimexpr0.1633\linewidth-0.5em\relax}{\raggedright\small\mbox{}18° 30′}\hspace{0.5em}}{}
\EdL{23}{AK07-J0086-L23}{\parbox[t]{\dimexpr0.4900\linewidth-0.5em\relax}{\raggedright\small\mbox{}}\hspace{0.5em}\parbox[t]{\dimexpr0.3267\linewidth-0.5em\relax}{\raggedright\small\mbox{}„ „ 14 Tagen}\hspace{0.5em}\parbox[t]{\dimexpr0.1633\linewidth-0.5em\relax}{\raggedright\small\mbox{}8° 38′}\hspace{0.5em}}{}
\EdL{24}{AK07-J0086-L24}{\parbox[t]{\dimexpr0.6533\linewidth-0.5em\relax}{\raggedright\small\mbox{}}\hspace{0.5em}\parbox[t]{\dimexpr0.1633\linewidth-0.5em\relax}{\raggedright\small\mbox{}Summa\,:}\hspace{0.5em}\parbox[t]{\dimexpr0.1633\linewidth-0.5em\relax}{\raggedright\small\mbox{}476° 8′}\hspace{0.5em}}{}
\EdL{25}{AK07-J0086-L25}{\parbox[t]{\dimexpr0.8167\linewidth-0.5em\relax}{\raggedright\small\mbox{}}\hspace{0.5em}\parbox[t]{\dimexpr0.1633\linewidth-0.5em\relax}{\raggedright\small\mbox{}— 360°}\hspace{0.5em}}{}
\EdL{26}{AK07-J0086-L26}{\parbox[t]{\dimexpr0.4900\linewidth-0.5em\relax}{\raggedright\small\mbox{}}\hspace{0.5em}\parbox[t]{\dimexpr0.1633\linewidth-0.5em\relax}{\raggedright\small\mbox{}Anomalie}\hspace{0.5em}\parbox[t]{\dimexpr0.1633\linewidth-0.5em\relax}{\raggedright\small\mbox{}=}\hspace{0.5em}\parbox[t]{\dimexpr0.1633\linewidth-0.5em\relax}{\raggedright\small\mbox{}116° 8′}\hspace{0.5em}}{}
\EdL{27}{AK07-J0086-L27}{\parbox[t]{\dimexpr0.5600\linewidth-0.5em\relax}{\raggedright\small\mbox{}Mittlere Länge wie die der Sonne}\hspace{0.5em}\parbox[t]{\dimexpr0.1400\linewidth-0.5em\relax}{\raggedright\small\mbox{}=}\hspace{0.5em}\parbox[t]{\dimexpr0.2800\linewidth-0.5em\relax}{\raggedright\small\mbox{}131° 1′ (s. p. 42 u. 53)}\hspace{0.5em}}{}
\EdL{28}{AK07-J0086-L28}{\parbox[t]{\dimexpr0.5600\linewidth-0.5em\relax}{\raggedright\small\mbox{}Sublimatio (Apog.) exam. für 116° 8′ (nach Tafel 48)}\hspace{0.5em}\parbox[t]{\dimexpr0.1400\linewidth-0.5em\relax}{\raggedright\small\mbox{}=}\hspace{0.5em}\parbox[t]{\dimexpr0.2800\linewidth-0.5em\relax}{\raggedright\small\mbox{}59° 9′}\hspace{0.5em}}{}
\EdL{29}{AK07-J0086-L29}{\parbox[t]{\dimexpr0.4200\linewidth-0.5em\relax}{\raggedright\small\mbox{}}\hspace{0.5em}\parbox[t]{\dimexpr0.1400\linewidth-0.5em\relax}{\raggedright\small\mbox{}Centrum}\hspace{0.5em}\parbox[t]{\dimexpr0.1400\linewidth-0.5em\relax}{\raggedright\small\mbox{}=}\hspace{0.5em}\parbox[t]{\dimexpr0.2800\linewidth-0.5em\relax}{\raggedright\small\mbox{}71° 52′}\hspace{0.5em}}{}
\EdL{30}{AK07-J0086-L30}{\parbox[t]{\dimexpr0.2450\linewidth-0.5em\relax}{\raggedright\small\mbox{}Centrum}\hspace{0.5em}\parbox[t]{\dimexpr0.2450\linewidth-0.5em\relax}{\raggedright\small\mbox{}71° 52′}\hspace{0.5em}\parbox[t]{\dimexpr0.2450\linewidth-0.5em\relax}{\raggedright\small\mbox{}Anomalie}\hspace{0.5em}\parbox[t]{\dimexpr0.2450\linewidth-0.5em\relax}{\raggedright\small\mbox{}116° 8′}\hspace{0.5em}}{}
\EdL{31}{AK07-J0086-L31}{\parbox[t]{\dimexpr0.2450\linewidth-0.5em\relax}{\raggedright\small\mbox{}Examin. Centri (nach Tafel 47)}\hspace{0.5em}\parbox[t]{\dimexpr0.2450\linewidth-0.5em\relax}{\raggedright\small\mbox{}2° 7′}\hspace{0.5em}\parbox[t]{\dimexpr0.2450\linewidth-0.5em\relax}{\raggedright\small\mbox{}Exam. Centri}\hspace{0.5em}\parbox[t]{\dimexpr0.2450\linewidth-0.5em\relax}{\raggedright\small\mbox{}2° 7′}\hspace{0.5em}}{}
\EdL{32}{AK07-J0086-L32}{\parbox[t]{\dimexpr0.2450\linewidth-0.5em\relax}{\raggedright\small\mbox{}Centrum aequat.}\hspace{0.5em}\parbox[t]{\dimexpr0.2450\linewidth-0.5em\relax}{\raggedright\small\mbox{}69° 45′}\hspace{0.5em}\parbox[t]{\dimexpr0.2450\linewidth-0.5em\relax}{\raggedright\small\mbox{}Argument. exam.}\hspace{0.5em}\parbox[t]{\dimexpr0.2450\linewidth-0.5em\relax}{\raggedright\small\mbox{}118° 15′}\hspace{0.5em}}{}
\EdL{33}{AK07-J0086-L33}{Zu diesem Argum. exam. findet man in Tafel 48 als Gleichung der Anomalie (Exa-}{}
\EdL{34}{AK07-J0086-L34}{minatio elheza) = 44° 43\textsuperscript{1}/\textsubscript{2}′, diese ist zum Centrum aequat. zu addieren, gibt 114° 28\textsuperscript{1}/\textsubscript{2}′,}{}
\EdL{35}{AK07-J0086-L35}{dies ist das Centrum ultimum\,; hiezu die Sublim. examin. (= 59° 9′) addiert, gibt}{}
\NoteRule
\EdN{N1}{AK07-J0086-N01}{\textsuperscript{1}) Dies will der in seinem ersten Teil etwas unklare Satz sagen\,: \textit{quod eorum in libro scriptum}}
\EdN{N2}{AK07-J0086-N02}{\textit{est argumentum atque eorundem et solis unum idemque est elwacat.}}
\EdN{N3}{AK07-J0086-N03}{\textsuperscript{2}) Ed. \Sperr{Nallino}, I, 113.}
\EdN{N4}{AK07-J0086-N04}{\textsuperscript{3}) Ibid. II, 243–244.}
\FrameHead{0087}{p. 55}{0087}{p. 55}
\EdL{1}{AK07-J0087-L01}{173° 37\textsuperscript{1}/\textsubscript{2}′\,; dies war die \textit{wahre Länge} der Venus zu der gegebenen Zeit, d. h. dieselbe}{}
\EdL{2}{AK07-J0087-L02}{befand sich in 23° 37\textsuperscript{1}/\textsubscript{2}′ der Jungfrau.}{}
\EdLp{3}{AK07-J0087-L03}{Nach den Tafeln \Sperr{Batt.}’s findet man die wahre Länge = 181° 17′. Dass hier der}{}
\EdL{4}{AK07-J0087-L04}{Unterschied nicht 10° wie bei Saturn, sondern nur etwa 7° 40′ beträgt, rührt davon}{}
\EdL{5}{AK07-J0087-L05}{her, dass die mittlere Anomalie der Venus bei \Sperr{Kh.} c. 5° grösser ist als bei \Sperr{Batt.}}{}
\EdLc{6}{AK07-J0087-L06}{\textbf{Kap. 12. Bestimmung der Länge der Mondknoten.}}{}
\EdLc{7}{AK07-J0087-L07}{(Text p. 12, Tafeln 19 u. 20.)}{}
\EdLp{8}{AK07-J0087-L08}{Es wird wie bei \Sperr{Batt.} nur der aufsteigende Knoten (Drachenkopf) berücksichtigt,}{}
\EdL{9}{AK07-J0087-L09}{weil die Länge des absteigenden (Drachenschwanzes) um 180° von jenem abweicht.}{}
\EdL{10}{AK07-J0087-L10}{Die mittlere Länge wird berechnet wie die der Sonne und des Mondes nach den}{}
\EdL{11}{AK07-J0087-L11}{Tafeln 19 u. 20 für jeden gegebenen Zeitpunkt, das Ergebnis wird dann von 360°}{}
\EdL{12}{AK07-J0087-L12}{abgezogen, weil die Bewegung der Knoten retrograd ist, der Rest ist die wahre Länge}{}
\EdL{13}{AK07-J0087-L13}{des Knotens. Ich unterlasse es hiefür ein Beispiel zu geben.}{}
\EdLc{14}{AK07-J0087-L14}{\textbf{Kap. 13. Über die Stationen und das Vor- und Rückwärtsschreiten der Planeten.}}{}
\EdLc{15}{AK07-J0087-L15}{(Text p. 12, Tafeln 27—56.)}{}
\EdLp{16}{AK07-J0087-L16}{\Sperr{Kh.} stellt für diese Bestimmungen keine eigenen Tafeln auf, wie \Sperr{Batt.} II, 138}{}
\EdL{17}{AK07-J0087-L17}{u. 139, sondern hat in den Tafeln für die Gleichungen der Planeten eine besondere}{}
\EdL{18}{AK07-J0087-L18}{Kolonne hiefür, betitelt \textit{Statio prima}\,; die Kolonne \textit{Statio secunda} bei \Sperr{Batt.}, die nur}{}
\EdL{19}{AK07-J0087-L19}{die Ergänzungen der Zahlen der \textit{Statio prima} zu 360° enthält, lässt \Sperr{Kh.} weg.}{}
\EdLp{20}{AK07-J0087-L20}{Die Regel ist folgende\,: Für irgend einen gegebenen Zeitpunkt bestimme man wie}{}
\EdL{21}{AK07-J0087-L21}{bei der Berechnung des wahren Ortes des Planeten sein \textit{Centrum aequat.} (\textit{examinat.}),}{}
\EdL{22}{AK07-J0087-L22}{suche zu diesem Argument die \textit{Statio prima}, subtrahiere sie von 360°, so hat man}{}
\EdL{23}{AK07-J0087-L23}{die \textit{Statio secunda}\,; dann bestimme man auch die \textit{Anomalia} (\textit{elheza}) \textit{examinata} des}{}
\EdL{24}{AK07-J0087-L24}{Planeten in diesem Zeitpunkte, ist diese einer der beiden Stationen gleich, so ist der}{}
\EdL{25}{AK07-J0087-L25}{Planet stationnär, ist sie aber zwischen den beiden, also grösser als die erste und}{}
\EdL{26}{AK07-J0087-L26}{kleiner als die zweite, so ist er rückwärtsgehend, im andern Falle (also kleiner als}{}
\EdL{27}{AK07-J0087-L27}{die erste, oder grösser als die zweite) vorwärtsschreitend.\textsuperscript{1})}{}
\EdLp{28}{AK07-J0087-L28}{Ganz gleich lautet die Regel bei \Sperr{Batt.} (I, 114), nur lässt dieser (wohl weil er}{}
\EdL{29}{AK07-J0087-L29}{es für selbstverständlich hielt) den Fall weg, wo die Anomal. exam. einer der beiden}{}
\EdL{30}{AK07-J0087-L30}{Stationen gleich ist.}{}
\EdLp{31}{AK07-J0087-L31}{\textit{Beispiel}\,: Wie verhielt es sich in dieser Beziehung mit Saturn am Mittag des}{}
\EdL{32}{AK07-J0087-L32}{15. Ṣafar des Jahres 540 d. H.\,?}{}
\EdLp{33}{AK07-J0087-L33}{Wir haben oben (p. 53) für Saturn gefunden\,:}{}
\EdL{34}{AK07-J0087-L34}{\parbox[t]{\dimexpr0.1960\linewidth-0.5em\relax}{\raggedright\small\mbox{}}\hspace{0.5em}\parbox[t]{\dimexpr0.3920\linewidth-0.5em\relax}{\raggedright\small\mbox{}Centrum aequatum}\hspace{0.5em}\parbox[t]{\dimexpr0.3920\linewidth-0.5em\relax}{\raggedright\small\mbox{}Anomalia aequata}\hspace{0.5em}}{}
\EdL{35}{AK07-J0087-L35}{\parbox[t]{\dimexpr0.1960\linewidth-0.5em\relax}{\raggedright\small\mbox{}}\hspace{0.5em}\parbox[t]{\dimexpr0.3920\linewidth-0.5em\relax}{\raggedright\small\mbox{}139° 31\textsuperscript{1}/\textsubscript{2}′}\hspace{0.5em}\parbox[t]{\dimexpr0.3920\linewidth-0.5em\relax}{\raggedright\small\mbox{}109° 24\textsuperscript{1}/\textsubscript{2}′}\hspace{0.5em}}{}
\NoteRule
\EdN{N1}{AK07-J0087-N01}{\textsuperscript{1}) Der zweite Satz bei \Sperr{Batt.}\,: \textit{si contra maior secunda statione et minor prima, directus}, hat}
\EdN{N2}{AK07-J0087-N02}{einen kleinen Fehler, denn so ist der Fall unmöglich, es muss einfach statt \textit{et} heissen \textit{aut}\,; der}
\EdN{N3}{AK07-J0087-N03}{Fehler befindet sich schon im arab. Text, wo statt \textit{wa} stehen sollte \textit{aw} (au). (Ed. \Sperr{Nallino}, I, 114}
\EdN{N4}{AK07-J0087-N04}{u. III, 173).}
\FrameHead{0088}{p. 56}{0088}{p. 56}
\EdLr{1}{AK07-J0088-L01}{Zu 139°\textsuperscript{1}) gehört in Tafel 31 die \textit{Statio prima} 115° 10′}{}
\EdLr{2}{AK07-J0088-L02}{also ist die \textit{Statio secunda} = 244° 50′}{}
\EdLp{3}{AK07-J0088-L03}{Nun ist die \textit{Anomal. aeq.} kleiner als die erste Station, also war Saturn an jenem}{}
\EdL{4}{AK07-J0088-L04}{Tage \textit{vorwärtsschreitend} (\textit{rechtläufig}). Die Regel ist wohl ohne weitern Kommentar}{}
\EdL{5}{AK07-J0088-L05}{zu verstehen.}{}
\EdLc{6}{AK07-J0088-L06}{\textbf{Kap. 14. Über die Zeit der Rechtläufigkeit und Rückläufigkeit.}}{}
\EdLc{7}{AK07-J0088-L07}{(Text p. 12—13, Tafeln 27—56.)}{}
\EdLp{8}{AK07-J0088-L08}{Will man die Zeit wissen, während welcher ein Planet rechtläufig oder rück-}{}
\EdL{9}{AK07-J0088-L09}{läufig ist, so hat man nach \Sperr{Kh.} so zu verfahren\,:}{}
\EdLp{10}{AK07-J0088-L10}{a) Der Planet sei \textit{rückläufig.} Will man nun wissen, wie lange er schon rück-}{}
\EdL{11}{AK07-J0088-L11}{läufig war, so nimmt man den Unterschied zwischen der \textit{Statio prima} und der}{}
\EdL{12}{AK07-J0088-L12}{\textit{Anomal. aeq.} und bestimmt, wie viel Zeit für den Planeten nötig war, um diesen}{}
\EdL{13}{AK07-J0088-L13}{Unterschied zu durchlaufen\,; will man aber wissen, wie lange er noch rückläufig}{}
\EdL{14}{AK07-J0088-L14}{bleibe, so nimmt man den Unterschied zwischen der \textit{Statio secunda} und der \textit{Anomal.}}{}
\EdL{15}{AK07-J0088-L15}{\textit{aeq.} und bestimmt, wie viel Zeit der Planet gebrauche, um diesen Unterschied zu}{}
\EdL{16}{AK07-J0088-L16}{durchlaufen. Um nun diese Zeiten zu finden, muss man die tägliche Bewegung der}{}
\EdL{17}{AK07-J0088-L17}{Anomalie des Planeten kennen, diese führt \Sperr{Batt.} in seiner Darstellung dieser Be-}{}
\EdL{18}{AK07-J0088-L18}{stimmung an. \Sperr{Kh.} spricht von diesen täglichen Bewegungen nicht, er sagt einfach\,:}{}
\EdL{19}{AK07-J0088-L19}{\textit{Deinde per argumentum} (d. i. hier = log. Schluss, Verfahren) \textit{ad hoc negotium pro-}}{}
\EdL{20}{AK07-J0088-L20}{\textit{prium, quanto tempore illud transeat, coniectabitur.} Ein mit der Sache nicht Vertrauter}{}
\EdL{21}{AK07-J0088-L21}{würde schwerlich herausfinden, was hiermit gemeint sei\,; da führen uns die klaren}{}
\EdL{22}{AK07-J0088-L22}{Regeln \Sperr{Batt.}’s auf den richtigen Weg\,: man muss die genannten Unterschiede zwischen}{}
\EdL{23}{AK07-J0088-L23}{den Stationen und der Anomal. aeq. durch die tägliche Bewegung der Anomalie des}{}
\EdL{24}{AK07-J0088-L24}{Planeten teilen\textsuperscript{2}), und erhält dann die Zahl der Tage, die der Planet schon rückläufig}{}
\EdL{25}{AK07-J0088-L25}{ist, bezw. noch rückläufig bleibt.}{}
\EdLp{26}{AK07-J0088-L26}{b) Der Planet sei \textit{rechtläufig.} Für diesen Fall sind die Regeln bei \Sperr{Batt.} sowohl}{}
\EdL{27}{AK07-J0088-L27}{als bei \Sperr{Kh.} unvollständig, bei letzterem noch dazu unverständlich. Man muss hier}{}
\EdL{28}{AK07-J0088-L28}{zwei Fälle unterscheiden, nämlich ob die Anomal. aeq. kleiner sei als die 1. Station,}{}
\EdL{29}{AK07-J0088-L29}{oder grösser als die 2.\,; in beiden Fällen kann man wieder fragen, wie lange bleibt}{}
\EdL{30}{AK07-J0088-L30}{der Planet noch rechtläufig, und wie lange war er schon rechtläufig. \Sperr{Batt.} betrachtet}{}
\EdL{31}{AK07-J0088-L31}{nur den ersten Fall und stellt nur die erste Frage, \Sperr{Kh.} betrachtet beide Fälle, stellt}{}
\EdL{32}{AK07-J0088-L32}{aber bei jedem nur eine Frage. Die vollständige Regel wäre folgende\,:}{}
\EdLp{33}{AK07-J0088-L33}{\textit{1. Fall}\,: Die \textit{Anomal. aeq.} ist \textit{kleiner} als die 1. Station. Will man wissen, wie}{}
\EdL{34}{AK07-J0088-L34}{lange der Planet noch rechtläufig bleibt, so nimmt man den Unterschied zwischen}{}
\EdL{35}{AK07-J0088-L35}{beiden, und dividiert ihn durch die tägliche Bewegung der Anomalie des Planeten\textsuperscript{3}).}{}
\NoteRule
\EdN{N1}{AK07-J0088-N01}{\textsuperscript{1}) Ich nehme nur die ganzen Grade.}
\EdN{N2}{AK07-J0088-N02}{\textsuperscript{2}) Diese Zahlen für die tägliche Bewegung der Planeten findet man auf folgenden Tafeln}
\EdN{N3}{AK07-J0088-N03}{(unter der Kolonnenüberschrift\,: \textit{Heza} (Anomalia oder Argumentum))\,: Tafel 9, Saturn = 57′ 8″\,;}
\EdN{N4}{AK07-J0088-N04}{Tafel 11, Jupiter = 54′ 9″\,; Tafel 13, Mars = 27′ 42″\,; Tafel 16, Venus = 37′ 0″\,; Tafel 18, Merkur}
\EdN{N5}{AK07-J0088-N05}{= 3° 6′ 24″.}
\EdN{N6}{AK07-J0088-N06}{\textsuperscript{3}) Diesen Fall berücksichtigt \Sperr{Batt.} am Schlusse des 46. Kap. (ed. \Sperr{Nallino}, I, 114) und}
\EdN{N7}{AK07-J0088-N07}{\Sperr{Kh.} am Schlusse seiner Regel (p. 13)\,: \textit{Quod si non procul a statione prima fuerit} . . . .}
\FrameHead{0089}{p. 57}{0089}{p. 57}
\EdL{1}{AK07-J0089-L01}{Zu diesem Falle gehört unser obiges Beispiel für den Planeten Saturn\,: Er ist, wie}{}
\EdL{2}{AK07-J0089-L02}{wir gefunden haben, rechtläufig, der Unterschied zwischen Statio I. und Anomal. aeq.}{}
\EdL{3}{AK07-J0089-L03}{ist 5° 45\textsuperscript{1}/\textsubscript{2}′, also bleibt er noch \Frac{5° 45\textsuperscript{1}/\textsubscript{2}′}{57′} = c. 6 Tage rechtläufig, wird dann stationnär}{}
\EdL{4}{AK07-J0089-L04}{und hierauf rückläufig. — Will man wissen, wie lange der Planet schon rechtläufig}{}
\EdL{5}{AK07-J0089-L05}{war, so addiert man die Anomal. aeq. zur 1. Station\textsuperscript{1}) und teilt die Summe durch}{}
\EdL{6}{AK07-J0089-L06}{57′\,; dies gibt in unserm Falle \Frac{224° 34\textsuperscript{1}/\textsubscript{2}′}{57′} = 236 Tage\,; im ganzen war er also 242}{}
\EdL{7}{AK07-J0089-L07}{Tage rechtläufig.}{}
\EdLp{8}{AK07-J0089-L08}{\textit{2. Fall}\,: Die \textit{Anomal. aeq.} ist \textit{grösser} als die 2. Station. Will man wissen, wie}{}
\EdL{9}{AK07-J0089-L09}{lange der Planet noch rechtläufig bleibt, so subtrahiert man die Anomal. aeq. von}{}
\EdL{10}{AK07-J0089-L10}{360°, addiert zum Reste die 1. Station, und dividiert das Resultat durch die tägliche}{}
\EdL{11}{AK07-J0089-L11}{Bewegung der Anomalie des Planeten. Will man wissen, wie lange der Planet schon}{}
\EdL{12}{AK07-J0089-L12}{rechtläufig war, so nimmt man den Unterschied zwischen der Anomal. aeq. und der}{}
\EdL{13}{AK07-J0089-L13}{2. Station, und teilt ihn durch die tägliche Bewegung der Anomalie des Planeten\textsuperscript{2}).}{}
\EdLp{14}{AK07-J0089-L14}{Dieses alles wird sich wohl leicht aus folgender Figur erkennen lassen, wenn}{}
\EdL{15}{AK07-J0089-L15}{wir derselben noch einige Bemerkungen beifügen\,:}{}
\EdLp{16}{AK07-J0089-L16}{\textit{E} sei die Erde, \textit{H} das Centrum}{}
\EdL{17}{AK07-J0089-L17}{des Epizykels, \textit{B} und \textit{C} sind die}{}
\EdL{18}{AK07-J0089-L18}{Stationen I und II, bei denen die}{}
\EdL{19}{AK07-J0089-L19}{Rückläufigkeit beginnt, bezw. auf-}{}
\EdL{20}{AK07-J0089-L20}{hört. Nach der Theorie des \Sperr{Ptol.},}{}
\EdL{21}{AK07-J0089-L21}{der auch \Sperr{Batt.} folgt, beginnt die}{}
\EdL{22}{AK07-J0089-L22}{Rückläufigkeit nicht etwa in dem}{}
\EdL{23}{AK07-J0089-L23}{Berührungspunkt \textit{D} der Tangente}{}
\EdL{24}{AK07-J0089-L24}{von \textit{E} an den Epizykel, sondern in}{}
\EdL{25}{AK07-J0089-L25}{einem Punkte \textit{B}, der so liegt, dass}{}
\EdL{26}{AK07-J0089-L26}{\textit{FB} ∶ \textit{BE} = dem Verhältnis der Geschwindigkeit des Epizykelzentrums zu derjenigen}{}
\EdL{27}{AK07-J0089-L27}{des Planeten im Epizykel ist\,; ist jenes Verhältnis (\textit{FB} ∶ \textit{BE}) grösser als dieses, so ist}{}
\EdL{28}{AK07-J0089-L28}{der Planet rückläufig, und dies findet statt in dem Bogen \textit{BC}, in dem Bogen \textit{CAB}}{}
\EdL{29}{AK07-J0089-L29}{ist er rechtläufig (wobei auch die geringe Zeit des Stillstandes inbegriffen ist). Die}{}
\EdL{30}{AK07-J0089-L30}{Stationen I und II sind aber nicht konstant, sondern ändern sich mit der Lage des}{}
\EdL{31}{AK07-J0089-L31}{Epizykelcentrums \textit{H} im \textit{Deferens}, d. h. mit dem sog. \textit{Centrum} des Planeten\,; daher}{}
\EdL{32}{AK07-J0089-L32}{muss man mit dem \textit{Centrum} (\textit{aequat.}) in die Tafeln der Recht- und Rückläufigkeit}{}
\EdL{33}{AK07-J0089-L33}{eintreten. Über die Berechnung der Stationen für jedes beliebige Centrum von 0°—180°}{}
\EdL{34}{AK07-J0089-L34}{vergleiche man \Sperr{Batt.}\textsuperscript{3}) und \Sperr{Ptol.}\textsuperscript{4})}{}
\NoteRule
\EdN{N1}{AK07-J0089-N01}{\textsuperscript{1}) Eigentlich\,: (360° — Stat. II.) + Anomal. aeq.}
\EdN{N2}{AK07-J0089-N02}{\textsuperscript{2}) Diesen Fall berüchsichtigt \Sperr{Batt.} gar nicht, \Sperr{Kh.} nur die 2. Frage dieses Falles mit der Stelle}
\EdN{N3}{AK07-J0089-N03}{(p. 13), die beginnt\,: \textit{Si vero de progrediente scire desideras.} Der Text des 46. Kap. bei \Sperr{Batt.} weicht}
\EdN{N4}{AK07-J0089-N04}{übrigens in der Ausgabe von 1645 an einigen Stellen von der Übersetzung \Sperr{Nallinos} und dem}
\EdN{N5}{AK07-J0089-N05}{arab. Texte ab\,; es ist wahrscheinlich, dass der letztere Auslassungen aufweist, die der arab. Text,}
\EdN{N6}{AK07-J0089-N06}{der \Sperr{Plato von Tivoli} vorlag, nicht hatte.}
\EdN{N7}{AK07-J0089-N07}{\textsuperscript{3}) Ed. \Sperr{Nallino}, II, 245—247.}
\EdN{N8}{AK07-J0089-N08}{\textsuperscript{4}) Ed. \Sperr{Heiberg}, II, 450—494\,; Übersetzung v. K. \Sperr{Manitius}, II, 267– 302.}
\FrameHead{0090}{p. 58}{0090}{p. 58}
\EdLp{1}{AK07-J0090-L01}{Die Werte der \textit{Statio prima} bei \Sperr{Kh.} stimmen mehr mit denen bei \Sperr{Batt.}, als}{}
\EdL{2}{AK07-J0090-L02}{mit denen bei \Sperr{Ptol.} überein, doch kommen zwischen \Sperr{Kh.} und \Sperr{Batt.} bisweilen}{}
\EdL{3}{AK07-J0090-L03}{Unterschiede von 1′ bis 2′ vor. Die grössern Abweichungen bei \Sperr{Ptol.} rühren davon}{}
\EdL{4}{AK07-J0090-L04}{her, dass dieser als Argument das \textit{mittlere Centrum}, die beiden arabischen Astronomen}{}
\EdL{5}{AK07-J0090-L05}{aber das \textit{Centrum aequatum} annehmen. Die Übereinstimmung zwischen \Sperr{Batt.} und}{}
\EdL{6}{AK07-J0090-L06}{\Sperr{Kh.} ist gewiss keine zufällige, ohne Zweifel hat \Sperr{Batt.} die Tafeln \Sperr{Kh.}’s gekannt und}{}
\EdL{7}{AK07-J0090-L07}{wohl auch benutzt, wenn er sie auch nie zitiert.}{}
\EdLc{8}{AK07-J0090-L08}{\textbf{Kap. 15. Über die Deklination der Sonne.}}{}
\EdLc{9}{AK07-J0090-L09}{(Text p. 13, Tafeln 21—26.)}{}
\EdLp{10}{AK07-J0090-L10}{Über diese Tafeln ist wenig zu bemerken. \textit{Porta} im Titel ist die Übersetzung}{}
\EdL{11}{AK07-J0090-L11}{des arabischen \textit{bāb} = Pforte, Kapitel. Die Kolonne in den Tafeln 21—26, die die}{}
\EdL{12}{AK07-J0090-L12}{Deklination der Sonne enthält, ist überschrieben\,: \textit{Obliquatio solis} = Neigung der}{}
\EdL{13}{AK07-J0090-L13}{Sonne (arab. \textit{mail al-shams})\,; wir finden dann von 0° bis 90° (bezw. 360°—270°) hin-}{}
\EdL{14}{AK07-J0090-L14}{zugefügt \textit{accedentis} (für \textit{ascendentis}), von 90° bis 180° (270°—180°) \textit{recedentis} (für}{}
\EdL{15}{AK07-J0090-L15}{\textit{descendentis}). Die \textit{grösste Deklination}, also \textit{Schiefe der Ekliptik}, gibt \Sperr{Kh.} zu 23° 51′ an.}{}
\EdL{16}{AK07-J0090-L16}{Der \textit{Sūrya-Siddhānta}\textsuperscript{1}) hat 24°, die Verfasser der „erprobten Tafeln“ geben sie zu}{}
\EdL{17}{AK07-J0090-L17}{23° 33′ an, \Sperr{Ptol.} hat den Wert 23° 51′ 20″, \Sperr{Batt.} 23° 35′\,; \Sperr{Kh.} schliesst sich also am}{}
\EdL{18}{AK07-J0090-L18}{nächsten an \Sperr{Ptol.} an, er lässt bloss die 20″ weg (vergl. oben p. 33). Es ist übrigens}{}
\EdL{19}{AK07-J0090-L19}{zu bemerken, dass eine Reihe arabischer Astronomen den ptolemäischen Wert eben-}{}
\EdL{20}{AK07-J0090-L20}{falls mit Weglassung der Sekunden zu 23° 51′ angeben. Die Werte \Sperr{Kh.}’s weichen}{}
\EdL{21}{AK07-J0090-L21}{von denen nach der Formel sin \textit{δ} = sin \textit{l} · sin \textit{e} (wo \textit{l} die Grade der Ekliptik vom}{}
\EdL{22}{AK07-J0090-L22}{Frühlingspunkt an gezählt, also die Länge der Sonne, und \textit{e} die Schiefe der Ekliptik}{}
\EdL{23}{AK07-J0090-L23}{bezeichnen) berechneten im Maximum um 2″ ab. Diese Deklinationswerte bilden}{}
\EdL{24}{AK07-J0090-L24}{ein schönes Zeugnis dafür, mit welcher Genauigkeit \Sperr{Kh.} bei der Abfassung seines}{}
\EdL{25}{AK07-J0090-L25}{Tafelwerkes vorgegangen ist.}{}
\EdLc{26}{AK07-J0090-L26}{\textbf{Kap. 16. Über die Breite des Mondes.}}{}
\EdLc{27}{AK07-J0090-L27}{(Text p. 14, Tafel 21—26, u. 19—20.)}{}
\EdLp{28}{AK07-J0090-L28}{Die Breite des Mondes wird nach der gleichen Formel bestimmt wie die Deklina-}{}
\EdL{29}{AK07-J0090-L29}{tion der Sonne, also \textit{sin latit. lunae} = \textit{sin des Argumentes der Breite} × \textit{sin der Schiefe}}{}
\EdL{30}{AK07-J0090-L30}{\textit{der Mondbahn.} Unter dem Argument der Breite versteht man die Entfernung des}{}
\EdL{31}{AK07-J0090-L31}{wahren Ortes des Mondes in der Ekliptik vom nördlichen Knoten (Drachenkopf).}{}
\EdL{32}{AK07-J0090-L32}{Dieses \textit{Argument der Breite} wird auf folgende Weise gefunden\,:}{}
\EdLp{33}{AK07-J0090-L33}{Man addiert zum wahren Ort des Mondes die mittlere Länge des Dıachenkopfes}{}
\EdL{34}{AK07-J0090-L34}{[und subtrahiert von dieser Summe, wenn sie 360° übersteigt, diese Grösse]\textsuperscript{2}).}{}
\EdLp{35}{AK07-J0090-L35}{\Sperr{Batt.} hat dieselbe Regel, dann aber noch folgende, die \Sperr{Kh.} weglässt\,: Man}{}
\EdL{36}{AK07-J0090-L36}{subtrahiert den wahren Ort des Drachenkopfes vom wahren Ort des Mondes, so hat}{}
\EdL{37}{AK07-J0090-L37}{man das Argument der Breite\textsuperscript{3}).}{}
\NoteRule
\EdN{N1}{AK07-J0090-N01}{\textsuperscript{1}) Ed. \Sperr{Burgess}, p. 201.}
\EdN{N2}{AK07-J0090-N02}{\textsuperscript{2}) Was in eckigen Klammern steht, fehlt bei \Sperr{Kh.}}
\EdN{N3}{AK07-J0090-N03}{\textsuperscript{3}) Ed. \Sperr{Nallino}, I, 76.}
\FrameHead{0091}{p. 59}{0091}{p. 59}
\EdLp{1}{AK07-J0091-L01}{Die letztere Regel kommt auch bei den Indern vor\textsuperscript{1}), und gilt in dieser Form}{}
\EdL{2}{AK07-J0091-L02}{auch für die drei obern Planeten\,; es ist wohl zu beachten, dass \Sperr{Kh.} die Regel nicht}{}
\EdL{3}{AK07-J0091-L03}{in der indischen Form (wenigstens des \textit{Sūrya-Siddhānta}) gibt.}{}
\EdLp{4}{AK07-J0091-L04}{Als grösste Breite (also Schiefe der Mondbahn zur Ekliptik) hat \Sperr{Kh.} beim Argu-}{}
\EdL{5}{AK07-J0091-L05}{ment 90° den Wert 4° 30′, während \Sperr{Batt.} wie \Sperr{Ptol.} 5° hat. Hier folgt nun \Sperr{Kh.} dem}{}
\EdL{6}{AK07-J0091-L06}{\textit{Sūrya-Siddhānta} (oder den Tafeln des \textit{Shāh}), wenigstens findet sich der gleiche Wert}{}
\EdL{7}{AK07-J0091-L07}{im erstern\textsuperscript{2}), er mag aber auch in die spätern \textit{Siddhāntas} übergegangen sein.}{}
\EdLp{8}{AK07-J0091-L08}{\textit{Beispiel}\,: Es ist die Breite des Mondes für Mittag den 15. Ṣafar des Jahres 540}{}
\EdL{9}{AK07-J0091-L09}{d. H. zu bestimmen.}{}
\ColHead{left column}
\EdL{10}{AK07-J0091-L10}{Wahrer Ort des Mondes (s. oben p. 45)\,:}{}
\EdL{11}{AK07-J0091-L11}{319° 2′\textsuperscript{3})}{}
\EdL{12}{AK07-J0091-L12}{279° 41′ = mittlere Länge des Drachen-}{}
\EdL{13}{AK07-J0091-L13}{kopfes}{}
\EdL{14}{AK07-J0091-L14}{598° 43′}{}
\EdL{15}{AK07-J0091-L15}{— 360°}{}
\EdL{16}{AK07-J0091-L16}{238° 43′ = Argument der Breite}{}
\EdL{17}{AK07-J0091-L17}{(nach der 1. Regel)}{}
\EdL{18}{AK07-J0091-L18}{319° 2′ = wahrer Ort des Mondes}{}
\EdL{19}{AK07-J0091-L19}{80° 19′ = korrig. Länge (= wahrer Ort)}{}
\EdL{20}{AK07-J0091-L20}{des Drachenkopfes}{}
\EdL{21}{AK07-J0091-L21}{238° 43′ = Argument der Breite}{}
\EdL{22}{AK07-J0091-L22}{(nach der 2. Regel)}{}
\ColHead{right column}
\EdL{23}{AK07-J0091-L23}{Mittlere Länge des Drachenkopfes\,:}{}
\EdL{24}{AK07-J0091-L24}{(nach den Tafeln 19 u. 20)}{}
\EdL{25}{AK07-J0091-L25}{am Ende des 510. Jahres 92° 39′\textsuperscript{3})}{}
\EdL{26}{AK07-J0091-L26}{für 29 Jahre 184° 42′}{}
\EdL{27}{AK07-J0091-L27}{„ 1 Monat 1° 35′}{}
\EdL{28}{AK07-J0091-L28}{„ 14 Tage 45′}{}
\EdL{29}{AK07-J0091-L29}{Summa\,: 279° 41′}{}
\EdL{30}{AK07-J0091-L30}{360° — 279° 41′ = 80° 19′ =}{}
\EdL{31}{AK07-J0091-L31}{korrigierte Länge, da der Drachen-}{}
\EdL{32}{AK07-J0091-L32}{kopf retrograd ist.}{}
\ColHead{full width}
\EdLp{33}{AK07-J0091-L33}{Mit diesem Argument tritt man in die Tafel 25 ein, und findet da in der letzten}{}
\EdL{34}{AK07-J0091-L34}{\parbox[t]{\dimexpr0.4900\linewidth-0.5em\relax}{\raggedright\small\mbox{}Kolonne für 238° (7\textsuperscript{s} 28°) die Latitudo Lunae}\hspace{0.5em}\parbox[t]{\dimexpr0.2450\linewidth-0.5em\relax}{\raggedright\small\mbox{}=}\hspace{0.5em}\parbox[t]{\dimexpr0.2450\linewidth-0.5em\relax}{\raggedright\small\mbox{}3° 48′ 57″}\hspace{0.5em}}{}
\EdL{35}{AK07-J0091-L35}{\parbox[t]{\dimexpr0.2450\linewidth-0.5em\relax}{\raggedright\small\mbox{}}\hspace{0.5em}\parbox[t]{\dimexpr0.2450\linewidth-0.5em\relax}{\raggedright\small\mbox{}Differenz für 43′}\hspace{0.5em}\parbox[t]{\dimexpr0.2450\linewidth-0.5em\relax}{\raggedright\small\mbox{}=}\hspace{0.5em}\parbox[t]{\dimexpr0.2450\linewidth-0.5em\relax}{\raggedright\small\mbox{}1′ 46″}\hspace{0.5em}}{}
\EdL{36}{AK07-J0091-L36}{\parbox[t]{\dimexpr0.2450\linewidth-0.5em\relax}{\raggedright\small\mbox{}}\hspace{0.5em}\parbox[t]{\dimexpr0.4900\linewidth-0.5em\relax}{\raggedright\small\mbox{}\textit{Breite} des Mondes zu der gegeb. Zeit\,:}\hspace{0.5em}\parbox[t]{\dimexpr0.2450\linewidth-0.5em\relax}{\raggedright\small\mbox{}3° 50′ 43″}\hspace{0.5em}}{}
\EdLp{37}{AK07-J0091-L37}{Nach den Bemerkungen, die \Sperr{Batt.} seinen Regeln noch beifügt, \Sperr{Kh.} aber weg-}{}
\EdL{38}{AK07-J0091-L38}{lässt, ist die Breite des Mondes in diesem Falle \textit{südlich} und \textit{wachsend,} d. h. der Mond}{}
\EdL{39}{AK07-J0091-L39}{geht noch weiter gegen Süden, was auch die Überschrift der letzten Kolonne, \textit{Latitudo}}{}
\EdL{40}{AK07-J0091-L40}{\textit{lunae recedentis,} andeutet.}{}
\EdLp{41}{AK07-J0091-L41}{Die Überschriften der beiden ersten Kolonnen der Tafeln 21—26\,: \textit{Computatio}}{}
\EdL{42}{AK07-J0091-L42}{\textit{minuenda et augmentum examinis et numeri} — \textit{Computatio augenda et diminutio exa-}}{}
\EdL{43}{AK07-J0091-L43}{\textit{minis et numeri} sind nicht ganz klar, sie wollen aber wahrscheinlich sagen\,: Wenn das}{}
\EdL{44}{AK07-J0091-L44}{Argument kleiner als 180° ist, was in der 1. Kolonne der Fall ist, so muss die}{}
\EdL{45}{AK07-J0091-L45}{gefundene Gleichung (\textit{Computatio}) von der mittlern Länge abgezogen werden, wenn}{}
\EdL{46}{AK07-J0091-L46}{es grösser als 180° ist, was in der 2. Kolonne der Fall ist, so muss die Gleichung}{}
\EdL{47}{AK07-J0091-L47}{addiert werden\,; ferner, wenn das Argument von 0° bis 180° zunimmt, so nimmt}{}
\EdL{48}{AK07-J0091-L48}{auch die Gleichung (\textit{examen}) zu, wenn das Argument aber von 180° bis 360° wächst,}{}
\EdL{49}{AK07-J0091-L49}{so nimmt die Gleichung wieder ab.}{}
\NoteRule
\EdN{N1}{AK07-J0091-N01}{\textsuperscript{1}) \textit{Sūrya-Siddhānta}, ed. \Sperr{Burgess}, p. 228.}
\EdN{N2}{AK07-J0091-N02}{\textsuperscript{2}) Ibid. p. 188—189.}
\EdN{N3}{AK07-J0091-N03}{\textsuperscript{3}) Die Sekunden runde ich zu Minuten ab oder auf.}
\FrameHead{0092}{p. 60}{0092}{p. 60}
\EdLc{1}{AK07-J0092-L01}{\textbf{Kap. 17. Über die Breite der Planeten.}}{}
\EdLc{2}{AK07-J0092-L02}{(Text p. 14, Tafeln 27–56.)}{}
\EdLp{3}{AK07-J0092-L03}{Der lateinische Titel dieses Kapitels ist unrichtig, dasselbe handelt im 1. Alinea}{}
\EdL{4}{AK07-J0092-L04}{über die drei obern, im 2. Alinea über die zwei untern Planeten. Die Theorie der}{}
\EdL{5}{AK07-J0092-L05}{Änderung der Breite der Planeten ist die schwierigste und dunkelste Partie (13. Buch)}{}
\EdL{6}{AK07-J0092-L06}{bei \Sperr{Ptol.}, so schwierig und dunkel, dass \Sperr{Batt.} darauf verzichtet hat, auf diese Theorie}{}
\EdL{7}{AK07-J0092-L07}{näher einzutreten, er gibt nur die Regeln, nach denen man die Breite der Planeten}{}
\EdL{8}{AK07-J0092-L08}{aus den Tafeln finden kann\textsuperscript{1}).}{}
\EdLp{9}{AK07-J0092-L09}{Aber dieses Dunkel hat \Sperr{Nallino}, nachdem \Sperr{Delambre} u. a. daran geschei-}{}
\EdL{10}{AK07-J0092-L10}{tert sind, in vorzüglicher Weise gelichtet, und dadurch die Theorie des \Sperr{Ptol.}}{}
\EdL{11}{AK07-J0092-L11}{unserm Verständnis näher gebracht\,; er hat aber auch versucht, dieselbe auf}{}
\EdL{12}{AK07-J0092-L12}{mathematischem Wege zu begründen, was ihm ebenfalls in schöner Weise gelungen}{}
\EdL{13}{AK07-J0092-L13}{ist. Wir können hier nicht näher auf diese Begründungen eintreten, sondern müssen}{}
\EdL{14}{AK07-J0092-L14}{den Leser auf die Arbeit \Sperr{Nallinos} verweisen\textsuperscript{2})\,; wir wollen hier nur kurz die Haupt-}{}
\EdL{15}{AK07-J0092-L15}{faktoren erwähnen, von denen die Breite der Planeten in ihren verschiedenen}{}
\EdL{16}{AK07-J0092-L16}{Stellungen abhängig ist· Die Breite der \textit{drei obern} Planeten hängt ab\,: erstens von}{}
\EdL{17}{AK07-J0092-L17}{der konstanten Neigung der Ebene des Deferens zur Ekliptik, und zweitens von}{}
\EdL{18}{AK07-J0092-L18}{der variabeln Neigung des Epizykels zur Ebene des Deferens. Die Breite der}{}
\EdL{19}{AK07-J0092-L19}{\textit{zwei untern} Planeten hängt ab\,: erstens von der Neigung der Ebene des Deferens}{}
\EdL{20}{AK07-J0092-L20}{zur Ekliptik, die \Sperr{Ptol.} für die untern Planeten als variabel annimmt, zweitens}{}
\EdL{21}{AK07-J0092-L21}{von der variabeln Neigung des Durchmessers Apogeum-Perigeum des Epizykels}{}
\EdL{22}{AK07-J0092-L22}{zur Ebene des Deferens, und drittens von der variabeln Neigung des auf dem}{}
\EdL{23}{AK07-J0092-L23}{genannten Durchmesser senkrecht stehenden zur Ebene der Ekliptik (dieser Durch-}{}
\EdL{24}{AK07-J0092-L24}{messer ist bei den drei obern Planeten stets zur Ebene der Ekliptik parallel). —}{}
\EdL{25}{AK07-J0092-L25}{Was wir einzig an den Darstellungen \Sperr{Nallinos} auszusetzen haben, ist, dass er die}{}
\EdL{26}{AK07-J0092-L26}{unklaren Figuren des \Sperr{Ptol.} nicht durch Beifügung perspektivischer Figuren anschau-}{}
\EdL{27}{AK07-J0092-L27}{licher zu machen versucht hat.}{}
\EdLp{28}{AK07-J0092-L28}{Vergleichen wir nun die Regeln, die \Sperr{Kh.} zur Bestimmung der Breite der}{}
\EdL{29}{AK07-J0092-L29}{Planeten gegeben hat, mit denen \Sperr{Batt.}’s, so sehen wir sofort, dass jene auf andere}{}
\EdL{30}{AK07-J0092-L30}{Theorien sich stützen müssen als diese\,; es ist auch wohl begreiflich, dass die}{}
\EdL{31}{AK07-J0092-L31}{Schwierigkeit der ptolemäischen Darstellung seine Nachfolger abgeschreckt hat, dem}{}
\EdL{32}{AK07-J0092-L32}{Wege des grossen Meisters zu folgen, und sie gezwungen hat, andere einfachere}{}
\EdL{33}{AK07-J0092-L33}{Methoden anzuwenden.}{}
\EdLp{34}{AK07-J0092-L34}{Sehen wir uns nun die Regel des \textit{Sūrya-Siddhānta} an\textsuperscript{3}), so finden wir auch}{}
\EdL{35}{AK07-J0092-L35}{hier keine direkte Übereinstimmung mit derjenigen \Sperr{Kh.}’s\,; aber dass beide von}{}
\EdL{36}{AK07-J0092-L36}{ähnlichen Gesichtspunkten aus vorgegangen sind, ist doch ersichtlich. Die Regel}{}
\EdL{37}{AK07-J0092-L37}{des \textit{Sūrya-Siddhānta} heisst\,: Der Sinus des Bogens (Differenz der wahren Länge des}{}
\EdL{38}{AK07-J0092-L38}{Planeten und der Länge des Knotens) multipliziert mit der maximalen Breite und}{}
\EdL{39}{AK07-J0092-L39}{geteilt durch den Radius (für diese Stellung des Planeten) gibt die Breite für die}{}
\NoteRule
\EdN{N1}{AK07-J0092-N01}{\textsuperscript{1}) \Sperr{Batt.}, ed. \Sperr{Nallino}, I, 115—116, II, 140—141.}
\EdN{N2}{AK07-J0092-N02}{\textsuperscript{2}) Ibid. II, 247—255.}
\EdN{N3}{AK07-J0092-N03}{\textsuperscript{3}) \textit{Transl.} Burgess, \textit{Journal of the American Orient. Soc.}, Vol. VI, p. 228.}
\FrameHead{0093}{p. 61}{0093}{p. 61}
\EdL{1}{AK07-J0093-L01}{betreffende Stellung (Zeit). \Sperr{Burgess} sagt im Kommentar\textsuperscript{1}) zu dieser Regel, dieselbe}{}
\EdL{2}{AK07-J0093-L02}{ergebe sich aus der Zusammensetzung zweier Proportionen, von denen die eine nur}{}
\EdL{3}{AK07-J0093-L03}{beim Monde Gültigkeit habe, die andere aber für die Planeten hinzukomme, nämlich\,:}{}
\EdL{4}{AK07-J0093-L04}{\parbox[t]{\dimexpr0.0891\linewidth-0.5em\relax}{\raggedright\small\mbox{}\textit{1. Prop.}}\hspace{0.5em}\parbox[t]{\dimexpr0.2673\linewidth-0.5em\relax}{\raggedright\small\mbox{}\textit{r} ∶ sin Dist. v. Knoten}\hspace{0.5em}\parbox[t]{\dimexpr0.0891\linewidth-0.5em\relax}{\raggedright\small\mbox{}=}\hspace{0.5em}\parbox[t]{\dimexpr0.5345\linewidth-0.5em\relax}{\raggedright\small\mbox{}grösste Breite ∶ Breite in der Dist. \textit{r}\textsuperscript{2})}\hspace{0.5em}}{}
\EdL{5}{AK07-J0093-L05}{\parbox[t]{\dimexpr0.0891\linewidth-0.5em\relax}{\raggedright\small\mbox{}\textit{2. Prop.}}\hspace{0.5em}\parbox[t]{\dimexpr0.2673\linewidth-0.5em\relax}{\raggedright\small\mbox{}var. \textit{r} ∶ \textit{r}}\hspace{0.5em}\parbox[t]{\dimexpr0.0891\linewidth-0.5em\relax}{\raggedright\small\mbox{}=}\hspace{0.5em}\parbox[t]{\dimexpr0.5345\linewidth-0.5em\relax}{\raggedright\small\mbox{}Breite in der Dist. \textit{r} ∶ Breite in der wahren Dist.}\hspace{0.5em}}{}
\EdLp{6}{AK07-J0093-L06}{Aus beiden Proportionen folgt durch Multiplikation\,:}{}
\EdLc{7}{AK07-J0093-L07}{var. \textit{r} ∶ sin. Dist. v. Knoten = grösste Breite ∶ Breite in der wahren Dist.}{}
\EdL{8}{AK07-J0093-L08}{also\,: Breite in der wahren Dist. = \Frac{sin Dist. v. Knoten × grösste Breite}{var. \textit{r}}.}{}
\EdLp{9}{AK07-J0093-L09}{Was die erste Proportion anbetrifft, so ergibt sich dieselbe aus der Regel der}{}
\EdL{10}{AK07-J0093-L10}{4 Grössen in der sphär. Trigonometrie\,: sin \textit{AB} ∶ sin \textit{AB′} = sin \textit{BC} ∶ sin \textit{B′C′}, die in}{}
\EdL{11}{AK07-J0093-L11}{die Formel\,:}{}
\EdLc{12}{AK07-J0093-L12}{sin \textit{α} = \Frac{sin \textit{B′C′}}{sin \textit{AB′}}}{}
\EdL{13}{AK07-J0093-L13}{übergeht, wenn \textit{AB} = 90° wird. Auf unsern Fall an-}{}
\EdL{14}{AK07-J0093-L14}{gewandt, ist \textit{A} einer der Knoten (Durchschnittspunkte des}{}
\EdL{15}{AK07-J0093-L15}{Deferens mit der Ekliptik), \textit{α} = Neigung des Deferens}{}
\EdL{16}{AK07-J0093-L16}{zur Ekliptik, \textit{BC} = maximale Breite, wenn \textit{AB} = 90°}{}
\EdL{17}{AK07-J0093-L17}{ist, \textit{B′C′} = Breite des Planeten bei seiner Stellung \textit{B′} im}{}
\EdL{18}{AK07-J0093-L18}{Deferens, \textit{AB′} = Differenz der wahren Länge des Planeten und der Länge des Knotens.}{}
\EdLp{19}{AK07-J0093-L19}{Die zweite Proportion ist jedenfalls nur angenähert richtig, wir können sie aus}{}
\EdL{20}{AK07-J0093-L20}{folgender Figur\,:}{}
\EdLp{21}{AK07-J0093-L21}{\textit{P} sei die Stellung des Planeten im Epizykel, \textit{A} seine Projektion auf die Ebene}{}
\EdL{22}{AK07-J0093-L22}{der Ekliptik, \textit{B} die Projektion des Epizykelzentrums auf dieselbe Ebene, so sind}{}
\EdL{23}{AK07-J0093-L23}{die Winkel \textit{HMB} und \textit{PEA} das, was im}{}
\EdL{24}{AK07-J0093-L24}{\textit{Sūrya-Siddhānta} genannt ist\,: Breite des}{}
\EdL{25}{AK07-J0093-L25}{Planeten in der Distanz \textit{r} (mittlere Di-}{}
\EdL{26}{AK07-J0093-L26}{stanz \textit{MH}), bezw. Breite des Planeten in}{}
\EdL{27}{AK07-J0093-L27}{der Distanz \textit{r′} (variable Distanz \textit{EP}).}{}
\EdL{28}{AK07-J0093-L28}{Wegen der geringen Neigung der Ebene}{}
\EdL{29}{AK07-J0093-L29}{des Epizykels zur Ekliptik nahmen}{}
\EdL{30}{AK07-J0093-L30}{die Inder die beiden Senkrechten \textit{PA}}{}
\EdL{31}{AK07-J0093-L31}{und \textit{HB} als gleich gross an, dann}{}
\EdL{32}{AK07-J0093-L32}{hat man\,:}{}
\EdLc{33}{AK07-J0093-L33}{sin \textit{HMB} = \Frac{\textit{HB}}{\textit{MH}} = \Frac{\textit{HB}}{\textit{r}}\,; sin \textit{PEA} = \Frac{\textit{PA}}{\textit{EP}} = \Frac{\textit{PA}}{\textit{r′}}}{}
\EdL{34}{AK07-J0093-L34}{also \Frac{sin \textit{HMB}}{sin \textit{PEA}} = \Frac{\textit{r′}}{\textit{r}},}{}
\NoteRule
\EdN{N1}{AK07-J0093-N01}{\textsuperscript{1}) Ibid. p. 229—230.}
\EdN{N2}{AK07-J0093-N02}{\textsuperscript{2}) Eigentlich sollte von den zwei letzten Gliedern der Proportion auch der Sinus genommen}
\EdN{N3}{AK07-J0093-N03}{werden, da diese Neigungen bei den Planetenbahnen aber nur gering sind, kann man ohne grossen}
\EdN{N4}{AK07-J0093-N04}{Fehler statt der Sinus die Bogen setzen.}
\FrameHead{0094}{p. 62}{0094}{p. 62}
\EdL{1}{AK07-J0094-L01}{oder wie oben bemerkt wurde, statt der Sinus die Bogen (Winkel) gesetzt\,:}{}
\EdLc{2}{AK07-J0094-L02}{\Frac{Breite in der mittleren Dist. \textit{r}}{Breite in der wahren Dist \textit{r′}} = \Frac{\textit{r′}}{\textit{r}},}{}
\EdL{3}{AK07-J0094-L03}{dies ist die 2. Proportion des \textit{Sūrya-Siddhānta.}}{}
\EdLp{4}{AK07-J0094-L04}{Es stimmen nun die Werte der Breite der Planeten nach \Sperr{Batt.} berechnet}{}
\EdL{5}{AK07-J0094-L05}{keineswegs mit den nach dem \textit{Sūrya-Siddhānta} berechneten, aber auch die letztern}{}
\EdL{6}{AK07-J0094-L06}{stimmen nicht mit den nach \Sperr{Kh.} berechneten. Wir haben oben schon bemerkt,}{}
\EdL{7}{AK07-J0094-L07}{dass die Regel \Sperr{Kh.}’s keine direkte Übereinstimmung mit der des \textit{Sūrya-Siddhānta}}{}
\EdL{8}{AK07-J0094-L08}{zeige, wir wollen nun nachzuweisen versuchen, wie weit sie übereinstimmen.}{}
\EdLp{9}{AK07-J0094-L09}{Wir haben gefunden, dass die Zahlen der letzten Kolonne in den Tafeln 27—56,}{}
\EdL{10}{AK07-J0094-L10}{die die Überschrift \textit{Latitudo secunda} trägt, erhalten werden nach der Formel\,:}{}
\EdLc{11}{AK07-J0094-L11}{sin \textit{lat. sec.} = sin \textit{argum.} × sin \textit{lat. maxim.},}{}
\EdL{12}{AK07-J0094-L12}{wo das \textit{argumentum} = wahre Länge des Planeten weniger Länge seines Knotens =}{}
\EdL{13}{AK07-J0094-L13}{Distanz des Planeten vom Knoten ist, und die \textit{lat. max.} (in der mittleren Distanz) =}{}
\EdL{14}{AK07-J0094-L14}{der \textit{lat. secunda} für 90° ist (bei Saturn = 5° nach \Sperr{Kh.})\,; es stimmt also diese \textit{lat.}}{}
\EdL{15}{AK07-J0094-L15}{\textit{secunda} mit dem Zähler des Wertes für die Breite in wahrer Distanz (s. p. 61)\textsuperscript{1}).}{}
\EdL{16}{AK07-J0094-L16}{Nun heisst es in der Regel \Sperr{Kh.}’s, man müsse die \textit{lat. secunda} durch die \textit{lat. prima}}{}
\EdL{17}{AK07-J0094-L17}{teilen, um die gesuchte Breite zu erhalten, aber nicht durch die \textit{lat. prima,} die dem}{}
\EdL{18}{AK07-J0094-L18}{vorigen Argumente (Distanz des Planeten vom Knoten) entspricht, sondern durch}{}
\EdL{19}{AK07-J0094-L19}{diejenige, die dem Argument \textit{Anomalia aequata} (\textit{elheza examinatum} (!) im Texte)}{}
\EdL{20}{AK07-J0094-L20}{entspricht\,; vergleicht man dies mit der im \textit{Sūrya-Siddhānta} für die Breite in der}{}
\EdL{21}{AK07-J0094-L21}{wahren Distanz gegebenen Formel, so müsste also diese \textit{lat. prima,} die der \textit{Anomalia}}{}
\EdL{22}{AK07-J0094-L22}{\textit{aequata} entspricht, gleich dem variabeln Radius \textit{r′} sein. Die Tafeln der Planeten}{}
\EdL{23}{AK07-J0094-L23}{(Nr. 27—56) enthalten aber dieses \textit{r′} nicht in absoluter Grösse, sondern geben sein}{}
\EdL{24}{AK07-J0094-L24}{Verhältnis zum mittleren Radius \textit{r} an, oder mit andern Worten\,: Die Zahlen der}{}
\EdL{25}{AK07-J0094-L25}{Kolonne \textit{lat. prima} repräsentieren die Werte für \textit{r′}, wenn \textit{r} = 1 angenommen wird.}{}
\EdLp{26}{AK07-J0094-L26}{Wie dieses \textit{r′} bestimmt worden ist, konnten wir nicht herausfinden, aber}{}
\EdL{27}{AK07-J0094-L27}{jedenfalls ist es unrichtig, wenn die Kolonne der \textit{lat. prima} mit Graden und}{}
\EdL{28}{AK07-J0094-L28}{Minuten überschrieben ist, \textit{r′} oder vielmehr \Frac{\textit{r′}}{\textit{r}} ist eine unbenannte Grösse, und es}{}
\EdL{29}{AK07-J0094-L29}{sollte daher statt \textit{Grade} und \textit{Minuten} heissen \textit{Partes} und \textit{Minutae,} dann erhält man}{}
\EdL{30}{AK07-J0094-L30}{auch für den Quotienten von \textit{lat. sec.} durch \textit{lat. prima} wieder Grade und Minuten.}{}
\EdLp{31}{AK07-J0094-L31}{Wir wollen nun genau den Regeln \Sperr{Kh.}’s entsprechend ein \textit{Beispiel} geben. Es}{}
\EdL{32}{AK07-J0094-L32}{sei die \textit{Breite} des \textit{Saturns} für den Anfang des 15. Ṣafar d. J. 540 d. H. zu bestimmen\,:}{}
\EdLp{33}{AK07-J0094-L33}{Die \textit{Anomalia aequata} (\textit{elheza examinata}) für diesen Zeitpunkt ist = 109° 24\textsuperscript{1}/\textsubscript{2}′}{}
\EdL{34}{AK07-J0094-L34}{(s. p. 53), dazu gehört nach Tafel 30 die latitudo prima 1° 15′, oder vielmehr 1\textsuperscript{\textit{p}} 15′}{}
\EdL{35}{AK07-J0094-L35}{= 1\textsuperscript{1}/\textsubscript{4}\,; der wahre Ort des Planeten für diesen Zeitpunkt ist = 27° 12′ (s. p. 53), die}{}
\EdL{36}{AK07-J0094-L36}{Länge des Knotens = 103° 12′ (s. p. 64), diese beiden addiert, gibt 130° 24′\,; diesem}{}
\EdL{37}{AK07-J0094-L37}{Argument entspricht in Tafel 31 die latitudo secunda 3° 48″ 27′, nun haben wir\,:}{}
\EdLc{38}{AK07-J0094-L38}{3° 48\textsuperscript{1}/\textsubscript{2}′ ∶ 1\textsuperscript{\textit{p}} 15′ = 228\textsuperscript{1}/\textsubscript{2}′ ∶ 1\textsuperscript{1}/\textsubscript{4} = 183′ = 3° 3′ = \textit{gesuchte Breite.}}{}
\NoteRule
\EdN{N1}{AK07-J0094-N01}{\textsuperscript{1}) Wenn man statt „grösste Breite“ „\textit{sin} grösste Breite“ setzt.}
\FrameHead{0095}{p. 63}{0095}{p. 63}
\EdLp{1}{AK07-J0095-L01}{Die Berechnung nach den Regeln \Sperr{Batt.}’s ergibt den Wert 2° 34′\,; die Berechnung}{}
\EdL{2}{AK07-J0095-L02}{nach den Daten \Sperr{Batt.}’s, aber nach den Regeln \Sperr{Kh.}’s ergibt 2° 28′.}{}
\EdLp{3}{AK07-J0095-L03}{Ob die Breite \textit{nördlich} oder \textit{südlich} sei, darüber sagt unser Text gar nichts,}{}
\EdL{4}{AK07-J0095-L04}{während sich \Sperr{Batt.} (I, 115—116) deutlich darüber ausspricht\,; nach seinen Regeln}{}
\EdL{5}{AK07-J0095-L05}{ist sie in unserm Beispiel für Saturn \textit{südlich}\,; ferner ist der Planet \textit{descendens,} d. h.}{}
\EdL{6}{AK07-J0095-L06}{von der Ekliptik aus gegen Süden gehend, was man auch in unserer Tafel 31 direkt}{}
\EdL{7}{AK07-J0095-L07}{ablesen kann\,: die Kolonne, der wir die \textit{lat. sec.} entnommen haben, ist überschrieben}{}
\EdL{8}{AK07-J0095-L08}{\textit{latitudo Saturni recedentis.} — Nach der Regel unseres Textes ist das Argument für}{}
\EdL{9}{AK07-J0095-L09}{die \textit{lat. sec.,} wie es scheint, vom aufsteigenden Knoten aus über das Perigeum des}{}
\EdL{10}{AK07-J0095-L10}{Deferens des Planeten gezählt, von demselben Knoten aus über das Apogeum gezählt}{}
\EdL{11}{AK07-J0095-L11}{wäre es = 360° — 130° 24′ = 229° 36′, aus welchem Werte leichter zu erkennen}{}
\EdL{12}{AK07-J0095-L12}{wäre, dass die Breite südlich ist. — Die hier gegebene Regel für die Bestimmung}{}
\EdL{13}{AK07-J0095-L13}{der \textit{lat. sec.} ist übrigens dieselbe, wie die zur Bestimmung der Breite des Mondes}{}
\EdL{14}{AK07-J0095-L14}{gegebene (s. oben p. 58—59).}{}
\EdLp{15}{AK07-J0095-L15}{Vergleicht man unsere Regel mit denjenigen bei \Sperr{Ptol.} und \Sperr{Batt.}, die ziemlich}{}
\EdL{16}{AK07-J0095-L16}{übereinstimmend sind, so sieht man, dass sie sich in Bezug auf die Argumente für}{}
\EdL{17}{AK07-J0095-L17}{die Breite und die Verhältniszahlen gerade entgegengesetzt verhalten\,: \Sperr{Ptol.} bestimmt}{}
\EdL{18}{AK07-J0095-L18}{die Breitenzahl mit Hilfe der \textit{Anomalia aequata} als Argument, die Verhältniszahl}{}
\EdL{19}{AK07-J0095-L19}{(mit der man die gefundene Breite multiplizieren muss) mit Hilfe der um 50° (beim}{}
\EdL{20}{AK07-J0095-L20}{Saturn) vermehrten wahren Länge des Planeten\textsuperscript{1}) als Argument\,; \Sperr{Kh.} dagegen bestimmt}{}
\EdL{21}{AK07-J0095-L21}{die Breitenzahl mit Hilfe des Abstandes des Planeten vom aufsteigenden Knoten}{}
\EdL{22}{AK07-J0095-L22}{(d. h. indirekt mit Hilfe der wahren Länge), die Verhältniszahl (\textit{latit. prima}) aber mit}{}
\EdL{23}{AK07-J0095-L23}{Hilfe der \textit{Anomalia aequata} als Argument. Diese verschiedenen Auffassungen bedingen}{}
\EdL{24}{AK07-J0095-L24}{dann natürlich auch eine ganz abweichende Einrichtung der Tafeln, wir wollen}{}
\EdL{25}{AK07-J0095-L25}{nur einen Hauptunterschied anführen\,: \Sperr{Ptol.} und \Sperr{Batt.} haben bei den obern Planeten}{}
\EdL{26}{AK07-J0095-L26}{beim Argument 180° das Maximum der Breite, bei 0° (360°) das Minimum\,; \Sperr{Kh.} dagegen}{}
\EdL{27}{AK07-J0095-L27}{hat bei 90° (270°) das Maximum der \textit{Latitudo secunda,} bei 0° und 180° das Minimum.}{}
\EdLp{28}{AK07-J0095-L28}{Für die \textit{untern} Planeten ist die Theorie der Bestimmung der \textit{Breite} und die}{}
\EdL{29}{AK07-J0095-L29}{Berechnung mit Hilfe der Tafeln bei \Sperr{Ptol.} und \Sperr{Batt.} noch bedeutend komplizierter}{}
\EdL{30}{AK07-J0095-L30}{als für die obern\,; man erlasse mir auf dieselbe hier einzutreten, ich verweise den}{}
\EdL{31}{AK07-J0095-L31}{Leser für das nähere auf die \Sperr{Nallino}’sche Ausgabe \Sperr{Batt.}’s\textsuperscript{2}).}{}
\EdLp{32}{AK07-J0095-L32}{Als \textit{Beispiel} geben wir die Berechnung der \textit{Breite} der \textit{Venus} für den Beginn}{}
\EdL{33}{AK07-J0095-L33}{des 15. Ṣafar 540 d. H. nach den Regeln und Tafeln \Sperr{Kh.}’s.}{}
\EdLp{34}{AK07-J0095-L34}{Die \textit{Anomalia} (\textit{argumentum} im Text) \textit{aequata} der Venus für diesen Zeitpunkt}{}
\EdL{35}{AK07-J0095-L35}{ist nach den Tafeln 15—17 berechnet = 118° 15′ (s. oben p. 54),}{}
\EdL{36}{AK07-J0095-L36}{dazu gehört in Tafel 48 die \textit{latit. prima} = 1\textsuperscript{\textit{p}} 15′.}{}
\EdL{37}{AK07-J0095-L37}{\parbox[t]{\dimexpr0.3267\linewidth-0.5em\relax}{\raggedright\small\mbox{}}\hspace{0.5em}\parbox[t]{\dimexpr0.3267\linewidth-0.5em\relax}{\raggedright\small\mbox{}Anomalia aequata}\hspace{0.5em}\parbox[t]{\dimexpr0.1633\linewidth-0.5em\relax}{\raggedright\small\mbox{}=}\hspace{0.5em}\parbox[t]{\dimexpr0.1633\linewidth-0.5em\relax}{\raggedright\small\mbox{}118° 15′}\hspace{0.5em}}{}
\EdL{38}{AK07-J0095-L38}{\parbox[t]{\dimexpr0.3267\linewidth-0.5em\relax}{\raggedright\small\mbox{}}\hspace{0.5em}\parbox[t]{\dimexpr0.3267\linewidth-0.5em\relax}{\raggedright\small\mbox{}Länge des Knotens}\hspace{0.5em}\parbox[t]{\dimexpr0.1633\linewidth-0.5em\relax}{\raggedright\small\mbox{}=}\hspace{0.5em}\parbox[t]{\dimexpr0.1633\linewidth-0.5em\relax}{\raggedright\small\mbox{}59° 27′}\hspace{0.5em}}{}
\EdL{39}{AK07-J0095-L39}{\parbox[t]{\dimexpr0.6533\linewidth-0.5em\relax}{\raggedright\small\mbox{}Mittlere Länge der Venus gleich der der Sonne}\hspace{0.5em}\parbox[t]{\dimexpr0.1633\linewidth-0.5em\relax}{\raggedright\small\mbox{}=}\hspace{0.5em}\parbox[t]{\dimexpr0.1633\linewidth-0.5em\relax}{\raggedright\small\mbox{}131° 1′}\hspace{0.5em}}{}
\EdL{40}{AK07-J0095-L40}{\parbox[t]{\dimexpr0.3267\linewidth-0.5em\relax}{\raggedright\small\mbox{}}\hspace{0.5em}\parbox[t]{\dimexpr0.4900\linewidth-0.5em\relax}{\raggedright\small\mbox{}Summa\,:}\hspace{0.5em}\parbox[t]{\dimexpr0.1633\linewidth-0.5em\relax}{\raggedright\small\mbox{}308° 43′}\hspace{0.5em}}{}
\EdL{41}{AK07-J0095-L41}{dazu gehört in Tafel 46 die \textit{latitud. sec.} 3° 54′.}{}
\NoteRule
\EdN{N1}{AK07-J0095-N01}{\textsuperscript{1}) \Sperr{Batt.} nimmt statt der wahren Länge das Centrum.}
\EdN{N2}{AK07-J0095-N02}{\textsuperscript{2}) I, 115—116, und II, 249—253.}
\FrameHead{0096}{p. 64}{0096}{p. 64}
\EdLp{1}{AK07-J0096-L01}{3° 54′ ∶ 1\textsuperscript{\textit{p}} 15′ = 234′ ∶ 1,25 = 187′ = 3° 7′ = \textit{Breite} der \textit{Venus} im gegebenen}{}
\EdL{2}{AK07-J0096-L02}{Zeitpunkt. Nach den Tafeln \Sperr{Batt.}’s findet man die Breite = 57′\,; dieser grosse}{}
\EdL{3}{AK07-J0096-L03}{Unterschied lässt vermuten, dass die Regel \Sperr{Kh.}’s nicht richtig sei\,; in der Tat ist es}{}
\EdL{4}{AK07-J0096-L04}{uns unklar, warum die \textit{Anomalia aequata} auch noch zur Summe der mittlern Länge}{}
\EdL{5}{AK07-J0096-L05}{und der Länge des Knotens addiert werden muss, lässt man jene weg, so ergibt}{}
\EdL{6}{AK07-J0096-L06}{sich als Argument für die lat. sec. 190° 28′ und dazu gehört in Tafel 50 die lat. 54\textsuperscript{1}/\textsubscript{2}′,}{}
\EdL{7}{AK07-J0096-L07}{also ist die \textit{Breite} =}{}
\EdLc{8}{AK07-J0096-L08}{54,5′ ∶ 1,25 = c. 44′.}{}
\EdLp{9}{AK07-J0096-L09}{Auch hier sagt unser Text über die Art der Breite nichts\,; nach \textit{Batt.’s} Regeln}{}
\EdL{10}{AK07-J0096-L10}{ist sie \textit{nördlich.}}{}
\EdLp{11}{AK07-J0096-L11}{Hiermit verlassen wir dieses schwierige Kapitel über die Breite der Planeten.}{}
\EdL{12}{AK07-J0096-L12}{In unserm Falle wurde diese Frage noch mehr erschwert durch den Umstand, dass}{}
\EdL{13}{AK07-J0096-L13}{auch die indischen astronomischen Werke in diesem Punkte nicht einig sind\,: Die}{}
\EdL{14}{AK07-J0096-L14}{Regel im \textit{Sūrya-Siddhānta,} die wir oben zitiert haben, stimmt nicht mit derjenigen}{}
\EdL{15}{AK07-J0096-L15}{in der \textit{Panchasiddhāntikā}\textsuperscript{1})\,; hier wird für Saturn und Merkur der Sinus der Distanz}{}
\EdL{16}{AK07-J0096-L16}{vom Knoten mit der Konstanten \Frac{9}{8} multipliziert (wenn der Kommentator recht hat),}{}
\EdL{17}{AK07-J0096-L17}{während er im \textit{Sūrya-Siddhānta} mit der variabeln Grösse \Frac{grösste Breite}{\textit{r′}} multipliziert wird.}{}
\EdLc{18}{AK07-J0096-L18}{\textbf{Kap. 18. Über Apogeum und Perigeum der Planeten.}}{}
\EdLc{19}{AK07-J0096-L19}{(Text p. 14—15.)}{}
\EdLp{20}{AK07-J0096-L20}{Diese technischen Ausdrücke werden vom Übersetzer durch \textit{sublimatio} (O und C}{}
\EdL{21}{AK07-J0096-L21}{haben bisweilen \textit{sullimatio}) und \textit{submissio} wiedergegeben. Wir haben oben schon}{}
\EdL{22}{AK07-J0096-L22}{(p. 46) die Apogeen der Planeten und der Sonne nach \Sperr{Kh.} angegeben, und auch über}{}
\EdL{23}{AK07-J0096-L23}{ihre Abweichungen von den Werten \Sperr{Batt.}’s und des \textit{Sūrya-Siddhānta} gesprochen\,;}{}
\EdL{24}{AK07-J0096-L24}{wir führen hier nur noch an, dass in Bezug auf den Mond die Bemerkung hinzu-}{}
\EdL{25}{AK07-J0096-L25}{gefügt wird, dass sein Apogeum beweglich sei und zwar betrage diese Bewegung}{}
\EdL{26}{AK07-J0096-L26}{im Tage 6′ 40″ 18‴, in der Stunde 16″ 42‴ 27\textsuperscript{\textit{IV}}\,; diese Bewegung, die die Ursache}{}
\EdL{27}{AK07-J0096-L27}{des anomalistischen Monats ist, finden wir bei \Sperr{Batt.} nicht erwähnt\,; nach den neuern}{}
\EdL{28}{AK07-J0096-L28}{Bestimmungen braucht die Apsidenlinie zu einer vollen Umdrehung 3231,466\textsuperscript{d}}{}
\EdL{29}{AK07-J0096-L29}{(= c. 9\textsuperscript{a}), dies macht im Tage eine Bewegung von 6′ 41″\,; vergleicht man in den}{}
\EdL{30}{AK07-J0096-L30}{obigen Zahlen die stündliche Bewegung mit der täglichen, so müsste letztere, wenn}{}
\EdL{31}{AK07-J0096-L31}{die erstere richtig ist, 6′ 40″ 58‴ sein, es soll also wahrscheinlich statt 18‴ heissen}{}
\EdL{32}{AK07-J0096-L32}{58‴, welche Verwechslung in Zahlen mit arabischen Buchstaben ausgedrückt leicht}{}
\EdL{33}{AK07-J0096-L33}{möglich ist\,; dann stimmen die Angaben \Sperr{Kh.}’s ausgezeichnet mit den modernen Werten.}{}
\EdLc{34}{AK07-J0096-L34}{\textbf{Kap. 19. Über die Knotenpunkte der Planeten.}}{}
\EdLc{35}{AK07-J0096-L35}{(Text p. 15.)}{}
\EdLp{36}{AK07-J0096-L36}{Die Längen des Drachenkopfes für die einzelnen Planeten sind\,:}{}
\EdL{37}{AK07-J0096-L37}{\parbox[t]{\dimexpr0.1960\linewidth-0.5em\relax}{\raggedright\small\mbox{}Saturn}\hspace{0.5em}\parbox[t]{\dimexpr0.1960\linewidth-0.5em\relax}{\raggedright\small\mbox{}Jupiter}\hspace{0.5em}\parbox[t]{\dimexpr0.1960\linewidth-0.5em\relax}{\raggedright\small\mbox{}Mars}\hspace{0.5em}\parbox[t]{\dimexpr0.1960\linewidth-0.5em\relax}{\raggedright\small\mbox{}Venus}\hspace{0.5em}\parbox[t]{\dimexpr0.1960\linewidth-0.5em\relax}{\raggedright\small\mbox{}Merkur}\hspace{0.5em}}{}
\EdL{38}{AK07-J0096-L38}{\parbox[t]{\dimexpr0.1960\linewidth-0.5em\relax}{\raggedright\small\mbox{}103° 12′}\hspace{0.5em}\parbox[t]{\dimexpr0.1960\linewidth-0.5em\relax}{\raggedright\small\mbox{}82° 1′}\hspace{0.5em}\parbox[t]{\dimexpr0.1960\linewidth-0.5em\relax}{\raggedright\small\mbox{}21° 54′ (44′)}\hspace{0.5em}\parbox[t]{\dimexpr0.1960\linewidth-0.5em\relax}{\raggedright\small\mbox{}59° 27′}\hspace{0.5em}\parbox[t]{\dimexpr0.1960\linewidth-0.5em\relax}{\raggedright\small\mbox{}21° 10.}\hspace{0.5em}}{}
\NoteRule
\EdN{N1}{AK07-J0096-N01}{\textsuperscript{1}) Ed. G. \Sperr{Thibaut}, p. 95 und 96 der engl. Übersetzung.}
\FrameHead{0097}{p. 65}{0097}{p. 65}
\EdLp{1}{AK07-J0097-L01}{Für den Mond bemerkt er wieder, dass seine Knotenpunkte beweglich seien,}{}
\EdL{2}{AK07-J0097-L02}{und zwar betrage die Bewegung derselben täglich 3′ 11″ 48‴, stündlich 7″ 57‴\,;}{}
\EdL{3}{AK07-J0097-L03}{multipliziert man die letztere Zahl mit 24, so erhält man 3′ 10″ 48‴, so muss es}{}
\EdL{4}{AK07-J0097-L04}{also heissen statt 3′ 11″ 48‴, dann stimmt es auch mit der Tafel 20, die für die}{}
\EdL{5}{AK07-J0097-L05}{tägliche Bewegung des Drachenkopfes 3′ 11″, für die stündliche 8″ hat, wie auch}{}
\EdL{6}{AK07-J0097-L06}{\Sperr{Batt.} Dass diese Bewegung der Knoten derjenigen der Apsiden entgegengesetzt sei,}{}
\EdL{7}{AK07-J0097-L07}{wird von \Sperr{Kh.} nicht erwähnt. Nach den neuern Bestimmungen braucht die Knoten-}{}
\EdL{8}{AK07-J0097-L08}{linie zu einer vollen Umdrehung 6798\textsuperscript{d},3355, dies macht im Tag eine Bewegung}{}
\EdL{9}{AK07-J0097-L09}{von 3′ 10″ 38‴, also wiederum eine sehr schöne Übereinstimmung. Die Tafeln}{}
\EdL{10}{AK07-J0097-L10}{\Sperr{Batt.}’s\textsuperscript{1}) stimmen mit denen \Sperr{Kh.}’s nur bis zu 4 Tagen, von hier an beginnen die}{}
\EdL{11}{AK07-J0097-L11}{Abweichungen, die in einem Jahr auf 1′ 14″, in 30 Jahren auf 34′ 46″ anwachsen.}{}
\EdLc{12}{AK07-J0097-L12}{\textbf{Kap. 20 und 21. Über die tägliche Bewegung des Mondes und der Sonne.}}{}
\EdLc{13}{AK07-J0097-L13}{(Text p. 15—16.)}{}
\EdLp{14}{AK07-J0097-L14}{Diese beiden Kapitel stehen hier nicht am richtigen Orte, sie sollten dem}{}
\EdL{15}{AK07-J0097-L15}{Kapitel 7 „Über die mittlere Bewegung der Planeten“ unmittelbar folgen, oder besser}{}
\EdL{16}{AK07-J0097-L16}{mit ihm vereinigt sein.}{}
\EdLp{17}{AK07-J0097-L17}{Die mittlere tägliche Bewegung des Mondes wird angegeben zu 13° 10′ 34″}{}
\EdL{18}{AK07-J0097-L18}{52‴ 48\textsuperscript{\textit{IV}}, die stündliche zu 32′ 56″ 47‴ 52\textsuperscript{\textit{IV}}\,; die letzten zwei Zahlen sind unrichtig,}{}
\EdL{19}{AK07-J0097-L19}{es muss heissen 27‴ 12\textsuperscript{\textit{IV}}.}{}
\EdLp{20}{AK07-J0097-L20}{Die mittlere tägliche Bewegung der Sonne wird angegeben zu 59′ 8″, die stünd-}{}
\EdL{21}{AK07-J0097-L21}{liche zu 2′ 27″ 50‴ 25\textsuperscript{\textit{IV}}\,; die Division von 59′ 8″ durch 24 ergibt aber nur 2′ 27″ 50‴,}{}
\EdL{22}{AK07-J0097-L22}{also sind wahrscheinlich bei 59′ 8″ zehn Tertien ausgefallen oder weggelassen}{}
\EdL{23}{AK07-J0097-L23}{worden.}{}
\EdLp{24}{AK07-J0097-L24}{Legen wir für die Bewegung der Sonne die obige mittlere stündliche Bewegung}{}
\EdL{25}{AK07-J0097-L25}{derselben von 2′ 27″ 50‴ 25\textsuperscript{\textit{IV}} zu Grunde, so erhalten wir für die mittlere Bewegung}{}
\EdL{26}{AK07-J0097-L26}{in 365 Tagen 359° 44′ 40″ 50‴. Dieser mittlern Bewegung entspricht eine Jahreslänge}{}
\EdL{27}{AK07-J0097-L27}{von 365\textsuperscript{d} 6\textsuperscript{h} 13\textsuperscript{m} 2\textsuperscript{s}, ungefähr dieselbe Grösse, die wir oben (p. 42) gefunden haben,}{}
\EdL{28}{AK07-J0097-L28}{und die uns also wieder zeigt, dass \Sperr{Kh.} für seine Sonnen- und Planetentafeln das}{}
\EdL{29}{AK07-J0097-L29}{siderische Jahr als Basis genommen hat. Diese Jahreslänge entspricht aber nicht ganz}{}
\EdL{30}{AK07-J0097-L30}{den Werten, die man in den indischen \textit{Siddhāntas} findet\textsuperscript{2}), sie ist um c. 25\textsuperscript{s} grösser}{}
\EdL{31}{AK07-J0097-L31}{als sie vom \textit{Sūrya-Siddhānta} und dem spätern \textit{Ārya-Siddhānta,} und um 53\textsuperscript{s} grösser}{}
\EdL{32}{AK07-J0097-L32}{als sie vom \textit{Brahma-Siddhānta} angegeben wird. Gehen wir aber umgekehrt von}{}
\EdL{33}{AK07-J0097-L33}{der Jahreslänge dieser \textit{Siddhāntas} aus, so kommen wir beim \textit{Sūrya-Siddhānta} (365\textsuperscript{d}}{}
\EdL{34}{AK07-J0097-L34}{6\textsuperscript{h} 12\textsuperscript{m} 36,5\textsuperscript{s}) auf eine mittlere tägliche Bewegung von 59′ 8″ 10‴ 10\textsuperscript{\textit{IV}}, beim \textit{Brahma-}}{}
\EdL{35}{AK07-J0097-L35}{\textit{Siddhānta} (365\textsuperscript{d} 6\textsuperscript{h} 12\textsuperscript{m} 9\textsuperscript{s}) auf eine solche von 59′ 8″ 10‴ 21\textsuperscript{\textit{IV}} 22\textsuperscript{\textit{V}}\,; dem erstern Werte}{}
\EdL{36}{AK07-J0097-L36}{entspricht eine stündliche Bewegung von 2′ 27″ 50‴ 25,4\textsuperscript{\textit{IV}}, dem zweiten eine solche}{}
\EdL{37}{AK07-J0097-L37}{von 2′ 27″ 50‴ 25,8\textsuperscript{\textit{IV}}\,; da \Sperr{Kh.} die stündliche Bewegung zu 2′ 27″ 50‴ 25\textsuperscript{\textit{IV}} angibt, so}{}
\EdL{38}{AK07-J0097-L38}{ist die Wahrscheinlichkeit etwas grösser, dass er dem \textit{Sūrya-Siddhānta} und dem}{}
\NoteRule
\EdN{N1}{AK07-J0097-N01}{\textsuperscript{1}) Ed. \Sperr{Nallino}, II, 19—22.}
\EdN{N2}{AK07-J0097-N02}{\textsuperscript{2}) \Sperr{Batt.} ed. \Sperr{Nallino}, I, 206.}
\FrameHead{0098}{p. 66}{0098}{p. 66}
\EdL{1}{AK07-J0098-L01}{spätern \textit{Ārya-Siddhānta} gefolgt sei, doch kann diese Frage aus den gegebenen Zahlen}{}
\EdL{2}{AK07-J0098-L02}{keineswegs endgiltig entschieden werden\textsuperscript{1}).}{}
\EdLp{3}{AK07-J0098-L03}{In einer Note, die nur das Ms. C enthält, werden die Daten des \Sperr{Ptol.} für}{}
\EdL{4}{AK07-J0098-L04}{die Jahreslänge, die monatliche, tägliche und stündliche mittlere Bewegung der}{}
\EdL{5}{AK07-J0098-L05}{Sonne richtig angegeben, ebenso die Bahn der Sonne in einem Jahre (von 365 ganzen}{}
\EdL{6}{AK07-J0098-L06}{Tagen) und in 18 Jahren\textsuperscript{2}). Was die Jahreslänge anbetrifft, so ist sie zu 365\textsuperscript{d} 14′ 48″}{}
\EdL{7}{AK07-J0098-L07}{angegeben, d. h. zu 365 + \Frac{14}{60} + \Frac{48}{3600} Tagen, dies macht in Stunden, Minuten und}{}
\EdL{8}{AK07-J0098-L08}{Sekunden verwandelt = 365\textsuperscript{d} 5\textsuperscript{h} 55\textsuperscript{m} 12\textsuperscript{s}, d. i. die Länge des tropischen Jahres des}{}
\EdL{9}{AK07-J0098-L09}{\Sperr{Ptol.} — Eigentümlich ist es, dass im ganzen Werke \Sperr{Kh.}’s keine Angabe über die}{}
\EdL{10}{AK07-J0098-L10}{Länge des zu Grunde gelegten siderischen Jahres zu finden ist.}{}
\EdLc{11}{AK07-J0098-L11}{\textbf{Kap. 22. Über das Erscheinen des Mondes am Abend des 29. Tages (des Monats).}}{}
\EdLc{12}{AK07-J0098-L12}{(Text p. 16—17, Tafel 57\textsuperscript{a}.)}{}
\EdLp{13}{AK07-J0098-L13}{Dieses Kapitel gehört zu den schwierigsten der arabischen Astronomie. \Sperr{Al-Bīrūnī}}{}
\EdL{14}{AK07-J0098-L14}{sagt hierüber in seiner \textit{Chronologie}\textsuperscript{3})\,: „Was das Erscheinen des Neumondes anbetrifft,}{}
\EdL{15}{AK07-J0098-L15}{so ist seine genaue Bestimmung weitläufig und schwierig (zu erklären) und bedarf}{}
\EdL{16}{AK07-J0098-L16}{komplizierter Rechnungen und vieler Tafeln\,; es möge genügen auf das hinzuweisen,}{}
\EdL{17}{AK07-J0098-L17}{was über diesen Gegenstand in den Tafeln \Sperr{Al-Battānīs} und \Sperr{Ḥabash}’ des Rechners}{}
\EdL{18}{AK07-J0098-L18}{gesagt ist“.}{}
\EdLp{19}{AK07-J0098-L19}{\Sperr{Batt.} behandelt diese Frage im 41. Kapitel\textsuperscript{4}), und G. \Sperr{Schiaparelli} und \Sperr{Nal-}}{}
\EdL{20}{AK07-J0098-L20}{\Sperr{lino} haben sie in ausgezeichneter Weise kommentiert\textsuperscript{5}). Ich kann hier nicht näher}{}
\EdL{21}{AK07-J0098-L21}{auf diesen Gegenstand eintreten, es möge genügen zu bemerken, dass die Regel \Sperr{Kh.}’s}{}
\EdL{22}{AK07-J0098-L22}{und seine Tafel wahrscheinlich auf anderer Grundlage beruhen müssen als die}{}
\EdL{23}{AK07-J0098-L23}{Darstellung \Sperr{Batt.}’s. Ich hebe nur einen Punkt hervor, der diese Behauptung zu}{}
\EdL{24}{AK07-J0098-L24}{stützen vermag. Nach \Sperr{Batt.} hängt das erste Erscheinen der Mondsichel nicht nur}{}
\EdL{25}{AK07-J0098-L25}{von dem Unterschied der wahren Längen von Sonne und Mond ab, sondern auch}{}
\EdL{26}{AK07-J0098-L26}{noch von der Breite des Mondes\,; er gibt nun}{}
\EdL{27}{AK07-J0098-L27}{(p. 87) die Vorschrift, man müsse das Quadrat}{}
\EdL{28}{AK07-J0098-L28}{des Unterschiedes der wahren Längen von Sonne}{}
\EdL{29}{AK07-J0098-L29}{und Mond zum Quadrat der Breite des Mondes}{}
\EdL{30}{AK07-J0098-L30}{addieren, so erhalte man angenähert das Quadrat}{}
\EdL{31}{AK07-J0098-L31}{der Distanz des Mondes von der Sonne\,; die Rich-}{}
\EdL{32}{AK07-J0098-L32}{tigkeit dieser Vorschrift ergibt sich ohne weiters aus beistehender Figur\,:}{}
\EdLc{33}{AK07-J0098-L33}{\textit{MS} \ensuremath{\backsim} \Root{\textit{CS}\textsuperscript{2} + \textit{MC}\textsuperscript{2}}.}{}
\EdLp{34}{AK07-J0098-L34}{\Sperr{Kh.} aber, der die Breite des Mondes allerdings auch berücksichtigt, sagt}{}
\EdL{35}{AK07-J0098-L35}{(p. 16)\,: „Wenn die Breite des Mondes nördlich ist, so ist sie zu der Differenz der}{}
\EdL{36}{AK07-J0098-L36}{wahren Längen von Sonne und Mond zu addieren, wenn sie südlich ist, von jener}{}
\NoteRule
\EdN{N1}{AK07-J0098-N01}{\textsuperscript{1}) Vergl. auch Kap. 38.}
\EdN{N2}{AK07-J0098-N02}{\textsuperscript{2}) Diese Angabe ist in den Graden fehlerhaft, es muss statt 1350 heissen 355.}
\EdN{N3}{AK07-J0098-N03}{\textsuperscript{3}) Ed. \Sperr{Sachau}, arab. Text p. 196, engl. Übers. p. 176.}
\EdN{N4}{AK07-J0098-N04}{\textsuperscript{4}) Ed. \Sperr{Nallino} I, 85—92.}
\EdN{N5}{AK07-J0098-N05}{\textsuperscript{5}) Ibid. I, 265–272.}
\FrameHead{0099}{p. 67}{0099}{p. 67}
\EdL{1}{AK07-J0099-L01}{Differenz zu subtrahieren.“ — Wie die Tafel 57\textsuperscript{a} zu gebrauchen ist, um die Zeit}{}
\EdL{2}{AK07-J0099-L02}{der ersten Sichtbarkeit der Mondsichel zu bestimmen, wird am besten aus einem}{}
\EdL{3}{AK07-J0099-L03}{Beispiel klar werden.}{}
\EdLp{4}{AK07-J0099-L04}{\textit{Beispiel\,:} Es ist zu untersuchen, ob am 29. \textit{Shaʿbān}\textsuperscript{1}) des Jahres 540 d. H. Abends}{}
\EdL{5}{AK07-J0099-L05}{die Mondsichel sichtbar wurde.}{}
\ColHead{left column}
\EdL{6}{AK07-J0099-L06}{\textit{Mittlere Länge der Sonne\,:}}{}
\EdL{7}{AK07-J0099-L07}{Am Ende des 510. Jahres = 38° 35′ 45″}{}
\EdL{8}{AK07-J0099-L08}{Zunahme für 29 Jahre = 49° 2′ 50″}{}
\EdL{9}{AK07-J0099-L09}{„ „ 7 Mon. = 204° 1′ 11″}{}
\EdL{10}{AK07-J0099-L10}{„ „ 29 Tage = 28° 34′ 57″}{}
\EdL{11}{AK07-J0099-L11}{„ „ 6 Std. = 14′ 47″}{}
\EdL{12}{AK07-J0099-L12}{320° 29′ 30″}{}
\EdL{13}{AK07-J0099-L13}{Apogeum der Sonne\,: 77° 55′}{}
\EdL{14}{AK07-J0099-L14}{\textit{Anomalie\,:} 242° 34′ 30″}{}
\EdL{15}{AK07-J0099-L15}{Hiezu findet man in Tafel 24 die Glei-}{}
\EdL{16}{AK07-J0099-L16}{chung 1° 57′ 35″, diese zu der mittlern}{}
\EdL{17}{AK07-J0099-L17}{Länge addiert, gibt für die \textit{wahre Länge}}{}
\EdL{18}{AK07-J0099-L18}{der Sonne rund 322° 27′.}{}
\ColHead{right column}
\EdL{19}{AK07-J0099-L19}{\textit{Mittlere Länge des Mondes\,:}}{}
\EdL{20}{AK07-J0099-L20}{Am Ende des 510. Jahres = 40° 56′ 38″}{}
\EdL{21}{AK07-J0099-L21}{Zunahme für 29 Jahre = 53° 24′ 16″}{}
\EdL{22}{AK07-J0099-L22}{„ „ 7 Mon. = 207° 30′ 20″}{}
\EdL{23}{AK07-J0099-L23}{„ „ 29 Tage = 22° 6′ 52″}{}
\EdL{24}{AK07-J0099-L24}{„ „ 6 Std. = 3° 17′ 39″}{}
\EdL{25}{AK07-J0099-L25}{327° 15′ 45″}{}
\EdL{26}{AK07-J0099-L26}{\textit{Anomalie des Mondes\,:}}{}
\EdL{27}{AK07-J0099-L27}{Am Ende des 510. Jahres = 61° 52′ 40″}{}
\EdL{28}{AK07-J0099-L28}{Zunahme für 29 Jahre = 348° 56′ 57″}{}
\EdL{29}{AK07-J0099-L29}{„ „ 7 Mon. = 184° 27′ 14″}{}
\EdL{30}{AK07-J0099-L30}{„ „ 29 Tage = 18° 53′ 5″}{}
\EdL{31}{AK07-J0099-L31}{„ „ 6 Std. = 3° 15′ 58″}{}
\EdL{32}{AK07-J0099-L32}{617° 25′ 54″}{}
\EdL{33}{AK07-J0099-L33}{—360°}{}
\EdL{34}{AK07-J0099-L34}{\textit{Anomalie\,:} 257° 25′ 54″}{}
\EdL{35}{AK07-J0099-L35}{Hiezu findet man in Tafel 24 die Glei-}{}
\EdL{36}{AK07-J0099-L36}{chung 4° 48′ 29″, diese zu der mittlern}{}
\EdL{37}{AK07-J0099-L37}{Länge addiert, gibt für die \textit{wahre Länge}}{}
\EdL{38}{AK07-J0099-L38}{des Mondes rund 332° 4′.}{}
\ColHead{full width}
\EdLp{39}{AK07-J0099-L39}{Differenz der beiden wahren Längen = 9° 37′. Wir haben nun die \textit{Breite} des}{}
\EdL{40}{AK07-J0099-L40}{Mondes für diese Zeit zu bestimmen\,:}{}
\ColHead{left column}
\EdL{41}{AK07-J0099-L41}{Wahre Länge des Mondes = 332° 4′.}{}
\ColHead{right column}
\EdL{42}{AK07-J0099-L42}{Länge des Drachenkopfes (n. Tafel 19 u. 20)\,:}{}
\EdL{43}{AK07-J0099-L43}{Am Ende des 510. Jahres = 92° 38′ 42″}{}
\EdL{44}{AK07-J0099-L44}{Zunahme für 29 Jahre = 184° 41′ 57″}{}
\EdL{45}{AK07-J0099-L45}{„ „ 7 Mon. = 10° 58′ 16″}{}
\EdL{46}{AK07-J0099-L46}{„ „ 29 Tage = 1° 32′ 14″}{}
\EdL{47}{AK07-J0099-L47}{„ „ 6 Std. = 48″}{}
\EdL{48}{AK07-J0099-L48}{289° 51′ 57″}{}
\NoteRule
\EdN{N1}{AK07-J0099-N01}{\textsuperscript{1}) Dass ich diesen Tag wähle, hat seinen Grund darin, weil die Bestimmung des Sichtbar-}
\EdN{N2}{AK07-J0099-N02}{werdens der Mondsichel eigentlich nur für diesen Tag von Bedeutung ist, denn mit dem Ende des}
\EdN{N3}{AK07-J0099-N03}{29. \textit{Shaʿbān} beginnt (Abends 6 Uhr) der Fastenmonat \textit{Ramaḍān.} Die Fasten beginnen aber nicht}
\EdN{N4}{AK07-J0099-N04}{etwa um 6 Uhr, sondern im Momente, wo die Mondsichel zuerst sichtbar wird. Dies kann längere}
\EdN{N5}{AK07-J0099-N05}{Zeit vor oder nach 6 Uhr stattfinden. Da nun das Wetter nicht immer gestattet, das Erscheinen}
\EdN{N6}{AK07-J0099-N06}{der Mondsichel zu beobachten, so müssen eben die Rechnung und die Tafeln zu Hilfe genommen}
\EdN{N7}{AK07-J0099-N07}{werden.}
\FrameHead{0100}{p. 68}{0100}{p. 68}
\EdL{1}{AK07-J0100-L01}{Wahre Länge weniger Länge des Drachenkopfes = 42° 12′. Dies ist das \textit{Argument}}{}
\EdL{2}{AK07-J0100-L02}{für die \textit{Breite,} dazu gehört in Tafel 22 die \textit{Breite} 3° 1′ 22″\,; diese Breite zu der vorhin}{}
\EdL{3}{AK07-J0100-L03}{gefundenen Differenz der Längen (9° 37″) addiert, denn die Breite des Mondes ist}{}
\EdL{4}{AK07-J0100-L04}{nördlich, gibt rund 12° 38′\,; diese Grösse wird die \textit{distantia definita} genannt. Mit der}{}
\EdL{5}{AK07-J0100-L05}{wahren Länge des Mondes wird nun in die Tafel 57\textsuperscript{a} eingetreten und zwar in fol-}{}
\EdL{6}{AK07-J0100-L06}{gender Weise\,: die wahre Länge wird ausgedrückt in Zeichen, Graden und Minuten,}{}
\EdL{7}{AK07-J0100-L07}{also ist sie in unserm Falle = 11\textsuperscript{s} 2° 4′, der Mond steht also im 2.° (der ersten Facies)\textsuperscript{1})}{}
\EdL{8}{AK07-J0100-L08}{der Fische\,; in unsern Tafeln entspricht der Kolonne \textit{Pisces} in der Zeile X die Zahl}{}
\EdL{9}{AK07-J0100-L09}{9° 21′, es wäre aber eher 9° 19′ zu nehmen, da der Ort des Mondes XXX° des Aqua-}{}
\EdL{10}{AK07-J0100-L10}{rius näher ist als X° der Fische, doch dies macht hier nichts aus\,; diese Zahl (9° 21′)}{}
\EdL{11}{AK07-J0100-L11}{ist nun kleiner als die oben gefundene Distanz, der Mond ist also weiter von der}{}
\EdL{12}{AK07-J0100-L12}{Sonne entfernt als 9° 21′, also muss an diesem Abend die Mondsichel um 6 Uhr}{}
\EdL{13}{AK07-J0100-L13}{sichtbar sein. Wäre die in den Tafeln gefundene Zahl grösser als die Distanz,}{}
\EdL{14}{AK07-J0100-L14}{so würde die Mondsichel am Abend des 29. \textit{Shaʿbān} nicht mehr sichtbar werden,}{}
\EdL{15}{AK07-J0100-L15}{die Fasten würden also erst am folgenden Abend, d. h. am Abend des 1. \textit{Ramaḍān}}{}
\EdL{16}{AK07-J0100-L16}{beginnen.}{}
\EdLp{17}{AK07-J0100-L17}{Die gegebene Berechnung gilt für den Meridian von \textit{Arin.} Im letzten Alinea fügt}{}
\EdL{18}{AK07-J0100-L18}{\Sperr{Kh.} die Bemerkung hinzu, dass für einen andern Ort als Arin der Zeitunterschied}{}
\EdL{19}{AK07-J0100-L19}{berücksichtigt werden müsse und dementsprechend die Längen von Sonne und Mond}{}
\EdL{20}{AK07-J0100-L20}{zu korrigieren seien.}{}
\EdLp{21}{AK07-J0100-L21}{Im mittlern Alinea wird angegeben, wie zu verfahren sei, wenn man wissen}{}
\EdL{22}{AK07-J0100-L22}{wolle, ob am Morgen des 27. Tages des Monats (also ein Tag vor der Konjunktion) die}{}
\EdL{23}{AK07-J0100-L23}{Mondsichel sichtbar werde oder nicht\,; das Verfahren ist natürlich dem vorigen analog.}{}
\EdLp{24}{AK07-J0100-L24}{Wie die Zahlen der Tafel 57\textsuperscript{a} gefunden worden sind, war uns unmöglich fest-}{}
\EdL{25}{AK07-J0100-L25}{zustellen. Die Regeln \Sperr{Kh.}’s weichen auch von denen des \textit{Sūrya-Siddhānta}\textsuperscript{2}) und der}{}
\EdL{26}{AK07-J0100-L26}{\textit{Panchasiddhāntikā}\textsuperscript{3}) ziemlich ab\,; in diesen beiden astronomischen Werken findet man}{}
\EdL{27}{AK07-J0100-L27}{noch die Angabe, die bei \Sperr{Kh.} fehlt, dass die Mondsichel sichtbar werde im Westen,}{}
\EdL{28}{AK07-J0100-L28}{wenn die Distanz derselben von der Sonne 12 oder mehr Aequatorialgrade betrage.}{}
\EdLp{29}{AK07-J0100-L29}{Die Tafel 57\textsuperscript{b} gehört nicht zu diesem Kapitel, ich habe sie an dieser Stelle}{}
\EdL{30}{AK07-J0100-L30}{untergebracht, um den kleinen Raum, den 57\textsuperscript{a} einnimmt, noch auszunützen. Wie der}{}
\EdL{31}{AK07-J0100-L31}{Leser sofort erkennen wird, ist sie eine Multiplikations-Tafel für Sexagesimal-Teile,}{}
\EdL{32}{AK07-J0100-L32}{und bedarf also keiner weitern Erklärung\,; sie ist im Text auch nirgends erwähnt.}{}
\EdLc{33}{AK07-J0100-L33}{\textbf{Kap. 23. Über die Auffindung des Sinus aus dem Bogen und umgekehrt.}}{}
\EdLc{34}{AK07-J0100-L34}{(Text p. 17—18, Tafel 58 u. 58\textsuperscript{a}.)}{}
\EdLp{35}{AK07-J0100-L35}{Über dieses trigonometrische Kapitel haben wir noch einige Bemerkungen zu}{}
\EdL{36}{AK07-J0100-L36}{dem hinzuzufügen, was schon \Sperr{Björnbo} über die Sinus- und Cotangententafel dieses}{}
\EdL{37}{AK07-J0100-L37}{Werkes veröffentlicht hat.\textsuperscript{4})}{}
\NoteRule
\EdN{N1}{AK07-J0100-N01}{\textsuperscript{1}) Ein Tierkreiszeichen zerfällt in drei Dekane (\textit{facies}), von denen jeder 10° umfasst.}
\EdN{N2}{AK07-J0100-N02}{\textsuperscript{2}) \textit{Transl.} \Sperr{Burgess}, p. 365 u. 370.}
\EdN{N3}{AK07-J0100-N03}{\textsuperscript{3}) Ed. \Sperr{Thibaut}, p. 38—39 der Übersetzung.}
\EdN{N4}{AK07-J0100-N04}{\textsuperscript{4}) \Sperr{Al-Chwarizmi}’s \textit{Trigonometriske Tavler.} \textit{Af Festskrift til} H. G. \Sperr{Zeuthen}, København, 1909.}
\FrameHead{0101}{p. 69}{0101}{p. 69}
\EdLp{1}{AK07-J0101-L01}{\Sperr{Kh.} nennt den Sinus eines Winkels \textit{elgeib elmustewi} = \textit{al-djaib al-mustawī,} und}{}
\EdL{2}{AK07-J0101-L02}{\Sperr{Athelhard} gibt dies durch \textit{elgeib planum} wieder (die meisten arabischen Wörter}{}
\EdL{3}{AK07-J0101-L03}{nimmt er als Neutra an)\,; \textit{planus} ist wohl eine Bedeutung von \textit{mustawī,} besser aber}{}
\EdL{4}{AK07-J0101-L04}{wäre hier \textit{rectus} = gerade, oder gerad aufgerichtet. \Sperr{Athelhard} kannte also (ums Jahr}{}
\EdL{5}{AK07-J0101-L05}{1126) die Übersetzung von \textit{djaib} durch das lateinische \textit{sinus} noch nicht, dagegen}{}
\EdL{6}{AK07-J0101-L06}{ersehen wir aus dem Madrider Ms. N, dass \Sperr{Robert von Chester}, der die bessere}{}
\EdL{7}{AK07-J0101-L07}{„Ordnung“ der \Sperr{Athelhard}’schen Übersetzung wohl etwa zwischen 1143 und 1150}{}
\EdL{8}{AK07-J0101-L08}{ausgeführt haben mag, das Wort \textit{sinus} gekannt hat. Wir wissen nicht, wann \Sperr{Gerhard}}{}
\EdL{9}{AK07-J0101-L09}{\Sperr{von Cremona}, der seine Übersetzungen meist nicht datiert hat, die Astronomie}{}
\EdL{10}{AK07-J0101-L10}{des \Sperr{Djābir b. Aflaḥ}, in der das Wort \textit{sinus} vorkommt, übersetzt hat, können also}{}
\EdL{11}{AK07-J0101-L11}{nicht entscheiden, ob er oder \Sperr{Robert von Chester} zuerst dieses Wort gebraucht}{}
\EdL{12}{AK07-J0101-L12}{habe, doch ist die Wahrscheinlichkeit für letztern grösser.}{}
\EdLp{13}{AK07-J0101-L13}{Den \textit{Sinus versus} nennt \Sperr{Athelhard} \textit{elgeib elmankūs} (oder \textit{maʿkūs}) = der um-}{}
\EdL{14}{AK07-J0101-L14}{gekehrte Sinus\,; er übersetzt aber \textit{mankūs} unrichtig mit \textit{diminutum,} indem er es mit}{}
\EdL{15}{AK07-J0101-L15}{\textit{manḳūṣ} verwechselt, das „vermindert“, „verkleinert“ bedeutet, und beinahe gleich}{}
\EdL{16}{AK07-J0101-L16}{ausgesprochen wird. Dies mag wohl ein Beweis dafür sein, dass man damals eine}{}
\EdL{17}{AK07-J0101-L17}{Sprache nicht sowohl aus Büchern erlernte, als vielmehr im sprachlichen Verkehr}{}
\EdL{18}{AK07-J0101-L18}{mit Menschen, daher bei Übersetzungen mehr auf das gesprochene Wort, als auf}{}
\EdL{19}{AK07-J0101-L19}{das geschriebene, bezw. gelesene vertraute\,; daher mögen wohl auch noch andere}{}
\EdL{20}{AK07-J0101-L20}{falsche Schreibweisen arabischer Wörter stammen.}{}
\EdLp{21}{AK07-J0101-L21}{Die unrichtigen Übersetzungen \Sperr{Athelhards}, \textit{mustawī} durch \textit{planus} und \textit{mankūs}}{}
\EdL{22}{AK07-J0101-L22}{durch \textit{diminutus,} hat \Sperr{Robert von Chester} nicht verbessert, er nennt den Sinus}{}
\EdL{23}{AK07-J0101-L23}{ebenfalls \textit{sinus planus} und den Sinus versus \textit{sinus diminutus.}}{}
\EdLp{24}{AK07-J0101-L24}{Wir haben noch besonders darauf aufmerksam zu machen, dass \Sperr{Kh.} in der}{}
\EdL{25}{AK07-J0101-L25}{Sinustafel in den vier Spalten, die mit \textit{Semitae numerorum} überschrieben sind, die}{}
\EdL{26}{AK07-J0101-L26}{Winkel aller vier Quadranten angibt, die denselben Sinus (abgesehen vom Vorzeichen)}{}
\EdL{27}{AK07-J0101-L27}{haben, so in der 6. Zahlenzeile\,:}{}
\EdL{28}{AK07-J0101-L28}{\parbox[t]{\dimexpr0.0980\linewidth-0.5em\relax}{\raggedright\small\mbox{}Sig.}\hspace{0.5em}\parbox[t]{\dimexpr0.0980\linewidth-0.5em\relax}{\raggedright\small\mbox{}Gr.}\hspace{0.5em}\parbox[t]{\dimexpr0.0980\linewidth-0.5em\relax}{\raggedright\small\mbox{}Sig.}\hspace{0.5em}\parbox[t]{\dimexpr0.0980\linewidth-0.5em\relax}{\raggedright\small\mbox{}Gr.}\hspace{0.5em}\parbox[t]{\dimexpr0.0980\linewidth-0.5em\relax}{\raggedright\small\mbox{}Sig.}\hspace{0.5em}\parbox[t]{\dimexpr0.0980\linewidth-0.5em\relax}{\raggedright\small\mbox{}Gr.}\hspace{0.5em}\parbox[t]{\dimexpr0.0980\linewidth-0.5em\relax}{\raggedright\small\mbox{}Sig.}\hspace{0.5em}\parbox[t]{\dimexpr0.2940\linewidth-0.5em\relax}{\raggedright\small\mbox{}Gr.}\hspace{0.5em}}{}
\EdL{29}{AK07-J0101-L29}{\parbox[t]{\dimexpr0.0980\linewidth-0.5em\relax}{\raggedright\small\mbox{}0}\hspace{0.5em}\parbox[t]{\dimexpr0.0980\linewidth-0.5em\relax}{\raggedright\small\mbox{}6}\hspace{0.5em}\parbox[t]{\dimexpr0.0980\linewidth-0.5em\relax}{\raggedright\small\mbox{}5}\hspace{0.5em}\parbox[t]{\dimexpr0.0980\linewidth-0.5em\relax}{\raggedright\small\mbox{}24}\hspace{0.5em}\parbox[t]{\dimexpr0.0980\linewidth-0.5em\relax}{\raggedright\small\mbox{}6}\hspace{0.5em}\parbox[t]{\dimexpr0.0980\linewidth-0.5em\relax}{\raggedright\small\mbox{}6}\hspace{0.5em}\parbox[t]{\dimexpr0.0980\linewidth-0.5em\relax}{\raggedright\small\mbox{}11}\hspace{0.5em}\parbox[t]{\dimexpr0.0980\linewidth-0.5em\relax}{\raggedright\small\mbox{}24}\hspace{0.5em}\parbox[t]{\dimexpr0.1960\linewidth-0.5em\relax}{\raggedright\small\mbox{}das heisst\,:}\hspace{0.5em}}{}
\EdL{30}{AK07-J0101-L30}{\parbox[t]{\dimexpr0.1960\linewidth-0.5em\relax}{\raggedright\small\mbox{}6°}\hspace{0.5em}\parbox[t]{\dimexpr0.1960\linewidth-0.5em\relax}{\raggedright\small\mbox{}174°}\hspace{0.5em}\parbox[t]{\dimexpr0.1960\linewidth-0.5em\relax}{\raggedright\small\mbox{}186°}\hspace{0.5em}\parbox[t]{\dimexpr0.3920\linewidth-0.5em\relax}{\raggedright\small\mbox{}354°}\hspace{0.5em}}{}
\EdLp{31}{AK07-J0101-L31}{Der Sinus ist angegeben in \textit{partes} (Radius = 60 solcher), oder wie der Text hat}{}
\EdL{32}{AK07-J0101-L32}{\textit{partitiones, dakaicae} = arab. \textit{daḳā’iḳ} = Minuten, und \textit{theniae} (mit dem Artikel \textit{eltheniae})}{}
\EdL{33}{AK07-J0101-L33}{= arab. \textit{thāniya} = Sekunden\textsuperscript{1}). So ist der Sinus von 6° = 6\textsuperscript{\textit{p}} 16′ 18″ = 6 + \Frac{16}{60} + \Frac{18}{60\textsuperscript{2}} =}{}
\EdL{34}{AK07-J0101-L34}{6,27166, dies gibt, den Radius = 1 gesetzt\,: sin 6° = 0,10453 (richtiger Wert = 0,104528).}{}
\EdLp{35}{AK07-J0101-L35}{Bei Abweichungen von den Tafeln \Sperr{Batt.}’s haben wir, wenn der Fehler nur 1″}{}
\EdL{36}{AK07-J0101-L36}{betrug, die Lesart der Manuskripte in die Tafel aufgenommen, bei einem Fehler von}{}
\EdL{37}{AK07-J0101-L37}{mehr als 1″ aber den richtigen Wert \Sperr{Batt.}’s, entgegen dem Verfahren \Sperr{Björnbo}’s}{}
\EdL{38}{AK07-J0101-L38}{in der oben genannten Schrift, der überall die Werte von C und O aufgenommen}{}
\EdL{39}{AK07-J0101-L39}{hat, auch wenn sie unrichtig waren.}{}
\NoteRule
\EdN{N1}{AK07-J0101-N01}{\textsuperscript{1}) Vergl. oben p. 36.}
\FrameHead{0102}{p. 70}{0102}{p. 70}
\EdLp{1}{AK07-J0102-L01}{Wir sind mit \Sperr{Björnbo} der Ansicht, \Sperr{Kh.} habe seine Sinustafel der \Sperr{Ptole-}}{}
\EdL{2}{AK07-J0102-L02}{\Sperr{mäischen} Sehnentafel durch Halbierung der Winkel und Sehnen entnommen,}{}
\EdL{3}{AK07-J0102-L03}{erwähnt er doch selbst am Schlusse des Kapitels den \textit{Almagest}\,; vielleicht rühren}{}
\EdL{4}{AK07-J0102-L04}{auch daher die ziemlich häufig vorkommenden Abweichungen von 1″, indem der}{}
\EdL{5}{AK07-J0102-L05}{eine bei der Halbierung ungerader Sekundenzahlen auf-, der andere abgerundet hat.}{}
\EdL{6}{AK07-J0102-L06}{Unmöglich wäre es freilich nicht, dass der indische Gelehrte, der ums Jahr 750 an}{}
\EdL{7}{AK07-J0102-L07}{den Hof \Sperr{Al-Manṣūrs} kam, in seinem \textit{Siddhānta} eine Sinustafel von Grad zu Grad}{}
\EdL{8}{AK07-J0102-L08}{fortschreitend mitgebracht hätte, die dann vom indischen Radius 3438 auf den Radius}{}
\EdL{9}{AK07-J0102-L09}{60 reduziert worden wäre. Doch kennen wir bis jetzt kein früheres Werk als den}{}
\EdL{10}{AK07-J0102-L10}{\textit{Siddhānta-Śiromani} des \Sperr{Bhāskara}, in dem die Berechnung des Sinus von 1° gelehrt}{}
\EdL{11}{AK07-J0102-L11}{wird. Übrigens darf nicht unerwähnt bleiben, dass ein wahrscheinlich etwas älterer}{}
\EdL{12}{AK07-J0102-L12}{Zeitgenosse \Sperr{Kh.}’s, \Sperr{Yaʿḳūb b. Ṭāriḳ}, zwei Werke verfasst hat, betitelt „über die}{}
\EdL{13}{AK07-J0102-L13}{Teilung des \textit{Kardadja}“ (d. h. des Sinus oder Bogens von 3° 45′), und „astronomische}{}
\EdL{14}{AK07-J0102-L14}{Tafeln nach dem \textit{Sindhind} verfasst von Grad zu Grad“.\textsuperscript{1}) Hieraus muss man, wenn}{}
\NoteRule
\EdN{N1}{AK07-J0102-N01}{\textsuperscript{1}) Vergl. H. \Sperr{Suter}, \textit{Die Mathematiker und Astronomen der Araber}, etc. p. 4.}
\FrameHead{0103}{p. 71}{0103}{p. 71}
\EdL{1}{AK07-J0103-L01}{diese Titel genau sind, den Schluss ziehen, dass schon vor \Sperr{Muh. b. Mūsā al-Kh.}}{}
\EdL{2}{AK07-J0103-L02}{Sinustafeln von Grad zu Grad fortschreitend existiert haben.}{}
\EdLp{3}{AK07-J0103-L03}{Was die Regeln zur Auffindung des Sinus und Sinus versus zu einem Winkel}{}
\EdL{4}{AK07-J0103-L04}{und umgekehrt betrifft, so werde ich wohl dem Texte nichts beizufügen brauchen.}{}
\EdL{5}{AK07-J0103-L05}{Sonderbar ist die Bemerkung in Alinea 4 „\textit{ut in sequenti tabula subnotatur}“, gleich-}{}
\EdL{6}{AK07-J0103-L06}{sam als ob zur Auffindung des Winkels aus dem Sinus eine andere Tafel notwendig}{}
\EdL{7}{AK07-J0103-L07}{wäre. Wir machen hier darauf aufmerksam, dass A. \Sperr{Björnbo} der Ansicht war\textsuperscript{1}),}{}
\EdL{8}{AK07-J0103-L08}{es könnte das Werk \Sperr{Kh.}’s noch eine Arcus-Sinustafel enthalten haben, wie sie bei}{}
\EdL{9}{AK07-J0103-L09}{\Sperr{Abū’l-Ḥasan ʿAlī} von Marokko vorkommt, der sie die „Tafel des Kh.“ nennt.}{}
\EdLp{10}{AK07-J0103-L10}{Ich schliesse hier noch eine Figur an, die am Schlusse dieses Kapitels, aber}{}
\EdL{11}{AK07-J0103-L11}{nur im Ms. N steht, und die jedenfalls dem Autographen \Sperr{Athelhards} nicht ange-}{}
\EdL{12}{AK07-J0103-L12}{hört hat.}{}
\EdLp{13}{AK07-J0103-L13}{Die Zahlen rings herum von I bis XII bedeuten die zwölf Zeichen (\textit{Signa}) zu}{}
\EdL{14}{AK07-J0103-L14}{je 30°\,; die sechs XX über dem horizontalen Durchmesser bezeichnen je einen}{}
\EdL{15}{AK07-J0103-L15}{Drittel des Radius, also 20 partes\,; unter dem horizontalen Durchmesser sollte statt}{}
\EdL{16}{AK07-J0103-L16}{Gr. (= Gradus) P. (= Partes) stehen\,; man kann dann aus der Figur ablesen\,: sin I}{}
\EdL{17}{AK07-J0103-L17}{= sin 30° = sin V = sin 150° = sin VII = sin 210° = sin XI = sin 330° = XXX P.\,;}{}
\EdL{18}{AK07-J0103-L18}{sin (I + II) = sin 60° = sin 120° = sin 240° = sin 300° = XXX P. + XXI P. LVII M.}{}
\EdL{19}{AK07-J0103-L19}{XLI S. = LI P. LVII M. XLI S.\,; sin (I+II+III) = sin 90° = sin 270° = LI P. LVII M.}{}
\EdL{20}{AK07-J0103-L20}{XLI S. + VIII P. II M. XIX S. = LX P. — Zu bemerken ist noch, dass in dieser Figur}{}
\EdL{21}{AK07-J0103-L21}{die Sinuslinien horizontal stehen, dann wäre auch die Übersetzung von \textit{djaib mustawī}}{}
\EdL{22}{AK07-J0103-L22}{durch \textit{sinus planus} eher zu verstehen.}{}
\EdLc{23}{AK07-J0103-L23}{\textbf{Kap. 24. Über die Auffindung der Breite eines beliebigen Ortes.}}{}
\EdLc{24}{AK07-J0103-L24}{(Text p. 18.)}{}
\EdLp{25}{AK07-J0103-L25}{Zuerst wird angegeben, dass, wenn die Sonne im Anfang des Widders oder}{}
\EdL{26}{AK07-J0103-L26}{der Wage (d. h. im Aequator) stehe, ihre Mittagshöhe (\textit{H}) von 90° zu subtrahieren}{}
\EdL{27}{AK07-J0103-L27}{sei, um die geogr. Breite des Ortes zu erhalten (\textit{φ} = 90° — \textit{H}).}{}
\EdLp{28}{AK07-J0103-L28}{Steht die Sonne in einem andern Zeichen, so subtrahiert man von ihrer Mittags-}{}
\EdL{29}{AK07-J0103-L29}{höhe ihre Deklination, wenn diese nördlich, oder addiert sie zu der Mittagshöhe,}{}
\EdL{30}{AK07-J0103-L30}{wenn sie südlich ist, und subtrahiert das Resultat von 90°. (\textit{φ} = 90° — (\textit{H} ∓ \textit{δ})).}{}
\EdLp{31}{AK07-J0103-L31}{Dann wird die Methode angegeben, mit dem Astrolabium die Polhöhe zu}{}
\EdL{32}{AK07-J0103-L32}{bestimmen, die gleich der geogr. Breite des Ortes ist.}{}
\EdLp{33}{AK07-J0103-L33}{Endlich wird die Methode behandelt, aus der grössten und kleinsten Höhe}{}
\EdL{34}{AK07-J0103-L34}{eines Circumpolarsternes die Höhe des Poles, d. h. die geogr. Breite zu bestimmen.}{}
\EdL{35}{AK07-J0103-L35}{(\textit{φ} = \textit{h′} + \Frac{\textit{h″} — \textit{h′}}{2} = \Frac{\textit{h′} + \textit{h″}}{2}).}{}
\EdLp{36}{AK07-J0103-L36}{Dieselben Bestimmungen findet man auch bei \Sperr{Batt.}\textsuperscript{2})}{}
\NoteRule
\EdN{N1}{AK07-J0103-N01}{\textsuperscript{1}) Vergl. die oben zitierte \textit{Festskrift til} H. G. \Sperr{Zeuthen}, p. 14, und v. \Sperr{Braunmühl}, \textit{Vor-}}
\EdN{N2}{AK07-J0103-N02}{\textit{lesungen über Geschichte der Trigonometrie}, Leipzig 1900, I, p. 84.}
\EdN{N3}{AK07-J0103-N03}{\textsuperscript{2}) Ed. \Sperr{Nallino}, I, 15, 29—30.}
\FrameHead{0104}{p. 72}{0104}{p. 72}
\EdLc{1}{AK07-J0104-L01}{\textbf{Kap. 25. Die Aufgänge der Tierkreiszeichen in der geraden Sphäre.}}{}
\EdLc{2}{AK07-J0104-L02}{(Text p. 18—19, Tafel 59—59\textsuperscript{b}).}{}
\EdLp{3}{AK07-J0104-L03}{Dieses Kapitel mit den zugehörigen Tafeln lehrt die \textit{Rectascension} (\textit{a}) jedes Grades}{}
\EdL{4}{AK07-J0104-L04}{der Ekliptik in der geraden Sphäre (d. h. am Aequator, oder wie unsere Tafeln haben\,:}{}
\EdL{5}{AK07-J0104-L05}{\textit{secundum terram Arin})\textsuperscript{1}) berechnen. Es geschieht dies nach derselben Formel wie}{}
\EdL{6}{AK07-J0104-L06}{bei \Sperr{Batt.}, nämlich}{}
\EdLc{7}{AK07-J0104-L07}{sin \textit{a} = \Frac{sin \textit{δ} · cos \textit{ε}}{cos \textit{δ} · sin \textit{ε}} · \textit{r}.}{}
\EdLp{8}{AK07-J0104-L08}{Da \Sperr{Batt.}, sowie auch \Sperr{Kh.} die Tangenten und Cotangenten wohl kennen, aber}{}
\EdL{9}{AK07-J0104-L09}{nur Tafeln derselben für den Radius (Gnomon) 12 haben, so benutzen sie nicht die}{}
\EdL{10}{AK07-J0104-L10}{einfachere Formel\,: sin \textit{a} = tg \textit{δ} · cot \textit{ε}, die der Formel sin \textit{a} = \Frac{tg \textit{b}}{tg \textit{B}} des rechtwinkligen}{}
\EdL{11}{AK07-J0104-L11}{sphärischen Dreieckes entspricht. In unserm Falle ist Winkel \textit{B} die Schiefe (\textit{ε}) der}{}
\EdL{12}{AK07-J0104-L12}{Ekliptik, \textit{c} ein bestimmter Bogen der Ekliptik, \textit{b} die Deklination (\textit{δ}) des Punktes \textit{A}}{}
\EdL{13}{AK07-J0104-L13}{der Ekliptik, \textit{a} der entsprechende Bogen des Aequa-}{}
\EdL{14}{AK07-J0104-L14}{tors, d. h. die \textit{Rectascension} des Punktes \textit{A}\,; obige}{}
\EdL{15}{AK07-J0104-L15}{Formel lautet also für unsern Fall\,:}{}
\EdLc{16}{AK07-J0104-L16}{sin \textit{a} = \Frac{tg \textit{δ}}{tg \textit{ε}} = tg \textit{δ} · cot \textit{ε} = \Frac{sin \textit{δ} · cos \textit{ε}}{cos \textit{δ} · sin \textit{ε}}.}{}
\EdL{17}{AK07-J0104-L17}{Diesen Ausdruck mussten die arabischen Mathema-}{}
\EdL{18}{AK07-J0104-L18}{tiker noch mit \textit{r} = 60\textsuperscript{\textit{p}} multiplizieren, weil ihre}{}
\EdL{19}{AK07-J0104-L19}{trigonometrischen Funktionen für diesen Radius be-}{}
\EdL{20}{AK07-J0104-L20}{rechnet waren. Wir geben hiefür ein \textit{Beispiel\,:}}{}
\EdLp{21}{AK07-J0104-L21}{Was für eine Rectascension entspricht dem Bogen 25° des Widders\,? Für}{}
\EdL{22}{AK07-J0104-L22}{diesen Fall müssen wir zuerst in Tafel 21 die Deklination (Obliquatio) \textit{δ} der Sonne}{}
\EdL{23}{AK07-J0104-L23}{für den 25.° des Widders suchen, sie ist 9° 50′ 21″, die Schiefe der Ekliptik nimmt}{}
\EdL{24}{AK07-J0104-L24}{\Sperr{Kh.} zu 23° 51′ an, die Berechnung nach der Formel sin \textit{a} = tg \textit{δ} · cot \textit{ε} ergibt (mit}{}
\EdL{25}{AK07-J0104-L25}{fünfstelligen Logarithmen berechnet) für \textit{a} den Wert 23° 5′ 50″, unsere Tafel 59 hat}{}
\EdL{26}{AK07-J0104-L26}{bei 25° Länge (in der Rubrik \textit{Horoscopus arietis}) die Zahl 113° 5′ 49″\,; da \Sperr{Kh.}, wie}{}
\EdL{27}{AK07-J0104-L27}{auch \Sperr{Batt.} in den entsprechenden Tafeln\textsuperscript{2}), die Rectascension von 0° des Steinbocks}{}
\EdL{28}{AK07-J0104-L28}{an zählt statt von 0° des Widders an, so muss man 3 Zeichen = 90° abziehen und}{}
\EdL{29}{AK07-J0104-L29}{erhält so \textit{a} = 23° 5′ 49″.\textsuperscript{3}) — \Sperr{Batt.}’s Werte weichen natürlich von denen \Sperr{Kh.}’s ab,}{}
\EdL{30}{AK07-J0104-L30}{da er die Schiefe der Ekliptik = 23° 35′ annimmt, auch begnügt sich \Sperr{Batt.} mit}{}
\EdL{31}{AK07-J0104-L31}{Graden und Minuten, er hat in unserm Beispiele eine Rectascension von 113° 8′}{}
\EdL{32}{AK07-J0104-L32}{(= 23° 8′).}{}
\EdLp{33}{AK07-J0104-L33}{Dass einige arabische Astronomen die Rectascension von 0° des Steinbocks}{}
\EdL{34}{AK07-J0104-L34}{an zählen, hat seinen Grund darin, weil dann gerade diese um 90° vermehrte Rect-}{}
\EdL{35}{AK07-J0104-L35}{ascension den Grad des Aequators angibt, der zu derselben Zeit aufgeht, da der}{}
\NoteRule
\EdN{N1}{AK07-J0104-N01}{\textsuperscript{1}) \textit{Arin} (\textit{Azin}, heute \textit{Udjain}, s. oben p. 33) wurde von den ältern arabischen Astronomen unter}
\EdN{N2}{AK07-J0104-N02}{den Aequator versetzt, und als die Mitte der bewohnten Erdhälfte betrachtet.}
\EdN{N3}{AK07-J0104-N03}{\textsuperscript{2}) Ed. \Sperr{Nallino}, II, 61—64.}
\EdN{N4}{AK07-J0104-N04}{\textsuperscript{3}) Wir haben bei sechs Proben keinen grössern Fehler als 1″—2″ getroffen.}
\FrameHead{0105}{p. 73}{0105}{p. 73}
\EdL{1}{AK07-J0105-L01}{betreffende Grad der Ekliptik kulminiert, das war für astrologische Zwecke wichtig}{}
\EdL{2}{AK07-J0105-L02}{zu wissen. In unserm Beispiele geht also der Punkt des Aequators mit der Rect-}{}
\EdL{3}{AK07-J0105-L03}{ascension 113° 5′ 49″ zu gleicher Zeit auf, da der Punkt der Ekliptik mit der Länge}{}
\EdL{4}{AK07-J0105-L04}{25° kulminiert.\textsuperscript{1})}{}
\EdLp{5}{AK07-J0105-L05}{\Sperr{Kh.} braucht den Ausdruck „Sinus der Ergänzung“ (d. h. Cosinus) noch nicht}{}
\EdL{6}{AK07-J0105-L06}{wie \Sperr{Batt.}, er sagt\,: „man subtrahiere den Bogen von 90° und suche zu der Differenz}{}
\EdL{7}{AK07-J0105-L07}{den Sinus.“ Der Begriff einer besondern Funktion war bei \Sperr{Kh.} also noch nicht so}{}
\EdL{8}{AK07-J0105-L08}{weit entwickelt wie bei \Sperr{Batt.}}{}
\EdLp{9}{AK07-J0105-L09}{Die Erklärung der Tafeln stimmt mit derjenigen bei \Sperr{Batt.}\textsuperscript{2}) fast durchweg}{}
\EdL{10}{AK07-J0105-L10}{überein. Wir haben noch einige Fehler unserer Tafeln zu erwähnen\,: Die Rect-}{}
\EdL{11}{AK07-J0105-L11}{ascensionen der drei ersten Grade des Krebses (Tafel 59\textsuperscript{a}) stimmen in den Sekunden}{}
\EdL{12}{AK07-J0105-L12}{nicht mit denjenigen des Steinbocks, und sollten doch dieselben sein vermehrt um}{}
\EdL{13}{AK07-J0105-L13}{180°\,; unsere Berechnung ergibt sogar in den Minuten noch teilweise Abweichungen,}{}
\EdL{14}{AK07-J0105-L14}{es können diese auch von Fehlern in den Deklinationen herrühren, denn bei Längen,}{}
\EdL{15}{AK07-J0105-L15}{die nicht mehr als 1°—2° von 90° bezw. 270° abweichen, bewirkt ein Unterschied}{}
\EdL{16}{AK07-J0105-L16}{von 1″ in der Deklination einen solchen von 1\textsuperscript{1}/\textsubscript{2}′ bis 2\textsuperscript{1}/\textsubscript{2}′ in der Rectascension\,;}{}
\EdL{17}{AK07-J0105-L17}{nimmt man z. B. die Deklination bei 89° bezw. 271° 23° 50′ 46″ an statt 23° 50′ 47″,}{}
\EdL{18}{AK07-J0105-L18}{wie die Tafel hat, so wird die Rectascension von 1° des Steinbocks = 1° 5′ 55″ statt}{}
\EdL{19}{AK07-J0105-L19}{1° 3′ 30″. Nach unsern Berechnungen ergeben sich für die Rectascensionen der drei}{}
\EdL{20}{AK07-J0105-L20}{ersten Grade des Steinbocks und des Krebses folgende Werte, wenn man die Deklina-}{}
\EdL{21}{AK07-J0105-L21}{tionen nimmt, wie sie in Tafel 23 und 24 angegeben sind\,:}{}
\EdL{22}{AK07-J0105-L22}{\parbox[t]{\dimexpr0.2450\linewidth-0.5em\relax}{\raggedright\small\mbox{}}\hspace{0.5em}\parbox[t]{\dimexpr0.2450\linewidth-0.5em\relax}{\raggedright\small\mbox{}Steinbock\,:}\hspace{0.5em}\parbox[t]{\dimexpr0.4900\linewidth-0.5em\relax}{\raggedright\small\mbox{}Krebs\,:}\hspace{0.5em}}{}
\EdL{23}{AK07-J0105-L23}{\parbox[t]{\dimexpr0.2450\linewidth-0.5em\relax}{\raggedright\small\mbox{}}\hspace{0.5em}\parbox[t]{\dimexpr0.2450\linewidth-0.5em\relax}{\raggedright\small\mbox{}1° 3′ 30″}\hspace{0.5em}\parbox[t]{\dimexpr0.4900\linewidth-0.5em\relax}{\raggedright\small\mbox{}181° 3′ 30″}\hspace{0.5em}}{}
\EdL{24}{AK07-J0105-L24}{\parbox[t]{\dimexpr0.2450\linewidth-0.5em\relax}{\raggedright\small\mbox{}}\hspace{0.5em}\parbox[t]{\dimexpr0.2450\linewidth-0.5em\relax}{\raggedright\small\mbox{}2° 10′ 34″}\hspace{0.5em}\parbox[t]{\dimexpr0.4900\linewidth-0.5em\relax}{\raggedright\small\mbox{}182° 10′ 34″}\hspace{0.5em}}{}
\EdL{25}{AK07-J0105-L25}{\parbox[t]{\dimexpr0.2450\linewidth-0.5em\relax}{\raggedright\small\mbox{}}\hspace{0.5em}\parbox[t]{\dimexpr0.2450\linewidth-0.5em\relax}{\raggedright\small\mbox{}3° 16′ 48″}\hspace{0.5em}\parbox[t]{\dimexpr0.4900\linewidth-0.5em\relax}{\raggedright\small\mbox{}183° 16′ 48″}\hspace{0.5em}}{}
\EdLp{26}{AK07-J0105-L26}{Von diesen Werten stimmt nur derjenige für 3° des Krebses mit dem in Tafel 59\textsuperscript{a}}{}
\EdL{27}{AK07-J0105-L27}{gegebenen bis auf 1″ überein.}{}
\EdLc{28}{AK07-J0105-L28}{\textbf{Kap. 26 u. 26\textsuperscript{a}. Mit wie viel Graden (des Aequators) an einem beliebigen Orte jedes}}{}
\EdLc{29}{AK07-J0105-L29}{\textbf{Tierkreiszeichen aufsteigt.}}{}
\EdLc{30}{AK07-J0105-L30}{(Text p. 20 u. 21, Tafel 59–59\textsuperscript{b}.)}{}
\EdLp{31}{AK07-J0105-L31}{Für Orte, die eine bestimmte Breite haben, also nicht im Aequator liegen, hat}{}
\EdL{32}{AK07-J0105-L32}{unser Werk keine besondern Tafeln für die Bestimmung der Aufsteigungen, erstens}{}
\EdL{33}{AK07-J0105-L33}{wohl aus dem Grunde, weil, wenn man viele Breiten berücksichtigen wollte, die Zahl}{}
\EdL{34}{AK07-J0105-L34}{der Tafeln sich sehr vergrössern würde, und zweitens vielleicht auch deshalb, weil}{}
\EdL{35}{AK07-J0105-L35}{diese sog. schiefen Aufsteigungen (\textit{ascensiones obliquae}) in ziemlich einfacher Weise}{}
\EdL{36}{AK07-J0105-L36}{aus den geraden abgeleitet werden können. \Sperr{Batt.} hat solche Tafeln für die Breiten}{}
\NoteRule
\EdN{N1}{AK07-J0105-N01}{\textsuperscript{1}) Wir würden sagen\,: Der Punkt der Ekliptik mit der Länge 25° geht zu gleicher Zeit auf}
\EdN{N2}{AK07-J0105-N02}{wie der Punkt des Aequators mit der Rectascension 23° 5′ 49″.}
\EdN{N3}{AK07-J0105-N03}{\textsuperscript{2}) Ed. \Sperr{Nallino}, I, 13—14.}
\FrameHead{0106}{p. 74}{0106}{p. 74}
\EdL{1}{AK07-J0106-L01}{von \textit{Mekka}, \textit{Baghdād}, \textit{Ḥarrān} und \textit{al-Raḳḳa}, und dann noch allgemein für alle Orte,}{}
\EdL{2}{AK07-J0106-L02}{deren Tageslängen 13\textsuperscript{h} bis 16\textsuperscript{h}, in Intervallen von je 15\textsuperscript{m}, betragen.\textsuperscript{1})}{}
\EdLp{3}{AK07-J0106-L03}{\Sperr{Batt.} gibt zwei Wege an für die Bestimmung der schiefen Aufsteigungen für}{}
\EdL{4}{AK07-J0106-L04}{irgend einen Ort\textsuperscript{2}), der zweite, dem auch teilweise \Sperr{Kh.} folgt, ist folgender\,:}{}
\EdLc{5}{AK07-J0106-L05}{\textit{Ascensio obliqua} für die Breite \textit{φ} = \textit{Ascensio recta} ∓ \textit{arc sin} (\textit{tg φ} · \textit{tg δ})}{}
\EdL{6}{AK07-J0106-L06}{wo \textit{δ} die Deklination des betreffenden Punktes der Ekliptik ist, dessen \textit{Ascensio obliqua}}{}
\EdL{7}{AK07-J0106-L07}{man sucht. Diese Formel ergibt sich aus folgender Figur, aus der sie auch \Sperr{Ptol.} mittels}{}
\EdL{8}{AK07-J0106-L08}{des Satzes des \Sperr{Menelaus} ableitet\textsuperscript{3})\,; \textit{HT} sei}{}
\EdL{9}{AK07-J0106-L09}{der Horizont, \textit{AQ} der Aequator, \textit{EK} die Ekliptik,}{}
\EdL{10}{AK07-J0106-L10}{♈ der Frühlingspunkt, \textit{P} der nördliche Pol}{}
\EdL{11}{AK07-J0106-L11}{des Aequators, \textit{PN} der Deklinationskreis des}{}
\EdL{12}{AK07-J0106-L12}{aufgehenden Punktes \textit{L} der Ekliptik\,; dann}{}
\EdL{13}{AK07-J0106-L13}{ist ♈\textit{N} die Rectascension und ♈\textit{D} die Ascen-}{}
\EdL{14}{AK07-J0106-L14}{sio obliqua des Punktes \textit{L} der Ekliptik, \textit{DN}}{}
\EdL{15}{AK07-J0106-L15}{ist die sog. \textit{Ascensional-Differenz.} Nun ist}{}
\EdLc{16}{AK07-J0106-L16}{♈\textit{D} = ♈\textit{N} — \textit{DN},}{}
\EdL{17}{AK07-J0106-L17}{\textit{DN} aber ergibt sich aus dem rechtwinkligen}{}
\EdL{18}{AK07-J0106-L18}{sphärischen Dreieck \textit{DNL} in folgender Weise\,:}{}
\EdL{19}{AK07-J0106-L19}{\textit{LN} ist die Deklination \textit{δ} des Punktes \textit{L},}{}
\EdL{20}{AK07-J0106-L20}{∠ \textit{D} = arcus \textit{AH} = 90°—\textit{φ}\,; also ist nach der}{}
\EdL{21}{AK07-J0106-L21}{oben angeführten Formel (sin \textit{a} = \Frac{tg \textit{b}}{tg \textit{B}}):}{}
\EdLc{22}{AK07-J0106-L22}{sin \textit{DN} = \Frac{tg \textit{LN}}{tg \textit{D}} = \Frac{tg \textit{δ}}{tg (90—\textit{φ})} = tg \textit{δ} · tg \textit{φ},}{}
\EdL{23}{AK07-J0106-L23}{also ♈\textit{D} = ♈\textit{N} — arc sin (tg \textit{δ} · tg \textit{φ})}{}
\EdL{24}{AK07-J0106-L24}{oder\,: \textit{Ascensio obliqua} = \textit{Ascensio recta} — \textit{arc sin} (\textit{tg δ} · \textit{tg φ}).}{}
\EdL{25}{AK07-J0106-L25}{Ist die Deklination \textit{δ} des Punktes \textit{L} südlich, so muss der Bogen, dessen Sinus}{}
\EdL{26}{AK07-J0106-L26}{tg \textit{δ} · tg \textit{φ} ist, zur \textit{Ascensio recta} addiert werden.}{}
\EdLp{27}{AK07-J0106-L27}{\Sperr{Batt.} gebraucht die Formel wieder in der Form\,: sin \textit{DN} = \Frac{sin \textit{φ} · sin \textit{δ}}{cos \textit{φ} · cos \textit{δ}},}{}
\EdL{28}{AK07-J0106-L28}{während merkwürdigerweise \Sperr{Kh.} tg \textit{φ} mit Hilfe des Gnomons bestimmt, oder aus der}{}
\EdL{29}{AK07-J0106-L29}{Schattentafel (Tafel der Cotangenten) entnimmt\textsuperscript{4}), er multipliziert dieses tg \textit{φ} aber}{}
\EdL{30}{AK07-J0106-L30}{nicht mit tg \textit{δ}, sondern in den verschiedenen Fällen mit gewissen Faktoren, deren}{}
\EdL{31}{AK07-J0106-L31}{Ursprung nicht recht klar ist, und die auch nur angenäherte Resultate ergeben.}{}
\EdL{32}{AK07-J0106-L32}{\Sperr{Kh.} gibt zuerst die Regel für die Bestimmung der \textit{Ascensio obliqua} für das ganze}{}
\EdL{33}{AK07-J0106-L33}{Zeichen des Widders, er sagt\,: Man multipliziere den Mittagsschatten der Sonne,}{}
\EdL{34}{AK07-J0106-L34}{wenn sie im 1.° des Widders (eigentlich 0.°) sich befindet, mit 114 und teile das}{}
\NoteRule
\EdN{N1}{AK07-J0106-N01}{\textsuperscript{1}) Ibid. II, 65—71.}
\EdN{N2}{AK07-J0106-N02}{\textsuperscript{2}) Ibid. II, 27 u. 187—189.}
\EdN{N3}{AK07-J0106-N03}{\textsuperscript{3}) Ed. \Sperr{Heiberg}, I, 121 ff., Übersetzg. v. K. \Sperr{Manitius}, I, 83—85.}
\EdN{N4}{AK07-J0106-N04}{\textsuperscript{4}) Die Schatten, besonders der Aequinoctial-Schatten (= tg \textit{φ}) kommen übrigens schon im}
\EdN{N5}{AK07-J0106-N05}{\textit{Sūrya-Siddhānta} vor, eine Tafel derselben findet sich allerdings darin nicht, es wird jeweilen sin.}
\EdN{N6}{AK07-J0106-N06}{durch cos. geteilt, oder der Schatten wird mit Hilfe des Gnomons für jeden Fall besonders bestimmt.}
\FrameHead{0107}{p. 75}{0107}{p. 75}
\EdL{1}{AK07-J0107-L01}{Produkt durch 115, so habe man (in Graden) das, was von der Rectascension sub-}{}
\EdL{2}{AK07-J0107-L02}{trahiert, bezw. zu ihr addiert werden müsse, um die \textit{Ascensio obliqua} zu erhalten.}{}
\EdL{3}{AK07-J0107-L03}{Wir geben im folgenden die Berechnung für \textit{Al-Raḳḳa} (geogr. Breite = 36°) nach}{}
\EdL{4}{AK07-J0107-L04}{dieser Regel\,:}{}
\EdLc{5}{AK07-J0107-L05}{cot 54° = tg 36° = 8\textsuperscript{\textit{p}} 44′ (nach Tafel 60)\,; 8\textsuperscript{\textit{p}} 44′ · \Frac{114}{115} = 8\textsuperscript{\textit{p}} 39′\,;}{}
\EdL{6}{AK07-J0107-L06}{dieses setzt \Sperr{Kh.} = 8° 39′ und zieht es von der Rectascension von 30° des Widders,}{}
\EdL{7}{AK07-J0107-L07}{die nach Tafel 59 = 117° 50′ — 90° = 27° 50′ ist, ab, so ergibt sich die \textit{Ascensio obliqua}}{}
\EdL{8}{AK07-J0107-L08}{für das Zeichen des Widders für \textit{Al-Raḳḳa} = 19° 11′. Nach den Tafel \Sperr{Batt.}’s erhält}{}
\EdL{9}{AK07-J0107-L09}{man 19° 21′, \Sperr{Ptol.} (l. c.) hat 19° 12′, da er ungefähr die gleiche Schiefe der Ekliptik}{}
\EdL{10}{AK07-J0107-L10}{annimmt wie \Sperr{Kh.} Berechnet man \textit{DN} nach der Formel \Sperr{Batt.}’s, aber mit den}{}
\EdL{11}{AK07-J0107-L11}{Werten \Sperr{Kh.}’s, so erhält man 8° 37\textsuperscript{1}/\textsubscript{2}′, der Faktor \Frac{114}{115} ist also ein wenig zu gross.}{}
\EdLp{12}{AK07-J0107-L12}{Würde man sin \textit{DN} vollständig mit der Cotangenten-Tafel \Sperr{Kh.}’s berechnen, so}{}
\EdL{13}{AK07-J0107-L13}{fände man zunächst in Tafel 21 die Deklination von 30° des Widders = 11° 39′ 50″}{}
\EdL{14}{AK07-J0107-L14}{= 11° 40′, dazu gehört in Tafel 60 die Tangente 2\textsuperscript{\textit{p}} 28′, also hätte man\,: sin \textit{DN}}{}
\EdL{15}{AK07-J0107-L15}{= 8\textsuperscript{\textit{p}} 44′ · 2\textsuperscript{\textit{p}} 28′\,; dies muss man aber, weil die Tangenten für den Gnomon 12, und}{}
\EdL{16}{AK07-J0107-L16}{die Sinus für den Radius 60 berechnet sind, durch 144 teilen und mit 60 multiplizieren,}{}
\EdL{17}{AK07-J0107-L17}{dies gibt\,:}{}
\EdLc{18}{AK07-J0107-L18}{sin \textit{DN} = \Frac{8\textsuperscript{\textit{p}} 44′ · 2\textsuperscript{\textit{p}} 28′ · 60}{144} = 8\textsuperscript{\textit{p}} 44′ · \Frac{148}{144} = 8\textsuperscript{\textit{p}} 59′\,;}{}
\EdL{19}{AK07-J0107-L19}{zu diesem Sinus gehört in Tafel 58 der Winkel 8° 37′. Um den Sinus zu erhalten,}{}
\EdL{20}{AK07-J0107-L20}{muss man also 8\textsuperscript{\textit{p}} 44′ mit dem Faktor \Frac{148}{144} multiplizieren\,; um sofort den Winkel zu}{}
\EdL{21}{AK07-J0107-L21}{erhalten, multipliziert \Sperr{Kh.} 8\textsuperscript{\textit{p}} 44′ mit dem Faktor \Frac{114}{115}\,; ob und wie dieser aus dem}{}
\EdL{22}{AK07-J0107-L22}{erstern erhalten worden sei, lässt sich schwerlich herausfinden, nach unserer Ansicht}{}
\EdL{23}{AK07-J0107-L23}{ist er auf empirischem Wege gefunden worden.}{}
\EdLp{24}{AK07-J0107-L24}{Will man die Ascensional-Differenz nur für einen bestimmten Teil des Zeichens}{}
\EdL{25}{AK07-J0107-L25}{haben, nimmt man einfach den entsprechenden Teil des Faktors, also z. B. für 15°}{}
\EdL{26}{AK07-J0107-L26}{nur \Frac{57}{115}, dies gibt also dann für Bogen \textit{DN} den Wert 4° 19\textsuperscript{1}/\textsubscript{2}′, also für die \textit{Ascensio}}{}
\EdL{27}{AK07-J0107-L27}{\textit{obliqua} von 15° des Widders den Wert 9° 26\textsuperscript{1}/\textsubscript{2}′ (\Sperr{Batt.} hat 9° 28′).}{}
\EdLp{28}{AK07-J0107-L28}{Für das Zeichen des Stiers hat \Sperr{Kh.} den Faktor \Frac{113}{116}\,; hier liegt jedenfalls ein}{}
\EdL{29}{AK07-J0107-L29}{Fehler vor, die Ascensional-Differenz wird zu gross, ein besserer Wert wäre \Frac{103}{126}. —}{}
\EdL{30}{AK07-J0107-L30}{Für das Zeichen der Zwillinge gibt \Sperr{Kh.} den Faktor \Frac{1}{3}, der ziemlich gut stimmt\,:}{}
\EdL{31}{AK07-J0107-L31}{man erhält für die \textit{Ascensio obliqua} dieses Zeichens mit diesem Faktor den Wert}{}
\EdL{32}{AK07-J0107-L32}{29° 21′, während \Sperr{Batt.} 29° 17′ hat. — Im \textit{Sūrya-Siddhānta} finden wir nichts von}{}
\EdL{33}{AK07-J0107-L33}{diesen Faktoren, dagegen in den Versen 60 u. 61 des 2. Kap.\textsuperscript{1}) die Regel\,:}{}
\EdLc{34}{AK07-J0107-L34}{sin Ascens.-Diff. = \Frac{sin \textit{δ} · tg \textit{φ} · 60}{12 · cos \textit{δ}},}{}
\EdL{35}{AK07-J0107-L35}{die mit der oben gegebenen Formel \Sperr{Batt.}’s übereinstimmt.}{}
\NoteRule
\EdN{N1}{AK07-J0107-N01}{\textsuperscript{1}) \textit{Transl.} \Sperr{Burgess}, p. 232—233. Dieselbe Regel findet sich auch in der \textit{Panchasiddhāntikā},}
\EdN{N2}{AK07-J0107-N02}{ed. G. \Sperr{Thibaut}, p. 27 der Übersetzung.}
\FrameHead{0108}{p. 76}{0108}{p. 76}
\EdLp{1}{AK07-J0108-L01}{Derselbe Faktor, der für den Widder gebraucht wird, gibt auch die Ascensional-}{}
\EdL{2}{AK07-J0108-L02}{Differenz für die Fische, Jungfrau und Wage, der Faktor für den Stier auch die-}{}
\EdL{3}{AK07-J0108-L03}{jenige für den Wassermann, Löwe und Skorpion, der Faktor für die Zwillinge auch}{}
\EdL{4}{AK07-J0108-L04}{die für den Steinbock, Krebs und Schützen. Bei Jungfrau, Wage, Löwe, Skorpion,}{}
\EdL{5}{AK07-J0108-L05}{Krebs und Schütze muss aber die Ascensional-Differenz zur Rectascension addiert}{}
\EdL{6}{AK07-J0108-L06}{werden.}{}
\EdLp{7}{AK07-J0108-L07}{Im Ms. O befindet sich (in unserm Text p. 21) eine zweite Regel für die}{}
\EdL{8}{AK07-J0108-L08}{Bestimmung der \textit{Ascensio obliqua,} die ich, da sie zu Kap. 26 gehört, aber an späterer}{}
\EdL{9}{AK07-J0108-L09}{Stelle eingeschoben ist, mit Kap. 26\textsuperscript{a} bezeichnet habe. Sie stimmt mit der erstern}{}
\EdL{10}{AK07-J0108-L10}{darin überein, dass sie sagt, man bestimme zuerst tg \textit{φ} (den Aequinoctial-Schatten)\,;}{}
\EdL{11}{AK07-J0108-L11}{dann multipliziert sie diesen aber nicht mit den oben genannten Faktoren, sondern}{}
\EdL{12}{AK07-J0108-L12}{verweist auf eine Tafel \textit{differentiarum,} die aber sowohl in O als in C und N fehlt,}{}
\EdL{13}{AK07-J0108-L13}{in der man für jeden beliebigen Grad der Ekliptik die Grösse finde, mit der man}{}
\EdL{14}{AK07-J0108-L14}{tg \textit{φ} multiplizieren müsse, um den Sinus der Ascensional-Differenz zu erhalten\,; diese}{}
\EdL{15}{AK07-J0108-L15}{Tafel muss also die Werte für tg \textit{δ} · \Frac{60}{144} für alle Grade der Ekliptik enthalten haben,}{}
\EdL{16}{AK07-J0108-L16}{und hätte also jedenfalls genauere Werte ergeben, als man sie mit Hilfe der obigen}{}
\EdL{17}{AK07-J0108-L17}{Faktoren erhält.}{}
\EdLc{18}{AK07-J0108-L18}{\textbf{Kap. 27. Bestimmung der Grösse der Zeitstunde\textsuperscript{1}) eines beliebigen Tages für}}{}
\EdLc{19}{AK07-J0108-L19}{\textbf{einen beliebigen Ort.}}{}
\EdLc{20}{AK07-J0108-L20}{(Text p. 20—21.)}{}
\EdLp{21}{AK07-J0108-L21}{Hiefür enthält unser Werk die kurze Regel\,: „Man teile den Tagebogen der}{}
\EdL{22}{AK07-J0108-L22}{Sonne durch 12, so hat man das Gesuchte.“ Dies ist allerdings erst die Länge einer}{}
\EdL{23}{AK07-J0108-L23}{Zeitstunde in Graden\,; um sie in Zeit zu verwandeln, hat man bekanntlich die}{}
\EdL{24}{AK07-J0108-L24}{Beziehungen\,: 15° = 1\textsuperscript{h}, 15′ = 1\textsuperscript{m}, 15″ = 1\textsuperscript{s}. Wie man aber den Tagebogen bestimme,}{}
\EdL{25}{AK07-J0108-L25}{darüber enthalten unsere Tafeln keine Angaben.\textsuperscript{2})}{}
\EdLc{26}{AK07-J0108-L26}{\textbf{Kap. 28—28\textsuperscript{b}. Wie man aus der Höhe der Sonne die beiden Schatten bestimmt}}{}
\EdLc{27}{AK07-J0108-L27}{\textbf{und umgekehrt.\textsuperscript{3})}}{}
\EdLc{28}{AK07-J0108-L28}{(Text p. 21—22, Tafel 60.)}{}
\EdLp{29}{AK07-J0108-L29}{Auch über dieses Kapitel hat schon \Sperr{Björnbo} in der oben (p. 68) zitierten}{}
\EdL{30}{AK07-J0108-L30}{Schrift gesprochen. Der Radius (\textit{Gnomon}) wird also = 12\textsuperscript{\textit{p}} angenommen, bekanntlich}{}
\EdL{31}{AK07-J0108-L31}{hat erst \Sperr{Ḥabash}\textsuperscript{4}), ein Zeitgenosse \Sperr{Batt.}’s, die Cotangenten für den Radius 60}{}
\EdL{32}{AK07-J0108-L32}{berechnet, und dies wahrscheinlich \Sperr{nach} dem Erscheinen der Tafeln \Sperr{Batt.}’s, sonst}{}
\EdL{33}{AK07-J0108-L33}{hätte dieser wohl diese Tafeln benutzt.}{}
\NoteRule
\EdN{N1}{AK07-J0108-N01}{\textsuperscript{1}) Wie bekannt, unterschieden die Griechen und dann auch die Araber Zeit- und Aequatorial-}
\EdN{N2}{AK07-J0108-N02}{stunden, die erstern waren der 12. Teil des Tagebogens, also ungleich je nach der geogr. Breite}
\EdN{N3}{AK07-J0108-N03}{und Jahreszeit, die letztern sind unsere heutigen gleichen Stunden.}
\EdN{N4}{AK07-J0108-N04}{\textsuperscript{2}) Vergl. \Sperr{Batt.} ed. \Sperr{Nallino}, I, 28–29.}
\EdN{N5}{AK07-J0108-N05}{\textsuperscript{3}) Diese Kap., sowie die Tafel fehlen in N.}
\EdN{N6}{AK07-J0108-N06}{\textsuperscript{4}) Vergl. \textit{Abhandlgn. z. Gesch. d. mathem. Wissensch.} X, 1900, p. 12.}
\FrameHead{0109}{p. 77}{0109}{p. 77}
\EdLp{1}{AK07-J0109-L01}{\Sperr{Kh.} nennt die Cotangente \textit{el-ẓill el-mustawī} = der ebene oder gerade Schatten,}{}
\EdL{2}{AK07-J0109-L02}{die Tangente \textit{el-ẓill el-mankūs} (oder \textit{maʿkūs}) = der verkehrte Schatten. Besser sind}{}
\EdL{3}{AK07-J0109-L03}{die Ausdrücke, die \Sperr{Batt.} braucht, nämlich\,: \textit{el-ẓill el-mabsūṭ} = der ausgebreitete,}{}
\EdL{4}{AK07-J0109-L04}{ebene Schatten, und \textit{el-ẓill el-muntaṣib} oder \textit{el-ḳā’im} = der aufgerichtete, senkrechte}{}
\EdL{5}{AK07-J0109-L05}{Schatten\,; sie werden auch \textit{zweiter,} bezw. \textit{erster} Schatten genannt.}{}
\EdLp{6}{AK07-J0109-L06}{Wie \Sperr{Batt.}\textsuperscript{1}), so lehrt auch \Sperr{Kh.} die Tangente und Cotangente aus dem Sinus}{}
\EdL{7}{AK07-J0109-L07}{und Cosinus berechnen, gebraucht also die Formeln\,:}{}
\EdLc{8}{AK07-J0109-L08}{cot \textit{α} = \Frac{cos \textit{α}}{sin \textit{α}} · 12 und tg \textit{α} = \Frac{sin \textit{α}}{cos \textit{α}} · 12.}{}
\EdLp{9}{AK07-J0109-L09}{Für die Bestimmung der Sonnenhöhe \textit{α} aus der Cotangente wird die Formel}{}
\EdL{10}{AK07-J0109-L10}{angegeben\,:}{}
\EdLc{11}{AK07-J0109-L11}{sin (90 — \textit{α}) = \Frac{cot \textit{α} · 60}{\Root{12\textsuperscript{2} + cot\textsuperscript{2}\textit{α}}},}{}
\EdL{12}{AK07-J0109-L12}{d. i. unsere heutige Formel\,: cos \textit{α} = \Frac{cot \textit{α}}{\Root{1 + cot\textsuperscript{2}\textit{α}}},}{}
\EdLp{13}{AK07-J0109-L13}{\Sperr{Batt.} gibt auch noch die Formeln\,: sin \textit{α} = \Frac{1}{\Root{1 + cot\textsuperscript{2}\textit{α}}}, sin \textit{α} = \Frac{tg \textit{α}}{\Root{1 + tg\textsuperscript{2}\textit{α}}},}{}
\EdL{14}{AK07-J0109-L14}{cos \textit{α} = \Frac{1}{\Root{1 + tg\textsuperscript{2}\textit{α}}}, welche bei \Sperr{Kh.} fehlen.}{}
\EdLp{15}{AK07-J0109-L15}{Am Schlusse verweist \Sperr{Kh.}, wie auch \Sperr{Batt.} tut, auf die Cotangententafel, aus}{}
\EdL{16}{AK07-J0109-L16}{der man direkt aus der Höhe den Schatten und umgekehrt bestimmen könne.}{}
\EdLp{17}{AK07-J0109-L17}{Wir finden also die erste Cotangententafel bei \Sperr{Kh.}, und nicht erst bei \Sperr{Batt.},}{}
\EdL{18}{AK07-J0109-L18}{bezw. \Sperr{Ḥabash}. Doch stehen wir entgegen \Sperr{Björnbo}\textsuperscript{2}) davon ab, hieraus weitere}{}
\EdL{19}{AK07-J0109-L19}{Schlüsse zu ziehen\,; denn es ist keineswegs unmöglich, ja sogar nicht unwahr-}{}
\EdL{20}{AK07-J0109-L20}{scheinlich, dass diese Cotangententafel sich nicht im ursprünglichen Werke \Sperr{Kh.}’s}{}
\EdL{21}{AK07-J0109-L21}{befunden habe, sondern erst von \Sperr{Maslama} in seine Bearbeitung aufgenommen}{}
\EdL{22}{AK07-J0109-L22}{worden sei, der ja auch einen Auszug aus den Tafeln \Sperr{Batt.}’s verfasst hat\,; denn die}{}
\EdL{23}{AK07-J0109-L23}{Regeln des Textes sind sehr ähnlich denjenigen \Sperr{Batt.}’s, nur etwas kürzer gefasst. —}{}
\EdL{24}{AK07-J0109-L24}{Was die Zahlenwerte anbetrifft, die wir in die Tafel aufgenommen haben, so gilt hier}{}
\EdL{25}{AK07-J0109-L25}{dasselbe, was wir oben (p. 69) bei Besprechung der Sinustafel gesagt haben.}{}
\EdLp{26}{AK07-J0109-L26}{Wir kommen hier noch auf einen Ausspruch \Sperr{Björnbos} in der zitierten Schrift}{}
\EdL{27}{AK07-J0109-L27}{zu sprechen, er sagt p. 17\,: „\Sperr{Abū’l-Wefā} endlich tat den letzten und grössten}{}
\EdL{28}{AK07-J0109-L28}{Schritt. Er warf beides, die griechische Sechzig-Teilung und die indische Zwölf-}{}
\EdL{29}{AK07-J0109-L29}{Teilung über Bord, setzte den Radius = 1, machte, ohne die Namen zu ändern, den}{}
\EdL{30}{AK07-J0109-L30}{ebenen Schatten zu dem, was wir verstehen unter Cotangens, den aufrechten Schatten}{}
\EdL{31}{AK07-J0109-L31}{zu dem, was wir unter Tangens verstehen, und schloss damit die Entwicklungs-}{}
\EdL{32}{AK07-J0109-L32}{geschichte der beiden Funktionen ab.“ Dies ist unserer Ansicht nach etwas zu weit}{}
\EdL{33}{AK07-J0109-L33}{gegangen\,; wenn auch \Sperr{Abū’l-Wefā} an einer Stelle seines Almagestes sagt\,: „hieraus}{}
\EdL{34}{AK07-J0109-L34}{ist klar, dass, wenn man den Radius = 1 annimmt, das Verhältnis des Sinus zum}{}
\EdL{35}{AK07-J0109-L35}{Cosinus gleich der Tangente, das Verhältnis des Cosinus zum Sinus gleich der}{}
\NoteRule
\EdN{N1}{AK07-J0109-N01}{\textsuperscript{1}) Ed. \Sperr{Nallino}, I, 22—23. u. 180—182.}
\EdN{N2}{AK07-J0109-N02}{\textsuperscript{2}) \textit{Al-Chwarizmi’s Trigonometriske Tavler}, p. 16–17.}
\FrameHead{0110}{p. 78}{0110}{p. 78}
\EdL{1}{AK07-J0110-L01}{Cotangente ist“\textsuperscript{1}), so ist es doch so viel als gewiss, dass weder er selbst noch irgend}{}
\EdL{2}{AK07-J0110-L02}{einer seiner Nachfolger diese Vereinfachung praktisch durchgeführt hat. Allerdings}{}
\EdL{3}{AK07-J0110-L03}{sagt \Sperr{Carra de Vaux}, \Sperr{Abū’l-Wefā} habe eine Tafel für Tangens und Cotangens}{}
\EdL{4}{AK07-J0110-L04}{sowohl für den Radius 1 als auch für den Radius 60 aufgestellt, allein dies steht nicht}{}
\EdL{5}{AK07-J0110-L05}{ganz fest, die Tafeln sind leider verloren gegangen, aber wenn auch solche Tafeln}{}
\EdL{6}{AK07-J0110-L06}{einmal vorhanden waren, so wurden sie nachher gewiss nie mehr benutzt.}{}
\EdLc{7}{AK07-J0110-L07}{\textbf{Kap. 29. Über elbuht und seine Auffindung.}}{}
\EdLc{8}{AK07-J0110-L08}{(Text p. 22, Tafeln 61—66.)}{}
\EdLp{9}{AK07-J0110-L09}{Unter \textit{el-buht} verstanden die arabischen und besonders die persischen\textsuperscript{2}) Astro-}{}
\EdL{10}{AK07-J0110-L10}{nomen die wahre oder ungleiche Bewegung der Sonne, des Mondes und der Planeten}{}
\EdL{11}{AK07-J0110-L11}{in einem Tage\,; das Wort kommt vom indischen \textit{bhukti,} das wörtlich „der Genuss,}{}
\EdL{12}{AK07-J0110-L12}{die Freude“ bedeutet\textsuperscript{3}), und das im Persischen in \textit{buht} überging, und so von den}{}
\EdL{13}{AK07-J0110-L13}{Arabern übernommen wurde\textsuperscript{4}).}{}
\EdLp{14}{AK07-J0110-L14}{Bekanntlich ist die Bewegung der Planeten eine ungleichförmige, im Perigeum}{}
\EdL{15}{AK07-J0110-L15}{schneller als im Apogeum, wir finden also bei 180° Anomalie die grösste, bei 0° die}{}
\EdL{16}{AK07-J0110-L16}{geringste wahre Bewegung. Unsere Tafeln 61—66 sind also leicht zu verstehen\,:}{}
\EdL{17}{AK07-J0110-L17}{Die Kolonne \textit{semitae numerorum} überschrieben enthält die Anomalien von 0°—180°}{}
\EdL{18}{AK07-J0110-L18}{(bezw. 360°—180°), die zweite Kolonne die wahre Bewegung der Sonne in einer Stunde,}{}
\EdL{19}{AK07-J0110-L19}{die dritte Kolonne die des Mondes in einer Stunde\textsuperscript{5}). Für die Sonne ist das Mini-}{}
\EdL{20}{AK07-J0110-L20}{mum der Bewegung 2′ 22″, das Maximum 2′ 34″, für den Mond dagegen 30′ 12″,}{}
\EdL{21}{AK07-J0110-L21}{bezw. 35′ 40″\,; \Sperr{Batt.} (l. c.) hat die Werte 2′ 23″, 2′ 33″ für die Sonne, 30′ 18″, 36′ 4″}{}
\EdL{22}{AK07-J0110-L22}{für den Mond.}{}
\EdLc{23}{AK07-J0110-L23}{\textbf{Kap. 30. u. 30\textsuperscript{a}. Über die scheinbaren Sonnen- und Monddurchmesser.}}{}
\EdLc{24}{AK07-J0110-L24}{(Text p. 22—23, Tafeln 61—66.)}{}
\EdLp{25}{AK07-J0110-L25}{Die scheinbaren Sonnen- und Monddurchmesser nennt \Sperr{Kh.} \textit{quantitas rotae solis,}}{}
\EdL{26}{AK07-J0110-L26}{bezw. \textit{rotae lunae.} Diese Durchmesser ändern sich natürlich mit der Entfernung}{}
\EdL{27}{AK07-J0110-L27}{dieser Himmelskörper von der Erde. Während nun \Sperr{Batt.} sich damit begnügt, die}{}
\EdL{28}{AK07-J0110-L28}{maximale, mittlere und minimale Grösse dieser Durchmesser anzugeben\textsuperscript{6}), gibt sich}{}
\EdL{29}{AK07-J0110-L29}{\Sperr{Kh.} die Mühe, diese Grössen für jeden Grad der Anomalie von 0°—180° zu bestimmen.}{}
\EdL{30}{AK07-J0110-L30}{In den Kolonnen 4 u. 5 der Tafeln 61—66 gibt er die scheinbaren Halbmesser (\textit{dimi-}}{}
\EdL{31}{AK07-J0110-L31}{\textit{dium rotae solis vel lunae}) von Sonne und Mond für die verschiedenen Anomalien.}{}
\EdL{32}{AK07-J0110-L32}{Es treten hier wiederum konstante Faktoren auf, mit denen er die wahren \textit{täglichen}}{}
\NoteRule
\EdN{N1}{AK07-J0110-N01}{\textsuperscript{1}) Vergl. \Sperr{Carra de Vaux}, \textit{L’Almageste d’Aboul-Wéfa}, im \textit{Journal asiatique}, Série 8, T. XIX,}
\EdN{N2}{AK07-J0110-N02}{p. 420.}
\EdN{N3}{AK07-J0110-N03}{\textsuperscript{2}) \Sperr{Ulūgh Beg} gebraucht den Ausdruck in seinen Tafeln.}
\EdN{N4}{AK07-J0110-N04}{\textsuperscript{3}) \textit{Sūrya-Siddhānta}, transl. \Sperr{Burgess}, p. 158\,: „\textit{When a planet has} \textit{\Sperr{enjoyed}} \textit{the whole troop}}
\EdN{N5}{AK07-J0110-N05}{\textit{of asterisms, it has made a complete revolution.}}
\EdN{N6}{AK07-J0110-N06}{\textsuperscript{4}) \Sperr{Batt.} kennt das Wort nicht, oder gebraucht es nicht, vergl. I, XLIII und II, 324.}
\EdN{N7}{AK07-J0110-N07}{\textsuperscript{5}) Für die fünf Planeten fehlen die entsprechenden Tafeln, wie auch bei \Sperr{Batt.}, für Sonne}
\EdN{N8}{AK07-J0110-N08}{und Mond gibt letzterer sie nur für Anomalien von je 6° zu 6° an, vergl. II, 88.}
\EdN{N9}{AK07-J0110-N09}{\textsuperscript{6}) Ed. \Sperr{Nallino}, I, 58.}
\FrameHead{0111}{p. 79}{0111}{p. 79}
\EdL{1}{AK07-J0111-L01}{Bewegungen von Sonne und Mond multipliziert, um die scheinbaren Durchmesser}{}
\EdL{2}{AK07-J0111-L02}{zu erhalten, und zwar ist der Faktor für die Sonne = 33′ = \CaseFrac{33}{60} = \CaseFrac{11}{20}, der für den}{}
\EdL{3}{AK07-J0111-L03}{Mond = 2′ 25″\textsuperscript{1}) = 2\CaseFrac{25}{60} ∶ 60 = \CaseFrac{145}{3600} = \CaseFrac{29}{720}. \Sperr{Batt.}\textsuperscript{2}) gibt für die Berechnung des}{}
\EdL{4}{AK07-J0111-L04}{Monddurchmessers eine ähnliche Regel\,; er sagt, man müsse die wahre \textit{stündliche}}{}
\EdL{5}{AK07-J0111-L05}{Bewegung des Mondes mit 5\CaseFrac{7}{8} multiplizieren und durch 6 teilen (also mit \CaseFrac{47}{48} multip-}{}
\EdL{6}{AK07-J0111-L06}{lizieren), um den scheinbaren Monddurchmesser zu erhalten\,; reduzieren wir den}{}
\EdL{7}{AK07-J0111-L07}{Faktor \Sperr{Kh.}’s auf die \textit{stündliche} wahre Bewegung, so wird er = \CaseFrac{29}{30}, der Unterschied}{}
\EdL{8}{AK07-J0111-L08}{beider Faktoren ist = \CaseFrac{1}{80}. Geben wir für Sonne und Mond je ein \textit{Beispiel\,:}}{}
\EdLp{9}{AK07-J0111-L09}{Die tägliche wahre Bewegung der Sonne für 90° der Anomalie ist nach Tafel 63}{}
\EdL{10}{AK07-J0111-L10}{= 2′ 28″ · 24, dieses mit \CaseFrac{11}{20} multipliziert, gibt für den Sonnendurchmesser 32′ 33\CaseFrac{1}{2}″,}{}
\EdL{11}{AK07-J0111-L11}{für den Halbmesser also 16′ 16\CaseFrac{3}{4}″\,; die Tafel hat 16′ 17″\textsuperscript{3}).}{}
\EdLp{12}{AK07-J0111-L12}{Die tägliche wahre Bewegung des Mondes für 90° der Anomalie ist nach der-}{}
\EdL{13}{AK07-J0111-L13}{selben Tafel 32′ 56″ · 24, dieses mit \CaseFrac{29}{720} multipliziert, gibt für den Monddurchmesser}{}
\EdL{14}{AK07-J0111-L14}{31′ 50″, für den Halbmesser also 15′ 55″\,; die Tafel hat 15′ 58″.}{}
\EdLp{15}{AK07-J0111-L15}{Man wird übrigens nicht immer den richtigen Wert der täglichen wahren}{}
\EdL{16}{AK07-J0111-L16}{Bewegung von Sonne und Mond finden, wenn man die stündliche mit 24 multipliziert,}{}
\EdL{17}{AK07-J0111-L17}{denn die letztern Werte sind in den Tafeln teils ab-, teils aufgerundet, und durch}{}
\EdL{18}{AK07-J0111-L18}{die Multiplikation mit 24 wird der Fehler soviel mal grösser. Der Text bemerkt}{}
\EdL{19}{AK07-J0111-L19}{daher auch, man finde die tägliche Bewegung, indem man den wahren Ort des}{}
\EdL{20}{AK07-J0111-L20}{Gestirns für den Mittag eines bestimmten Tages und ebenso für den Mittag des fol-}{}
\EdL{21}{AK07-J0111-L21}{genden Tages bestimme, der Unterschied beider Örter sei der \textit{buht} dieses Gestirns.}{}
\EdLp{22}{AK07-J0111-L22}{Die oben genannten Regeln \Sperr{Kh.}’s zur Bestimmung der scheinbaren Durchmesser}{}
\EdL{23}{AK07-J0111-L23}{von Sonne und Mond scheinen den indischen Werken entnommen zu sein. Der}{}
\EdL{24}{AK07-J0111-L24}{\textit{Sūrya-Siddhānta} enthält allerdings dieselben nicht in dieser Form, und seine Werte}{}
\EdL{25}{AK07-J0111-L25}{für die mittlern Durchmesser von Sonne und Mond (32′ 25″ bezw. 32′)\textsuperscript{4}) stimmen}{}
\EdL{26}{AK07-J0111-L26}{auch nicht ganz mit denen \Sperr{Kh.}’s (32′ 34″ und 31′ 50″), die Unterschiede sind aller-}{}
\EdL{27}{AK07-J0111-L27}{dings gering. Aber \Sperr{Al-Bīrūnī} hat in seinem Werk über Indien\textsuperscript{5}) eine Stelle, die}{}
\EdL{28}{AK07-J0111-L28}{uns beweist, dass \Sperr{Kh.} in diesem Punkte sich auf indische Quellen gestützt hat\,; er}{}
\EdL{29}{AK07-J0111-L29}{sagt, dass die Methode der Berechnung der scheinbaren Durchmesser von Sonne}{}
\EdL{30}{AK07-J0111-L30}{und Mond, wie sie der \textit{Khaṇḍakhādyaka}\textsuperscript{6}) des \Sperr{Brahmagupta} enthalte, die gleiche}{}
\EdL{31}{AK07-J0111-L31}{sei, wie sie in den Tafeln \Sperr{Kh.}’s gefunden werde. In der Tat kannten die Araber}{}
\EdL{32}{AK07-J0111-L32}{schon früh eine Übersetzung dieses Werkes unter dem Namen \textit{Al-Arkand,} die später}{}
\EdL{33}{AK07-J0111-L33}{von \Sperr{Al-Bīrūnī} selbst verbessert worden ist, wie er im Verzeichnis seiner Schriften}{}
\EdL{34}{AK07-J0111-L34}{angibt. Es unterliegt also wohl keinem Zweifel, dass \Sperr{Kh.} auch dieses Werk \Sperr{Brahma-}}{}
\EdL{35}{AK07-J0111-L35}{\Sperr{guptas} gekannt und benutzt habe. (Vergl. weiter unten den Schluss von Kap. 33).}{}
\NoteRule
\EdN{N1}{AK07-J0111-N01}{\textsuperscript{1}) Der Text hat unrichtig 2′ 16″.}
\EdN{N2}{AK07-J0111-N02}{\textsuperscript{2}) Ibid. I, 97.}
\EdN{N3}{AK07-J0111-N03}{\textsuperscript{3}) Bei andern Werten der Anomalie stimmt freilich die Rechnung nicht immer so schön.}
\EdN{N4}{AK07-J0111-N04}{\textsuperscript{4}) Transl. \Sperr{Burgess}, p. 269—272.}
\EdN{N5}{AK07-J0111-N05}{\textsuperscript{5}) \textit{Alberunis India}, English Edition by E. C. \Sperr{Sachau}, London 1888 und 1910, II, 79 (arab.}
\EdN{N6}{AK07-J0111-N06}{Text, London 1887, p. 241).}
\EdN{N7}{AK07-J0111-N07}{\textsuperscript{6}) Ein Werk, das \Sperr{Brahmagupta} nach seinem \textit{Siddhānta}, und auf etwas andern Grundsätzen}
\EdN{N8}{AK07-J0111-N08}{aufgebaut, geschrieben hat.}
\FrameHead{0112}{p. 80}{0112}{p. 80}
\EdLp{1}{AK07-J0112-L01}{Interessant ist das lateinische Wort \textit{stupor,} mit dem \Sperr{Athelhard} das persische}{}
\EdL{2}{AK07-J0112-L02}{\textit{buht} wiedergibt, wir können es wohl am bestem mit „Stocken“, „Anhalten“ (in der}{}
\EdL{3}{AK07-J0112-L03}{Bewegung) übersetzen, denn wenn das Gestirn bald schneller, bald langsamer geht,}{}
\EdL{4}{AK07-J0112-L04}{so stockt es gleichsam bisweilen in seiner Bewegung\textsuperscript{1})\,; bei den Planeten passt}{}
\EdL{5}{AK07-J0112-L05}{dieser Ausdruck allerdings besser als bei Sonne und Mond.}{}
\EdLp{6}{AK07-J0112-L06}{Über die Bedeutung der 6. u. 7. Kolonne der Tafeln 61—66, die überschrieben}{}
\EdL{7}{AK07-J0112-L07}{sind mit \textit{primum dimidium rotae draconis} und \textit{secundum dimidium rotae draconis,}}{}
\EdL{8}{AK07-J0112-L08}{komme ich am Ende von Kap. 33 zu sprechen.}{}
\EdLc{9}{AK07-J0112-L09}{\textbf{Die Tafeln der Tagesgleichung (Zeitgleichung).}}{}
\EdLc{10}{AK07-J0112-L10}{(Tafeln 67 und 68.)}{}
\EdLp{11}{AK07-J0112-L11}{Der Text des vorliegenden Werkes scheint, besonders gegen das Ende hin,}{}
\EdL{12}{AK07-J0112-L12}{mangelhaft zu sein, für diese Tafeln (67 u. 68) z. B. findet sich im Text kein erklä-}{}
\EdL{13}{AK07-J0112-L13}{rendes Kapitel. Die Tafeln entsprechen den Kolonnen bei \Sperr{Batt.}\textsuperscript{2}), die mit \textit{Aequatio}}{}
\EdL{14}{AK07-J0112-L14}{\textit{nychthemerōn} überschrieben sind, nur beginnen die Tafeln \Sperr{Batt.}’s mit dem 1. Grad}{}
\EdL{15}{AK07-J0112-L15}{des Steinbocks, bei \Sperr{Kh.} mit dem 1. Grad des Widders. Ein weiterer Unterschied}{}
\EdL{16}{AK07-J0112-L16}{ist der, dass \Sperr{Batt.} die \textit{Zeitgleichung} (also Unterschied zwischen \textit{wahrer} und \textit{mittlerer}}{}
\EdL{17}{AK07-J0112-L17}{Zeit) in Graden und Minuten, \Sperr{Kh.} dagegen in Zeit-Minuten und Sekunden angibt.}{}
\EdL{18}{AK07-J0112-L18}{Sie stimmen dagegen darin überein, dass sie die Zeitgleichung Null annehmen, wenn}{}
\EdL{19}{AK07-J0112-L19}{die wahre Zeit am meisten hinter der mittlern nachgeht, zu allen andern Zeiten muss}{}
\EdL{20}{AK07-J0112-L20}{also dann die wahre Zeit der mittlern vorgehen, d. h. die \textit{Zeitgleichung,} die die Tafel}{}
\EdL{21}{AK07-J0112-L21}{angibt, muss stets \textit{subtrahiert} werden\textsuperscript{3}).}{}
\EdLp{22}{AK07-J0112-L22}{\textit{Beispiel\,:} Die Zeitgleichung ist \textit{Null,} wenn die Sonne im 22.° des Wassermanns}{}
\EdL{23}{AK07-J0112-L23}{steht (\Sperr{Batt.} hat im 18.° u. 19.° Null, im 22.° 0° 1′)\,; die arabischen Astronomen nahmen}{}
\EdL{24}{AK07-J0112-L24}{also bei diesem Stand der Sonne die wahre Zeit gleich der mittlern an. Wenn die}{}
\EdL{25}{AK07-J0112-L25}{Sonne im 9.° bis 10.° des Skorpions steht, so erreicht die Zeitgleichung ihr \textit{Maximum}}{}
\EdL{26}{AK07-J0112-L26}{mit 34\textsuperscript{m} 28\textsuperscript{s}, d. h. die wahre Zeit geht der mittlern um 34\textsuperscript{m} 28\textsuperscript{s} vor (\Sperr{Batt.} hat 7° 54′}{}
\EdL{27}{AK07-J0112-L27}{= 31\textsuperscript{m} 36\textsuperscript{s}). Vergleichen wir diese Werte mit den heutigen, so entspricht der ara-}{}
\EdL{28}{AK07-J0112-L28}{bischen Zeitgleichung Null unsere Zeitgleichung von + 14\textsuperscript{m} 31\textsuperscript{s} am 12. Febr., und der}{}
\EdL{29}{AK07-J0112-L29}{arabischen Zeitgleichung 34\textsuperscript{m} 28\textsuperscript{s} unsere Zeitgleichung von — 16\textsuperscript{m} 18\textsuperscript{s} am 18. Novbr.\,;}{}
\EdL{30}{AK07-J0112-L30}{der Unterschied dieser beiden ist 30\textsuperscript{m} 49\textsuperscript{s}, also ist der Wert \Sperr{Batt.}’s der Wahrheit}{}
\EdL{31}{AK07-J0112-L31}{näher als derjenige \Sperr{Kh.}’s.}{}
\EdLp{32}{AK07-J0112-L32}{Wem sich \Sperr{Kh.} in dieser Frage angeschlossen hat, können wir nicht entscheiden\,;}{}
\EdL{33}{AK07-J0112-L33}{das Maximum der Zeitgleichung bei \Sperr{Kh.} stimmt besser mit dem bei \Sperr{Ptol.} (33\textsuperscript{m} 30\textsuperscript{s})}{}
\EdL{34}{AK07-J0112-L34}{als mit dem bei \Sperr{Batt.} Im \textit{Sūrya-Siddhānta} haben wir keine Werte der Zeitgleichung}{}
\EdL{35}{AK07-J0112-L35}{gefunden\textsuperscript{4}). — Diese Tafeln der Zeitgleichung werden am Schlusse des folgenden}{}
\NoteRule
\EdN{N1}{AK07-J0112-N01}{\textsuperscript{1}) Man vergleiche das \Sperr{Ovid}’sche\,: \textit{stupuit Ixionis orbis}, oder auch\,: \textit{stupuerunt verba palato}}
\EdN{N2}{AK07-J0112-N02}{(\Sperr{Georges}’ latein. Wörterbuch).}
\EdN{N3}{AK07-J0112-N03}{\textsuperscript{2}) Ed. \Sperr{Nallino}, II, p. 61—64.}
\EdN{N4}{AK07-J0112-N04}{\textsuperscript{3}) Vergl. hiefür \Sperr{Batt.}, ed. \Sperr{Nallino}, I, p. 221—222 (nach \Sperr{Schiaparelli}).}
\EdN{N5}{AK07-J0112-N05}{\textsuperscript{4}) Vergl. \textit{Sūrya-Siddhānta}, transl. \Sperr{Burgess}, p. 220—221.}
\FrameHead{0113}{p. 81}{0113}{p. 81}
\EdL{1}{AK07-J0113-L01}{Kapitels erwähnt, wo es sich um die Umrechnung von Äquatorialstunden in ungleiche}{}
\EdL{2}{AK07-J0113-L02}{Stunden und umgekehrt handelt.}{}
\EdLc{3}{AK07-J0113-L03}{\textbf{Kap. 31. Über Konjunktion und Opposition von Sonne und Mond.}}{}
\EdLc{4}{AK07-J0113-L04}{(Text p. 23—25, Tafeln 69—72.)}{}
\EdLp{5}{AK07-J0113-L05}{Mit Hilfe dieser Tafeln sollen nach den Regeln des Textes die Zeiten der Kon-}{}
\EdL{6}{AK07-J0113-L06}{junktionen und Oppositionen von Sonne und Mond bestimmt werden. Pag. 25 des}{}
\EdL{7}{AK07-J0113-L07}{Textes finden wir die Bemerkung\,: \textit{Haec autem omnia pro situ Cordubae dicta sunt.}}{}
\EdL{8}{AK07-J0113-L08}{Dies beweist, dass \Sperr{Maslama al-Madjrīṭī} diese Tafeln entweder von sich aus hin-}{}
\EdL{9}{AK07-J0113-L09}{zugefügt, oder aus denen \Sperr{Kh.}’s für den Meridian von Cordova umgerechnet hat.}{}
\EdL{10}{AK07-J0113-L10}{Ganz ähnlich eingerichtete, aber in den Zahlenwerten etwas abweichende Tafeln}{}
\EdL{11}{AK07-J0113-L11}{finden sich bei \Sperr{Batt.}\textsuperscript{1})\,; sie enthalten aber nur die \textit{anni collecti,} es fehlen die \textit{anni}}{}
\EdL{12}{AK07-J0113-L12}{\textit{plani} (\textit{singuli}) und die \textit{menses}\,; \Sperr{Nallino} hat sie mit Recht zu den unechten Tafeln}{}
\EdL{13}{AK07-J0113-L13}{gerechnet. \Sperr{Batt.} hat für den Meridian von \textit{Al-Raḳḳa} Tafeln der Syzygien in römi-}{}
\EdL{14}{AK07-J0113-L14}{schen Jahren aufgestellt\textsuperscript{2}), die nach dem gleichen Prinzip eingerichtet sind wie die}{}
\EdL{15}{AK07-J0113-L15}{unsrigen, nur schreiten die \textit{anni collecti} um je 24 Jahre vorwärts statt 30, und sie}{}
\EdL{16}{AK07-J0113-L16}{beginnen mit dem 1. \textit{Ādhār} 879 der Seleucid. Aera = 1. März 568 n. Chr.}{}
\EdLp{17}{AK07-J0113-L17}{Die Bestimmung einer \textit{Konjunktion} oder \textit{Opposition} von Sonne und Mond beruht}{}
\EdL{18}{AK07-J0113-L18}{natürlich in erster Linie auf der Berechnung der wahren Längen dieser Gestirne für}{}
\EdL{19}{AK07-J0113-L19}{den betreffenden Zeitpunkt. Findet man diese Längen gleich, so sind diese beiden}{}
\EdL{20}{AK07-J0113-L20}{Gestirne dannzumal in Konjunktion, sind sie um 180° verschieden, so sind sie in}{}
\EdL{21}{AK07-J0113-L21}{Opposition. Stimmen sie nicht ganz überein, so kann man aus dem Unterschied}{}
\EdL{22}{AK07-J0113-L22}{der Längen, und aus den wahren stündlichen Bewegungen der beiden Gestirne auf die}{}
\EdL{23}{AK07-J0113-L23}{Zeit schliessen, die es noch braucht, bis sie in die wahre Konjunktion bezw. Oppo-}{}
\EdL{24}{AK07-J0113-L24}{sition kommen. Die Begründung dieser Korrektur ist bei \Sperr{Batt.}\textsuperscript{3}) ziemlich kompliziert}{}
\EdL{25}{AK07-J0113-L25}{dargestellt, er gibt aber auch eine einfachere Darstellung\textsuperscript{4}), die mit derjenigen des}{}
\EdL{26}{AK07-J0113-L26}{\Sperr{Ptol.}\textsuperscript{5}) übereinstimmt, und der auch \Sperr{Kh.} folgt. — Die Herstellung der Tafeln sowohl}{}
\EdL{27}{AK07-J0113-L27}{als auch die Berechnung der Konjunktionen und Oppositionen mit Hilfe derselben}{}
\EdL{28}{AK07-J0113-L28}{ist nach der muhammedanischen Zeitrechnung weit einfacher als nach der seleuci-}{}
\EdL{29}{AK07-J0113-L29}{dischen oder julianischen, weil die Konjunktionen auf den letzten oder vorletzten}{}
\EdL{30}{AK07-J0113-L30}{Tag des Monats, die Oppositionen stets in dessen Mitte fallen.}{}
\EdLp{31}{AK07-J0113-L31}{Das Prinzip, auf dem die Einrichtung der Tafeln und besonders der 3. Kolonne,}{}
\EdL{32}{AK07-J0113-L32}{\textit{Dies et horae} oder nur \textit{Horae} überschrieben, beruht, ist das, die Zeit zu bestimmen,}{}
\EdL{33}{AK07-J0113-L33}{welche von einem bestimmten Zeitpunkt, bezw. von einer bestimmten Stellung der}{}
\EdL{34}{AK07-J0113-L34}{Gestirne (\textit{Radix}) an nötig ist bis zu der gesuchten Konjunktion, d. h. bis die beiden}{}
\EdL{35}{AK07-J0113-L35}{Gestirne die gleiche Länge haben\,; so wird in Tafel 69 z. B., in der Kolonne \textit{Dies et}}{}
\EdL{36}{AK07-J0113-L36}{\textit{horae} überschrieben, die Zahl 29\textsuperscript{d} 0\textsuperscript{h} 14\textsuperscript{m} 46\textsuperscript{s}, die der Zahl 31 in der Kolonne \textit{Anni}}{}
\NoteRule
\EdN{N1}{AK07-J0113-N01}{\textsuperscript{1}) Ed. \Sperr{Nallino}, II, 306—307.}
\EdN{N2}{AK07-J0113-N02}{\textsuperscript{2}) Ibid. II, 84—86.}
\EdN{N3}{AK07-J0113-N03}{\textsuperscript{3}) Ibid. I, 93—94\,; vergl. dazu die Erklärungen von \Sperr{Schiaparelli}, I, 273—274.}
\EdN{N4}{AK07-J0113-N04}{\textsuperscript{4}) Ibid. I, 94, Z. 26—34.}
\EdN{N5}{AK07-J0113-N05}{\textsuperscript{5}) Ed. \Sperr{Heiberg}, I, 472—475, Übers. v. K. \Sperr{Manitius}, I, 342—349.}
\FrameHead{0114}{p. 82}{0114}{p. 82}
\EdL{1}{AK07-J0114-L01}{\textit{collecti} überschrieben entspricht, bedeuten, dass es nach Verfluss von 30 ganzen}{}
\EdL{2}{AK07-J0114-L02}{arabischen Jahren\textsuperscript{1}) noch 29\textsuperscript{d} 0\textsuperscript{h} 14\textsuperscript{m} 46\textsuperscript{s} bedarf bis zur nächsten Konjunktion von}{}
\EdL{3}{AK07-J0114-L03}{Sonne und Mond\,; ebenso entsprechend in den andern Tafeln für die einzelnen Jahre}{}
\EdL{4}{AK07-J0114-L04}{und Monate, z. B. nach Ablauf des 1. Jahres nach dem 30. braucht es bis zur nächsten}{}
\EdL{5}{AK07-J0114-L05}{Konjunktion nach Tafel 71 noch 8\textsuperscript{h} 48\textsuperscript{m} 7\textsuperscript{s}. — Die 4. Kolonne gibt die mittlere Länge}{}
\EdL{6}{AK07-J0114-L06}{von Sonne und Mond jeweilen für den betreffenden Zeitpunkt an, die 5. Kolonne}{}
\EdL{7}{AK07-J0114-L07}{enthält die Anomalie des Mondes, und die 6. das Argument der Breite des Mondes}{}
\EdL{8}{AK07-J0114-L08}{für jenen Zeitpunkt. Die vorderste Kolonne enthält die \textit{Signa} (oder \textit{Notae}) der be-}{}
\EdL{9}{AK07-J0114-L09}{treffenden Jahre, bezw. Monate\,; diese wie die letzte (6.) Kolonne brauchen wir zur}{}
\EdL{10}{AK07-J0114-L10}{Bestimmung der Konjunktionen und Oppositionen nicht.}{}
\EdLp{11}{AK07-J0114-L11}{Wir geben im folgenden ein \textit{Beispiel} einer Konjunktionsberechnung nach den}{}
\EdL{12}{AK07-J0114-L12}{Regeln des Textes, obgleich es hier vielleicht mehr als an irgend einer andern Stelle}{}
\EdL{13}{AK07-J0114-L13}{notwendig gewesen wäre, von diesem schlecht übersetzten, oder uns wenigstens in}{}
\EdL{14}{AK07-J0114-L14}{sehr schlechtem Latein überlieferten Kapitel eine verständliche deutsche Übersetzung}{}
\EdL{15}{AK07-J0114-L15}{zu geben. Wir wählen dafür den Zeitpunkt der Sonnenfinsternis, die \Sperr{Batt.} im 30.}{}
\EdL{16}{AK07-J0114-L16}{Kapitel\textsuperscript{2}) anführt, und deren Elemente \Sperr{Schiaparelli} nach den Tafeln \Sperr{Batt.}’s nach-}{}
\EdL{17}{AK07-J0114-L17}{gerechnet hat\textsuperscript{3}).}{}
\EdLp{18}{AK07-J0114-L18}{Zu welcher Zeit hat am Schlusse des Monats \textit{Muḥarram} d. J. 288 d. H. (= 22./23.}{}
\EdL{19}{AK07-J0114-L19}{Jan. 901) die \textit{Konjunktion} von Sonne und Mond (der \textit{Neumond}) stattgefunden\,?}{}
\EdLp{20}{AK07-J0114-L20}{Nach den Tafeln 69—72 findet man\,:}{}
\EdL{21}{AK07-J0114-L21}{\parbox[t]{\dimexpr0.4900\linewidth-0.5em\relax}{\raggedright\small\mbox{}}\hspace{0.5em}\parbox[t]{\dimexpr0.1633\linewidth-0.5em\relax}{\raggedright\small\mbox{}Longit. media}\hspace{0.5em}\parbox[t]{\dimexpr0.1633\linewidth-0.5em\relax}{\raggedright\small\mbox{}Anomalia}\hspace{0.5em}\parbox[t]{\dimexpr0.1633\linewidth-0.5em\relax}{\raggedright\small\mbox{}Argumentum}\hspace{0.5em}}{}
\EdL{22}{AK07-J0114-L22}{\parbox[t]{\dimexpr0.3267\linewidth-0.5em\relax}{\raggedright\small\mbox{}}\hspace{0.5em}\parbox[t]{\dimexpr0.1633\linewidth-0.5em\relax}{\raggedright\small\mbox{}Dies et horae}\hspace{0.5em}\parbox[t]{\dimexpr0.1633\linewidth-0.5em\relax}{\raggedright\small\mbox{}solis et lunae}\hspace{0.5em}\parbox[t]{\dimexpr0.1633\linewidth-0.5em\relax}{\raggedright\small\mbox{}lunae}\hspace{0.5em}\parbox[t]{\dimexpr0.1633\linewidth-0.5em\relax}{\raggedright\small\mbox{}latitud. lunae}\hspace{0.5em}}{}
\EdL{23}{AK07-J0114-L23}{\parbox[t]{\dimexpr0.3267\linewidth-0.5em\relax}{\raggedright\small\mbox{}Am Ende des 270. J.}\hspace{0.5em}\parbox[t]{\dimexpr0.1633\linewidth-0.5em\relax}{\raggedright\small\mbox{}29\textsuperscript{d} 2\textsuperscript{h} 4\textsuperscript{m} 51\textsuperscript{s}}\hspace{0.5em}\parbox[t]{\dimexpr0.1633\linewidth-0.5em\relax}{\raggedright\small\mbox{}123° 49′ 42″}\hspace{0.5em}\parbox[t]{\dimexpr0.1633\linewidth-0.5em\relax}{\raggedright\small\mbox{}252° 31′ 18″}\hspace{0.5em}\parbox[t]{\dimexpr0.1633\linewidth-0.5em\relax}{\raggedright\small\mbox{}30° 19′ 50″}\hspace{0.5em}}{}
\EdL{24}{AK07-J0114-L24}{\parbox[t]{\dimexpr0.3267\linewidth-0.5em\relax}{\raggedright\small\mbox{}für 17 Jahre}\hspace{0.5em}\parbox[t]{\dimexpr0.1633\linewidth-0.5em\relax}{\raggedright\small\mbox{}0 5 43 50}\hspace{0.5em}\parbox[t]{\dimexpr0.1633\linewidth-0.5em\relax}{\raggedright\small\mbox{}177 30 35}\hspace{0.5em}\parbox[t]{\dimexpr0.1633\linewidth-0.5em\relax}{\raggedright\small\mbox{}226 38 40}\hspace{0.5em}\parbox[t]{\dimexpr0.1633\linewidth-0.5em\relax}{\raggedright\small\mbox{}136 48 10}\hspace{0.5em}}{}
\EdL{25}{AK07-J0114-L25}{\parbox[t]{\dimexpr0.3267\linewidth-0.5em\relax}{\raggedright\small\mbox{}für Muḥarram}\hspace{0.5em}\parbox[t]{\dimexpr0.1633\linewidth-0.5em\relax}{\raggedright\small\mbox{}0 0 0 0}\hspace{0.5em}\parbox[t]{\dimexpr0.1633\linewidth-0.5em\relax}{\raggedright\small\mbox{}0 0 0}\hspace{0.5em}\parbox[t]{\dimexpr0.1633\linewidth-0.5em\relax}{\raggedright\small\mbox{}0 0 0}\hspace{0.5em}\parbox[t]{\dimexpr0.1633\linewidth-0.5em\relax}{\raggedright\small\mbox{}0 0 0}\hspace{0.5em}}{}
\EdL{26}{AK07-J0114-L26}{\parbox[t]{\dimexpr0.3267\linewidth-0.5em\relax}{\raggedright\small\mbox{}für den 0. Muḥ. 288\,:}\hspace{0.5em}\parbox[t]{\dimexpr0.1633\linewidth-0.5em\relax}{\raggedright\small\mbox{}29\textsuperscript{d} 7\textsuperscript{h} 48\textsuperscript{m} 41\textsuperscript{s}}\hspace{0.5em}\parbox[t]{\dimexpr0.1633\linewidth-0.5em\relax}{\raggedright\small\mbox{}301° 20′ 17″}\hspace{0.5em}\parbox[t]{\dimexpr0.1633\linewidth-0.5em\relax}{\raggedright\small\mbox{}119° 9′ 58″\textsuperscript{4})}\hspace{0.5em}\parbox[t]{\dimexpr0.1633\linewidth-0.5em\relax}{\raggedright\small\mbox{}167° 8′ 0″}\hspace{0.5em}}{}
\EdLp{27}{AK07-J0114-L27}{Sonne und Mond haben also die gleiche mittlere Länge und zwar 301° 20′ 17″}{}
\EdL{28}{AK07-J0114-L28}{am 29. Muḥ. des J. 288 d. H. 7\textsuperscript{h} 49\textsuperscript{m} Abends\textsuperscript{5}). Man muss nun die wahren Örter von}{}
\EdL{29}{AK07-J0114-L29}{Sonne und Mond bestimmen für diese Zeit nach Kap. 8 u. 9 (Tafeln 21—26)\,:}{}
\EdL{30}{AK07-J0114-L30}{\parbox[t]{\dimexpr0.2450\linewidth-0.5em\relax}{\raggedright\small\mbox{}}\hspace{0.5em}\parbox[t]{\dimexpr0.4900\linewidth-0.5em\relax}{\raggedright\small\mbox{}Mittlere Länge d. Sonne\,:}\hspace{0.5em}\parbox[t]{\dimexpr0.2450\linewidth-0.5em\relax}{\raggedright\small\mbox{}Mittlere Länge d. Mondes\,:}\hspace{0.5em}}{}
\EdL{31}{AK07-J0114-L31}{\parbox[t]{\dimexpr0.2450\linewidth-0.5em\relax}{\raggedright\small\mbox{}}\hspace{0.5em}\parbox[t]{\dimexpr0.4900\linewidth-0.5em\relax}{\raggedright\small\mbox{}301° 20′ 17″}\hspace{0.5em}\parbox[t]{\dimexpr0.2450\linewidth-0.5em\relax}{\raggedright\small\mbox{}301° 20′ 17″}\hspace{0.5em}}{}
\EdL{32}{AK07-J0114-L32}{\parbox[t]{\dimexpr0.4900\linewidth-0.5em\relax}{\raggedright\small\mbox{}Gleichung d. Sonne}\hspace{0.5em}\parbox[t]{\dimexpr0.4900\linewidth-0.5em\relax}{\raggedright\small\mbox{}Gleichung d. Mondes}\hspace{0.5em}}{}
\EdL{33}{AK07-J0114-L33}{\parbox[t]{\dimexpr0.2450\linewidth-0.5em\relax}{\raggedright\small\mbox{}n. Tafel 25\textsuperscript{6})}\hspace{0.5em}\parbox[t]{\dimexpr0.2450\linewidth-0.5em\relax}{\raggedright\small\mbox{}1° 30′ 39″}\hspace{0.5em}\parbox[t]{\dimexpr0.2450\linewidth-0.5em\relax}{\raggedright\small\mbox{}n. Tafel 24\textsuperscript{7})}\hspace{0.5em}\parbox[t]{\dimexpr0.2450\linewidth-0.5em\relax}{\raggedright\small\mbox{}4° 15′ 45″}\hspace{0.5em}}{}
\EdL{34}{AK07-J0114-L34}{\parbox[t]{\dimexpr0.2450\linewidth-0.5em\relax}{\raggedright\small\mbox{}}\hspace{0.5em}\parbox[t]{\dimexpr0.4900\linewidth-0.5em\relax}{\raggedright\small\mbox{}302° 50′ 56″}\hspace{0.5em}\parbox[t]{\dimexpr0.2450\linewidth-0.5em\relax}{\raggedright\small\mbox{}297° 4′ 32″}\hspace{0.5em}}{}
\NoteRule
\EdN{N1}{AK07-J0114-N01}{\textsuperscript{1}) Das 31 bedeutet den 0. \textit{Muḥarram} des Jahres 31 d. H.}
\EdN{N2}{AK07-J0114-N02}{\textsuperscript{2}) Ed. \Sperr{Nallino}, I, 56—57.}
\EdN{N3}{AK07-J0114-N03}{\textsuperscript{3}) Ibid. I, 228—229.}
\EdN{N4}{AK07-J0114-N04}{\textsuperscript{4}) Eigentlich 479° 9′ 58″, davon werden aber 360° = eine ganze Umdrehung abgezogen.}
\EdN{N5}{AK07-J0114-N05}{\textsuperscript{5}) Die Tage der Tafeln beginnen am Mittag bei Kulmination der Sonne. Die Sekunden habe}
\EdN{N6}{AK07-J0114-N06}{ich zu 1 Min. aufgerundet.}
\EdN{N7}{AK07-J0114-N07}{\textsuperscript{6}) Diese Gleichung entspricht dem Argument 223° 25′ 17″, das man erhält, wenn man von der}
\EdN{N8}{AK07-J0114-N08}{mittlern Länge der Sonne die Länge des Apogeums (77° 55′) abzieht.}
\EdN{N9}{AK07-J0114-N09}{\textsuperscript{7}) Diese Gleichung entspricht dem Argument 119° 9′ 58″ (119° 10′), nämlich der Anomalie des}
\EdN{N10}{AK07-J0114-N10}{Mondes.}
\FrameHead{0115}{p. 83}{0115}{p. 83}
\EdLp{1}{AK07-J0115-L01}{Der Mond is also um 302° 50′ 56″ — 297° 4′ 32″ = 5° 46′ 24″ hinter der Sonne}{}
\EdL{2}{AK07-J0115-L02}{zurück, d. h. die Sonne \textit{praecessit,} wie es im Texte heisst, die wahre Konjunktion}{}
\EdL{3}{AK07-J0115-L03}{findet also erst nach der oben gefundenen Zeit statt, und zwar, wenn die Sonne als}{}
\EdL{4}{AK07-J0115-L04}{stillstehend betrachtet würde, um so viele Stunden, als die wahre stündliche Bewe-}{}
\EdL{5}{AK07-J0115-L05}{gung des Mondes in der Längendifferenz 5° 46′ 24″ enthalten ist\,; da aber die Länge}{}
\EdL{6}{AK07-J0115-L06}{der Sonne sich auch ändert, kommt eine Korrektur zur letztern Differenz hinzu, und}{}
\EdL{7}{AK07-J0115-L07}{zwar der 12. Teil derselben = 28′ 52″, die Längendifferenz ist nun also 6° 15′ 16″\,;}{}
\EdL{8}{AK07-J0115-L08}{nun ist\,: 6° 15′ 16″ ∶ 34′ 20″\textsuperscript{1}) = 10\textsuperscript{h} 55\textsuperscript{m} 48\textsuperscript{s}}{}
\EdL{9}{AK07-J0115-L09}{die \textit{wahre Konjunktion} hat also stattgefunden am 29. Muḥ. 288 d. H., 18\textsuperscript{h} 44\textsuperscript{m} 29\textsuperscript{s} nach}{}
\EdL{10}{AK07-J0115-L10}{Kulmination der Sonne, d. i. am 23. Januar 901, 6\textsuperscript{h} 44\textsuperscript{m} Morgens\textsuperscript{2}).}{}
\EdLp{11}{AK07-J0115-L11}{Um nun die Örter von Sonne und Mond bei der wahren Konjunktion zu erhalten,}{}
\EdL{12}{AK07-J0115-L12}{muss man nach den gegebenen Vorschriften 6° 15′ 16″ zum wahren Ort des Mondes}{}
\EdL{13}{AK07-J0115-L13}{addieren, zu dem der Sonne aber nur den 12. Teil der Längendifferenz (oder den}{}
\EdL{14}{AK07-J0115-L14}{13. Teil von 6° 15′ 16″), also 28′ 52″, so muss man natürlich dieselben Grössen erhalten\,:}{}
\EdL{15}{AK07-J0115-L15}{\parbox[t]{\dimexpr0.2450\linewidth-0.5em\relax}{\raggedright\small\mbox{}Ort der Sonne\,:}\hspace{0.5em}\parbox[t]{\dimexpr0.2450\linewidth-0.5em\relax}{\raggedright\small\mbox{}302° 50′ 56″}\hspace{0.5em}\parbox[t]{\dimexpr0.2450\linewidth-0.5em\relax}{\raggedright\small\mbox{}Ort des Mondes\,:}\hspace{0.5em}\parbox[t]{\dimexpr0.2450\linewidth-0.5em\relax}{\raggedright\small\mbox{}297° 4′ 32″}\hspace{0.5em}}{}
\EdL{16}{AK07-J0115-L16}{\parbox[t]{\dimexpr0.2450\linewidth-0.5em\relax}{\raggedright\small\mbox{}}\hspace{0.5em}\parbox[t]{\dimexpr0.4900\linewidth-0.5em\relax}{\raggedright\small\mbox{}+ 28′ 52″}\hspace{0.5em}\parbox[t]{\dimexpr0.2450\linewidth-0.5em\relax}{\raggedright\small\mbox{}+ 6° 15′ 16″}\hspace{0.5em}}{}
\EdL{17}{AK07-J0115-L17}{\parbox[t]{\dimexpr0.2450\linewidth-0.5em\relax}{\raggedright\small\mbox{}}\hspace{0.5em}\parbox[t]{\dimexpr0.4900\linewidth-0.5em\relax}{\raggedright\small\mbox{}303° 19′ 48″}\hspace{0.5em}\parbox[t]{\dimexpr0.2450\linewidth-0.5em\relax}{\raggedright\small\mbox{}303° 19′ 48″}\hspace{0.5em}}{}
\EdLp{18}{AK07-J0115-L18}{Dies ist also die \textit{wahre Länge} von Sonne und Mond zur Zeit ihrer wahren Kon-}{}
\EdL{19}{AK07-J0115-L19}{junktion am 29. Muḥ. 288, Morgens 6\textsuperscript{h} 44\textsuperscript{m}.}{}
\EdLp{20}{AK07-J0115-L20}{Der Grund, warum zu der Länge der Sonne \CaseFrac{1}{12} und zu der Länge des Mondes}{}
\EdL{21}{AK07-J0115-L21}{\CaseFrac{13}{12} der Längendifferenz addiert werden müssen (im Falle, wo der Mond der Sonne}{}
\EdL{22}{AK07-J0115-L22}{voraus ist, müssen sie subtrahiert werden), ist jedenfalls der, dass sich die mittlern}{}
\EdL{23}{AK07-J0115-L23}{stündlichen Bewegungen von Sonne und Mond nahezu wie 1 ∶ 13 verhalten\textsuperscript{3}). Es sei}{}
\EdL{24}{AK07-J0115-L24}{\textit{K} der Ort der wahren Konjunktion, \textit{S} der Ort der Sonne, \textit{M} derjenige des Mondes,}{}
\EdL{25}{AK07-J0115-L25}{nach der Korrektur durch ihre Gleichungen, so kommen auf \textit{MK} 13 Teile, auf \textit{SK}}{}
\EdL{26}{AK07-J0115-L26}{ein Teil, also auf \textit{MS} 12 Teile, man teilt also die Längendifferenz \textit{MS} durch 12, so}{}
\EdL{27}{AK07-J0115-L27}{hat man einen Teil, addiert diesen zur Länge der Sonne, und sein dreizehnfaches}{}
\EdL{28}{AK07-J0115-L28}{zur Länge des Mondes, so hat man die gleiche wahre Länge beider Gestirne zur}{}
\EdL{29}{AK07-J0115-L29}{Zeit der Konjunktion, und \textit{MK} geteilt durch die wahre stündliche Bewegung des}{}
\EdL{30}{AK07-J0115-L30}{Mondes gibt die Zahl der Stunden an, die noch die beiden Gestirne von den Stel-}{}
\NoteRule
\EdN{N1}{AK07-J0115-N01}{\textsuperscript{1}) Dies ist nach Tafel 64 die wahre stündliche Bewegung des Mondes für das Argument}
\EdN{N2}{AK07-J0115-N02}{(Anomalie) 119° 10′.}
\EdN{N3}{AK07-J0115-N03}{\textsuperscript{2}) Die Tafeln \Sperr{Schrams} (\textit{Kalendariograph. u. chronolog. Tafeln}, Leipzig, 1908) ergeben den}
\EdN{N4}{AK07-J0115-N04}{23. Jan. 6\textsuperscript{h} 43\textsuperscript{m} Morgens für \textit{Greenwich}, also für \textit{Cordova}, das c. 5° westlicher liegt, 6\textsuperscript{h} 23\textsuperscript{m}, der Fehler}
\EdN{N5}{AK07-J0115-N05}{wäre also c. 21\textsuperscript{m}\,; \Sperr{Schram} bemerkt aber, dass seine Tafeln die Syzygien mit einer Genauigkeit}
\EdN{N6}{AK07-J0115-N06}{von etwa einer halben Stunde ergeben, also könnte sich der Fehler bis auf Null reduzieren. Für}
\EdN{N7}{AK07-J0115-N07}{\textit{Antiochia} flndet \Sperr{Schiaparelli} die Zeit der wahren Konjunktion = 9\textsuperscript{h} 20\textsuperscript{m} Morgens\,; da \textit{Antiochia}}
\EdN{N8}{AK07-J0115-N08}{c. 41° östlich von \textit{Cordova} liegt, so würde die Zeit für Antiochia nach unsern Tafeln 9\textsuperscript{h} 29\textsuperscript{m} aus-}
\EdN{N9}{AK07-J0115-N09}{machen, also nur 9\textsuperscript{m} Unterschied gegenüber den Tafeln \Sperr{Batt.}’s.}
\EdN{N10}{AK07-J0115-N10}{\textsuperscript{3}) Mittlere stündliche Bewegung d. Mondes 32′ 56″, der Sonne 2′ 28″, die Division ergibt 13,3.}
\FrameHead{0116}{p. 84}{0116}{p. 84}
\EdL{1}{AK07-J0116-L01}{lungen \textit{M} und \textit{S} an bis zur Konjunktion \textit{K} zu laufen haben. Da das Verhältnis der}{}
\EdL{2}{AK07-J0116-L02}{stündlichen Bewegungen der beiden Gestirne nicht konstant ist, so gibt dies Verfahren}{}
\EdL{3}{AK07-J0116-L03}{nicht ganz genaue Resultate, und \Sperr{Batt.} bemerkt daher auch, sein weitläufigeres}{}
\EdL{4}{AK07-J0116-L04}{Verfahren sei das richtigere.}{}
\EdLp{5}{AK07-J0116-L05}{Es werden nun noch zwei andere Vorschriften gemacht, die für die Bestimmung}{}
\EdL{6}{AK07-J0116-L06}{der Konjunktion ohne Belang sind, aber für die später kommende Berechnung einer}{}
\EdL{7}{AK07-J0116-L07}{allfällig eintretenden Finsternis gebraucht werden. \textit{Erstens} wird (p. 24) bemerkt, die}{}
\EdL{8}{AK07-J0116-L08}{Gleichung des Mondes sei auch vom Argument der Breite zu subtrahieren (bezw.}{}
\EdL{9}{AK07-J0116-L09}{dazu zu addieren), dies gibt in unserm Falle 167° 8′ — 4° 15′ 45″ = 162° 52′ 15″,}{}
\EdL{10}{AK07-J0116-L10}{ferner sei zu dieser Grösse die korrigierte Längendifferenz (6° 15′ 16″) zu addieren, dies}{}
\EdL{11}{AK07-J0116-L11}{gibt\,: 169° 7′ 31″ als korrigiertes Argument der Breite. \textit{Zweitens} wird (p. 24—25) bei-}{}
\EdL{12}{AK07-J0116-L12}{gefügt, dass man auch die Anomalie des Mondes noch für die gefundene Zeitkorrektur}{}
\EdL{13}{AK07-J0116-L13}{von 10\textsuperscript{h} 55\textsuperscript{m} 48\textsuperscript{s} zu korrigieren habe\,; dies geschieht dadurch, dass man diese Grösse}{}
\EdL{14}{AK07-J0116-L14}{mit 32′ 40″, d. h. der mittlern stündlichen Bewegung der Anomalie des Mondes multi-}{}
\EdL{15}{AK07-J0116-L15}{pliziert, man erhält dadurch, wenn man statt 10\textsuperscript{h} 55\textsuperscript{m} 48\textsuperscript{s} den Dezimalbruch 10,93}{}
\EdL{16}{AK07-J0116-L16}{nimmt, 5° 57′ 3″, dies zur Anomalie (119° 10′) addiert\textsuperscript{1}) gibt 125° 7′ = die korrigierte}{}
\EdL{17}{AK07-J0116-L17}{Anomalie (im Texte\,: \textit{argumentum ratum}).}{}
\EdLp{18}{AK07-J0116-L18}{Am Schlusse wird noch angegeben, wie man die Zeit der Konjunktion zu}{}
\EdL{19}{AK07-J0116-L19}{korrigieren habe, wenn der Beobachtungsort westlich oder östlich von Cordova liege.\textsuperscript{2})}{}
\EdL{20}{AK07-J0116-L20}{Dann wird noch für den Fall, dass man Aequatorialstunden in ungleiche Stunden,}{}
\EdL{21}{AK07-J0116-L21}{oder umgekehrt verwandeln wolle, darauf hingewiesen, dass man hiefür die Tafeln}{}
\EdL{22}{AK07-J0116-L22}{der Zeitgleichung (67—68) benutzen müsse.}{}
\EdLp{23}{AK07-J0116-L23}{Wir geben auch noch ein \textit{Beispiel} einer \textit{Oppositionsberechnung.}}{}
\EdLp{24}{AK07-J0116-L24}{Zu welcher Zeit hat um die Mitte des Monats \textit{Rabīʿ II} (zählt 29 Tage) d. J. 146}{}
\EdL{25}{AK07-J0116-L25}{d. H. die \textit{Opposition} von Sonne und Mond (der \textit{Vollmond}) stattgefunden\,?\textsuperscript{3})}{}
\EdLp{26}{AK07-J0116-L26}{Nach den Tafeln 70—72 findet man\,:}{}
\EdL{27}{AK07-J0116-L27}{\parbox[t]{\dimexpr0.4900\linewidth-0.5em\relax}{\raggedright\small\mbox{}}\hspace{0.5em}\parbox[t]{\dimexpr0.1633\linewidth-0.5em\relax}{\raggedright\small\mbox{}Longit. media}\hspace{0.5em}\parbox[t]{\dimexpr0.1633\linewidth-0.5em\relax}{\raggedright\small\mbox{}Anomalia}\hspace{0.5em}\parbox[t]{\dimexpr0.1633\linewidth-0.5em\relax}{\raggedright\small\mbox{}Argumentum}\hspace{0.5em}}{}
\EdL{28}{AK07-J0116-L28}{\parbox[t]{\dimexpr0.3267\linewidth-0.5em\relax}{\raggedright\small\mbox{}}\hspace{0.5em}\parbox[t]{\dimexpr0.1633\linewidth-0.5em\relax}{\raggedright\small\mbox{}Dies et horae}\hspace{0.5em}\parbox[t]{\dimexpr0.1633\linewidth-0.5em\relax}{\raggedright\small\mbox{}solis et lunae}\hspace{0.5em}\parbox[t]{\dimexpr0.1633\linewidth-0.5em\relax}{\raggedright\small\mbox{}lunae}\hspace{0.5em}\parbox[t]{\dimexpr0.1633\linewidth-0.5em\relax}{\raggedright\small\mbox{}latitud. lunae}\hspace{0.5em}}{}
\EdL{29}{AK07-J0116-L29}{\parbox[t]{\dimexpr0.3267\linewidth-0.5em\relax}{\raggedright\small\mbox{}Am Ende des 120. Jahr}\hspace{0.5em}\parbox[t]{\dimexpr0.1633\linewidth-0.5em\relax}{\raggedright\small\mbox{}14\textsuperscript{d} 6\textsuperscript{h} 34\textsuperscript{m} 9\textsuperscript{s}}\hspace{0.5em}\parbox[t]{\dimexpr0.1633\linewidth-0.5em\relax}{\raggedright\small\mbox{}279° 28′ 28″}\hspace{0.5em}\parbox[t]{\dimexpr0.1633\linewidth-0.5em\relax}{\raggedright\small\mbox{}29° 12′ 6″}\hspace{0.5em}\parbox[t]{\dimexpr0.1633\linewidth-0.5em\relax}{\raggedright\small\mbox{}67° 12′ 57″}\hspace{0.5em}}{}
\EdL{30}{AK07-J0116-L30}{\parbox[t]{\dimexpr0.3267\linewidth-0.5em\relax}{\raggedright\small\mbox{}für 25 Jahre}\hspace{0.5em}\parbox[t]{\dimexpr0.1633\linewidth-0.5em\relax}{\raggedright\small\mbox{}4 11 21}\hspace{0.5em}\parbox[t]{\dimexpr0.1633\linewidth-0.5em\relax}{\raggedright\small\mbox{}91 37 55}\hspace{0.5em}\parbox[t]{\dimexpr0.1633\linewidth-0.5em\relax}{\raggedright\small\mbox{}185 3 54}\hspace{0.5em}\parbox[t]{\dimexpr0.1633\linewidth-0.5em\relax}{\raggedright\small\mbox{}201 10 50}\hspace{0.5em}}{}
\EdL{31}{AK07-J0116-L31}{\parbox[t]{\dimexpr0.3267\linewidth-0.5em\relax}{\raggedright\small\mbox{}für 4 Monate}\hspace{0.5em}\parbox[t]{\dimexpr0.1633\linewidth-0.5em\relax}{\raggedright\small\mbox{}\underline{1} 14 12 9\textsuperscript{4})}\hspace{0.5em}\parbox[t]{\dimexpr0.1633\linewidth-0.5em\relax}{\raggedright\small\mbox{}87 18 58}\hspace{0.5em}\parbox[t]{\dimexpr0.1633\linewidth-0.5em\relax}{\raggedright\small\mbox{}77 27 4}\hspace{0.5em}\parbox[t]{\dimexpr0.1633\linewidth-0.5em\relax}{\raggedright\small\mbox{}92 0 42}\hspace{0.5em}}{}
\EdL{32}{AK07-J0116-L32}{\parbox[t]{\dimexpr0.3267\linewidth-0.5em\relax}{\raggedright\small\mbox{}für den 0. Rabīʿ II, 146\,:}\hspace{0.5em}\parbox[t]{\dimexpr0.1633\linewidth-0.5em\relax}{\raggedright\small\mbox{}14\textsuperscript{d} 0\textsuperscript{h} 57\textsuperscript{m} 39\textsuperscript{s}}\hspace{0.5em}\parbox[t]{\dimexpr0.1633\linewidth-0.5em\relax}{\raggedright\small\mbox{}458° 25′ 21″}\hspace{0.5em}\parbox[t]{\dimexpr0.1633\linewidth-0.5em\relax}{\raggedright\small\mbox{}291° 43′ 4″}\hspace{0.5em}\parbox[t]{\dimexpr0.1633\linewidth-0.5em\relax}{\raggedright\small\mbox{}360° 24′ 29″}\hspace{0.5em}}{}
\EdL{33}{AK07-J0116-L33}{\parbox[t]{\dimexpr0.4900\linewidth-0.5em\relax}{\raggedright\small\mbox{}}\hspace{0.5em}\parbox[t]{\dimexpr0.3267\linewidth-0.5em\relax}{\raggedright\small\mbox{}= 98° 25′ 21″}\hspace{0.5em}\parbox[t]{\dimexpr0.1633\linewidth-0.5em\relax}{\raggedright\small\mbox{}= 0° 24′ 29″}\hspace{0.5em}}{}
\EdLp{34}{AK07-J0116-L34}{Man bestimmt nun die wahren Längen von Sonne und Mond nach den Tafeln}{}
\EdL{35}{AK07-J0116-L35}{21—26\,:}{}
\NoteRule
\EdN{N1}{AK07-J0116-N01}{\textsuperscript{1}) Diese Korrektur ist zu subtrahieren, wenn der Mond der Sonne voraus ist.}
\EdN{N2}{AK07-J0116-N02}{\textsuperscript{2}) Der Text ist hier fehlerhaft, es muss \textit{oriente} mit \textit{occidente} vertauscht werden.}
\EdN{N3}{AK07-J0116-N03}{\textsuperscript{3}) Der 1. Rabīʿ II d. J. 146 entspricht dem 17. Juni 763.}
\EdN{N4}{AK07-J0116-N04}{\textsuperscript{4}) \underline{1} bedeutet, dass dieser Tag negativ zu nehmen, also von der Summe der übrigen Tage und}
\EdN{N5}{AK07-J0116-N05}{Stunden zu subtrahieren ist.}
\FrameHead{0117}{p. 85}{0117}{p. 85}
\EdL{1}{AK07-J0117-L01}{\parbox[t]{\dimexpr0.1960\linewidth-0.5em\relax}{\raggedright\small\mbox{}}\hspace{0.5em}\parbox[t]{\dimexpr0.3920\linewidth-0.5em\relax}{\raggedright\small\mbox{}Mittlere Länge der Sonne\,:}\hspace{0.5em}\parbox[t]{\dimexpr0.3920\linewidth-0.5em\relax}{\raggedright\small\mbox{}Mittlere Länge des Mondes\,:\textsuperscript{1})}\hspace{0.5em}}{}
\EdL{2}{AK07-J0117-L02}{\parbox[t]{\dimexpr0.2450\linewidth-0.5em\relax}{\raggedright\small\mbox{}}\hspace{0.5em}\parbox[t]{\dimexpr0.4900\linewidth-0.5em\relax}{\raggedright\small\mbox{}98° 25′ 21″}\hspace{0.5em}\parbox[t]{\dimexpr0.2450\linewidth-0.5em\relax}{\raggedright\small\mbox{}98° 25′ 21″}\hspace{0.5em}}{}
\EdL{3}{AK07-J0117-L03}{\parbox[t]{\dimexpr0.4900\linewidth-0.5em\relax}{\raggedright\small\mbox{}Gleichung der Sonne}\hspace{0.5em}\parbox[t]{\dimexpr0.4900\linewidth-0.5em\relax}{\raggedright\small\mbox{}Gleichung des Mondes}\hspace{0.5em}}{}
\EdL{4}{AK07-J0117-L04}{\parbox[t]{\dimexpr0.2450\linewidth-0.5em\relax}{\raggedright\small\mbox{}n. Tafel 21\,:}\hspace{0.5em}\parbox[t]{\dimexpr0.2450\linewidth-0.5em\relax}{\raggedright\small\mbox{}0° 45′ 44″}\hspace{0.5em}\parbox[t]{\dimexpr0.2450\linewidth-0.5em\relax}{\raggedright\small\mbox{}n. Tafel 23\,:}\hspace{0.5em}\parbox[t]{\dimexpr0.2450\linewidth-0.5em\relax}{\raggedright\small\mbox{}4° 36′ 36″}\hspace{0.5em}}{}
\EdL{5}{AK07-J0117-L05}{\parbox[t]{\dimexpr0.2450\linewidth-0.5em\relax}{\raggedright\small\mbox{}}\hspace{0.5em}\parbox[t]{\dimexpr0.4900\linewidth-0.5em\relax}{\raggedright\small\mbox{}97° 39′ 37″}\hspace{0.5em}\parbox[t]{\dimexpr0.2450\linewidth-0.5em\relax}{\raggedright\small\mbox{}103° 1′ 57″}\hspace{0.5em}}{}
\EdLp{6}{AK07-J0117-L06}{Der Mond ist also um 5° 22′ 20″ der Sonne voraus, die Opposition hat also}{}
\EdL{7}{AK07-J0117-L07}{schon stattgefunden. Man dividiert die Differenz der Längen durch 12, dies gibt}{}
\EdL{8}{AK07-J0117-L08}{26′ 52″, addiert dieses zu 5° 22′ 20″, gibt 5° 49′ 12″\,; diese Grösse durch die stündliche}{}
\EdL{9}{AK07-J0117-L09}{wahre Bewegung des Mondes, die nach Tafel 63 = 31′ 50″ ist, geteilt, gibt 10,97\textsuperscript{h} =}{}
\EdL{10}{AK07-J0117-L10}{10\textsuperscript{h} 58\textsuperscript{m}, diese Zeitdifferenz von der Zeit, zu der die mittlere Opposition stattgefunden}{}
\EdL{11}{AK07-J0117-L11}{hat, abgezogen, gibt als Zeit der wahren Opposition rund 13\textsuperscript{d}14\textsuperscript{h}, d. h. diese \textit{Opposi-}}{}
\EdL{12}{AK07-J0117-L12}{\textit{tion} hat stattgefunden am 13. Rabīʿ II, 146, 14\textsuperscript{h} nach 12 Uhr Mittags, d. h. nach unserer}{}
\EdL{13}{AK07-J0117-L13}{(bürgerlichen) Zeitrechnung (Beginn des Tages um Mitternacht) am 30. Juni 763,}{}
\EdL{14}{AK07-J0117-L14}{2 Uhr Morgens\textsuperscript{2}).}{}
\EdLp{15}{AK07-J0117-L15}{Um die wahren Längen von Sonne und Mond bei der Opposition zu erhalten,}{}
\EdL{16}{AK07-J0117-L16}{muss man von dem oben gefundenen Wert der Länge der Sonne den 12. Teil der}{}
\EdL{17}{AK07-J0117-L17}{Differenz der Längen, also 26′ 52″ subtrahieren, von der Länge des Mondes aber}{}
\EdL{18}{AK07-J0117-L18}{das 13 fache, d. h. 5° 49′ 12″, so erhält man\,:}{}
\EdL{19}{AK07-J0117-L19}{\parbox[t]{\dimexpr0.2450\linewidth-0.5em\relax}{\raggedright\small\mbox{}Länge der Sonne\,:}\hspace{0.5em}\parbox[t]{\dimexpr0.2450\linewidth-0.5em\relax}{\raggedright\small\mbox{}97° 39′ 37″}\hspace{0.5em}\parbox[t]{\dimexpr0.2450\linewidth-0.5em\relax}{\raggedright\small\mbox{}Länge des Mondes\,:}\hspace{0.5em}\parbox[t]{\dimexpr0.2450\linewidth-0.5em\relax}{\raggedright\small\mbox{}103° 1′ 57″}\hspace{0.5em}}{}
\EdL{20}{AK07-J0117-L20}{\parbox[t]{\dimexpr0.2450\linewidth-0.5em\relax}{\raggedright\small\mbox{}}\hspace{0.5em}\parbox[t]{\dimexpr0.4900\linewidth-0.5em\relax}{\raggedright\small\mbox{}— 26′ 52″}\hspace{0.5em}\parbox[t]{\dimexpr0.2450\linewidth-0.5em\relax}{\raggedright\small\mbox{}— 5° 49′ 12″}\hspace{0.5em}}{}
\EdL{21}{AK07-J0117-L21}{\parbox[t]{\dimexpr0.2450\linewidth-0.5em\relax}{\raggedright\small\mbox{}}\hspace{0.5em}\parbox[t]{\dimexpr0.4900\linewidth-0.5em\relax}{\raggedright\small\mbox{}97° 12′ 45″}\hspace{0.5em}\parbox[t]{\dimexpr0.2450\linewidth-0.5em\relax}{\raggedright\small\mbox{}97° 12′ 45″}\hspace{0.5em}}{}
\EdLp{22}{AK07-J0117-L22}{Diese Zahl gibt die wahre Länge der Sonne zur Zeit der Opposition an, die}{}
\EdL{23}{AK07-J0117-L23}{des Mondes ist also um 180° grösser, d. h. = 277° 12′ 45″.}{}
\EdLp{24}{AK07-J0117-L24}{Nach den Vorschriften des Textes müssen nun noch das Argument der Breite}{}
\EdL{25}{AK07-J0117-L25}{des Mondes und seine Anomalie korrigiert werden.}{}
\ColHead{left column}
\EdL{26}{AK07-J0117-L26}{Argument der Breite = 360° 24′ 29″}{}
\EdL{27}{AK07-J0117-L27}{— 5° 49′ 12″}{}
\EdL{28}{AK07-J0117-L28}{354° 35′ 17″}{}
\EdL{29}{AK07-J0117-L29}{Dazu addiert die Gleichung}{}
\EdL{30}{AK07-J0117-L30}{des Mondes\,: 4° 36′ 36″}{}
\EdL{31}{AK07-J0117-L31}{korrig. Argument der Breite\,: 359° 11′ 53 ″}{}
\ColHead{right column}
\EdL{32}{AK07-J0117-L32}{Anomalie des Mondes = 291° 43′ 4″}{}
\EdL{33}{AK07-J0117-L33}{Korrektur entsprechend der}{}
\EdL{34}{AK07-J0117-L34}{Zeitkorrektur von 10\textsuperscript{h} 58\textsuperscript{m}\textsuperscript{3}) 5° 58′ 21″}{}
\EdL{35}{AK07-J0117-L35}{korrigierte Anomalie\textsuperscript{4})\,: 285° 44′ 43″}{}
\NoteRule
\EdN{N1}{AK07-J0117-N01}{\textsuperscript{1}) Für die weitern Rechnungen behält man diese mit der mittlern Länge der Sonne überein-}
\EdN{N2}{AK07-J0117-N02}{stimmende Länge des Mondes bei, obgleich sie bei der Opposition in Wirklichkeit um 180° grösser}
\EdN{N3}{AK07-J0117-N03}{(bezw. kleiner) ist.}
\EdN{N4}{AK07-J0117-N04}{\textsuperscript{2}) Die Tafeln \Sperr{Schrams} ergeben\,: am 30. Juni 763, 1\textsuperscript{h} 41\textsuperscript{m} Morgens für Greenwich, also für}
\EdN{N5}{AK07-J0117-N05}{Cordova 1\textsuperscript{h} 21\textsuperscript{m}, der Unterschied würde also hier c. 40\textsuperscript{m} ausmachen, nach der oben angeführten}
\EdN{N6}{AK07-J0117-N06}{Bemerkung \Sperr{Schrams} könnte sich derselbe auf 10\textsuperscript{m} reduzieren.}
\EdN{N7}{AK07-J0117-N07}{\textsuperscript{3}) Wie oben muss man nach der Regel (p. 24) die Zeitkorrektur mit 32′ 40″ multiplizieren.}
\EdN{N8}{AK07-J0117-N08}{\textsuperscript{4}) Von \Sperr{Batt.} \textit{Anomalia aequata} genannt.}
\FrameHead{0118}{p. 86}{0118}{p. 86}
\EdLp{1}{AK07-J0118-L01}{Aus diesen Beispielen ersieht man, dass \Sperr{Kh.} (bezw. \Sperr{Al-Madjrīṭī}) in der}{}
\EdL{2}{AK07-J0118-L02}{Bestimmung der Syzygien im grossen und ganzen dem \Sperr{Ptol.} gefolgt ist, wie es auch}{}
\EdL{3}{AK07-J0118-L03}{\Sperr{Batt.} getan hat. Die Bestimmungen des \textit{Sūrya-Siddhānta} sind in dieser Hinsicht}{}
\EdL{4}{AK07-J0118-L04}{unvollständig und dunkel\textsuperscript{1}).}{}
\EdLc{5}{AK07-J0118-L05}{\textbf{Kap. 33. Wie die Verfinsterung des Mondes bestimmt wird.\textsuperscript{2})}}{}
\EdLc{6}{AK07-J0118-L06}{(Text p. 26—28, Tafeln 73—76.)}{}
\EdLp{7}{AK07-J0118-L07}{Da Kap. 32 an dieser Stelle als eingeschoben erscheint\textsuperscript{3}), indem es gar nicht zu}{}
\EdL{8}{AK07-J0118-L08}{dem Gegenstande passt, der in Kap. 31 u. 33 behandelt wird, nämlich die Berechnung}{}
\EdL{9}{AK07-J0118-L09}{der Syzygien und Finsternisse, so schliessen wir hier unmittelbar Kap. 33 an.}{}
\EdLp{10}{AK07-J0118-L10}{Die Regeln \Sperr{Batt.}’s zur Bestimmung der Zeit und Grösse der Finsternisse sind}{}
\EdL{11}{AK07-J0118-L11}{viel komplizierter als bei \Sperr{Kh.}, er beschäftigt sich damit in den Kap. 43 und 44\textsuperscript{4})\,;}{}
\EdL{12}{AK07-J0118-L12}{dafür sind die Tafeln \Sperr{Kh.}’s etwas weiter ausgeführt, \Sperr{Batt.} hat verschiedenes der}{}
\EdL{13}{AK07-J0118-L13}{Rechnung überlassen, was bei \Sperr{Kh.} aus den Tafeln abgelesen werden kann.}{}
\EdLp{14}{AK07-J0118-L14}{Zur Berechnung einer \textit{Mondfinsternis} werden zuerst bestimmt und notiert\,: die}{}
\EdL{15}{AK07-J0118-L15}{Zeit der Opposition, die wahren Örter von Sonne und Mond, und das korrigierte}{}
\EdL{16}{AK07-J0118-L16}{Argument der Breite. Für unser oben (p. 84—85) gegebenes Beispiel einer \textit{Opposition}}{}
\EdL{17}{AK07-J0118-L17}{sind diese Grössen folgende\,:}{}
\EdL{18}{AK07-J0118-L18}{\parbox[t]{\dimexpr0.3267\linewidth-0.5em\relax}{\raggedright\small\mbox{}}\hspace{0.5em}\parbox[t]{\dimexpr0.3267\linewidth-0.5em\relax}{\raggedright\small\mbox{}Wahrer Ort der Sonne}\hspace{0.5em}\parbox[t]{\dimexpr0.3267\linewidth-0.5em\relax}{\raggedright\small\mbox{}Korrig. Argument}\hspace{0.5em}}{}
\EdL{19}{AK07-J0118-L19}{\parbox[t]{\dimexpr0.3267\linewidth-0.5em\relax}{\raggedright\small\mbox{}Zeit der Opposition\,:}\hspace{0.5em}\parbox[t]{\dimexpr0.3267\linewidth-0.5em\relax}{\raggedright\small\mbox{}(bezw. des Mondes):}\hspace{0.5em}\parbox[t]{\dimexpr0.3267\linewidth-0.5em\relax}{\raggedright\small\mbox{}der Breite\,:}\hspace{0.5em}}{}
\EdL{20}{AK07-J0118-L20}{\parbox[t]{\dimexpr0.3267\linewidth-0.5em\relax}{\raggedright\small\mbox{}13. Rabīʿ II, 146, 14\textsuperscript{h}}\hspace{0.5em}\parbox[t]{\dimexpr0.3267\linewidth-0.5em\relax}{\raggedright\small\mbox{}97° 12′ 45″ (277° 12′ 45″)}\hspace{0.5em}\parbox[t]{\dimexpr0.3267\linewidth-0.5em\relax}{\raggedright\small\mbox{}359° 11′ 53″}\hspace{0.5em}}{}
\EdLp{21}{AK07-J0118-L21}{Wir müssen nun zuerst untersuchen, ob der Mond in Erdnähe oder Erdferne}{}
\EdL{22}{AK07-J0118-L22}{sei, denn davon hängt die Grösse des Erdschattens, und damit auch die Dauer und}{}
\EdL{23}{AK07-J0118-L23}{Grösse der Finsternis ab. Nach einer Bemerkung am Schlusse, die nur in den}{}
\EdL{24}{AK07-J0118-L24}{Mss. C u. N\textsuperscript{5}) steht, ist der Mond in \textit{spatio longinquo} (Erdferne), wenn seine Länge}{}
\EdL{25}{AK07-J0118-L25}{89°—91° beträgt, in \textit{spatio propinquo} (Erdnähe), wenn sie 269°—271° beträgt. In unserm}{}
\EdL{26}{AK07-J0118-L26}{Falle beträgt sie c. 277° 13′, das wird schon nicht mehr als Erdnähe aufgefasst,}{}
\EdL{27}{AK07-J0118-L27}{wenigstens in Bezug auf unsere Berechnung\,; in diesem Falle müssen wir also nach}{}
\EdL{28}{AK07-J0118-L28}{der Vorschrift mit dem Argument der Breite (359° 12′)\textsuperscript{6}) sowohl in die Tafeln \textit{pro}}{}
\EdL{29}{AK07-J0118-L29}{\textit{spatio propinquo} als auch in die \textit{pro spatio longinquo} eintreten. Wir finden nun in}{}
\EdL{30}{AK07-J0118-L30}{Tafel 76 zu diesem Argument der Breite (mit Hilfe der Interpolation) folgende Daten\,:}{}
\NoteRule
\EdN{N1}{AK07-J0118-N01}{\textsuperscript{1}) Transl. \Sperr{Burgess}, p. 275.}
\EdN{N2}{AK07-J0118-N02}{\textsuperscript{2}) Der Titel des Textes spricht von Sonne und Mond, es werden aber in diesem Kapitel tat-}
\EdN{N3}{AK07-J0118-N03}{sächlich nur die Mondfinsternisse behandelt, die Sonnenfinsternisse folgen im Kap. 35, daher habe}
\EdN{N4}{AK07-J0118-N04}{ich den Titel in diesen Sinne abgeändert.}
\EdN{N5}{AK07-J0118-N05}{\textsuperscript{3}) Es befindet sich nur in Ms. C an dieser Stelle, in N erst auf fol. 56\textsuperscript{r}, d. h. an richtiger}
\EdN{N6}{AK07-J0118-N06}{Stelle, in O fehlt es ganz.}
\EdN{N7}{AK07-J0118-N07}{\textsuperscript{4}) Ed. \Sperr{Nallino}, I, 96—113, u. 275—282.}
\EdN{N8}{AK07-J0118-N08}{\textsuperscript{5}) In N ist der Wortlaut von C durch eine spätere Hand korrigiert worden, s. p. 28, Note 1;}
\EdN{N9}{AK07-J0118-N09}{ich folge in meiner Darstellung dem Ms. C, dessen Fassung ich als die richtige betrachte.}
\EdN{N10}{AK07-J0118-N10}{\textsuperscript{6}) Die Tafeln \Sperr{Batt.}’s (II, 90) haben in der 1. Kolonne die Breite selbst.}
\FrameHead{0119}{p. 87}{0119}{p. 87}
\EdL{1}{AK07-J0119-L01}{\parbox[t]{\dimexpr0.2450\linewidth-0.5em\relax}{\raggedright\small\mbox{}}\hspace{0.5em}\parbox[t]{\dimexpr0.2450\linewidth-0.5em\relax}{\raggedright\small\mbox{}Digiti der Verfinsterung\,:}\hspace{0.5em}\parbox[t]{\dimexpr0.2450\linewidth-0.5em\relax}{\raggedright\small\mbox{}Min. u. Sec. des Casus\,:}\hspace{0.5em}\parbox[t]{\dimexpr0.2450\linewidth-0.5em\relax}{\raggedright\small\mbox{}Min. u. Sec. der halben Mora.}\hspace{0.5em}}{}
\EdL{2}{AK07-J0119-L02}{\parbox[t]{\dimexpr0.2450\linewidth-0.5em\relax}{\raggedright\small\mbox{}}\hspace{0.5em}\parbox[t]{\dimexpr0.2450\linewidth-0.5em\relax}{\raggedright\small\mbox{}20\textsuperscript{dig.} 14′}\hspace{0.5em}\parbox[t]{\dimexpr0.2450\linewidth-0.5em\relax}{\raggedright\small\mbox{}34′ 45″}\hspace{0.5em}\parbox[t]{\dimexpr0.2450\linewidth-0.5em\relax}{\raggedright\small\mbox{}27′ 8″}\hspace{0.5em}}{}
\EdL{3}{AK07-J0119-L03}{in Tafel 74 folgende\,:}{}
\EdL{4}{AK07-J0119-L04}{\parbox[t]{\dimexpr0.2450\linewidth-0.5em\relax}{\raggedright\small\mbox{}}\hspace{0.5em}\parbox[t]{\dimexpr0.2450\linewidth-0.5em\relax}{\raggedright\small\mbox{}19\textsuperscript{dig.} 16′}\hspace{0.5em}\parbox[t]{\dimexpr0.2450\linewidth-0.5em\relax}{\raggedright\small\mbox{}29′ 39″}\hspace{0.5em}\parbox[t]{\dimexpr0.2450\linewidth-0.5em\relax}{\raggedright\small\mbox{}21′ 2″}\hspace{0.5em}}{}
\EdL{5}{AK07-J0119-L05}{\parbox[t]{\dimexpr0.2450\linewidth-0.5em\relax}{\raggedright\small\mbox{}Differenzen\,:}\hspace{0.5em}\parbox[t]{\dimexpr0.2450\linewidth-0.5em\relax}{\raggedright\small\mbox{}0\textsuperscript{dig.} 58′}\hspace{0.5em}\parbox[t]{\dimexpr0.2450\linewidth-0.5em\relax}{\raggedright\small\mbox{}5′ 6″}\hspace{0.5em}\parbox[t]{\dimexpr0.2450\linewidth-0.5em\relax}{\raggedright\small\mbox{}6′ 6″}\hspace{0.5em}}{}
\EdLp{6}{AK07-J0119-L06}{Dieses sind die Unterschiede der drei Daten in Erdnähe und Erdferne. Was}{}
\EdL{7}{AK07-J0119-L07}{bedeuten nun diese Daten\,? Aus unserm Texte erfahren wir darüber nichts\,; \Sperr{Batt.}}{}
\EdL{8}{AK07-J0119-L08}{belehrt uns aber hierüber im 43. Kap.\textsuperscript{1}) Er wie auch \Sperr{Kh.} verstehen unter \textit{Digiti} der}{}
\EdL{9}{AK07-J0119-L09}{Verfinsterung sowohl bei \textit{partialer} als auch \textit{totaler} Mondfinsternis die Distanz vom}{}
\EdL{10}{AK07-J0119-L10}{Rande des Erdschattens bis zu dem zu ihm konvexen Rande der Mondscheibe, in}{}
\EdL{11}{AK07-J0119-L11}{solchen Teilen, von denen 12 auf den Monddurchmesser kommen.}{}
\EdLp{12}{AK07-J0119-L12}{Ist \textit{S} der Mittelpunkt des Erdschattens, der in der Ekliptik liegen muss, \textit{M} der}{}
\EdL{13}{AK07-J0119-L13}{Mittelpunkt des Mondes (in Fig. 1 bei partialer, in Fig. 2 bei totaler Finsternis),}{}
\EdL{14}{AK07-J0119-L14}{so ist in Fig. 1 \textit{AM} in Digiti (Zollen) ausgedrückt die Grösse der Verfinsterung, in}{}
\EdL{15}{AK07-J0119-L15}{\parbox[t]{\dimexpr0.4900\linewidth-0.5em\relax}{\raggedright\small\mbox{}Fig. 1.}\hspace{0.5em}\parbox[t]{\dimexpr0.4900\linewidth-0.5em\relax}{\raggedright\small\mbox{}Fig. 2.}\hspace{0.5em}}{}
\EdL{16}{AK07-J0119-L16}{Fig. 2 \textit{AB}. Berechnet wird dieselbe in folgender Weise (nach der einfachen Vor-}{}
\EdL{17}{AK07-J0119-L17}{schrift \Sperr{Batt.}’s)\,: Man addiert die beiden Radien des Schattens und des Mondes (in}{}
\EdL{18}{AK07-J0119-L18}{beiden Fig. sei dies \textit{SC}), subtrahiert von dieser Summe die Breite des Mondes (\textit{SM}),}{}
\EdL{19}{AK07-J0119-L19}{so hat man die Grösse der Verfinsterung in Minuten (in Fig. 1 \textit{AM,} in Fig. 2 \textit{CM} = \textit{AB}).}{}
\EdL{20}{AK07-J0119-L20}{Diese multipliziert man mit 12, und teilt das Produkt durch den Monddurchmesser}{}
\EdL{21}{AK07-J0119-L21}{in Minuten, so hat man die Verfinsterung in \textit{Digiti.}}{}
\EdLp{22}{AK07-J0119-L22}{Was die Ausdrücke \textit{Casus} und \textit{Mora} anbetrifft, so ist der erstere die Übersetzung}{}
\EdL{23}{AK07-J0119-L23}{des arab. \textit{suḳūṭ,} das wir hier mit „Einfallen“ übersetzen müssen\,; mit diesem Worte}{}
\EdL{24}{AK07-J0119-L24}{geben die arab. Astronomen das ptolemäische \textit{ἔμπτωσις} = das „Hineinfallen“ wieder\,;}{}
\EdL{25}{AK07-J0119-L25}{\textit{Casus} bedeutet also die Phase von Beginn der Finsternis bis zum Beginn der Totalität.}{}
\EdL{26}{AK07-J0119-L26}{Der zweite (Mora) ist die Übersetzung des arab. \textit{mukth} (od. \textit{makth}) = „Verweilen“,}{}
\EdL{27}{AK07-J0119-L27}{dies ist die Übersetzung des ptolemäischen \textit{μονή}\,; \textit{Mora} bedeutet also die Phase}{}
\EdL{28}{AK07-J0119-L28}{der totalen Verfinsterung. — Im Text dieses Kapitels kommt noch \textit{respiratio} vor als}{}
\EdL{29}{AK07-J0119-L29}{Übersetzung des arab. \textit{indjilā’} = „Erscheinen“, das wiederum dem ptolemäischen}{}
\NoteRule
\EdN{N1}{AK07-J0119-N01}{\textsuperscript{1}) Ed. \Sperr{Nallino}, I, 97, Z. 25—98, Z. 5.}
\FrameHead{0120}{p. 88}{0120}{p. 88}
\EdL{1}{AK07-J0120-L01}{\textit{ἀναπλήρωσις} = das „Wiedervollmachen“ (des Mondes) entspricht\,; besser wäre \textit{emersio}}{}
\EdL{2}{AK07-J0120-L02}{statt \textit{respiratio.}}{}
\EdLp{3}{AK07-J0120-L03}{Für die Berechnung der Minuten des \textit{Casus} und der \textit{Mora} geben wir ebenfalls}{}
\EdL{4}{AK07-J0120-L04}{die Regeln \Sperr{Batt.}’s (I, 98, Z. 5—14) und zeigen ihre Richtigkeit an der Hand folgender}{}
\EdL{5}{AK07-J0120-L05}{zwei Figuren. \Sperr{Batt.} sagt zuerst\,: Vom Quadrat der Summe der beiden Radien (\textit{MS}\textsuperscript{2}}{}
\EdL{6}{AK07-J0120-L06}{in Fig. 3) subtrahiere das Quadrat der Breite des Mondes (\textit{SN}\textsuperscript{2}), so gibt die Wurzel}{}
\EdL{7}{AK07-J0120-L07}{aus der Differenz (also \textit{MN}), durch die stündliche Bewegung des Mondes geteilt, die}{}
\EdL{8}{AK07-J0120-L08}{Zeit vom Beginn der Finsternis bis zur Mitte der totalen Verfinsterung (also \textit{Casus}}{}
\EdL{9}{AK07-J0120-L09}{und \textit{halbe Mora} zusammen). — Zweitens sagt er\,: Wenn totale Verfinsterung stattfindet}{}
\EdL{10}{AK07-J0120-L10}{(und man will die Zeit dieser allein kennen), so subtrahiere man den Monddurch-}{}
\EdL{11}{AK07-J0120-L11}{messer (\textit{MT} in Fig. 4) von der Summe der beiden Radien (\textit{ST})\textsuperscript{1}), vom Quadrate dieser}{}
\EdL{12}{AK07-J0120-L12}{Differenz (\textit{SM}\textsuperscript{2}) subtrahiere man das Quadrat der Breite des Mondes (\textit{SN}\textsuperscript{2}), so gibt die}{}
\EdL{13}{AK07-J0120-L13}{Wurzel aus der Differenz (also \textit{MN}), durch die stündliche Bewegung des Mondes geteilt,}{}
\EdL{14}{AK07-J0120-L14}{die halbe Zeit der totalen Verfinsterung (also die \textit{halbe Mora}). Subtrahiert man diese}{}
\EdL{15}{AK07-J0120-L15}{von der nach der ersten Regel erhaltenen Zeit, so erhält man die Zeit des \textit{Casus.}}{}
\EdL{16}{AK07-J0120-L16}{\parbox[t]{\dimexpr0.4900\linewidth-0.5em\relax}{\raggedright\small\mbox{}Fig. 3.}\hspace{0.5em}\parbox[t]{\dimexpr0.4900\linewidth-0.5em\relax}{\raggedright\small\mbox{}Fig. 4.}\hspace{0.5em}}{}
\EdLp{17}{AK07-J0120-L17}{Kehren wir nun zu den oben erhaltenen Differenzen zwischen den Daten der}{}
\EdL{18}{AK07-J0120-L18}{Erdnähe und Erdferne zurück. Diese erfahren nun eine Korrektur, die von der}{}
\EdL{19}{AK07-J0120-L19}{Änderung des Monddurchmessers in seiner Bewegung vom Apogeum zum Perigeum}{}
\EdL{20}{AK07-J0120-L20}{herrührt, und die mit Hilfe der zweitletzten Kolonne von Tafel 74 (bezw. 73) bestimmt}{}
\EdL{21}{AK07-J0120-L21}{wird, überschrieben \textit{Tabula proportionis.} Man tritt mit der \textit{Anomalia aequata} (korrig.}{}
\EdL{22}{AK07-J0120-L22}{Anomalie), die in unserm Falle rund 285° 45′ ist (s. oben p. 85), in diese Kolonne}{}
\EdL{23}{AK07-J0120-L23}{ein, der zugehörige Wert, \textit{dakaicae superationum} genannt, ist 20′ 43″, mit dieser Zahl,}{}
\EdL{24}{AK07-J0120-L24}{die also = 20\CaseFrac{43}{60} ∶ 60, d. h. in einem Dezimalbruch = 0,345 ist, multipliziert man}{}
\EdL{25}{AK07-J0120-L25}{die obigen drei Differenzen, und erhält die Werte\,:}{}
\EdL{26}{AK07-J0120-L26}{\parbox[t]{\dimexpr0.3267\linewidth-0.5em\relax}{\raggedright\small\mbox{}0\textsuperscript{dig.} 20′}\hspace{0.5em}\parbox[t]{\dimexpr0.3267\linewidth-0.5em\relax}{\raggedright\small\mbox{}1′ 46″}\hspace{0.5em}\parbox[t]{\dimexpr0.3267\linewidth-0.5em\relax}{\raggedright\small\mbox{}2′ 6″}\hspace{0.5em}}{}
\EdL{27}{AK07-J0120-L27}{diese addiert man zu den entsprechenden Werten der zweiten Reihe (für die Erd-}{}
\EdL{28}{AK07-J0120-L28}{ferne) und erhält so die korrigierten Werte\,:}{}
\EdL{29}{AK07-J0120-L29}{\parbox[t]{\dimexpr0.3267\linewidth-0.5em\relax}{\raggedright\small\mbox{}Digiti der Verfinsterung\,:}\hspace{0.5em}\parbox[t]{\dimexpr0.3267\linewidth-0.5em\relax}{\raggedright\small\mbox{}Min. u. Sec. des Casus\,:}\hspace{0.5em}\parbox[t]{\dimexpr0.3267\linewidth-0.5em\relax}{\raggedright\small\mbox{}Min. u. Sec. der halben Mora\,:}\hspace{0.5em}}{}
\EdL{30}{AK07-J0120-L30}{\parbox[t]{\dimexpr0.3267\linewidth-0.5em\relax}{\raggedright\small\mbox{}19\textsuperscript{dig.} 36′}\hspace{0.5em}\parbox[t]{\dimexpr0.3267\linewidth-0.5em\relax}{\raggedright\small\mbox{}31′ 25″}\hspace{0.5em}\parbox[t]{\dimexpr0.3267\linewidth-0.5em\relax}{\raggedright\small\mbox{}23′ 8″}\hspace{0.5em}}{}
\NoteRule
\EdN{N1}{AK07-J0120-N01}{\textsuperscript{1}) Dafür hätte er auch einfacher sagen können\,: den Mondradius (\textit{MP}) vom Schattenradius}
\EdN{N2}{AK07-J0120-N02}{(\textit{SP}).}
\FrameHead{0121}{p. 89}{0121}{p. 89}
\EdLp{1}{AK07-J0121-L01}{Nun widersprechen sich zwei Stellen in unserm Texte\,: p. 26 heisst es\,: \textit{In}}{}
\EdL{2}{AK07-J0121-L02}{\textit{supposita autem tabula tam casus quam morae tempora per dakaicas et elthenias}}{}
\EdL{3}{AK07-J0121-L03}{\textit{horarum determinata sunt}\,; das Wort \textit{horarum} hat aber nur O, C und N nicht.}{}
\EdL{4}{AK07-J0121-L04}{C und N fügen nun einen Zusatz zu diesem Kapitel hinzu, betitelt\,: \textit{Scientia examina-}}{}
\EdL{5}{AK07-J0121-L05}{\textit{tionis lunarium defectuum} (s. p. 27). In diesem Zuzatz findet sich die Vorschrift,}{}
\EdL{6}{AK07-J0121-L06}{diese gefundenen Minuten und Secunden für den Casus und die Mora um einen}{}
\EdL{7}{AK07-J0121-L07}{Zwölftel zu vermehren, und die Summe durch die stündliche Bewegung des Mondes}{}
\EdL{8}{AK07-J0121-L08}{zu teilen, so erhalte man die Zeit der partialen und der totalen Verfinsterung.}{}
\EdL{9}{AK07-J0121-L09}{C fasst also die gefundenen Min. und Sec. als Bogen-Min. und Sec. auf, O dagegen}{}
\EdL{10}{AK07-J0121-L10}{als Zeit-Min. und Sec. Wir müssen nach Vergleichung mit den Vorschriften \Sperr{Batt.}’s,}{}
\EdL{11}{AK07-J0121-L11}{die mit denen von C ziemlich identisch sind, die Darstellung des letztern Ms. als}{}
\EdL{12}{AK07-J0121-L12}{die richtige annehmen. Aber es besteht auch noch ein Unterschied zwischen den}{}
\EdL{13}{AK07-J0121-L13}{Vorschriften des Ms. C und denen \Sperr{Batt.}’s, der letztere verlangt die oben erwähnte}{}
\EdL{14}{AK07-J0121-L14}{Addition eines Zwölftels zu den Werten des Casus und der Mora nicht. Wir ver-}{}
\EdL{15}{AK07-J0121-L15}{fahren nun nach den Vorschriften des Ms. C und finden\,:}{}
\EdL{16}{AK07-J0121-L16}{\parbox[t]{\dimexpr0.1960\linewidth-0.5em\relax}{\raggedright\small\mbox{}Digiti der Verfinsterung\,:}\hspace{0.5em}\parbox[t]{\dimexpr0.3920\linewidth-0.5em\relax}{\raggedright\small\mbox{}Min. u. Sec. des Casus\,:}\hspace{0.5em}\parbox[t]{\dimexpr0.3920\linewidth-0.5em\relax}{\raggedright\small\mbox{}Min. u. Sec. der halben Mora\,:}\hspace{0.5em}}{}
\EdL{17}{AK07-J0121-L17}{\parbox[t]{\dimexpr0.3920\linewidth-0.5em\relax}{\raggedright\small\mbox{}19\textsuperscript{dig.} 36′}\hspace{0.5em}\parbox[t]{\dimexpr0.3920\linewidth-0.5em\relax}{\raggedright\small\mbox{}31′ 25″}\hspace{0.5em}\parbox[t]{\dimexpr0.1960\linewidth-0.5em\relax}{\raggedright\small\mbox{}23′ 8″}\hspace{0.5em}}{}
\EdL{18}{AK07-J0121-L18}{\parbox[t]{\dimexpr0.1960\linewidth-0.5em\relax}{\raggedright\small\mbox{}}\hspace{0.5em}\parbox[t]{\dimexpr0.1960\linewidth-0.5em\relax}{\raggedright\small\mbox{}\CaseFrac{1}{12} davon\,:}\hspace{0.5em}\parbox[t]{\dimexpr0.1960\linewidth-0.5em\relax}{\raggedright\small\mbox{}2′ 37″}\hspace{0.5em}\parbox[t]{\dimexpr0.1960\linewidth-0.5em\relax}{\raggedright\small\mbox{}\CaseFrac{1}{12} davon\,:}\hspace{0.5em}\parbox[t]{\dimexpr0.1960\linewidth-0.5em\relax}{\raggedright\small\mbox{}1′ 55″}\hspace{0.5em}}{}
\EdL{19}{AK07-J0121-L19}{\parbox[t]{\dimexpr0.3920\linewidth-0.5em\relax}{\raggedright\small\mbox{}}\hspace{0.5em}\parbox[t]{\dimexpr0.3920\linewidth-0.5em\relax}{\raggedright\small\mbox{}34′ 2″\textsuperscript{1})}\hspace{0.5em}\parbox[t]{\dimexpr0.1960\linewidth-0.5em\relax}{\raggedright\small\mbox{}25′ 3″\textsuperscript{1})}\hspace{0.5em}}{}
\EdLp{20}{AK07-J0121-L20}{Teilen wir nun die Min. und Sec. des Casus und der Mora durch 32′ 9″, d. i. die}{}
\EdL{21}{AK07-J0121-L21}{stündliche Bewegung des Mondes bei einer Anomalie von 285° 45′, so erhalten wir}{}
\EdL{22}{AK07-J0121-L22}{für die Zeit des \textit{Casus} 1\textsuperscript{h} 3\textsuperscript{m} 30\textsuperscript{s}, für die Zeit der \textit{halben Mora} 46\textsuperscript{m} 45\textsuperscript{s}. Wir haben}{}
\EdL{23}{AK07-J0121-L23}{nun folgende Daten für diese Finsternis\,:}{}
\EdLp{24}{AK07-J0121-L24}{1) \textit{Totale,} aber nicht centrale Verfinsterung, da die Zahl der Digiti grösser ist}{}
\EdL{25}{AK07-J0121-L25}{als 12 — 2) Grösse der Verfinsterung 19\textsuperscript{dig.} 36′ — 3) Dauer der partialen Verfinsterung}{}
\EdL{26}{AK07-J0121-L26}{(\textit{Casus} + \textit{Emersio}) = 2 (1\textsuperscript{h} 3\textsuperscript{m} 30\textsuperscript{s}) = 2\textsuperscript{h} 7\textsuperscript{m} — 4) Dauer der totalen Verfinsterung (\textit{Mora})}{}
\EdL{27}{AK07-J0121-L27}{= 2 (46\textsuperscript{m} 45\textsuperscript{s}) = 1\textsuperscript{h} 33\textsuperscript{m} 30\textsuperscript{s} — 5) Dauer der \textit{ganzen} Finsternis 3\textsuperscript{h} 40\textsuperscript{m} 30\textsuperscript{s}.}{}
\EdLp{28}{AK07-J0121-L28}{Die Opposition hat, wie wir oben (p. 85) gefunden haben, am 13. Rabīʿ II, 146,14\textsuperscript{h}}{}
\EdL{29}{AK07-J0121-L29}{nach 12 Uhr Mittags = 30. Juni 763, 2 Uhr Morgens stattgefunden, dies ist die Zeit}{}
\EdL{30}{AK07-J0121-L30}{der Mitte der Finsternis, also war der Beginn derselben, wenn wir die Sekunden}{}
\EdL{31}{AK07-J0121-L31}{vernachlässigen, am 30. Juni 763, 0\textsuperscript{h} 10\textsuperscript{m} Morgens, das Ende 3\textsuperscript{h} 50\textsuperscript{m} Morgens. — Nach}{}
\EdL{32}{AK07-J0121-L32}{\Sperr{Oppolzer}\textsuperscript{2}) ergeben sich folgende Daten (für Greenwich)\,: Mitte der Finsternis}{}
\EdL{33}{AK07-J0121-L33}{30. Juni 763, 1\textsuperscript{h} 42\textsuperscript{m} Morgens\,; Grösse = 22,3 Zoll\,; Dauer der partialen Verfinsterung 2\textsuperscript{h}\,;}{}
\EdL{34}{AK07-J0121-L34}{Dauer der Totalität 1\textsuperscript{h} 44\textsuperscript{m}\,; Dauer der ganzen Finsternis 3\textsuperscript{h} 44\textsuperscript{m}. Es ist dies gewiss}{}
\EdL{35}{AK07-J0121-L35}{ein schönes Ergebnis mit Hilfe dieser über 1000 Jahre alten Tafeln der Araber!}{}
\EdLp{36}{AK07-J0121-L36}{Wir haben hier noch einiges über die Bestimmung der Grösse des \textit{Erdschat-}}{}
\EdL{37}{AK07-J0121-L37}{\textit{tens} bei Mondfinsternissen hinzuzufügen. \Sperr{Al-Bīrūnī} macht in der oben (p. 79)}{}
\EdL{38}{AK07-J0121-L38}{zitierten Stelle noch die weitere Bemerkung, dass die Methode der Berechnung der}{}
\NoteRule
\EdN{N1}{AK07-J0121-N01}{\textsuperscript{1}) Die entsprechenden Zahlen nach den Vorschriften u. Tafeln \Sperr{Batt.}’s sind\,: 31′ 41″ u. 25′ 0″.}
\EdN{N2}{AK07-J0121-N02}{\textsuperscript{2}) \textit{Kanon der Finsternisse} (Denkschriften der Wiener Akademie, mathem. Klasse, LII. Bd.),}
\EdN{N3}{AK07-J0121-N03}{p. 355, Nr. 3050 der Mondfinsternisse.}
\FrameHead{0122}{p. 90}{0122}{p. 90}
\EdL{1}{AK07-J0122-L01}{Grösse des Erdschattens bei Mondfinsternissen ebenfalls dieselbe sei im \textit{Khaṇḍakhād-}}{}
\EdL{2}{AK07-J0122-L02}{\textit{yaka} wie bei \Sperr{Kh.} Von dieser Methode finden wir aber in unserm Werke gar nichts,}{}
\EdL{3}{AK07-J0122-L03}{wir haben nur in den Tafeln 61—66 je zwei Kolonnen (6. u. 7.), die überschrieben}{}
\EdL{4}{AK07-J0122-L04}{sind\,: \textit{Primum dimidium rotae draconis,} und \textit{Secundum dimidium rotae draconis,} mit}{}
\EdL{5}{AK07-J0122-L05}{denen wir uns sogleich nachher beschäftigen werden. Dass jene Regeln für die}{}
\EdL{6}{AK07-J0122-L06}{Bestimmung der Grösse des Erdschattens in diesem Werke fehlen, kann seinen Grund}{}
\EdL{7}{AK07-J0122-L07}{entweder darin haben, dass \Sperr{Al-Bīrūnī} eine andere Ausgabe der Tafel \Sperr{Kh.}’s vor}{}
\EdL{8}{AK07-J0122-L08}{sich gehabt hat als \Sperr{Maslama al-Madjrīṭī}, oder, was wahrscheinlicher ist, dass}{}
\EdL{9}{AK07-J0122-L09}{das vorliegende Werk stark gekürzt worden ist, ob durch \Sperr{Maslama}, oder den}{}
\EdL{10}{AK07-J0122-L10}{Übersetzer, oder erst durch spätere Abschreiber, bleibt unentschieden.}{}
\EdLp{11}{AK07-J0122-L11}{In Bezug auf die 6. und 7. Kolonne der Tafeln 61—66 habe ich nun folgendes}{}
\EdL{12}{AK07-J0122-L12}{gefunden\,: Die Zahlen der 6. Kolonne sind das 5-fache der entsprechenden Werte}{}
\EdL{13}{AK07-J0122-L13}{der stündlichen ungleichen Bewegung (\textit{buht}) der Sonne, die in Kolonne 2 verzeichnet}{}
\EdL{14}{AK07-J0122-L14}{sind, die Zahlen der 7. Kolonne sind das 1\textsuperscript{3}/\textsubscript{5}-fache der entsprechenden Werte der}{}
\EdL{15}{AK07-J0122-L15}{stündlichen ungleichen Bewegung des Mondes, die in Kolonne 3 gegeben sind, sub-}{}
\EdL{16}{AK07-J0122-L16}{trahiert man die Werte der Kolonne 6 von den entsprechenden der Kolonne 7, so}{}
\EdL{17}{AK07-J0122-L17}{erhält man jeweilen die den Anomalien der Kolonne 1 entsprechenden Werte für}{}
\EdL{18}{AK07-J0122-L18}{den Radius der Erdschattens. Wenn also \Sperr{Kh.} diese Berechnungen nach dem \textit{Khaṇ-}}{}
\EdL{19}{AK07-J0122-L19}{\textit{ḍakhādyaka} ausgeführt hat, wie \Sperr{Al-Bīrūnī} berichtet, so muss die Regel des indischen}{}
\EdL{20}{AK07-J0122-L20}{Werkes ungefähr so gelautet haben\,: Man subtrahiere das 5-fache des stündlichen}{}
\EdL{21}{AK07-J0122-L21}{\textit{Buht} der Sonne vom 1\textsuperscript{3}/\textsubscript{5}-fachen des stündlichen \textit{Buht} des Mondes, so hat man den}{}
\EdL{22}{AK07-J0122-L22}{scheinbaren Radius des Erdschattens. \textit{Beispiel\,:} Bei einer Anomalie von 3\textsuperscript{sig.} = 90°}{}
\EdL{23}{AK07-J0122-L23}{(Tafel 63 unterste Zeile) ist der stündliche \textit{Buht} der Sonne (\textit{motus solis in hora}) =}{}
\EdL{24}{AK07-J0122-L24}{2′ 28″, derjenige des Mondes = 32′ 56″\,; die erstere Zahl mit 5 multipliziert gibt 12′ 20″,}{}
\EdL{25}{AK07-J0122-L25}{und dies ist die entsprechenden Zahl der 6. Kolonne, d. h. das \textit{Primum dimidium}}{}
\EdL{26}{AK07-J0122-L26}{\textit{rotae draconis}\,; die zweite Zahl mit 1\textsuperscript{3}/\textsubscript{5} multlpliziert gibt 52′ 42″, die entsprechende}{}
\EdL{27}{AK07-J0122-L27}{Zahl der 7. Kolonne, d. h. das \textit{Secundum dimidium rotae draconis} ist = 52′ 47″\textsuperscript{1})\,; der}{}
\EdL{28}{AK07-J0122-L28}{scheinbare Halbmesser des Erdschattens ist also bei dieser Anomalie = 52′ 47″ — 12′ 20″}{}
\EdL{29}{AK07-J0122-L29}{= 40′ 27″.}{}
\EdLp{30}{AK07-J0122-L30}{\Sperr{Ptol.}\textsuperscript{2}) und nach ihm \Sperr{Batt.}\textsuperscript{3}) geben an, dass der Radius des Erdschattens das}{}
\EdL{31}{AK07-J0122-L31}{c. 2\textsuperscript{3}/\textsubscript{5}-fache von dem des scheinbaren Mondradius sei\,; bei \Sperr{Kh.} ist dieser Faktor}{}
\EdL{32}{AK07-J0122-L32}{nicht ganz konstant, nach verschiedenen Proben, die wir gemacht haben, schwanken}{}
\EdL{33}{AK07-J0122-L33}{seine Werte zwischen 2,5 und 2,6, weichen also von dem \Sperr{Ptol.}’schen Werte wenig}{}
\EdL{34}{AK07-J0122-L34}{ab. Dies führt uns auf die Vermutung, \Sperr{Brahmagupta} möchte in seinem \textit{Khaṇḍa-}}{}
\EdL{35}{AK07-J0122-L35}{\textit{khādyaka} in dieser Frage vielleicht mehr dem \Sperr{Ptol.} gefolgt sein als seinen indischen}{}
\EdL{36}{AK07-J0122-L36}{Vorläufern, wenigstens zeigen die Regeln des \textit{Sūrya-Siddhānta}\textsuperscript{4}) und der \textit{Pancha-}}{}
\EdL{37}{AK07-J0122-L37}{\textit{siddhāntikā}\textsuperscript{5}) für die Bestimmung der Grösse des Erdschattens keine grosse Ver-}{}
\NoteRule
\EdN{N1}{AK07-J0122-N01}{\textsuperscript{1}) Die Abweichung beträgt 5″, es ist dies die grösste, die wir bei unsern Proben gefunden haben.}
\EdN{N2}{AK07-J0122-N02}{\textsuperscript{2}) Ed. \Sperr{Heiberg}, I, 421—423, Übersetzung von \Sperr{Manitius}, I, p. 309—310.}
\EdN{N3}{AK07-J0122-N03}{\textsuperscript{3}) Ed. \Sperr{Nallino}, I, 59—60.}
\EdN{N4}{AK07-J0122-N04}{\textsuperscript{4}) Transl. \Sperr{Burgess}, p. 273.}
\EdN{N5}{AK07-J0122-N05}{\textsuperscript{5}) Ed. \Sperr{Thibaut}, p. 63—65 der Übersetzung.}
\FrameHead{0123}{p. 91}{0123}{p. 91}
\EdL{1}{AK07-J0123-L01}{wandschaft mit der \Sperr{Kh.}’s. Im ersten Werke werden die wahren Durchmesser von}{}
\EdL{2}{AK07-J0123-L02}{Erde, Sonne und Mond und die mittlern und wahren Bewegungen der beiden letztern}{}
\EdL{3}{AK07-J0123-L03}{Himmelskörper zur Berechnung herbeigezogen, und das Resultat gibt den wahren}{}
\EdL{4}{AK07-J0123-L04}{Durchmesser des Erdschattens in \textit{Yojanas}\textsuperscript{1}) an, die allerdings dann leicht in Minuten}{}
\EdL{5}{AK07-J0123-L05}{verwandelt werden können. Im zweiten Werke heisst die Regel\,: „Man multipliziere die}{}
\EdL{6}{AK07-J0123-L06}{wahre Distanz des Mondes mit 36, und teile dieses Produkt durch die wahre Distanz}{}
\EdL{7}{AK07-J0123-L07}{der Sonne multipliziert mit 90 und geteilt durch 276, diesen Quotienten subtrahiere}{}
\EdL{8}{AK07-J0123-L08}{man von 36, so habe man den Durchmesser des Erdschattens (in \textit{Yojanas}).“}{}
\EdLc{9}{AK07-J0123-L09}{\textbf{Kap. 34. Über die Parallaxen des Mondes in Länge und Breite.}}{}
\EdLc{10}{AK07-J0123-L10}{(Text p. 28—29, Tafeln 77 u. 77\textsuperscript{a}.)}{}
\EdLp{11}{AK07-J0123-L11}{Die Regel \Sperr{Kh.}’s für die Bestimmung der \textit{Parallaxen} des Mondes in Länge und}{}
\EdL{12}{AK07-J0123-L12}{Breite sind, besonders die erstere, meiner Ansicht nach etwas verdorben, zum min-}{}
\EdL{13}{AK07-J0123-L13}{desten recht unklar. Zur Orientierung schicken wir voraus, dass, da die Parallaxen-}{}
\EdL{14}{AK07-J0123-L14}{bestimmung in diesem Werke nur zur Berechnung der Sonnenfinsternisse gebraucht}{}
\EdL{15}{AK07-J0123-L15}{wird, bei der Bestimmung der Parallaxe in der Länge der Mond in der Ekliptik}{}
\EdL{16}{AK07-J0123-L16}{angenommen, seine Breite also gleich Null gesetzt wird. Ferner ist zu beachten, dass,}{}
\EdL{17}{AK07-J0123-L17}{wenn der Mond zur Zeit der Konjunktion im Meridian des Beobachtungsortes sich}{}
\EdL{18}{AK07-J0123-L18}{befindet, seine Parallaxe in der Länge gleich Null ist.}{}
\EdLp{19}{AK07-J0123-L19}{Was \Sperr{Batt.} anbetrifft, so ist seine Behandlung dieses Gegenstandes eine sehr}{}
\EdL{20}{AK07-J0123-L20}{eingehende und komplizierte\textsuperscript{2})\,; auf das einzelne kann ich hier nicht eintreten, ich}{}
\EdL{21}{AK07-J0123-L21}{bemerke nur zusammenfassend, dass er in erster Linie die \textit{Zenitdistanz z} des Mondes}{}
\EdL{22}{AK07-J0123-L22}{braucht, die er nach der Formel findet\,:}{}
\EdLc{23}{AK07-J0123-L23}{cos \textit{z} = sin \textit{h} = \Frac{sin \textit{EH} · sin \textit{MB}}{sin \textit{EB}} ,}{}
\EdL{24}{AK07-J0123-L24}{die sich aus den Dreiecken \textit{EHB} und \textit{MAB}}{}
\EdL{25}{AK07-J0123-L25}{nach der Regel der vier Grössen sofort er-}{}
\EdL{26}{AK07-J0123-L26}{gibt. \textit{EH} ist die Ekliptikhöhe, aus der Dekli-}{}
\EdL{27}{AK07-J0123-L27}{nation der Sonne und der geographischen}{}
\EdL{28}{AK07-J0123-L28}{Breite des Beobachtungsortes leicht zu}{}
\EdL{29}{AK07-J0123-L29}{finden\,; \textit{E} ist das sog. \textit{Medium Coelum,}}{}
\EdL{30}{AK07-J0123-L30}{\textit{EB} die Entfernung des \textit{Medium Coelum}}{}
\EdL{31}{AK07-J0123-L31}{vom \textit{Aszendenten B, MB} die Entfernung}{}
\EdL{32}{AK07-J0123-L32}{des Mondes vom Aszendenten. Dann be-}{}
\EdL{33}{AK07-J0123-L33}{rechnet \Sperr{Batt.} die sog. \textit{Winkel} der \textit{Breite}}{}
\EdL{34}{AK07-J0123-L34}{und \textit{Länge,} d. h. den Winkel, den die}{}
\EdL{35}{AK07-J0123-L35}{Ekliptik mit dem durch den Mond ge-}{}
\EdL{36}{AK07-J0123-L36}{henden Höhenkreis bildet, und sein Kom-}{}
\EdL{37}{AK07-J0123-L37}{plement (also die Winkel \textit{EMZ} und \textit{ZMP,}}{}
\EdL{38}{AK07-J0123-L38}{wenn \textit{Z} der Zenit und \textit{P} der Pol der}{}
\NoteRule
\EdN{N1}{AK07-J0123-N01}{\textsuperscript{1}) Ein \textit{Yojana} = 8 arab. Meilen = c. 16000 Meter.}
\EdN{N2}{AK07-J0123-N02}{\textsuperscript{2}) Ed. \Sperr{Nallino}, I, 76—84, u. 251—265.}
\FrameHead{0124}{p. 92}{0124}{p. 92}
\EdL{1}{AK07-J0124-L01}{Ekliptik ist). Nachdem er dann aus den Tafeln (II, 93—94) mit Hilfe der Zenit-}{}
\EdL{2}{AK07-J0124-L02}{distanz als Argument und verschiedener Korrekturen die Parallaxe des Mondes im}{}
\EdL{3}{AK07-J0124-L03}{Höhenkreis \textit{ZA} bestimmt hat, multipliziert er diese Parallaxe sowohl mit dem Sinus}{}
\EdL{4}{AK07-J0124-L04}{des Winkels der Länge als auch mit dem des Winkels der Breite, das erste Produkt}{}
\EdL{5}{AK07-J0124-L05}{gibt die \textit{Parallaxe} des Mondes in der \textit{Länge}, das zweite diejenige in der \textit{Breite.}}{}
\EdL{6}{AK07-J0124-L06}{Nach \Sperr{Batt.} verhalten sich also die Parallaxen des Mondes in Länge und Breite wie}{}
\EdL{7}{AK07-J0124-L07}{die Sinus der Winkel der Länge und Breite, oder anders ausgedrückt, er zerlegt die}{}
\EdL{8}{AK07-J0124-L08}{Parallaxe der Höhe in ihre Komponenten nach der Länge und Breite.}{}
\EdLp{9}{AK07-J0124-L09}{Die Regeln des \textit{Sūrya-Siddhānta} über die Parallaxenbestimmung befinden sich}{}
\EdL{10}{AK07-J0124-L10}{im 5. Kapitel desselben\textsuperscript{1}). Sie weichen im Prinzip nicht wesentlich von denjenigen}{}
\EdL{11}{AK07-J0124-L11}{\Sperr{Batt.}’s ab, es werden auch trigonometrische Beziehungen benutzt, die sog. \textit{Meridian-}}{}
\EdL{12}{AK07-J0124-L12}{\textit{distanz} und \textit{Ekliptik-Zenitdistanz} berechnet, und aus diesen dann die Parallaxen in}{}
\EdL{13}{AK07-J0124-L13}{Länge und Breite bestimmt.}{}
\EdLp{14}{AK07-J0124-L14}{Die Regeln \Sperr{Kh.}’s sind viel einfacher\,; da er aber nur den Gebrauch der Tafeln}{}
\EdL{15}{AK07-J0124-L15}{erklärt und gar keine trigonometrische Ableitungen gibt, so ist es schwer zu ent-}{}
\EdL{16}{AK07-J0124-L16}{scheiden, auf welchem Wege er die Zahlen seiner Tafeln gefunden hat. Sein Prinzip}{}
\EdL{17}{AK07-J0124-L17}{ist kurz ausgedrückt folgendes\,: Die \textit{Parallaxe} in der \textit{Länge} ist eine Funktion der}{}
\EdL{18}{AK07-J0124-L18}{Abweichung des Mondes vom Meridian nach Osten oder Westen, d. h. nach unsrer}{}
\EdL{19}{AK07-J0124-L19}{heutigen Ausdrucksweise eine \textit{Funktion} des \textit{Stundenwinkels,} die \textit{Parallaxe} in der \textit{Breite}}{}
\EdL{20}{AK07-J0124-L20}{ist eine Funktion seiner \textit{Zenitdistanz.} Die Zahlen in der ersten Kolonne (Tafel 77)}{}
\EdL{21}{AK07-J0124-L21}{bedeuten also zweierlei\,: Handelt es sich um die Bestimmung der \textit{Parallaxe} in der}{}
\EdL{22}{AK07-J0124-L22}{\textit{Länge,} zo bezeichnen sie den Abstand des Mondes vom Meridian des Beobachtungs-}{}
\EdL{23}{AK07-J0124-L23}{ortes auf dem Aequator (oder einem Parallelkreis) gemessen, also den \textit{Stundenwinkel}}{}
\EdL{24}{AK07-J0124-L24}{des Mondes\,; handelt es sich um die Bestimmung der \textit{Parallaxe} in der \textit{Breite,} so}{}
\EdL{25}{AK07-J0124-L25}{bedeuten sie die Zenitdistanz des Mondes, und zwar seine sog. \textit{Ekliptik-Zenitdistanz,}}{}
\EdL{26}{AK07-J0124-L26}{d. h. wenn er im \textit{Medium Coelum} ist. — In dieser Hinsicht scheint sich also \Sperr{Kh.}}{}
\EdL{27}{AK07-J0124-L27}{mehr den Indern als \Sperr{Batt.} zu nähern.}{}
\EdLp{28}{AK07-J0124-L28}{Der Text von Kap. 34 verlangt zuerst die Bestimmung des \textit{Aszendenten} (\textit{Horo-}}{}
\EdL{29}{AK07-J0124-L29}{\textit{scopus}) für die gegebene Zeit der Konjunktion, d. h. die Länge des Punktes der}{}
\EdL{30}{AK07-J0124-L30}{Ekliptik, der, um unser oben (p. 82—84) gegebenes \textit{Beispiel} für eine \textit{Konjunktion}}{}
\EdL{31}{AK07-J0124-L31}{von Sonne und Mond wieder aufzunehmen, am 29. Muḥarram des J. 288 d. H. um}{}
\EdL{32}{AK07-J0124-L32}{18\textsuperscript{h} 44\textsuperscript{m} 29\textsuperscript{s} aufgeht. Der Text sagt gar nichts darüber, wie dieser Aszendent bestimmt}{}
\EdL{33}{AK07-J0124-L33}{werde\,; er könnte auch mit den Tafeln, die dieses Werk enthält, unserer Ansicht}{}
\EdL{34}{AK07-J0124-L34}{nach nicht gefunden werden, es fehlen nämlich die Tafeln der schiefen Aufsteigungen,}{}
\EdL{35}{AK07-J0124-L35}{von denen wir wenigstens diejenigen für die Breite von Cordova kennen sollten\,;}{}
\EdL{36}{AK07-J0124-L36}{immerhin könnte er auf dem Wege der Rechnung mit Hilfe der Aszensional-Differenz}{}
\EdL{37}{AK07-J0124-L37}{bestimmt werden. Wir haben daher den \textit{Aszendenten} nach den Vorschriften \Sperr{Batt.}’s\textsuperscript{2})}{}
\EdL{38}{AK07-J0124-L38}{bestimmt und ihn = 297° 3′ gefunden. Auch die Regel für die Auffindung des}{}
\EdL{39}{AK07-J0124-L39}{\textit{Medium Coelum} scheint bei \Sperr{Kh.} fehlerhaft zu sein, zum mindesten ist es unklar, von}{}
\EdL{40}{AK07-J0124-L40}{was man die 90° zu subtrahieren habe, ob vom \textit{Aszendenten} selbst, oder von seiner}{}
\NoteRule
\EdN{N1}{AK07-J0124-N01}{\textsuperscript{1}) \textit{Transl.} \Sperr{Burgess}, p. 285—295.}
\EdN{N2}{AK07-J0124-N02}{\textsuperscript{2}) Ed. \Sperr{Nallino}, I, 73.}
\FrameHead{0125}{p. 93}{0125}{p. 93}
\EdL{1}{AK07-J0125-L01}{\textit{Rectascension} oder von seiner \textit{Ascensio obliqua}\,; wir haben daher dasselbe ebenfalls}{}
\EdL{2}{AK07-J0125-L02}{nach \Sperr{Batt.} bestimmt, und = 228° 40′ gefunden. Nun sagen die Regeln \Sperr{Kh.}’s, man}{}
\EdL{3}{AK07-J0125-L03}{müsse sowohl mit der \textit{wahren Länge} des Mondes (= 303° 20′), als auch mit der des}{}
\EdL{4}{AK07-J0125-L04}{\textit{Medium Coelum} (= 228° 40′) in die Tafeln 59—59\textsuperscript{b}, mit denen man die Längen in}{}
\EdL{5}{AK07-J0125-L05}{Rectascensionen und umgekehrt verwandelt, eintreten\,; tut man dies, so findet man für}{}
\EdL{6}{AK07-J0125-L06}{die Rectascension des Mondes 305° 43′, für die des Medium Coelum 226° 10′, die}{}
\EdL{7}{AK07-J0125-L07}{Differenz beider = 79° 33′ ist der \textit{Stundenwinkel} des Mondes zu der gegebenen Zeit.\textsuperscript{1})}{}
\EdL{8}{AK07-J0125-L08}{Zu besserem Verständnis gebe ich hier bei-}{}
\EdL{9}{AK07-J0125-L09}{stehende Figur, in welcher, damit sie mit}{}
\EdL{10}{AK07-J0125-L10}{unserem Beispiel besser stimme, der auf der}{}
\EdL{11}{AK07-J0125-L11}{hintern Halbkugel liegende Schnittpunkt von}{}
\EdL{12}{AK07-J0125-L12}{Aequator und Ekliptik als Frühlingspunkt,}{}
\EdL{13}{AK07-J0125-L13}{der auf der vordern liegende also als Herbst-}{}
\EdL{14}{AK07-J0125-L14}{punkt \textit{W} angenommen werden musste\,; dann}{}
\EdL{15}{AK07-J0125-L15}{ist also (wenn \textit{HT} = Horizont, \textit{AQ} = Aequa-}{}
\EdL{16}{AK07-J0125-L16}{tor, \textit{EK} = Ekliptik, \textit{L} = Aszendent, \textit{M} =}{}
\EdL{17}{AK07-J0125-L17}{Mond) 180° + \textit{WM} die Länge des Mondes,}{}
\EdL{18}{AK07-J0125-L18}{180° + \textit{WR} seine Rectascension, 180° + \textit{WL}}{}
\EdL{19}{AK07-J0125-L19}{die Länge des Aszendenten, 180° + \textit{WN} seine}{}
\EdL{20}{AK07-J0125-L20}{Rectascension, 180° — \textit{EW} die Länge des}{}
\EdL{21}{AK07-J0125-L21}{Medium Coelum, 180° — \textit{AW} seine Rect-}{}
\EdL{22}{AK07-J0125-L22}{ascension, mithin ist 180° + \textit{WR} — (180° — \textit{AW})}{}
\EdL{23}{AK07-J0125-L23}{= \textit{WR} + \textit{AW} = \textit{AR} = ∠ \textit{APR} = \textit{s} = Stunden-}{}
\EdL{24}{AK07-J0125-L24}{winkel des Mondes. — Mit dem Stundenwinkel des Mondes (79° 33′) tritt man nun}{}
\EdL{25}{AK07-J0125-L25}{in die Tafel der Parallaxen (Tafel 77) ein und findet daselbst für die \textit{Längenparallaxe}}{}
\EdL{26}{AK07-J0125-L26}{c. 1\textsuperscript{h} 33\textsuperscript{1}/\textsubscript{2}\textsuperscript{m}\,; diese Grösse multipliziert man mit der ungleichen stündlichen Bewegung}{}
\EdL{27}{AK07-J0125-L27}{des Mondes, die für die Anomalie 125° 7′ = 34′ 34″ ist, man erhält so c. 54′. — \Sperr{Kh.}}{}
\EdL{28}{AK07-J0125-L28}{gibt also in seinen Tafeln die Parallaxe der Länge in Zeit an, \Sperr{Batt.} dagegen in}{}
\EdL{29}{AK07-J0125-L29}{Kreisteilen\,; in unserm Falle ist also der scheinbare Ort des Mondes, wie er sich}{}
\EdL{30}{AK07-J0125-L30}{unsern Augen zeigt, von dem wahren, wie er sich aus der Berechnung und den}{}
\EdL{31}{AK07-J0125-L31}{Tafeln ergeben hat, um 54′ entfernt, und zwar muss man nach den Vorschriften}{}
\EdL{32}{AK07-J0125-L32}{des Textes diese Parallaxe zum wahren Ort des Mondes (303° 19′ 48″) addieren, um}{}
\EdL{33}{AK07-J0125-L33}{den scheinbaren Ort des Mondes zu erhalten, dieser ist also 304° 13′ 48″.}{}
\EdLp{34}{AK07-J0125-L34}{Die Regel für die Auffinduug der \textit{Parallaxe} in der \textit{Breite} ist nun ein wenig}{}
\EdL{35}{AK07-J0125-L35}{einfacher, da die Fragen des \textit{Aszendenten} und des \textit{Medium Coelum} für uns erledigt}{}
\EdL{36}{AK07-J0125-L36}{sind. Es wird nun verlangt, die \textit{Obliquatio} des \textit{Medium Coelum,} d. h. die Deklination}{}
\EdL{37}{AK07-J0125-L37}{des Mondes (= der der Sonne) zu jener Zeit zu bestimmen, diese ist nach Tafel 22}{}
\EdL{38}{AK07-J0125-L38}{= — 19° 44′ 38″, die Breite von Cordova nehmen wir rund 37° 50′ an, dies gibt für die}{}
\EdL{39}{AK07-J0125-L39}{Ekliptik-Zenitdistanz \textit{EZ} = 57° 34′ 38″, unser Text nennt dies das \textit{interstitium medii}}{}
\NoteRule
\EdN{N1}{AK07-J0125-N01}{\textsuperscript{1}) Auf trigonometrischem Wege finden wir diesen Stundenwinkel = 79° 36′.}
\EdN{N2}{AK07-J0125-N02}{\textsuperscript{2}) Denn der Mond befindet sich zwischen Ascendens und Medium Coelum, oder ist, wie sich}
\EdN{N3}{AK07-J0125-N03}{\Sperr{Batt.} (I, 81 u. 83) deutlicher ausdrückt, vom Ascendens weniger als 90° entfernt.}
\FrameHead{0126}{p. 94}{0126}{p. 94}
\EdL{1}{AK07-J0126-L01}{\textit{coeli huius horoscopi et imminentis capitibus nostris.} Nun ist die \textit{Breite} des Mondes zu}{}
\EdL{2}{AK07-J0126-L02}{bestimmen, das korrigierte Argument derselben ist nach p. 84 = 169° 7\textsuperscript{1}/\textsubscript{2}′, dazu gehört}{}
\EdL{3}{AK07-J0126-L03}{nach Tafel 26 die Breite 50′ 53″, sie ist nördlich, ist also von der vorhin gefundenen}{}
\EdL{4}{AK07-J0126-L04}{Ekliptik-Zenitdistanz zu subtrahieren, dies gibt 56° 43′ 45″\,; mit diesem Argument}{}
\EdL{5}{AK07-J0126-L05}{tritt man in Tafel 77 ein, und findet daselbst die \textit{Parallaxe} der \textit{Breite} = 40′ 43″,}{}
\EdL{6}{AK07-J0126-L06}{und zwar ist sie \textit{südlich} zu nehmen, weil das \textit{Medium Coelum} auch südlich vom}{}
\EdL{7}{AK07-J0126-L07}{Zenit von Cordova gelegen ist.}{}
\EdLp{8}{AK07-J0126-L08}{Nach den Tafeln \Sperr{Theons} (bei \Sperr{Batt.} II, 98) findet man mit einer Korrektur}{}
\EdL{9}{AK07-J0126-L09}{(s. \Sperr{Batt.} I, 83) die Parallaxe der Länge c. 35′, die der Breite c. 47′\,; ungefähr dieselben}{}
\EdL{10}{AK07-J0126-L10}{Resultate ergeben sich nach den Regeln, die \Sperr{Batt.} I, 79—81 aufstellt. Wir sehen}{}
\EdL{11}{AK07-J0126-L11}{also, dass die Parallaxe der Länge bei \Sperr{Kh.} wesentlich grösser ist als bei \Sperr{Batt.},}{}
\EdL{12}{AK07-J0126-L12}{auch die in der Breite weicht um 6′ ab\,; \Sperr{Batt.} bemerkt aber selbst, dass diese}{}
\EdL{13}{AK07-J0126-L13}{\Sperr{Theon}’schen Tafeln keine ganz befriedigende Resultate ergeben können.}{}
\EdLc{14}{AK07-J0126-L14}{\textbf{Kap. 35. Über die Verfinsterung der Sonne.}}{}
\EdLc{15}{AK07-J0126-L15}{(Text p. 29, Tafel 78.)}{}
\EdLp{16}{AK07-J0126-L16}{Dieses Kapitel ist sehr kurz und jedenfalls unvollständig. Aus Tafel 78 ersehen}{}
\EdL{17}{AK07-J0126-L17}{wir, dass \Sperr{Kh.} die Grenzen des Argumentes der Breite, innerhalb welcher eine Fin-}{}
\EdL{18}{AK07-J0126-L18}{sternis überhaupt stattfinden kann, folgendermassen annimmt\,: von 172° 49′ bis}{}
\EdL{19}{AK07-J0126-L19}{187° 30′, von 0° bis 7° 11′ und 352° 30′ bis 360°\,; \Sperr{Batt.} hat (II, 88) folgende Zahlen\,:}{}
\EdL{20}{AK07-J0126-L20}{159° 44′ bis 191° 16′, 0° bis 20° 16′ und 348° 44′ bis 360°.}{}
\EdLp{21}{AK07-J0126-L21}{Es wird in diesem Kapitel nichts darüber gesagt, wie man die Tafeln \textit{pro spatio}}{}
\EdL{22}{AK07-J0126-L22}{\textit{longinquo} und \textit{pro spatio propinquo} zu benutzen habe, wir nehmen an, dass dies auf}{}
\EdL{23}{AK07-J0126-L23}{ähnliche Weise geschehen müsse wie bei den Tafeln der Mondfinsternisse und wie}{}
\EdL{24}{AK07-J0126-L24}{bei \Sperr{Batt.}\textsuperscript{1}), doch fehlen in Tafel 78 die \textit{Tabulae proportionis,} die sich in den Mond-}{}
\EdL{25}{AK07-J0126-L25}{tafeln finden\,; wahrscheinlich aber sind dieselben Proportionaltafeln zu benutzen}{}
\EdL{26}{AK07-J0126-L26}{wie dort (Tafeln 73—76).}{}
\EdLp{27}{AK07-J0126-L27}{Da sich nun das Argument der Breite bei unserer Konjunktion (s. p. 82—84),}{}
\EdL{28}{AK07-J0126-L28}{nämlich 169° 7\textsuperscript{1}/\textsubscript{2}′, nicht in Tafel 78 vorfindet, so müssen wir schliessen, es finde bei}{}
\EdL{29}{AK07-J0126-L29}{dieser Konjunktion keine Sonnenfinsternis für Cordova statt. Nach den Tafeln \Sperr{Batt.}’s}{}
\EdL{30}{AK07-J0126-L30}{kann aber eine solche stattfinden für andere Orte, und wir wollen nun die weitere}{}
\EdL{31}{AK07-J0126-L31}{Rechnung nach diesen Tafeln durchführen, wir nehmen aber die bis jetzt nach}{}
\EdL{32}{AK07-J0126-L32}{den Tafeln \Sperr{Kh.}’s gefundenen Werte\,:}{}
\EdL{33}{AK07-J0126-L33}{\parbox[t]{\dimexpr0.2450\linewidth-0.5em\relax}{\raggedright\small\mbox{}}\hspace{0.5em}\parbox[t]{\dimexpr0.2450\linewidth-0.5em\relax}{\raggedright\small\mbox{}Breite des Mondes}\hspace{0.5em}\parbox[t]{\dimexpr0.4900\linewidth-0.5em\relax}{\raggedright\small\mbox{}50′ 53″,}\hspace{0.5em}}{}
\EdL{34}{AK07-J0126-L34}{\parbox[t]{\dimexpr0.2450\linewidth-0.5em\relax}{\raggedright\small\mbox{}}\hspace{0.5em}\parbox[t]{\dimexpr0.2450\linewidth-0.5em\relax}{\raggedright\small\mbox{}Parallaxe der Breite}\hspace{0.5em}\parbox[t]{\dimexpr0.4900\linewidth-0.5em\relax}{\raggedright\small\mbox{}40′ 43″.}\hspace{0.5em}}{}
\EdLp{35}{AK07-J0126-L35}{Da die \textit{Breite nördlich} ist, die \textit{Parallaxe} aber \textit{südlich} zu nehmen ist, müssen}{}
\EdL{36}{AK07-J0126-L36}{wir die Parallaxe von der Breite subtrahieren, dies gibt 10′ 10″, dies ist die schein-}{}
\EdL{37}{AK07-J0126-L37}{bare Breite des Mondes bei der für uns sichtbaren Konjunktion\textsuperscript{2}). Für diese \textit{schein-}}{}
\NoteRule
\EdN{N1}{AK07-J0126-N01}{\textsuperscript{1}) Ed. \Sperr{Nallino}, I, 111, und II, 91.}
\EdN{N2}{AK07-J0126-N02}{\textsuperscript{2}) Nach \Sperr{Batt.} findet man\,: scheinbare Breite = 58′ 25″ — 47′ = 11′ 25″\,; nach den \Sperr{Theon}’schen}
\EdN{N3}{AK07-J0126-N03}{Tafeln lässt sich aber die Parallaxe nicht genau bestimmen, es wäre also möglich, dass doch die}
\EdN{N4}{AK07-J0126-N04}{beiden Werte für die scheinbare Breite nahezu übereinstimmen könnten.}
\FrameHead{0127}{p. 95}{0127}{p. 95}
\EdL{1}{AK07-J0127-L01}{\textit{bare Breite} finden wir nun nach den Tafeln \Sperr{Batt.}’s (II, 91) folgende Daten\,:}{}
\EdL{2}{AK07-J0127-L02}{\parbox[t]{\dimexpr0.4900\linewidth-0.5em\relax}{\raggedright\small\mbox{}in maxima longinquitate lunae\,:}\hspace{0.5em}\parbox[t]{\dimexpr0.4900\linewidth-0.5em\relax}{\raggedright\small\mbox{}in minima longinquitate lunae\,:}\hspace{0.5em}}{}
\EdL{3}{AK07-J0127-L03}{\parbox[t]{\dimexpr0.2450\linewidth-0.5em\relax}{\raggedright\small\mbox{}Digiti diametri solaris\,:}\hspace{0.5em}\parbox[t]{\dimexpr0.2450\linewidth-0.5em\relax}{\raggedright\small\mbox{}Quantitas incidentiae\,:}\hspace{0.5em}\parbox[t]{\dimexpr0.2450\linewidth-0.5em\relax}{\raggedright\small\mbox{}Digiti diametri solaris\,:}\hspace{0.5em}\parbox[t]{\dimexpr0.2450\linewidth-0.5em\relax}{\raggedright\small\mbox{}Quantitas incidentiae\,:}\hspace{0.5em}}{}
\EdL{4}{AK07-J0127-L04}{\parbox[t]{\dimexpr0.2450\linewidth-0.5em\relax}{\raggedright\small\mbox{}7\textsuperscript{dig.} 42′}\hspace{0.5em}\parbox[t]{\dimexpr0.2450\linewidth-0.5em\relax}{\raggedright\small\mbox{}29′ 15″}\hspace{0.5em}\parbox[t]{\dimexpr0.2450\linewidth-0.5em\relax}{\raggedright\small\mbox{}8\textsuperscript{dig.} 48′}\hspace{0.5em}\parbox[t]{\dimexpr0.2450\linewidth-0.5em\relax}{\raggedright\small\mbox{}32′ 26″}\hspace{0.5em}}{}
\EdLp{5}{AK07-J0127-L05}{Die Differenzen der entsprechende Werte sind\,:}{}
\EdLc{6}{AK07-J0127-L06}{1\textsuperscript{dig.} 6′ und 3′ 11″.}{}
\EdLp{7}{AK07-J0127-L07}{Von diesen muss man nach den Vorschriften (I, 111) Proportionalteile nehmen,}{}
\EdL{8}{AK07-J0127-L08}{und zwar müssen dieselben im Verhältnis 46\textsuperscript{1}/\textsubscript{3} ∶ 60 verkleinert werden (die Zahl 46\textsuperscript{1}/\textsubscript{3},}{}
\EdL{9}{AK07-J0127-L09}{genauer 46′ 17″, wird in der \textit{Tabula correctionis} (II, 89) gefunden, als entsprechend}{}
\EdL{10}{AK07-J0127-L10}{der Anomalie 125° 7′ des Mondes). Führt man diese Korrekturen aus, so werden}{}
\EdL{11}{AK07-J0127-L11}{jene Differenzen verwandelt in 51′, bezw. 2′ 27″, diese Werte zu den Daten für die}{}
\EdL{12}{AK07-J0127-L12}{\textit{maxima longinquitas lunae} (= \textit{pro spatio longinquo}) addiert, gibt}{}
\EdLc{13}{AK07-J0127-L13}{8\textsuperscript{dig.} 33′ und 31′ 42″.}{}
\EdLp{14}{AK07-J0127-L14}{Die Minuten der \textit{quantitas incidentiae} (in unsern Tafeln \textit{casus}) werden durch die}{}
\EdL{15}{AK07-J0127-L15}{ungleiche stündliche Bewegung des Mondes bei der Anomalie von 125° 7′, also durch}{}
\EdL{16}{AK07-J0127-L16}{34′ 34″ dividiert, dies gibt c. 55\textsuperscript{m}\,; wir haben nun folgende Resultate\,:}{}
\EdLp{17}{AK07-J0127-L17}{1) Für die Zeit der wahren Konjunktion haben wir oben (p. 83) gefunden\,:}{}
\EdL{18}{AK07-J0127-L18}{18\textsuperscript{h} 44\textsuperscript{m} 29\textsuperscript{s}, nach den Vorschriften unsers Textes und auch \Sperr{Batt.}’s (I, 105 und 110)}{}
\EdL{19}{AK07-J0127-L19}{müssen wir hievon die Längen-Parallaxe in Zeit, also 1\textsuperscript{h} 33\textsuperscript{1}/\textsubscript{2}\textsuperscript{m}, subtrahieren, dies}{}
\EdL{20}{AK07-J0127-L20}{gibt für die Zeit der \textit{scheinbaren Konjunktion}, d. h. der \textit{Mitte} der Finsternis rund}{}
\EdL{21}{AK07-J0127-L21}{17\textsuperscript{h} 11\textsuperscript{m}, d. h. dieselbe fand statt am 23. Jan. 901, um 5\textsuperscript{h} 11\textsuperscript{m} Morgens (für Cordova).}{}
\EdLp{22}{AK07-J0127-L22}{2) Der \textit{Beginn} der Finsternis fand statt um 5\textsuperscript{h} 11\textsuperscript{m} — 55\textsuperscript{m} = 4\textsuperscript{h} 16\textsuperscript{m}, das \textit{Ende}}{}
\EdL{23}{AK07-J0127-L23}{um 5\textsuperscript{h} 11\textsuperscript{m} + 55\textsuperscript{m} = 6\textsuperscript{h} 6\textsuperscript{m}.}{}
\EdLp{24}{AK07-J0127-L24}{3) Das \textit{Maximum} der Verfinsterung betrug 8\textsuperscript{dig.} 33′, also etwas mehr als \textsuperscript{2}/\textsubscript{3} des}{}
\EdL{25}{AK07-J0127-L25}{Sonnendurchmessers.}{}
\EdLp{26}{AK07-J0127-L26}{Ich bemerke ausdrücklich, dass diese Daten auf Genauigkeit keinen Anspruch}{}
\EdL{27}{AK07-J0127-L27}{machen dürfen, da sie teilweise nach den Tafeln \Sperr{Kh.}’s, teilweise nach denen \Sperr{Batt.}’s}{}
\EdL{28}{AK07-J0127-L28}{berechnet sind. Es wird aber gleichwohl von Interesse sein, die Resultate mit den-}{}
\EdL{29}{AK07-J0127-L29}{jenigen zu vergleichen, die man für diese Finsternis aus dem Kanon \Sperr{Oppolzers}}{}
\EdL{30}{AK07-J0127-L30}{erhält\,: Für Cordova war die grösste Phase (Mitte) der Finsternis um 5\textsuperscript{h} 14\textsuperscript{m} Morgens,}{}
\EdL{31}{AK07-J0127-L31}{und das Ende um 6\textsuperscript{h} 0\textsuperscript{m}. Da aber die Sonne in Cordova am 23. Jan. 901 erst um}{}
\EdL{32}{AK07-J0127-L32}{7\textsuperscript{h} aufging, war die Finsternis daselbst nicht sichtbar\textsuperscript{1}). Also wiederum ein recht}{}
\EdL{33}{AK07-J0127-L33}{schönes Resultat\,!}{}
\EdLp{34}{AK07-J0127-L34}{Vergleichen wir diese Daten noch mit den Angaben \Sperr{Batt.}’s über diese Finsternis\,;}{}
\EdL{35}{AK07-J0127-L35}{er sagt\textsuperscript{2}), die Mitte der Finsternis habe in \textit{Antiochia} 3\textsuperscript{h} 40\textsuperscript{m} vor dem Beginn des 23. Jan.}{}
\EdL{36}{AK07-J0127-L36}{901 (Anfang nach den arab. Astronomen Mittags 12 Uhr) stattgefunden, also nach}{}
\NoteRule
\EdN{N1}{AK07-J0127-N01}{\textsuperscript{1}) Ich verdanke die Berechnung für \textit{Cordova} nach dem Kanon \Sperr{Oppolzers} der Güte des}
\EdN{N2}{AK07-J0127-N02}{Herrn \Sperr{Broger}, Assistenten an der Züricher Sternwarte, wofür ich ihm an dieser Stelle den besten}
\EdN{N3}{AK07-J0127-N03}{Dank ausspreche.}
\EdN{N4}{AK07-J0127-N04}{\textsuperscript{2}) Ed. \Sperr{Nallino}, I, 56—57.}
\FrameHead{0128}{p. 96}{0128}{p. 96}
\EdL{1}{AK07-J0128-L01}{unserer bürgerlichen Zeitrechnung am 23. Jan. 8\textsuperscript{h} 20\textsuperscript{m} Morgens. Nun liegt \textit{Antiochia}}{}
\EdL{2}{AK07-J0128-L02}{c. 41° östlich von \textit{Cordova,} dies macht in Zeit 2\textsuperscript{h} 44\textsuperscript{m}, dies zu 5\textsuperscript{h} 11\textsuperscript{m} addiert gibt 7\textsuperscript{h} 55\textsuperscript{m}}{}
\EdL{3}{AK07-J0128-L03}{als die Zeit der Mitte der Finsternis für \textit{Antiochia}\,; der Unterschied gegenüber der}{}
\EdL{4}{AK07-J0128-L04}{Angabe \Sperr{Batt.}’s ist also c. 25\textsuperscript{m}. Die Grösse der Verfinsterung, sagt \Sperr{Batt.}, habe wenig}{}
\EdL{5}{AK07-J0128-L05}{mehr als die Hälfte, oder in \textit{Raḳḳa} etwas weniger als zwei Drittel betragen. Auch}{}
\EdL{6}{AK07-J0128-L06}{diese Angaben weichen wenig von denen unserer Berechnung ab.}{}
\EdLp{7}{AK07-J0128-L07}{Wie wir oben bemerkt haben, sind die beiden Kapitel 34 und 35 sehr unklar}{}
\EdL{8}{AK07-J0128-L08}{und höchst wahrscheinlich unvollständig und teilweise unrichtig in der Übersetzung}{}
\EdL{9}{AK07-J0128-L09}{wiedergegeben. Wir werden kaum fehlgehen, wenn wir annehmen, dass \Sperr{Athelhard}}{}
\EdL{10}{AK07-J0128-L10}{diese etwas schwierigen astronomischen Partien nicht recht verstanden habe, und}{}
\EdL{11}{AK07-J0128-L11}{spätere Abschreiber natürlich noch viel weniger, daraus ergaben sich teils falsche}{}
\EdL{12}{AK07-J0128-L12}{Auffassungen, teils Auslassungen, von denen ich zwei hier erwähnen und zu korri-}{}
\EdL{13}{AK07-J0128-L13}{gieren versuchen will. In Kap. 34 (p. 28) heisst es\,: \textit{erit hoc productum horis coniunc-}}{}
\EdL{14}{AK07-J0128-L14}{\textit{tionis addendum}\,; statt \textit{horis coniunctionis} muss es wohl heissen \textit{loco coniunctionis,} oder}{}
\EdL{15}{AK07-J0128-L15}{besser \textit{loco vero lunae}\,; denn es handelt sich hier nicht speziell um eine Konjunktion,}{}
\EdL{16}{AK07-J0128-L16}{sondern allgemein um den durch die \textit{Parallaxe} veränderten Ort des Mondes. Und}{}
\EdL{17}{AK07-J0128-L17}{weiter heisst es ebenda\,: \textit{Quodque inde surget, horam eclipsis supra locum pro quo in-}}{}
\EdL{18}{AK07-J0128-L18}{\textit{quiris ex diversitate aspectus mutatam ponit}\,; statt \textit{horam eclipsis} muss es heissen\,:}{}
\EdL{19}{AK07-J0128-L19}{\textit{locum lunae apparentem in ecliptica}\textsuperscript{1}). In Kap. 35 ist der Passus vom Anfang an bis}{}
\EdL{20}{AK07-J0128-L20}{\textit{Si enim fuerit locus coniunctionis} ganz ungenügend für die Bestimmung einer Sonnen-}{}
\EdL{21}{AK07-J0128-L21}{finsternis\,; man wird es mir aber nicht verargen, wenn ich auf einen Versuch den}{}
\EdL{22}{AK07-J0128-L22}{richtigen Text herzustellen verzichte.}{}
\EdLc{23}{AK07-J0128-L23}{\textbf{Kap. 32 u. 36. Über die Verteilung der zwölf astrologischen Häuser.}}{}
\EdLc{24}{AK07-J0128-L24}{(Text p. 25 u. 29—30, Tafeln 79—90.)}{}
\EdLp{25}{AK07-J0128-L25}{Diese beiden Kapitel und das nächstfolgende letzte beschäftigen sich mit rein}{}
\EdL{26}{AK07-J0128-L26}{\textit{astrologischen} Fragen. Kapitel 32 steht nur in den Mss. C u. N und zwar im erstern}{}
\EdL{27}{AK07-J0128-L27}{an unrichtiger Stelle\textsuperscript{2}), es sollte dem Kap. 36 unmittelbar vorangehen, denn es}{}
\EdL{28}{AK07-J0128-L28}{enthält die Vorschriften für die Auffindung der Anfänge der zwölf Häuser auf der}{}
\EdL{29}{AK07-J0128-L29}{Ekliptik, Kap. 36 dagegen die Anleitung zur Benutzung der zugehörigen Tafeln.}{}
\EdLp{30}{AK07-J0128-L30}{Das arabische \textit{taswiya} = Verteilung, Einteilung, Anordnung ist vom Übersetzer}{}
\EdL{31}{AK07-J0128-L31}{teils durch \textit{aequatio,} teils durch das richtigere \textit{dispositio} wiedergegeben worden. Es}{}
\EdL{32}{AK07-J0128-L32}{handelt sich um die Verteilung der zwölf astrologischen Häuser auf die Himmelskugel.}{}
\EdL{33}{AK07-J0128-L33}{Es existierten verschiedene Arten dieser Verteilung, wofür ich auf \Sperr{Batt.}\textsuperscript{3}) verweise\,;}{}
\EdL{34}{AK07-J0128-L34}{alle haben das gemeinsame, dass sich sechs grösste Kreise auf der Himmelskugel in}{}
\EdL{35}{AK07-J0128-L35}{einem Durchmesser schneiden, wodurch dieselbe in zwölf sphärische Zweiecke geteilt}{}
\EdL{36}{AK07-J0128-L36}{wird. Die Einteilung \Sperr{Batt.}’s, die wahrscheinlich auch die \Sperr{Ptolemäische} im \textit{Quadri-}}{}
\EdL{37}{AK07-J0128-L37}{\textit{partitum} ist, und der auch \Sperr{Kh.} folgt, ist folgende\,: Die sechs grössten Kreise gehen}{}
\NoteRule
\EdN{N1}{AK07-J0128-N01}{\textsuperscript{1}) Vergl. \Sperr{Batt.}, ed. \Sperr{Nallino}, I, 83—84.}
\EdN{N2}{AK07-J0128-N02}{\textsuperscript{2}) S. oben p. 86, Note 3.}
\EdN{N3}{AK07-J0128-N03}{\textsuperscript{3}) Ed. \Sperr{Nallino}, I, 246—249.}
\FrameHead{0129}{p. 97}{0129}{p. 97}
\EdL{1}{AK07-J0129-L01}{durch die Pole des Aequators, sind also Meridiane oder Deklinationskreise\,; der erste,}{}
\EdL{2}{AK07-J0129-L02}{durch den \textit{Aszendenten} gehende, markiert den Anfang des I. Hauses\,; um nun den}{}
\EdL{3}{AK07-J0129-L03}{Punkt auf der Ekliptik zu erhalten, durch den der zweite Kreis (vom Aszendenten}{}
\EdL{4}{AK07-J0129-L04}{aus gegen Osten nach der Ordnung der Tierkreiszeichen gezählt) geht, mit dem also}{}
\EdL{5}{AK07-J0129-L05}{das II. Haus beginnt, muss man den Nachtbogen des Aszendenten durch 6 teilen\textsuperscript{1}),}{}
\EdL{6}{AK07-J0129-L06}{das Resultat zur Rectascension des Aszendenten addieren, und zu dieser neuen}{}
\EdL{7}{AK07-J0129-L07}{Rectascension in den Tafeln 59—59\textsuperscript{b} die Länge suchen, so hat man den Grad des}{}
\EdL{8}{AK07-J0129-L08}{Anfangs des II. Hauses, und um 180° vermehrt den Grad des Anfangs des VIII. Hauses,}{}
\EdL{9}{AK07-J0129-L09}{der der Nadir (\textit{Naẓīr}) zu dem des II. ist. Will man den Anfang des III. (bezw. IX.)}{}
\EdL{10}{AK07-J0129-L10}{Hauses haben, so addiert man zu der vorhin gefundenen Rectascension wiederum}{}
\EdL{11}{AK07-J0129-L11}{den 6. Teil des Nachtbogens, so ist die Summe die Rectascension des Anfangs des}{}
\EdL{12}{AK07-J0129-L12}{III. bezw. IX.) Hauses\,; das gleiche gilt auch noch für die Berechnung des Anfangs}{}
\EdL{13}{AK07-J0129-L13}{des IV. (bezw. X.) Hauses. Um den Anfang des V. (bezw. XI.) Hauses zu erhalten,}{}
\EdL{14}{AK07-J0129-L14}{muss man zur Rectascension des IV. Hauses den 6. Teil des Tagebogens des Aszen-}{}
\EdL{15}{AK07-J0129-L15}{denten (= \CaseFrac{360° — Nachtbogen}{6} = 60° — \CaseFrac{Nachtbogen}{6}) addieren, und zu dieser Rectascension}{}
\EdL{16}{AK07-J0129-L16}{wieder die Länge suchen, diese ist der Anfang des V. (bezw. XI.) Hauses. In gleicher}{}
\EdL{17}{AK07-J0129-L17}{Weise verfährt man nun auch für die Auffindung der Anfänge des VI. u. VII. Hauses\,;}{}
\EdL{18}{AK07-J0129-L18}{wenn die Berechnung richtig durchgeführt ist, so muss der Anfang des VII. Hauses}{}
\EdL{19}{AK07-J0129-L19}{der Nadir des Anfangs des I. Hauses, d. h. des Aszendenten sein.}{}
\EdLp{20}{AK07-J0129-L20}{Man sieht also hieraus, dass diese zwölf Häuser von ungleicher Grösse waren,}{}
\EdL{21}{AK07-J0129-L21}{immerhin sind das I., II., III., VII., VIII., IX. unter sich kongruent, ebenso das IV., V.,}{}
\EdL{22}{AK07-J0129-L22}{VI., X., XI., XII. unter sich. Das Eigentümliche ist, dass sowohl auf den Tagebogen,}{}
\EdL{23}{AK07-J0129-L23}{wie auch auf den Nachtbogen kleinere und grössere Häuser fallen. Jede geographische}{}
\EdL{24}{AK07-J0129-L24}{Breite hat ihre besondere Häusergrösse, und diese wechselt wieder mit der Jahres-}{}
\EdL{25}{AK07-J0129-L25}{zeit\,; nur am Aequator sind alle Häuser zu jeder Jahreszeit gleich gross\textsuperscript{2}).}{}
\EdLp{26}{AK07-J0129-L26}{Im Kap. 32 wird dann noch angegeben wie man die Grösse (in Graden) einer}{}
\EdL{27}{AK07-J0129-L27}{\textit{Zeitstunde,} also sowohl einer \textit{Tages-} als einer \textit{Nachtstunde} bestimmen könne\,; man}{}
\EdL{28}{AK07-J0129-L28}{braucht dazu die Tafeln der schiefen Aufsteigungen für die betreffende Breite, diese}{}
\EdL{29}{AK07-J0129-L29}{Tafeln fehlen aber, wie früher schon bemerkt wurde, in diesem Werke. Wir haben}{}
\EdL{30}{AK07-J0129-L30}{die Grösse der Zeitstunde für Cordova, um eine Probe für die Anfänge der zwölf}{}
\EdL{31}{AK07-J0129-L31}{Häuser zu machen, angenähert nach den Tafeln \Sperr{Batt.}’s (II, 67) bestimmt, und für}{}
\EdL{32}{AK07-J0129-L32}{den Morgen des 23. Jan. 901, d. h. für den Aszendenten 297° 3′ (= 27° 3′ des Steinbocks)}{}
\EdL{33}{AK07-J0129-L33}{= 12° 4′ gefunden\textsuperscript{3})\,; man kann allerdings dieselbe auch trigonometrisch berechnen,}{}
\EdL{34}{AK07-J0129-L34}{indem man den Stundenwinkel der Sonne bei ihrem Aufgang an dem betreffenden}{}
\EdL{35}{AK07-J0129-L35}{Tage in Cordova sucht. Die Länge der Doppel-Tagesstunde ist also 24° 8′, die Länge}{}
\EdL{36}{AK07-J0129-L36}{der Doppel-Nachtstunde ist also = 35° 52′.}{}
\NoteRule
\EdN{N1}{AK07-J0129-N01}{\textsuperscript{1}) Dies gibt nach \Sperr{Kh.} die \textit{tempora duarum horarum nocturnarum}\,; unter \textit{tempora} sind hier}
\EdN{N2}{AK07-J0129-N02}{Aequatorialgrade verstanden.}
\EdN{N3}{AK07-J0129-N03}{\textsuperscript{2}) Zur Zeit der Tag- und Nachtgleiche ist dies auch in den andern Breiten der Fall.}
\EdN{N4}{AK07-J0129-N04}{\textsuperscript{3}) Diesem Werte liegt für Cordova die geogr. Breite 38° 30′ zu Grunde\,; wir haben nämlich}
\EdN{N5}{AK07-J0129-N05}{nachträglich in den Tafeln der \textit{proiectio radiorum} gefunden, dass \Sperr{Al-Madjrīṭī} sie so gross an-}
\EdN{N6}{AK07-J0129-N06}{genommen hat.}
\FrameHead{0130}{p. 98}{0130}{p. 98}
\EdLp{1}{AK07-J0130-L01}{Nach den Vorschriften unseres Textes (Kap. 32) hat man zur Rectascension des}{}
\EdL{2}{AK07-J0130-L02}{Aszendenten, die nach Tafel 59 = 29° 10′ ist, die Länge der Doppel-Nachtstunde zu}{}
\EdL{3}{AK07-J0130-L03}{addieren, dies gibt 65° 2′\,; zu dieser Rectascension findet man in derselben Tafel die}{}
\EdL{4}{AK07-J0130-L04}{Länge 3° 1′ der Fische, dies ist der Anfang des II. Hauses vom Aszendenten aus.}{}
\EdL{5}{AK07-J0130-L05}{Addiert man zu der Rectascension 65° 2′ wieder den Betrag der Doppel-Nachtstunde,}{}
\EdL{6}{AK07-J0130-L06}{so erhält man 100° 54′, dazu gehört in derselben Tafel 59 die Länge 11° 54′ des}{}
\EdL{7}{AK07-J0130-L07}{Widders, dies ist der Anfang des III. Hauses. Zu 100° 54′ wieder 35° 52′ addiert,}{}
\EdL{8}{AK07-J0130-L08}{gibt 136° 46′, dazu gehört in Tafel 59\textsuperscript{a} die Länge 19° 19′ des Stiers, dies ist der}{}
\EdL{9}{AK07-J0130-L09}{Anfang des IV. Hauses. — Nun addiert man zu 136° 46′ die Länge der Doppel-}{}
\EdL{10}{AK07-J0130-L10}{Tagesstunde (= 24° 8′), dies gibt 160° 54′, dazu gehört in derselben Tafel 59\textsuperscript{a} die Länge}{}
\EdL{11}{AK07-J0130-L11}{12° 25′ der Zwillinge, dies ist der Anfang des V. Hauses\,; zu 160° 54′ wieder 24° 8′}{}
\EdL{12}{AK07-J0130-L12}{addiert, gibt 185° 2′, dazu gehört in derselben Tafel die Länge 4° 36′ des Krebses,}{}
\EdL{13}{AK07-J0130-L13}{dies ist der Anfang des VI. Hauses\,; zu 185° 2′ wieder 24° 8′ addiert, gibt 209° 10′,}{}
\EdL{14}{AK07-J0130-L14}{dazu gehört in derselben Tafel die Länge 27° 3′ des Krebses, dies ist der Anfang}{}
\EdL{15}{AK07-J0130-L15}{des VII. Hauses und dies ist der Nadir zum Aszendenten, = 27° 3′ des Steinbocks.}{}
\EdLp{16}{AK07-J0130-L16}{Wir haben also für die Anfänge der zwölf Häuser folgende Werte\,:}{}
\EdL{17}{AK07-J0130-L17}{\parbox[t]{\dimexpr0.1633\linewidth-0.5em\relax}{\raggedright\small\mbox{}I. u. VII.}\hspace{0.5em}\parbox[t]{\dimexpr0.1633\linewidth-0.5em\relax}{\raggedright\small\mbox{}II. u. VIII.}\hspace{0.5em}\parbox[t]{\dimexpr0.1633\linewidth-0.5em\relax}{\raggedright\small\mbox{}III. u. IX.}\hspace{0.5em}\parbox[t]{\dimexpr0.1633\linewidth-0.5em\relax}{\raggedright\small\mbox{}IV. u. X.}\hspace{0.5em}\parbox[t]{\dimexpr0.1633\linewidth-0.5em\relax}{\raggedright\small\mbox{}V. u. XI.}\hspace{0.5em}\parbox[t]{\dimexpr0.1633\linewidth-0.5em\relax}{\raggedright\small\mbox{}VI. u. XII.}\hspace{0.5em}}{}
\EdL{18}{AK07-J0130-L18}{\parbox[t]{\dimexpr0.1633\linewidth-0.5em\relax}{\raggedright\small\mbox{}27° 3′}\hspace{0.5em}\parbox[t]{\dimexpr0.1633\linewidth-0.5em\relax}{\raggedright\small\mbox{}3° 1′}\hspace{0.5em}\parbox[t]{\dimexpr0.1633\linewidth-0.5em\relax}{\raggedright\small\mbox{}11° 54′}\hspace{0.5em}\parbox[t]{\dimexpr0.1633\linewidth-0.5em\relax}{\raggedright\small\mbox{}19° 19′}\hspace{0.5em}\parbox[t]{\dimexpr0.1633\linewidth-0.5em\relax}{\raggedright\small\mbox{}12° 25′}\hspace{0.5em}\parbox[t]{\dimexpr0.1633\linewidth-0.5em\relax}{\raggedright\small\mbox{}4° 36′}\hspace{0.5em}}{}
\EdL{19}{AK07-J0130-L19}{\parbox[t]{\dimexpr0.1633\linewidth-0.5em\relax}{\raggedright\small\mbox{}des Steinbocks}\hspace{0.5em}\parbox[t]{\dimexpr0.1633\linewidth-0.5em\relax}{\raggedright\small\mbox{}der Fische}\hspace{0.5em}\parbox[t]{\dimexpr0.1633\linewidth-0.5em\relax}{\raggedright\small\mbox{}des Widders}\hspace{0.5em}\parbox[t]{\dimexpr0.1633\linewidth-0.5em\relax}{\raggedright\small\mbox{}des Stiers}\hspace{0.5em}\parbox[t]{\dimexpr0.1633\linewidth-0.5em\relax}{\raggedright\small\mbox{}der Zwillinge}\hspace{0.5em}\parbox[t]{\dimexpr0.1633\linewidth-0.5em\relax}{\raggedright\small\mbox{}des Krebses}\hspace{0.5em}}{}
\EdL{20}{AK07-J0130-L20}{\parbox[t]{\dimexpr0.1633\linewidth-0.5em\relax}{\raggedright\small\mbox{}bezw.}\hspace{0.5em}\parbox[t]{\dimexpr0.1633\linewidth-0.5em\relax}{\raggedright\small\mbox{}bezw.}\hspace{0.5em}\parbox[t]{\dimexpr0.1633\linewidth-0.5em\relax}{\raggedright\small\mbox{}bezw.}\hspace{0.5em}\parbox[t]{\dimexpr0.1633\linewidth-0.5em\relax}{\raggedright\small\mbox{}bezw.}\hspace{0.5em}\parbox[t]{\dimexpr0.1633\linewidth-0.5em\relax}{\raggedright\small\mbox{}bezw.}\hspace{0.5em}\parbox[t]{\dimexpr0.1633\linewidth-0.5em\relax}{\raggedright\small\mbox{}bezw.}\hspace{0.5em}}{}
\EdL{21}{AK07-J0130-L21}{\parbox[t]{\dimexpr0.1633\linewidth-0.5em\relax}{\raggedright\small\mbox{}des Krebses}\hspace{0.5em}\parbox[t]{\dimexpr0.1633\linewidth-0.5em\relax}{\raggedright\small\mbox{}der Jungfrau}\hspace{0.5em}\parbox[t]{\dimexpr0.1633\linewidth-0.5em\relax}{\raggedright\small\mbox{}der Wage}\hspace{0.5em}\parbox[t]{\dimexpr0.1633\linewidth-0.5em\relax}{\raggedright\small\mbox{}des Skorpions}\hspace{0.5em}\parbox[t]{\dimexpr0.1633\linewidth-0.5em\relax}{\raggedright\small\mbox{}des Schützen}\hspace{0.5em}\parbox[t]{\dimexpr0.1633\linewidth-0.5em\relax}{\raggedright\small\mbox{}des Steinbocks}\hspace{0.5em}}{}
\EdLp{22}{AK07-J0130-L22}{Wir finden nun in unserer Tafel 88 für den Aszendenten 27° 3′ des Steinbocks fol-}{}
\EdL{23}{AK07-J0130-L23}{gende Werte für die Anfänge der zwölf Häuser (mit Weglassung der Tierkreiszeichen):}{}
\EdL{24}{AK07-J0130-L24}{\parbox[t]{\dimexpr0.1633\linewidth-0.5em\relax}{\raggedright\small\mbox{}I. u. VII.}\hspace{0.5em}\parbox[t]{\dimexpr0.1633\linewidth-0.5em\relax}{\raggedright\small\mbox{}II. u. VIII.}\hspace{0.5em}\parbox[t]{\dimexpr0.1633\linewidth-0.5em\relax}{\raggedright\small\mbox{}III. u. IX.}\hspace{0.5em}\parbox[t]{\dimexpr0.1633\linewidth-0.5em\relax}{\raggedright\small\mbox{}IV. u. X.}\hspace{0.5em}\parbox[t]{\dimexpr0.1633\linewidth-0.5em\relax}{\raggedright\small\mbox{}V. u. XI.}\hspace{0.5em}\parbox[t]{\dimexpr0.1633\linewidth-0.5em\relax}{\raggedright\small\mbox{}VI. u. XII.}\hspace{0.5em}}{}
\EdL{25}{AK07-J0130-L25}{\parbox[t]{\dimexpr0.1633\linewidth-0.5em\relax}{\raggedright\small\mbox{}27° 3′}\hspace{0.5em}\parbox[t]{\dimexpr0.1633\linewidth-0.5em\relax}{\raggedright\small\mbox{}3° 8′}\hspace{0.5em}\parbox[t]{\dimexpr0.1633\linewidth-0.5em\relax}{\raggedright\small\mbox{}12° 2′}\hspace{0.5em}\parbox[t]{\dimexpr0.1633\linewidth-0.5em\relax}{\raggedright\small\mbox{}19° 31′}\hspace{0.5em}\parbox[t]{\dimexpr0.1633\linewidth-0.5em\relax}{\raggedright\small\mbox{}12° 33′}\hspace{0.5em}\parbox[t]{\dimexpr0.1633\linewidth-0.5em\relax}{\raggedright\small\mbox{}4° 36′}\hspace{0.5em}}{}
\EdLp{26}{AK07-J0130-L26}{Die Unterschiede bewegen sich zwischen 0′ und 12′, sie mögen in erster Linie}{}
\EdL{27}{AK07-J0130-L27}{von einer abweichenden Länge der Zeitstunde herrühren, es können aber auch}{}
\EdL{28}{AK07-J0130-L28}{Schreibfehler in den Zahlen die Ursache sein.}{}
\EdLp{29}{AK07-J0130-L29}{Für die Stellung eines Horoskops bei einer Geburt etc. waren diese Tafeln von}{}
\EdL{30}{AK07-J0130-L30}{grosser Wichtigkeit, und ersetzten umständliche Berechnungen und nicht immer}{}
\EdL{31}{AK07-J0130-L31}{mögliche Beobachtungen\,: kannte man nach den Tafeln 4—56 die mittlern und}{}
\EdL{32}{AK07-J0130-L32}{wahren Örter von Sonne, Mond und der fünf Planeten, so konnte man nach den}{}
\EdL{33}{AK07-J0130-L33}{Tafeln 79—90 sofort bestimmen, in welchem der zwölf Häuser sich jeder dieser}{}
\EdL{34}{AK07-J0130-L34}{Himmelskörper zu der betreffenden Zeit befand.}{}
\EdLc{35}{AK07-J0130-L35}{\textbf{Kap. 37. Über die Örter der Strahlenwerfung.}}{}
\EdLc{36}{AK07-J0130-L36}{(Text p. 30—31, Tafeln 91—114.)}{}
\EdLp{37}{AK07-J0130-L37}{Von ebenso grosser oder noch grösserer Wichtigkeit wie die Kenntnis der Ver-}{}
\EdL{38}{AK07-J0130-L38}{teilung der zwölf Häuser war in der Astrologie die Bestimmung der \textit{Proiectiones}}{}
\EdL{39}{AK07-J0130-L39}{\textit{radiorum,} d. h. derjenigen Punkte der Ekliptik, die von einem bestimmten Punkt am}{}
\EdL{40}{AK07-J0130-L40}{Himmel, z. B. von einem Planeten, um 60°, 90° oder 120° abstehen, oder also mit}{}
\EdL{41}{AK07-J0130-L41}{ihm im \textit{Sextil-}, \textit{Quadrat-} oder \textit{Trigonalschein} stehen. Zu diesen drei \textit{Aspekten} kommen}{}
\FrameHead{0131}{p. 99}{0131}{p. 99}
\EdL{1}{AK07-J0131-L01}{dann noch die \textit{Konjunktion} und \textit{Opposition} (\textit{Gegenschein}) hinzu. Es ist charakteristisch,}{}
\EdL{2}{AK07-J0131-L02}{dass \Sperr{Kh.} gerade diesem Kapitel den grössten Raum seines Tafelwerkes gewidmet hat\,:}{}
\EdL{3}{AK07-J0131-L03}{Text und Tafeln umfassen zusammen 26 Seiten\,; dies illustriert besser als irgend}{}
\EdL{4}{AK07-J0131-L04}{welche andere Betrachtungen oder Untersuchungen die ja schon bekannte Tatsache,}{}
\EdL{5}{AK07-J0131-L05}{dass die ganze wissenschaftliche Arbeit des Altertums\textsuperscript{1}) und Mittelalters auf dem}{}
\EdL{6}{AK07-J0131-L06}{Gebiete der Astronomie nicht Selbstzweck war, sondern nur die Hilfsmittel liefern}{}
\EdL{7}{AK07-J0131-L07}{musste für die Bedürfnisse der Astrologie. Übrigens ist ja bekannt, dass noch \textit{Regio-}}{}
\EdL{8}{AK07-J0131-L08}{\textit{montanus} ein eigenes Werk über einen verwandten astrologischen Gegenstand ver-}{}
\EdL{9}{AK07-J0131-L09}{fasst hat\textsuperscript{2}), in welchem er ebenfalls (Problema 31) über die Frage der Bestimmung}{}
\EdL{10}{AK07-J0131-L10}{des \textit{Ortes} der \textit{Strahlenwerfung} (\textit{radiatio}) handelt.}{}
\EdLp{11}{AK07-J0131-L11}{Auch für die Bestimmung dieser Örter gab es verschiedene Wege, für deren}{}
\EdL{12}{AK07-J0131-L12}{Details ich wieder auf \Sperr{Batt.}\textsuperscript{3}) verweisen muss. Die einfachsten waren diejenigen,}{}
\EdL{13}{AK07-J0131-L13}{die die Breite der Planeten vernachlässigten, sie also in der Ekliptik liegend annahmen\,;}{}
\EdL{14}{AK07-J0131-L14}{dann war es natürlich das nächstliegende, vom Planeten aus 60° (bezw. 90° od. 120°)}{}
\EdL{15}{AK07-J0131-L15}{in der Ekliptik vor oder rückwärts zu gehen, so hatte man die beiden Punkte (den}{}
\EdL{16}{AK07-J0131-L16}{linken und rechten) des Sextilscheins (bezw. Quadrat- oder Trigonalscheins)\,: dieser}{}
\EdL{17}{AK07-J0131-L17}{primitive Weg wurde z. B. von \Sperr{Al-Ḳabīṣī} (\Sperr{Alcabitius}) und \Sperr{Abraham} b. \Sperr{ʿEzrā}}{}
\EdL{18}{AK07-J0131-L18}{eingeschlagen\textsuperscript{4}). Andere, wie z. B. der griechische Astrolog \Sperr{Dorotheus}, sollen nach}{}
\EdL{19}{AK07-J0131-L19}{\Sperr{Ibn Hibintā} den Ort der Strahlenwerfung nicht auf der Ekliptik, sondern auf dem}{}
\EdL{20}{AK07-J0131-L20}{Aequator gemessen haben, d. h. sie betrachteten als Radius des Sextilscheins den}{}
\EdL{21}{AK07-J0131-L21}{Bogen des Aequators, der gleich der Differenz\,: Rectascension des Planeten minus}{}
\EdL{22}{AK07-J0131-L22}{Rectascention von [Länge des Planeten ± 60°] oder umgekehrt war\,; sie kamen auf}{}
\EdL{23}{AK07-J0131-L23}{diese Weise zum gleichen Punkt der Ekliptik wie die vorhin genannten. — \Sperr{Ptol.}}{}
\EdL{24}{AK07-J0131-L24}{schlug im \textit{Quadripartitum} einen etwas komplizierteren Weg ein\textsuperscript{5}), vernachlässigte}{}
\EdL{25}{AK07-J0131-L25}{aber ebenfalls die Breite der Planeten. — Der erste, der diese berücksichtigt hat,}{}
\EdL{26}{AK07-J0131-L26}{scheint \Sperr{Batt.} gewesen zu sein, der die Formel aufgestellt hat\,:}{}
\EdLc{27}{AK07-J0131-L27}{crd. \textit{x} = \Root{\textit{R} · \Frac{crd.\textsuperscript{2} 60° — crd.\textsuperscript{2} \textit{b}}{cos \textit{b}}} ,}{}
\EdL{28}{AK07-J0131-L28}{wo \textit{b} die Breite des Planeten und \textit{x} der Bogen der Ekliptik ist, der der Hypotenuse}{}
\EdL{29}{AK07-J0131-L29}{60° (Radius des Sextilscheins) entspricht\,; diese Formel tritt schon implicite in der}{}
\EdL{30}{AK07-J0131-L30}{Aufgabe auf, die Entfernung zweier Sterne aus ihren Breiten und Längen zu bestimmen,}{}
\EdL{31}{AK07-J0131-L31}{und kann, wie \Sperr{Nallino} gezeigt hat\textsuperscript{6}), leicht in die Formel}{}
\EdLc{32}{AK07-J0131-L32}{cos \textit{x} = \Frac{cos 60°}{cos \textit{b}}}{}
\EdL{33}{AK07-J0131-L33}{übergeführt werden, die aus dem rechtwinkligen sphärischen}{}
\EdL{34}{AK07-J0131-L34}{Dreiecke \textit{ABC} nach der Formel cos \textit{c} = cos \textit{a} · cos \textit{b}, die}{}
\EdL{35}{AK07-J0131-L35}{\Sperr{Batt.} noch nicht gekannt hat, sich sofort ergibt.}{}
\NoteRule
\EdN{N1}{AK07-J0131-N01}{\textsuperscript{1}) Wenigstens von \Sperr{Ptol.} an.}
\EdN{N2}{AK07-J0131-N02}{\textsuperscript{2}) \textit{Tabulae directionum profectionumque}, Augsburg 1490, Tübingen 1550.}
\EdN{N3}{AK07-J0131-N03}{\textsuperscript{3}) Ed. \Sperr{Nallino}, I, 307—313.}
\EdN{N4}{AK07-J0131-N04}{\textsuperscript{4}) Ibid. I, 309.}
\EdN{N5}{AK07-J0131-N05}{\textsuperscript{5}) Ibid. I, 311.}
\EdN{N6}{AK07-J0131-N06}{\textsuperscript{6}) Ibid. I, 307—308.}
\FrameHead{0132}{p. 100}{0132}{p. 100}
\EdLp{1}{AK07-J0132-L01}{Welchen Weg hat nun \Sperr{Kh.} eingeschlagen\,? Nach dem, was \Sperr{Nallino}\textsuperscript{1}) aus}{}
\EdL{2}{AK07-J0132-L02}{\Sperr{Ibn Hibintā} zitiert, müsste man schliessen, \Sperr{Kh.} sei ungefähr so verfahren, wie}{}
\EdL{3}{AK07-J0132-L03}{oben von \Sperr{Dorotheus} berichtet wurde, nur habe er die Rectascension des Aequator-}{}
\EdL{4}{AK07-J0132-L04}{bogens wieder in Ekliptikgrade verwandelt, und dann diesen Ekliptikbogen als Radius}{}
\EdL{5}{AK07-J0132-L05}{des Sextilscheins genommen. Die uns hier vorliegenden Tafeln weichen aber von}{}
\EdL{6}{AK07-J0132-L06}{den Tafeln \Sperr{Kh.}’s, wie sie \Sperr{Ibn Hibintā} uns überliefert hat, und von denen \Sperr{Nallino}}{}
\EdL{7}{AK07-J0132-L07}{mir gütigst eine Probe übersandt hat, sowohl in der Anordnung als im Inhalt wesent-}{}
\EdL{8}{AK07-J0132-L08}{lich ab\,; erstens gehen sie nur von 10° zu 10° der Zeichen, die unsrigen aber von}{}
\EdL{9}{AK07-J0132-L09}{5° zu 5°\,; zweitens geben sie jeweilen nicht die ganze Rectascension vom Frühlings-}{}
\EdL{10}{AK07-J0132-L10}{punkt an gerechnet, sondern nur die Rectascension jedes Zeichens für sich\textsuperscript{2})\,; auch}{}
\EdL{11}{AK07-J0132-L11}{die Zahlen für den Sextil-, Quadrat- und Trigonalschein sind andere. Unsere Tafeln}{}
\EdL{12}{AK07-J0132-L12}{nehmen die \textit{Ascensiones obliquae} für die Breite von \textit{Cordova,} die \Sperr{Maslama} zu 38° 30′}{}
\EdL{13}{AK07-J0132-L13}{angibt, zur Grundlage\,; die Tafeln \Sperr{Kh.}’s für die \textit{proiectiones radiorum} sind also eben-}{}
\EdL{14}{AK07-J0132-L14}{falls für Cordova umgearbeitet worden. Aber nun erwächst für uns ein schwieriges}{}
\EdL{15}{AK07-J0132-L15}{Problem\,; die Ascensiones obliquae stimmen jeweilen nur mit den Zahlen der ersten}{}
\EdL{16}{AK07-J0132-L16}{drei Kolonnen, nachher weichen sie ziemlich stark ab\,; geben wir hiefür ein Beispiel\,:}{}
\EdLp{17}{AK07-J0132-L17}{Tafel 91, erste Zeile, hat folgende Zahlen\,:}{}
\EdL{18}{AK07-J0132-L18}{\parbox[t]{\dimexpr0.1400\linewidth-0.5em\relax}{\raggedright\small\mbox{}Gradus}\hspace{0.5em}\parbox[t]{\dimexpr0.1400\linewidth-0.5em\relax}{\raggedright\small\mbox{}Aries}\hspace{0.5em}\parbox[t]{\dimexpr0.1400\linewidth-0.5em\relax}{\raggedright\small\mbox{}Taurus}\hspace{0.5em}\parbox[t]{\dimexpr0.1400\linewidth-0.5em\relax}{\raggedright\small\mbox{}Gemini}\hspace{0.5em}\parbox[t]{\dimexpr0.1400\linewidth-0.5em\relax}{\raggedright\small\mbox{}Cancer}\hspace{0.5em}\parbox[t]{\dimexpr0.1400\linewidth-0.5em\relax}{\raggedright\small\mbox{}Leo}\hspace{0.5em}\parbox[t]{\dimexpr0.1400\linewidth-0.5em\relax}{\raggedright\small\mbox{}Virgo}\hspace{0.5em}}{}
\EdL{19}{AK07-J0132-L19}{\parbox[t]{\dimexpr0.1400\linewidth-0.5em\relax}{\raggedright\small\mbox{}aequales}\hspace{0.5em}\parbox[t]{\dimexpr0.1400\linewidth-0.5em\relax}{\raggedright\small\mbox{}(Libra)}\hspace{0.5em}\parbox[t]{\dimexpr0.1400\linewidth-0.5em\relax}{\raggedright\small\mbox{}(Scorpius)}\hspace{0.5em}\parbox[t]{\dimexpr0.1400\linewidth-0.5em\relax}{\raggedright\small\mbox{}(Sagittar.)}\hspace{0.5em}\parbox[t]{\dimexpr0.1400\linewidth-0.5em\relax}{\raggedright\small\mbox{}(Capricorn.)}\hspace{0.5em}\parbox[t]{\dimexpr0.1400\linewidth-0.5em\relax}{\raggedright\small\mbox{}(Aquar.)}\hspace{0.5em}\parbox[t]{\dimexpr0.1400\linewidth-0.5em\relax}{\raggedright\small\mbox{}(Pisces)}\hspace{0.5em}}{}
\EdL{20}{AK07-J0132-L20}{\parbox[t]{\dimexpr0.1400\linewidth-0.5em\relax}{\raggedright\small\mbox{}5°}\hspace{0.5em}\parbox[t]{\dimexpr0.1400\linewidth-0.5em\relax}{\raggedright\small\mbox{}2° 58′}\hspace{0.5em}\parbox[t]{\dimexpr0.1400\linewidth-0.5em\relax}{\raggedright\small\mbox{}21° 50′}\hspace{0.5em}\parbox[t]{\dimexpr0.1400\linewidth-0.5em\relax}{\raggedright\small\mbox{}45° 20′}\hspace{0.5em}\parbox[t]{\dimexpr0.1400\linewidth-0.5em\relax}{\raggedright\small\mbox{}76° 0′}\hspace{0.5em}\parbox[t]{\dimexpr0.1400\linewidth-0.5em\relax}{\raggedright\small\mbox{}106° 32′}\hspace{0.5em}\parbox[t]{\dimexpr0.1400\linewidth-0.5em\relax}{\raggedright\small\mbox{}131° 9′}\hspace{0.5em}}{}
\EdLp{21}{AK07-J0132-L21}{Die Ascensiones obliquae der ersten 5 Grade jedes Zeichens sind aber für die}{}
\EdL{22}{AK07-J0132-L22}{Breite 38° 30′\,:}{}
\EdL{23}{AK07-J0132-L23}{\parbox[t]{\dimexpr0.1400\linewidth-0.5em\relax}{\raggedright\small\mbox{}}\hspace{0.5em}\parbox[t]{\dimexpr0.1400\linewidth-0.5em\relax}{\raggedright\small\mbox{}3° 0′}\hspace{0.5em}\parbox[t]{\dimexpr0.1400\linewidth-0.5em\relax}{\raggedright\small\mbox{}21° 55′}\hspace{0.5em}\parbox[t]{\dimexpr0.1400\linewidth-0.5em\relax}{\raggedright\small\mbox{}45° 5′}\hspace{0.5em}\parbox[t]{\dimexpr0.1400\linewidth-0.5em\relax}{\raggedright\small\mbox{}75° 18′}\hspace{0.5em}\parbox[t]{\dimexpr0.1400\linewidth-0.5em\relax}{\raggedright\small\mbox{}111° 24′}\hspace{0.5em}\parbox[t]{\dimexpr0.1400\linewidth-0.5em\relax}{\raggedright\small\mbox{}149° 2′}\hspace{0.5em}}{}
\EdLp{24}{AK07-J0132-L24}{Diese Abweichungen in den letzten drei bis vier Kolonnen zeigen sich überall,}{}
\EdL{25}{AK07-J0132-L25}{woher rühren sie\,? Es ist mir bis jetzt nicht gelungen, dies zu erklären, die Art und}{}
\EdL{26}{AK07-J0132-L26}{Weise, wie die Tafeln berechnet worden sind, ist mir also noch dunkel, doch habe}{}
\EdL{27}{AK07-J0132-L27}{ich aus den, wenn auch sehr kurzen und oft durch die verdorbene lateinische Über-}{}
\EdL{28}{AK07-J0132-L28}{setzung entstellten Regeln des Textes herausgefunden, wie man sie zur Bestimmung}{}
\EdL{29}{AK07-J0132-L29}{des Ortes der Strahlenwerfung benutzen muss. Ich gebe hiefür drei \textit{Beispiele}\,:}{}
\EdLp{30}{AK07-J0132-L30}{1. Zur Berechnung des Ortes braucht man in erster Linie das \textit{Horoskop}\textsuperscript{3}),}{}
\EdL{31}{AK07-J0132-L31}{d. h. den Grad der Ekliptik, der zur bestimmten Zeit (der Berechnung) aufgeht, er}{}
\EdL{32}{AK07-J0132-L32}{sei 55° (= 25° des Stiers), in zweiter Linie die \textit{Länge} des Planeten, sie sei 65° (5° der}{}
\EdL{33}{AK07-J0132-L33}{Zwillinge). Man sucht nun in den Tafeln zuerst das Feld (\textit{area}) auf, das dem Horo-}{}
\EdL{34}{AK07-J0132-L34}{skop 25° des Stiers entspricht, es ist das Feld auf Tafel 94, das mit \textit{Viginti quinque}}{}
\EdL{35}{AK07-J0132-L35}{\textit{gradus tauri} überzchrieben ist, dann tritt man mit der Länge des Planeten, 5° der}{}
\EdL{36}{AK07-J0132-L36}{Zwillinge, in dieses Feld ein, d. h. man sucht 5° in der ersten Kolonne, und geht}{}
\EdL{37}{AK07-J0132-L37}{horizontal weiter bis zur Kolonne, die oben mit \SmallStack{Gemini}{Sagittar.} überschrieben ist, hier findet}{}
\NoteRule
\EdN{N1}{AK07-J0132-N01}{\textsuperscript{1}) Ibid. I, 312.}
\EdN{N2}{AK07-J0132-N02}{\textsuperscript{2}) d. h. z. B. die Rectascension des Widders, also von 30°, ist 27° 54′, diejenige von Widder}
\EdN{N3}{AK07-J0132-N03}{und Stier zusammen, also von 60°, ist 57° 48′, daher ist diejenige des Stiers allein gleich der Diffe}
\EdN{N4}{AK07-J0132-N04}{renz dieser Zahlen, d. i. = 29° 54′.}
\EdN{N5}{AK07-J0132-N05}{\textsuperscript{3}) \textit{Horoscopus} = \textit{Ascendens.}}
\FrameHead{0133}{p. 101}{0133}{p. 101}
\EdL{1}{AK07-J0133-L01}{man die Zahl 26° 45′, zu dieser muss man die Zahl, die am Kopf dieses Feldes bei}{}
\EdL{2}{AK07-J0133-L02}{\textit{Exagonum} (Sextilschein) steht, also 53° 30′ addieren\textsuperscript{1}), dies gibt 80° 15′\,; diese Zahl}{}
\EdL{3}{AK07-J0133-L03}{sucht man in demselben Feld unter den Zahlen der sechs Zeilen und sechs Kolonnen,}{}
\EdL{4}{AK07-J0133-L04}{und findet die nächst höhere 80° 23′ in der 5. Zeile, 4. Kolonne, dazu gehört (sie wird}{}
\EdL{5}{AK07-J0133-L05}{als eine \textit{ascensio obliqua} aufgefasst) die Länge 25° des Krebses, die Interpolation ergibt}{}
\EdL{6}{AK07-J0133-L06}{für die Differenz 80° 23′ — 80° 15′ = 8′ die Zahl 7′, also ist der Ort (= Länge) des}{}
\EdL{7}{AK07-J0133-L07}{\textit{linken Sextilscheins} = 24° 53′ des Krebses\textsuperscript{2}). Will man den Ort des rechten Sextil-}{}
\EdL{8}{AK07-J0133-L08}{scheins haben, so muss man den Ort des linken Trigonalscheins suchen, und zu ihm}{}
\EdL{9}{AK07-J0133-L09}{180° addieren.}{}
\EdLp{10}{AK07-J0133-L10}{2. \textit{Horoskop} = 317° (= 17° des Wassermanns), \textit{Länge} des Planeten = 138°}{}
\EdL{11}{AK07-J0133-L11}{(= 18° des Löwen). Im Felde von Tafel 111, überschrieben \textit{Quindecim gradus Aquarii,}}{}
\EdL{12}{AK07-J0133-L12}{findet man zu 18° des Löwen mit Zuhilfenahme der Interpolation die Zahl 14° 16′,}{}
\EdL{13}{AK07-J0133-L13}{dazu die Zahl, die am Kopfe des Feldes bei \textit{Exagonum} steht, 49° 28′ addiert, gibt}{}
\EdL{14}{AK07-J0133-L14}{63° 44′, zu dieser Zahl findet man in demselben Feld mit Hilfe der Interpolation die}{}
\EdL{15}{AK07-J0133-L15}{Länge 3° 40′ des Skorpions. Da nun aber das Horoskop nicht 15°, sondern 17° des}{}
\EdL{16}{AK07-J0133-L16}{Wassermanns, so muss man noch den Sextilschein für 20° des Wassermanns bestimmen,}{}
\EdL{17}{AK07-J0133-L17}{diesen findet man nach dem gleichen Verfahren im nächsten Felde (Tafel 112) = 3° 55′}{}
\EdL{18}{AK07-J0133-L18}{des Skorpions\,; die Differenz der beiden Werte ist also für 5° des Horoskops = 15′,}{}
\EdL{19}{AK07-J0133-L19}{für 2° (von 15° auf 17°) also = 6′, dies zu 3° 40′ addiert, gibt für den gesuchten Ort}{}
\EdL{20}{AK07-J0133-L20}{des \textit{linken Sextilscheins} 3° 46′ des Skorpions.}{}
\EdLp{21}{AK07-J0133-L21}{3. \textit{Horoskop} = 335° (= 5° der Fische), \textit{Länge} des Planeten = 115° (= 25° des}{}
\EdL{22}{AK07-J0133-L22}{Krebses). Im Felde von Tafel 113, überschrieben \textit{Quinque gradus Piscium} findet man}{}
\EdL{23}{AK07-J0133-L23}{zu 25° des Krebses die Zahl 117° 13′, dazu die Zahl, die bei \textit{Exagonum} steht, 48° 19′}{}
\EdL{24}{AK07-J0133-L24}{addiert, gibt 165° 32′\,; diese Zahl findet man in diesem Felde nicht mehr, die höchste}{}
\EdL{25}{AK07-J0133-L25}{desselben ist 144° 58′, diese subtrahiert man von der erhaltenen, bleibt 20° 34′, dazu}{}
\EdL{26}{AK07-J0133-L26}{findet man in der 1.—2. Kolonne desselben Feldes mit Hilfe der Interpolation die}{}
\EdL{27}{AK07-J0133-L27}{Länge 3° 45′ der Wage als Ort des \textit{linken Sextilscheins.}}{}
\EdLp{28}{AK07-J0133-L28}{Hat man die \textit{linken} Sextil-, Quadrat- und Trigonalscheine auf diese Weise}{}
\EdL{29}{AK07-J0133-L29}{bestimmt, so findet man die \textit{rechten} nach den Regeln\,:}{}
\EdL{30}{AK07-J0133-L30}{\parbox[t]{\dimexpr0.1960\linewidth-0.5em\relax}{\raggedright\small\mbox{}Der rechte}\hspace{0.5em}\parbox[t]{\dimexpr0.1960\linewidth-0.5em\relax}{\raggedright\small\mbox{}Trigonalschein}\hspace{0.5em}\parbox[t]{\dimexpr0.3920\linewidth-0.5em\relax}{\raggedright\small\mbox{}ist der Nadir des linken}\hspace{0.5em}\parbox[t]{\dimexpr0.1960\linewidth-0.5em\relax}{\raggedright\small\mbox{}Sextilscheins.}\hspace{0.5em}}{}
\EdL{31}{AK07-J0133-L31}{\parbox[t]{\dimexpr0.1960\linewidth-0.5em\relax}{\raggedright\small\mbox{}„ „}\hspace{0.5em}\parbox[t]{\dimexpr0.1960\linewidth-0.5em\relax}{\raggedright\small\mbox{}Quadratschein}\hspace{0.5em}\parbox[t]{\dimexpr0.3920\linewidth-0.5em\relax}{\raggedright\small\mbox{}„ „ „ „ „}\hspace{0.5em}\parbox[t]{\dimexpr0.1960\linewidth-0.5em\relax}{\raggedright\small\mbox{}Quadratscheins.}\hspace{0.5em}}{}
\EdL{32}{AK07-J0133-L32}{\parbox[t]{\dimexpr0.1960\linewidth-0.5em\relax}{\raggedright\small\mbox{}„ „}\hspace{0.5em}\parbox[t]{\dimexpr0.1960\linewidth-0.5em\relax}{\raggedright\small\mbox{}Sextilschein}\hspace{0.5em}\parbox[t]{\dimexpr0.3920\linewidth-0.5em\relax}{\raggedright\small\mbox{}„ „ „ „ „}\hspace{0.5em}\parbox[t]{\dimexpr0.1960\linewidth-0.5em\relax}{\raggedright\small\mbox{}Trigonalscheins.}\hspace{0.5em}}{}
\EdLp{33}{AK07-J0133-L33}{So ist also im 1. Beispiel der rechte Trigonalschein = 24° 53′ des Steinbocks.}{}
\EdL{34}{AK07-J0133-L34}{In den gegebenen drei Beispielen sind alle Fälle zur Darstellung gebracht, die bei}{}
\EdL{35}{AK07-J0133-L35}{diesen Bestimmungen eventuell eintreten können und die in den Regeln des Textes}{}
\EdL{36}{AK07-J0133-L36}{berücksichtigt sind.}{}
\EdLp{37}{AK07-J0133-L37}{Es gehört wohl die Methode \Sperr{Kh.}’s zu der von \textit{Regiomontan} (l. c.) genannten}{}
\NoteRule
\EdN{N1}{AK07-J0133-N01}{\textsuperscript{1}) Dies natürlich dann, wenn man den Sextilschein sucht\,; sucht man den Quadrat- oder}
\EdN{N2}{AK07-J0133-N02}{Trigonalschein, so muss man die Zahl addieren, die am Kopf des Feldes bei \textit{Tetragonum,} bezw.}
\EdN{N3}{AK07-J0133-N03}{\textit{Trigonum} steht.}
\EdN{N4}{AK07-J0133-N04}{\textsuperscript{2}) Dies ist also c. 10° weniger als sich ergeben würde, wenn man einfach vom Ort des Pla-}
\EdN{N5}{AK07-J0133-N05}{neten auf der Ekliptik 60° vorwärts zählen würde, man käme so zu 5° des Löwen.}
\FrameHead{0134}{p. 102}{0134}{p. 102}
\EdL{1}{AK07-J0134-L01}{Kategorie, von der er sagt\,: \textit{Alii autem exequuntur id negotium per ascensiones quidem}}{}
\EdL{2}{AK07-J0134-L02}{\textit{rectas stella meridianum tenente, per obliquas autem ascensiones regionis, si in oriente}}{}
\EdL{3}{AK07-J0134-L03}{\textit{extiterit, aut per descensiones si in occidente}\,; nur unterscheidet \Sperr{Kh.} den Fall, wo der}{}
\EdL{4}{AK07-J0134-L04}{Stern im Meridian ist, nicht besonders.}{}
\EdLp{5}{AK07-J0134-L05}{Man muss zugestehen, dass diese Tafeln \Sperr{Kh.}’s für die Astrologen zur Bestim-}{}
\EdL{6}{AK07-J0134-L06}{mung der verschiedenen Aspekten sehr bequem waren, nur mögen sie in Bezug}{}
\EdL{7}{AK07-J0134-L07}{auf die Genauigkeit etwas zu wünschen übrig gelassen haben.}{}
\EdLp{8}{AK07-J0134-L08}{Damit ist der Kommentar zum Texte und zu den Tafeln, soweit diese in den}{}
\EdL{9}{AK07-J0134-L09}{Handschriften C, O und N übereinstimmen, beendigt. Das Ms. O enthält aber noch}{}
\EdL{10}{AK07-J0134-L10}{zwei Tafeln am Schlusse (fol. 158\textsuperscript{r} u. 158\textsuperscript{v}), die \Sperr{Björnbo} gar nicht in seine Abschrift}{}
\EdL{11}{AK07-J0134-L11}{aufgenommen hatte, wahrscheinlich übersah er den Satz\,: \textit{Explicit ezich,} etc., der am}{}
\EdL{12}{AK07-J0134-L12}{Schlusse von Tafel 116 steht, und hielt deshalb diese beiden Tafeln für von einem}{}
\EdL{13}{AK07-J0134-L13}{Abschreiber beigefügte Anhängsel, die nicht zum vorliegenden Werke gehörten. Es}{}
\EdL{14}{AK07-J0134-L14}{ist zuzugeben, dass dies wohl möglich wäre, ich halte es aber für wahrscheinlicher,}{}
\EdL{15}{AK07-J0134-L15}{dass, da diese Tafeln doch einigermassen im Zusammenhang mit den vorhergehenden}{}
\EdL{16}{AK07-J0134-L16}{astrologischen Tafeln stehen, sie den Schluss der Bearbeitung der Tafeln \Sperr{Kh.}’s durch}{}
\EdL{17}{AK07-J0134-L17}{\Sperr{Maslama} bildeten, zumal sich die zweite Tafel (als astrologische Rose) auch unter}{}
\EdL{18}{AK07-J0134-L18}{den \textit{Tabulae spuriae} \Sperr{Batt.}’s befindet, die alle von \Sperr{Maslama} herstammen. Ich gebe}{}
\EdL{19}{AK07-J0134-L19}{daher in den folgenden zwei (nicht numerierten) Schlusskapiteln eine kurze Be-}{}
\EdL{20}{AK07-J0134-L20}{schreibung und Erklärung dieser beiden Tafeln.}{}
\EdLc{21}{AK07-J0134-L21}{\textbf{Über die Umwandlung der Geburtsjahre.}}{}
\EdLc{22}{AK07-J0134-L22}{(Tafel 115.)}{}
\EdLp{23}{AK07-J0134-L23}{Die Tafel 115 lehrt den Ort am Himmel bestimmen, den die Sonne nach Ver-}{}
\EdL{24}{AK07-J0134-L24}{lauf einer Anzahl voller Jahre, z. B. am 35. Geburtstage eines Menschen, am Himmel}{}
\EdL{25}{AK07-J0134-L25}{einnimmt, oder mit andern Worten, wie viel Grade (oder wie viel Zeit) sie noch zu}{}
\EdL{26}{AK07-J0134-L26}{durchlaufen hat, bis sie wieder bei demselben Punkte (Fixstern) angekommen ist, bei}{}
\EdL{27}{AK07-J0134-L27}{dem sie bei der Geburt stand. Vom rein astronomischen Standpunkt aus aufgefasst,}{}
\EdL{28}{AK07-J0134-L28}{kann man diese Tafel auch betrachten als Tafel der Umwandlung der bürgerlichen}{}
\EdL{29}{AK07-J0134-L29}{(julianischen) Jahre in siderische Jahre. \Sperr{Batt.} hat (II, 188) eine ähnliche Tafel für}{}
\EdL{30}{AK07-J0134-L30}{Umwandlung der bürgerlichen Jahre in tropische, die \textit{anni collecti} schreiten aber}{}
\EdL{31}{AK07-J0134-L31}{von 20 zu 20 fort, statt wie bei unserer Tafel von 10 zu 10, und die Überschüsse}{}
\EdL{32}{AK07-J0134-L32}{über die vollen bürgerlichen Jahre sind nur in Kreisteilen angegeben, während sie}{}
\EdL{33}{AK07-J0134-L33}{in unsrer Tafel auch in Zeitmass angegeben sind.}{}
\EdLp{34}{AK07-J0134-L34}{\textit{Beispiel\,:} Es ist Jemand geboren am 4. Januar 962 zu einer gewissen Stunde,}{}
\EdL{35}{AK07-J0134-L35}{wie viel Zeit muss er noch warten von derselben Stunde des 4. Jan. 997 an, bis die}{}
\EdL{36}{AK07-J0134-L36}{Sonne an demselben Ort des Himmels, d. h. bei demselben Fixstern angekommen}{}
\EdL{37}{AK07-J0134-L37}{ist, bei dem sie im Zeitpunkt der Geburt stand\,?}{}
\EdLp{38}{AK07-J0134-L38}{Die Tafel 115 ergibt für 30 Jahre (in der Kolonne \textit{anni deceni}) 7\textsuperscript{d} 18\textsuperscript{h} 4\textsuperscript{m}, für 5}{}
\EdL{39}{AK07-J0134-L39}{Jahre (in der Kolonne \textit{anni singulares}) 1\textsuperscript{d} 7\textsuperscript{h} 1\textsuperscript{m}, zusammen 9\textsuperscript{d} 1\textsuperscript{h} 5\textsuperscript{m}. Also erst nach}{}
\NoteRule
\EdN{N1}{AK07-J0134-N01}{\textsuperscript{1}) Ed. \Sperr{Nallino}, II, 299.}
\FrameHead{0135}{p. 103}{0135}{p. 103}
\EdL{1}{AK07-J0135-L01}{9\textsuperscript{d} 1\textsuperscript{h} 5\textsuperscript{m} von seinem 35. Geburtstage an steht die Sonne bei demselben Fixstern, wie}{}
\EdL{2}{AK07-J0135-L02}{bei seiner Geburt. Diese Zeit erfährt aber noch eine Korrektur durch die Bestimmung}{}
\EdL{3}{AK07-J0135-L03}{des \textit{wahren} Ortes der Sonne, denn der gefundene Ort ist nur der \textit{mittlere}\,; ich werde}{}
\EdL{4}{AK07-J0135-L04}{aber hierauf nicht näher eintreten, sondern verweise dafür auf \Sperr{Batt.} I, 148. Die}{}
\EdL{5}{AK07-J0135-L05}{zweite Hälfte der Tafel gibt die von der Sonne noch zu durchlaufende Bahn in Graden,}{}
\EdL{6}{AK07-J0135-L06}{Minuten und Sekunden an. — Diese \textit{Revolutiones} oder \textit{Conversiones annorum nativi-}}{}
\EdL{7}{AK07-J0135-L07}{\textit{tatum,} und ebenso die \textit{Revolutiones annorum mundi} (oder \textit{civitatum}) spielten in der}{}
\EdL{8}{AK07-J0135-L08}{Astrologie des Altertums und Mittelalters eine grosse Rolle\textsuperscript{1}).}{}
\EdLp{9}{AK07-J0135-L09}{Es ist hier noch zu erwähnen, dass in dieser Tafel der Überschuss des siderischen}{}
\EdL{10}{AK07-J0135-L10}{Jahres über das bürgerliche zu 6\textsuperscript{h} 12\textsuperscript{m}, oder wenn wir die 93° 2′ 15″ der letzten Kolonne}{}
\EdL{11}{AK07-J0135-L11}{in Zeit verwandeln, genauer zu 6\textsuperscript{h} 12\textsuperscript{m} 9\textsuperscript{s} angenommen ist, während nach Kap. 20 u. 21}{}
\EdL{12}{AK07-J0135-L12}{6\textsuperscript{h} 13\textsuperscript{m} 2\textsuperscript{s} zu erwarten gewesen wären. Der Wert von 6\textsuperscript{h} 12\textsuperscript{m} 9\textsuperscript{s} entspricht nach den}{}
\EdL{13}{AK07-J0135-L13}{Angaben verschiedener arabischer Astronomen der Länge des siderischen Jahres im}{}
\EdL{14}{AK07-J0135-L14}{arabischen \textit{Sindhind,} es scheint also auch \Sperr{Maslama al-Madjrīṭī} diesen Wert}{}
\EdL{15}{AK07-J0135-L15}{angenommen zu haben, und hierin also nicht \Sperr{Kh.} gefolgt zu sein. \textsuperscript{2})}{}
\EdLc{16}{AK07-J0135-L16}{\textbf{Über die Domicilia, Exaltationes, Decani, etc. der Planeten.}}{}
\EdLc{17}{AK07-J0135-L17}{(Tafel 116.)}{}
\EdLp{18}{AK07-J0135-L18}{Diese Tafel ist im ganzen identisch mit der \textit{astrologischen Rose} bei \Sperr{Batt.}, wie}{}
\EdL{19}{AK07-J0135-L19}{ich oben schon erwähnt habe\,; in einzelnen Punkten aber, die noch zu nennen sind,}{}
\EdL{20}{AK07-J0135-L20}{weichen die beiden von einander ab.}{}
\EdLp{21}{AK07-J0135-L21}{Die 1. Kolonne\textsuperscript{3}) enthält die zwölf Zeichen des Tierkreises, die in der genannten}{}
\EdL{22}{AK07-J0135-L22}{Rose ausserhalb des äussersten Kreises stehen, der dem entsprechend in zwölf Sek-}{}
\EdL{23}{AK07-J0135-L23}{toren geteilt ist. In der 2. Kolonne stehen die Herren (\textit{Domini}) der entsprechenden}{}
\EdL{24}{AK07-J0135-L24}{Häuser (\textit{Domicilia}) der 1. Kolonne, sie wäre also eher mit \textit{Domini} statt \textit{Domicilia} zu}{}
\EdL{25}{AK07-J0135-L25}{überschreiben\,; es ist also der Widder (\textit{Aries}) das Haus des \textit{Mars,} oder \textit{Mars} ist der}{}
\EdL{26}{AK07-J0135-L26}{Hausherr des \textit{Widders, Venus} der des \textit{Stiers} (\textit{Taurus}) u. s. w. Sonne und Mond haben}{}
\EdL{27}{AK07-J0135-L27}{je nur \textit{ein} Haus, die fünf Planeten aber jeder deren \textit{zwei,} so ist ausser dem \textit{Widder}}{}
\EdL{28}{AK07-J0135-L28}{auch noch der \textit{Skorpion} das Haus des \textit{Mars,} ausser dem \textit{Stier} auch noch die \textit{Wage}}{}
\EdL{29}{AK07-J0135-L29}{(\textit{Libra}) das Haus der \textit{Venus.} — Diese \textit{Domini} (arab. \textit{ṣāḥib} oder \textit{rabb}) stehen in der}{}
\EdL{30}{AK07-J0135-L30}{Rose \Sperr{Batt.}’s im zweitinnersten Kreise.}{}
\EdLp{31}{AK07-J0135-L31}{In der 4. Kolonne, überschrieben \textit{Alienatio,} stehen die Planeten, die in dem}{}
\EdL{32}{AK07-J0135-L32}{entsprechenden Tierkreiszeichen im sog. \textit{Schaden} (\textit{detrimentum,} arab. \textit{wabāl}) sind\,;}{}
\EdL{33}{AK07-J0135-L33}{ein Planet ist in dem Haus im Schaden, das um 180° von seinem Domicilium ent-}{}
\EdL{34}{AK07-J0135-L34}{fernt ist, oder das der Nadir seines Domiciliums ist, es ist dann also diesem gleich-}{}
\EdL{35}{AK07-J0135-L35}{sam entfremdet, daher der Name \textit{Alienatio,} den ich sonst bei keinem mittelalterlichen}{}
\EdL{36}{AK07-J0135-L36}{Astrologen bis jetzt gefunden habe. So ist \textit{Venus} im \textit{Widder} und \textit{Skorpion} im \textit{Schaden,}}{}
\EdL{37}{AK07-J0135-L37}{weil ihre Häuser die Wage und der Stier sind, \textit{Mars} im \textit{Stier} und in der \textit{Wage} im}{}
\NoteRule
\EdN{N1}{AK07-J0135-N01}{\textsuperscript{1}) Vergl. \Sperr{Batt.}, ed. \Sperr{Nallino}, I, 304—305.}
\EdN{N2}{AK07-J0135-N02}{\textsuperscript{2}) S. oben p. 42 u. 65.}
\EdN{N3}{AK07-J0135-N03}{\textsuperscript{3}) Ich habe die Kolonnen der obern Hälfte der Tafel der Bequemlichkeit des Zitierens halber}
\EdN{N4}{AK07-J0135-N04}{mit Nummern versehen, die im Ms. fehlen.}
\FrameHead{0136}{p. 104}{0136}{p. 104}
\EdL{1}{AK07-J0136-L01}{\textit{Schaden,} weil seine Häuser der Skorpion und der Widder sind. — Diese \textit{Alienationes}}{}
\EdL{2}{AK07-J0136-L02}{fehlen in der Rose \Sperr{Batt.}’s.}{}
\EdLp{3}{AK07-J0136-L03}{Die Kolonnen 3 und 5 enthalten die sog. \textit{Erhöhungen} (\textit{Dignitates, Exaltationes,}}{}
\EdL{4}{AK07-J0136-L04}{arab. \textit{sharaf,} pl. \textit{ashrāf}) und \textit{Tiefstände} (\textit{Deiectiones, Casus,} arab. \textit{hubūṭ}), oder vielmehr}{}
\EdL{5}{AK07-J0136-L05}{die Planeten, die in dem entsprechenden Zeichen der 1. Kolonne in der Erhöhung}{}
\EdL{6}{AK07-J0136-L06}{oder im Tiefstand (Fall) sind. So ist z. B. \textit{Jupiter} in der \textit{Erhöhung} im \textit{Krebs} (\textit{Cancer}),}{}
\EdL{7}{AK07-J0136-L07}{d. h. in diesem Zeichen hat er die höchste Macht, den grössten Einfluss, dagegen}{}
\EdL{8}{AK07-J0136-L08}{ist er im \textit{Tiefstand} oder im \textit{Fall} im \textit{Steinbock} (\textit{Capricornus}), welches Zeichen um 180°}{}
\EdL{9}{AK07-J0136-L09}{vom Krebs absteht, in diesem Zeichen hat er den geringsten oder gar keinen Einfluss.}{}
\EdL{10}{AK07-J0136-L10}{\textit{Saturn} ist in der \textit{Wage} in der \textit{Erhöhung,} dagegen im \textit{Widder} im \textit{Fall.} Im Löwen}{}
\EdL{11}{AK07-J0136-L11}{und Wassermann (\textit{Aquarius}) steht kein Planet in der Erhöhung, und keiner im Fall,}{}
\EdL{12}{AK07-J0136-L12}{im Stier keiner im Fall, im Skorpion keiner in der Erhöhung. Ausser Sonne, Mond}{}
\EdL{13}{AK07-J0136-L13}{und den fünf Planeten haben auch die Mondknoten, Drachenkopf und Drachen-}{}
\EdL{14}{AK07-J0136-L14}{schwanz ihre Exaltationes und Casus. Der 3. Kolonne entspricht der innerste Kreis}{}
\EdL{15}{AK07-J0136-L15}{der Rose \Sperr{Batt.}’s, die 5. Kolonne fehlt in derselben. In der Rose ist auch noch}{}
\EdL{16}{AK07-J0136-L16}{der Grad des Zeichens angegeben, in welchem der Planet das Maximum seiner Macht}{}
\EdL{17}{AK07-J0136-L17}{erreicht (z. B. Venus im 27.° der Fische), diese Grade fehlen in unserer Tafel. Auch}{}
\EdL{18}{AK07-J0136-L18}{den Überschriften dieser Kolonnen, \textit{Regnum} für \textit{Dignitas} oder \textit{Exaltatio,} und \textit{Servitus}}{}
\EdL{19}{AK07-J0136-L19}{für \textit{Casus} oder \textit{Deiectio} begegnen wir hier zum ersten mal.}{}
\EdLp{20}{AK07-J0136-L20}{Die Kolonnen 6 u. 7 enthalten die \textit{Herren} oder \textit{Regenten} (hier \textit{Judices} genannt)}{}
\EdL{21}{AK07-J0136-L21}{der sog. \textit{Triplizitäten} (hier \textit{Trigonalia} genannt), d. h. je dreier Zeichen des Tierkreises,}{}
\EdL{22}{AK07-J0136-L22}{die die Ecken eines gleichseitigen Dreieckes bilden. So bilden z. B. \textit{Widder, Löwe}}{}
\EdL{23}{AK07-J0136-L23}{und \textit{Schütze} eine \textit{Triplizität} (d. h. stehen unter einander im Trigonalschein) und ihre}{}
\EdL{24}{AK07-J0136-L24}{Herren (\textit{Judices, Domini}) sind \textit{Sonne, Jupiter} u. \textit{Saturn} am \textit{Tage, Jupiter, Sonne, Saturn}}{}
\EdL{25}{AK07-J0136-L25}{bei \textit{Nacht.} Eigentlich sind nur die ersten beiden die Herren, der dritte, Saturn, nur}{}
\EdL{26}{AK07-J0136-L26}{ihr (etwas untergeordneter) Genosse. Diese \textit{Judices} (\textit{Domini}) der Triplizitäten finden}{}
\EdL{27}{AK07-J0136-L27}{sich in der astrologischen Rose im zweitäussersten Ring und zwar auch für Tag}{}
\EdL{28}{AK07-J0136-L28}{und für Nacht.}{}
\EdLp{29}{AK07-J0136-L29}{Die untere Hälfte der Tafel enhält links die sog. \textit{Facies} (\textit{Decani,} arab. \textit{wadjh,}}{}
\EdL{30}{AK07-J0136-L30}{pl. \textit{wudjūh}) und ihre Herren (\textit{Domini}). Jedes Zeichen des Tierkreises wird nämlich}{}
\EdL{31}{AK07-J0136-L31}{in drei gleiche Teile von je 10° geteilt, und jedem dieser Teile ein Planet (Sonne}{}
\EdL{32}{AK07-J0136-L32}{und Mond inbegriffen) als Herr zugeteilt. Diese Verteilung wurde nicht bei allen}{}
\EdL{33}{AK07-J0136-L33}{Völkern, die sich mit Astrologie beschäftigt haben, nach dem gleichen Prinzip durch-}{}
\EdL{34}{AK07-J0136-L34}{geführt\,; die hier gegebene Verteilung ist die ägyptisch-griechische\textsuperscript{1}), die von den}{}
\EdL{35}{AK07-J0136-L35}{meisten arabischen Astrologen adoptiert worden ist\,: dem ersten Decan des Widders}{}
\EdL{36}{AK07-J0136-L36}{wird sein Hausherr Mars zugeteilt, dem zweiten die Sonne, dem dritten die Venus,}{}
\EdL{37}{AK07-J0136-L37}{dem ersten Decan des Stiers Merkur, u. s. f., stets in der alten Reihefolge dieser}{}
\EdL{38}{AK07-J0136-L38}{Himmelskörper\,: Saturn, Jupiter, Mars, Sonne, Venus, Merkur, Mond. — Die astro-}{}
\EdL{39}{AK07-J0136-L39}{logische Rose enthält die Facies im 3. (mittlern) Kreisring, und zwar in gleicher}{}
\EdL{40}{AK07-J0136-L40}{Weise verteilt wie in unserer Tafel.}{}
\EdLp{41}{AK07-J0136-L41}{Rechts in der untern Hälfte der Tafel sind die \textit{Grenzen} (\textit{Termini,} hier \textit{Fines}}{}
\NoteRule
\EdN{N1}{AK07-J0136-N01}{\textsuperscript{1}) \Sperr{Ptolemäus} ausgenommen, der die \textit{Facie}s nicht kennt.}
\FrameHead{0137}{p. 105}{0137}{p. 105}
\EdL{1}{AK07-J0137-L01}{genannt, arab. \textit{ḥadd,} pl. \textit{ḥudūd}) verzeichnet. Jedes Tierkreiszeichen wird auch in}{}
\EdL{2}{AK07-J0137-L02}{fünf, im allgemeinen ungleiche Teile geteilt, und jedem einer der fünf Planeten}{}
\EdL{3}{AK07-J0137-L03}{(Sonne und Mond fallen hier weg) als Herr zugeteilt. So wird den ersten 6° des}{}
\EdL{4}{AK07-J0137-L04}{Widders \textit{Jupiter,} den nächsten 6° \textit{Venus,} den folgenden 8° \textit{Merkur,} den weitern 5°}{}
\EdL{5}{AK07-J0137-L05}{\textit{Mars} und den letzten 5° \textit{Saturn} zugeteilt. Auch hier existierten verschiedene Ver-}{}
\EdL{6}{AK07-J0137-L06}{teilungssysteme, das hier befolgte weicht von dem der Rose \Sperr{Batt.}’s (im äussersten}{}
\EdL{7}{AK07-J0137-L07}{Kreisring) ab, und scheint das von \Sperr{Nallino}\textsuperscript{1}) den \Sperr{Aegyptern} zugeschriebene}{}
\EdL{8}{AK07-J0137-L08}{System zu sein, das auch von den meisten arabischen Astrologen angenommen wurde.}{}
\EdL{9}{AK07-J0137-L09}{Der Titel dieses Teils der Tafel ist\,: \textit{Fines secundum Medos}\,; der Verfasser schreibt}{}
\EdL{10}{AK07-J0137-L10}{also diese Verteilungsmethode den Medern zu, meint er damit Babylonier oder Perser\,?}{}
\EdL{11}{AK07-J0137-L11}{\Sperr{Batt.} nennt in seinen Königstafeln (II, 2) einen \Sperr{Darius}, den Vorgänger des \Sperr{Cyrus},}{}
\EdL{12}{AK07-J0137-L12}{einen Meder (\textit{Medus}), ebenso \Sperr{Al-Bīrūnī} in seiner Chronologie. \textsuperscript{2})}{}
\EdLp{13}{AK07-J0137-L13}{Für weitere Details in Bezug auf diese astrologischen Begriffe verweise ich}{}
\EdL{14}{AK07-J0137-L14}{ausser auf die lateinischen Ausgaben griechischer und arabischer Astrologen, wie}{}
\EdL{15}{AK07-J0137-L15}{\Sperr{Ptol.}, \Sperr{Albumasar}, \Sperr{Alcabitius}, \Sperr{Albohazen}, etc. noch auf folgende Schriften\,:}{}
\EdL{16}{AK07-J0137-L16}{\Sperr{Batt.}, ed. \Sperr{Nallino}, II, 309—314\,; \Sperr{Loth}, \textit{Al-Kindi als Astrolog}, in \textit{Morgenländische}}{}
\EdL{17}{AK07-J0137-L17}{\textit{Forschungen,} Festschrift f. \Sperr{H. L. Fleischer}, Leipzig 1875\,; \Sperr{Dieterici}, \textit{Die Propädeutik}}{}
\EdL{18}{AK07-J0137-L18}{\textit{der Araber im 10. Jahrh.}, Berlin 1865, p. 46—85\,; \Sperr{Boll}, \textit{Sphaera,} Leipzig 1903\,; \Sperr{Bouché-}}{}
\EdL{19}{AK07-J0137-L19}{\Sperr{Leclercq}, \textit{L’astrologie grecque,} Paris 1899\,; \Sperr{Uhlemann}, \textit{Grundzüge der Astronomie}}{}
\EdL{20}{AK07-J0137-L20}{\textit{und Astrologie der Alten, besonders der Aegypter,} Leipzig 1857.}{}
\EdLp{21}{AK07-J0137-L21}{Wenn wir am Schlusse zu überblicken versuchen, was sich uns über die Quellen,}{}
\EdL{22}{AK07-J0137-L22}{denen \Sperr{Kh.} gefolgt ist, aus dem Studium dieses Werkes ergeben hat, so müssen wir}{}
\EdL{23}{AK07-J0137-L23}{eingestehen, dass das Resultat kein entschiedenes, nur in wenigen Punkten sicher}{}
\EdL{24}{AK07-J0137-L24}{stehendes ist. Gewiss ist, dass er in der Bestimmung der \textit{mittlern Örter} der Planeten}{}
\EdL{25}{AK07-J0137-L25}{das \textit{indisch-persische} Verfahren adoptiert, d. h. das siderische Jahr zu Grunde gelegt}{}
\EdL{26}{AK07-J0137-L26}{hat\,; ferner dass er in der Grösse der \textit{Schiefe} der \textit{Ekliptik} dem \Sperr{Ptol.} gefolgt ist.}{}
\EdL{27}{AK07-J0137-L27}{Sehr wahrscheinlich ist\textsuperscript{3}), dass er in der Berechnung der \textit{Gleichungen} der \textit{Planeten}}{}
\EdL{28}{AK07-J0137-L28}{den Weg der \textit{Perser,} d. h. der \textit{Tafeln} des \textit{Shāh,} gegangen ist, die wiederum auf den}{}
\EdL{29}{AK07-J0137-L29}{\textit{Siddhāntas} der \textit{Inder,} auf welchen ist nicht bestimmt zu sagen, basieren. \textit{Persische}}{}
\EdL{30}{AK07-J0137-L30}{Spuren bemerken wir ferner in den Tafeln der \textit{wahren} oder \textit{ungleichen Bewegung} der}{}
\EdL{31}{AK07-J0137-L31}{Sonne und des Mondes, wie die Beibehaltung des Ausdruckes \textit{buht} beweist, den \Sperr{Batt.}}{}
\EdL{32}{AK07-J0137-L32}{nicht kennt. In den Tafeln der \textit{Konjunktionen} u. \textit{Oppositionen} (\textit{Syzygien}) ist \Sperr{Kh.} (oder}{}
\EdL{33}{AK07-J0137-L33}{\Sperr{Maslama?}) wahrscheinlich hauptsächlich \Sperr{Ptol.} gefolgt, in der Berechnung der \textit{Fin-}}{}
\EdL{34}{AK07-J0137-L34}{\textit{sternisse} und der \textit{Mond-Parallaxe} finden wir wieder Annäherung an die \textit{Inder.} Dass sich}{}
\EdL{35}{AK07-J0137-L35}{\Sperr{Kh.}, mehr an die indirekten persischen Quellen als an die direkten indischen gehalten}{}
\EdL{36}{AK07-J0137-L36}{haben mag, ist wahrscheinlich und auch leicht erklärlich, da er eben ein geborner}{}
\EdL{37}{AK07-J0137-L37}{Perser war\,; es ist auch möglich, dass er in der ersten Ausgabe seiner Tafeln sich}{}
\NoteRule
\EdN{N1}{AK07-J0137-N01}{\textsuperscript{1}) \Sperr{Batt.} II, 309.}
\EdN{N2}{AK07-J0137-N02}{\textsuperscript{2}) Ed. \Sperr{Sachau}, engl. Übersetzung, p. 101.}
\EdN{N3}{AK07-J0137-N03}{\textsuperscript{3}) S. oben p. 32.}
\FrameHead{0138}{p. 106}{0138}{p. 106}
\EdL{1}{AK07-J0138-L01}{noch mehr an die indischen \textit{Siddhāntas,} in der zweiten mehr an die \textit{Tafeln} des \textit{Shāh,}}{}
\EdL{2}{AK07-J0138-L02}{an \Sperr{Ptol.} und vielleicht auch an die \textit{erprobten Tafeln} der Astronomen \Sperr{Al-Ma’mūns}}{}
\EdL{3}{AK07-J0138-L03}{angelehnt hat.}{}
\EdLp{4}{AK07-J0138-L04}{Ebenso ist auch unsere Kenntnis von den Veränderungen, die \Sperr{Maslama al-}}{}
\EdL{5}{AK07-J0138-L05}{\Sperr{Madjrīṭī} an den Tafeln \Sperr{Kh.}’s vorgenommen hat, in gewissen Punkten mangelhaft.}{}
\EdL{6}{AK07-J0138-L06}{Sicher können wir nur behaupten, dass \Sperr{Maslama} in den Tafeln für die Bestim-}{}
\EdL{7}{AK07-J0138-L07}{mung der \textit{mittlern Örter} der Planeten die \textit{persische Aera} \Sperr{Yezdegirds} durch die}{}
\EdL{8}{AK07-J0138-L08}{\textit{muhammedanische} ersetzt hat, und dass er die Tafeln der \textit{Syzygien, Finsternisse} und}{}
\EdL{9}{AK07-J0138-L09}{der \textit{Mond-Parallaxe} (wenn diese überhaupt im ursprünglichen Werke \Sperr{Kh.}’s vorhanden}{}
\EdL{10}{AK07-J0138-L10}{waren), ebenso die Tafeln der \textit{Verteilung} der zwölf \textit{astrologischen Häuser} und die-}{}
\EdL{11}{AK07-J0138-L11}{jenigen der \textit{Proiectiones radiorum} für den Meridian und die Breite von \textit{Cordova} um-}{}
\EdL{12}{AK07-J0138-L12}{gearbeitet hat. Was die letztern Tafeln anbetrifft, so hat er nach dem, was wir aus}{}
\EdL{13}{AK07-J0138-L13}{\Sperr{Ibn Hibintā} wissen (s. oben p. 100) noch einige weitere Änderungen daran}{}
\EdL{14}{AK07-J0138-L14}{vorgenommen. Von \Sperr{Maslama} selbst hinzugefügt sind vielleicht (sicher steht dies}{}
\EdL{15}{AK07-J0138-L15}{ebenfalls nicht) einige der Tafeln 1 bis 3\textsuperscript{a} (s. oben p. 40—41), die Tafel 60 der \textit{Co-}}{}
\EdL{16}{AK07-J0138-L16}{\textit{tangenten,} die Tafeln 67—68 der \textit{Zeitgleichung,} und die zwei letzten Tafeln 115 u. 116}{}
\EdL{17}{AK07-J0138-L17}{über die \textit{Umwandlung} der \textit{Geburtsjahre} und über die \textit{Domicilia, Exaltationes} etc.}{}
\EdLp{18}{AK07-J0138-L18}{Um das Urteil darüber, was das Werk \Sperr{Kh.}’s ursprünglich enthalten haben}{}
\EdL{19}{AK07-J0138-L19}{möchte und was nicht, noch schwieriger zu machen, kommt endlich noch hinzu,}{}
\EdL{20}{AK07-J0138-L20}{dass der (oder die) Übersetzer, vielleicht auch spätere Bearbeiter und Ordner}{}
\EdL{21}{AK07-J0138-L21}{(\Sperr{Robert v. Chester}) der Übersetzung, oder auch Abschreiber, Zusätze zu dem}{}
\EdL{22}{AK07-J0138-L22}{Texte und auch zu einigen Tafeln gemacht haben, und dass diese Zusätze nicht}{}
\EdL{23}{AK07-J0138-L23}{immer so leicht als solche zu erkennen sind, wie dies bei einigen Stellen des Textes}{}
\EdL{24}{AK07-J0138-L24}{der Fall ist, die oben p. 35 u. 37 besprochen worden sind.}{}
\EdLp{25}{AK07-J0138-L25}{Aber trotz dieser mehrfachen Mit- und Abänderungsarbeit an diesen Tafeln}{}
\EdL{26}{AK07-J0138-L26}{behaupten dieselben auch in der Form, wie sie uns vorliegen, neben denen \Sperr{Batt.}’s}{}
\EdL{27}{AK07-J0138-L27}{gleichwohl einen nicht zu unterschätzenden Wert\,: sie bilden sehr wahrscheinlich die}{}
\EdL{28}{AK07-J0138-L28}{Grundlage für spätere westarabische Tafeln, und damit vielleicht auch für gewisse}{}
\EdL{29}{AK07-J0138-L29}{Tafeln des christlichen Mittelalters\,; direkte Beweise können wir allerdings für diese}{}
\EdL{30}{AK07-J0138-L30}{Ansicht noch nicht beibringen, sie mögen aber wohl noch gefunden werden\,; vielleicht}{}
\EdL{31}{AK07-J0138-L31}{stützen sich gerade die in der Einleitung erwähnten Tafeln \Sperr{Rogers von}}{}
\EdL{32}{AK07-J0138-L32}{\Sperr{Hereford} auf die vorliegenden. Sie zeigen uns ferner sehr deutlich den grossen Fort-}{}
\EdL{33}{AK07-J0138-L33}{schritt, den die arabische Astronomie von \Sperr{Kh.} bis auf \Sperr{Batt.} gemacht hat\,: \Sperr{hier} tri-}{}
\EdL{34}{AK07-J0138-L34}{gonometrische Formeln und mathematische Beweise für die Probleme, \Sperr{dort} nur nackte,}{}
\EdL{35}{AK07-J0138-L35}{kurze Regeln ohne jede Erläuterung, und ganz wenige Spuren trigonometrischer}{}
\EdL{36}{AK07-J0138-L36}{Ableitung. Wenn also durch Veröffentlichung dieser Tafeln einerseits die Grösse der}{}
\EdL{37}{AK07-J0138-L37}{Arbeit des bedeutendsten arabischen Astronomen um so glänzender hervortritt, so}{}
\EdL{38}{AK07-J0138-L38}{beweisen sie wiederum andrerseits, dass man schon zu jener Zeit auch mit ganz}{}
\EdL{39}{AK07-J0138-L39}{elementaren mathematischen Hilfsmitteln astronomische Tafeln herzustellen verstand,}{}
\EdL{40}{AK07-J0138-L40}{die, wie wir gesehen haben, nicht geringe Ansprüche auf Genauigkeit zu befriedigen}{}
\EdL{41}{AK07-J0138-L41}{im stande waren.}{}
\FrameHead{0139}{p. [107]}{0139}{p. [107]}
\EdLc{1}{AK07-J0139-L01}{TAFELN}{}
\FrameHead{0140}{p. [108]}{0140}{p. [108]}
\Withheld{blank}
\FrameHead{0264}{p. [232]}{0264}{p. [232]}
\Withheld{blank}
\FrameHead{0265}{p. [233]}{0265}{p. [233]}
\EdLc{1}{AK07-J0265-L01}{NACHTRÄGE UND BERICHTIGUNGEN}{}
\EdLp{2}{AK07-J0265-L02}{\textbf{1.} Zu Kap. 10 meines Kommentars (Bestimmung des wahren Ortes der drei}{}
\EdL{3}{AK07-J0265-L03}{obern Planeten) habe ich einige Zusätze und Berichtigungen zu geben, auf die ich}{}
\EdL{4}{AK07-J0265-L04}{durch eine briefliche Mitteilung des Herren C. F. \Sperr{Pechüle}, w. Astronomen des kgl.}{}
\EdL{5}{AK07-J0265-L05}{Observatoriums in Kopenhagen, aufmerksam gemacht worden bin.}{}
\EdLp{6}{AK07-J0265-L06}{\Sperr{Nallino} (Ed. \Sperr{Batt.}, II, 241) hat zuerst darauf hingewiesen, dass die Werte der}{}
\EdL{7}{AK07-J0265-L07}{\textit{Prosthaphäresis} der Länge (Gleichung des \textit{Centrums}), die \Sperr{Ptol.} in der 3. Kol. seiner}{}
\EdL{8}{AK07-J0265-L08}{Tafeln für die wahre Länge der Planeten\textsuperscript{1}) gibt, nicht der Formel (p. 48)}{}
\EdLc{9}{AK07-J0265-L09}{tg aeq. centri = \Frac{2\textit{e}\textsubscript{1} sin \textit{m}}{\Root{1 − (\textit{e}\textsubscript{1} sin \textit{m})\textsuperscript{2}} + \textit{e}\textsubscript{1} cos \textit{m}}}{}
\EdL{10}{AK07-J0265-L10}{entsprechen, sondern der einfachern\,:}{}
\EdLc{11}{AK07-J0265-L11}{tg aeq. centri = \Frac{\textit{EN}}{\textit{HN}} = \Frac{2\textit{e}\textsubscript{1} sin \textit{m}}{1 + 2\textit{e}\textsubscript{1} cos \textit{m}}}{}
\EdL{12}{AK07-J0265-L12}{die sich aus der Fig. von p. 47 des Kommentars ergibt, wenn man das Centrum \textit{H}}{}
\EdL{13}{AK07-J0265-L13}{des Epizykels im \textit{Aequans} (den wir weggelassen haben) annimmt, und dann den}{}
\EdL{14}{AK07-J0265-L14}{Radius \textit{HE′} des Aequans = 1 setzt. Da dies aber nur angenähert richtig ist, indem}{}
\EdL{15}{AK07-J0265-L15}{das Centrum \textit{H} sich nicht im Aequans, sondern im \textit{Deferens} bewegt, korrigiert \Sperr{Ptol.}}{}
\EdL{16}{AK07-J0265-L16}{die nach dieser Formel gefundenen Werte durch die Differenzenwerte der 4. Kol.\,;}{}
\EdL{17}{AK07-J0265-L17}{diese korrigierten Werte hat dann \Sperr{Batt.} in die 3. Kol. seiner Tafeln der Gleichungen}{}
\EdL{18}{AK07-J0265-L18}{der Planeten (Ed. \Sperr{Nallino}, II, 108—135) aufgenommen.}{}
\EdLp{19}{AK07-J0265-L19}{Da nun die Werte der Gleichung des Centrums bei \Sperr{Kh.} wesentlich von denen}{}
\EdL{20}{AK07-J0265-L20}{bei \Sperr{Ptol.-Batt.} abweichen, habe ich den Schluss gezogen (p. 50), es möchte die}{}
\EdL{21}{AK07-J0265-L21}{eine Hälfte der halbierten Gleichung der Anomalie bei Kh. zur Korrektur der Gleichung}{}
\EdL{22}{AK07-J0265-L22}{des Centrums benutzt worden sein. Es scheint dies aber in der Tat nicht der Fall}{}
\EdL{23}{AK07-J0265-L23}{zu sein\,; Herr \Sperr{Pechüle} hat mir mitgeteilt, dass die Werte der Gleichung des Centrums}{}
\EdL{24}{AK07-J0265-L24}{bei \Sperr{Kh.} warhscheinlich nach der Formel \textit{C}\textsubscript{\textit{α}} = \textit{C}\textsubscript{90} · sin \textit{α} berechnet seien, wo \textit{C}\textsubscript{\textit{α}} die}{}
\EdL{25}{AK07-J0265-L25}{Gleichung des Centrums für das Argument \textit{α}, \textit{C}\textsubscript{90} diejenige für das Argument 90°,}{}
\EdL{26}{AK07-J0265-L26}{d. h. die maximale Gleichung (= 2\textit{e}) bedeutet. Besser wäre freilich nach meiner}{}
\NoteRule
\EdN{N1}{AK07-J0265-N01}{\textsuperscript{1}) Ausgabe von \Sperr{Heiberg}, II, 436—444, Übers. v. \Sperr{Manitius}, II, 261–265.}
\FrameHead{0266}{p. 234}{0266}{p. 234}
\EdL{1}{AK07-J0266-L01}{Ansicht\,: sin \textit{C}\textsubscript{\textit{α}} = 2\textit{e} · sin \textit{α}\,; man kann ja allerdings bei ganz kleinen Winkeln statt}{}
\EdL{2}{AK07-J0266-L02}{der Sinus die Bogen setzen, allein es handelt sich doch hier um Winkel von 8°—9°.}{}
\EdL{3}{AK07-J0266-L03}{Berechnet man die Gleichungen des Centrums nach der genannten Formel, so findet}{}
\EdL{4}{AK07-J0266-L04}{man das Maximum der Abweichungen zwischen den so berechneten Werten und}{}
\EdL{5}{AK07-J0266-L05}{denen der Tafeln \Sperr{Kh.’s} = 3\textsuperscript{1}/\textsubscript{2}′—4′. — Ist die Ansicht \Sperr{Pechüle’s} richtig, und}{}
\EdL{6}{AK07-J0266-L06}{es scheint mir dies wirklich der Fall zu sein, so nähme also \Sperr{Kh.} in der Formel}{}
\EdL{7}{AK07-J0266-L07}{(s. Fig. p. 47)}{}
\EdLc{8}{AK07-J0266-L08}{sin \textit{E′HE} = sin aeq. centri = \Frac{\textit{EN}}{\textit{HE}}}{}
\EdL{9}{AK07-J0266-L09}{\textit{HE} = 1 (bezw. = 60\textsuperscript{p}) an\,; dann ist, da \textit{EN} = 2\textit{e}\textsubscript{1} sin \textit{m},}{}
\EdLc{10}{AK07-J0266-L10}{sin aeq. centri = 2\textit{e}\textsubscript{1} sin \textit{m}.}{}
\EdLp{11}{AK07-J0266-L11}{Nach dieser Formel lassen sich auch die Exzentrizitäten der Planetenbahnen}{}
\EdL{12}{AK07-J0266-L12}{bestimmen, die sich bei \Sperr{Kh.} nirgends verzeichnet finden, weder im Text noch in}{}
\EdL{13}{AK07-J0266-L13}{den Tafeln\,; für Saturn hat z. B. \Sperr{Kh.} die maximale Gleichung 8°36′ bei 90° Centrum,}{}
\EdL{14}{AK07-J0266-L14}{also ist sin 8°36′ = 2\textit{e}\textsubscript{1} sin 90° = 2\textit{e}\textsubscript{1}, dies giebt 2\textit{e}\textsubscript{1} = 0,1495 = 8\textsuperscript{p}58′, \Sperr{Ptol.} hat}{}
\EdL{15}{AK07-J0266-L15}{6\textsuperscript{p}50′. Die Exzentrizitäten der Bahnen von Jupiter, Mars und Venus bei \Sperr{Kh.} stimmen,}{}
\EdL{16}{AK07-J0266-L16}{so berechnet, mit denen bei \Sperr{Ptol.-Batt.} ziemlich überein, bei denen für Saturn}{}
\EdL{17}{AK07-J0266-L17}{und Merkur dagegen zeigen sich stärkere Abweichungen, also bei Saturn mehr als 2\textsuperscript{p}.}{}
\EdLp{18}{AK07-J0266-L18}{Was die Berechnung der Gleichung der \textit{Anomalie} anbetrifft, so habe ich p. 49}{}
\EdL{19}{AK07-J0266-L19}{die Formel gegeben\,:}{}
\EdLc{20}{AK07-J0266-L20}{tg \textit{RET} = tg anom. = \Frac{\textit{r} sin \textit{ψ}}{Dist. centri epic. + \textit{r} cos \textit{ψ}}}{}
\EdL{21}{AK07-J0266-L21}{wo \textit{r} der Radius des Epizykels in Dezimalteilen des Radius des Deferens ausgedrückt,}{}
\EdL{22}{AK07-J0266-L22}{und \textit{ψ} die Anomalie bedeutet. Herr \Sperr{Pechüle} teilte mir mit, dass bei \Sperr{Ptol.} die}{}
\EdL{23}{AK07-J0266-L23}{Werte der Gleichungen der Anomalie (in Kol. 6 seiner Tafeln für die wahre Länge}{}
\EdL{24}{AK07-J0266-L24}{der Planeten) fast genau der Formel}{}
\EdLc{25}{AK07-J0266-L25}{tg \textit{B}\textsubscript{\textit{α}} = \Frac{tg \textit{B}\textsubscript{90} · sin \textit{α}}{1 + tg \textit{B}\textsubscript{90} · cos \textit{α}}}{}
\EdL{26}{AK07-J0266-L26}{entsprechen, wo \textit{Bα} die Gleichung der Anomalie für das Argument (Anomalie) \textit{α}}{}
\EdL{27}{AK07-J0266-L27}{bedeutet. Ich teilte ihm mit, dass die Werte des \Sperr{Ptol.} mit der Formel noch besser}{}
\EdL{28}{AK07-J0266-L28}{stimmen würden, wenn man statt tg \textit{B}\textsubscript{90} \textit{sin aeq. anom. max.} setzen würde, also z. B.}{}
\EdL{29}{AK07-J0266-L29}{für Saturn sin \textit{B}\textsubscript{96} statt tg \textit{B}\textsubscript{90}, denn bei 96° hat \Sperr{Ptol.} für Saturn das Maximum}{}
\EdL{30}{AK07-J0266-L30}{der Anomalie (ebenso Kh. bei 95°—98°), weil nicht bei 90°, sondern bei 96° Anomalie}{}
\EdL{31}{AK07-J0266-L31}{die Linie \textit{EP} (s. Fig. p. 49) den Epizykel berührt, so dass dann das Dreieck \textit{EHP}}{}
\EdL{32}{AK07-J0266-L32}{bei \textit{P} rechtwinklig wird, und sin \textit{HEP} = \Frac{\textit{HP}}{\textit{HE}} wird, also sin \textit{HEP} = \textit{HP} = \textit{r}, wenn}{}
\EdL{33}{AK07-J0266-L33}{\textit{HE}, d. h. der Radius des Deferens, = 1 gesetzt wird\,; die Formel würde dann also}{}
\EdL{34}{AK07-J0266-L34}{lauten\,:}{}
\EdLc{35}{AK07-J0266-L35}{tg \textit{B}\textsubscript{\textit{α}} = \Frac{\textit{r} · sin \textit{α}}{1 + \textit{r} · cos \textit{α}}}{}
\EdL{36}{AK07-J0266-L36}{also ganz entsprechend der Formel für die Gleichung des Centrums, wenn an Stelle}{}
\EdL{37}{AK07-J0266-L37}{von 2\textit{e}\textsubscript{1} d. h. der Exzentrizität der Planetenbahn, der Radius \textit{r} des Epizykels gesetzt}{}
\FrameHead{0267}{p. 235}{0267}{p. 235}
\EdL{1}{AK07-J0267-L01}{wird. — Herr \Sperr{Pechüle} hat nun auch die Werte der Gleichung der Anomalie in}{}
\EdL{2}{AK07-J0267-L02}{den Tafeln \Sperr{Kh.’s} auf diese Formel geprüft, und im allgemeinen ziemliche Überein-}{}
\EdL{3}{AK07-J0267-L03}{stimmung gefunden (grösste Abweichungen 4′—5′), nur bei Mars zeigten sich Diffe-}{}
\EdL{4}{AK07-J0267-L04}{renzen bis zu 21′ (bei einer Anomalie von 160°—170°). \Sperr{Kh.} scheint also in der}{}
\EdL{5}{AK07-J0267-L05}{oben gegebenen Formel für tg \textit{RET} die Distanz des Centrums des Epizykels von der}{}
\EdL{6}{AK07-J0267-L06}{Erde als konstant, also = 1 (bezw. 60\textsuperscript{p}) angenommen zu haben, denn dann ver-}{}
\EdL{7}{AK07-J0267-L07}{wandelt sie sich in}{}
\EdLc{8}{AK07-J0267-L08}{tg aeq. anom. = \Frac{\textit{r} sin \textit{ψ}}{1 + \textit{r} · cos \textit{ψ}}}{}
\EdL{9}{AK07-J0267-L09}{Nur bei Mars, der eine ziemlich grosse Exzentrizität hat, scheint er die Variabilität}{}
\EdL{10}{AK07-J0267-L10}{dieser Distanz etwas berücksichtigt zu haben, in welcher Weise aber ist kaum aus-}{}
\EdL{11}{AK07-J0267-L11}{findig zu machen.}{}
\EdLp{12}{AK07-J0267-L12}{Mit diesen Zusätzen zum 10. Kap. sind aber keineswegs die Gründe enthüllt,}{}
\EdL{13}{AK07-J0267-L13}{weshalb \Sperr{Kh.} das mittlere Apogeum der Planeten um die halbe Gleichung der}{}
\EdL{14}{AK07-J0267-L14}{Anomalie vermindert hat.}{}
\EdLp{15}{AK07-J0267-L15}{\textbf{2.} P. 48, Zeile 10 v. o. muss es statt „weil er aus dem Centrum des Aequans}{}
\EdL{16}{AK07-J0267-L16}{gemessen gedacht wird“, heissen\,: „weil er die Lage des Centrums des Epizykels in}{}
\EdL{17}{AK07-J0267-L17}{Bezug auf das Apogeum des Exzenters angibt“.}{}
\EdLp{18}{AK07-J0267-L18}{\textbf{3.} P. 46, etc. habe ich das astronomische Werk des Inders \Sperr{Varāhamihira}}{}
\EdL{19}{AK07-J0267-L19}{in der von \Sperr{Thibaut} nach englischer Transskriptionsweise gebrauchten Form \textit{Pan-}}{}
\EdL{20}{AK07-J0267-L20}{\textit{chasiddhāntikā} wiedergegeben\,; nach der jetzt von den Orientalisten allgemein adop-}{}
\EdL{21}{AK07-J0267-L21}{tierten Transskriptionsart sollte es \textit{Pañcasiddhāntikā} geschrieben sein.}{}
\EdLp{22}{AK07-J0267-L22}{\textbf{4.} Ich habe p. 42 (Kap. 7) den mittlern Lauf der Sonne in 360 Tagen so be-}{}
\EdL{23}{AK07-J0267-L23}{rechnet, dass ich den mittlern Lauf in 30 Tagen mit 12 multipliziert habe\,; dies ist}{}
\EdL{24}{AK07-J0267-L24}{nicht ganz richtig, denn, wenn die 5″ beim mittlern Lauf in 30 Tagen (= 29°34′5″)}{}
\EdL{25}{AK07-J0267-L25}{noch einen Bruchteil hätten, der weggeworfen wurde, so würde der dadurch ent-}{}
\EdL{26}{AK07-J0267-L26}{standene Fehler durch Multiplikation mit 12 so viel mal grösser\,; wir geben daher}{}
\EdL{27}{AK07-J0267-L27}{im folgenden eine andere Bestimmung\,:}{}
\EdL{28}{AK07-J0267-L28}{\parbox[t]{\dimexpr0.1708\linewidth-0.5em\relax}{\raggedright\small\mbox{}Mittellauf}\hspace{0.5em}\parbox[t]{\dimexpr0.0594\linewidth-0.5em\relax}{\raggedright\small\mbox{}in}\hspace{0.5em}\parbox[t]{\dimexpr0.1559\linewidth-0.5em\relax}{\raggedright\small\mbox{}1 Jahr zu}\hspace{0.5em}\parbox[t]{\dimexpr0.0742\linewidth-0.5em\relax}{\raggedright\small\mbox{}354}\hspace{0.5em}\parbox[t]{\dimexpr0.1559\linewidth-0.5em\relax}{\raggedright\small\mbox{}Tagen}\hspace{0.5em}\parbox[t]{\dimexpr0.0445\linewidth-0.5em\relax}{\raggedright\small\mbox{}=}\hspace{0.5em}\parbox[t]{\dimexpr0.1782\linewidth-0.5em\relax}{\raggedright\small\mbox{}348°54′13″}\hspace{0.5em}\parbox[t]{\dimexpr0.1411\linewidth-0.5em\relax}{\raggedright\small\mbox{}(s. T. 4)}\hspace{0.5em}}{}
\EdL{29}{AK07-J0267-L29}{\parbox[t]{\dimexpr0.1708\linewidth-0.5em\relax}{\raggedright\small\mbox{}—}\hspace{0.5em}\parbox[t]{\dimexpr0.2153\linewidth-0.5em\relax}{\raggedright\small\mbox{}-}\hspace{0.5em}\parbox[t]{\dimexpr0.0742\linewidth-0.5em\relax}{\raggedright\small\mbox{}11}\hspace{0.5em}\parbox[t]{\dimexpr0.1559\linewidth-0.5em\relax}{\raggedright\small\mbox{}—}\hspace{0.5em}\parbox[t]{\dimexpr0.0445\linewidth-0.5em\relax}{\raggedright\small\mbox{}=}\hspace{0.5em}\parbox[t]{\dimexpr0.1782\linewidth-0.5em\relax}{\raggedright\small\mbox{}10°50′30″}\hspace{0.5em}\parbox[t]{\dimexpr0.1411\linewidth-0.5em\relax}{\raggedright\small\mbox{}(s. T. 5)}\hspace{0.5em}}{}
\EdL{30}{AK07-J0267-L30}{\parbox[t]{\dimexpr0.1708\linewidth-0.5em\relax}{\raggedright\small\mbox{}—}\hspace{0.5em}\parbox[t]{\dimexpr0.2153\linewidth-0.5em\relax}{\raggedright\small\mbox{}-}\hspace{0.5em}\parbox[t]{\dimexpr0.0742\linewidth-0.5em\relax}{\raggedright\small\mbox{}6}\hspace{0.5em}\parbox[t]{\dimexpr0.1559\linewidth-0.5em\relax}{\raggedright\small\mbox{}Stunden}\hspace{0.5em}\parbox[t]{\dimexpr0.0445\linewidth-0.5em\relax}{\raggedright\small\mbox{}=}\hspace{0.5em}\parbox[t]{\dimexpr0.1782\linewidth-0.5em\relax}{\raggedright\small\mbox{}14′47″}\hspace{0.5em}\parbox[t]{\dimexpr0.1411\linewidth-0.5em\relax}{\raggedright\small\mbox{}(s. T. 5)}\hspace{0.5em}}{}
\EdL{31}{AK07-J0267-L31}{\parbox[t]{\dimexpr0.1708\linewidth-0.5em\relax}{\raggedright\small\mbox{}—}\hspace{0.5em}\parbox[t]{\dimexpr0.2153\linewidth-0.5em\relax}{\raggedright\small\mbox{}-}\hspace{0.5em}\parbox[t]{\dimexpr0.0742\linewidth-0.5em\relax}{\raggedright\small\mbox{}12}\hspace{0.5em}\parbox[t]{\dimexpr0.1559\linewidth-0.5em\relax}{\raggedright\small\mbox{}Minuten}\hspace{0.5em}\parbox[t]{\dimexpr0.0445\linewidth-0.5em\relax}{\raggedright\small\mbox{}=}\hspace{0.5em}\parbox[t]{\dimexpr0.1782\linewidth-0.5em\relax}{\raggedright\small\mbox{}30″}\hspace{0.5em}\parbox[t]{\dimexpr0.1411\linewidth-0.5em\relax}{\raggedright\small\mbox{}(s. T. 5)}\hspace{0.5em}}{}
\EdL{32}{AK07-J0267-L32}{\parbox[t]{\dimexpr0.1708\linewidth-0.5em\relax}{\raggedright\small\mbox{}Mittellauf}\hspace{0.5em}\parbox[t]{\dimexpr0.0594\linewidth-0.5em\relax}{\raggedright\small\mbox{}in}\hspace{0.5em}\parbox[t]{\dimexpr0.3861\linewidth-0.5em\relax}{\raggedright\small\mbox{}365\textsuperscript{d} 6\textsuperscript{h} 12\textsuperscript{m}}\hspace{0.5em}\parbox[t]{\dimexpr0.0445\linewidth-0.5em\relax}{\raggedright\small\mbox{}=}\hspace{0.5em}\parbox[t]{\dimexpr0.3192\linewidth-0.5em\relax}{\raggedright\small\mbox{}360° 0′ 0″}\hspace{0.5em}}{}
\EdL{33}{AK07-J0267-L33}{d. h. \Sperr{Kh.} nahm die Länge des Sternjahres zu 365\textsuperscript{d} 6\textsuperscript{h} 12\textsuperscript{m} an. Dies würde mit dem}{}
\EdL{34}{AK07-J0267-L34}{stimmen, was uns \Sperr{Al-Bīrūnī} in seiner Chronologie (ed. Sachau, arab. Text p. 11,}{}
\EdL{35}{AK07-J0267-L35}{engl. Übers. p. 12—13) über das persische Sonnenjahr berichtet\,; an einer andern}{}
\EdL{36}{AK07-J0267-L36}{Stelle desselben Werkes (arab. Text p. 52, engl. Übers. p. 61—62) fügt er dann aller-}{}
\EdL{37}{AK07-J0267-L37}{dings dieser Jahreslänge noch 9\textsuperscript{s} zu, und erhält so den Wert, den auch der \textit{Brahma-}}{}
\EdL{38}{AK07-J0267-L38}{\textit{Siddhānta} und der arabische \textit{Sindhind} haben. Es scheint also, dass auch in diesem}{}
\EdL{39}{AK07-J0267-L39}{Punkte \Sperr{Kh.} wahrscheinlich den \textit{persischen} Tafeln des \textit{Shāh} gefolgt war (vergl. auch}{}
\FrameHead{0268}{p. 236}{0268}{p. 236}
\EdL{1}{AK07-J0268-L01}{\Sperr{Batt.} ed. \Sperr{Nallino}, I, 206—209). — Diese Bemerkung ist auch in Berücksichtigung}{}
\EdL{2}{AK07-J0268-L02}{zu ziehen bei meinen Angaben p. 65 und 103.}{}
\EdLp{3}{AK07-J0268-L03}{\textbf{5.} Auf Tafel 92 in den Noten soll \textit{S... elmirech} wahrscheinlich heissen \textit{Secun-}}{}
\EdL{4}{AK07-J0268-L04}{\textit{dum Emirethium} (vergl. auch Tafel 7), und sich auf die astronomischen Tafeln des}{}
\EdL{5}{AK07-J0268-L05}{\Sperr{Eumathius} (d. i. wahrscheinlich \Sperr{Ammonius}) beziehen, die von \Sperr{Al-Zarḳālī}}{}
\EdL{6}{AK07-J0268-L06}{verbessert worden sind, von welcher Verbesserung sich ein Ms. in München (Nr. 853}{}
\EdL{7}{AK07-J0268-L07}{des Katal. v. Aumer) befindet, dessen Schluss lautet\,: „Beendigt ist der \textit{Kānōn} des}{}
\EdL{8}{AK07-J0268-L08}{\Sperr{Eumathius}, verbessert von \Sperr{Abū Isḥāḳ al-Naḳḳāsh}, bekannt unter dem Namen}{}
\EdL{9}{AK07-J0268-L09}{\Sperr{Al-Zarḳālah}, in der Mitte des \textit{Dhu’l-Ḥidjdja} 655 (Ende 1257)“. — Ammonius lebte}{}
\EdL{10}{AK07-J0268-L10}{zu Alexandria am Ende des 5. Jahrh. n. Chr. (vergl. auch \Sperr{Batt.} ed. \Sperr{Nallino}, I,}{}
\EdL{11}{AK07-J0268-L11}{XXXV, N. 5). — Ich habe auch an \textit{al-mirrīkh} = Mars gedacht, ebenso an \textit{al-mizādj}}{}
\EdL{12}{AK07-J0268-L12}{= Temperament, astrol. Natur eines Gestirns (vergl. \Sperr{Batt.} ed. \Sperr{Nallino}, II, 352—353),}{}
\EdL{13}{AK07-J0268-L13}{aber beides gibt an jener Stelle (Tafel 92) keinen Sinn.}{}
\EdLp{14}{AK07-J0268-L14}{\textbf{6.} Als diese Arbeit schon zur Hälfte gedruckt war, erhielt ich von Herrn}{}
\EdL{15}{AK07-J0268-L15}{Prof. C. H. \Sperr{Haskins} in Cambridge (Mass.) die Mitteilung, dass im Ms. 283 des \textit{Cor-}}{}
\EdL{16}{AK07-J0268-L16}{\textit{pus Christi Coll.} in Oxford eine Abschrift von \Sperr{Athelhards} Übersetzung der Tafeln}{}
\EdL{17}{AK07-J0268-L17}{\Sperr{Kb.’s} vorhanden sei, die aber wahrscheinlich ziemlich unvollständig sei. Er bemerkt}{}
\EdL{18}{AK07-J0268-L18}{noch, dass, soweit er es feststellen konnte, dieser Text eher mit dem des Oxforder}{}
\EdL{19}{AK07-J0268-L19}{Ms. \textit{Auct. F. I. 9} als mit dem der Mss. von Chartres und Madrid übereinstimme. Es}{}
\EdL{20}{AK07-J0268-L20}{wäre allerdings von Interesse, dieses Ms. nachträglich noch zur Vergleichung heran-}{}
\EdL{21}{AK07-J0268-L21}{zuziehen\,; ich muss aber diese Arbeit solchen Gelehrten überlassen, denen dieses}{}
\EdL{22}{AK07-J0268-L22}{Ms. leichter zugänglich ist als mir, also in erster Linie den englischen.}{}
\EdLp{23}{AK07-J0268-L23}{\textbf{7.} Es wurde in Kap. 34 vergessen, den Leser darauf aufmerksam zu machen,}{}
\EdL{24}{AK07-J0268-L24}{dass in den Tafeln 77 und 77\textsuperscript{a} die Parallaxe der Breite in Kreisteilen, nicht in}{}
\EdL{25}{AK07-J0268-L25}{Zeitteilen wie diejenige der Länge, angegeben ist. — In Tafel 77\textsuperscript{a} ist die Überschrift}{}
\EdL{26}{AK07-J0268-L26}{\textit{Longitudo et Latitudo} irreführend, es sind nämlich die in diesen Kolonnen für die}{}
\EdL{27}{AK07-J0268-L27}{Argumente 91—150 verzeichneten Werte nur Parallaxen der Länge.}{}
\FrameHead{0269}{unnumbered}{0269}{unnumbered}
\EdLc{1}{AK07-J0269-L01}{\textbf{Nachtrag zu den Berichtigungen.}}{}
\EdLp{2}{AK07-J0269-L02}{S. 47, Z. 14 v. u. soll es statt „Die obern Planeten vol-}{}
\EdL{3}{AK07-J0269-L03}{lenden . . . . Zeiten“ heissen\,: Die beiden obersten Planeten}{}
\EdL{4}{AK07-J0269-L04}{(Jupiter und Saturn) vollenden ihren Umlauf im Epizykel}{}
\EdL{5}{AK07-J0269-L05}{ungefähr in einem Jahre, Mars dagegen in ca. 2 Jahren}{}
\EdL{6}{AK07-J0269-L06}{(780 Tagen)\,; die Centren ihrer Epizykeln vollenden ihren Um-}{}
\EdL{7}{AK07-J0269-L07}{lauf im Deferens in verschiedenen Zeiten, die mit den heu-}{}
\EdL{8}{AK07-J0269-L08}{tigen siderischen Umlaufszeiten ziemlich übereinstimmen.}{}
\FrameHead{0270}{blank}{0270}{blank}
\Withheld{blank}
\FrameHead{0271}{p. [237]}{0271}{p. [237]}
\EdLc{1}{AK07-J0271-L01}{\textbf{Druckfehler.}}{}
\EdLc{2}{AK07-J0271-L02}{(Bei Angabe der Zeilenzahl sind die Fuss-Noten nicht berücksichtigt).}{}
\EdL{3}{AK07-J0271-L03}{\mbox{\parbox[t]{\dimexpr0.0605\linewidth-0.5em\relax}{\raggedright\small\mbox{}Pag.}\hspace{0.5em}\parbox[t]{\dimexpr0.0581\linewidth-0.5em\relax}{\raggedright\small\mbox{}3,}\hspace{0.5em}\parbox[t]{\dimexpr0.9315\linewidth-0.5em\relax}{\raggedright\small\mbox{}Note 32 lies kiteb hea (librum canonum N) statt librum hea (canonum N).}\hspace{0.5em}}}{}
\EdL{4}{AK07-J0271-L04}{\mbox{\parbox[t]{\dimexpr0.0605\linewidth-0.5em\relax}{\raggedright\small\mbox{}—}\hspace{0.5em}\parbox[t]{\dimexpr0.0581\linewidth-0.5em\relax}{\raggedright\small\mbox{}6,}\hspace{0.5em}\parbox[t]{\dimexpr0.9315\linewidth-0.5em\relax}{\raggedright\small\mbox{}Z. 10 v. u. lies \textsuperscript{26}) statt \textsuperscript{21}).}\hspace{0.5em}}}{}
\EdL{5}{AK07-J0271-L05}{\mbox{\parbox[t]{\dimexpr0.0605\linewidth-0.5em\relax}{\raggedright\small\mbox{}—}\hspace{0.5em}\parbox[t]{\dimexpr0.0581\linewidth-0.5em\relax}{\raggedright\small\mbox{}17,}\hspace{0.5em}\parbox[t]{\dimexpr0.9315\linewidth-0.5em\relax}{\raggedright\small\mbox{}Note 8 muss zwischen CO und alius eingefügt werden\,: \textit{sed supra} planum \textit{scrips.} C elmus-}\hspace{0.5em}}}{}
\EdL{6}{AK07-J0271-L06}{\mbox{\parbox[t]{\dimexpr0.1185\linewidth-0.5em\relax}{\raggedright\small\mbox{}}\hspace{0.5em}\parbox[t]{\dimexpr0.9315\linewidth-0.5em\relax}{\raggedright\small\mbox{}tewi, \textit{supra} diminutum elmaruz.}\hspace{0.5em}}}{}
\EdL{7}{AK07-J0271-L07}{\mbox{\parbox[t]{\dimexpr0.0605\linewidth-0.5em\relax}{\raggedright\small\mbox{}—}\hspace{0.5em}\parbox[t]{\dimexpr0.0581\linewidth-0.5em\relax}{\raggedright\small\mbox{}22,}\hspace{0.5em}\parbox[t]{\dimexpr0.9315\linewidth-0.5em\relax}{\raggedright\small\mbox{}Z. 9 v. o. lies \textsuperscript{11a}) statt \textsuperscript{10a}).}\hspace{0.5em}}}{}
\EdL{8}{AK07-J0271-L08}{\mbox{\parbox[t]{\dimexpr0.0605\linewidth-0.5em\relax}{\raggedright\small\mbox{}—}\hspace{0.5em}\parbox[t]{\dimexpr0.0581\linewidth-0.5em\relax}{\raggedright\small\mbox{}26,}\hspace{0.5em}\parbox[t]{\dimexpr0.9315\linewidth-0.5em\relax}{\raggedright\small\mbox{}Z. 1 v. u. lies \textsuperscript{39}) und \textsuperscript{40}) statt \textsuperscript{38}) und \textsuperscript{39}), ebenso in den Noten (letzte Zeile) \textsuperscript{40}) statt \textsuperscript{39}).}\hspace{0.5em}}}{}
\EdL{9}{AK07-J0271-L09}{\mbox{\parbox[t]{\dimexpr0.0605\linewidth-0.5em\relax}{\raggedright\small\mbox{}—}\hspace{0.5em}\parbox[t]{\dimexpr0.0581\linewidth-0.5em\relax}{\raggedright\small\mbox{}—,}\hspace{0.5em}\parbox[t]{\dimexpr0.9315\linewidth-0.5em\relax}{\raggedright\small\mbox{}Z. 8 v. u. lies tabulam statt tabulum.}\hspace{0.5em}}}{}
\EdL{10}{AK07-J0271-L10}{\mbox{\parbox[t]{\dimexpr0.0605\linewidth-0.5em\relax}{\raggedright\small\mbox{}—}\hspace{0.5em}\parbox[t]{\dimexpr0.0581\linewidth-0.5em\relax}{\raggedright\small\mbox{}31,}\hspace{0.5em}\parbox[t]{\dimexpr0.9315\linewidth-0.5em\relax}{\raggedright\small\mbox{}Z. 15 v. o. (Anfang) lies tetragoni statt trigoni.}\hspace{0.5em}}}{}
\EdL{11}{AK07-J0271-L11}{\mbox{\parbox[t]{\dimexpr0.0605\linewidth-0.5em\relax}{\raggedright\small\mbox{}—}\hspace{0.5em}\parbox[t]{\dimexpr0.0581\linewidth-0.5em\relax}{\raggedright\small\mbox{}49,}\hspace{0.5em}\parbox[t]{\dimexpr0.9315\linewidth-0.5em\relax}{\raggedright\small\mbox{}Z. 11 v. u. lies p. VII stat p. I.}\hspace{0.5em}}}{}
\EdL{12}{AK07-J0271-L12}{\mbox{\parbox[t]{\dimexpr0.0605\linewidth-0.5em\relax}{\raggedright\small\mbox{}—}\hspace{0.5em}\parbox[t]{\dimexpr0.0581\linewidth-0.5em\relax}{\raggedright\small\mbox{}50,}\hspace{0.5em}\parbox[t]{\dimexpr0.9315\linewidth-0.5em\relax}{\raggedright\small\mbox{}Z. 7 v. u. lies 167° statt 67°.}\hspace{0.5em}}}{}
\EdL{13}{AK07-J0271-L13}{\mbox{\parbox[t]{\dimexpr0.0605\linewidth-0.5em\relax}{\raggedright\small\mbox{}—}\hspace{0.5em}\parbox[t]{\dimexpr0.0581\linewidth-0.5em\relax}{\raggedright\small\mbox{}62,}\hspace{0.5em}\parbox[t]{\dimexpr0.9315\linewidth-0.5em\relax}{\raggedright\small\mbox{}Z. 2 \textbf{v}. u. lies 48′27″ statt 48″27′.}\hspace{0.5em}}}{}
\EdL{14}{AK07-J0271-L14}{\mbox{\parbox[t]{\dimexpr0.0605\linewidth-0.5em\relax}{\raggedright\small\mbox{}—}\hspace{0.5em}\parbox[t]{\dimexpr0.0581\linewidth-0.5em\relax}{\raggedright\small\mbox{}66,}\hspace{0.5em}\parbox[t]{\dimexpr0.9315\linewidth-0.5em\relax}{\raggedright\small\mbox{}Note 1 lies p. 103 statt Kap. 38.}\hspace{0.5em}}}{}
\EdL{15}{AK07-J0271-L15}{\mbox{\parbox[t]{\dimexpr0.0605\linewidth-0.5em\relax}{\raggedright\small\mbox{}—}\hspace{0.5em}\parbox[t]{\dimexpr0.0581\linewidth-0.5em\relax}{\raggedright\small\mbox{}92,}\hspace{0.5em}\parbox[t]{\dimexpr0.9315\linewidth-0.5em\relax}{\raggedright\small\mbox{}Z. 19 v. u. lies so statt zo.}\hspace{0.5em}}}{}
\EdL{16}{AK07-J0271-L16}{\mbox{\parbox[t]{\dimexpr0.0605\linewidth-0.5em\relax}{\raggedright\small\mbox{}—}\hspace{0.5em}\parbox[t]{\dimexpr0.0581\linewidth-0.5em\relax}{\raggedright\small\mbox{}100,}\hspace{0.5em}\parbox[t]{\dimexpr0.9315\linewidth-0.5em\relax}{\raggedright\small\mbox{}Z. 9 und 10 v. o. lies Ascensio obliqua statt Rectascension.}\hspace{0.5em}}}{}
\EdL{17}{AK07-J0271-L17}{\mbox{\parbox[t]{\dimexpr0.0605\linewidth-0.5em\relax}{\raggedright\small\mbox{}—}\hspace{0.5em}\parbox[t]{\dimexpr0.0581\linewidth-0.5em\relax}{\raggedright\small\mbox{}120,}\hspace{0.5em}\parbox[t]{\dimexpr0.9315\linewidth-0.5em\relax}{\raggedright\small\mbox{}(Tab. 9) lies in der Abteilung „Heza Saturni supradictum“ überschrieben bei 60 (Dies)}\hspace{0.5em}}}{}
\EdL{18}{AK07-J0271-L18}{\mbox{\parbox[t]{\dimexpr0.1185\linewidth-0.5em\relax}{\raggedright\small\mbox{}}\hspace{0.5em}\parbox[t]{\dimexpr0.9315\linewidth-0.5em\relax}{\raggedright\small\mbox{}in der Kol. der „Signa“ 1 statt 0.}\hspace{0.5em}}}{}
\EdL{19}{AK07-J0271-L19}{\mbox{\parbox[t]{\dimexpr0.0605\linewidth-0.5em\relax}{\raggedright\small\mbox{}—}\hspace{0.5em}\parbox[t]{\dimexpr0.0581\linewidth-0.5em\relax}{\raggedright\small\mbox{}131}\hspace{0.5em}\parbox[t]{\dimexpr0.9315\linewidth-0.5em\relax}{\raggedright\small\mbox{}(Tab. 20) N. 7 lies in statt tn.}\hspace{0.5em}}}{}
\EdL{20}{AK07-J0271-L20}{\mbox{\parbox[t]{\dimexpr0.0605\linewidth-0.5em\relax}{\raggedright\small\mbox{}—}\hspace{0.5em}\parbox[t]{\dimexpr0.0581\linewidth-0.5em\relax}{\raggedright\small\mbox{}145}\hspace{0.5em}\parbox[t]{\dimexpr0.9315\linewidth-0.5em\relax}{\raggedright\small\mbox{}(Tab. 34) N. 1 lies Jauzegarhu statt Jauzezarhu.}\hspace{0.5em}}}{}
\FrameHead{0272}{p. [238]}{0272}{p. [238]}
\Withheld{blank}
\FrameHead{0273}{p. [239]}{0273}{p. [239]}
\EdLc{1}{AK07-J0273-L01}{GLOSSAR}{}
\EdLc{2}{AK07-J0273-L02}{UND}{}
\EdLc{3}{AK07-J0273-L03}{NAMEN- UND SACHREGISTER}{}
\FrameHead{0274}{p. [240]}{0274}{p. [240]}
\Withheld{blank}
\FrameHead{0275}{p. [241]}{0275}{p. [241]}
\EdLh{1}{AK07-J0275-L01}{Glossar.}{}
\EdLc{2}{AK07-J0275-L02}{I. Arabische (und einige persische) Wörter.}{}
\ColHead{left column}
\EdL{3}{AK07-J0275-L03}{\textbf{A.}}{}
\EdL{4}{AK07-J0275-L04}{\textit{Aakrab} s. \textit{Akrab}.}{}
\EdL{5}{AK07-J0275-L05}{\textit{Aalemet} = \textit{ʿalāmāt} (pl. v. \textit{ʿalāma}) \textit{Notulae},}{}
\EdL{6}{AK07-J0275-L06}{Zeichen, Marke, Signatur, hier die Zahlen}{}
\EdL{7}{AK07-J0275-L07}{der Wochentage (1 . . . 7 = Sonntag bis Sams-}{}
\EdL{8}{AK07-J0275-L08}{tag), T. 69.}{}
\EdL{9}{AK07-J0275-L09}{\textit{Aard} s. \textit{Ard}.}{}
\EdL{10}{AK07-J0275-L10}{\textit{Abaad} = \textit{abʿad}, weiter entfernt (\textit{longinquior}),}{}
\EdL{11}{AK07-J0275-L11}{T. 73, 74.}{}
\EdL{12}{AK07-J0275-L12}{\textit{Achir} = \textit{ākhir}, (der) letzte, zweite, T. 2, 3 etc.}{}
\EdL{13}{AK07-J0275-L13}{\textit{Acofra} = \textit{aṣ-ṣufr} (\textit{Ta’rīkh aṣ-ṣufr}), spanische}{}
\EdL{14}{AK07-J0275-L14}{Aera, auch einfach Aera (im Text \textit{Era}) ge-}{}
\EdL{15}{AK07-J0275-L15}{nannt, 3, 35, T. 1.}{}
\EdL{16}{AK07-J0275-L16}{\textit{Acrab} s. \textit{Akrab}.}{}
\EdL{17}{AK07-J0275-L17}{\textit{Aded} u. \textit{Adet} = \textit{ʿadad}, Zahl, T. 21 etc., viel-}{}
\EdL{18}{AK07-J0275-L18}{leicht bedeutet es auch bisweilen den Plural}{}
\EdL{19}{AK07-J0275-L19}{\textit{aʿdād}.}{}
\EdL{20}{AK07-J0275-L20}{\textit{Adel} (?) T. 10.}{}
\EdL{21}{AK07-J0275-L21}{\textit{Adhel} = \textit{al-dhel}, s. \textit{Dhel}.}{}
\EdL{22}{AK07-J0275-L22}{\textit{Aguze} = \textit{adjzā’} (pl. v. \textit{djuz}), Teile, T. 73.}{}
\EdL{23}{AK07-J0275-L23}{\textit{Ahemar} = \textit{aḥmar}, rot, hier auch Name des}{}
\EdL{24}{AK07-J0275-L24}{Planeten Mars, neben \textit{Mirrīkh}, T. 39—42.}{}
\EdL{25}{AK07-J0275-L25}{\textit{Ajuze} s. \textit{Aguze}.}{}
\EdL{26}{AK07-J0275-L26}{\textit{ʿAḳd}, Knoten (der Planetenbahnen); \textit{ʿaḳd al-}}{}
\EdL{27}{AK07-J0275-L27}{\textit{shamālī}, der nördliche Knoten (der Mond-}{}
\EdL{28}{AK07-J0275-L28}{bahn) 41.}{}
\EdL{29}{AK07-J0275-L29}{\textit{Akrab}, 1) = \textit{aḳrab}, näher (\textit{propinquior}) T. 75,}{}
\EdL{30}{AK07-J0275-L30}{76; 2) = \textit{ʿaḳrab}, Skorpion (8. Zeichen des}{}
\EdL{31}{AK07-J0275-L31}{Tierkreises) T. 59\textsuperscript{b}.}{}
\EdL{32}{AK07-J0275-L32}{\textit{Anāhīd}, pers. Name des Planeten Venus, XXII,}{}
\EdL{33}{AK07-J0275-L33}{T. 45.}{}
\EdL{34}{AK07-J0275-L34}{\textit{Arafa} (?) 15 (N. 20).}{}
\EdL{35}{AK07-J0275-L35}{\textit{Ard} = \textit{ʿarḍ}, Breite eines Gestirns, oder geogr.}{}
\EdL{36}{AK07-J0275-L36}{Breite eines Ortes, 19 (N. 1), T. 22, 79.}{}
\EdL{37}{AK07-J0275-L37}{\textit{Arhik} T. 15 (N. 7) s. \textit{Achir}.}{}
\EdL{38}{AK07-J0275-L38}{\textit{Arin} = \textit{Uzain} = \textit{Udjain}, VII, XVIII, 1, 8, 19,}{}
\EdL{39}{AK07-J0275-L39}{33, etc.}{}
\ColHead{right column}
\EdL{40}{AK07-J0275-L40}{\textit{Artifa} = \textit{irtifāʿ}, Höhe (eines Gestirns), 18, 21,}{}
\EdL{41}{AK07-J0275-L41}{22, etc. T. 60.}{}
\EdL{42}{AK07-J0275-L42}{\textit{Aṣbaʿ}, Finger, Zoll, ein arab. Längenmass, 43.}{}
\EdL{43}{AK07-J0275-L43}{\textit{Asceir} s. \textit{Sceir}.}{}
\EdL{44}{AK07-J0275-L44}{\textit{Ascowa} s. \textit{Shoa}.}{}
\EdL{45}{AK07-J0275-L45}{\textit{Ashoa} s. \textit{Shoa}.}{}
\EdL{46}{AK07-J0275-L46}{\textit{Asmir} = \textit{as-sair}, oder \textit{al-masīr}, der Lauf,}{}
\EdL{47}{AK07-J0275-L47}{die Bahn (eines Gestirns), XXI, T. 9.}{}
\EdL{48}{AK07-J0275-L48}{\textit{Athofen} = \textit{aṭ-ṭūfān}, die Sündflut (\textit{diluvium}) 5,}{}
\EdL{49}{AK07-J0275-L49}{35, T. 1.}{}
\EdL{50}{AK07-J0275-L50}{\textit{Aub}, verdorben aus \textit{Aug}, s. dieses.}{}
\EdL{51}{AK07-J0275-L51}{\textit{Audel}, wahrsch. verdorben aus \textit{ʿarḍ el}, Breite}{}
\EdL{52}{AK07-J0275-L52}{des . . ., oder Breite der . . . T. 10.}{}
\EdL{53}{AK07-J0275-L53}{\textit{Aug} = \textit{audj} (\textit{awdj}), \textit{Apogeum} = Erdferne}{}
\EdL{54}{AK07-J0275-L54}{eines Planeten, 9, 10, 14\,; \textit{aug el-musteri}}{}
\EdL{55}{AK07-J0275-L55}{\textit{muadala} = korrigiertes Apogeum des Ju-}{}
\EdL{56}{AK07-J0275-L56}{piter, T. 36\,; s. auch \textit{Muadala}.}{}
\EdL{57}{AK07-J0275-L57}{\textit{Autarid} s. \textit{Otari[d]}.}{}
\EdL{58}{AK07-J0275-L58}{\textit{Avvel} s. \textit{Awel}.}{}
\EdL{59}{AK07-J0275-L59}{\textit{Avtarid} s. \textit{Otari[d]}.}{}
\EdL{60}{AK07-J0275-L60}{\textit{Awel} = \textit{awwel}, (der) erste, T. 2, 3, etc.}{}
\EdL{61}{AK07-J0275-L61}{\textit{Azabaa} = \textit{aṣābiʿ} (pl. v. \textit{aṣbaʿ}), Finger (\textit{digiti}),}{}
\EdL{62}{AK07-J0275-L62}{Mass für den Sonnen- und Monddurch-}{}
\EdL{63}{AK07-J0275-L63}{messer, auch für den Erdschatten bei Mond-}{}
\EdL{64}{AK07-J0275-L64}{finsternissen, T. 73.}{}
\EdL{65}{AK07-J0275-L65}{\textit{Azofar} s. \textit{Acofra}.}{}
\EdL{66}{AK07-J0275-L66}{\textbf{B.}}{}
\EdL{67}{AK07-J0275-L67}{\textit{Bāb}, Pforte, Kapitel, 58.}{}
\EdL{68}{AK07-J0275-L68}{\textit{Baheram} u. \textit{Bahram} = \textit{Bahrām}, pers. Name}{}
\EdL{69}{AK07-J0275-L69}{für den Planeten Mars, XXII, T. 39—42.}{}
\EdL{70}{AK07-J0275-L70}{\textit{Batiae}, wahrsch. das fem. zu \textit{baṭī’} = \textit{baṭī’a},}{}
\EdL{71}{AK07-J0275-L71}{langsam, T. 24.}{}
\EdL{72}{AK07-J0275-L72}{\textit{Bei ed halt} s. das folgende.}{}
\EdL{73}{AK07-J0275-L73}{\textit{Beied hakht} = \textit{Bēdukht}, pers. Name für den}{}
\EdL{74}{AK07-J0275-L74}{Planeten Venus, XXII, T. 45.}{}
\FrameHead{0276}{p. 242}{0276}{p. 242}
\ColHead{left column}
\EdL{1}{AK07-J0276-L01}{\textit{Bergis} = \textit{Birdjis}, pers. Name für den Planeten}{}
\EdL{2}{AK07-J0276-L02}{Jupiter, T. 33, 36.}{}
\EdL{3}{AK07-J0276-L03}{\textit{Bi}, mit dem Artikel verbunden \textit{bil}, arab.}{}
\EdL{4}{AK07-J0276-L04}{Praepos., in, auf, mit, z. B. \textit{bilfelek}, in (od.}{}
\EdL{5}{AK07-J0276-L05}{auf) der Sphäre (s. \textit{felek}).}{}
\EdL{6}{AK07-J0276-L06}{\textit{Bhukti} (sanskr.) Genuss, Freude, 78.}{}
\EdL{7}{AK07-J0276-L07}{\textit{Bood} = \textit{buʿd}, Entfernung, T. 73.}{}
\EdL{8}{AK07-J0276-L08}{\textit{Buht} (pers.), in der Astron. die wahre (un-}{}
\EdL{9}{AK07-J0276-L09}{gleiche) tägliche Bewegung der Sonne, des}{}
\EdL{10}{AK07-J0276-L10}{Mondes und der fünf Planeten, 22, 78, 90.}{}
\EdL{11}{AK07-J0276-L11}{\textit{Buiut} = \textit{buyūt} (pl. v. \textit{bait}), Häuser, inbesondere}{}
\EdL{12}{AK07-J0276-L12}{die zwölf astrologischen Häuser, 25, T. 79.}{}
\EdL{13}{AK07-J0276-L13}{\textit{Burug} = \textit{burūdj} (pl. v. \textit{burdj}), Burgen, Häuser,}{}
\EdL{14}{AK07-J0276-L14}{insbesondere die zwölf Zeichen (Häuser)}{}
\EdL{15}{AK07-J0276-L15}{des Tierkreises, 11 (N. 40), 18 (N. 23), T. 79.}{}
\EdL{16}{AK07-J0276-L16}{\textit{Buruga} s. \textit{Burug}.}{}
\EdL{17}{AK07-J0276-L17}{\textit{Burus} s. \textit{Burug}.}{}
\EdL{18}{AK07-J0276-L18}{\textbf{C.}}{}
\EdL{19}{AK07-J0276-L19}{\textit{Caat} = \textit{sāʿāt} (pl. v. \textit{sāʿa}), Stunden, vielleicht}{}
\EdL{20}{AK07-J0276-L20}{auch der Singular, T. 5.}{}
\EdL{21}{AK07-J0276-L21}{\textit{Caratan} = \textit{saraṭān}, Krebs, hier Zeichen des}{}
\EdL{22}{AK07-J0276-L22}{Krebses im Tierkreis, T. 59\textsuperscript{a}.}{}
\EdL{23}{AK07-J0276-L23}{\textit{Care} = \textit{sarīʿ}, schnell, T. 24.}{}
\EdL{24}{AK07-J0276-L24}{\textit{Carne} (?) T. 9.}{}
\EdL{25}{AK07-J0276-L25}{\textit{Carneli} (?) T. 13.}{}
\EdL{26}{AK07-J0276-L26}{\textit{Catar} = \textit{saṭar} oder \textit{saṭr}, Linie, Zeile, Kolonne,}{}
\EdL{27}{AK07-J0276-L27}{XIV\,; pl. \textit{cotor} = \textit{suṭūr}\,; \textit{catar eladed} = \textit{saṭar}}{}
\EdL{28}{AK07-J0276-L28}{\textit{al-ʿadad}, Kolonne der Zahl, T. 21, 22, 25—30,}{}
\EdL{29}{AK07-J0276-L29}{32, etc.}{}
\EdL{30}{AK07-J0276-L30}{\textit{Catur} s. \textit{Catar}.}{}
\EdL{31}{AK07-J0276-L31}{\textit{Cauz} = \textit{ḳaus} (\textit{ḳaws}), Bogen (eines Kreises),}{}
\EdL{32}{AK07-J0276-L32}{auch Zeichen des Schützen im Tierkreis,}{}
\EdL{33}{AK07-J0276-L33}{T. 59\textsuperscript{b}.}{}
\EdL{34}{AK07-J0276-L34}{\textit{Cavsc} s. \textit{Cauz}.}{}
\EdL{35}{AK07-J0276-L35}{\textit{Ciicuth} s. \textit{Zukut}.}{}
\EdL{36}{AK07-J0276-L36}{\textit{Ciikuf} s. \textit{Zukut}.}{}
\EdL{37}{AK07-J0276-L37}{\textit{Cotoreladet} = \textit{suṭūr al-ʿadad}, Kolonnen der}{}
\EdL{38}{AK07-J0276-L38}{Zahl, 27 (N. 16)\,; s. auch \textit{Catar}.}{}
\EdL{39}{AK07-J0276-L39}{\textit{Cukut} s. \textit{Zukut}.}{}
\EdL{40}{AK07-J0276-L40}{\textit{Cumbula} = \textit{sumbula}, Ähre, hier Tierkreis-}{}
\EdL{41}{AK07-J0276-L41}{zeichen der Jungfrau, T. 59\textsuperscript{b}.}{}
\EdL{42}{AK07-J0276-L42}{\textit{Cuzuf} s. \textit{Kucuf}.}{}
\EdL{43}{AK07-J0276-L43}{\textbf{D.}}{}
\EdL{44}{AK07-J0276-L44}{\textit{Dakaicae} = \textit{daḳā’iḳ} (pl. v. \textit{daḳīḳa}) Minuten}{}
\EdL{45}{AK07-J0276-L45}{(60. Teil des Grades oder der Stunde) XXIII,}{}
\EdL{46}{AK07-J0276-L46}{XXV\,; 3, 4, 36, etc.}{}
\EdL{47}{AK07-J0276-L47}{\textit{Dakaik} s. \textit{Dakaicae}.}{}
\EdL{48}{AK07-J0276-L48}{\textit{Dakeca} s. \textit{Dakaicae}.}{}
\EdL{49}{AK07-J0276-L49}{\textit{Darega} = \textit{daradja} (pl. \textit{daradj}) Grad (der 360.}{}
\EdL{50}{AK07-J0276-L50}{Teil des Kreisumfanges) 11 (N. 40).}{}
\ColHead{right column}
\EdL{51}{AK07-J0276-L51}{\textit{Dareis}, verdorben aus \textit{daregis}, Dat. plur. v.}{}
\EdL{52}{AK07-J0276-L52}{\textit{darega}, s. dieses.}{}
\EdL{53}{AK07-J0276-L53}{\textit{Dehi} = \textit{dalw} (\textit{dalu}), grosser Krug, Eimer, Tier-}{}
\EdL{54}{AK07-J0276-L54}{kreiszeichen des Wassermanns (\textit{Aquarius})}{}
\EdL{55}{AK07-J0276-L55}{T. 59.}{}
\EdL{56}{AK07-J0276-L56}{\textit{Delidhe} s. \textit{Thelitha}.}{}
\EdL{57}{AK07-J0276-L57}{\textit{Dhel} = \textit{ẓill} Schatten, dann die trigon. Funk-}{}
\EdL{58}{AK07-J0276-L58}{tion der \textit{Cotangens}, T. 60.}{}
\EdL{59}{AK07-J0276-L59}{\textit{Dhirāʿ}, Elle (ein Längenmass) 43.}{}
\EdL{60}{AK07-J0276-L60}{\textit{Djebr}, Wiederherstellung\,; \textit{al-djebr wa’l-muḳā-}}{}
\EdL{61}{AK07-J0276-L61}{\textit{bala} = Algebra, VIII\,; s. auch \textit{Muḳābala}.}{}
\EdL{62}{AK07-J0276-L62}{\textit{Dulheia} u. \textit{Dulheiga} = \textit{Dhū’l-ḥidjdja}, der}{}
\EdL{63}{AK07-J0276-L63}{zwölfte Monat der Muhammedaner, Monat}{}
\EdL{64}{AK07-J0276-L64}{der Wallfahrt, XIV, XV, 2, T. 2, 3, etc.}{}
\EdL{65}{AK07-J0276-L65}{\textit{Dulkada} = \textit{Dhū’l-ḳaʿda}, der elfte Monat der}{}
\EdL{66}{AK07-J0276-L66}{Muhammedaner, XIV, XV, 2, T. 2, 3, etc.}{}
\EdL{67}{AK07-J0276-L67}{\textit{Dulkarnain} = \textit{Dhū’l-ḳarnain}, der zweigehörnte}{}
\EdL{68}{AK07-J0276-L68}{(\textit{bicornus}), Beiname Alexanders des Grossen}{}
\EdL{69}{AK07-J0276-L69}{bei den Arabern 3, 35.}{}
\EdL{70}{AK07-J0276-L70}{\textit{Dulkeda} u. \textit{Dulkidda} s. \textit{Dulkada}.}{}
\EdL{71}{AK07-J0276-L71}{\textbf{E.}}{}
\EdL{72}{AK07-J0276-L72}{\textit{Ecumbula} s. \textit{Cumbula}.}{}
\EdL{73}{AK07-J0276-L73}{\textit{Ekucufet} s. \textit{Kucufet}.}{}
\EdL{74}{AK07-J0276-L74}{\textit{Eneid} s. \textit{Anāhīd}.}{}
\EdL{75}{AK07-J0276-L75}{\textit{Etdeiard} s. \textit{Jazdeiard}.}{}
\EdL{76}{AK07-J0276-L76}{\textit{Ezdeiard} s. \textit{Jazdeiard}.}{}
\EdL{77}{AK07-J0276-L77}{\textit{Ezed} = \textit{asad}, Löwe, Tierkreiszeichen des Lö-}{}
\EdL{78}{AK07-J0276-L78}{wen, T. 59\textsuperscript{a}.}{}
\EdL{79}{AK07-J0276-L79}{\textit{Ezeig} s. \textit{Zich}.}{}
\EdL{80}{AK07-J0276-L80}{\textit{Ezich} s. \textit{Zich}.}{}
\EdL{81}{AK07-J0276-L81}{\textbf{F.}}{}
\EdL{82}{AK07-J0276-L82}{\textit{Felek} = \textit{falak}, Sphäre, Himmelsgewölbe, kreis-}{}
\EdL{83}{AK07-J0276-L83}{förmige Bahn eines Himmelkörpers um die}{}
\EdL{84}{AK07-J0276-L84}{Erde, 7, 18, 25, etc.\,; \textit{felek elmustakim} s.}{}
\EdL{85}{AK07-J0276-L85}{\textit{Mustakim}\,; \textit{felek eltedewir} s. \textit{Tedewir}.}{}
\EdL{86}{AK07-J0276-L86}{\textit{Feler} s. \textit{Felek}.}{}
\EdL{87}{AK07-J0276-L87}{\textit{Fi}, mit dem Artikel verbunden \textit{fil} (= \textit{fi’l}), arab.}{}
\EdL{88}{AK07-J0276-L88}{Praepos., in auf, z. B. \textit{fiaard} = \textit{fī ʿarḍ}, in}{}
\EdL{89}{AK07-J0276-L89}{Breite (mit folg. Genitiv) T. 79 (Hier ist hinter}{}
\EdL{90}{AK07-J0276-L90}{\textit{fiaard} ausgefallen \textit{Cordubae}) — \textit{filcaat} und}{}
\EdL{91}{AK07-J0276-L91}{\textit{filacaat} = \textit{fī’l-sāʿa}, in der Stunde, oder}{}
\EdL{92}{AK07-J0276-L92}{\textit{fī’l-sāʿāt}, in den Stunden, T. 5 — \textit{filcinun} =}{}
\EdL{93}{AK07-J0276-L93}{\textit{fī’l-sinīn}, in den Jahren, T. 69 — \textit{fil dakaik}}{}
\EdL{94}{AK07-J0276-L94}{\textit{elcaat} = \textit{fī daḳā’iḳ al-sāʿa}, in den Minuten}{}
\EdL{95}{AK07-J0276-L95}{der Stunde, T. 5 — \textit{filiaum} = \textit{fī’l-yaum}, im}{}
\EdL{96}{AK07-J0276-L96}{Tage T. 5 ‐‐ \textit{filshohor} = \textit{fī’l-shuhūr}, in den}{}
\EdL{97}{AK07-J0276-L97}{Monaten, T. 9 — \textit{fil zahat} s. \textit{filcaat}.}{}
\EdL{98}{AK07-J0276-L98}{\textit{Fudhul} = \textit{fuḍūl} (pl. v. \textit{faḍl}), Überschüsse, Dif-}{}
\EdL{99}{AK07-J0276-L99}{ferenzen, 27 (N. 19), T. 73.}{}
\FrameHead{0277}{p. 243}{0277}{p. 243}
\ColHead{left column}
\EdL{1}{AK07-J0277-L01}{\textbf{G.}}{}
\EdL{2}{AK07-J0277-L02}{\textit{Gauze} s. \textit{Jauze}.}{}
\EdL{3}{AK07-J0277-L03}{\textit{Gauzehar} = \textit{djauzahar} (pers.), die Knoten-}{}
\EdL{4}{AK07-J0277-L04}{punkte der Planetenbahnen mit der Ekliptik,}{}
\EdL{5}{AK07-J0277-L05}{insbesondere die Mondknoten (s. auch \textit{ʿAḳd})}{}
\EdL{6}{AK07-J0277-L06}{15, 41, T. 19, 20.}{}
\EdL{7}{AK07-J0277-L07}{\textit{Gavzehar} s. \textit{Gauzehar}.}{}
\EdL{8}{AK07-J0277-L08}{\textit{Gedi} = \textit{djady} (od. \textit{djadi}), Bock, Böcklein, Tier-}{}
\EdL{9}{AK07-J0277-L09}{kreiszeichen des Steinbocks, T. 59.}{}
\EdL{10}{AK07-J0277-L10}{\textit{Gedval} = \textit{djadwal}, Tafel, Tabelle (astron.) 28,}{}
\EdL{11}{AK07-J0277-L11}{T. 60—66, 69‐‐71.}{}
\EdL{12}{AK07-J0277-L12}{\textit{Geib} = \textit{djaib}, die trigon. Funktion des \textit{Sinus},}{}
\EdL{13}{AK07-J0277-L13}{17—19, 21, 22, 69, 71, T. 58, 58\textsuperscript{a}.}{}
\EdL{14}{AK07-J0277-L14}{\textit{Genubi} = \textit{djanūbī}, südlich, 13, 16, 18, 29 (N. 1,}{}
\EdL{15}{AK07-J0277-L15}{6, 14, 15)\,; \textit{genubia} = \textit{djanūbīya} (fem). 29}{}
\EdL{16}{AK07-J0277-L16}{(N. 1, 6, 15).}{}
\EdL{17}{AK07-J0277-L17}{\textit{Genubium} s. d. vorhergehende.}{}
\EdL{18}{AK07-J0277-L18}{\textit{Gidher} = \textit{djidhr}, Wurzel, insbesondere Qua-}{}
\EdL{19}{AK07-J0277-L19}{dratwurzel, 22.}{}
\EdL{20}{AK07-J0277-L20}{\textit{Gumedi} s. \textit{Jumedi}.}{}
\EdL{21}{AK07-J0277-L21}{\textbf{H.}}{}
\EdL{22}{AK07-J0277-L22}{\textit{Haazem} (?) XXII, 20 (N. 3).}{}
\EdL{23}{AK07-J0277-L23}{\textit{Hacib} = \textit{ḥāsib}, Rechner, XXII, 20 (N. 3).}{}
\EdL{24}{AK07-J0277-L24}{\textit{Hacil} = \textit{ḥāṣil}, Ergebnis, Resultat (z. B. einer}{}
\EdL{25}{AK07-J0277-L25}{Rechnung) 11.}{}
\EdL{26}{AK07-J0277-L26}{\textit{Ḥadd} (pl. \textit{ḥudūd}), Grenzen, Teile der Tierkreis-}{}
\EdL{27}{AK07-J0277-L27}{zeichen (lat. \textit{termini}, \textit{fines}) 105.}{}
\EdL{28}{AK07-J0277-L28}{\textit{Hadhid} = \textit{ḥaḍīḍ}, unterster Teil, in der Astron.}{}
\EdL{29}{AK07-J0277-L29}{\textit{Perigeum} eines Planeten, in diesem Werke}{}
\EdL{30}{AK07-J0277-L30}{übersetzt mit \textit{submissio}, 14 (N. 28).}{}
\EdL{31}{AK07-J0277-L31}{\textit{Hamel} = \textit{ḥamal}, Widder, Tierkreiszeichen des}{}
\EdL{32}{AK07-J0277-L32}{Widders (\textit{Aries}) T. 59.}{}
\EdL{33}{AK07-J0277-L33}{\textit{Harket} = \textit{ḥaraka}, Bewegung\,; \textit{harket elscems}}{}
\EdL{34}{AK07-J0277-L34}{= \textit{ḥarakat al-shams}, Bewegung der Sonne,}{}
\EdL{35}{AK07-J0277-L35}{T. 61.}{}
\EdL{36}{AK07-J0277-L36}{\textit{Hea} = \textit{hai’a}, äussere Erscheinung (des Him-}{}
\EdL{37}{AK07-J0277-L37}{mels)\,; \textit{ʿilm al-hai’a}, Wissenschaft der Him-}{}
\EdL{38}{AK07-J0277-L38}{melskunde, Astronomie\,: s. auch \textit{Kiteb}.}{}
\EdL{39}{AK07-J0277-L39}{\textit{Heca} s. \textit{Heza}.}{}
\EdL{40}{AK07-J0277-L40}{\textit{Helhacib} s. \textit{Hacib}.}{}
\EdL{41}{AK07-J0277-L41}{\textit{Hemar} s. \textit{Ahemar}.}{}
\EdL{42}{AK07-J0277-L42}{\textit{Hemeli} s. \textit{Shemeli}.}{}
\EdL{43}{AK07-J0277-L43}{\textit{Heza} = \textit{ḥāṣṣa} u. \textit{khāṣṣa}, die dem Planeten}{}
\EdL{44}{AK07-J0277-L44}{eigentümliche Bewegung im Epizykel, die}{}
\EdL{45}{AK07-J0277-L45}{sog. \textit{Anomalie}, in diesem Werk durch \textit{Ar-}}{}
\EdL{46}{AK07-J0277-L46}{\textit{gumentum} wiedergegeben XIV, XXI, 9—14}{}
\EdL{47}{AK07-J0277-L47}{(u. N.), 28 (N. 1) 46, 53\,; T. 6—9, 11, 13, 15—18,}{}
\EdL{48}{AK07-J0277-L48}{etc.}{}
\EdL{49}{AK07-J0277-L49}{\textit{Higera} = \textit{Hidjra}, Flucht Muhammeds aus}{}
\EdL{50}{AK07-J0277-L50}{Mekka, Beginn der muhammed. Aera, 1, 3,}{}
\EdL{51}{AK07-J0277-L51}{33, etc. T. 1.}{}
\ColHead{right column}
\EdL{52}{AK07-J0277-L52}{\textit{Ḥiṣṣa}, Anteil, Loos, Argument, mit dem man}{}
\EdL{53}{AK07-J0277-L53}{in eine Tafel eintritt, 46.}{}
\EdL{54}{AK07-J0277-L54}{\textit{Hormi} u. \textit{Horom}, verdorben aus \textit{Hormuzd}}{}
\EdL{55}{AK07-J0277-L55}{od. \textit{Hormuz} = \textit{Ormuzd}, pers. Name für den}{}
\EdL{56}{AK07-J0277-L56}{Planeten Jupiter, XXII, T. 33, 36.}{}
\EdL{57}{AK07-J0277-L57}{\textit{Hout} = \textit{ḥūt}, Fisch, Tierkreiszeichen der Fische,}{}
\EdL{58}{AK07-J0277-L58}{T. 59.}{}
\EdL{59}{AK07-J0277-L59}{\textit{Hubūṭ}, Fall, Fallen, in der Astrologie der Tief-}{}
\EdL{60}{AK07-J0277-L60}{stand, Fall (\textit{casus}) eines Planeten, 104.}{}
\EdL{61}{AK07-J0277-L61}{\textbf{I.}}{}
\EdL{62}{AK07-J0277-L62}{\textit{Iblīs}, Satan, Teufel, XXI.}{}
\EdL{63}{AK07-J0277-L63}{\textit{Idhnascer} = \textit{ithnā ʿashar}, zwölf, T. 79.}{}
\EdL{64}{AK07-J0277-L64}{\textit{Igstimaat} = \textit{idjtimāʿāt} (pl. v. \textit{idjtimāʿ}) s. \textit{Istima}.}{}
\EdL{65}{AK07-J0277-L65}{\textit{Ihtilef} = \textit{ikhtilāf}, Verschiedenheit T. 73\,; \textit{elihtilef}}{}
\EdL{66}{AK07-J0277-L66}{\textit{lilmanshar} = \textit{al-ikhtilāf lil-manẓar}, die}{}
\EdL{67}{AK07-J0277-L67}{Verschiedenheit des Anblicks (eines Ge-}{}
\EdL{68}{AK07-J0277-L68}{stirns), d. h. seine \textit{Parallaxe} 28 (N. 3).}{}
\EdL{69}{AK07-J0277-L69}{\textit{Ihtilief}, verdorben aus \textit{Ihtilef}, s. dieses.}{}
\EdL{70}{AK07-J0277-L70}{\textit{Ingile} = \textit{indjilā’}, Verschwinden, insbesondere}{}
\EdL{71}{AK07-J0277-L71}{der Ausgang oder das Ende einer Finsternis,}{}
\EdL{72}{AK07-J0277-L72}{die Phase von der Totalität bis zum Ende}{}
\EdL{73}{AK07-J0277-L73}{(lat. \textit{Emersio}), 26, 87.}{}
\EdL{74}{AK07-J0277-L74}{\textit{Ingine} s. \textit{Ingile}.}{}
\EdL{75}{AK07-J0277-L75}{\textit{Iṣbaʿ} s. \textit{Aṣbaʿ}.}{}
\EdL{76}{AK07-J0277-L76}{\textit{Ishemz} verdorben aus \textit{eshemz} = \textit{el-shems} s.}{}
\EdL{77}{AK07-J0277-L77}{\textit{Shems}.}{}
\EdL{78}{AK07-J0277-L78}{\textit{Isticbelet} = \textit{istiḳbālāt} (pl. v. \textit{istiḳbāl}) s. \textit{Istikbel}.}{}
\EdL{79}{AK07-J0277-L79}{\textit{Istikbel} = \textit{istiḳbāl}, Opposition (zweier Gestirne),}{}
\EdL{80}{AK07-J0277-L80}{insbesondere von Sonne und Mond, 23, T. 70.}{}
\EdL{81}{AK07-J0277-L81}{\textit{Istima} = \textit{idjtimāʿ}, Konjunktion (zweier Ge-}{}
\EdL{82}{AK07-J0277-L82}{stirne), insbesondere von Sonne und Mond,}{}
\EdL{83}{AK07-J0277-L83}{23, T. 69.}{}
\EdL{84}{AK07-J0277-L84}{\textit{Istimaz} s. \textit{Istima}.}{}
\EdL{85}{AK07-J0277-L85}{\textbf{J.}}{}
\EdL{86}{AK07-J0277-L86}{\textit{Jauze} = \textit{djauzā’} (oder \textit{djawzā’}), das Tierkreis-}{}
\EdL{87}{AK07-J0277-L87}{zeichen der Zwillinge, T. 59\textsuperscript{a}\,; auch das}{}
\EdL{88}{AK07-J0277-L88}{Sternbild des Orion.}{}
\EdL{89}{AK07-J0277-L89}{\textit{Jauzehar} s. \textit{Gauzehar}.}{}
\EdL{90}{AK07-J0277-L90}{\textit{Jazdeiard} = \textit{Yezdegird} III., letzter König der}{}
\EdL{91}{AK07-J0277-L91}{Sassaniden, XV, XVI, T. 1.}{}
\EdL{92}{AK07-J0277-L92}{\textit{Jenubi} s. \textit{Genubi}.}{}
\EdL{93}{AK07-J0277-L93}{\textit{Jumedi elachir} = \textit{Djumādā al-ākhira}, \textit{Djumādā}}{}
\EdL{94}{AK07-J0277-L94}{\textit{II.}, der sechste Monat des muhammed.}{}
\EdL{95}{AK07-J0277-L95}{Jahres, XIV, XV, T. 2, 3, etc.}{}
\EdL{96}{AK07-J0277-L96}{\textit{Jumedi elule} = \textit{Djumādā al-ūlā}, \textit{Djumādā I},}{}
\EdL{97}{AK07-J0277-L97}{der fünfte Monat des muhammed. Jahres,}{}
\EdL{98}{AK07-J0277-L98}{XIV, XV, T. 2, 3, etc.}{}
\FrameHead{0278}{p. 244}{0278}{p. 244}
\ColHead{left column}
\EdL{1}{AK07-J0278-L01}{\textbf{K.}}{}
\EdL{2}{AK07-J0278-L02}{\textit{Kamar} = \textit{ḳamar}, Mond, T. 2\textsuperscript{a}, 6—8, 21—26, etc.}{}
\EdL{3}{AK07-J0278-L03}{\textit{Kamarie} = \textit{ḳamarīya}, fem. v. \textit{ḳamarī}, auf den}{}
\EdL{4}{AK07-J0278-L04}{Mond bezüglich, T. 73, 74.}{}
\EdL{5}{AK07-J0278-L05}{\textit{Kamene}, verdorben aus \textit{Kamarie}.}{}
\EdL{6}{AK07-J0278-L06}{\textit{Karnain} s. \textit{Dulkarnain}.}{}
\EdL{7}{AK07-J0278-L07}{\textit{Kawekib} (auch \textit{Kavvekib}) = \textit{kawākib} (pl. v.}{}
\EdL{8}{AK07-J0278-L08}{\textit{kaukab} od. \textit{kawkab}), die Gestirne, insbe-}{}
\EdL{9}{AK07-J0278-L09}{sondere die sieben Planeten, T. 55 (N.).}{}
\EdL{10}{AK07-J0278-L10}{\textit{Keiwen} = \textit{Kaiwān}, pers. Name für den Pla-}{}
\EdL{11}{AK07-J0278-L11}{neten Saturn, XXII, T. 27—29.}{}
\EdL{12}{AK07-J0278-L12}{\textit{Ketib} = \textit{Kātib}, Schreiber, dann hauptsächlich}{}
\EdL{13}{AK07-J0278-L13}{in Spanien und Nordafrika gebräuchlicher}{}
\EdL{14}{AK07-J0278-L14}{Name für den Planeten Merkur, T. 51, 53.}{}
\EdL{15}{AK07-J0278-L15}{\textit{Kiteb} = \textit{Kitāb}, Buch\,; \textit{Kiteb hea} = \textit{kitāb al-hai’a},}{}
\EdL{16}{AK07-J0278-L16}{Buch der Astronomie, 3 (N. 32)\,; es ist hiermit}{}
\EdL{17}{AK07-J0278-L17}{vielleicht die Astronomie des \textit{Djābir b. Aflaḥ}}{}
\EdL{18}{AK07-J0278-L18}{gemeint, die im Ms. des Escorial den Titel}{}
\EdL{19}{AK07-J0278-L19}{trägt \textit{kitāb al-hai’a}.}{}
\EdL{20}{AK07-J0278-L20}{\textit{Kucuf} u. \textit{Kuzuf} = \textit{kusūf}, Finsternis, Ver-}{}
\EdL{21}{AK07-J0278-L21}{finsterung (v. Sonne und Mond) 27, T. 73, 74.}{}
\EdL{22}{AK07-J0278-L22}{\textit{Kucufet} = \textit{kusūfāt} (pl. v. \textit{kusūf}) s. \textit{Kucuf}.}{}
\EdL{23}{AK07-J0278-L23}{\textit{Kuscus}, wahrscheinl. = \textit{ḳaus}, Bogen, 9\,; s. auch}{}
\EdL{24}{AK07-J0278-L24}{\textit{Cauz}.}{}
\EdL{25}{AK07-J0278-L25}{\textit{Kuszuf} s. \textit{Kucuf}.}{}
\EdL{26}{AK07-J0278-L26}{\textbf{L.}}{}
\EdL{27}{AK07-J0278-L27}{\textit{Li}, mit dem Artikel verbunden \textit{lil}, arab.}{}
\EdL{28}{AK07-J0278-L28}{Praepos., zu, für, wegen, in Bezug auf, be-}{}
\EdL{29}{AK07-J0278-L29}{zeichnet auch einfach den Dativ\,; auch}{}
\EdL{30}{AK07-J0278-L30}{Konjunktion\,: dass, damit. Z.B.\textit{libood elabaad}}{}
\EdL{31}{AK07-J0278-L31}{= \textit{lil-buʿd al-abʿad}, für die weiteste Ent-}{}
\EdL{32}{AK07-J0278-L32}{fernung (\textit{pro spatio longinquo}) T. 73 — \textit{li-}}{}
\EdL{33}{AK07-J0278-L33}{\textit{burug elhamel} = \textit{liburdj al-ḥamal} für das}{}
\EdL{34}{AK07-J0278-L34}{Zeichen des Widders, T. 79 — \textit{lieztekim} =}{}
\EdL{35}{AK07-J0278-L35}{\textit{liyastaḳīm}, dass er bestimme, reguliere . . .}{}
\EdL{36}{AK07-J0278-L36}{XXI, T. 9 — \textit{lilkawekib} = \textit{lil-kawākib}, für}{}
\EdL{37}{AK07-J0278-L37}{die Gestirne, od. in Bezug auf die Gestirne,}{}
\EdL{38}{AK07-J0278-L38}{T. 91 — \textit{lilmanshar} = \textit{lil-manẓar}, in Bezug}{}
\EdL{39}{AK07-J0278-L39}{auf den Anblick (Ort des Erscheinens) 28,}{}
\EdL{40}{AK07-J0278-L40}{s. auch \textit{Manshar} — \textit{liwacat alcuzuf} = \textit{liwa-}}{}
\EdL{41}{AK07-J0278-L41}{\textit{saṭ al-kusūf}, für die Mitte der Finsternis,}{}
\EdL{42}{AK07-J0278-L42}{27, s. auch \textit{Wacat}.}{}
\EdL{43}{AK07-J0278-L43}{\textbf{M.}}{}
\EdL{44}{AK07-J0278-L44}{\textit{Mabsūṭ}, ausgebreitet, (\textit{planus}), 36, 77\,; s. auch}{}
\EdL{45}{AK07-J0278-L45}{\textit{Mobzota}.}{}
\EdL{46}{AK07-J0278-L46}{\textit{Magesti} (auch \textit{Megesti}) = \textit{Madjisṭī}, der Alma-}{}
\EdL{47}{AK07-J0278-L47}{gest (die grosse Syntaxis) des \Sperr{Ptolemäus}}{}
\EdL{48}{AK07-J0278-L48}{2, 18, etc.}{}
\ColHead{right column}
\EdL{49}{AK07-J0278-L49}{\textit{Magmoa} = \textit{madjmūʿ}, Summe, Gesamtheit\,; auch}{}
\EdL{50}{AK07-J0278-L50}{Partizip = gesammelt, (fem. \textit{madjmūʿa}) T.}{}
\EdL{51}{AK07-J0278-L51}{15 (N. 1)\,; \textit{filcinun elmagmoa} = \textit{fī’l-sinīn}}{}
\EdL{52}{AK07-J0278-L52}{\textit{al-madjmūʿa}, in den (zu 30) vereinigten}{}
\EdL{53}{AK07-J0278-L53}{Jahren (\textit{in annis collectis}) T. 69.}{}
\EdL{54}{AK07-J0278-L54}{\textit{Mahwar} = \textit{miḥwar}, Achse, Erdachse, hier Pol}{}
\EdL{55}{AK07-J0278-L55}{derselben, 18.}{}
\EdL{56}{AK07-J0278-L56}{\textit{Maimoa} s. \textit{Magmoa}.}{}
\EdL{57}{AK07-J0278-L57}{\textit{Makuz} = \textit{maʿkūs}, verkehrt, 22 (N. 6), 69\,; s. auch}{}
\EdL{58}{AK07-J0278-L58}{\textit{Mankuz}.}{}
\EdL{59}{AK07-J0278-L59}{\textit{Mandhar} = \textit{manẓar}, Anblick, Ort des Erschei-}{}
\EdL{60}{AK07-J0278-L60}{nens, 28 (N. 3), T. 77\,; s. auch \textit{Ihtilef}.}{}
\EdL{61}{AK07-J0278-L61}{\textit{Mankut} s. \textit{Mankuz}.}{}
\EdL{62}{AK07-J0278-L62}{\textit{Mankuz}, 1) = \textit{mankūs}, verkehrt, 17, 22, 69, 77\,;}{}
\EdL{63}{AK07-J0278-L63}{2) = \textit{manḳūṣ}, vermindert (\textit{diminutus}), 17, 69.}{}
\EdL{64}{AK07-J0278-L64}{\textit{Manshar} s. \textit{Mandhar}.}{}
\EdL{65}{AK07-J0278-L65}{\textit{Marech} = \textit{Mirrīkh}, der Planet Mars, T. 13, 14,}{}
\EdL{66}{AK07-J0278-L66}{39, 41—44.}{}
\EdL{67}{AK07-J0278-L67}{\textit{Mareh} s. \textit{Marech}.}{}
\EdL{68}{AK07-J0278-L68}{\textit{Markez} = \textit{markaz}, Mittelpunkt, Centrum, 48\,;}{}
\EdL{69}{AK07-J0278-L69}{\textit{Tadil elmarkez} = \textit{Taʿdīl al-markaz}, Glei-}{}
\EdL{70}{AK07-J0278-L70}{chung des Centrums, T. 54.}{}
\EdL{71}{AK07-J0278-L71}{\textit{Maruz}, verdorben aus \textit{Mankuz}, s. dieses.}{}
\EdL{72}{AK07-J0278-L72}{\textit{Masīr}, Lauf, Bahn, XXI\,; \textit{al-masīr al-wasaṭ},}{}
\EdL{73}{AK07-J0278-L73}{der mittlere Lauf, die mittlere Bewegung,}{}
\EdL{74}{AK07-J0278-L74}{41.}{}
\EdL{75}{AK07-J0278-L75}{\textit{Matale} = \textit{maṭāliʿ} (pl. v. \textit{maṭlaʿ}), Aufsteigungen}{}
\EdL{76}{AK07-J0278-L76}{der Gestirne (\textit{Ascensiones rectae et obliquae}),}{}
\EdL{77}{AK07-J0278-L77}{in diesem Werke mit \textit{Horoscopi} übersetzt,}{}
\EdL{78}{AK07-J0278-L78}{18 (N. 23), T. 59—59\textsuperscript{b}, 91 (N. 1), 95 (N. 1), 97}{}
\EdL{79}{AK07-J0278-L79}{(N. 4).}{}
\EdL{80}{AK07-J0278-L80}{\textit{Matarech} s. \textit{Matareh}.}{}
\EdL{81}{AK07-J0278-L81}{\textit{Matareh} = \textit{maṭāriḥ} (pl. v. \textit{maṭraḥ}) s. \textit{Matrak}.}{}
\EdL{82}{AK07-J0278-L82}{\textit{Matrak} = \textit{maṭraḥ}, Ort der Projektion\,; \textit{matrak}}{}
\EdL{83}{AK07-J0278-L83}{\textit{elshoa} = \textit{maṭraḥ al-shuʿāʿ}, Ort der Strahlen-}{}
\EdL{84}{AK07-J0278-L84}{werfung, 31 (N. 6), T. 91\,; s. auch \textit{Shoa}.}{}
\EdL{85}{AK07-J0278-L85}{\textit{Maudho} = \textit{maudūʿ}, Ort, Lage\,; \textit{maudho Jauze-}}{}
\EdL{86}{AK07-J0278-L86}{\textit{garhu} = \textit{mauḍūʿ djauzaharhu}, Ort seines}{}
\EdL{87}{AK07-J0278-L87}{(aufsteigenden) Knotens, T. 34.}{}
\EdL{88}{AK07-J0278-L88}{\textit{Medek} = \textit{mudaḳḳ}, fein zerstossen, in ganz}{}
\EdL{89}{AK07-J0278-L89}{kleine Teile geteilt, Minute, 7 (N. 20).}{}
\EdL{90}{AK07-J0278-L90}{\textit{Medte} = \textit{madda} (od. \textit{mudda}), Ausdehnung,}{}
\EdL{91}{AK07-J0278-L91}{Dauer, T. 13 (N. 17).}{}
\EdL{92}{AK07-J0278-L92}{\textit{Megil} s. \textit{Meil}.}{}
\EdL{93}{AK07-J0278-L93}{\textit{Meil} = \textit{mail}, Neigung, Deklination jedes Punk-}{}
\EdL{94}{AK07-J0278-L94}{tes der Ekliptik zum Aequator\,; 13, 14, 58,}{}
\EdL{95}{AK07-J0278-L95}{T. 21 (N. 4)\,; \textit{al-mail al-kullī}, die ganze Nei-}{}
\EdL{96}{AK07-J0278-L96}{gung, d. h. die Schiefe der Ekliptik.}{}
\EdL{97}{AK07-J0278-L97}{\textit{Mezen} = \textit{mīzān}, Wage, das Tierkreiszeichen}{}
\EdL{98}{AK07-J0278-L98}{der Wage, T. 59\textsuperscript{b}.}{}
\EdL{99}{AK07-J0278-L99}{\textit{Mirech}, T. 92, s. Nachtrag 5.}{}
\EdL{100}{AK07-J0278-L100}{\textit{Mirrīkh}, der Planet Mars, 236.}{}
\FrameHead{0279}{p. 245}{0279}{p. 245}
\ColHead{left column}
\EdL{1}{AK07-J0279-L01}{\textit{Mizādj}, Temperament, astrol. Natur eines Ge-}{}
\EdL{2}{AK07-J0279-L02}{stirns, s. Nachtrag 5.}{}
\EdL{3}{AK07-J0279-L03}{\textit{Mobzota} = \textit{mabsūṭa} (fem. v. \textit{mabsūṭ}), ausge-}{}
\EdL{4}{AK07-J0279-L04}{breitet, eben (\textit{planus}), T. 15 (N. 2).}{}
\EdL{5}{AK07-J0279-L05}{\textit{Muadala} u. \textit{Muaddala} = \textit{muʿaddala} (fem. v.}{}
\EdL{6}{AK07-J0279-L06}{\textit{muʿaddal}), ausgeglichen, korrigiert, 13 (N. 2),}{}
\EdL{7}{AK07-J0279-L07}{T. 36, 54.}{}
\EdL{8}{AK07-J0279-L08}{\textit{Muʿāmalāt} (pl. v. \textit{muʿāmala}), kaufmännische}{}
\EdL{9}{AK07-J0279-L09}{oder Geschäftsrechnungen, IX.}{}
\EdL{10}{AK07-J0279-L10}{\textit{Muchre} = \textit{mu’akhkhara}, Abschluss, Ende, hier}{}
\EdL{11}{AK07-J0279-L11}{östliches und westliches Ende der bewohn-}{}
\EdL{12}{AK07-J0279-L12}{ten Erdhälfte, 9.}{}
\EdL{13}{AK07-J0279-L13}{\textit{Muharram} = \textit{Muḥarram}, der erste Monat des}{}
\EdL{14}{AK07-J0279-L14}{muhammed. Jahres, XIV, XV, 1, 2, 34, 37,}{}
\EdL{15}{AK07-J0279-L15}{etc. T. 2, 3, etc.}{}
\EdL{16}{AK07-J0279-L16}{\textit{Muhit} = \textit{muḥīṭ}, Umfang, Peripherie, 18.}{}
\EdL{17}{AK07-J0279-L17}{\textit{Muht} s. \textit{Mukdh}.}{}
\EdL{18}{AK07-J0279-L18}{\textit{Mukaam} = \textit{muḳām}, Aufenthalt, Station (der}{}
\EdL{19}{AK07-J0279-L19}{Planeten), 12, 13 (N. 8, 12, 14), T. 54.}{}
\EdL{20}{AK07-J0279-L20}{\textit{Muḳābala}, Gegenüberstellung, Ausgleichung,}{}
\EdL{21}{AK07-J0279-L21}{VIII\,; s. auch \textit{Djebr}.}{}
\EdL{22}{AK07-J0279-L22}{\textit{Mukatil} = \textit{muḳātil}, Kämpfer, Mörder, Beiname}{}
\EdL{23}{AK07-J0279-L23}{des Planeten Saturn, wahrscheinlich nur}{}
\EdL{24}{AK07-J0279-L24}{in Spanien und Nordafrika gebräuchlich,}{}
\EdL{25}{AK07-J0279-L25}{T. 27—29.}{}
\EdL{26}{AK07-J0279-L26}{\textit{Mukdh} = \textit{mukth} (od. \textit{makth}), Verweilen, Auf-}{}
\EdL{27}{AK07-J0279-L27}{enthalt (\textit{mora}), die Zeit, während welcher}{}
\EdL{28}{AK07-J0279-L28}{eine Finsternis total ist, 26, 27, 87.}{}
\EdL{29}{AK07-J0279-L29}{\textit{Mult} s. \textit{Muhit}.}{}
\EdL{30}{AK07-J0279-L30}{\textit{Munegim} = \textit{munadjdjim}, Astronom, Astrolog,}{}
\EdL{31}{AK07-J0279-L31}{XXII, 20 (N. 3).}{}
\EdL{32}{AK07-J0279-L32}{\textit{Mustakim} = \textit{mustaḳīm}, gerade\,; \textit{elfelek elmu-}}{}
\EdL{33}{AK07-J0279-L33}{\textit{stakim} = \textit{al-falak al-mustaḳīm}, die gerade}{}
\EdL{34}{AK07-J0279-L34}{Sphäre, d. h. die Himmelssphäre, wie sie}{}
\EdL{35}{AK07-J0279-L35}{am Aequator erscheint, 18, 20, 25, etc. T.}{}
\EdL{36}{AK07-J0279-L36}{59—59\textsuperscript{b}.}{}
\EdL{37}{AK07-J0279-L37}{\textit{Musteri} = \textit{Mushtarī}, der Planet Jupiter, T. 11,}{}
\EdL{38}{AK07-J0279-L38}{12, 33—38.}{}
\EdL{39}{AK07-J0279-L39}{\textit{Mustewi} = \textit{mustawī}, eben, gerade, 69, 77\,; fem.}{}
\EdL{40}{AK07-J0279-L40}{\textit{mustewia} = \textit{mustawiya}, 21.}{}
\EdL{41}{AK07-J0279-L41}{\textit{Muth} s. \textit{Mukdh}.}{}
\EdL{42}{AK07-J0279-L42}{\textbf{N.}}{}
\EdL{43}{AK07-J0279-L43}{\textit{Nadair} s. das folg.}{}
\EdL{44}{AK07-J0279-L44}{\textit{Nadir} = \textit{naẓīr}, der dem Zenit (\textit{samt al-ra’s})}{}
\EdL{45}{AK07-J0279-L45}{gegenüberliegende Punkt des Himmels,}{}
\EdL{46}{AK07-J0279-L46}{dann überhaupt der Gegenpunkt irgend}{}
\EdL{47}{AK07-J0279-L47}{eines Punktes am Himmel XX, 97, 98, 101.}{}
\EdL{48}{AK07-J0279-L48}{\textit{Nāhīd} s. \textit{Anāhīd}.}{}
\EdL{49}{AK07-J0279-L49}{\textit{Nuzf} = \textit{nuṣf} u. \textit{niṣf}, Hälfte (\textit{dimidium}),}{}
\EdL{50}{AK07-J0279-L50}{T. 73.}{}
\ColHead{right column}
\EdL{51}{AK07-J0279-L51}{\textbf{O.}}{}
\EdL{52}{AK07-J0279-L52}{\textit{Otari[d]} = \textit{ʿUṭārid}, der Planet Merkur, XXI,}{}
\EdL{53}{AK07-J0279-L53}{T. 17, 18, 51—56.}{}
\EdL{54}{AK07-J0279-L54}{\textit{Ozari} s. das vorhergeh.}{}
\EdL{55}{AK07-J0279-L55}{\textbf{R.}}{}
\EdL{56}{AK07-J0279-L56}{\textit{Rabb}, Herr, Meister, 103.}{}
\EdL{57}{AK07-J0279-L57}{\textit{Raba} s. \textit{Rabia}.}{}
\EdL{58}{AK07-J0279-L58}{\textit{Rabe} = pl. v. \textit{Raba} oder \textit{Rabia} s. dieses.}{}
\EdL{59}{AK07-J0279-L59}{\textit{Rabe elachir} = \textit{Rabīʿ al-ākhir}, der vierte mu-}{}
\EdL{60}{AK07-J0279-L60}{hammed. Monat, XIV, XV, 34, 84, 86, 89, T.}{}
\EdL{61}{AK07-J0279-L61}{2, 3, etc}{}
\EdL{62}{AK07-J0279-L62}{\textit{Rabe elawel} = \textit{Rabīʿ al-awwal}, der dritte mu-}{}
\EdL{63}{AK07-J0279-L63}{hammed. Monat, XIV, XV, 34, T. 2, 3, etc.}{}
\EdL{64}{AK07-J0279-L64}{\textit{Rabia} = \textit{rābiʿ}, fem. \textit{rābiʿa}, (der, die) vierte\,;}{}
\EdL{65}{AK07-J0279-L65}{in der Rechnung mit Sexagesimalbrüchen}{}
\EdL{66}{AK07-J0279-L66}{die \textit{partes quartae}, d. h. die Vielfachen von}{}
\EdL{67}{AK07-J0279-L67}{\Frac{1}{60\textsuperscript{4}}, 15 (N. 26).}{}
\EdL{68}{AK07-J0279-L68}{\textit{Race} = \textit{ra’s}, Kopf\,; wahrscheinl. soll \textit{race} Ab-}{}
\EdL{69}{AK07-J0279-L69}{lativ von \textit{rac} sein\,; \textit{Race elhamel} = \textit{ra’s al-}}{}
\EdL{70}{AK07-J0279-L70}{\textit{ḥamal}, Kopf, Anfang (des Zeichens) des}{}
\EdL{71}{AK07-J0279-L71}{Widders, 20.}{}
\EdL{72}{AK07-J0279-L72}{\textit{Rageb} s. \textit{Regeb}.}{}
\EdL{73}{AK07-J0279-L73}{\textit{Ramadham} = \textit{Ramaḍān}, der neunte muham-}{}
\EdL{74}{AK07-J0279-L74}{med. Monat, der Fastenmonat XIV, XV, 34,}{}
\EdL{75}{AK07-J0279-L75}{37, 67, 68, T. 2, 3, etc.}{}
\EdL{76}{AK07-J0279-L76}{\textit{Raz} s. \textit{Race}.}{}
\EdL{77}{AK07-J0279-L77}{\textit{Regeb} = \textit{Radjab}, der siebente muhammed.}{}
\EdL{78}{AK07-J0279-L78}{Monat XIV, XV, 34, 38, 39, T. 2, 3, etc.}{}
\EdL{79}{AK07-J0279-L79}{\textbf{S.}}{}
\EdL{80}{AK07-J0279-L80}{\textit{Saaben} s. \textit{Shabin}.}{}
\EdL{81}{AK07-J0279-L81}{\textit{Safar} s. \textit{Zafar}.}{}
\EdL{82}{AK07-J0279-L82}{\textit{Ṣāḥib}, Herr, Hausherr, Besitzer, in der Astrol.}{}
\EdL{83}{AK07-J0279-L83}{\textit{Dominus} eines Zeichens, oder eines Teils}{}
\EdL{84}{AK07-J0279-L84}{eines Zeichens, 103.}{}
\EdL{85}{AK07-J0279-L85}{\textit{Sana kabīsa}, Schaltjahr, 34.}{}
\EdL{86}{AK07-J0279-L86}{\textit{Scabin} s. \textit{Shabin}.}{}
\EdL{87}{AK07-J0279-L87}{\textit{Scawal} s. \textit{Shawel}.}{}
\EdL{88}{AK07-J0279-L88}{\textit{Sceir} = \textit{sair}, Gang, Lauf, Bahn, XXI.}{}
\EdL{89}{AK07-J0279-L89}{\textit{Sewil} s. \textit{Shawel}.}{}
\EdL{90}{AK07-J0279-L90}{\textit{Shabin} = \textit{Shaʿbān}, der achte muhammed.}{}
\EdL{91}{AK07-J0279-L91}{Monat XIV, XV, 34, 67, 68, T. 2, 3, etc.}{}
\EdL{92}{AK07-J0279-L92}{\textit{Sharaf} (pl \textit{ashrāf}), Höhe, Erhöhung, Auszeich-}{}
\EdL{93}{AK07-J0279-L93}{nung\,; in der mittelalterl. Astrologie \textit{exal-}}{}
\EdL{94}{AK07-J0279-L94}{\textit{tatio} XXI, 104.}{}
\EdL{95}{AK07-J0279-L95}{\textit{Shawel} (auch \textit{Shavvel}) = \textit{Shawwāl}, der zehnte}{}
\EdL{96}{AK07-J0279-L96}{muhammed. Monat XIV, XV, 34, 39, T. 2, 3,}{}
\EdL{97}{AK07-J0279-L97}{etc.}{}
\FrameHead{0280}{p. 246}{0280}{p. 246}
\ColHead{left column}
\EdL{1}{AK07-J0280-L01}{\textit{Shemeli} = \textit{shamālī}, nördlich, 13, 16, 18, 29}{}
\EdL{2}{AK07-J0280-L02}{(N. 5, 12, 13).}{}
\EdL{3}{AK07-J0280-L03}{\textit{Shemelia} = \textit{shamālīya}, fem. des vorhergeh.}{}
\EdL{4}{AK07-J0280-L04}{29 (N. 5, 13).}{}
\EdL{5}{AK07-J0280-L05}{\textit{Shems} = \textit{Shams}, Sonne, 58, T. 21—26.}{}
\EdL{6}{AK07-J0280-L06}{\textit{Shoa} = \textit{shuʿāʿ} (pl. \textit{shuʿuʿ} und \textit{shuʿāʿāt}), Strahl}{}
\EdL{7}{AK07-J0280-L07}{eines Gestirns, z. B. der Sonne, 31 (N. 6),}{}
\EdL{8}{AK07-J0280-L08}{T. 91\,; \textit{matrak elshoa} s. \textit{Matrak}.}{}
\EdL{9}{AK07-J0280-L09}{\textit{Shohor} = \textit{shuhūr}, pl. v. \textit{shahr}, Monat, T. 9\,;}{}
\EdL{10}{AK07-J0280-L10}{s. auch \textit{filshohor} (unter \textit{fi}).}{}
\EdL{11}{AK07-J0280-L11}{\textit{Ṣifr}, leer, Null, 35.}{}
\EdL{12}{AK07-J0280-L12}{\textit{Sohor} s. \textit{Shohor}.}{}
\EdL{13}{AK07-J0280-L13}{\textit{Ṣūra}, Gestalt, Figur\,; \textit{ṣūrat al-arḍ}, Abbildung}{}
\EdL{14}{AK07-J0280-L14}{(Figur) der Erde, IX.}{}
\EdL{15}{AK07-J0280-L15}{\textit{Szukuf} s. \textit{Zukut}.}{}
\EdL{16}{AK07-J0280-L16}{\textit{Szukut} s. \textit{Zukut}.}{}
\EdL{17}{AK07-J0280-L17}{\textbf{T.}}{}
\EdL{18}{AK07-J0280-L18}{\textit{Ṭabaḳāt} (pl. v. \textit{ṭabaḳa}), Klassen\,; \textit{ṭabaḳāt al-}}{}
\EdL{19}{AK07-J0280-L19}{\textit{umam}, Klassen der Völker, IX\,; s. auch}{}
\EdL{20}{AK07-J0280-L20}{\textit{Umam}.}{}
\EdL{21}{AK07-J0280-L21}{\textit{Tabeb} = \textit{ṭabīb}, Gelehrter, Arzt XXII, 20.}{}
\EdL{22}{AK07-J0280-L22}{\textit{Tadil} = \textit{taʿdīl}, Gleichung (der Sonne, des Mon-}{}
\EdL{23}{AK07-J0280-L23}{des und der Planeten), wird im lat. Text}{}
\EdL{24}{AK07-J0280-L24}{durch \textit{aequatio}, \textit{coaequatio} und \textit{examinatio}}{}
\EdL{25}{AK07-J0280-L25}{wiedergegeben, 9, 10, 46 etc. T. 21—56.}{}
\EdL{26}{AK07-J0280-L26}{\textit{Tarbea} = \textit{tarbīʿ}, Geviert- oder Quadratschein}{}
\EdL{27}{AK07-J0280-L27}{(Astrol.) T. 91.}{}
\EdL{28}{AK07-J0280-L28}{\textit{Ta’rīkh}, Zeitrechnung, Chronik, Datum, IX\,;}{}
\EdL{29}{AK07-J0280-L29}{\textit{Ta’rīkh al-ṣufr}, die spanische Zeitrechnung,}{}
\EdL{30}{AK07-J0280-L30}{35\,; s. auch \textit{Acofra}.}{}
\EdL{31}{AK07-J0280-L31}{\textit{Tedewir} = \textit{tadwīr}, Umdrehung (in einem}{}
\EdL{32}{AK07-J0280-L32}{Kreise)\,; \textit{felek eltedewir} = \textit{falak al-tadwīr},}{}
\EdL{33}{AK07-J0280-L33}{der Epizykel, T. 9 (N. 20).}{}
\EdL{34}{AK07-J0280-L34}{\textit{Tethlith} = \textit{tathlīth}, Gedritt- oder Trigonal-}{}
\EdL{35}{AK07-J0280-L35}{schein (Astrol.) T. 91.}{}
\EdL{36}{AK07-J0280-L36}{\textit{Tezdiz} = \textit{tasdīs}, Gesechst- oder Hexagonal-}{}
\EdL{37}{AK07-J0280-L37}{schein (Astrol.) T. 91.}{}
\EdL{38}{AK07-J0280-L38}{\textit{Tezweit} s. \textit{Tezwiet}.}{}
\EdL{39}{AK07-J0280-L39}{\textit{Tezwiet} = \textit{taswiya}, Gleichstellung, Anordnung,}{}
\EdL{40}{AK07-J0280-L40}{Verteilung, 25, 96, T. 79.}{}
\EdL{41}{AK07-J0280-L41}{\textit{Thaur}, Stier, das Tierkreiszeichen des Stiers}{}
\EdL{42}{AK07-J0280-L42}{(\textit{Taurus}) T. 59\textsuperscript{a}.}{}
\EdL{43}{AK07-J0280-L43}{\textit{Thedewir} s. \textit{Tedewir}.}{}
\EdL{44}{AK07-J0280-L44}{\textit{Thelide} u. \textit{Thelidhe} = pl. v. \textit{Thelidha}, s. d.}{}
\EdL{45}{AK07-J0280-L45}{folg.}{}
\EdL{46}{AK07-J0280-L46}{\textit{Thelitha} = \textit{thālitha}, fem. v. \textit{thālith}, (die) dritte,}{}
\EdL{47}{AK07-J0280-L47}{in der Rechnung mit Sexagesimalbrüchen}{}
\EdL{48}{AK07-J0280-L48}{die \textit{partes tertiae}, d. h. die Vielfachen von}{}
\EdL{49}{AK07-J0280-L49}{\Frac{1}{60\textsuperscript{3}}, 15 (N.).}{}
\EdL{50}{AK07-J0280-L50}{\textit{Theni} = \textit{thānī}, (der) zweite, 12.}{}
\ColHead{right column}
\EdL{51}{AK07-J0280-L51}{\textit{Thenia} = \textit{thāniya}, fem. des vorhergeh., Se-}{}
\EdL{52}{AK07-J0280-L52}{kunde, in der Rechnung mit Sexagesimal-}{}
\EdL{53}{AK07-J0280-L53}{brüchen die \textit{partes secundae}, d. h. die Viel-}{}
\EdL{54}{AK07-J0280-L54}{fachen von \Frac{1}{60\textsuperscript{2}}, XXIII, XXV, 11, 14, 15, 36,}{}
\EdL{55}{AK07-J0280-L55}{69. T. 4. 5, etc.}{}
\EdL{56}{AK07-J0280-L56}{\textit{Thuid} s. \textit{Thuld}.}{}
\EdL{57}{AK07-J0280-L57}{\textit{Thuld} = \textit{thulth}, ein Drittel, 9.}{}
\EdL{58}{AK07-J0280-L58}{\textit{Thur} s. \textit{Thaur}.}{}
\EdL{59}{AK07-J0280-L59}{\textbf{U.}}{}
\EdL{60}{AK07-J0280-L60}{\textit{Udjain} s. \textit{Arin}.}{}
\EdL{61}{AK07-J0280-L61}{\textit{Ule} = \textit{ūlā}, fem. v. \textit{awwal}, (die) erste, T. 2, 3,}{}
\EdL{62}{AK07-J0280-L62}{etc.}{}
\EdL{63}{AK07-J0280-L63}{\textit{Umam} (pl. v. \textit{umma}), Völker, Nationen.}{}
\EdL{64}{AK07-J0280-L64}{\textit{Uzain}, arab. Form für das ind. \textit{Udjain}, s. \textit{Arin}.}{}
\EdL{65}{AK07-J0280-L65}{\textbf{W.}}{}
\EdL{66}{AK07-J0280-L66}{\textit{Wa}, arab. Praepos., und\,; z. B. \textit{wakhod} =}{}
\EdL{67}{AK07-J0280-L67}{\textit{wakhudh}, und nimm XXII, 27 (N. 18)\,; \textit{wa-}}{}
\EdL{68}{AK07-J0280-L68}{\textit{hezatu} = \textit{wakhāṣṣathu}, und seine Anomalie}{}
\EdL{69}{AK07-J0280-L69}{T. 6—8, s. auch \textit{Heza}.}{}
\EdL{70}{AK07-J0280-L70}{\textit{Wabāl}, Schaden, 161.}{}
\EdL{71}{AK07-J0280-L71}{\textit{Wacat} = \textit{wasaṭ} (od. \textit{wasṭ}), Mitte, Mittel, mitt-}{}
\EdL{72}{AK07-J0280-L72}{lere Bewegung (der Sonne, des Mondes und}{}
\EdL{73}{AK07-J0280-L73}{der Planeten), im Text und in den Tafeln}{}
\EdL{74}{AK07-J0280-L74}{meistens durch \textit{medialitas} wiedergegeben,}{}
\EdL{75}{AK07-J0280-L75}{7—12 etc., 41, T. 4—20, etc. \textit{wacat alkamar}}{}
\EdL{76}{AK07-J0280-L76}{= \textit{wasaṭ al-ḳamar}, Mittellauf (mittlere Be-}{}
\EdL{77}{AK07-J0280-L77}{wegung) des Mondes, T. 6—8.}{}
\EdL{78}{AK07-J0280-L78}{\textit{Wadjh}, pl. \textit{wudjūh}, vordere Seite, Antlitz, Art}{}
\EdL{79}{AK07-J0280-L79}{und Weise, etc.\,; in der Astrologie der dritte}{}
\EdL{80}{AK07-J0280-L80}{Teil (10°) eines Tierkreiszeichens (lat. \textit{de-}}{}
\EdL{81}{AK07-J0280-L81}{\textit{canus}, auch \textit{facies}) 104.}{}
\EdL{82}{AK07-J0280-L82}{\textit{Wakhod mebibelche} XXII, 27 (N. 18).}{}
\EdL{83}{AK07-J0280-L83}{\textit{Wazat} s. \textit{Wacat}.}{}
\EdL{84}{AK07-J0280-L84}{\textbf{Y.}}{}
\EdL{85}{AK07-J0280-L85}{\textit{Yle} s. \textit{Ule}.}{}
\EdL{86}{AK07-J0280-L86}{\textbf{Z.}}{}
\EdL{87}{AK07-J0280-L87}{\textit{Zafar} = \textit{Ṣafar}, der zweite muhammed. Monat,}{}
\EdL{88}{AK07-J0280-L88}{XIV, XV, 2, 34, 42, 43, 45 etc. T. 2, 3, 4 etc.}{}
\EdL{89}{AK07-J0280-L89}{\textit{Zahat} s. \textit{Caat}.}{}
\EdL{90}{AK07-J0280-L90}{\textit{Zare} s. \textit{Care}.}{}
\EdL{91}{AK07-J0280-L91}{\textit{Zenia} s. \textit{Thenia}.}{}
\EdL{92}{AK07-J0280-L92}{\textit{Zich} = \textit{zīdj} (pers. \textit{zīg}) astronomisches Tafel-}{}
\EdL{93}{AK07-J0280-L93}{werk, VIII, 1, 32, etc. T. 116\,; \textit{zīdj al-shāh}}{}
\EdL{94}{AK07-J0280-L94}{oder \textit{al-shahriyār}, die Tafeln des \Sperr{Shāh}}{}
\EdL{95}{AK07-J0280-L95}{(des Königs) = die persischen Tafeln, 32,}{}
\FrameHead{0281}{p. 247}{0281}{p. 247}
\ColHead{left column}
\EdL{1}{AK07-J0281-L01}{33, 52, etc.\,; s. auch Namen- und Sachregister}{}
\EdL{2}{AK07-J0281-L02}{unter „Tafeln“.}{}
\EdL{3}{AK07-J0281-L03}{\textit{Zīg} s. \textit{Zich}.}{}
\EdL{4}{AK07-J0281-L04}{\textit{Ẓill}, Schatten, die trigon. Funktion der \textit{Tangens}}{}
\EdL{5}{AK07-J0281-L05}{und \textit{Cotangens} 77\,; s. auch \textit{Dhel}.}{}
\EdL{6}{AK07-J0281-L06}{\textit{Zohal} = \textit{Zuḥal}, der Planet Saturn, T. 9, 10,}{}
\EdL{7}{AK07-J0281-L07}{27—32.}{}
\EdL{8}{AK07-J0281-L08}{\textit{Zohara} = \textit{Zuhara}, der Planet Venus, XXII,}{}
\EdL{9}{AK07-J0281-L09}{T. 45, 47, 49, 50.}{}
\ColHead{right column}
\EdL{10}{AK07-J0281-L10}{\textit{Zohari} = \textit{zuharī}, zur Venus gehörig, T. 15,}{}
\EdL{11}{AK07-J0281-L11}{16.}{}
\EdL{12}{AK07-J0281-L12}{\textit{Zukut} = \textit{suḳūṭ}, (das) Fallen, Hineinfallen\,; in}{}
\EdL{13}{AK07-J0281-L13}{der Astron. die Phase einer Finsternis vom}{}
\EdL{14}{AK07-J0281-L14}{Beginn bis zur Totalität, in diesem Werke}{}
\EdL{15}{AK07-J0281-L15}{durch \textit{casus} wiedergegeben, sonst durch}{}
\EdL{16}{AK07-J0281-L16}{\textit{immersio}, 26, 27, 87, T. 73, 78.}{}
\EdL{17}{AK07-J0281-L17}{\textit{Zukuz} s. \textit{Zukut}.}{}
\ColHead{full width}
\EdLc{18}{AK07-J0281-L18}{II. Ägyptische (koptische) Monatsnamen (Tab. 2\textsuperscript{b}).}{}
\EdL{19}{AK07-J0281-L19}{(Die koptischen Monatsnamen sind in diesen Tafeln dem oberägyptischen (thebanischen)}{}
\EdL{20}{AK07-J0281-L20}{Dialekt entnommen, in Klammern habe ich auch die Formen des unterägyptischen (memphi-}{}
\EdLc{21}{AK07-J0281-L21}{tischen) Dialektes beigefügt).}{}
\EdL{22}{AK07-J0281-L22}{\parbox[t]{\dimexpr0.0540\linewidth-0.5em\relax}{\raggedright\small\mbox{}1.}\hspace{0.5em}\parbox[t]{\dimexpr0.5580\linewidth-0.5em\relax}{\raggedright\small\mbox{}Tut = thowt (thōwt)}\hspace{0.5em}\parbox[t]{\dimexpr0.1170\linewidth-0.5em\relax}{\raggedright\small\mbox{}griech.}\hspace{0.5em}\parbox[t]{\dimexpr0.1710\linewidth-0.5em\relax}{\raggedright\small\mbox{}θώθ}\hspace{0.5em}}{}
\EdL{23}{AK07-J0281-L23}{\parbox[t]{\dimexpr0.0540\linewidth-0.5em\relax}{\raggedright\small\mbox{}2.}\hspace{0.5em}\parbox[t]{\dimexpr0.5580\linewidth-0.5em\relax}{\raggedright\small\mbox{}Beba = paape (paopi)}\hspace{0.5em}\parbox[t]{\dimexpr0.1170\linewidth-0.5em\relax}{\raggedright\small\mbox{}—}\hspace{0.5em}\parbox[t]{\dimexpr0.1710\linewidth-0.5em\relax}{\raggedright\small\mbox{}φαωφί}\hspace{0.5em}}{}
\EdL{24}{AK07-J0281-L24}{\parbox[t]{\dimexpr0.0540\linewidth-0.5em\relax}{\raggedright\small\mbox{}3.}\hspace{0.5em}\parbox[t]{\dimexpr0.5580\linewidth-0.5em\relax}{\raggedright\small\mbox{}Hettur = hatōr (athōr)}\hspace{0.5em}\parbox[t]{\dimexpr0.1170\linewidth-0.5em\relax}{\raggedright\small\mbox{}—}\hspace{0.5em}\parbox[t]{\dimexpr0.1710\linewidth-0.5em\relax}{\raggedright\small\mbox{}ἀθύρ}\hspace{0.5em}}{}
\EdL{25}{AK07-J0281-L25}{\parbox[t]{\dimexpr0.0540\linewidth-0.5em\relax}{\raggedright\small\mbox{}4.}\hspace{0.5em}\parbox[t]{\dimexpr0.5580\linewidth-0.5em\relax}{\raggedright\small\mbox{}Keihec = kiakh (khoiak)}\hspace{0.5em}\parbox[t]{\dimexpr0.1170\linewidth-0.5em\relax}{\raggedright\small\mbox{}—}\hspace{0.5em}\parbox[t]{\dimexpr0.1710\linewidth-0.5em\relax}{\raggedright\small\mbox{}χοιάκ}\hspace{0.5em}}{}
\EdL{26}{AK07-J0281-L26}{\parbox[t]{\dimexpr0.0540\linewidth-0.5em\relax}{\raggedright\small\mbox{}5.}\hspace{0.5em}\parbox[t]{\dimexpr0.5580\linewidth-0.5em\relax}{\raggedright\small\mbox{}Tuba = tōbe (tōbi)}\hspace{0.5em}\parbox[t]{\dimexpr0.1170\linewidth-0.5em\relax}{\raggedright\small\mbox{}—}\hspace{0.5em}\parbox[t]{\dimexpr0.1710\linewidth-0.5em\relax}{\raggedright\small\mbox{}τυβί}\hspace{0.5em}}{}
\EdL{27}{AK07-J0281-L27}{\parbox[t]{\dimexpr0.0540\linewidth-0.5em\relax}{\raggedright\small\mbox{}6.}\hspace{0.5em}\parbox[t]{\dimexpr0.5580\linewidth-0.5em\relax}{\raggedright\small\mbox{}Amushir = \textsuperscript{e}mshir (mekhir)}\hspace{0.5em}\parbox[t]{\dimexpr0.1170\linewidth-0.5em\relax}{\raggedright\small\mbox{}—}\hspace{0.5em}\parbox[t]{\dimexpr0.1710\linewidth-0.5em\relax}{\raggedright\small\mbox{}μεχίρ}\hspace{0.5em}}{}
\EdL{28}{AK07-J0281-L28}{\parbox[t]{\dimexpr0.0540\linewidth-0.5em\relax}{\raggedright\small\mbox{}7.}\hspace{0.5em}\parbox[t]{\dimexpr0.5580\linewidth-0.5em\relax}{\raggedright\small\mbox{}Burimihez = parmhāt (phamenōth)}\hspace{0.5em}\parbox[t]{\dimexpr0.1170\linewidth-0.5em\relax}{\raggedright\small\mbox{}—}\hspace{0.5em}\parbox[t]{\dimexpr0.1710\linewidth-0.5em\relax}{\raggedright\small\mbox{}φαμενώθ}\hspace{0.5em}}{}
\EdL{29}{AK07-J0281-L29}{\parbox[t]{\dimexpr0.0540\linewidth-0.5em\relax}{\raggedright\small\mbox{}8.}\hspace{0.5em}\parbox[t]{\dimexpr0.5580\linewidth-0.5em\relax}{\raggedright\small\mbox{}Burmuda = parmūte (pharmūthi)}\hspace{0.5em}\parbox[t]{\dimexpr0.1170\linewidth-0.5em\relax}{\raggedright\small\mbox{}—}\hspace{0.5em}\parbox[t]{\dimexpr0.1710\linewidth-0.5em\relax}{\raggedright\small\mbox{}φαρμουθί}\hspace{0.5em}}{}
\EdL{30}{AK07-J0281-L30}{\parbox[t]{\dimexpr0.0540\linewidth-0.5em\relax}{\raggedright\small\mbox{}9.}\hspace{0.5em}\parbox[t]{\dimexpr0.5580\linewidth-0.5em\relax}{\raggedright\small\mbox{}Besenz = pashons (pakhōn)}\hspace{0.5em}\parbox[t]{\dimexpr0.1170\linewidth-0.5em\relax}{\raggedright\small\mbox{}—}\hspace{0.5em}\parbox[t]{\dimexpr0.1710\linewidth-0.5em\relax}{\raggedright\small\mbox{}παχών}\hspace{0.5em}}{}
\EdL{31}{AK07-J0281-L31}{\parbox[t]{\dimexpr0.0540\linewidth-0.5em\relax}{\raggedright\small\mbox{}10.}\hspace{0.5em}\parbox[t]{\dimexpr0.5580\linewidth-0.5em\relax}{\raggedright\small\mbox{}Buba = paōne (paōni)}\hspace{0.5em}\parbox[t]{\dimexpr0.1170\linewidth-0.5em\relax}{\raggedright\small\mbox{}—}\hspace{0.5em}\parbox[t]{\dimexpr0.1710\linewidth-0.5em\relax}{\raggedright\small\mbox{}παϋνί}\hspace{0.5em}}{}
\EdL{32}{AK07-J0281-L32}{\parbox[t]{\dimexpr0.0540\linewidth-0.5em\relax}{\raggedright\small\mbox{}11.}\hspace{0.5em}\parbox[t]{\dimexpr0.5580\linewidth-0.5em\relax}{\raggedright\small\mbox{}Abib = epep (epēp)}\hspace{0.5em}\parbox[t]{\dimexpr0.1170\linewidth-0.5em\relax}{\raggedright\small\mbox{}—}\hspace{0.5em}\parbox[t]{\dimexpr0.1710\linewidth-0.5em\relax}{\raggedright\small\mbox{}ἐπιφί}\hspace{0.5em}}{}
\EdL{33}{AK07-J0281-L33}{\parbox[t]{\dimexpr0.0540\linewidth-0.5em\relax}{\raggedright\small\mbox{}12.}\hspace{0.5em}\parbox[t]{\dimexpr0.5580\linewidth-0.5em\relax}{\raggedright\small\mbox{}Muzre = mesōrē (mesōrē)}\hspace{0.5em}\parbox[t]{\dimexpr0.1170\linewidth-0.5em\relax}{\raggedright\small\mbox{}—}\hspace{0.5em}\parbox[t]{\dimexpr0.1710\linewidth-0.5em\relax}{\raggedright\small\mbox{}μεσορή}\hspace{0.5em}}{}
\EdLc{34}{AK07-J0281-L34}{III. Persische Monatsnamen (Tab. 2\textsuperscript{a}).}{}
\EdLc{35}{AK07-J0281-L35}{(\textit{māh} = Monat)}{}
\EdL{36}{AK07-J0281-L36}{\parbox[t]{\dimexpr0.0581\linewidth-0.5em\relax}{\raggedright\small\mbox{}1.}\hspace{0.5em}\parbox[t]{\dimexpr0.2228\linewidth-0.5em\relax}{\raggedright\small\mbox{}Affrordinme}\hspace{0.5em}\parbox[t]{\dimexpr0.3391\linewidth-0.5em\relax}{\raggedright\small\mbox{}= Farwardīn-māh.}\hspace{0.5em}}{}
\EdL{37}{AK07-J0281-L37}{\parbox[t]{\dimexpr0.0581\linewidth-0.5em\relax}{\raggedright\small\mbox{}2.}\hspace{0.5em}\parbox[t]{\dimexpr0.2228\linewidth-0.5em\relax}{\raggedright\small\mbox{}Azdihestme}\hspace{0.5em}\parbox[t]{\dimexpr0.3391\linewidth-0.5em\relax}{\raggedright\small\mbox{}= Ardībihist-māh.}\hspace{0.5em}}{}
\EdL{38}{AK07-J0281-L38}{\parbox[t]{\dimexpr0.0581\linewidth-0.5em\relax}{\raggedright\small\mbox{}3.}\hspace{0.5em}\parbox[t]{\dimexpr0.2228\linewidth-0.5em\relax}{\raggedright\small\mbox{}Chordezme}\hspace{0.5em}\parbox[t]{\dimexpr0.3391\linewidth-0.5em\relax}{\raggedright\small\mbox{}= Khurdādh-māh.}\hspace{0.5em}}{}
\EdL{39}{AK07-J0281-L39}{\parbox[t]{\dimexpr0.0581\linewidth-0.5em\relax}{\raggedright\small\mbox{}4.}\hspace{0.5em}\parbox[t]{\dimexpr0.2228\linewidth-0.5em\relax}{\raggedright\small\mbox{}Tiierme}\hspace{0.5em}\parbox[t]{\dimexpr0.3391\linewidth-0.5em\relax}{\raggedright\small\mbox{}= Tīr-māh.}\hspace{0.5em}}{}
\EdL{40}{AK07-J0281-L40}{\parbox[t]{\dimexpr0.0581\linewidth-0.5em\relax}{\raggedright\small\mbox{}5.}\hspace{0.5em}\parbox[t]{\dimexpr0.2228\linewidth-0.5em\relax}{\raggedright\small\mbox{}Merdezme}\hspace{0.5em}\parbox[t]{\dimexpr0.3391\linewidth-0.5em\relax}{\raggedright\small\mbox{}= Murdādh-māh.}\hspace{0.5em}}{}
\EdL{41}{AK07-J0281-L41}{\parbox[t]{\dimexpr0.0581\linewidth-0.5em\relax}{\raggedright\small\mbox{}6.}\hspace{0.5em}\parbox[t]{\dimexpr0.2228\linewidth-0.5em\relax}{\raggedright\small\mbox{}Shaharirme}\hspace{0.5em}\parbox[t]{\dimexpr0.3391\linewidth-0.5em\relax}{\raggedright\small\mbox{}= Shahrīwar-māh.}\hspace{0.5em}}{}
\EdL{42}{AK07-J0281-L42}{\parbox[t]{\dimexpr0.0581\linewidth-0.5em\relax}{\raggedright\small\mbox{}7.}\hspace{0.5em}\parbox[t]{\dimexpr0.2228\linewidth-0.5em\relax}{\raggedright\small\mbox{}Meherime}\hspace{0.5em}\parbox[t]{\dimexpr0.3391\linewidth-0.5em\relax}{\raggedright\small\mbox{}= Mihr-māh.}\hspace{0.5em}}{}
\EdL{43}{AK07-J0281-L43}{\parbox[t]{\dimexpr0.0581\linewidth-0.5em\relax}{\raggedright\small\mbox{}8.}\hspace{0.5em}\parbox[t]{\dimexpr0.2228\linewidth-0.5em\relax}{\raggedright\small\mbox{}Abenme}\hspace{0.5em}\parbox[t]{\dimexpr0.3391\linewidth-0.5em\relax}{\raggedright\small\mbox{}= Ābān-māh.}\hspace{0.5em}}{}
\EdL{44}{AK07-J0281-L44}{\parbox[t]{\dimexpr0.0581\linewidth-0.5em\relax}{\raggedright\small\mbox{}9.}\hspace{0.5em}\parbox[t]{\dimexpr0.2228\linewidth-0.5em\relax}{\raggedright\small\mbox{}Adarme}\hspace{0.5em}\parbox[t]{\dimexpr0.3391\linewidth-0.5em\relax}{\raggedright\small\mbox{}= Ādhar-māh.}\hspace{0.5em}}{}
\EdL{45}{AK07-J0281-L45}{\parbox[t]{\dimexpr0.0581\linewidth-0.5em\relax}{\raggedright\small\mbox{}10.}\hspace{0.5em}\parbox[t]{\dimexpr0.2228\linewidth-0.5em\relax}{\raggedright\small\mbox{}Dime}\hspace{0.5em}\parbox[t]{\dimexpr0.3391\linewidth-0.5em\relax}{\raggedright\small\mbox{}= Dai-māh.}\hspace{0.5em}}{}
\EdL{46}{AK07-J0281-L46}{\parbox[t]{\dimexpr0.0581\linewidth-0.5em\relax}{\raggedright\small\mbox{}11.}\hspace{0.5em}\parbox[t]{\dimexpr0.2228\linewidth-0.5em\relax}{\raggedright\small\mbox{}Bahamanme}\hspace{0.5em}\parbox[t]{\dimexpr0.3391\linewidth-0.5em\relax}{\raggedright\small\mbox{}= Bahman-māh.}\hspace{0.5em}}{}
\EdL{47}{AK07-J0281-L47}{\parbox[t]{\dimexpr0.0581\linewidth-0.5em\relax}{\raggedright\small\mbox{}12.}\hspace{0.5em}\parbox[t]{\dimexpr0.2228\linewidth-0.5em\relax}{\raggedright\small\mbox{}Izfindarme}\hspace{0.5em}\parbox[t]{\dimexpr0.3391\linewidth-0.5em\relax}{\raggedright\small\mbox{}= Isfandārmadh-māh.}\hspace{0.5em}}{}
\FrameHead{0282}{p. 248}{0282}{p. 248}
\EdLc{1}{AK07-J0282-L01}{IV. Syrische (Seleucid.) Monatsnamen (Tab. 3\textsuperscript{a}).}{}
\EdLc{2}{AK07-J0282-L02}{(\textit{Nomina mensium Romanorum}).}{}
\EdL{3}{AK07-J0282-L03}{\mbox{\parbox[t]{\dimexpr0.0508\linewidth-0.5em\relax}{\raggedright\small\mbox{}1.}\hspace{0.5em}\parbox[t]{\dimexpr0.2203\linewidth-0.5em\relax}{\raggedright\small\mbox{}Tisrin elawel}\hspace{0.5em}\parbox[t]{\dimexpr0.2966\linewidth-0.5em\relax}{\raggedright\small\mbox{}= Teshrīn q\textsuperscript{e}dīm (I.)}\hspace{0.5em}\parbox[t]{\dimexpr0.1949\linewidth-0.5em\relax}{\raggedright\small\mbox{}= October}\hspace{0.5em}\parbox[t]{\dimexpr0.2373\linewidth-0.5em\relax}{\raggedright\small\mbox{}(des julian. Jahres).}\hspace{0.5em}}}{}
\EdL{4}{AK07-J0282-L04}{\mbox{\parbox[t]{\dimexpr0.0508\linewidth-0.5em\relax}{\raggedright\small\mbox{}2.}\hspace{0.5em}\parbox[t]{\dimexpr0.2203\linewidth-0.5em\relax}{\raggedright\small\mbox{}Tisrin elachir}\hspace{0.5em}\parbox[t]{\dimexpr0.2966\linewidth-0.5em\relax}{\raggedright\small\mbox{}= Teshrīn ḥ\textsuperscript{e}rai (II.)}\hspace{0.5em}\parbox[t]{\dimexpr0.1949\linewidth-0.5em\relax}{\raggedright\small\mbox{}= November}\hspace{0.5em}\parbox[t]{\dimexpr0.2373\linewidth-0.5em\relax}{\raggedright\small\mbox{}( — — ).}\hspace{0.5em}}}{}
\EdL{5}{AK07-J0282-L05}{\mbox{\parbox[t]{\dimexpr0.0508\linewidth-0.5em\relax}{\raggedright\small\mbox{}3.}\hspace{0.5em}\parbox[t]{\dimexpr0.2203\linewidth-0.5em\relax}{\raggedright\small\mbox{}Keinun elawel}\hspace{0.5em}\parbox[t]{\dimexpr0.2966\linewidth-0.5em\relax}{\raggedright\small\mbox{}= Kānūn q\textsuperscript{e}dīm (I.)}\hspace{0.5em}\parbox[t]{\dimexpr0.1949\linewidth-0.5em\relax}{\raggedright\small\mbox{}= December}\hspace{0.5em}\parbox[t]{\dimexpr0.2373\linewidth-0.5em\relax}{\raggedright\small\mbox{}( — — ).}\hspace{0.5em}}}{}
\EdL{6}{AK07-J0282-L06}{\mbox{\parbox[t]{\dimexpr0.0508\linewidth-0.5em\relax}{\raggedright\small\mbox{}4.}\hspace{0.5em}\parbox[t]{\dimexpr0.2203\linewidth-0.5em\relax}{\raggedright\small\mbox{}Keinun elachir}\hspace{0.5em}\parbox[t]{\dimexpr0.2966\linewidth-0.5em\relax}{\raggedright\small\mbox{}= Kānūn ḥ\textsuperscript{e}rai (II.)}\hspace{0.5em}\parbox[t]{\dimexpr0.1949\linewidth-0.5em\relax}{\raggedright\small\mbox{}= Januarius}\hspace{0.5em}\parbox[t]{\dimexpr0.2373\linewidth-0.5em\relax}{\raggedright\small\mbox{}( — — ).}\hspace{0.5em}}}{}
\EdL{7}{AK07-J0282-L07}{\mbox{\parbox[t]{\dimexpr0.0508\linewidth-0.5em\relax}{\raggedright\small\mbox{}5.}\hspace{0.5em}\parbox[t]{\dimexpr0.2203\linewidth-0.5em\relax}{\raggedright\small\mbox{}Sabat}\hspace{0.5em}\parbox[t]{\dimexpr0.2966\linewidth-0.5em\relax}{\raggedright\small\mbox{}= Sh\textsuperscript{e}bāṭ}\hspace{0.5em}\parbox[t]{\dimexpr0.1949\linewidth-0.5em\relax}{\raggedright\small\mbox{}= Februarius}\hspace{0.5em}\parbox[t]{\dimexpr0.2373\linewidth-0.5em\relax}{\raggedright\small\mbox{}( — — ).}\hspace{0.5em}}}{}
\EdL{8}{AK07-J0282-L08}{\mbox{\parbox[t]{\dimexpr0.0508\linewidth-0.5em\relax}{\raggedright\small\mbox{}6.}\hspace{0.5em}\parbox[t]{\dimexpr0.2203\linewidth-0.5em\relax}{\raggedright\small\mbox{}Adar}\hspace{0.5em}\parbox[t]{\dimexpr0.2966\linewidth-0.5em\relax}{\raggedright\small\mbox{}= Ādār}\hspace{0.5em}\parbox[t]{\dimexpr0.1949\linewidth-0.5em\relax}{\raggedright\small\mbox{}= Martius}\hspace{0.5em}\parbox[t]{\dimexpr0.2373\linewidth-0.5em\relax}{\raggedright\small\mbox{}( — — ).}\hspace{0.5em}}}{}
\EdL{9}{AK07-J0282-L09}{\mbox{\parbox[t]{\dimexpr0.0508\linewidth-0.5em\relax}{\raggedright\small\mbox{}7.}\hspace{0.5em}\parbox[t]{\dimexpr0.2203\linewidth-0.5em\relax}{\raggedright\small\mbox{}Nizan}\hspace{0.5em}\parbox[t]{\dimexpr0.2966\linewidth-0.5em\relax}{\raggedright\small\mbox{}= Nīsān}\hspace{0.5em}\parbox[t]{\dimexpr0.1949\linewidth-0.5em\relax}{\raggedright\small\mbox{}= Aprilis}\hspace{0.5em}\parbox[t]{\dimexpr0.2373\linewidth-0.5em\relax}{\raggedright\small\mbox{}( — — ).}\hspace{0.5em}}}{}
\EdL{10}{AK07-J0282-L10}{\mbox{\parbox[t]{\dimexpr0.0508\linewidth-0.5em\relax}{\raggedright\small\mbox{}8.}\hspace{0.5em}\parbox[t]{\dimexpr0.2203\linewidth-0.5em\relax}{\raggedright\small\mbox{}Iiar (Iyar)}\hspace{0.5em}\parbox[t]{\dimexpr0.2966\linewidth-0.5em\relax}{\raggedright\small\mbox{}= Iyār}\hspace{0.5em}\parbox[t]{\dimexpr0.1949\linewidth-0.5em\relax}{\raggedright\small\mbox{}= Maius}\hspace{0.5em}\parbox[t]{\dimexpr0.2373\linewidth-0.5em\relax}{\raggedright\small\mbox{}( — — ).}\hspace{0.5em}}}{}
\EdL{11}{AK07-J0282-L11}{\mbox{\parbox[t]{\dimexpr0.0508\linewidth-0.5em\relax}{\raggedright\small\mbox{}9.}\hspace{0.5em}\parbox[t]{\dimexpr0.2203\linewidth-0.5em\relax}{\raggedright\small\mbox{}Heziren}\hspace{0.5em}\parbox[t]{\dimexpr0.2966\linewidth-0.5em\relax}{\raggedright\small\mbox{}= Ḥ\textsuperscript{e}zīrān}\hspace{0.5em}\parbox[t]{\dimexpr0.1949\linewidth-0.5em\relax}{\raggedright\small\mbox{}= Junius}\hspace{0.5em}\parbox[t]{\dimexpr0.2373\linewidth-0.5em\relax}{\raggedright\small\mbox{}( — — ).}\hspace{0.5em}}}{}
\EdL{12}{AK07-J0282-L12}{\mbox{\parbox[t]{\dimexpr0.0508\linewidth-0.5em\relax}{\raggedright\small\mbox{}10.}\hspace{0.5em}\parbox[t]{\dimexpr0.2203\linewidth-0.5em\relax}{\raggedright\small\mbox{}Temuz}\hspace{0.5em}\parbox[t]{\dimexpr0.2966\linewidth-0.5em\relax}{\raggedright\small\mbox{}= Tāmūz}\hspace{0.5em}\parbox[t]{\dimexpr0.1949\linewidth-0.5em\relax}{\raggedright\small\mbox{}= Julius}\hspace{0.5em}\parbox[t]{\dimexpr0.2373\linewidth-0.5em\relax}{\raggedright\small\mbox{}( — — ).}\hspace{0.5em}}}{}
\EdL{13}{AK07-J0282-L13}{\mbox{\parbox[t]{\dimexpr0.0508\linewidth-0.5em\relax}{\raggedright\small\mbox{}11.}\hspace{0.5em}\parbox[t]{\dimexpr0.2203\linewidth-0.5em\relax}{\raggedright\small\mbox{}Ab}\hspace{0.5em}\parbox[t]{\dimexpr0.2966\linewidth-0.5em\relax}{\raggedright\small\mbox{}= Āb}\hspace{0.5em}\parbox[t]{\dimexpr0.1949\linewidth-0.5em\relax}{\raggedright\small\mbox{}= Augustus}\hspace{0.5em}\parbox[t]{\dimexpr0.2373\linewidth-0.5em\relax}{\raggedright\small\mbox{}( — — ).}\hspace{0.5em}}}{}
\EdL{14}{AK07-J0282-L14}{\mbox{\parbox[t]{\dimexpr0.0508\linewidth-0.5em\relax}{\raggedright\small\mbox{}12.}\hspace{0.5em}\parbox[t]{\dimexpr0.2203\linewidth-0.5em\relax}{\raggedright\small\mbox{}Elul}\hspace{0.5em}\parbox[t]{\dimexpr0.2966\linewidth-0.5em\relax}{\raggedright\small\mbox{}= Īlūl}\hspace{0.5em}\parbox[t]{\dimexpr0.1949\linewidth-0.5em\relax}{\raggedright\small\mbox{}= September}\hspace{0.5em}\parbox[t]{\dimexpr0.2373\linewidth-0.5em\relax}{\raggedright\small\mbox{}( — — ).}\hspace{0.5em}}}{}
\FrameHead{0283}{p. [249]}{0283}{p. [249]}
\EdLh{1}{AK07-J0283-L01}{Namen- und Sachregister.}{}
\ColHead{left column}
\EdL{2}{AK07-J0283-L02}{\textbf{A.}}{}
\EdL{3}{AK07-J0283-L03}{Abacus X.}{}
\EdL{4}{AK07-J0283-L04}{ʿAbbās b. Saʿīd al-Djauharī 32.}{}
\EdL{5}{AK07-J0283-L05}{Abraham b. ʿEzrā 41, 51, 99.}{}
\EdL{6}{AK07-J0283-L06}{Abū’l-Ḥasan ʿAlī v. Marokko 71.}{}
\EdL{7}{AK07-J0283-L07}{Abū Isḥāḳ al-Naḳḳāsh s. al-Zarḳālī.}{}
\EdL{8}{AK07-J0283-L08}{Abū Naṣr Manṣūr s. Manṣur b. ʿAlī b. ʿIrāḳ.}{}
\EdL{9}{AK07-J0283-L09}{Abū’l-Wefā 77, 78.}{}
\EdL{10}{AK07-J0283-L10}{Adam 35, T. 1.}{}
\EdL{11}{AK07-J0283-L11}{Adelhard s. Athelhard.}{}
\EdL{12}{AK07-J0283-L12}{Adiectiones XVIII.}{}
\EdL{13}{AK07-J0283-L13}{Aequans 47, 48, 50, 233.}{}
\EdL{14}{AK07-J0283-L14}{Aequatio s. Examinatio.}{}
\EdL{15}{AK07-J0283-L15}{— domorum, s. Häuser, astrolog.}{}
\EdL{16}{AK07-J0283-L16}{— nychthemerōn 80.}{}
\EdL{17}{AK07-J0283-L17}{Aequatorial-Stunden (gleiche) 24, 76, 81, 84.}{}
\EdL{18}{AK07-J0283-L18}{Aera Arabum (od. der \textit{Hidjra}) VII, IX, 1—3,}{}
\EdL{19}{AK07-J0283-L19}{33—39, 106\,; — Augusti (od. Alexandrina)}{}
\EdL{20}{AK07-J0283-L20}{XVI\,; — Coptorum XV\,; — Diocletiani XV\,;}{}
\EdL{21}{AK07-J0283-L21}{— Domini (christl., julian.) XVII, 3, 4, 35,}{}
\EdL{22}{AK07-J0283-L22}{36, 102\,; — Nabonassaris XVI\,; — Philippi}{}
\EdL{23}{AK07-J0283-L23}{Aridaei XVI\,; — Romanorum (Seleucid.)}{}
\EdL{24}{AK07-J0283-L24}{2—4, 6, 7, 34—40, 81\,; — spanische (\textit{acofra})}{}
\EdL{25}{AK07-J0283-L25}{XVIII, 3, 35\,; — Yezdegirds VII, VIII, XV,}{}
\EdL{26}{AK07-J0283-L26}{33, 37, 106. T. 1.}{}
\EdL{27}{AK07-J0283-L27}{Aḥmed b. al-Muthannā 51.}{}
\EdL{28}{AK07-J0283-L28}{Albohazen 105.}{}
\EdL{29}{AK07-J0283-L29}{Albumasar 105.}{}
\EdL{30}{AK07-J0283-L30}{Alcabitius s. Ḳabīṣī.}{}
\EdL{31}{AK07-J0283-L31}{Algebra (al-Khwārizmīs) VIII, XII.}{}
\EdL{32}{AK07-J0283-L32}{Alienatio 103, 104, T. 116.}{}
\EdL{33}{AK07-J0283-L33}{Alexander, der Grosse, 5, 34, 35, 38, T. 1.}{}
\EdL{34}{AK07-J0283-L34}{Almagest, der, des Ptolemäus XI, 2, 16, 18, 33,}{}
\EdL{35}{AK07-J0283-L35}{44, 70, etc.}{}
\EdL{36}{AK07-J0283-L36}{Altitudo s. Höhe.}{}
\EdL{37}{AK07-J0283-L37}{Ammonius s. Emirethius.}{}
\EdL{38}{AK07-J0283-L38}{ἀναπλήρωσις = das Wiedervollmachen (des}{}
\EdL{39}{AK07-J0283-L39}{Mondes, od. der Sonne bei Finsternissen) 88.}{}
\ColHead{right column}
\EdL{40}{AK07-J0283-L40}{Andersen, D. VI.}{}
\EdL{41}{AK07-J0283-L41}{Anni collecti XV, XVII, 4—7, 23, 36, etc. T. 2,}{}
\EdL{42}{AK07-J0283-L42}{3, 4, 6, etc.}{}
\EdL{43}{AK07-J0283-L43}{— plani (vel expansi, singuli, residui) XV,}{}
\EdL{44}{AK07-J0283-L44}{4—7, 23, 36, etc. T. 2, 3, 4, 6, etc.}{}
\EdL{45}{AK07-J0283-L45}{— vagi (der Ägypter) 38.}{}
\EdL{46}{AK07-J0283-L46}{Annus Arabum, Beschreibung desselbcn, 2, 34\,;}{}
\EdL{47}{AK07-J0283-L47}{Verwandlung desselben in römische Jahre,}{}
\EdL{48}{AK07-J0283-L48}{6—7, 38—40, T. 3, 3\textsuperscript{a}.}{}
\EdL{49}{AK07-J0283-L49}{— bissextus od. bissextilis (Schaltjahr) XVII,}{}
\EdL{50}{AK07-J0283-L50}{XVIII, 6—7, 34, 35, T· 3\textsuperscript{a}.}{}
\EdL{51}{AK07-J0283-L51}{— Persarum, s. Aera Yezdegirds.}{}
\EdL{52}{AK07-J0283-L52}{— Romanorum vel Alexandri, Beschreibung}{}
\EdL{53}{AK07-J0283-L53}{desselben, 2—4, 34—36\,; Auffindung desselben}{}
\EdL{54}{AK07-J0283-L54}{6—7, 38—40, T. 1, 3.}{}
\EdL{55}{AK07-J0283-L55}{Anomalia s. Argumentum.}{}
\EdL{56}{AK07-J0283-L56}{Antiochia 83 (N. 2), 95, 96.}{}
\EdL{57}{AK07-J0283-L57}{Apogeum s. Sublimatio.}{}
\EdL{58}{AK07-J0283-L58}{Apsidenlinie 65.}{}
\EdL{59}{AK07-J0283-L59}{Argumentum (Anomalia, \textit{Heza}) XIV, 9—12,}{}
\EdL{60}{AK07-J0283-L60}{43—46, 48—55, 67, 82—85, 93, T. 6—9, 11, 13,}{}
\EdL{61}{AK07-J0283-L61}{15—18, 27—56, 69—72.}{}
\EdL{62}{AK07-J0283-L62}{— (examinatum, od. aequatum) 11—13, 14,}{}
\EdL{63}{AK07-J0283-L63}{52—57, 62—64.}{}
\EdL{64}{AK07-J0283-L64}{— (Grösse, mit der man in eine Tafel eintritt)}{}
\EdL{65}{AK07-J0283-L65}{19, 20, 55, 92.}{}
\EdL{66}{AK07-J0283-L66}{— (latitudinis lunae) 23, 24, 26, 27, 29, 58, 59,}{}
\EdL{67}{AK07-J0283-L67}{68, 82, 84—86, 94, T. 69—76, 78.}{}
\EdL{68}{AK07-J0283-L68}{— (Verfahren, log. Schluss) 13, 56.}{}
\EdL{69}{AK07-J0283-L69}{— (ratum) 25, 84.}{}
\EdL{70}{AK07-J0283-L70}{Arin VII, XVIII, 1, 8, 9, 17—19, 33, 41, 43, 68, 72.}{}
\EdL{71}{AK07-J0283-L71}{Arkand, das Buch al-, s. Khaṇḍakhādyaka.}{}
\EdL{72}{AK07-J0283-L72}{Āryabhaṭṭa (od. Āryabhaṭa) 42.}{}
\EdL{73}{AK07-J0283-L73}{Ārya-Siddhānta s. Siddhāntas.}{}
\EdL{74}{AK07-J0283-L74}{Ascendens s. Horoscopus.}{}
\EdL{75}{AK07-J0283-L75}{Ascensional-Differenz 74—76, 92.}{}
\EdL{76}{AK07-J0283-L76}{Ascensiones obliquae XII, 20, 21, 73—76, 92, 93,}{}
\EdL{77}{AK07-J0283-L77}{100—102.}{}
\FrameHead{0284}{p. 250}{0284}{p. 250}
\ColHead{left column}
\EdL{1}{AK07-J0284-L01}{Ascensiones rectae (= Rectascension) 17—20,}{}
\EdL{2}{AK07-J0284-L02}{72—76, 93, 99, 100, 102.}{}
\EdL{3}{AK07-J0284-L03}{Asín, Miguel, XI, XVII.}{}
\EdL{4}{AK07-J0284-L04}{Astrolabium IX, X, XX, 18, 71.}{}
\EdL{5}{AK07-J0284-L05}{Aszendent s. Ascendens.}{}
\EdL{6}{AK07-J0284-L06}{Athelhard v. Bath VI—X, XII, XIII, XXIIl}{}
\EdL{7}{AK07-J0284-L07}{(N. 2), 1, 37, 69, 71, 80, 96, 236, T. 116.}{}
\EdL{8}{AK07-J0284-L08}{Aug s. Sublimatio.}{}
\EdL{9}{AK07-J0284-L09}{Augustus (Kaiser) XVI, XVII.}{}
\EdL{10}{AK07-J0284-L10}{Azin s. Arin.}{}
\EdL{11}{AK07-J0284-L11}{\textbf{B.}}{}
\EdL{12}{AK07-J0284-L12}{Baghdād 74.}{}
\EdL{13}{AK07-J0284-L13}{al-Battānī, Muḥ. b. Djābir b. Sinān (und seine}{}
\EdL{14}{AK07-J0284-L14}{astron. Tafeln, ed. Nallino) IX, XIII, XV,}{}
\EdL{15}{AK07-J0284-L15}{XVI, XVIII, XXI, XXIII, XXIV, 32, 34, 35,}{}
\EdL{16}{AK07-J0284-L16}{38‐‐50, 52, 54—60, 62—66, 69, 71—97, 99, 100,}{}
\EdL{17}{AK07-J0284-L17}{102—106, 233, 234, 236.}{}
\EdL{18}{AK07-J0284-L18}{Bernelinus X.}{}
\EdL{19}{AK07-J0284-L19}{Bernhelmus s. das vorhergeh.}{}
\EdL{20}{AK07-J0284-L20}{Besthorn, R. O., V, XXV.}{}
\EdL{21}{AK07-J0284-L21}{Bewegung, mittlere (der Sonne) XIV, XVIII,}{}
\EdL{22}{AK07-J0284-L22}{7—9, 23, 33, 36. 41—44, 65, 67, 82, 84, 85, 105,}{}
\EdL{23}{AK07-J0284-L23}{106, 235, T. 4, 5.}{}
\EdL{24}{AK07-J0284-L24}{— mittlere (des Mondes) XIV, XVIII, 7, 8, 10,}{}
\EdL{25}{AK07-J0284-L25}{23, 33, 36, 41—43, 45, 46, 65, 67, 82, 84, 85,}{}
\EdL{26}{AK07-J0284-L26}{105, 106, T. 6, 7, 8.}{}
\EdL{27}{AK07-J0284-L27}{— mittlere (der obern Planeten) XIV, XVIII,}{}
\EdL{28}{AK07-J0284-L28}{7, 8, 10—12, 36, 41—44, 51—54, 63, 64, 105,}{}
\EdL{29}{AK07-J0284-L29}{106, T. 9—14.}{}
\EdL{30}{AK07-J0284-L30}{— mittlere (des Drachenkopfes) XIV, XIX, 7,}{}
\EdL{31}{AK07-J0284-L31}{8, 12, T. 19, 20.}{}
\EdL{32}{AK07-J0284-L32}{— mittlere (der Anomalie der untern Pla-}{}
\EdL{33}{AK07-J0284-L33}{neten) 54, T. 15—18.}{}
\EdL{34}{AK07-J0284-L34}{— wahre (d. h. ungleiche, des Mondes) 22, 23,}{}
\EdL{35}{AK07-J0284-L35}{28, 78, 79, 81, 83, 105, T. 61—66.}{}
\EdL{36}{AK07-J0284-L36}{— wahre (der Sonne) 22, 23, 78, 79, 81, 83, 105,}{}
\EdL{37}{AK07-J0284-L37}{T. 61—66.}{}
\EdL{38}{AK07-J0284-L38}{Bhāskara, 70.}{}
\EdL{39}{AK07-J0284-L39}{Billeter, G. Vl.}{}
\EdL{40}{AK07-J0284-L40}{al-Bīrūnī, Muḥ. b. Aḥmed, 37, 38, 44, 49, 51,}{}
\EdL{41}{AK07-J0284-L41}{66, 79, 89, 90, 105, 235.}{}
\EdL{42}{AK07-J0284-L42}{Björnbo, A. V, VII, X, XI, XIII, XXIV, 68,}{}
\EdL{43}{AK07-J0284-L43}{69—71, 76, 77, 102.}{}
\EdL{44}{AK07-J0284-L44}{Boll, F. 105.}{}
\EdL{45}{AK07-J0284-L45}{Boncompagni, B. VIII.}{}
\EdL{46}{AK07-J0284-L46}{Bouché-Leclercq, A. 105.}{}
\EdL{47}{AK07-J0284-L47}{Brahmagupta 32, 42, 45, 79, 90.}{}
\EdL{48}{AK07-J0284-L48}{v. Braunmühl, A. 71.}{}
\EdL{49}{AK07-J0284-L49}{Breite, des Mondes 14, 27, 29, 58, 59, 66—68,}{}
\EdL{50}{AK07-J0284-L50}{86, 88, 94, T. 21—26.}{}
\EdL{51}{AK07-J0284-L51}{— der Planeten, 14, 60—64, S. 27—56.}{}
\ColHead{right column}
\EdL{52}{AK07-J0284-L52}{Breite, scheinbare, des Mondes (durch die Pa-}{}
\EdL{53}{AK07-J0284-L53}{rallaxe hervorgebracht) 94.}{}
\EdL{54}{AK07-J0284-L54}{— geograph. eines Ortes (Bestimmung der-}{}
\EdL{55}{AK07-J0284-L55}{selben) 18, 28, 71.}{}
\EdL{56}{AK07-J0284-L56}{Brevis circulus s. Epicyklus.}{}
\EdL{57}{AK07-J0284-L57}{Broger, M. 95 (N. 1).}{}
\EdL{58}{AK07-J0284-L58}{Bukhtanaṣsar s. Nabonassar.}{}
\EdL{59}{AK07-J0284-L59}{Burgess, E. 42, 45, 46, 51, 52, 58—61, 68, 75,}{}
\EdL{60}{AK07-J0284-L60}{78—80, 86, 90, 92.}{}
\EdL{61}{AK07-J0284-L61}{\textbf{C.}}{}
\EdL{62}{AK07-J0284-L62}{Caldei 9, 11, 43.}{}
\EdL{63}{AK07-J0284-L63}{Caput draconis s. Knotenpunkte (des Mondes).}{}
\EdL{64}{AK07-J0284-L64}{Carra de Vaux 78.}{}
\EdL{65}{AK07-J0284-L65}{Casae s. Häuser.}{}
\EdL{66}{AK07-J0284-L66}{Casus 26, 27, 87—89, 95, s. auch Tiefstand.}{}
\EdL{67}{AK07-J0284-L67}{Cauda draconis s. Knotenpunkte (des Mondes).}{}
\EdL{68}{AK07-J0284-L68}{Celada (Kardinal) XI.}{}
\EdL{69}{AK07-J0284-L69}{Centrum (astron.) 11, 48—54, 57, 58. s. auch}{}
\EdL{70}{AK07-J0284-L70}{Examinatio centri.}{}
\EdL{71}{AK07-J0284-L71}{— examinatum vel aequatum 11, 12, 52—55, 57.}{}
\EdL{72}{AK07-J0284-L72}{— ultimum 11, 53, 54.}{}
\EdL{73}{AK07-J0284-L73}{Chaldei s. Caldei.}{}
\EdL{74}{AK07-J0284-L74}{Chasles, M. VII, X, 43, 49, 50.}{}
\EdL{75}{AK07-J0284-L75}{Cheikho, L. IX (N. 3).}{}
\EdL{76}{AK07-J0284-L76}{Cifra s. Null.}{}
\EdL{77}{AK07-J0284-L77}{Circulus aequinoctialis 18—21, 25.}{}
\EdL{78}{AK07-J0284-L78}{— excentricus 26, 44, 47.}{}
\EdL{79}{AK07-J0284-L79}{Circumpolarsterne 71.}{}
\EdL{80}{AK07-J0284-L80}{Coaequatio s. Examinatio.}{}
\EdL{81}{AK07-J0284-L81}{Computatio XIV, 59, T. 21—26.}{}
\EdL{82}{AK07-J0284-L82}{Computus (Gerlandi) X.}{}
\EdL{83}{AK07-J0284-L83}{Conversio annorum s. Umwandlung.}{}
\EdL{84}{AK07-J0284-L84}{Corduba (Cordova) VII, XII, 25, 33, 81, 83—85,}{}
\EdL{85}{AK07-J0284-L85}{92—97, 100, 106.}{}
\EdL{86}{AK07-J0284-L86}{Cotangens s. Umbra.}{}
\EdL{87}{AK07-J0284-L87}{Curtze, M. XXIII.}{}
\EdL{88}{AK07-J0284-L88}{\textbf{D.}}{}
\EdL{89}{AK07-J0284-L89}{Darius Medus 105.}{}
\EdL{90}{AK07-J0284-L90}{Decani s. Facies.}{}
\EdL{91}{AK07-J0284-L91}{Defectus s. Verfinsterung.}{}
\EdL{92}{AK07-J0284-L92}{Deferens 16, 17, 57, 60, 61, 233, 234.}{}
\EdL{93}{AK07-J0284-L93}{Deiectio s. Tiefstand.}{}
\EdL{94}{AK07-J0284-L94}{Deklination (Obliquatio) der Sonne 13, 18, 19,}{}
\EdL{95}{AK07-J0284-L95}{58, 91, T. 21—26.}{}
\EdL{96}{AK07-J0284-L96}{— irgend eines Punktes der Ekliptik 74, 75,}{}
\EdL{97}{AK07-J0284-L97}{s. auch Ekliptik.}{}
\EdL{98}{AK07-J0284-L98}{Delambre 60.}{}
\EdL{99}{AK07-J0284-L99}{Detrimentum s. Alienatio.}{}
\EdL{100}{AK07-J0284-L100}{Dieterici, F. 105.}{}
\EdL{101}{AK07-J0284-L101}{Digiti = Teile des Gnomon 20—22, T. 60.}{}
\FrameHead{0285}{p. 251}{0285}{p. 251}
\ColHead{left column}
\EdL{1}{AK07-J0285-L01}{Digiti = Teile des Mond- bezw. Sonnendurch-}{}
\EdL{2}{AK07-J0285-L02}{messers 26, 27, 87—89, 95.}{}
\EdL{3}{AK07-J0285-L03}{Dignitates s. Erhöhungen.}{}
\EdL{4}{AK07-J0285-L04}{Diluvium 3, 35, T. 1.}{}
\EdL{5}{AK07-J0285-L05}{Diocletianus (Kaiser) XV, T. 1.}{}
\EdL{6}{AK07-J0285-L06}{Dispositio (casarum duodecim) s. Häuser,}{}
\EdL{7}{AK07-J0285-L07}{astrol.}{}
\EdL{8}{AK07-J0285-L08}{Distantia definita 16, 68.}{}
\EdL{9}{AK07-J0285-L09}{Distanz, geozentr. des Epizykelzentrums 48, 49.}{}
\EdL{10}{AK07-J0285-L10}{Diversitas aspectus vel aspectuum s. Parallaxe.}{}
\EdL{11}{AK07-J0285-L11}{— respectus s. dasselbe.}{}
\EdL{12}{AK07-J0285-L12}{Djābir b. Aflaḥ 69.}{}
\EdL{13}{AK07-J0285-L13}{Domicilia 103, 106, T. 116.}{}
\EdL{14}{AK07-J0285-L14}{Domini 31, 103, 104, T. 116.}{}
\EdL{15}{AK07-J0285-L15}{Dorotheus 99, 100.}{}
\EdL{16}{AK07-J0285-L16}{Dozy, R. 36.}{}
\EdL{17}{AK07-J0285-L17}{Drachenkopf s. Knotenpunkte (des Mondes).}{}
\EdL{18}{AK07-J0285-L18}{Drachenschwanz s. dasselbe.}{}
\EdL{19}{AK07-J0285-L19}{Dracones s. dasselbe.}{}
\EdL{20}{AK07-J0285-L20}{Durchmesser (rota) scheinbarer (des Mondes)}{}
\EdL{21}{AK07-J0285-L21}{23, 78—79, 90, T. 61—66.}{}
\EdL{22}{AK07-J0285-L22}{— scheinbarer (der Sonne) 22—23, 78—79, T.}{}
\EdL{23}{AK07-J0285-L23}{61—66.}{}
\EdL{24}{AK07-J0285-L24}{— des Erdschattens 90, 91, T. 61—66.}{}
\EdL{25}{AK07-J0285-L25}{\textbf{E.}}{}
\EdL{26}{AK07-J0285-L26}{Eclipsis s. Verfinsterung.}{}
\EdL{27}{AK07-J0285-L27}{Ekliptik 44, 47, 58—61, 63, 72—74, 76, 96—101\,;}{}
\EdL{28}{AK07-J0285-L28}{Schiefe derselben 19, 33, 58, 72, 105.}{}
\EdL{29}{AK07-J0285-L29}{Ekliptik-Zenitdistanz 92, 93.}{}
\EdL{30}{AK07-J0285-L30}{Emersio s. Respiratio.}{}
\EdL{31}{AK07-J0285-L31}{Emirethius (?), vielleicht Eumathius (wahr-}{}
\EdL{32}{AK07-J0285-L32}{scheinl. = Ammonius), dessen Kanon}{}
\EdL{33}{AK07-J0285-L33}{(astron. Tafeln) von al-Zarḳālī verbessert}{}
\EdL{34}{AK07-J0285-L34}{worden sein soll (vergl. Nachtrag 5) T. 7, 92.}{}
\EdL{35}{AK07-J0285-L35}{ἔμπτωσις = das Hineinfallen 87.}{}
\EdL{36}{AK07-J0285-L36}{Epicyclus XX, XXI, 44, 47—49, 51, 57, 60, 61,}{}
\EdL{37}{AK07-J0285-L37}{234\,; Radius desselben beim Mond 46, bei}{}
\EdL{38}{AK07-J0285-L38}{den Planeten 48.}{}
\EdL{39}{AK07-J0285-L39}{Era s. Aera.}{}
\EdL{40}{AK07-J0285-L40}{Erdschatten s. Umbra terrae.}{}
\EdL{41}{AK07-J0285-L41}{Erhöhungen 103, 104, 106, T. 116.}{}
\EdL{42}{AK07-J0285-L42}{Eumathius s. Emirethius.}{}
\EdL{43}{AK07-J0285-L43}{Evectio 45.}{}
\EdL{44}{AK07-J0285-L44}{Exagonum s. Sextilschein.}{}
\EdL{45}{AK07-J0285-L45}{Exaltatio s. Erhöhung.}{}
\EdL{46}{AK07-J0285-L46}{Examen s. das folg.}{}
\EdL{47}{AK07-J0285-L47}{Examinatio (Gleichung) der Sonne, des Mondes}{}
\EdL{48}{AK07-J0285-L48}{und der Planeten XIV, 9, 10, 24, 43—46, 82,}{}
\EdL{49}{AK07-J0285-L49}{83, T. 21—26.}{}
\EdL{50}{AK07-J0285-L50}{— argumenti (= anomaliae, \textit{elheza}) 10, 11,}{}
\EdL{51}{AK07-J0285-L51}{48—54, 233—235, T. 27—56.}{}
\ColHead{right column}
\EdL{52}{AK07-J0285-L52}{Examinatio centri 11, 12, 48—54, 233, 234, T.}{}
\EdL{53}{AK07-J0285-L53}{27—56.}{}
\EdL{54}{AK07-J0285-L54}{Exzentrizität, der Planetenbahnen 47, 48, 234.}{}
\EdL{55}{AK07-J0285-L55}{— der Sonnenbahn 45.}{}
\EdL{56}{AK07-J0285-L56}{\textbf{F.}}{}
\EdL{57}{AK07-J0285-L57}{Facies (= Decanus), der dritte Teil (10°) eines}{}
\EdL{58}{AK07-J0285-L58}{Tierkreiszeichens, 16, 68, 103, 104, T. 116.}{}
\EdL{59}{AK07-J0285-L59}{Fall s. Tiefstand.}{}
\EdL{60}{AK07-J0285-L60}{Feria = Wochentag 1, 5—7, 36, 37.}{}
\EdL{61}{AK07-J0285-L61}{Fines s. Grenzen.}{}
\EdL{62}{AK07-J0285-L62}{Finsternisse s. Verfinsterung.}{}
\EdL{63}{AK07-J0285-L63}{Fleischer, H. L. 105.}{}
\EdL{64}{AK07-J0285-L64}{\textbf{G.}}{}
\EdL{65}{AK07-J0285-L65}{Gegenschein 99.}{}
\EdL{66}{AK07-J0285-L66}{Geib s. Sinus.}{}
\EdL{67}{AK07-J0285-L67}{Gerbert (Pabst Sylvester II.) X.}{}
\EdL{68}{AK07-J0285-L68}{Gerhard von Cremona VIII, 69.}{}
\EdL{69}{AK07-J0285-L69}{Gerlandus X.}{}
\EdL{70}{AK07-J0285-L70}{Ghoarizmi s. Khwārizmī.}{}
\EdL{71}{AK07-J0285-L71}{Gleichung der Anomalie und des Centrums}{}
\EdL{72}{AK07-J0285-L72}{(bei Battānī) 48, 50.}{}
\EdL{73}{AK07-J0285-L73}{— des Apogeums 50.}{}
\EdL{74}{AK07-J0285-L74}{— der Apsis 51.}{}
\EdL{75}{AK07-J0285-L75}{— der Konjunktion 51.}{}
\EdL{76}{AK07-J0285-L76}{Gleichungen (der Sonne, des Mondes, der}{}
\EdL{77}{AK07-J0285-L77}{Anomalie, des Centrums) s. Examinatio.}{}
\EdL{78}{AK07-J0285-L78}{— (der Planeten) 33, 49, 55, 105\,; s. auch}{}
\EdL{79}{AK07-J0285-L79}{Examinatio (argumenti und centri).}{}
\EdL{80}{AK07-J0285-L80}{Gnomon 72, 74—76.}{}
\EdL{81}{AK07-J0285-L81}{Goebel, J. XXII (N. 1).}{}
\EdL{82}{AK07-J0285-L82}{Gonzalez XVlI (N.)}{}
\EdL{83}{AK07-J0285-L83}{Greenwich 83, 85.}{}
\EdL{84}{AK07-J0285-L84}{Grenzen (ungleiche Teile der Tierkreiszeichen)}{}
\EdL{85}{AK07-J0285-L85}{104, 105, T. 116.}{}
\EdL{86}{AK07-J0285-L86}{Gründe der Tafeln des Khwārizmī 41, 51.}{}
\EdL{87}{AK07-J0285-L87}{\textbf{H.}}{}
\EdL{88}{AK07-J0285-L88}{Ḥabash, der Rechner, 66, 76, 77.}{}
\EdL{89}{AK07-J0285-L89}{Halbierung der Gleichung (der Planeten) 49—52.}{}
\EdL{90}{AK07-J0285-L90}{Ḥarrān 74.}{}
\EdL{91}{AK07-J0285-L91}{Haskins, C. H. VII, X, 236.}{}
\EdL{92}{AK07-J0285-L92}{Häuser, die 12 astrologischen, 25\,; Verteilung}{}
\EdL{93}{AK07-J0285-L93}{(aequatio, dispositio) derselben 29, 30, 96—98,}{}
\EdL{94}{AK07-J0285-L94}{106, T. 79—90.}{}
\EdL{95}{AK07-J0285-L95}{Heiberg, J. L. XIII, XXIV, XXV, 44, 57, 74, 81,}{}
\EdL{96}{AK07-J0285-L96}{90, 233.}{}
\EdL{97}{AK07-J0285-L97}{Henry, C. XXIII.}{}
\EdL{98}{AK07-J0285-L98}{Hermes Trismegistos XXII.}{}
\EdL{99}{AK07-J0285-L99}{Hermannus Contractus X.}{}
\EdL{100}{AK07-J0285-L100}{— Dalmata XIII, XVIII (N.).}{}
\FrameHead{0286}{p. 252}{0286}{p. 252}
\ColHead{left column}
\EdL{1}{AK07-J0286-L01}{Heza s. Argumentum.}{}
\EdL{2}{AK07-J0286-L02}{Höhe der Sonne (artifa solis) 21—22, 76, 77.}{}
\EdL{3}{AK07-J0286-L03}{Horae aequales s. Aequatorial-Stunden\,; — in-}{}
\EdL{4}{AK07-J0286-L04}{aequales s. Zeitstunden.}{}
\EdL{5}{AK07-J0286-L05}{— diurnae s. Tagesstunden.}{}
\EdL{6}{AK07-J0286-L06}{— nocturnae s. Nachtstunden.}{}
\EdL{7}{AK07-J0286-L07}{Horoscopi signorum (= Aufgänge der Tier-}{}
\EdL{8}{AK07-J0286-L08}{kreiszeichen) s. Signum, und Glossar, s. v.}{}
\EdL{9}{AK07-J0286-L09}{Matale.}{}
\EdL{10}{AK07-J0286-L10}{Horoscopus (Ascendens), der zu einer be-}{}
\EdL{11}{AK07-J0286-L11}{stimmten Zeit aufgehende Punkt der Eklip-}{}
\EdL{12}{AK07-J0286-L12}{tik, 25, 28—31, 91—94, 97, 98, 100, 101, T. 79.}{}
\EdL{13}{AK07-J0286-L13}{Houtsma, M. Th. XXV.}{}
\EdL{14}{AK07-J0286-L14}{\textbf{I.}}{}
\EdL{15}{AK07-J0286-L15}{Ibn Abī Uṣaibiʿa VII.}{}
\EdL{16}{AK07-J0286-L16}{— al-Adamī 33, 49.}{}
\EdL{17}{AK07-J0286-L17}{— ʿEzrā s. Abraham b. ʿEzrā.}{}
\EdL{18}{AK07-J0286-L18}{— Hibintā 43, 46, 99, 100, 106.}{}
\EdL{19}{AK07-J0286-L19}{— al-Ḳifṭī VII, 33.}{}
\EdL{20}{AK07-J0286-L20}{— Ṣāʿid s. Ḳāḍī Ṣāʿid.}{}
\EdL{21}{AK07-J0286-L21}{— Yūnus 47.}{}
\EdL{22}{AK07-J0286-L22}{Ibrāhīm al-Fazārī VIII, 32, 33.}{}
\EdL{23}{AK07-J0286-L23}{Ideler, L. XVIII, 36.}{}
\EdL{24}{AK07-J0286-L24}{Immersio s. Casus.}{}
\EdL{25}{AK07-J0286-L25}{Intercalation (des Schalttages) XV, 34.}{}
\EdL{26}{AK07-J0286-L26}{Interstitium 24, 29, 93.}{}
\EdL{27}{AK07-J0286-L27}{Isidor von Sevilla XVIlI.}{}
\EdL{28}{AK07-J0286-L28}{\textbf{J.}}{}
\EdL{29}{AK07-J0286-L29}{Jahr, arabisches, s. Annus Arabum.}{}
\EdL{30}{AK07-J0286-L30}{— römisches (seleucid.) s. Annus Romanorum.}{}
\EdL{31}{AK07-J0286-L31}{— siderisches (Sternjahr) 42, 65, 66, 102, 105,}{}
\EdL{32}{AK07-J0286-L32}{235.}{}
\EdL{33}{AK07-J0286-L33}{— tropisches 42, 66, 102.}{}
\EdL{34}{AK07-J0286-L34}{Johannes de Sacrobosco XXIII (N. 1).}{}
\EdL{35}{AK07-J0286-L35}{— von Sevilla Xl.}{}
\EdL{36}{AK07-J0286-L36}{Jordanus Nemorarius XI.}{}
\EdL{37}{AK07-J0286-L37}{Judices (= Domini) 104, T. 116.}{}
\EdL{38}{AK07-J0286-L38}{\textbf{K.}}{}
\EdL{39}{AK07-J0286-L39}{al-Ḳabīṣī 99, 105.}{}
\EdL{40}{AK07-J0286-L40}{Ḳāḍī Ṣāʿid VII, IX, 33.}{}
\EdL{41}{AK07-J0286-L41}{Kardadja, Teilung desselben 70.}{}
\EdL{42}{AK07-J0286-L42}{Karpinski, L. C. XII.}{}
\EdL{43}{AK07-J0286-L43}{Kaurezmi s. Khwārizmī.}{}
\EdL{44}{AK07-J0286-L44}{Khaṇḍakhādyaka 79, 90.}{}
\EdL{45}{AK07-J0286-L45}{al-Khwārizmī, Muh. b. Mūsā VII—X, XII,}{}
\EdL{46}{AK07-J0286-L46}{XIII, XV, XXI, XXIl, XXIV, 1, 32—35, 38—42,}{}
\EdL{47}{AK07-J0286-L47}{45—47, 49—60, 62—66, 68—75, 77—81, 86, 87,}{}
\EdL{48}{AK07-J0286-L48}{90—97, 99—103, 105, 106, 233—236, T. 116.}{}
\ColHead{right column}
\EdL{49}{AK07-J0286-L49}{al-Kindī 105.}{}
\EdL{50}{AK07-J0286-L50}{Knotenlinie, des Mondes 65.}{}
\EdL{51}{AK07-J0286-L51}{Knotenpunkte, des Mondes 12, 14, 15, 27, 41,}{}
\EdL{52}{AK07-J0286-L52}{55, 65, 67, 68.}{}
\EdL{53}{AK07-J0286-L53}{— der Planeten 15, 64.}{}
\EdL{54}{AK07-J0286-L54}{Konjunktion und Opposition (Syzygien) der}{}
\EdL{55}{AK07-J0286-L55}{Sonne und des Mondes XIX, 23—29, 33,}{}
\EdL{56}{AK07-J0286-L56}{81—86, 89, 92, 94—96, 106, T. 69—72.}{}
\EdL{57}{AK07-J0286-L57}{— scheinbare, von Sonne und Mond 94, 95.}{}
\EdL{58}{AK07-J0286-L58}{Kreis, exzentrischer s. Circulus excentricus.}{}
\EdL{59}{AK07-J0286-L59}{\textbf{L.}}{}
\EdL{60}{AK07-J0286-L60}{Länge, mittlere s. Bewegung, mittlere.}{}
\EdL{61}{AK07-J0286-L61}{Länge, wahre, des Mondes 10, 14, 16, 33, 45,}{}
\EdL{62}{AK07-J0286-L62}{46, 58, 59, 66—68, 81—86, 93, T. 6—8, 21—26.}{}
\EdL{63}{AK07-J0286-L63}{— — der Sonne 9, 10, 16, 33, 43—45, 66, 67,}{}
\EdL{64}{AK07-J0286-L64}{81—86, T. 4, 5, 21—26.}{}
\EdL{65}{AK07-J0286-L65}{— — der obern Planeten 10, 11, 33, 46—54,}{}
\EdL{66}{AK07-J0286-L66}{60, 62, T. 9—14, 27—44.}{}
\EdL{67}{AK07-J0286-L67}{— — der untern Planeten 12, 33, 54, 55, 60,}{}
\EdL{68}{AK07-J0286-L68}{62, T. 15—18, 45—56.}{}
\EdL{69}{AK07-J0286-L69}{— — des Drachenkopfes 12, 55, 58, 59, 67, 68,}{}
\EdL{70}{AK07-J0286-L70}{T. 19, 20.}{}
\EdL{71}{AK07-J0286-L71}{Latitudo lunae, etc., s. Breite.}{}
\EdL{72}{AK07-J0286-L72}{— regionum (habitationum) s. Breite, geogr.}{}
\EdL{73}{AK07-J0286-L73}{Liber Embadorum X.}{}
\EdL{74}{AK07-J0286-L74}{Lippert, J. 33.}{}
\EdL{75}{AK07-J0286-L75}{London, Meridian von, XVIII.}{}
\EdL{76}{AK07-J0286-L76}{Longinquitas maxima lunae s. Spatium lon-}{}
\EdL{77}{AK07-J0286-L77}{ginquum.}{}
\EdL{78}{AK07-J0286-L78}{— minima lunae s. Spatium propinquum.}{}
\EdL{79}{AK07-J0286-L79}{Loth, O. 105.}{}
\EdL{80}{AK07-J0286-L80}{\textbf{M.}}{}
\EdL{81}{AK07-J0286-L81}{al-Madjrīṭī s. Maslama.}{}
\EdL{82}{AK07-J0286-L82}{al-Ma’mūn, Khalife VIII, 32, 43, 47, 106.}{}
\EdL{83}{AK07-J0286-L83}{Manitius, K. 44, 57, 74, 81, 90, 233.}{}
\EdL{84}{AK07-J0286-L84}{al-Manṣūr, Khalife 70.}{}
\EdL{85}{AK07-J0286-L85}{— b. ʿAlī b. ʿIrāḳ 49.}{}
\EdL{86}{AK07-J0286-L86}{Maslama b. Aḥmed al-Madjrīṭī VII—X, XIIl,}{}
\EdL{87}{AK07-J0286-L87}{XV, XXl, XXII, 32—35, 39, 40, 77, 81, 86, 90,}{}
\EdL{88}{AK07-J0286-L88}{100, 102, 103, 105, 106.}{}
\EdL{89}{AK07-J0286-L89}{Medialitas s. Bewegung, mittlere.}{}
\EdL{90}{AK07-J0286-L90}{Mediocritas s. dasselbe.}{}
\EdL{91}{AK07-J0286-L91}{Medium coelum 25, 28, 29, 91—94.}{}
\EdL{92}{AK07-J0286-L92}{Medius cursus s. Bewegung, mittlere.}{}
\EdL{93}{AK07-J0286-L93}{Megistophiel XXII.}{}
\EdL{94}{AK07-J0286-L94}{Mekka 74.}{}
\EdL{95}{AK07-J0286-L95}{Menelaus, Satz des, 74.}{}
\EdL{96}{AK07-J0286-L96}{Mephistopheles XXII.}{}
\EdL{97}{AK07-J0286-L97}{Meridian-Distanz 92.}{}
\FrameHead{0287}{p. 253}{0287}{p. 253}
\ColHead{left column}
\EdL{1}{AK07-J0287-L01}{Monat, anomalistischer 64.}{}
\EdL{2}{AK07-J0287-L02}{Monate, arabische, Beginn derselben 2, 5, 36,}{}
\EdL{3}{AK07-J0287-L03}{37, T. 2\,; Namen derselben XIV, XV, 2, 34,}{}
\EdL{4}{AK07-J0287-L04}{T. 2, 3, 4, etc. und Glossar.}{}
\EdL{5}{AK07-J0287-L05}{— ägyptische od. koptische 37, 38, T. 2\textsuperscript{b}\,; Na-}{}
\EdL{6}{AK07-J0287-L06}{men derselben T. 2\textsuperscript{b} und Glossar.}{}
\EdL{7}{AK07-J0287-L07}{— persische 37, T. 2\textsuperscript{a}\,; Namen derselben T. 2\textsuperscript{a}}{}
\EdL{8}{AK07-J0287-L08}{und Glossar.}{}
\EdL{9}{AK07-J0287-L09}{— römische (seleucid.), Beginn derselben 6,}{}
\EdL{10}{AK07-J0287-L10}{7, 38—40, T. 3\textsuperscript{a}\,; Namen derselben T. 3\textsuperscript{a} und}{}
\EdL{11}{AK07-J0287-L11}{Glossar.}{}
\EdL{12}{AK07-J0287-L12}{Mond s. Bewegung, Breite, Länge, Sublimatio,}{}
\EdL{13}{AK07-J0287-L13}{Durchmesser, Examinatio, Verfinsterung\,;}{}
\EdL{14}{AK07-J0287-L14}{Erscheinen desselben s. Neumond.}{}
\EdL{15}{AK07-J0287-L15}{Mondbahn, Schiefe derselben zur Ekliptik 59.}{}
\EdL{16}{AK07-J0287-L16}{Mondknoten, Länge derselben, s. Länge des}{}
\EdL{17}{AK07-J0287-L17}{Drachenkopfes.}{}
\EdL{18}{AK07-J0287-L18}{μονή = das Bleiben, Verweilen s. Mora.}{}
\EdL{19}{AK07-J0287-L19}{Mora 26, 27, 87—89.}{}
\EdL{20}{AK07-J0287-L20}{Motus verus vel inaequalis s. Bewegung, wahre.}{}
\EdL{21}{AK07-J0287-L21}{Muḥammed b. Mūsā al-Khwārizmī s. al-Khwā-}{}
\EdL{22}{AK07-J0287-L22}{rizmī.}{}
\EdL{23}{AK07-J0287-L23}{Muḥammed b. Mūsā b. Shākir VIII.}{}
\EdL{24}{AK07-J0287-L24}{Multiplikationstafel für Sexagesimalteile 68,}{}
\EdL{25}{AK07-J0287-L25}{T. 57\textsuperscript{b}.}{}
\EdL{26}{AK07-J0287-L26}{\textbf{N.}}{}
\EdL{27}{AK07-J0287-L27}{Nabonassar XVI, 35, T. 1.}{}
\EdL{28}{AK07-J0287-L28}{Nabukodonossor s. das vorhergeh.}{}
\EdL{29}{AK07-J0287-L29}{Nachtbogen, des Ascendenten 97.}{}
\EdL{30}{AK07-J0287-L30}{Nachtstunden 20, 25, 97, 98.}{}
\EdL{31}{AK07-J0287-L31}{Nallino, C. A. IX, XIII, XV, XVI, XVIII (N.),}{}
\EdL{32}{AK07-J0287-L32}{XXI—XXIII (N.), XXIV, 32, 34, 38, 43—48,}{}
\EdL{33}{AK07-J0287-L33}{54—58, 60, 63, 65, 66, 71—73, 76—78, 80—82,}{}
\EdL{34}{AK07-J0287-L34}{86, 87, 90—92, 94—96, 99, 100, 102, 103, 105,}{}
\EdL{35}{AK07-J0287-L35}{233, 236\,; s. auch al-Battānī.}{}
\EdL{36}{AK07-J0287-L36}{Neumond, sein Erscheinen, 16, 17, 66—68, T. 57\textsuperscript{a}.}{}
\EdL{37}{AK07-J0287-L37}{Notae, od. Notulae (auch Signa) 5, 34, 37, 39, 82.}{}
\EdL{38}{AK07-J0287-L38}{Null (\textit{cifra}) 8, 10, 36\,; Zeichen dafür XXIII, 8.}{}
\EdL{39}{AK07-J0287-L39}{\textbf{O.}}{}
\EdL{40}{AK07-J0287-L40}{Obliquatio s. Deklination.}{}
\EdL{41}{AK07-J0287-L41}{Ocreat, N. XXIII (N. 2).}{}
\EdL{42}{AK07-J0287-L42}{Ophiel XXII.}{}
\EdL{43}{AK07-J0287-L43}{Oppolzer, Th. 89. 95.}{}
\EdL{44}{AK07-J0287-L44}{Opposition s. Konjunktion.}{}
\EdL{45}{AK07-J0287-L45}{Ort der Strahlenwerfung s. Radiatio stellarum.}{}
\EdL{46}{AK07-J0287-L46}{Örter, mittlere der Planeten s. Bewegung,}{}
\EdL{47}{AK07-J0287-L47}{mittlere.}{}
\EdL{48}{AK07-J0287-L48}{— wahre der Planeten s. Länge, wahre.}{}
\EdL{49}{AK07-J0287-L49}{Ortus iuxta circulum aequinoctialem s. Ascen-}{}
\EdL{50}{AK07-J0287-L50}{siones rectae.}{}
\ColHead{right column}
\EdL{51}{AK07-J0287-L51}{\textbf{P.}}{}
\EdL{52}{AK07-J0287-L52}{Panchasiddhāntikā (= das Buch der fünf}{}
\EdL{53}{AK07-J0287-L53}{Siddhāntas) 46, 64, 68, 75, 90, 235.}{}
\EdL{54}{AK07-J0287-L54}{Parallaxe des Mondes in Länge 28, 29, 91—93,}{}
\EdL{55}{AK07-J0287-L55}{95, 96, 106, T. 77, 77\textsuperscript{a}.}{}
\EdL{56}{AK07-J0287-L56}{— — in Breite 28, 29, 93, 94, 96, 106, T. 77, 77\textsuperscript{a}.}{}
\EdL{57}{AK07-J0287-L57}{Pechüle, C. F. VI, 233—235.}{}
\EdL{58}{AK07-J0287-L58}{Perigeum s. Submissio.}{}
\EdL{59}{AK07-J0287-L59}{Petrus de Dacia XXIII (N. 2).}{}
\EdL{60}{AK07-J0287-L60}{Philippus Aridaeus XVI, T. 1.}{}
\EdL{61}{AK07-J0287-L61}{Planeten s. Bewegung, Breite, Länge, Subli-}{}
\EdL{62}{AK07-J0287-L62}{matio, Examinatio, Exzentrizität, Stationen.}{}
\EdL{63}{AK07-J0287-L63}{Planisphaerium, des Ptolemäus IX, XIII.}{}
\EdL{64}{AK07-J0287-L64}{Plato von Tivoli 57.}{}
\EdL{65}{AK07-J0287-L65}{Polhöhe s. Breite, geogr.}{}
\EdL{66}{AK07-J0287-L66}{Praezession des Frühlingspunktes 42, 46, 50,}{}
\EdL{67}{AK07-J0287-L67}{52, 54.}{}
\EdL{68}{AK07-J0287-L68}{Proiectio lucis stellarum s. Radiatio stellarum.}{}
\EdL{69}{AK07-J0287-L69}{— radiorum s. dasselbe.}{}
\EdL{70}{AK07-J0287-L70}{Prosneusis 45.}{}
\EdL{71}{AK07-J0287-L71}{Prosthaphaeresis der Länge 49, 233.}{}
\EdL{72}{AK07-J0287-L72}{Ptolemäus IX, XIIl, XXIV, 2, 16, 18, 33, 35, 38,}{}
\EdL{73}{AK07-J0287-L73}{44—47, 49, 50, 57—60, 63, 66, 70, 74, 80, 81,}{}
\EdL{74}{AK07-J0287-L74}{86, 87, 90. 96, 99, 104—106, 233, 234.}{}
\EdL{75}{AK07-J0287-L75}{Puncta (= Minuta) 36.}{}
\EdL{76}{AK07-J0287-L76}{\textbf{Q.}}{}
\EdL{77}{AK07-J0287-L77}{Quadratschein (Tetragonum) 30, 31, 98, 99, 101,}{}
\EdL{78}{AK07-J0287-L78}{T. 91—114.}{}
\EdL{79}{AK07-J0287-L79}{Quadripartitum 96, 99.}{}
\EdL{80}{AK07-J0287-L80}{Quantitas incidentiae s. Casus.}{}
\EdL{81}{AK07-J0287-L81}{\textbf{R.}}{}
\EdL{82}{AK07-J0287-L82}{Radiatio stellarum (exagonalis, tetragonalis,}{}
\EdL{83}{AK07-J0287-L83}{trigonalis) 30, 31, 98—102, 106, T. 91—114.}{}
\EdL{84}{AK07-J0287-L84}{— dextra 31, 101.}{}
\EdL{85}{AK07-J0287-L85}{— sinistra 31, 101.}{}
\EdL{86}{AK07-J0287-L86}{Radius des Erdschattens s. Durchmesser.}{}
\EdL{87}{AK07-J0287-L87}{Radix XVI, 7, 33, 36, 41, 81, T. 3, 4, 6, 9, etc.}{}
\EdL{88}{AK07-J0287-L88}{al-Raḳḳa 41, 74, 75, 81, 96.}{}
\EdL{89}{AK07-J0287-L89}{Rechtläufigkeit der Planeten, s. Vorwärts-}{}
\EdL{90}{AK07-J0287-L90}{schreiten.}{}
\EdL{91}{AK07-J0287-L91}{Rectascension s. Ascensiones rectae.}{}
\EdL{92}{AK07-J0287-L92}{Regel der vier Grössen 61.}{}
\EdL{93}{AK07-J0287-L93}{Regenten s. Judices.}{}
\EdL{94}{AK07-J0287-L94}{Regiomontanus, Joh. 99, 101.}{}
\EdL{95}{AK07-J0287-L95}{Regnum s. Erhöhung.}{}
\EdL{96}{AK07-J0287-L96}{Respiratio 26, 28, 87, 89.}{}
\EdL{97}{AK07-J0287-L97}{Revolutiones annorum s. Umwandlung.}{}
\EdL{98}{AK07-J0287-L98}{Robert von Chester VIII, XI—XIII, XVIII, 1,}{}
\EdL{99}{AK07-J0287-L99}{69, 106.}{}
\FrameHead{0288}{p. 254}{0288}{p. 254}
\ColHead{left column}
\EdL{1}{AK07-J0288-L01}{Robertus Ketenensis od. Retinensis s. Robert}{}
\EdL{2}{AK07-J0288-L02}{von Chester.}{}
\EdL{3}{AK07-J0288-L03}{Rodbertus Cestrensis s Robert von Chester.}{}
\EdL{4}{AK07-J0288-L04}{Roger von Hereford XII, 106.}{}
\EdL{5}{AK07-J0288-L05}{Rose, astrologische 102—105.}{}
\EdL{6}{AK07-J0288-L06}{Rota = circumferentia, peripheria, sphaera 4,}{}
\EdL{7}{AK07-J0288-L07}{7\,; s. auch Durchmesser.}{}
\EdL{8}{AK07-J0288-L08}{— draconis s. Durchmesser des Erdschattens.}{}
\EdL{9}{AK07-J0288-L09}{Rudolf von Brügge XIII, XVIII.}{}
\EdL{10}{AK07-J0288-L10}{Rückläufigkeit s. das folg.}{}
\EdL{11}{AK07-J0288-L11}{Rückwärtsschreiten, der Planeten 12, 55, 56\,;}{}
\EdL{12}{AK07-J0288-L12}{Zeit desselben 12, 13, 56—58.}{}
\EdL{13}{AK07-J0288-L13}{Rumor (?) 31.}{}
\EdL{14}{AK07-J0288-L14}{\textbf{S.}}{}
\EdL{15}{AK07-J0288-L15}{Sachau, E. 37, 38, 49, 66, 79, 105, 235.}{}
\EdL{16}{AK07-J0288-L16}{Savasorda X.}{}
\EdL{17}{AK07-J0288-L17}{Schaden s. Alienatio.}{}
\EdL{18}{AK07-J0288-L18}{Schaltjahre der Araber 34, 37\,; der julian. Zeit-}{}
\EdL{19}{AK07-J0288-L19}{rechnung s. Annus bissextus.}{}
\EdL{20}{AK07-J0288-L20}{Schatten s. Umbra.}{}
\EdL{21}{AK07-J0288-L21}{Schiaparelli, Giov. XXIV, 66, 80—83.}{}
\EdL{22}{AK07-J0288-L22}{Schiefe der Ekliptik s. Ekliptik.}{}
\EdL{23}{AK07-J0288-L23}{— der Mondbahn s. Mondbahn.}{}
\EdL{24}{AK07-J0288-L24}{Schram, R. 83, 85 (N. 2).}{}
\EdL{25}{AK07-J0288-L25}{Semita numerorum (Kolonne der Zahlen) 69,}{}
\EdL{26}{AK07-J0288-L26}{78, T. 21, 22 etc., 58, 58\textsuperscript{a}, 61—66.}{}
\EdL{27}{AK07-J0288-L27}{Sened b. ʿAlī 32.}{}
\EdL{28}{AK07-J0288-L28}{Servitus s. Tiefstand.}{}
\EdL{29}{AK07-J0288-L29}{Sexagesimalteile s. Multiplikationstafel.}{}
\EdL{30}{AK07-J0288-L30}{Sextilschein (Exagonum) 30, 31, 98—101, T.}{}
\EdL{31}{AK07-J0288-L31}{91—114.}{}
\EdL{32}{AK07-J0288-L32}{Seybold, C. F. XXI.}{}
\EdL{33}{AK07-J0288-L33}{Siddhāntas VIII, 32, 59, 65, 70, 105, 106\,; Ārya-S.}{}
\EdL{34}{AK07-J0288-L34}{65, 66\,; Brahma-S. 32, 65, 235\,; Śiromani-S. 70\,;}{}
\EdL{35}{AK07-J0288-L35}{Sūrya-S. 32, 33, 42, 43, 45, 46, 51, 52, 58—60,}{}
\EdL{36}{AK07-J0288-L36}{62, 64, 65, 68, 74, 75, 78—80, 86, 90, 92.}{}
\EdL{37}{AK07-J0288-L37}{Signa s. Notae.}{}
\EdL{38}{AK07-J0288-L38}{Signum, Tierkreiszeichen 14, 15, 41, 71\,; Auf-}{}
\EdL{39}{AK07-J0288-L39}{gänge derselben in der geraden Sphäre}{}
\EdL{40}{AK07-J0288-L40}{18—20, 72, 73, T. 59, 59\textsuperscript{a}\,; s. auch Ascensiones}{}
\EdL{41}{AK07-J0288-L41}{rectae\,; in der schiefen Sphäre 20, 21, 73—76\,;}{}
\EdL{42}{AK07-J0288-L42}{s. auch Ascensiones obliquae.}{}
\EdL{43}{AK07-J0288-L43}{Sindhind 32, 33, 43, 46, 49, 70, 103, 235.}{}
\EdL{44}{AK07-J0288-L44}{Sinus, Auffindung desselben 17, 18, 68—71, T.}{}
\EdL{45}{AK07-J0288-L45}{58 und 58\textsuperscript{a}.}{}
\EdL{46}{AK07-J0288-L46}{— der Ergänzung (Cosinus) 73, 77.}{}
\EdL{47}{AK07-J0288-L47}{— diminutus (versus) 17, 69, 71.}{}
\EdL{48}{AK07-J0288-L48}{— planus (rectus) 17, 21, 22, 69, 71, 77.}{}
\EdL{49}{AK07-J0288-L49}{Sonne s. Bewegung, Länge, Sublimatio, Dekli-}{}
\EdL{50}{AK07-J0288-L50}{nation, Durchmesser, Examinatio, Verfin-}{}
\EdL{51}{AK07-J0288-L51}{sterung.}{}
\ColHead{right column}
\EdL{52}{AK07-J0288-L52}{Spatium longinquum (remotum) 26—28, 86 87,}{}
\EdL{53}{AK07-J0288-L53}{94, 95.}{}
\EdL{54}{AK07-J0288-L54}{— propinquum 26—28, 86, 87, 94, 95.}{}
\EdL{55}{AK07-J0288-L55}{Stationen, der Planeten 12, 13, 55—58, T. 27—56.}{}
\EdL{56}{AK07-J0288-L56}{Steinschneider, M. XIII, 41.}{}
\EdL{57}{AK07-J0288-L57}{Sternjahr s. Jahr, siderisches.}{}
\EdL{58}{AK07-J0288-L58}{Stillstand der Planeten, s. Stationen.}{}
\EdL{59}{AK07-J0288-L59}{Stunden, ungleiche, s. Zeitstunden.}{}
\EdL{60}{AK07-J0288-L60}{Stundenwinkel 92, 93.}{}
\EdL{61}{AK07-J0288-L61}{Stupor 22—24, 80.}{}
\EdL{62}{AK07-J0288-L62}{Sublimatio (Apogeum) des Mondes 14, 15, 64,}{}
\EdL{63}{AK07-J0288-L63}{88\,; der Planeten 14, 15, 46—50, 52, 53, 60, 63,}{}
\EdL{64}{AK07-J0288-L64}{64, 235\,; der Sonne 9, 14, 33, 43, 44, 46, 47,}{}
\EdL{65}{AK07-J0288-L65}{67, 82.}{}
\EdL{66}{AK07-J0288-L66}{— examinata vel definita 10, 11, 49, 50, 53, 54,}{}
\EdL{67}{AK07-J0288-L67}{T. 27—56.}{}
\EdL{68}{AK07-J0288-L68}{Submissio (Perigeum) 14, 15, 47, 63, 64, 88.}{}
\EdL{69}{AK07-J0288-L69}{Sullimatio s. Sublimatio.}{}
\EdL{70}{AK07-J0288-L70}{Syzygien s. Konjunktion und Opposition.}{}
\EdL{71}{AK07-J0288-L71}{\textbf{T.}}{}
\EdL{72}{AK07-J0288-L72}{Tabula (vel pagina) differentiarum 21, 76.}{}
\EdL{73}{AK07-J0288-L73}{— correctionis 95.}{}
\EdL{74}{AK07-J0288-L74}{Tabulae proportionis 26, 27, 88, 94, T. 73—75.}{}
\EdL{75}{AK07-J0288-L75}{— spuriae (al-Battānīs) XV, 40, 102.}{}
\EdL{76}{AK07-J0288-L76}{Tafeln (astron.)\,: die alphonsinischen 35\,; die}{}
\EdL{77}{AK07-J0288-L77}{des Battānī s. al-Battānī\,; die erprobten 32,}{}
\EdL{78}{AK07-J0288-L78}{46, 47, 58\,; die persischen (des Shâh) 32, 45,}{}
\EdL{79}{AK07-J0288-L79}{49, 51, 52, 105, 106, 235\,; die nach dem Sind-}{}
\EdL{80}{AK07-J0288-L80}{hind verfassten VIII, 70; die Theon’schen 94.}{}
\EdL{81}{AK07-J0288-L81}{Tagebogen des Ascendenten 97.}{}
\EdL{82}{AK07-J0288-L82}{Tagesstunden 20, 25, 76, 97, 98.}{}
\EdL{83}{AK07-J0288-L83}{Tangens s. Umbra.}{}
\EdL{84}{AK07-J0288-L84}{Tempora = Aequatorialgrade 97.}{}
\EdL{85}{AK07-J0288-L85}{Termini s. Grenzen.}{}
\EdL{86}{AK07-J0288-L86}{Tetragonum s. Quadratschein.}{}
\EdL{87}{AK07-J0288-L87}{Thābit b. Ḳurra XI.}{}
\EdL{88}{AK07-J0288-L88}{Thebit s. Thābit.}{}
\EdL{89}{AK07-J0288-L89}{Thibaut, G. 42, 46, 64, 68, 75, 90, 235.}{}
\EdL{90}{AK07-J0288-L90}{Thoth (ägypt. Gottheit) XXII.}{}
\EdL{91}{AK07-J0288-L91}{Tiefstand 104, T. 116.}{}
\EdL{92}{AK07-J0288-L92}{Tierkreiszeichen s. Signum.}{}
\EdL{93}{AK07-J0288-L93}{Trepidation der Fixsterne 42.}{}
\EdL{94}{AK07-J0288-L94}{Trigonalia s. Triplizitäten.}{}
\EdL{95}{AK07-J0288-L95}{Trigonalschein 30, 31, 98, 99, 101, T. 91—114.}{}
\EdL{96}{AK07-J0288-L96}{Trigonum s. das vorhergeh.}{}
\EdL{97}{AK07-J0288-L97}{Triplizitäten 104.}{}
\EdL{98}{AK07-J0288-L98}{\textbf{U.}}{}
\EdL{99}{AK07-J0288-L99}{Udjain s. Arin.}{}
\EdL{100}{AK07-J0288-L100}{Uhlemann 105.}{}
\FrameHead{0289}{p. 255}{0289}{p. 255}
\ColHead{left column}
\EdL{1}{AK07-J0289-L01}{Ujjayinī s. Arin.}{}
\EdL{2}{AK07-J0289-L02}{Ulūgh Beg 78.}{}
\EdL{3}{AK07-J0289-L03}{Umbra meridiana et aequinoctialis 20, 21, 74, 76.}{}
\EdL{4}{AK07-J0289-L04}{— mankūs (versa) = Tangens 22, 76—78.}{}
\EdL{5}{AK07-J0289-L05}{— mustewia (plana) = Cotangens 21, 22, 68,}{}
\EdL{6}{AK07-J0289-L06}{76—78, 106, T. 60.}{}
\EdL{7}{AK07-J0289-L07}{— terrae 26, 86—91\,; s. auch Durchmesser des}{}
\EdL{8}{AK07-J0289-L08}{Erdschattens.}{}
\EdL{9}{AK07-J0289-L09}{Umwandlung (conversio, revolutio) der Ge-}{}
\EdL{10}{AK07-J0289-L10}{burtsjahre (annorum nativitatum) 102, 103,}{}
\EdL{11}{AK07-J0289-L11}{106, T. 115.}{}
\EdL{12}{AK07-J0289-L12}{— der Weltjahre (annorum mundi) 103.}{}
\EdL{13}{AK07-J0289-L13}{\textbf{V.}}{}
\EdL{14}{AK07-J0289-L14}{Varāha-Mihira 46, 235.}{}
\EdL{15}{AK07-J0289-L15}{Verfinsterung des Mondes 26—28, 86‐‐90, T.}{}
\EdL{16}{AK07-J0289-L16}{73—76.}{}
\EdL{17}{AK07-J0289-L17}{— der Sonne 28, 29, 94—96, T. 78.}{}
\EdL{18}{AK07-J0289-L18}{Verteilung der 12 astrol. Häuser, s. Häuser.}{}
\EdL{19}{AK07-J0289-L19}{Vollmond 84.}{}
\EdL{20}{AK07-J0289-L20}{Vorwärtsschreiten der Planeten 12, 55, 56\,;}{}
\EdL{21}{AK07-J0289-L21}{Zeit desselben 12, 13, 56—58.}{}
\ColHead{right column}
\EdL{22}{AK07-J0289-L22}{\textbf{W.}}{}
\EdL{23}{AK07-J0289-L23}{Washington, Meridian von, 52.}{}
\EdL{24}{AK07-J0289-L24}{Winkel der Breite und Länge 91, 92.}{}
\EdL{25}{AK07-J0289-L25}{Wüstenfeld, F. XIII.}{}
\EdL{26}{AK07-J0289-L26}{\textbf{Y.}}{}
\EdL{27}{AK07-J0289-L27}{Yaḥyā b. Abī Manṣūr 32.}{}
\EdL{28}{AK07-J0289-L28}{Yaʿḳūb b. Ṭāriḳ 32, 70.}{}
\EdL{29}{AK07-J0289-L29}{Yezdegird III., letzter König der Sassaniden,}{}
\EdL{30}{AK07-J0289-L30}{VII, VIII, XV, 33, 35, 37, 106, T. 1, s. auch}{}
\EdL{31}{AK07-J0289-L31}{Aera.}{}
\EdL{32}{AK07-J0289-L32}{Ysidorus s. Isidor.}{}
\EdL{33}{AK07-J0289-L33}{Yojana 91.}{}
\EdL{34}{AK07-J0289-L34}{\textbf{Z.}}{}
\EdL{35}{AK07-J0289-L35}{al-Zarḳālī 236.}{}
\EdL{36}{AK07-J0289-L36}{Zeit, mittlere 80.}{}
\EdL{37}{AK07-J0289-L37}{— wahre 80.}{}
\EdL{38}{AK07-J0289-L38}{Zeitgleichung 25, 80, 84, 106, T. 67—68.}{}
\EdL{39}{AK07-J0289-L39}{Zeitstunden (ungleiche) 20, 21, 25, 76, 81, 84,}{}
\EdL{40}{AK07-J0289-L40}{97, 98\,; s. auch Nachtstunden und Tages-}{}
\EdL{41}{AK07-J0289-L41}{stunden.}{}
\EdL{42}{AK07-J0289-L42}{Zenitdistanz (des Mondes) 92.}{}
\EdL{43}{AK07-J0289-L43}{Zentrum s. Centrum.}{}
\EdL{44}{AK07-J0289-L44}{Zeuthen, H. G. VII, 68, 71.}{}
\EdL{45}{AK07-J0289-L45}{Zodiacus 29\,; s. auch Ekliptik.}{}
\FrameHead{0290}{p. [256]}{0290}{p. [256]}
\Withheld{blank}
\FrameHead{0291}{blank}{0291}{blank}
\Withheld{blank}
\FrameHead{0292}{blank}{0292}{blank}
\Withheld{blank}
\FrameHead{0293}{library matter}{0293}{library matter}
\Withheld{library matter only: A tick in pencil at the upper left.; Nothing of the edition is printed on this page. The pastedown of the rear board: the edges of the cloth turned over the board show through the paper along the sides and at the corners.}
\FrameHead{0294}{binding}{0294}{binding}
\Withheld{the binding: The back cover of the binding, of dark, textured cloth like the front cover (JP2 0001), without lettering; photographed in low light, at the size of the image of the front cover.}
\FrameHead{0295}{matter of the digitisation}{0295}{matter of the digitisation}
\Withheld{matter of the digitisation: The colour card of the digitisation: a chart of 24 colour patches («24ColorCard») in a holder labelled «APRIL2013», a white patch below it, a ruler in inches at the left and another along the foot, and a card marked «RIGHT» holding a printed form in Chinese and English (the fields «Name», «Post», «Unit», «No:» and «Date:» and an empty box for a photograph), on a black ground. Not part of the book; the navigation PDF of the Internet Archive has no page for this image, as for the colour card before the book (JP2 0000).}
\clearpage\pdfbookmark[1]{critical notes}{critical}{\large\bfseries Critical notes}\par\medskip\noindent These notes are not part of the transcription. They record errors of the print and propose corrections; the transcription keeps the print.\par\medskip
\noindent\textbf{AK07-CN001}\enspace\hyperlink{AK07-J0008-L01}{AK07-J0008-L01}\enspace(wrong punctuation)\par\noindent printed: Pag,\enspace proposed: Pag.\par\noindent the head of the column of pages, the abbreviation of Pagina, has a full stop on the first page of the contents («Pag.», JP2 0007, L02 and L13) and a comma here\par\noindent confidence: high\par\noindent impact: none\par\medskip
\noindent\textbf{AK07-CN002}\enspace\hyperlink{AK07-J0015-N02}{AK07-J0015-N02}\enspace(turned letter)\par\noindent printed: vərschiedener\enspace proposed: verschiedener\par\noindent the first e of the word is printed upside down, a turned sort; the word is «verschiedener», of different hands, as «verschiedene Hände» in L03\par\noindent confidence: high\par\noindent impact: none\par\medskip
\noindent\textbf{AK07-CN003}\enspace\hyperlink{AK07-J0018-L37}{AK07-J0018-L37}\enspace(misprint)\par\noindent printed: Zuzätze\enspace proposed: Zusätze\par\noindent the word is «Zusätze», additions, the additions and marginal glosses of N that the page goes on to list; the second z is a misprint for s\par\noindent confidence: high\par\noindent impact: none\par\medskip
\noindent\textbf{AK07-CN004}\enspace\hyperlink{AK07-J0025-L06}{AK07-J0025-L06}\enspace(misprint)\par\noindent printed: bestimmen\enspace proposed: bestimmten\par\noindent «unter einem bestimmen Stichwort», under a particular headword: the participle is «bestimmten»; the t has dropped out\par\noindent confidence: high\par\noindent impact: none\par\medskip
\noindent\textbf{AK07-CN005}\enspace\hyperlink{AK07-J0025-L41}{AK07-J0025-L41}\enspace(misprint)\par\noindent printed: griechiche\enspace proposed: griechische\par\noindent «der griechiche Hermes», the Greek Hermes: the adjective is «griechische»; the s before ch has dropped out\par\noindent confidence: high\par\noindent impact: none\par\medskip
\noindent\textbf{AK07-CN006}\enspace\hyperlink{AK07-J0026-L08}{AK07-J0026-L08}\enspace(stray point)\par\noindent printed: Wort\textsuperscript{1}). .\enspace proposed: Wort\textsuperscript{1}).\par\noindent a second point, printed after the full stop that ends the paragraph, has no function; a raised spacing sort or a point that has slipped in the forme\par\noindent confidence: medium\par\noindent impact: none\par\medskip
\noindent\textbf{AK07-CN007}\enspace\hyperlink{AK07-J0027-L24}{AK07-J0027-L24}\enspace(missing point)\par\noindent printed: z. B\enspace proposed: z. B.\par\noindent the abbreviation of «zum Beispiel» is printed «z. B.» elsewhere (N06 of this page); the point after the B, at the end of the line, is not printed\par\noindent confidence: high\par\noindent impact: none\par\medskip
\noindent\textbf{AK07-CN048}\enspace\hyperlink{AK07-J0065-L35}{AK07-J0065-L35}\enspace(wrong punctuation)\par\noindent printed: der Name „Ptolomaeus“,\enspace proposed: der Name „Ptolomaeus“.\par\noindent the sentence ends at the name: the next line begins a new sentence with a capital, «Man sieht also»; a comma is printed after the closing quotation mark\par\noindent confidence: high\par\noindent impact: none\par\medskip
\noindent\textbf{AK07-CN049}\enspace\hyperlink{AK07-J0067-L11}{AK07-J0067-L11}\enspace(misprint)\par\noindent printed: Das dieser Zusatz\enspace proposed: Dass dieser Zusatz\par\noindent the clause «dieser Zusatz … nicht gestanden, auch von Maslama nicht hinzugefügt worden sein kann» is the subject of «beweisen gewisse Sätze desselben» (L13) and needs the conjunction «dass», as in «dass die Mss. CMN» (L17); one s is printed\par\noindent confidence: high\par\noindent impact: none\par\medskip
\noindent\textbf{AK07-CN050}\enspace\hyperlink{AK07-J0073-L23}{AK07-J0073-L23}\enspace(turned letter)\par\noindent printed: Beweguug\enspace proposed: Bewegung\par\noindent «die mittlere Bewegung», as in L19 and the heading of chapter 7 (L08); the n before the final g is printed as u, a turned n\par\noindent confidence: high\par\noindent impact: none\par\medskip
\noindent\textbf{AK07-CN051}\enspace\hyperlink{AK07-J0081-L27}{AK07-J0081-L27}\enspace(damaged letter)\par\noindent printed: deı\enspace proposed: der\par\noindent «die Gründe der Halbierung der Gleichung»: the article in the genitive; the r is printed without its arm, so that the word reads «deı»\par\noindent confidence: high\par\noindent impact: none\par\medskip
\noindent\textbf{AK07-CN052}\enspace\hyperlink{AK07-J0081-L35}{AK07-J0081-L35}\enspace(missing accent)\par\noindent printed: précisement\enspace proposed: précisément\par\noindent the same sentence of Chasles, quoted again on p. 50 (JP2 0082, L02), has «précisément», the usual French spelling; here the acute accent on the second e is missing\par\noindent confidence: high\par\noindent impact: none\par\medskip
\noindent\textbf{AK07-CN282}\enspace\hyperlink{AK07-J0081-L29}{AK07-J0081-L29}\enspace(wrong numeral)\par\noindent printed: (p. I der Einleitung)\enspace proposed: (p. VII der Einleitung)\par\noindent the list of misprints, p. [237] (JP2 0271, L11), reads «p. VII» for «p. I»; the introduction is printed on pp. VII-XXV (contents, p. [III], JP2 0007)\par\noindent confidence: high\par\noindent impact: none\par\medskip
\noindent\textbf{AK07-CN053}\enspace\hyperlink{AK07-J0089-N02}{AK07-J0089-N02}\enspace(misprint)\par\noindent printed: berüchsichtigt\enspace proposed: berücksichtigt\par\noindent the verb is «berücksichtigen», «Diesen Fall berücksichtigt Batt. am Schlusse des 46. Kap.» (p. 56, JP2 0088, N06); an h is printed for the k\par\noindent confidence: high\par\noindent impact: none\par\medskip
\noindent\textbf{AK07-CN054}\enspace\hyperlink{AK07-J0090-L33}{AK07-J0090-L33}\enspace(damaged letter)\par\noindent printed: Dıachenkopfes\enspace proposed: Drachenkopfes\par\noindent «die mittlere Länge des Drachenkopfes», as «Drachenkopf» in L31 and L36; the r is printed without its arm, as the r of «der» on p. 49 (JP2 0081, L27)\par\noindent confidence: high\par\noindent impact: none\par\medskip
\noindent\textbf{AK07-CN055}\enspace\hyperlink{AK07-J0092-L16}{AK07-J0092-L16}\enspace(raised point)\par\noindent printed: ist· Die\enspace proposed: ist. Die\par\noindent the sentence ends after «abhängig ist», and «Die Breite» begins the next one; the full stop is printed at mid-height instead of on the base line\par\noindent confidence: high\par\noindent impact: none\par\medskip
\noindent\textbf{AK07-CN056}\enspace\hyperlink{AK07-J0094-L02}{AK07-J0094-L02}\enspace(missing point)\par\noindent printed: Breite in der wahren Dist r′\enspace proposed: Breite in der wahren Dist. r′\par\noindent the abbreviation has its point in the numerator of the same fraction, «Breite in der mittleren Dist. r», and on p. 61 (JP2 0093, L04-L08)\par\noindent confidence: high\par\noindent impact: none\par\medskip
\noindent\textbf{AK07-CN057}\enspace\hyperlink{AK07-J0094-L37}{AK07-J0094-L37}\enspace(wrong signs)\par\noindent printed: 3° 48″ 27′\enspace proposed: 3° 48′ 27″\par\noindent the minutes and the seconds have exchanged their signs: the next line divides 3° 48½′, that is 3° 48′ 30″, rounded from 3° 48′ 27″\par\noindent confidence: high\par\noindent impact: none on the result\par\medskip
\noindent\textbf{AK07-CN058}\enspace\hyperlink{AK07-J0096-L38}{AK07-J0096-L38}\enspace(missing sign)\par\noindent printed: 21° 10.\enspace proposed: 21° 10′.\par\noindent the other longitudes of the list have the sign of the minutes, «103° 12′», «82° 1′», and the Latin gives «dakaicae X» (p. 15, JP2 0047, L12); the prime after 10 is not printed\par\noindent confidence: high\par\noindent impact: none\par\medskip
\noindent\textbf{AK07-CN283}\enspace\hyperlink{AK07-J0098-N01}{AK07-J0098-N01}\enspace(wrong reference)\par\noindent printed: Vergl. auch Kap. 38.\enspace proposed: Vergl. auch p. 103.\par\noindent the list of misprints, p. [237] (JP2 0271, L14), reads «p. 103» for «Kap. 38»: the text has 37 chapters, and on p. 103 (JP2 0135, L09-L15) Suter takes up the length of the year again, with a note back to pp. 42 and 65\par\noindent confidence: high\par\noindent impact: none\par\medskip
\noindent\textbf{AK07-CN059}\enspace\hyperlink{AK07-J0100-L03}{AK07-J0100-L03}\enspace(wrong sign)\par\noindent printed: (9° 37″)\enspace proposed: (9° 37′)\par\noindent the difference of the two true longitudes, 332° 4′ − 322° 27′, is 9° 37′, as p. 67 (JP2 0099, L39) gives it; the sign of the seconds is printed for that of the minutes\par\noindent confidence: high\par\noindent impact: none\par\medskip
\noindent\textbf{AK07-CN060}\enspace\hyperlink{AK07-J0109-N01}{AK07-J0109-N01}\enspace(stray point)\par\noindent printed: 22—23. u. 180—182.\enspace proposed: 22—23 u. 180—182.\par\noindent the two page ranges are joined by «u.», as in «27 u. 187—189» (p. 74, JP2 0106, N02) and «67 u. 243» (p. 38, JP2 0070, N01); a full stop is printed after 23\par\noindent confidence: high\par\noindent impact: none\par\medskip
\noindent\textbf{AK07-CN061}\enspace\hyperlink{AK07-J0110-N05}{AK07-J0110-N05}\enspace(missing quotation mark)\par\noindent printed: complete revolution.\enspace proposed: complete revolution.“\par\noindent the quotation opened by „ before «When» in N04 is not closed; the other quotations of the commentary close with “, as «Funktionen ab.“» (p. 77, JP2 0109, L32) and «Cotangente ist“» (L01)\par\noindent confidence: medium\par\noindent impact: none\par\medskip
\noindent\textbf{AK07-CN062}\enspace\hyperlink{AK07-J0112-L02}{AK07-J0112-L02}\enspace(misprint)\par\noindent printed: am bestem\enspace proposed: am besten\par\noindent the superlative after «am» ends in -en, «am besten», as «am meisten» (L19); the n is printed as m\par\noindent confidence: high\par\noindent impact: none\par\medskip
\noindent\textbf{AK07-CN063}\enspace\hyperlink{AK07-J0115-L01}{AK07-J0115-L01}\enspace(misprint)\par\noindent printed: Der Mond is also\enspace proposed: Der Mond ist also\par\noindent the verb of the sentence is «ist», as in «ist also» (L18) and «enthalten ist» (L05); the t is not printed\par\noindent confidence: high\par\noindent impact: none\par\medskip
\noindent\textbf{AK07-CN064}\enspace\hyperlink{AK07-J0115-N07}{AK07-J0115-N07}\enspace(wrong sort)\par\noindent printed: flndet\enspace proposed: findet\par\noindent the verb is «findet»; the ligature fl, whose second stroke rises to the height of the d, is printed in the place of the ligature fi\par\noindent confidence: high\par\noindent impact: none\par\medskip
\noindent\textbf{AK07-CN065}\enspace\hyperlink{AK07-J0118-N04}{AK07-J0118-N04}\enspace(misprint)\par\noindent printed: in diesen Sinne\enspace proposed: in diesem Sinne\par\noindent the dative after «in» in the sense «accordingly», «in diesem Sinne», as «in diesem Kapitel» (N02) and «in diesem Falle» (L27); the m of the ending is printed as n\par\noindent confidence: high\par\noindent impact: none\par\medskip
\noindent\textbf{AK07-CN066}\enspace\hyperlink{AK07-J0121-L05}{AK07-J0121-L05}\enspace(misprint)\par\noindent printed: Zuzatz\enspace proposed: Zusatz\par\noindent «einen Zusatz zu diesem Kapitel» (L04), the same word; the s is printed as z, as in «Zuzätze» on p. XIV (JP2 0018, L37, AK07-CN003)\par\noindent confidence: high\par\noindent impact: none\par\medskip
\noindent\textbf{AK07-CN067}\enspace\hyperlink{AK07-J0122-L18}{AK07-J0122-L18}\enspace(misprint)\par\noindent printed: den Radius der Erdschattens\enspace proposed: den Radius des Erdschattens\par\noindent the genitive of «der Erdschatten» is «des Erdschattens», as in L01, L06, L22, L28, L30 and L37; «der» is printed for «des»\par\noindent confidence: high\par\noindent impact: none\par\medskip
\noindent\textbf{AK07-CN068}\enspace\hyperlink{AK07-J0122-L25}{AK07-J0122-L25}\enspace(misprint)\par\noindent printed: die entsprechenden Zahl\enspace proposed: die entsprechende Zahl\par\noindent the nominative singular after «die», as «die entsprechende Zahl der 7. Kolonne» (L26-L27); an n too many\par\noindent confidence: high\par\noindent impact: none\par\medskip
\noindent\textbf{AK07-CN069}\enspace\hyperlink{AK07-J0122-L26}{AK07-J0122-L26}\enspace(wrong sort)\par\noindent printed: multlpliziert\enspace proposed: multipliziert\par\noindent «multipliziert», as in L24; the second i is printed as l\par\noindent confidence: high\par\noindent impact: none\par\medskip
\noindent\textbf{AK07-CN070}\enspace\hyperlink{AK07-J0123-L01}{AK07-J0123-L01}\enspace(misprint)\par\noindent printed: wandschaft\enspace proposed: wandtschaft\par\noindent «Ver-wandschaft», divided across the pages: the noun of «verwandt» is «Verwandtschaft»; the t after the d is not printed\par\noindent confidence: medium\par\noindent impact: none\par\medskip
\noindent\textbf{AK07-CN071}\enspace\hyperlink{AK07-J0123-L11}{AK07-J0123-L11}\enspace(misprint)\par\noindent printed: Die Regel Kh.’s … sind\enspace proposed: Die Regeln Kh.’s … sind\par\noindent the verb «sind» and «besonders die erstere», the rule for the longitude as against that for the latitude, ask for the plural; the n is not printed\par\noindent confidence: medium\par\noindent impact: none\par\medskip
\noindent\textbf{AK07-CN072}\enspace\hyperlink{AK07-J0124-L22}{AK07-J0124-L22}\enspace(misprint)\par\noindent printed: Länge, zo bezeichnen\enspace proposed: Länge, so bezeichnen\par\noindent the correlative «so» after the conditional clause «Handelt es sich um …», as in «Breite, so bedeuten sie» (L24-L25); the s is printed as z\par\noindent confidence: high\par\noindent impact: none\par\medskip
\noindent\textbf{AK07-CN073}\enspace\hyperlink{AK07-J0125-L32}{AK07-J0125-L32}\enspace(missing note number)\par\noindent printed: addieren, um\enspace proposed: addieren\textsuperscript{2}), um\par\noindent the note 2 at the foot of the page, «Denn der Mond befindet sich zwischen Ascendens und Medium Coelum», gives the reason why the parallax is added to the true place of the moon; no reference to it is printed in the text\par\noindent confidence: medium\par\noindent impact: none on the text\par\medskip
\noindent\textbf{AK07-CN074}\enspace\hyperlink{AK07-J0125-L34}{AK07-J0125-L34}\enspace(turned letter)\par\noindent printed: Auffinduug\enspace proposed: Auffindung\par\noindent «die Regel für die Auffindung», as «die Regel für die Auffindung des Medium Coelum» (p. 92, JP2 0124, L38-L39); the n before the g is printed as u, a turned n\par\noindent confidence: high\par\noindent impact: none\par\medskip
\noindent\textbf{AK07-CN075}\enspace\hyperlink{AK07-J0127-L05}{AK07-J0127-L05}\enspace(misprint)\par\noindent printed: der entsprechende Werte\enspace proposed: der entsprechenden Werte\par\noindent after the article of the genitive plural the adjective takes the ending -en, as in «das 5-fache der entsprechenden Werte» (p. 90, JP2 0122, L12); the final n is not printed\par\noindent confidence: high\par\noindent impact: none\par\medskip
\noindent\textbf{AK07-CN076}\enspace\hyperlink{AK07-J0129-L12}{AK07-J0129-L12}\enspace(missing parenthesis)\par\noindent printed: III. bezw. IX.) Hauses\enspace proposed: III. (bezw. IX.) Hauses\par\noindent the closing parenthesis after «IX.» has no opening one; «des III. (bezw. IX.)» in L09, «des IV. (bezw. X.)» and «des V. (bezw. XI.)» in L13\par\noindent confidence: high\par\noindent impact: none\par\medskip
\noindent\textbf{AK07-CN077}\enspace\hyperlink{AK07-J0131-L22}{AK07-J0131-L22}\enspace(misprint)\par\noindent printed: Rectascention\enspace proposed: Rectascension\par\noindent the word is spelled «Rectascension» throughout, as in L21 at the end of the line above; the s before the i is printed as t\par\noindent confidence: high\par\noindent impact: none\par\medskip
\noindent\textbf{AK07-CN078}\enspace\hyperlink{AK07-J0132-L35}{AK07-J0132-L35}\enspace(wrong sort)\par\noindent printed: überzchrieben\enspace proposed: überschrieben\par\noindent «überschrieben», as in L37; the s is printed as z\par\noindent confidence: high\par\noindent impact: none\par\medskip
\noindent\textbf{AK07-CN079}\enspace\hyperlink{AK07-J0132-N03}{AK07-J0132-N03}\enspace(missing hyphen)\par\noindent printed: Diffe\enspace proposed: Diffe-\par\noindent «Differenz» is divided between N03 and N04 without the hyphen of the division\par\noindent confidence: high\par\noindent impact: none\par\medskip
\noindent\textbf{AK07-CN284}\enspace\hyperlink{AK07-J0132-L09}{AK07-J0132-L09}\enspace(wrong word)\par\noindent printed: die ganze Rectascension\enspace proposed: die ganze Ascensio obliqua\par\noindent the list of misprints, p. [237] (JP2 0271, L16), reads «Ascensio obliqua» for «Rectascension» in L09 and L10; footnote 2 (N02-N04), whose numbers are right ascensions, is left as it stands\par\noindent confidence: high\par\noindent impact: none\par\medskip
\noindent\textbf{AK07-CN285}\enspace\hyperlink{AK07-J0132-L10}{AK07-J0132-L10}\enspace(wrong word)\par\noindent printed: nur die Rectascension\enspace proposed: nur die Ascensio obliqua\par\noindent the list of misprints, p. [237] (JP2 0271, L16), reads «Ascensio obliqua» for «Rectascension» in L09 and L10 (critical note AK07-CN284)\par\noindent confidence: high\par\noindent impact: none\par\medskip
\noindent\textbf{AK07-CN080}\enspace\hyperlink{AK07-J0134-L18}{AK07-J0134-L18}\enspace(missing note number)\par\noindent printed: Batt.’s befindet\enspace proposed: Batt.’s\textsuperscript{1}) befindet\par\noindent the note 1, «Ed. Nallino, II, 299», the page of the «rosa astrologica» among the «Tabulae spuriae», has no reference in the text; it belongs to the second table «als astrologische Rose» among the «Tabulae spuriae Batt.’s»\par\noindent confidence: medium\par\noindent impact: none on the text\par\medskip
\noindent\textbf{AK07-CN081}\enspace\hyperlink{AK07-J0136-L29}{AK07-J0136-L29}\enspace(misprint)\par\noindent printed: enhält\enspace proposed: enthält\par\noindent «enthält», as in «Die astrologische Rose enthält die Facies» (L38-L39); the t is not printed\par\noindent confidence: high\par\noindent impact: none\par\medskip
\noindent\textbf{AK07-CN082}\enspace\hyperlink{AK07-J0137-L35}{AK07-J0137-L35}\enspace(stray comma)\par\noindent printed: Kh., mehr\enspace proposed: Kh. mehr\par\noindent the comma divides the subject «Kh.» of the clause «Dass sich Kh. … gehalten haben mag» from the rest of it; no comma belongs there\par\noindent confidence: medium\par\noindent impact: none\par\medskip
\noindent\textbf{AK07-CN274}\enspace\hyperlink{AK07-J0265-L24}{AK07-J0265-L24}\enspace(misprint)\par\noindent printed: warhscheinlich\enspace proposed: wahrscheinlich\par\noindent the h and the r are exchanged; «wahrscheinlich» elsewhere in Suter's text, as on p. 102 (JP2 0134, L14)\par\noindent confidence: high\par\noindent impact: none\par\medskip
\noindent\textbf{AK07-CN275}\enspace\hyperlink{AK07-J0268-L17}{AK07-J0268-L17}\enspace(wrong letter)\par\noindent printed: \Sperr{Kb.’s}\enspace proposed: \Sperr{Kh.’s}\par\noindent Suter's abbreviation of al-Khwārizmī's name is «Kh.» throughout, as «Kh.’s» on p. 235 (JP2 0267, L02); a b is printed for the h\par\noindent confidence: high\par\noindent impact: none\par\medskip
\noindent\textbf{AK07-CN276}\enspace\hyperlink{AK07-J0271-L11}{AK07-J0271-L11}\enspace(misprint)\par\noindent printed: stat p. I.\enspace proposed: statt p. I.\par\noindent «statt» in every other entry of the list; one t is printed\par\noindent confidence: high\par\noindent impact: none\par\medskip
\noindent\textbf{AK07-CN277}\enspace\hyperlink{AK07-J0271-L13}{AK07-J0271-L13}\enspace(wrong fount)\par\noindent printed: Z. 2 \textbf{v}. u.\enspace proposed: Z. 2 v. u.\par\noindent the abbreviation «v. u.» is set in roman in every other entry, as «Z. 7 v. u.» in the line above; the v is printed in bold type\par\noindent confidence: high\par\noindent impact: none\par\medskip
\noindent\textbf{AK07-CN288}\enspace\hyperlink{AK07-J0276-L11}{AK07-J0276-L11}\enspace(misprint)\par\noindent printed: inbesondere\enspace proposed: insbesondere\par\noindent the adverb «insbesondere», as in the entry «Burug» two lines below (L14); the s is missing\par\noindent confidence: high\par\noindent impact: none\par\medskip
\noindent\textbf{AK07-CN289}\enspace\hyperlink{AK07-J0276-L88}{AK07-J0276-L88}\enspace(missing comma)\par\noindent printed: Praepos., in auf, z. B.\enspace proposed: Praepos., in, auf, z. B.\par\noindent the meanings of the preposition are listed with commas, as «Praepos., in, auf, mit, z. B.» in the entry «Bi» (L04); no comma is printed after «in»\par\noindent confidence: medium\par\noindent impact: none\par\medskip
\noindent\textbf{AK07-CN290}\enspace\hyperlink{AK07-J0278-L85}{AK07-J0278-L85}\enspace(missing accent)\par\noindent printed: \textit{maudūʿ}\enspace proposed: \textit{mauḍūʿ}\par\noindent the word is printed «mauḍūʿ», with the dot under the d (ḍād), in the next line of the same entry (L86); here the dot is missing\par\noindent confidence: high\par\noindent impact: none\par\medskip
\noindent\textbf{AK07-CN291}\enspace\hyperlink{AK07-J0279-L61}{AK07-J0279-L61}\enspace(missing point)\par\noindent printed: 2, 3, etc\enspace proposed: 2, 3, etc.\par\noindent the entries of the months end with «etc.», as «T. 2, 3, etc.» in the next entry (L63); the point after «etc» is not printed\par\noindent confidence: medium\par\noindent impact: none\par\medskip
\noindent\textbf{AK07-CN292}\enspace\hyperlink{AK07-J0279-L92}{AK07-J0279-L92}\enspace(missing point)\par\noindent printed: (pl \textit{ashrāf})\enspace proposed: (pl. \textit{ashrāf})\par\noindent the abbreviation «pl.» has its point everywhere else in the glossary, as «(pl. v. muʿāmala)» on this page (L08) and «(pl. ḥudūd)» on p. 243 (JP2 0277, L26)\par\noindent confidence: high\par\noindent impact: none\par\medskip
\noindent\textbf{AK07-CN293}\enspace\hyperlink{AK07-J0280-L55}{AK07-J0280-L55}\enspace(wrong punctuation)\par\noindent printed: 69. T. 4. 5, etc.\enspace proposed: 69, T. 4, 5, etc.\par\noindent the references of an entry are divided by commas, also before the tables, as «41, T. 4—20, etc.» (L75) and «58, T. 21—26» (L05); points are printed after 69 and after 4\par\noindent confidence: high\par\noindent impact: none\par\medskip
\noindent\textbf{AK07-CN294}\enspace\hyperlink{AK07-J0280-L70}{AK07-J0280-L70}\enspace(wrong reference)\par\noindent printed: 161\enspace proposed: 103\par\noindent «wabāl», the detriment of a planet, is explained on p. 103 (JP2 0135, L32: «im sog. Schaden (detrimentum, arab. wabāl)»), with the other astrological terms that the glossary refers to pp. 103-105 («Rabb», «Ṣāḥib» 103; «Hubūṭ», «Sharaf», «Wadjh» 104; «Ḥadd» 105); p. 161 is Tab. 50, a table of Venus, without the word\par\noindent confidence: medium\par\noindent impact: none\par\medskip
\noindent\textbf{AK07-CN295}\enspace\hyperlink{AK07-J0281-L37}{AK07-J0281-L37}\enspace(missing letter)\par\noindent printed: Ardībihist-māh.\enspace proposed: Ardībihisht-māh.\par\noindent the second Persian month is Ardībihisht (Middle Persian Ardwahišt), which ends in sh and t; the glossary writes this sound sh, as in «Shahrīwar-māh» (L41) and «Shabin = Shaʿbān» (p. 245, JP2 0279, L90)\par\noindent confidence: medium\par\noindent impact: none\par\medskip
\noindent\textbf{AK07-CN296}\enspace\hyperlink{AK07-J0283-L08}{AK07-J0283-L08}\enspace(missing accent)\par\noindent printed: s. Manṣur b. ʿAlī\enspace proposed: s. Manṣūr b. ʿAlī\par\noindent the name is printed «Manṣūr» earlier in the same line and «Abū Naṣr Manṣūr b. ʿAlī b. ʿIrāḳ» in the commentary, p. 49 (JP2 0081, L25); the macron of the u is missing\par\noindent confidence: high\par\noindent impact: none\par\medskip
\noindent\textbf{AK07-CN297}\enspace\hyperlink{AK07-J0283-L26}{AK07-J0283-L26}\enspace(wrong punctuation)\par\noindent printed: 106. T. 1.\enspace proposed: 106, T. 1.\par\noindent the references of an entry are divided by commas, also before the tables, as «38, T. 1.» in the entry «Alexander» (L33) and «104, T. 116.» (L32); a point is printed after 106\par\noindent confidence: high\par\noindent impact: none\par\medskip
\noindent\textbf{AK07-CN298}\enspace\hyperlink{AK07-J0283-L33}{AK07-J0283-L33}\enspace(wrong reference)\par\noindent printed: 5, 34, 35, 38, T. 1.\enspace proposed: 3, 34, 35, 38, T. 1.\par\noindent p. 5 of the text (JP2 0037) does not name Alexander; p. 3 (JP2 0035) names him twice, in the intervals between the eras («usque ad elkarnain, id est Alexandrum», «ab Alexandro vero usque ad Acofra»), which the commentary discusses on p. 35 (JP2 0067)\par\noindent confidence: medium\par\noindent impact: none\par\medskip
\noindent\textbf{AK07-CN299}\enspace\hyperlink{AK07-J0283-L46}{AK07-J0283-L46}\enspace(misprint)\par\noindent printed: desselbcn\enspace proposed: desselben\par\noindent «Beschreibung desselben», as «desselben» in the next line (L47) and twice in L53; the e is printed as a c, without its bar\par\noindent confidence: high\par\noindent impact: none\par\medskip
\noindent\textbf{AK07-CN300}\enspace\hyperlink{AK07-J0283-L50}{AK07-J0283-L50}\enspace(raised point)\par\noindent printed: T· 3\textsuperscript{a}.\enspace proposed: T. 3\textsuperscript{a}.\par\noindent the abbreviation «T.» has its point on the base line, as «T. 3, 3a.» two lines above (L48); here the point is printed at mid-height\par\noindent confidence: high\par\noindent impact: none\par\medskip
\noindent\textbf{AK07-CN309}\enspace\hyperlink{AK07-J0284-L06}{AK07-J0284-L06}\enspace(wrong sort)\par\noindent printed: XXIIl\enspace proposed: XXIII\par\noindent Athelhard is named in note 2 of p. XXIII (JP2 0027); the last sign of the numeral is a lowercase l, with a flag on the left only at its top, where the two I before it have serifs on both sides\par\noindent confidence: high\par\noindent impact: none\par\medskip
\noindent\textbf{AK07-CN301}\enspace\hyperlink{AK07-J0284-L22}{AK07-J0284-L22}\enspace(wrong punctuation)\par\noindent printed: 36. 41—44\enspace proposed: 36, 41—44\par\noindent the references of the entry are divided by commas; a point is printed after 36\par\noindent confidence: high\par\noindent impact: none\par\medskip
\noindent\textbf{AK07-CN302}\enspace\hyperlink{AK07-J0284-L39}{AK07-J0284-L39}\enspace(wrong sort)\par\noindent printed: Vl.\enspace proposed: VI.\par\noindent Billeter is thanked in the foreword, p. VI (JP2 0010, L02); the second sign of the numeral is a lowercase l, with the sloping serif of the l, where the I of the roman numerals has level serifs\par\noindent confidence: high\par\noindent impact: none\par\medskip
\noindent\textbf{AK07-CN303}\enspace\hyperlink{AK07-J0284-L51}{AK07-J0284-L51}\enspace(wrong letter)\par\noindent printed: S. 27—56\enspace proposed: T. 27—56\par\noindent the latitudes of the planets stand in a column of each of their tables, Tab. 27-56 (pp. 138-167); the index cites the tables with «T.» and has no other «S.» before a number\par\noindent confidence: high\par\noindent impact: none\par\medskip
\noindent\textbf{AK07-CN304}\enspace\hyperlink{AK07-J0284-L58}{AK07-J0284-L58}\enspace(missing accent)\par\noindent printed: Bukhtanaṣsar\enspace proposed: Bukhtanaṣṣar\par\noindent the Arabic name has the emphatic s doubled; the index marks both letters of a doubled emphatic consonant, as «Naḳḳāsh» (p. [249], JP2 0283, L07), and the glossary «wakhāṣṣathu» and «ḥāṣṣa»; here only the first s has its dot (the commentary, p. 35, JP2 0067, N01, writes «Bukhtanassar» without dots)\par\noindent confidence: high\par\noindent impact: none\par\medskip
\noindent\textbf{AK07-CN305}\enspace\hyperlink{AK07-J0284-L69}{AK07-J0284-L69}\enspace(wrong punctuation)\par\noindent printed: 58. s. auch\enspace proposed: 58, s. auch\par\noindent the entry continues with «s. auch Examinatio centri»; a comma stands before «s. auch» in «95, s. auch Tiefstand.» (L66)\par\noindent confidence: medium\par\noindent impact: none\par\medskip
\noindent\textbf{AK07-CN306}\enspace\hyperlink{AK07-J0284-L92}{AK07-J0284-L92}\enspace(wrong reference)\par\noindent printed: Deferens 16, 17,\enspace proposed: Deferens 47, 48,\par\noindent the deferent is introduced on pp. 47 and 48 (JP2 0079, L14-L35; JP2 0080, L11, L24), which the entry omits; pp. 16 and 17 (JP2 0048-0049, the Latin text) do not name it; 47 and 48 are the 16th and 17th pages of the commentary, which begins on p. 32\par\noindent confidence: medium\par\noindent impact: none\par\medskip
\noindent\textbf{AK07-CN307}\enspace\hyperlink{AK07-J0285-L82}{AK07-J0285-L82}\enspace(wrong sort)\par\noindent printed: XVlI\enspace proposed: XVII\par\noindent Gonzalez is cited in a note of p. XVII (JP2 0021); the third sign of the numeral is a lowercase l, with the sloping serif of the l, where the I beside it has level serifs, as in «Vl» for VI on p. 250 (JP2 0284, L39)\par\noindent confidence: high\par\noindent impact: none\par\medskip
\noindent\textbf{AK07-CN308}\enspace\hyperlink{AK07-J0285-L82}{AK07-J0285-L82}\enspace(missing point)\par\noindent printed: (N.)\enspace proposed: (N.).\par\noindent the entries of the index end with a point, also after a parenthesis, as «XXII (N. 1).» in the entry above (L81) and «XVIII (N.).» (L100)\par\noindent confidence: high\par\noindent impact: none\par\medskip
\noindent\textbf{AK07-CN310}\enspace\hyperlink{AK07-J0286-L27}{AK07-J0286-L27}\enspace(wrong sort)\par\noindent printed: XVIlI\enspace proposed: XVIII\par\noindent Isidore of Seville is named on p. XVIII (JP2 0022); the fourth sign of the numeral is a lowercase l, with a flag on the left only at its top, where the I on either side of it has serifs on both sides\par\noindent confidence: high\par\noindent impact: none\par\medskip
\noindent\textbf{AK07-CN311}\enspace\hyperlink{AK07-J0286-L35}{AK07-J0286-L35}\enspace(wrong sort)\par\noindent printed: Xl\enspace proposed: XI\par\noindent John of Seville is named on p. XI (JP2 0015); the second sign of the numeral is a lowercase l, with a flag on the left only at its top, where the I of «XI» in the next line (L36) has serifs on both sides\par\noindent confidence: high\par\noindent impact: none\par\medskip
\noindent\textbf{AK07-CN312}\enspace\hyperlink{AK07-J0286-L46}{AK07-J0286-L46}\enspace(wrong sort)\par\noindent printed: XXIl\enspace proposed: XXII\par\noindent the numeral is read XXII among the references XXI, XXII, XXIV of al-Khwārizmī; its last sign is a lowercase l, with a flag on the left only at its top\par\noindent confidence: high\par\noindent impact: none\par\medskip
\noindent\textbf{AK07-CN313}\enspace\hyperlink{AK07-J0286-L45}{AK07-J0286-L45}\enspace(missing accent)\par\noindent printed: Muh.\enspace proposed: Muḥ.\par\noindent the abbreviation of Muḥammad has the dot under the h in the index, as «al-Battānī, Muḥ. b. Djābir» and «al-Bīrūnī, Muḥ. b. Aḥmed» (p. 250, JP2 0284, L13 and L40); here the dot is missing\par\noindent confidence: high\par\noindent impact: none\par\medskip
\noindent\textbf{AK07-CN314}\enspace\hyperlink{AK07-J0286-L86}{AK07-J0286-L86}\enspace(wrong sort)\par\noindent printed: XIIl\enspace proposed: XIII\par\noindent the numeral is read XIII between «VII—X» and «XV» of the references to Maslama; its last sign is a lowercase l, with a flag on the left only at its top\par\noindent confidence: high\par\noindent impact: none\par\medskip
\noindent\textbf{AK07-CN315}\enspace\hyperlink{AK07-J0286-L87}{AK07-J0286-L87}\enspace(wrong sort)\par\noindent printed: XXl\enspace proposed: XXI\par\noindent the numeral is read XXI between XV and XXII of the references to Maslama; its last sign is a lowercase l, with a flag on the left only at its top\par\noindent confidence: high\par\noindent impact: none\par\medskip
\noindent\textbf{AK07-CN316}\enspace\hyperlink{AK07-J0287-L43}{AK07-J0287-L43}\enspace(wrong punctuation)\par\noindent printed: 89. 95.\enspace proposed: 89, 95.\par\noindent the references of an entry are divided by commas; Oppolzer is cited on pp. 89 and 95 (JP2 0121, 0127); a point is printed after 89\par\noindent confidence: high\par\noindent impact: none\par\medskip
\noindent\textbf{AK07-CN317}\enspace\hyperlink{AK07-J0287-L59}{AK07-J0287-L59}\enspace(wrong note number)\par\noindent printed: XXIII (N. 2)\enspace proposed: XXIII (N. 1)\par\noindent «Petri Philomeni de Dacia in Algorismum vulgarem Johannis de Sacrobosco Commentarius» is cited in note 1 of p. XXIII (JP2 0027, N01), as the index has «Johannes de Sacrobosco XXIII (N. 1)» (p. 252, JP2 0286, L34); note 2 cites Ocreatus\par\noindent confidence: high\par\noindent impact: none\par\medskip
\noindent\textbf{AK07-CN318}\enspace\hyperlink{AK07-J0287-L72}{AK07-J0287-L72}\enspace(wrong sort)\par\noindent printed: XIIl\enspace proposed: XIII\par\noindent the references of Ptolemy run IX, XIII, XXIV; the last sign of the numeral is a lowercase l, with a flag on the left only at its top, where the I before it has serifs on both sides\par\noindent confidence: high\par\noindent impact: none\par\medskip
\noindent\textbf{AK07-CN319}\enspace\hyperlink{AK07-J0287-L74}{AK07-J0287-L74}\enspace(wrong punctuation)\par\noindent printed: 90. 96,\enspace proposed: 90, 96,\par\noindent the references of the entry are divided by commas; a point is printed after 90\par\noindent confidence: high\par\noindent impact: none\par\medskip
\noindent\textbf{AK07-CN320}\enspace\hyperlink{AK07-J0288-L03}{AK07-J0288-L03}\enspace(missing point)\par\noindent printed: Cestrensis s Robert\enspace proposed: Cestrensis s. Robert\par\noindent the abbreviation «s.» (siehe) has its point everywhere else in the index, as «Retinensis s. Robert von Chester» two lines above (L01); here no point is printed after the s\par\noindent confidence: high\par\noindent impact: none\par\medskip
\noindent\textbf{AK07-CN321}\enspace\hyperlink{AK07-J0288-L52}{AK07-J0288-L52}\enspace(missing comma)\par\noindent printed: 86 87,\enspace proposed: 86, 87,\par\noindent the references of an entry are divided by commas, as «86, 87, 94, 95» in the next entry (L54)\par\noindent confidence: high\par\noindent impact: none\par\medskip
\noindent\textbf{AK07-CN322}\enspace\hyperlink{AK07-J0288-L78}{AK07-J0288-L78}\enspace(wrong accent)\par\noindent printed: Shâh\enspace proposed: Shāh\par\noindent the name is printed «Shāh», with the macron, in the commentary (p. 32, JP2 0064; p. 52, JP2 0084) and in the glossary (p. 246, JP2 0280, L94); the index alone has the circumflex\par\noindent confidence: medium\par\noindent impact: none\par\medskip
\end{document}
